% !TEX program = xelatex
% =====================================================================
%  基础随机分析  ——  主文件
%  编译: xelatex -> xelatex (或 latexmk -xelatex stochastic_analysis.tex)
%  只编译某一章: 取消下方 \includeonly 的注释
% =====================================================================
\documentclass[UTF8,a4paper,oneside,zihao=-4,linespread=1.45,fontset=none]{ctexbook}

% \includeonly{chapters/chapter1}

% ---------------------------------------------------------------------
%  版面
% ---------------------------------------------------------------------
\usepackage[a4paper,top=2.6cm,bottom=2.6cm,left=2.7cm,right=2.7cm,headheight=15pt]{geometry}

% ---------------------------------------------------------------------
%  数学宏包 (须在 unicode-math 之前载入)
% ---------------------------------------------------------------------
\usepackage{amsmath,amsthm,mathtools}

% ---------------------------------------------------------------------
%  字体
%  西文与数学: TeX Gyre Pagella / Pagella Gyre Math
%  中文: 宋体; 加粗 -> 黑体; 斜体 -> 楷体
%  依次尝试 macOS / Windows 字体, 找不到时退回 TeX Live 自带的 Fandol 字体
% ---------------------------------------------------------------------
\usepackage{unicode-math}
% 按文件名载入 (TeX Live / MacTeX 自带), 不依赖系统字体注册
\setmainfont{texgyrepagella}[Extension=.otf, UprightFont=*-regular,
  BoldFont=*-bold, ItalicFont=*-italic, BoldItalicFont=*-bolditalic]
\setsansfont{texgyreheros}[Extension=.otf, UprightFont=*-regular,
  BoldFont=*-bold, ItalicFont=*-italic, BoldItalicFont=*-bolditalic]
\setmonofont{texgyrecursor}[Extension=.otf, UprightFont=*-regular,
  BoldFont=*-bold, ItalicFont=*-italic, BoldItalicFont=*-bolditalic]
\setmathfont{texgyrepagella-math.otf}

\newcommand\SAsong{FandolSong-Regular.otf}
\newcommand\SAhei{FandolHei-Regular.otf}
\newcommand\SAkai{FandolKai-Regular.otf}
\IfFontExistsTF{SimSun}{\renewcommand\SAsong{SimSun}}{}
\IfFontExistsTF{Songti SC}{\renewcommand\SAsong{Songti SC}}{}
\IfFontExistsTF{SimHei}{\renewcommand\SAhei{SimHei}}{}
\IfFontExistsTF{STHeiti}{\renewcommand\SAhei{STHeiti}}{}
\IfFontExistsTF{Heiti SC}{\renewcommand\SAhei{Heiti SC}}{}
\IfFontExistsTF{KaiTi}{\renewcommand\SAkai{KaiTi}}{}
\IfFontExistsTF{STKaiti}{\renewcommand\SAkai{STKaiti}}{}
\IfFontExistsTF{Kaiti SC}{\renewcommand\SAkai{Kaiti SC}}{}

\setCJKmainfont{\SAsong}[BoldFont=\SAhei, ItalicFont=\SAkai, BoldItalicFont=\SAhei]
\setCJKsansfont{\SAhei}
\setCJKmonofont{\SAsong}
\setCJKfamilyfont{zhsong}{\SAsong}
\setCJKfamilyfont{zhhei}{\SAhei}
\setCJKfamilyfont{zhkai}{\SAkai}
\providecommand\songti{\CJKfamily{zhsong}}
\providecommand\heiti{\CJKfamily{zhhei}}
\providecommand\kaishu{\CJKfamily{zhkai}}

% ---------------------------------------------------------------------
%  颜色: #ef476f #ffd166 #06d6a0 #118ab2 #073b4c
% ---------------------------------------------------------------------
\usepackage{xcolor}
\definecolor{SARed}{HTML}{EF476F}
\definecolor{SAYellow}{HTML}{FFD166}
\definecolor{SAGreen}{HTML}{06D6A0}
\definecolor{SABlue}{HTML}{118AB2}
\definecolor{SADark}{HTML}{073B4C}
\colorlet{SAGreenText}{SAGreen!60!SADark}   % 绿色作文字时加深以保证可读
\colorlet{SAYellowText}{SAYellow!55!SADark}

% ---------------------------------------------------------------------
%  绘图
% ---------------------------------------------------------------------
\usepackage{tikz}
\usetikzlibrary{arrows.meta,calc,positioning,decorations.pathreplacing,
  patterns,shapes.geometric,intersections,fit,backgrounds}
\usepackage{pgfplots}
\pgfplotsset{compat=1.17}
\tikzset{
  >={Stealth[length=2.2mm]},
  every picture/.style={line cap=round,line join=round},
}
\pgfplotsset{
  SAaxis/.style={axis lines=middle, axis line style={->,SADark!80},
    tick style={SADark!60}, every tick label/.append style={font=\small},
    label style={font=\small}},
}

% ---------------------------------------------------------------------
%  其他常用宏包
% ---------------------------------------------------------------------
\usepackage{booktabs,array,multirow}
\usepackage[shortlabels]{enumitem}
\setlist{itemsep=2pt,topsep=4pt,parsep=0pt}
\setlist[enumerate,1]{label=(\arabic*)}
\usepackage{fancyhdr}
\usepackage{needspace}  % 避免节标题落在页底
\usepackage{caption}
\captionsetup{font=small,labelfont={bf,color=SADark},labelsep=space}

% ---------------------------------------------------------------------
%  章节样式
% ---------------------------------------------------------------------
\ctexset{
  chapter = {
    name        = {第,章},
    number      = \arabic{chapter},
    format      = \raggedright,
    nameformat  = \large\bfseries\color{SARed},
    titleformat = \huge\bfseries\color{SADark},
    aftername   = \par\vspace{0.6em},
    beforeskip  = 10pt,
    afterskip   = 28pt,
    aftertitle  = {\par\vspace{0.5em}{\color{SABlue}\rule{\linewidth}{0.8pt}}},
  },
  section = {
    format     = \Large\bfseries\color{SADark}\raggedright,
    numberformat = \color{SABlue},
    beforeskip = 22pt plus 4pt minus 2pt,
    afterskip  = 12pt,
    break      = {\Needspace{7\baselineskip}},
  },
  subsection = {
    format     = \large\bfseries\color{SADark}\raggedright,
    numberformat = \color{SABlue},
    beforeskip = 16pt plus 3pt minus 2pt,
    afterskip  = 8pt,
    break      = {\Needspace{5\baselineskip}},
  },
  subsubsection = {
    format = \normalsize\bfseries\color{SADark}\raggedright,
    beforeskip = 12pt, afterskip = 6pt,
  },
  paragraph = {format = \bfseries\color{SADark}, afterskip = 1em},
}
\setcounter{secnumdepth}{3}
\setcounter{tocdepth}{2}

\pagestyle{fancy}
\fancyhf{}
\fancyhead[L]{\small\color{SADark!70}\leftmark}
\fancyhead[R]{\small\color{SADark!70}\thepage}
\renewcommand{\headrulewidth}{0.4pt}
\renewcommand{\headrule}{{\color{SABlue!50}\hrule width\headwidth height\headrulewidth}}
\fancypagestyle{plain}{\fancyhf{}\fancyfoot[C]{\small\color{SADark!70}\thepage}\renewcommand{\headrulewidth}{0pt}}
\renewcommand{\chaptermark}[1]{\markboth{第\,\thechapter\,章\quad #1}{}}

% ---------------------------------------------------------------------
%  定理类环境 (amsthm 编号 + tcolorbox 轻量外框)
%  所有编号环境共用一个计数器, 按章编号: 定义 1.1, 定理 1.2, 例 1.3, ...
% ---------------------------------------------------------------------
\usepackage[most]{tcolorbox}

\newtheoremstyle{SAstyle}%
  {0pt}{0pt}%            上下间距 (由外框控制)
  {\normalfont}{}%       正文字体, 缩进
  {\bfseries}{.}%        标题字体, 标题后标点
  {0.6em}%               标题后空白
  {\color{SAheadcolor}\thmname{#1}\thmnumber{\ #2}\thmnote{\ (#3)}}
\colorlet{SAheadcolor}{SADark}
\theoremstyle{SAstyle}

\newtheorem{theorem}{定理}[chapter]
\newtheorem{definition}[theorem]{定义}
\newtheorem{lemma}[theorem]{引理}
\newtheorem{proposition}[theorem]{命题}
\newtheorem{corollary}[theorem]{推论}
\newtheorem{example}[theorem]{例}
\newtheorem*{remark}{注}
\newtheorem*{intuition}{直观}

\tcbset{
  SAbox/.style={
    enhanced, breakable, sharp corners, frame hidden,
    colback=#1!5!white,
    borderline west={2.2pt}{0pt}{#1},
    left=9pt, right=8pt, top=5pt, bottom=5pt,
    before skip=10pt plus 2pt, after skip=10pt plus 2pt,
  },
  SAplain/.style={
    enhanced, breakable, sharp corners, frame hidden, colback=white,
    borderline west={1.4pt}{0pt}{#1},
    left=9pt, right=0pt, top=2pt, bottom=2pt,
    before skip=8pt plus 2pt, after skip=8pt plus 2pt,
  },
}
\tcolorboxenvironment{definition} {SAbox=SABlue,  before upper={\colorlet{SAheadcolor}{SABlue}}}
\tcolorboxenvironment{theorem}    {SAbox=SARed,   before upper={\colorlet{SAheadcolor}{SARed}}}
\tcolorboxenvironment{lemma}      {SAbox=SADark,  before upper={\colorlet{SAheadcolor}{SADark}}}
\tcolorboxenvironment{proposition}{SAbox=SADark,  before upper={\colorlet{SAheadcolor}{SADark}}}
\tcolorboxenvironment{corollary}  {SAbox=SADark,  before upper={\colorlet{SAheadcolor}{SADark}}}
\tcolorboxenvironment{example}    {SAplain=SAGreen, before upper={\colorlet{SAheadcolor}{SAGreenText}}}
\tcolorboxenvironment{remark}     {SAplain=SAYellow, before upper={\colorlet{SAheadcolor}{SAYellowText}}}
\tcolorboxenvironment{intuition}  {SAbox=SAYellow, before upper={\colorlet{SAheadcolor}{SAYellowText}}}

% 证明环境: 黑体标题, 不加框
\makeatletter
\renewenvironment{proof}[1][证明]{\par
  \pushQED{\qed}\normalfont\topsep6pt plus 2pt\relax
  \trivlist\item[\hskip\labelsep{\bfseries\color{SADark}#1.}]\ignorespaces
}{\popQED\endtrivlist\@endpefalse}
\makeatother
\renewcommand{\qedsymbol}{\tikz[baseline=0pt]\fill[SABlue] (0,0) rectangle (0.55em,0.55em);}

% ---------------------------------------------------------------------
%  全局行内公式 displaystyle, 并加大行间最小间隙避免高公式互相挤压
% ---------------------------------------------------------------------
\AtBeginDocument{\everymath{\displaystyle}}
\setlength{\emergencystretch}{3em}   % 避免长行内公式导致溢出
\setlength{\lineskiplimit}{3pt}
\setlength{\lineskip}{5pt}
\setlength{\jot}{6pt}
\setlength{\abovedisplayskip}{8pt plus 2pt minus 2pt}
\setlength{\belowdisplayskip}{8pt plus 2pt minus 2pt}

% ---------------------------------------------------------------------
%  符号约定 (详见 "符号约定" 一节)
%  概率与期望: 空心; σ-代数、集族、信息流: \mathcal;
%  随机向量: 加粗斜体 \rvec{X}; 矩阵: 加粗正体 \mat{A}
% ---------------------------------------------------------------------
\renewcommand{\P}{\mathbb{P}}
\newcommand{\E}{\mathbb{E}}
\newcommand{\R}{\mathbb{R}}
\newcommand{\N}{\mathbb{N}}
\newcommand{\Z}{\mathbb{Z}}
\newcommand{\Q}{\mathbb{Q}}
\newcommand{\1}{\mathbb{1}}                 % 示性函数 \1_A
\newcommand{\rvec}[1]{\symbfit{#1}}         % 随机向量
\newcommand{\mat}[1]{\symbf{#1}}            % 矩阵
\newcommand{\dd}{\mathop{}\!\mathrm{d}}     % 微分号
\newcommand{\T}{\top}                       % 转置
\DeclareMathOperator{\Var}{Var}
\DeclareMathOperator{\Cov}{Cov}
\DeclareMathOperator{\sgn}{sgn}
\DeclarePairedDelimiter{\abs}{\lvert}{\rvert}
\DeclarePairedDelimiter{\norm}{\lVert}{\rVert}
\DeclarePairedDelimiter{\ip}{\langle}{\rangle}
\newcommand{\iid}{\text{i.i.d.}}
\newcommand{\as}{\text{a.s.}}

% ---------------------------------------------------------------------
%  超链接 (最后载入)
% ---------------------------------------------------------------------
\usepackage[colorlinks=true,linkcolor=SABlue,citecolor=SAGreenText,urlcolor=SABlue,
  bookmarksnumbered=true]{hyperref}

% =====================================================================
\title{\color{SADark}\bfseries 基础随机分析}
\author{}
\date{}

\begin{document}

\frontmatter
\begin{titlepage}
  \centering
  \vspace*{0.28\textheight}
  {\color{SABlue}\rule{0.75\textwidth}{1pt}}\par\vspace{1.2em}
  {\Huge\bfseries\color{SADark} 基础随机分析}\par\vspace{0.8em}
  {\large\color{SADark!80} 从测度论概率到随机微分方程}\par\vspace{1.0em}
  {\color{SABlue}\rule{0.75\textwidth}{1pt}}\par
  \vfill
  {\large\color{SADark!80}\today}\par
  \vspace*{0.08\textheight}
\end{titlepage}

\tableofcontents
% =====================================================================
%  符号约定
% =====================================================================
\chapter*{符号约定}
\addcontentsline{toc}{chapter}{符号约定}
\markboth{符号约定}{}

全书统一使用下列记号. 标点一律使用英文标点 (顿号除外). 字体约定: 概率与期望用空心字母; $\sigma$-代数、集族与信息流用花体 ($\mathcal{F}$); 随机积分中被积过程构成的空间用书写体 ($\mathscr{L}$); 随机向量用加粗斜体, 矩阵用加粗正体.

\subsection*{基本记号 (第~1~章起)}

\begin{center}
\small\renewcommand{\arraystretch}{1.1}
\begin{tabular}{>{\raggedright\arraybackslash}p{0.36\textwidth} >{\raggedright\arraybackslash}p{0.52\textwidth}}
\toprule
\textbf{记号} & \textbf{含义} \\
\midrule
$\Omega$, $\omega$ & 样本空间、样本点 \\
$\mathcal{F}$, $\mathcal{G}$, $\mathcal{H}$ & $\sigma$-代数 (一般用花体大写字母表示集族、$\sigma$-代数) \\
$\P$, $\E$ & 概率测度、数学期望 (空心字母) \\
$X, Y, Z$ & 随机变量 (大写斜体字母) \\
$\rvec{X} = (X_1,\dots,X_n)^\T$ & 随机向量 (加粗斜体) \\
$\mat{A}$, $\mat{\Sigma}$ & 矩阵 (加粗正体), $\mat{A}^\T$ 为转置 \\
$\1_A$ & 事件 $A$ 的示性函数 \\
$\mathcal{B}(\R)$, $\mathcal{B}(\R^n)$ & Borel $\sigma$-代数 \\
$\sigma(X)$, $\sigma(\mathcal{C})$ & 由随机变量 $X$、集族 $\mathcal{C}$ 生成的 $\sigma$-代数 \\
$\E[X;A]$ & $\E[X\1_A]$ \\
$\Var(X)$, $\Cov(X,Y)$ & 方差、协方差 \\
\bottomrule
\end{tabular}
\end{center}

\subsection*{分布、收敛与其他 (第~1~章起)}

\begin{center}
\small\renewcommand{\arraystretch}{1.1}
\begin{tabular}{>{\raggedright\arraybackslash}p{0.36\textwidth} >{\raggedright\arraybackslash}p{0.52\textwidth}}
\toprule
\textbf{记号} & \textbf{含义} \\
\midrule
$N(\mu,\sigma^2)$, $N(\rvec{\mu},\mat{\Sigma})$ & 正态分布、多维正态分布 \\
$\varphi$, $\Phi$ & 标准正态分布的密度、分布函数 \\
$X\overset{d}{=}Y$ & $X$ 与 $Y$ 同分布 \\
$\xrightarrow{\as}$, $\xrightarrow{\P}$, $\xrightarrow{L^p}$, $\xrightarrow{d}$ & 几乎必然收敛、依概率收敛、$L^p$ 收敛、依分布收敛 \\
$L^p(\Omega,\mathcal{F},\P)$ & $p$ 次可积随机变量全体 (按几乎处处相等等同) \\
$\ip{X,Y}$, $\norm{X}_p$ & $L^2$ 内积、$L^p$ 范数 \\
$a\wedge b$, $a\vee b$ & $\min\{a,b\}$、$\max\{a,b\}$ \\
$x^+$, $x^-$ & $\max\{x,0\}$、$\max\{-x,0\}$ \\
$\sgn(x)$ & 符号函数: $x>0$ 时为 $1$, $x\le0$ 时为 $-1$ \\
$\as$ & 几乎必然 (almost surely) \\
$\iid$ & 独立同分布 \\
\bottomrule
\end{tabular}
\end{center}

\subsection*{条件期望、随机过程与鞅 (第~2,~3~章)}

\begin{center}
\small\renewcommand{\arraystretch}{1.1}
\begin{tabular}{>{\raggedright\arraybackslash}p{0.36\textwidth} >{\raggedright\arraybackslash}p{0.52\textwidth}}
\toprule
\textbf{记号} & \textbf{含义} \\
\midrule
$\E[X\,|\,\mathcal{G}]$, $\E[X\,|\,Y]$ & 条件期望 \\
$\P(A\,|\,\mathcal{G})$ & 条件概率 \\
$\{\mathcal{F}_t\}_{t\ge 0}$ & 信息流 ($\sigma$-代数流) \\
$\mathcal{F}^X_t=\sigma(X_s:s\le t)$ & 过程 $X$ 的自然信息流 \\
$\mathcal{F}_{t+}=\bigcap_{s>t}\mathcal{F}_s$ & 右极限 $\sigma$-代数 \\
$P_{s,t}$, $P_t$, $P_tf$ & 转移概率 (核)、时齐转移核、$\int f(y)P_t(x,\dd y)$ \\
$\tau$, $\sigma$ (在停时语境中) & 停时 \\
$\mathcal{F}_\tau$ & 停时 $\tau$ 之前的信息 \\
$X^\tau_t=X_{t\wedge\tau}$ & 停止过程 \\
$(H\cdot X)$ & 鞅变换 (离散) 或随机积分 (连续) \\
$\ip{M}$, $\ip{M,N}$ & (可料) 二次变差、交互变差 (与 $L^2$ 内积 $\ip{X,Y}$ 由上下文区分) \\
$\Delta X_n=X_n-X_{n-1}$ & 增量 \\
\bottomrule
\end{tabular}
\end{center}

\subsection*{布朗运动与随机分析 (第~4--6~章)}

\begin{center}
\small\renewcommand{\arraystretch}{1.1}
\begin{tabular}{>{\raggedright\arraybackslash}p{0.36\textwidth} >{\raggedright\arraybackslash}p{0.52\textwidth}}
\toprule
\textbf{记号} & \textbf{含义} \\
\midrule
$B$, $\rvec{B}$ & (一维、$d$ 维) 标准布朗运动 \\
$M_t=\max_{s\le t}B_s$ & 布朗运动的最大值过程 \\
$\tau_a=\inf\{t:B_t=a\}$ & 首达时 \\
$C[0,\infty)$, $\mathcal{C}$, $\mathbb{W}$ & 连续函数空间、坐标 $\sigma$-代数、Wiener 测度 \\
$\P_x$, $\E_x$ & 从 $x$ 出发的布朗运动 (或扩散过程) 的概率与期望 \\
$p_t(x,y)$ & 热核 (布朗运动的转移密度) \\
$\mathscr{S}_T$ & 简单过程全体 \\
$\mathscr{L}^2_T$, $\norm{H}_T$ & 满足 $\E\int_0^TH_s^2\dd s<\infty$ 的循序可测过程、相应的范数 \\
$\mathscr{L}^2_{\mathrm{loc}}$ & 满足 $\int_0^TH_s^2\dd s<\infty$ $\as$ 的循序可测过程 \\
$\int_0^tH_s\dd B_s=(H\cdot B)_t$ & It\^o 积分 \\
$\int_0^tY_s\circ\dd X_s$ & Stratonovich 积分 \\
$\mathscr{E}(\theta\cdot B)$ & 随机指数 $\exp\big(\int\theta\dd B-\frac12\int\theta^2\dd s\big)$ \\
$\widetilde{\P}$, $\widetilde{\E}$ & 变换后的 (例如风险中性) 概率测度与期望 \\
$L_t$, $L^a_t$ & 布朗运动在 $0$、在 $a$ 处的局部时 \\
$C^{1,2}$, $C^2_c$ & 关于 $t$ 一阶、关于 $x$ 二阶连续可微的函数; 二阶连续可微且支撑紧的函数 \\
$\partial_t$, $\partial_x$, $\partial_{xx}$, $\partial_i$, $\Delta$ & 偏导数、Laplace 算子 \\
$\mathscr{C}^2_T$ & 满足 $\E\sup_{t\le T}\abs{X_t}^2<\infty$ 的连续适应过程 \\
$X^{s,x}$, $X^x$ & 从时刻 $s$ (或 $0$)、位置 $x$ 出发的随机微分方程的解 \\
$\mathcal{A}$, $L$, $L^*$ & 无穷小生成元、微分算子 $b\partial_x+\frac12\sigma^2\partial_{xx}$ 及其伴随 \\
$s(x)$ & 一维扩散的尺度函数 \\
\bottomrule
\end{tabular}
\end{center}


\mainmatter
% =====================================================================
%  第 1 章  现代概率论基础
% =====================================================================
\chapter{现代概率论基础}
\label{ch:1}

随机分析研究的对象是随时间演化的随机现象, 例如股票价格、花粉颗粒的运动以及信号中的噪声. 要对这些对象做严格的微积分运算, 首先需要一个足够牢固的概率论框架: 它既要能够处理 ``无穷多次'' 试验, 又要允许我们对随机量取极限、交换极限与期望. 这个框架就是 Kolmogorov 在 1933 年建立的测度论概率论.

本章依次回答以下问题: 什么是事件与概率 (1.1 节)? 什么是随机变量和它的分布 (1.2 节)? 如何定义并计算期望, 何时可以交换极限与期望 (1.3 节)? 随机变量序列有哪些收敛方式 (1.4 节)? 多个随机变量如何联合描述, 特别是多元正态分布 (1.5 节)? 最后, 平方可积随机变量构成一个 Hilbert 空间, 这一几何观点 (1.6 节) 将直接引出第~2~章的条件期望.

% =====================================================================
\section{概率的公理化结构}
\label{sec:1-1}

设想把一枚均匀硬币无限次地抛掷下去, 并问: ``正面出现的频率最终趋于 $\frac{1}{2}$'' 这件事发生的概率是多少? 为了让这个问题有意义, 需要先说清楚三件事.

\begin{enumerate}
\item \textbf{所有可能的结果.} 用 $1$ 表示正面, $0$ 表示反面, 一次完整的试验结果是一个无穷 $0$-$1$ 序列 $\omega=(\omega_1,\omega_2,\dots)$, 全体结果构成集合 $\Omega=\{0,1\}^{\N}$. 这是一个不可数集.
\item \textbf{哪些子集算作 ``事件''.} 记 $S_n(\omega)=\omega_1+\cdots+\omega_n$ 为前 $n$ 次中正面的次数. 我们关心的事件是
\[
A=\Big\{\omega:\ \lim_{n\to\infty}\frac{S_n(\omega)}{n}=\frac{1}{2}\Big\}.
\]
把极限的定义逐层写开 (对任意 $m$, 存在 $N$, 使得对所有 $n\ge N$ 都有误差小于 $\frac{1}{m}$), 得到
\begin{equation}\label{eq:1-freq-event}
A=\bigcap_{m=1}^{\infty}\ \bigcup_{N=1}^{\infty}\ \bigcap_{n=N}^{\infty}\Big\{\omega:\ \Big|\frac{S_n(\omega)}{n}-\frac{1}{2}\Big|<\frac{1}{m}\Big\}.
\end{equation}
最里层的集合只与前 $n$ 次抛掷有关, 它的概率可以用排列组合算出; 但 $A$ 本身是由它们经过\emph{可列次}交与并得到的. 事件的全体因此必须对可列次集合运算封闭.
\item \textbf{如何给事件赋予概率.} 只涉及前 $n$ 次抛掷的事件, 可以按 ``每种结果的概率为 $2^{-n}$'' 计算; 对 $A$ 这样的事件, 则需要概率能够 ``穿过'' 可列次集合运算. 这就是可列可加性.
\end{enumerate}

这三点分别对应样本空间 $\Omega$、$\sigma$-代数 $\mathcal{F}$ 和概率测度 $\P$. 本节先介绍这三个对象, 再介绍处理 $\sigma$-代数的基本工具 $\pi$-$\lambda$ 定理, 最后讨论独立性, 以及刻画 ``无穷多次发生'' 的 Borel-Cantelli 引理与 Kolmogorov 0-1 律.

% ---------------------------------------------------------------------
\subsection{样本空间与 \texorpdfstring{$\sigma$}{sigma}-代数}

\begin{definition}[$\sigma$-代数与可测空间]\label{def:sigma-algebra}
设 $\Omega$ 是非空集合, 其元素 $\omega$ 称为\textbf{样本点}, $\Omega$ 称为\textbf{样本空间}. 记 $2^\Omega$ 为 $\Omega$ 的全体子集. 子集族 $\mathcal{F}\subseteq 2^\Omega$ 称为 $\Omega$ 上的 \textbf{$\sigma$-代数} (或 $\sigma$-域), 如果
\begin{enumerate}
\item $\Omega\in\mathcal{F}$;
\item 对补运算封闭: 若 $A\in\mathcal{F}$, 则 $A^c=\Omega\setminus A\in\mathcal{F}$;
\item 对可列并封闭: 若 $A_1,A_2,\dots\in\mathcal{F}$, 则 $\bigcup_{n=1}^\infty A_n\in\mathcal{F}$.
\end{enumerate}
$\mathcal{F}$ 中的元素称为\textbf{事件} (或可测集), 二元组 $(\Omega,\mathcal{F})$ 称为\textbf{可测空间}.
\end{definition}

\begin{proposition}\label{prop:sigma-closure}
设 $\mathcal{F}$ 是 $\sigma$-代数, 则
\begin{enumerate}
\item $\varnothing\in\mathcal{F}$;
\item $\mathcal{F}$ 对可列交封闭: 若 $A_1,A_2,\dots\in\mathcal{F}$, 则 $\bigcap_{n=1}^\infty A_n\in\mathcal{F}$;
\item $\mathcal{F}$ 对有限并、有限交以及差运算封闭.
\end{enumerate}
\end{proposition}

\begin{proof}
(1) 由定义中的 (1)(2), $\varnothing=\Omega^c\in\mathcal{F}$.

(2) 由补运算封闭, $A_n^c\in\mathcal{F}$; 由可列并封闭, $\bigcup_n A_n^c\in\mathcal{F}$; 再由 De Morgan 公式与补运算封闭,
\[
\bigcap_{n=1}^\infty A_n=\Big(\bigcup_{n=1}^\infty A_n^c\Big)^c\in\mathcal{F}.
\]

(3) 对 $A_1,\dots,A_k\in\mathcal{F}$, 令 $A_n=\varnothing$ ($n>k$), 则 $\bigcup_{n=1}^k A_n=\bigcup_{n=1}^\infty A_n\in\mathcal{F}$; 令 $A_n=\Omega$ ($n>k$) 并利用 (2) 可得有限交封闭. 最后 $A\setminus B=A\cap B^c\in\mathcal{F}$.
\end{proof}

\begin{intuition}
$\sigma$-代数是 ``可以谈论其概率的事件'' 的全体. 三条公理保证: 能谈论 $A$, 就能谈论 ``$A$ 不发生''; 能谈论每个 $A_n$, 就能谈论 ``至少一个 $A_n$ 发生'' 和 ``所有 $A_n$ 都发生''. 在随机过程理论中, $\sigma$-代数还有第二层含义: 它表示\emph{信息}. 一个事件属于 $\mathcal{F}$, 意味着掌握信息 $\mathcal{F}$ 的观察者能判断该事件是否发生.
\end{intuition}

\begin{example}\label{ex:sigma-examples}
\begin{enumerate}
\item \textbf{平凡 $\sigma$-代数} $\{\varnothing,\Omega\}$: 最小的 $\sigma$-代数, 代表 ``毫无信息''.
\item \textbf{幂集} $2^\Omega$: 最大的 $\sigma$-代数, 代表 ``完全信息''.
\item 对给定的 $A\subseteq\Omega$, $\{\varnothing,A,A^c,\Omega\}$ 是 $\sigma$-代数, 代表 ``只知道 $A$ 是否发生''.
\item \textbf{由有限分割生成的 $\sigma$-代数.} 设 $\Omega=D_1\cup\cdots\cup D_k$, 其中 $D_i$ 非空且两两不交. 令 $\mathcal{F}$ 为所有形如 $\bigcup_{i\in I}D_i$ ($I\subseteq\{1,\dots,k\}$, 约定 $I=\varnothing$ 时为 $\varnothing$) 的集合全体. 补集对应 $I$ 的补集, 并集对应指标集的并集, 所以 $\mathcal{F}$ 是 $\sigma$-代数, 共有 $2^k$ 个元素. 它代表的信息是: 观察者知道 $\omega$ 落在哪一块 $D_i$ 中, 但不知道块内的具体位置 (图~\ref{fig:partition}).
\end{enumerate}
\end{example}

\begin{figure}[htbp]
\centering
\begin{tikzpicture}[scale=1]
  \fill[SABlue!18] (0,0) rectangle (2.4,1.6);
  \fill[SABlue!18] (2.4,1.6) -- (5,1.6) -- (5,3) -- (3.1,3) -- cycle;
  \draw[SADark,thick] (0,0) rectangle (5,3);
  \draw[SADark] (0,1.6) -- (5,1.6);
  \draw[SADark] (2.4,0) -- (2.4,1.6);
  \draw[SADark] (2.4,1.6) -- (3.1,3);
  \node at (1.2,0.8) {$D_3$};
  \node at (3.7,0.8) {$D_4$};
  \node at (1.3,2.3) {$D_1$};
  \node at (4.1,2.3) {$D_2$};
  \node[SADark] at (5.35,2.8) {$\Omega$};
  \fill[SARed] (0.55,0.45) circle (1.5pt) node[right,font=\small] {$\omega$};
  \node[align=left,font=\small,anchor=west] at (5.8,1.5)
    {阴影部分 $D_2\cup D_3\in\mathcal{F}$.\\[2pt]
     观察者只知道 $\omega\in D_3$,\\[2pt]
     不知道 $\omega$ 在 $D_3$ 中的位置.};
\end{tikzpicture}
\caption{由分割 $\{D_1,D_2,D_3,D_4\}$ 生成的 $\sigma$-代数有 $2^4=16$ 个元素}
\label{fig:partition}
\end{figure}

既然幂集 $2^\Omega$ 是最大的 $\sigma$-代数, 为什么不总取 $\mathcal{F}=2^\Omega$? 下面的例子说明: 当 $\Omega$ 不可数时, 幂集太大, 以至于连 ``$[0,1)$ 上的均匀分布'' 都无法定义在全体子集上.

\begin{example}[Vitali 集]\label{ex:vitali}
对 $x\in[0,1)$ 与 $q\in[0,1)$, 记 $x\oplus q$ 为 $x+q$ 的小数部分 (即模 $1$ 的平移), 对 $V\subseteq[0,1)$ 记 $V\oplus q=\{v\oplus q: v\in V\}$. 假设存在定义在 $[0,1)$ 的\emph{全体}子集上的函数 $\P$, 满足 $\P([0,1))=1$、可列可加性 (定义~\ref{def:probability} 中的条件 (2)), 以及平移不变性 $\P(V\oplus q)=\P(V)$. 我们来导出矛盾.

在 $[0,1)$ 上定义等价关系: $x\sim y$ 当且仅当 $x-y\in\Q$. 由选择公理, 可以从每个等价类中恰好取出一个元素, 组成集合 $V$. 记 $Q=\Q\cap[0,1)$, 它是可列集.
\begin{itemize}
\item \emph{$\{V\oplus q\}_{q\in Q}$ 两两不交.} 若 $x\in(V\oplus q)\cap(V\oplus r)$, 则存在 $v,v'\in V$ 使 $x=v\oplus q=v'\oplus r$, 于是 $v-v'$ 等于 $r-q$ 加上某个整数, 是有理数. 从而 $v\sim v'$. 由于 $V$ 在每个等价类中只取一个元素, $v=v'$, 进而 $q-r$ 是整数. 又因为 $q,r\in[0,1)$, 所以 $q=r$.
\item \emph{$\bigcup_{q\in Q}(V\oplus q)=[0,1)$.} 任取 $x\in[0,1)$, 设 $v\in V$ 是 $x$ 所在等价类的代表元, 则 $x-v\in\Q\cap(-1,1)$. 若 $x-v\ge 0$, 取 $q=x-v$; 否则取 $q=x-v+1$. 两种情况都有 $q\in Q$ 且 $x=v\oplus q$.
\end{itemize}
由可列可加性与平移不变性,
\[
1=\P([0,1))=\sum_{q\in Q}\P(V\oplus q)=\sum_{q\in Q}\P(V).
\]
右边是可列个相同非负数之和: 若 $\P(V)=0$, 和为 $0$; 若 $\P(V)>0$, 和为 $+\infty$. 两者都不等于 $1$, 矛盾.
\end{example}

因此我们只能在比 $2^\Omega$ 小的 $\sigma$-代数上定义概率. 最常用的办法是先指定一族 ``基本'' 集合 (比如区间), 然后取包含它们的最小 $\sigma$-代数.

\begin{lemma}\label{lem:intersection-sigma}
设 $\{\mathcal{F}_i\}_{i\in I}$ 是 $\Omega$ 上的一族 $\sigma$-代数 ($I$ 为任意非空指标集), 则 $\bigcap_{i\in I}\mathcal{F}_i$ 也是 $\sigma$-代数.
\end{lemma}

\begin{proof}
记 $\mathcal{G}=\bigcap_{i\in I}\mathcal{F}_i$. 对每个 $i$ 都有 $\Omega\in\mathcal{F}_i$, 故 $\Omega\in\mathcal{G}$. 若 $A\in\mathcal{G}$, 则对每个 $i$ 都有 $A\in\mathcal{F}_i$, 从而 $A^c\in\mathcal{F}_i$, 故 $A^c\in\mathcal{G}$. 若 $A_1,A_2,\dots\in\mathcal{G}$, 则对每个 $i$ 都有 $\bigcup_n A_n\in\mathcal{F}_i$, 故 $\bigcup_n A_n\in\mathcal{G}$.
\end{proof}

\begin{definition}[生成的 $\sigma$-代数]\label{def:generated-sigma}
设 $\mathcal{C}\subseteq 2^\Omega$ 是任意子集族. 所有包含 $\mathcal{C}$ 的 $\sigma$-代数之交 (这样的 $\sigma$-代数至少有 $2^\Omega$ 一个, 由引理~\ref{lem:intersection-sigma}, 交仍是 $\sigma$-代数) 称为由 $\mathcal{C}$ \textbf{生成的 $\sigma$-代数}, 记为 $\sigma(\mathcal{C})$. 它是包含 $\mathcal{C}$ 的\emph{最小} $\sigma$-代数: 若 $\mathcal{G}$ 是 $\sigma$-代数且 $\mathcal{C}\subseteq\mathcal{G}$, 则 $\sigma(\mathcal{C})\subseteq\mathcal{G}$.
\end{definition}

``最小性'' 给出了一个反复使用的证明模式: 要证 $\sigma(\mathcal{C})\subseteq\mathcal{G}$, 只需验证 $\mathcal{G}$ 是 $\sigma$-代数且 $\mathcal{C}\subseteq\mathcal{G}$. 例如 $\sigma(\{A\})=\{\varnothing,A,A^c,\Omega\}$, 而例~\ref{ex:sigma-examples} (4) 中的 $\sigma$-代数就是 $\sigma(\{D_1,\dots,D_k\})$.

\begin{definition}[Borel $\sigma$-代数]\label{def:borel}
$\R$ 上由全体开集生成的 $\sigma$-代数称为 \textbf{Borel $\sigma$-代数}, 记为 $\mathcal{B}(\R)$, 其元素称为 Borel 集. 类似地, 对 $\R$ 的子集 $E$ (如 $E=[0,1]$), $\mathcal{B}(E)=\{B\cap E: B\in\mathcal{B}(\R)\}$.
\end{definition}

开集、闭集、单点集、可列集以及各种区间都是 Borel 集. 实际使用中最方便的是下面这组生成元.

\begin{proposition}\label{prop:borel-generator}
$\mathcal{B}(\R)=\sigma\big(\{(-\infty,x]: x\in\R\}\big)$.
\end{proposition}

\begin{proof}
记 $\mathcal{C}=\{(-\infty,x]: x\in\R\}$.

``$\supseteq$'': 每个 $(-\infty,x]$ 是闭集, 是开集 $(x,\infty)$ 的补集, 故属于 $\mathcal{B}(\R)$. 由 $\sigma(\mathcal{C})$ 的最小性, $\sigma(\mathcal{C})\subseteq\mathcal{B}(\R)$.

``$\subseteq$'': 只需证明每个开集属于 $\sigma(\mathcal{C})$. 对 $a<b$, 有 $(a,b]=(-\infty,b]\setminus(-\infty,a]\in\sigma(\mathcal{C})$, 进而
\[
(a,b)=\bigcup_{n>1/(b-a)}\Big(a,\,b-\frac{1}{n}\Big]\in\sigma(\mathcal{C}).
\]
设 $G$ 是开集. 对每个 $x\in G$, 存在 $\varepsilon>0$ 使 $(x-\varepsilon,x+\varepsilon)\subseteq G$, 再取有理数 $p,q$ 满足 $x-\varepsilon<p<x<q<x+\varepsilon$, 便有 $x\in(p,q)\subseteq G$. 因此 $G$ 是所有满足 $(p,q)\subseteq G$ 的有理端点区间 $(p,q)$ 之并. 这样的区间至多可列个, 故 $G\in\sigma(\mathcal{C})$. 由最小性, $\mathcal{B}(\R)\subseteq\sigma(\mathcal{C})$.
\end{proof}

\begin{remark}
同样的论证表明, $\mathcal{B}(\R)$ 也可以由 $\{(a,b): a<b\}$、$\{(-\infty,x]: x\in\Q\}$ 等生成. 生成元越 ``简单'', 验证某个性质对全体 Borel 集成立时就越方便, 这正是 1.1.3 小节 $\pi$-$\lambda$ 定理的用武之地.
\end{remark}

% ---------------------------------------------------------------------
\subsection{概率测度}

\begin{definition}[概率测度, Kolmogorov 公理]\label{def:probability}
设 $(\Omega,\mathcal{F})$ 是可测空间. 函数 $\P:\mathcal{F}\to[0,1]$ 称为\textbf{概率测度}, 如果
\begin{enumerate}
\item 规范性: $\P(\Omega)=1$;
\item 可列可加性: 若 $A_1,A_2,\dots\in\mathcal{F}$ 两两不交, 则
$\P\Big(\bigcup_{n=1}^\infty A_n\Big)=\sum_{n=1}^\infty\P(A_n)$.
\end{enumerate}
三元组 $(\Omega,\mathcal{F},\P)$ 称为\textbf{概率空间}. 若 $\P(A)=1$, 称事件 $A$ \textbf{几乎必然} (almost surely, 简记 $\as$) 发生; 若 $\P(N)=0$, 称 $N$ 为\textbf{零概率集}.
\end{definition}

\begin{remark}
若把值域 $[0,1]$ 换成 $[0,\infty]$, 并把规范性换成 $\mu(\varnothing)=0$, 得到一般的\textbf{测度} $\mu$. 本书用到的非概率测度主要是 $\R$ 上的 \textbf{Lebesgue 测度} $\lambda$, 它满足 $\lambda((a,b])=b-a$.
\end{remark}

\begin{proposition}[概率的基本性质]\label{prop:prob-basic}
设 $(\Omega,\mathcal{F},\P)$ 是概率空间, $A,B,A_1,A_2,\dots\in\mathcal{F}$. 则
\begin{enumerate}
\item $\P(\varnothing)=0$;
\item 有限可加性: 若 $A_1,\dots,A_k$ 两两不交, 则 $\P\big(\bigcup_{n=1}^k A_n\big)=\sum_{n=1}^k\P(A_n)$;
\item $\P(A^c)=1-\P(A)$;
\item 若 $A\subseteq B$, 则 $\P(B\setminus A)=\P(B)-\P(A)$, 特别地 $\P(A)\le\P(B)$ (单调性);
\item $\P(A\cup B)=\P(A)+\P(B)-\P(A\cap B)$;
\item 可列次可加性: $\P\Big(\bigcup_{n=1}^\infty A_n\Big)\le\sum_{n=1}^\infty\P(A_n)$.
\end{enumerate}
\end{proposition}

\begin{proof}
(1) 取 $A_1=\Omega$, $A_n=\varnothing$ ($n\ge2$), 它们两两不交. 由可列可加性, $1=\P(\Omega)=1+\sum_{n\ge2}\P(\varnothing)$. 右端的级数各项相同且非负, 只能都为 $0$.

(2) 令 $A_n=\varnothing$ ($n>k$), 由可列可加性和 (1) 即得.

(3) $A$ 与 $A^c$ 不交且并为 $\Omega$, 由 (2), $\P(A)+\P(A^c)=1$.

(4) $B$ 是不交并 $A\cup(B\setminus A)$, 由 (2), $\P(B)=\P(A)+\P(B\setminus A)$. 由于 $\P(B\setminus A)\ge0$, 得 $\P(A)\le\P(B)$.

(5) 把 $A\cup B$ 写成不交并 $A\cup(B\setminus(A\cap B))$, 再由 (2) 和 (4),
\begin{align*}
\P(A\cup B)&=\P(A)+\P\big(B\setminus(A\cap B)\big)\\
&=\P(A)+\P(B)-\P(A\cap B).
\end{align*}

(6) 令 $B_1=A_1$, $B_n=A_n\setminus\bigcup_{k<n}A_k$ ($n\ge2$). 则 $B_n$ 两两不交, $B_n\subseteq A_n$, 且 $\bigcup_n B_n=\bigcup_n A_n$ (对任意 $\omega\in\bigcup_n A_n$, 取满足 $\omega\in A_n$ 的最小指标 $n$, 则该样本点属于 $B_n$). 由可列可加性和单调性,
\[
\P\Big(\bigcup_{n=1}^\infty A_n\Big)=\sum_{n=1}^\infty\P(B_n)\le\sum_{n=1}^\infty\P(A_n). \qedhere
\]
\end{proof}

可列可加性最重要的推论是概率与单调集合序列的极限可以交换. 记号 $A_n\uparrow A$ 表示 $A_1\subseteq A_2\subseteq\cdots$ 且 $A=\bigcup_n A_n$; $A_n\downarrow A$ 表示 $A_1\supseteq A_2\supseteq\cdots$ 且 $A=\bigcap_n A_n$.

\begin{theorem}[概率的连续性]\label{thm:continuity}
设 $(\Omega,\mathcal{F},\P)$ 是概率空间.
\begin{enumerate}
\item 下连续性: 若 $A_n\uparrow A$, 则 $\P(A_n)\uparrow\P(A)$.
\item 上连续性: 若 $A_n\downarrow A$, 则 $\P(A_n)\downarrow\P(A)$.
\end{enumerate}
\end{theorem}

\begin{proof}
(1) 令 $B_1=A_1$, $B_n=A_n\setminus A_{n-1}$ ($n\ge2$). 由于 $A_n$ 单调增加, $B_n$ 两两不交, 且 $A_n=\bigcup_{k=1}^n B_k$, $A=\bigcup_{k=1}^\infty B_k$. 由可列可加性与有限可加性,
\[
\P(A)=\sum_{k=1}^\infty\P(B_k)=\lim_{n\to\infty}\sum_{k=1}^n\P(B_k)=\lim_{n\to\infty}\P(A_n).
\]
由单调性, $\P(A_n)$ 单调不减, 所以收敛是单调的.

(2) 由 De Morgan 公式, $A_n^c\uparrow A^c$. 由 (1), $\P(A_n^c)\uparrow\P(A^c)$, 即 $1-\P(A_n)\uparrow1-\P(A)$.
\end{proof}

\begin{intuition}
概率是 ``连续'' 的集合函数: 集合单调地逼近 $A$, 概率就逼近 $\P(A)$. 反过来, 可列可加性等价于 ``有限可加性 + 下连续性'': 若 $\P$ 有限可加且下连续, 对两两不交的 $A_n$, 令 $C_n=\bigcup_{k\le n}A_k\uparrow\bigcup_k A_k$, 则 $\P\big(\bigcup_k A_k\big)=\lim_n\P(C_n)=\lim_n\sum_{k\le n}\P(A_k)$. 后文一切极限定理 (单调收敛定理、控制收敛定理) 的根源都在这里.
\end{intuition}

\begin{example}[离散概率空间]\label{ex:discrete-space}
设 $\Omega$ 是有限集或可列集, $p:\Omega\to[0,1]$ 满足 $\sum_{\omega\in\Omega}p(\omega)=1$. 取 $\mathcal{F}=2^\Omega$, 令 $\P(A)=\sum_{\omega\in A}p(\omega)$. 由于非负项级数可以任意重排和分组求和, 对两两不交的 $A_n$ 有
$\P\big(\bigcup_n A_n\big)=\sum_{\omega\in\bigcup_n A_n}p(\omega)=\sum_n\sum_{\omega\in A_n}p(\omega)=\sum_n\P(A_n)$, 故 $\P$ 是概率测度. 掷骰子、有限次抛硬币等古典概型都属于此类.
\end{example}

当 $\Omega$ 不可数时, 我们通常只能在一小族 ``基本集合'' 上直接写出概率 (比如区间的长度、只依赖前 $n$ 次抛掷的事件的概率), 然后需要把它扩张到这族集合生成的 $\sigma$-代数上. 这个扩张由下面的测度论定理保证.

\begin{definition}[代数]\label{def:algebra}
子集族 $\mathcal{A}\subseteq2^\Omega$ 称为\textbf{代数} (或域), 如果 $\Omega\in\mathcal{A}$, 且 $\mathcal{A}$ 对补运算和有限并封闭.
\end{definition}

\begin{theorem}[Carath\'eodory 扩张定理]\label{thm:caratheodory}
设 $\mathcal{A}$ 是 $\Omega$ 上的代数, $\P_0:\mathcal{A}\to[0,1]$ 满足 $\P_0(\Omega)=1$, 并且在 $\mathcal{A}$ 上可列可加 (即只要两两不交的 $A_n\in\mathcal{A}$ 满足 $\bigcup_n A_n\in\mathcal{A}$, 就有 $\P_0(\bigcup_n A_n)=\sum_n\P_0(A_n)$). 则存在 $\sigma(\mathcal{A})$ 上\emph{唯一}的概率测度 $\P$, 使得 $\P$ 在 $\mathcal{A}$ 上与 $\P_0$ 一致.
\end{theorem}

这是纯测度论的结果, 存在性的证明从略 (见~\cite{durrett} 附录 A 或~\cite{billingsley} 第 3 节); 唯一性将在推论~\ref{cor:uniqueness} 中由 $\pi$-$\lambda$ 定理给出.

\begin{example}\label{ex:basic-spaces}
\begin{enumerate}
\item \textbf{$[0,1]$ 上的均匀分布.} 取 $\Omega=[0,1]$, $\mathcal{F}=\mathcal{B}([0,1])$. 形如 $(a,b]\cap[0,1]$ 的区间的有限不交并构成代数, 在其上令 $\P_0$ 为总长度. 可以验证 $\P_0$ 在此代数上可列可加, 由 Carath\'eodory 扩张定理~\ref{thm:caratheodory}, 得到 $\mathcal{B}([0,1])$ 上唯一的概率测度 $\P$, 满足 $\P((a,b])=b-a$. 它就是 $[0,1]$ 上的 Lebesgue 测度. 由例~\ref{ex:vitali}, 它无法再扩张到 $[0,1]$ 的全体子集上.
\item \textbf{无穷次抛硬币.} 取 $\Omega=\{0,1\}^{\N}$. 对 $a_1,\dots,a_n\in\{0,1\}$, 称 $\{\omega:\omega_1=a_1,\dots,\omega_n=a_n\}$ 为\textbf{柱集}, 规定其概率为 $2^{-n}$. 柱集的有限不交并构成代数, 并且可以验证上述赋值在其上可列可加. 由 Carath\'eodory 扩张定理, 得到 $\mathcal{F}=\sigma(\text{柱集})$ 上的概率测度 $\P$. 引言中的事件 \eqref{eq:1-freq-event} 由柱集经可列次运算得到, 所以属于 $\mathcal{F}$, 其概率有意义 (强大数定律断言它等于 $1$).
\end{enumerate}
\end{example}

% ---------------------------------------------------------------------
\subsection{\texorpdfstring{$\pi$-$\lambda$}{pi-lambda} 定理}

$\sigma$-代数通常没有显式的描述: 我们无法 ``列出'' 全体 Borel 集. 于是, 要证明某个性质对 $\sigma(\mathcal{C})$ 中的每个集合成立, 一般采用 ``好集合原理'': 把满足该性质的集合收集起来, 记为 $\mathcal{G}$, 然后证明 $\mathcal{G}\supseteq\sigma(\mathcal{C})$. 困难在于, $\mathcal{G}$ 常常只对 ``不交并'' 和 ``差'' 封闭 (因为概率只在这些运算下有好的公式), 未必是 $\sigma$-代数. $\pi$-$\lambda$ 定理正是为解决这一困难而设计的.

\begin{definition}[$\pi$-系与 $\lambda$-系]\label{def:pi-lambda}
设 $\mathcal{P},\mathcal{L}\subseteq2^\Omega$.
\begin{enumerate}
\item 称 $\mathcal{P}$ 为 \textbf{$\pi$-系}, 如果它对有限交封闭: $A,B\in\mathcal{P}\Rightarrow A\cap B\in\mathcal{P}$.
\item 称 $\mathcal{L}$ 为 \textbf{$\lambda$-系} (或 Dynkin 系), 如果
  \begin{enumerate}[(a)]
  \item $\Omega\in\mathcal{L}$;
  \item 对真差封闭: 若 $A,B\in\mathcal{L}$ 且 $A\subseteq B$, 则 $B\setminus A\in\mathcal{L}$;
  \item 对单调增极限封闭: 若 $A_n\in\mathcal{L}$ 且 $A_n\uparrow A$, 则 $A\in\mathcal{L}$.
  \end{enumerate}
\end{enumerate}
\end{definition}

\begin{example}\label{ex:pi-lambda}
\begin{enumerate}
\item $\{(-\infty,x]: x\in\R\}$ 是 $\pi$-系, 因为 $(-\infty,x]\cap(-\infty,y]=(-\infty,x\wedge y]$. 加上空集后, 无穷次抛硬币的柱集全体也是 $\pi$-系.
\item 每个 $\sigma$-代数既是 $\pi$-系又是 $\lambda$-系. $\lambda$-系对补运算封闭, 因为 $A^c=\Omega\setminus A$.
\item $\lambda$-系未必是 $\pi$-系. 取 $\Omega=\{1,2,3,4\}$,
\[
\mathcal{L}=\big\{\varnothing,\Omega,\{1,2\},\{3,4\},\{1,3\},\{2,4\}\big\}.
\]
$\mathcal{L}$ 中的真包含关系只有 $\varnothing\subseteq X$ 与 $X\subseteq\Omega$ 两类, 其差分别为 $X$ 与 $X^c$, 都在 $\mathcal{L}$ 中; $\mathcal{L}$ 是有限族, 单调增序列最终为常数. 所以 $\mathcal{L}$ 是 $\lambda$-系. 但 $\{1,2\}\cap\{1,3\}=\{1\}\notin\mathcal{L}$, 它不是 $\pi$-系.
\end{enumerate}
\end{example}

\begin{lemma}\label{lem:pi-and-lambda}
若 $\mathcal{F}$ 既是 $\pi$-系又是 $\lambda$-系, 则 $\mathcal{F}$ 是 $\sigma$-代数.
\end{lemma}

\begin{proof}
由 $\lambda$-系的性质 (a), $\Omega\in\mathcal{F}$; 由 (a)(b), $A^c=\Omega\setminus A\in\mathcal{F}$. 对 $A,B\in\mathcal{F}$, 由补运算封闭与 $\pi$-系性质, $A\cup B=(A^c\cap B^c)^c\in\mathcal{F}$, 归纳可知 $\mathcal{F}$ 对有限并封闭. 对 $A_1,A_2,\dots\in\mathcal{F}$, 令 $C_n=\bigcup_{k=1}^n A_k\in\mathcal{F}$, 则 $C_n\uparrow\bigcup_k A_k$, 由性质 (c), $\bigcup_k A_k\in\mathcal{F}$.
\end{proof}

\begin{theorem}[Dynkin $\pi$-$\lambda$ 定理]\label{thm:pi-lambda}
设 $\mathcal{P}$ 是 $\pi$-系, $\mathcal{L}$ 是 $\lambda$-系, 且 $\mathcal{P}\subseteq\mathcal{L}$. 则 $\sigma(\mathcal{P})\subseteq\mathcal{L}$.
\end{theorem}

\begin{proof}
与引理~\ref{lem:intersection-sigma} 的证明完全相同, 任意多个 $\lambda$-系之交仍是 $\lambda$-系. 令 $\ell(\mathcal{P})$ 为包含 $\mathcal{P}$ 的所有 $\lambda$-系之交, 它是包含 $\mathcal{P}$ 的最小 $\lambda$-系, 且 $\ell(\mathcal{P})\subseteq\mathcal{L}$. 下面证明 $\ell(\mathcal{P})$ 也是 $\pi$-系. 一旦证得, 由引理~\ref{lem:pi-and-lambda}, $\ell(\mathcal{P})$ 是包含 $\mathcal{P}$ 的 $\sigma$-代数, 于是
\[
\sigma(\mathcal{P})\subseteq\ell(\mathcal{P})\subseteq\mathcal{L}.
\]

对集合 $A\subseteq\Omega$, 令 $\mathcal{G}_A=\{B\subseteq\Omega: A\cap B\in\ell(\mathcal{P})\}$.

\emph{第一步: 若 $A\in\ell(\mathcal{P})$, 则 $\mathcal{G}_A$ 是 $\lambda$-系.} 逐条验证:
(a) $A\cap\Omega=A\in\ell(\mathcal{P})$, 故 $\Omega\in\mathcal{G}_A$.
(b) 设 $B,C\in\mathcal{G}_A$, $B\subseteq C$. 则 $A\cap B\subseteq A\cap C$ 都属于 $\ell(\mathcal{P})$, 而
$A\cap(C\setminus B)=(A\cap C)\setminus(A\cap B)$, 由 $\ell(\mathcal{P})$ 对真差封闭, $C\setminus B\in\mathcal{G}_A$.
(c) 设 $B_n\in\mathcal{G}_A$, $B_n\uparrow B$. 则 $A\cap B_n\uparrow A\cap B$, 由 $\ell(\mathcal{P})$ 对单调增极限封闭, $B\in\mathcal{G}_A$.

\emph{第二步: 若 $A\in\mathcal{P}$, 则 $\ell(\mathcal{P})\subseteq\mathcal{G}_A$.} 因为 $\mathcal{P}$ 是 $\pi$-系, 对任意 $B\in\mathcal{P}$ 有 $A\cap B\in\mathcal{P}\subseteq\ell(\mathcal{P})$, 即 $\mathcal{P}\subseteq\mathcal{G}_A$. 由第一步 ($A\in\mathcal{P}\subseteq\ell(\mathcal{P})$), $\mathcal{G}_A$ 是 $\lambda$-系; 由 $\ell(\mathcal{P})$ 的最小性, $\ell(\mathcal{P})\subseteq\mathcal{G}_A$, 即
\[
A\in\mathcal{P},\ B\in\ell(\mathcal{P})\ \Longrightarrow\ A\cap B\in\ell(\mathcal{P}).
\]

\emph{第三步: 若 $B\in\ell(\mathcal{P})$, 则 $\ell(\mathcal{P})\subseteq\mathcal{G}_B$.} 由第二步, 对任意 $A\in\mathcal{P}$ 有 $A\cap B\in\ell(\mathcal{P})$, 即 $\mathcal{P}\subseteq\mathcal{G}_B$. 由第一步, $\mathcal{G}_B$ 是 $\lambda$-系; 由最小性, $\ell(\mathcal{P})\subseteq\mathcal{G}_B$. 因此对任意 $B,C\in\ell(\mathcal{P})$ 都有 $B\cap C\in\ell(\mathcal{P})$, 即 $\ell(\mathcal{P})$ 是 $\pi$-系.
\end{proof}

\begin{figure}[htbp]
\centering
\begin{tikzpicture}
  \draw[SADark,thick,fill=SABlue!6] (0,0) ellipse (4.6 and 1.9);
  \draw[SADark,thick,fill=SABlue!16] (-0.6,0) ellipse (3.2 and 1.35);
  \draw[SADark,thick,fill=SARed!18] (-1.6,0) ellipse (1.5 and 0.75);
  \node at (-1.6,0) {$\mathcal{P}$\ ($\pi$-系)};
  \node[align=center] at (1.25,0) {$\sigma(\mathcal{P})=\ell(\mathcal{P})$};
  \node[align=center] at (3.65,0) {$\mathcal{L}$\\ ($\lambda$-系)};
\end{tikzpicture}
\caption{$\pi$-$\lambda$ 定理: 包含 $\pi$-系 $\mathcal{P}$ 的最小 $\lambda$-系恰为 $\sigma(\mathcal{P})$, 因而任何包含 $\mathcal{P}$ 的 $\lambda$-系都包含 $\sigma(\mathcal{P})$}
\label{fig:pi-lambda}
\end{figure}

\begin{intuition}
$\pi$-$\lambda$ 定理把 ``对 $\sigma(\mathcal{P})$ 验证性质'' 化归为两件容易的事: (i) 在一个对交封闭的简单集族 $\mathcal{P}$ 上验证性质; (ii) 验证 ``好集合'' 对真差和单调增极限封闭. 后者通常可以由概率的有限可加性和连续性直接得到. 证明中的关键是 $\ell(\mathcal{P})$ 对交封闭, 而这需要 $\mathcal{P}$ 本身对交封闭.
\end{intuition}

\begin{remark}
测度论教材中还有一个形式相近的 \textbf{单调类定理}: 若 $\mathcal{A}$ 是代数, 则包含 $\mathcal{A}$ 的最小单调类 (对单调增、单调减极限封闭的集族) 等于 $\sigma(\mathcal{A})$. 它与 $\pi$-$\lambda$ 定理的作用相同, 但要求生成元构成代数, 条件更强. 本书统一使用 $\pi$-$\lambda$ 定理.
\end{remark}

下面是 $\pi$-$\lambda$ 定理最典型的应用.

\begin{corollary}[概率测度的唯一性]\label{cor:uniqueness}
设 $\P_1,\P_2$ 是 $(\Omega,\mathcal{F})$ 上的两个概率测度, $\mathcal{P}\subseteq\mathcal{F}$ 是 $\pi$-系. 若 $\P_1$ 与 $\P_2$ 在 $\mathcal{P}$ 上相等, 则它们在 $\sigma(\mathcal{P})$ 上相等.
\end{corollary}

\begin{proof}
令 $\mathcal{L}=\{A\in\mathcal{F}: \P_1(A)=\P_2(A)\}$. 由假设, $\mathcal{P}\subseteq\mathcal{L}$. 验证 $\mathcal{L}$ 是 $\lambda$-系:
(a) $\P_1(\Omega)=1=\P_2(\Omega)$;
(b) 若 $A,B\in\mathcal{L}$, $A\subseteq B$, 则由命题~\ref{prop:prob-basic} (4), $\P_1(B\setminus A)=\P_1(B)-\P_1(A)=\P_2(B)-\P_2(A)=\P_2(B\setminus A)$;
(c) 若 $A_n\in\mathcal{L}$, $A_n\uparrow A$, 则由概率的连续性 (定理~\ref{thm:continuity}), $\P_1(A)=\lim_n\P_1(A_n)=\lim_n\P_2(A_n)=\P_2(A)$.
由 $\pi$-$\lambda$ 定理~\ref{thm:pi-lambda}, $\sigma(\mathcal{P})\subseteq\mathcal{L}$.
\end{proof}

特别地, 由于代数是 $\pi$-系, Carath\'eodory 扩张定理~\ref{thm:caratheodory} 中扩张的唯一性由此推出; 由命题~\ref{prop:borel-generator} 和例~\ref{ex:pi-lambda} (1), $(\R,\mathcal{B}(\R))$ 上的概率测度由它在全体 $(-\infty,x]$ 上的值唯一确定. 这正是 1.2 节 ``分布函数决定分布'' 的来源.

\begin{example}[$\pi$-系条件不可去掉]
取 $\Omega=\{1,2,3,4\}$, $\mathcal{F}=2^\Omega$, $\mathcal{C}=\{\{1,2\},\{1,3\}\}$. 令 $\P_1$ 为均匀分布, $\P_2$ 给 $1$ 和 $4$ 各分配概率 $\frac{1}{2}$. 则
$\P_1(\{1,2\})=\P_2(\{1,2\})=\frac{1}{2}$, $\P_1(\{1,3\})=\P_2(\{1,3\})=\frac{1}{2}$, 但 $\{1\}=\{1,2\}\cap\{1,3\}\in\sigma(\mathcal{C})$, 而 $\P_1(\{1\})=\frac{1}{4}\ne\frac{1}{2}=\P_2(\{1\})$. 这是因为 $\mathcal{C}$ 不是 $\pi$-系, 两个测度在 $\mathcal{C}$ 中的集合上一致, 却在它们的交集上不一致.
\end{example}

% ---------------------------------------------------------------------
\subsection{集合序列的极限与 Borel-Cantelli 第一引理}

许多重要事件形如 ``某事无穷多次发生'', 例如 ``随机游动无穷多次回到原点'', ``布朗运动在任意小的时间内无穷多次穿过零点''. 描述这类事件需要集合序列的上、下极限.

\begin{definition}[集合序列的上、下极限]\label{def:limsup-sets}
设 $A_1,A_2,\dots\subseteq\Omega$. 定义
\[
\limsup_{n\to\infty}A_n=\bigcap_{m=1}^\infty\bigcup_{n=m}^\infty A_n,\qquad
\liminf_{n\to\infty}A_n=\bigcup_{m=1}^\infty\bigcap_{n=m}^\infty A_n.
\]
若二者相等, 称 $\{A_n\}$ 的极限存在, 记为 $\lim_n A_n$.
\end{definition}

\begin{proposition}\label{prop:limsup-meaning}
\begin{enumerate}
\item $\omega\in\limsup_n A_n$ 当且仅当有无穷多个 $n$ 使 $\omega\in A_n$. 因此 $\limsup_n A_n$ 记作 $\{A_n\ \text{i.o.}\}$ (infinitely often), 读作 ``$A_n$ 无穷多次发生''.
\item $\omega\in\liminf_n A_n$ 当且仅当存在 $m$, 使得对所有 $n\ge m$ 都有 $\omega\in A_n$. 读作 ``$A_n$ 终将一直发生''.
\item $\liminf_n A_n\subseteq\limsup_n A_n$, 且 $\big(\limsup_n A_n\big)^c=\liminf_n A_n^c$.
\item $\1_{\limsup_n A_n}=\limsup_n\1_{A_n}$, $\1_{\liminf_n A_n}=\liminf_n\1_{A_n}$, 其中 $\1_A$ 是示性函数: $\omega\in A$ 时 $\1_A(\omega)=1$, 否则为 $0$.
\end{enumerate}
\end{proposition}

\begin{proof}
(1) $\omega\in\bigcap_m\bigcup_{n\ge m}A_n$ 等价于: 对每个 $m$, 存在 $n\ge m$ 使 $\omega\in A_n$. 这又等价于使 $\omega\in A_n$ 的 $n$ 不能有上界, 即有无穷多个.

(2) $\omega\in\bigcup_m\bigcap_{n\ge m}A_n$ 等价于: 存在 $m$, 对所有 $n\ge m$, $\omega\in A_n$.

(3) 若 $\omega$ 从某项起一直属于 $A_n$, 它当然属于无穷多个 $A_n$, 故第一个包含关系成立. 由 De Morgan 公式,
\[
\Big(\bigcap_{m}\bigcup_{n\ge m}A_n\Big)^c=\bigcup_m\Big(\bigcup_{n\ge m}A_n\Big)^c=\bigcup_m\bigcap_{n\ge m}A_n^c.
\]

(4) 数列 $\1_{A_n}(\omega)$ 只取 $0,1$ 两个值. 它的上极限为 $1$ 当且仅当它无穷多次取 $1$, 由 (1) 即当且仅当 $\omega\in\limsup_n A_n$. 下极限为 $1$ 当且仅当它从某项起恒为 $1$, 由 (2) 即当且仅当 $\omega\in\liminf_n A_n$.
\end{proof}

\begin{figure}[htbp]
\centering
\begin{tikzpicture}[xscale=0.6,yscale=0.8]
  \foreach \n in {1,...,14} { \node[font=\small,SADark!70] at (\n,3.6) {\n}; }
  \node[font=\small,SADark!70,anchor=east] at (0.3,3.6) {$n$};
  \foreach \y/\lab in {2.6/{\omega_1},1.6/{\omega_2},0.6/{\omega_3}} {
    \draw[SADark!25] (0.6,\y) -- (14.6,\y);
    \node[anchor=east] at (0.3,\y) {$\lab$};
  }
  % omega_1: eventually always
  \foreach \n in {2,4,5,7,8,9,10,11,12,13,14} { \fill[SABlue] (\n,2.6) circle (3.2pt); }
  \foreach \n in {1,3,6} { \draw[SABlue] (\n,2.6) circle (3.2pt); }
  % omega_2: infinitely often, infinitely often not
  \foreach \n in {1,3,4,7,10,13} { \fill[SARed] (\n,1.6) circle (3.2pt); }
  \foreach \n in {2,5,6,8,9,11,12,14} { \draw[SARed] (\n,1.6) circle (3.2pt); }
  % omega_3: finitely many
  \foreach \n in {1,2,5} { \fill[SAGreenText] (\n,0.6) circle (3.2pt); }
  \foreach \n in {3,4,6,7,8,9,10,11,12,13,14} { \draw[SAGreenText] (\n,0.6) circle (3.2pt); }
  \node[font=\small] at (15.6,2.6) {$\cdots$};
  \node[font=\small] at (15.6,1.6) {$\cdots$};
  \node[font=\small] at (15.6,0.6) {$\cdots$};
  \node[anchor=west,font=\small] at (16.2,2.6) {$\in\liminf A_n$};
  \node[anchor=west,font=\small] at (16.2,1.6) {$\in\limsup A_n$, $\notin\liminf A_n$};
  \node[anchor=west,font=\small] at (16.2,0.6) {$\notin\limsup A_n$};
\end{tikzpicture}
\caption{三个样本点关于 $A_1,A_2,\dots$ 的归属 (实心: $\omega\in A_n$; 空心: $\omega\notin A_n$). $\omega_1$ 从某项起一直在 $A_n$ 中; $\omega_2$ 无穷多次进入、也无穷多次离开; $\omega_3$ 只进入有限次}
\label{fig:limsup}
\end{figure}

\begin{proposition}[集合的 Fatou 不等式]\label{prop:fatou-sets}
设 $A_n\in\mathcal{F}$. 则
\[
\P\Big(\liminf_{n\to\infty}A_n\Big)\le\liminf_{n\to\infty}\P(A_n)\le\limsup_{n\to\infty}\P(A_n)\le\P\Big(\limsup_{n\to\infty}A_n\Big).
\]
\end{proposition}

\begin{proof}
令 $B_m=\bigcap_{n\ge m}A_n$, 则 $B_m\uparrow\liminf_n A_n$, 且对每个 $n\ge m$ 有 $B_m\subseteq A_n$, 从而 $\P(B_m)\le\inf_{n\ge m}\P(A_n)$. 由概率的连续性 (定理~\ref{thm:continuity}),
\[
\P\Big(\liminf_{n\to\infty}A_n\Big)=\lim_{m\to\infty}\P(B_m)\le\lim_{m\to\infty}\inf_{n\ge m}\P(A_n)=\liminf_{n\to\infty}\P(A_n).
\]
中间的不等式是数列上下极限的基本性质. 对 $C_m=\bigcup_{n\ge m}A_n\downarrow\limsup_n A_n$ 同理, $\P(C_m)\ge\sup_{n\ge m}\P(A_n)$, 由上连续性得到最后一个不等式.
\end{proof}

\begin{theorem}[Borel-Cantelli 第一引理]\label{thm:bc1}
设 $A_n\in\mathcal{F}$. 若 $\sum_{n=1}^\infty\P(A_n)<\infty$, 则 $\P(A_n\ \text{i.o.})=0$.
\end{theorem}

\begin{proof}
对每个 $m$, $\limsup_n A_n\subseteq\bigcup_{n\ge m}A_n$. 由单调性与可列次可加性 (命题~\ref{prop:prob-basic} (4)(6)),
\begin{align*}
\P\Big(\limsup_{n\to\infty}A_n\Big)
&\le\P\Big(\bigcup_{n=m}^\infty A_n\Big)\\
&\le\sum_{n=m}^\infty\P(A_n).
\end{align*}
右端是收敛级数的余项, 当 $m\to\infty$ 时趋于 $0$. 而左端与 $m$ 无关, 故为 $0$.
\end{proof}

\begin{intuition}
$\sum_n\P(A_n)$ 是 ``发生次数'' $\sum_n\1_{A_n}$ 的期望 (见 1.3 节). 期望有限, 发生次数自然几乎必然有限. 这里第一引理\emph{不需要任何独立性假设}.
\end{intuition}

\begin{example}\label{ex:bc1}
设 $\P(A_n)=\frac{1}{n^2}$, 则不论 $A_n$ 之间如何相关, 几乎必然只有有限个 $A_n$ 发生. 又如在 $[0,1]$ 上的均匀分布下, 令 $A_n=[0,\frac{1}{n^2}]$, 则 $\limsup_n A_n=\{0\}$, 概率确实为 $0$.
\end{example}

% ---------------------------------------------------------------------
\subsection{独立性与 Kolmogorov 0-1 律}

\begin{definition}[事件的独立性]\label{def:indep-events}
事件 $A,B$ 称为\textbf{独立}的, 如果 $\P(A\cap B)=\P(A)\P(B)$. 一族事件 $\{A_i\}_{i\in I}$ 称为 (相互) 独立的, 如果对任意有限个互不相同的指标 $i_1,\dots,i_k\in I$,
\[
\P(A_{i_1}\cap\cdots\cap A_{i_k})=\P(A_{i_1})\cdots\P(A_{i_k}).
\]
\end{definition}

\begin{example}[两两独立不蕴含相互独立]\label{ex:pairwise}
抛两次均匀硬币, 令 $A=$ ``第一次正面'', $B=$ ``第二次正面'', $C=$ ``两次结果相同''. 则 $\P(A)=\P(B)=\P(C)=\frac{1}{2}$, 且
\[
\P(A\cap B)=\P(A\cap C)=\P(B\cap C)=\P(\text{两次都是正面})=\frac{1}{4},
\]
所以三者两两独立. 但 $\P(A\cap B\cap C)=\frac{1}{4}\ne\frac{1}{8}$, 它们不相互独立. 原因是任意两个事件的结果决定了第三个.
\end{example}

在随机过程理论中, 我们更常讨论的是 ``信息'' 之间的独立性, 即 $\sigma$-代数之间的独立性.

\begin{definition}[$\sigma$-代数的独立性]\label{def:indep-sigma}
设 $\{\mathcal{F}_i\}_{i\in I}$ 是一族包含于 $\mathcal{F}$ 的 $\sigma$-代数. 称它们 (相互) \textbf{独立}, 如果对任意有限个互不相同的指标 $i_1,\dots,i_k$ 以及任意 $A_j\in\mathcal{F}_{i_j}$ ($j=1,\dots,k$),
\[
\P(A_1\cap\cdots\cap A_k)=\P(A_1)\cdots\P(A_k).
\]
\end{definition}

直接验证 $\sigma$-代数的独立性需要检查其中的全部事件, 这通常做不到. 下面的判别准则把检查范围缩小到生成元.

\begin{theorem}[独立性的 $\pi$-系判别准则]\label{thm:indep-pi}
设 $\mathcal{P}_1,\dots,\mathcal{P}_n\subseteq\mathcal{F}$ 都是 $\pi$-系, 且 $\Omega\in\mathcal{P}_i$ ($i=1,\dots,n$). 若对任意 $A_i\in\mathcal{P}_i$ 都有
\begin{equation}\label{eq:indep-product}
\P(A_1\cap\cdots\cap A_n)=\P(A_1)\cdots\P(A_n),
\end{equation}
则 $\sigma(\mathcal{P}_1),\dots,\sigma(\mathcal{P}_n)$ 相互独立.
\end{theorem}

\begin{proof}
\emph{第一步: 化归为 ``$n$ 个指标全部出现'' 的情形.} 定义~\ref{def:indep-sigma} 要求对任意子族成立乘积公式. 对子族 $\{i_1,\dots,i_k\}$, 把不在其中的指标 $i$ 对应的事件取为 $\Omega$, 由于 $\P(\Omega)=1$, 全指标的乘积公式就退化为子族的乘积公式. 所以只需对 $A_i\in\sigma(\mathcal{P}_i)$ ($i=1,\dots,n$) 证明 \eqref{eq:indep-product}.

\emph{第二步: 把 $\mathcal{P}_1$ 换成 $\sigma(\mathcal{P}_1)$.} 固定 $A_2\in\mathcal{P}_2,\dots,A_n\in\mathcal{P}_n$, 记 $D=A_2\cap\cdots\cap A_n$, $c=\P(A_2)\cdots\P(A_n)$, 令
\[
\mathcal{L}=\{A\in\mathcal{F}:\ \P(A\cap D)=\P(A)\,c\}.
\]
由假设 \eqref{eq:indep-product}, $\mathcal{P}_1\subseteq\mathcal{L}$. 验证 $\mathcal{L}$ 是 $\lambda$-系:
(a) $\Omega\in\mathcal{P}_1\subseteq\mathcal{L}$.
(b) 若 $A,B\in\mathcal{L}$, $A\subseteq B$, 则 $(B\setminus A)\cap D=(B\cap D)\setminus(A\cap D)$, 于是
\begin{align*}
\P\big((B\setminus A)\cap D\big)&=\P(B\cap D)-\P(A\cap D)\\
&=\big(\P(B)-\P(A)\big)c=\P(B\setminus A)\,c.
\end{align*}
(c) 若 $A_k\in\mathcal{L}$, $A_k\uparrow A$, 则 $A_k\cap D\uparrow A\cap D$, 由概率的连续性 (定理~\ref{thm:continuity}),
$\P(A\cap D)=\lim_k\P(A_k\cap D)=\lim_k\P(A_k)\,c=\P(A)\,c$.
由 $\pi$-$\lambda$ 定理~\ref{thm:pi-lambda}, $\sigma(\mathcal{P}_1)\subseteq\mathcal{L}$, 即 \eqref{eq:indep-product} 对 $A_1\in\sigma(\mathcal{P}_1)$, $A_i\in\mathcal{P}_i$ ($i\ge2$) 成立.

\emph{第三步: 逐个替换.} 现在把 $\sigma(\mathcal{P}_1),\mathcal{P}_2,\dots,\mathcal{P}_n$ 看作新的一组 $\pi$-系 ($\sigma$-代数当然是 $\pi$-系, 且包含 $\Omega$), 它们满足 \eqref{eq:indep-product}. 对第二个指标重复第二步的论证, 可把 $\mathcal{P}_2$ 换成 $\sigma(\mathcal{P}_2)$. 依此类推 $n$ 次, 即得 \eqref{eq:indep-product} 对所有 $A_i\in\sigma(\mathcal{P}_i)$ 成立.
\end{proof}

\begin{corollary}\label{cor:indep-events-sigma}
事件 $\{A_i\}_{i\in I}$ 相互独立, 当且仅当 $\sigma$-代数 $\{\sigma(\{A_i\})\}_{i\in I}=\{\{\varnothing,A_i,A_i^c,\Omega\}\}_{i\in I}$ 相互独立. 特别地, 把独立事件族中任意一部分事件换成它们的补事件, 所得事件族仍然独立.
\end{corollary}

\begin{proof}
``当'' 的方向是显然的. 对 ``仅当'' 的方向, 由于独立性只涉及有限子族, 只需对有限个指标 $i_1,\dots,i_n$ 验证. 令 $\mathcal{P}_j=\{A_{i_j},\Omega\}$, 它是包含 $\Omega$ 的 $\pi$-系. 对 $B_j\in\mathcal{P}_j$, 乘积公式 \eqref{eq:indep-product} 中取 $B_j=\Omega$ 的那些因子等于 $1$, 剩下的恰好是事件 $\{A_{i_1},\dots,A_{i_n}\}$ 某个子族的乘积公式, 由事件的独立性成立. 由定理~\ref{thm:indep-pi}, $\sigma(\mathcal{P}_1),\dots,\sigma(\mathcal{P}_n)$ 独立.
\end{proof}

\begin{corollary}[分组原理]\label{cor:grouping}
设 $\{\mathcal{F}_i\}_{i\in I}$ 相互独立, $I=\bigcup_{j\in J}I_j$ 是 $I$ 的一个分割 (各 $I_j$ 两两不交). 令 $\mathcal{G}_j=\sigma\big(\bigcup_{i\in I_j}\mathcal{F}_i\big)$, 则 $\{\mathcal{G}_j\}_{j\in J}$ 相互独立.
\end{corollary}

\begin{proof}
对每个 $j$, 令 $\mathcal{P}_j$ 为所有形如 $B_1\cap\cdots\cap B_m$ 的集合全体, 其中 $m\ge1$, $B_l\in\mathcal{F}_{i_l}$, 而 $i_1,\dots,i_m$ 是 $I_j$ 中互不相同的指标. $\mathcal{P}_j$ 对交封闭 (两个这样的交集之交仍是这种形式: 对同一个 $\mathcal{F}_i$ 中的两个集合先求交, 结果仍在 $\mathcal{F}_i$ 中), 包含 $\Omega$, 且 $\bigcup_{i\in I_j}\mathcal{F}_i\subseteq\mathcal{P}_j\subseteq\mathcal{G}_j$, 故 $\sigma(\mathcal{P}_j)=\mathcal{G}_j$. 任取有限个不同的 $j_1,\dots,j_n$ 以及 $C_k\in\mathcal{P}_{j_k}$, 每个 $C_k$ 都是来自 $I_{j_k}$ 中不同指标的 $\mathcal{F}_i$ 的事件之交; 由于不同的 $I_{j_k}$ 不交, $C_1\cap\cdots\cap C_n$ 是来自互不相同指标的事件之交. 由 $\{\mathcal{F}_i\}$ 的独立性,
$\P(C_1\cap\cdots\cap C_n)$ 等于所有这些事件概率之积, 按 $k$ 重新分组即为 $\P(C_1)\cdots\P(C_n)$. 由定理~\ref{thm:indep-pi}, $\mathcal{G}_{j_1},\dots,\mathcal{G}_{j_n}$ 独立.
\end{proof}

\begin{intuition}
分组原理说: 把相互独立的信息源分成互不重叠的几组, 每组 ``打包'' 后得到的信息仍然相互独立. 例如在无穷次抛硬币中, 令 $\mathcal{F}_n=\sigma(\{\omega_n=1\})$, 则 ``由前 $n$ 次决定的事件'' 与 ``由第 $n$ 次以后决定的事件'' 相互独立. 这一事实在后文定义马尔可夫性和独立增量时反复出现.
\end{intuition}

有了独立性, Borel-Cantelli 第一引理就有了逆命题.

\begin{theorem}[Borel-Cantelli 第二引理]\label{thm:bc2}
设事件 $A_1,A_2,\dots$ 相互独立. 若 $\sum_{n=1}^\infty\P(A_n)=\infty$, 则 $\P(A_n\ \text{i.o.})=1$.
\end{theorem}

\begin{proof}
由命题~\ref{prop:limsup-meaning} (3), 只需证明 $\P\big(\liminf_n A_n^c\big)=0$, 而由可列次可加性, 只需证明对每个 $m$, $\P\big(\bigcap_{n\ge m}A_n^c\big)=0$.

由推论~\ref{cor:indep-events-sigma}, $A_1^c,A_2^c,\dots$ 也相互独立. 利用初等不等式 $1-x\le e^{-x}$ ($x\in\R$), 对 $N\ge m$,
\begin{align*}
\P\Big(\bigcap_{n=m}^N A_n^c\Big)&=\prod_{n=m}^N\big(1-\P(A_n)\big)\\
&\le\prod_{n=m}^N e^{-\P(A_n)}=\exp\Big(-\sum_{n=m}^N\P(A_n)\Big).
\end{align*}
由假设, 当 $N\to\infty$ 时右端趋于 $0$. 又 $\bigcap_{n=m}^N A_n^c\downarrow\bigcap_{n\ge m}A_n^c$ ($N\to\infty$), 由上连续性 (定理~\ref{thm:continuity}), $\P\big(\bigcap_{n\ge m}A_n^c\big)=0$.
\end{proof}

\begin{intuition}
两个 Borel-Cantelli 引理合起来给出一个 ``0-1 二分法'': 对独立事件, $\P(A_n\ \text{i.o.})$ 只能是 $0$ 或 $1$, 取决于 $\sum_n\P(A_n)$ 是否收敛. 级数是否收敛, 刻画的是 $\P(A_n)$ 趋于 $0$ 的速度够不够快.
\end{intuition}

\begin{example}\label{ex:bc2}
\begin{enumerate}
\item 设 $A_n$ 独立, $\P(A_n)=\frac{1}{n}$. 由于调和级数发散, 几乎必然有无穷多个 $A_n$ 发生. 若改为 $\P(A_n)=\frac{1}{n^2}$, 则由第一引理几乎必然只有有限个发生.
\item \textbf{独立性不可去掉.} 取 $\P(A)=\frac{1}{2}$, 令所有 $A_n=A$. 则 $\sum_n\P(A_n)=\infty$, 但 $\P(A_n\ \text{i.o.})=\P(A)=\frac{1}{2}$.
\item \textbf{无穷猴子定理.} 在无穷次抛硬币中, 令 $B_k$ 为 ``第 $3k-2,3k-1,3k$ 次的结果依次为正、正、反''. 由分组原理 (推论~\ref{cor:grouping}), $B_1,B_2,\dots$ 相互独立, 且 $\P(B_k)=\frac{1}{8}$. 由第二引理, 这个图案几乎必然出现无穷多次. 同样的论证适用于任何有限长的图案.
\end{enumerate}
\end{example}

最后介绍一个深刻的结果: 由 ``无穷远处'' 的独立信息决定的事件, 其概率只能是 $0$ 或 $1$.

\begin{definition}[尾 $\sigma$-代数]\label{def:tail}
设 $\mathcal{F}_1,\mathcal{F}_2,\dots$ 是 $\sigma$-代数. 令 $\mathcal{T}_n=\sigma\big(\bigcup_{k>n}\mathcal{F}_k\big)$, 称
\[
\mathcal{T}=\bigcap_{n=1}^\infty\mathcal{T}_n
\]
为\textbf{尾 $\sigma$-代数}, 其中的事件称为\textbf{尾事件}.
\end{definition}

尾事件是这样的事件, 对每个 $n$, 只要知道第 $n$ 个以后的全部信息就能判断它是否发生, 改变有限个 $\mathcal{F}_k$ 中的信息不影响它是否发生.

\begin{example}\label{ex:tail-events}
\begin{enumerate}
\item 若 $A_k\in\mathcal{F}_k$, 则 $\{A_k\ \text{i.o.}\}$ 是尾事件: 对每个 $n$, $\limsup_k A_k=\bigcap_{m>n}\bigcup_{k\ge m}A_k\in\mathcal{T}_n$ (因为 $\bigcup_{k\ge m}A_k$ 关于 $m$ 单调递减, 去掉前面有限个 $m$ 不改变交集).
\item 无穷次抛硬币中, 取 $\mathcal{F}_k=\sigma(\{\omega_k=1\})$. 引言中的事件 $A$ (频率趋于 $\frac{1}{2}$) 是尾事件. 事实上, 对固定的 $n$, 前 $n$ 次抛掷对 $\frac{S_m}{m}$ 的贡献不超过 $\frac{n}{m}\to0$, 所以
$A=\big\{\lim_m\frac{1}{m}\sum_{k=n+1}^m\omega_k=\frac{1}{2}\big\}$, 按 \eqref{eq:1-freq-event} 的方式展开, 可知 $A\in\mathcal{T}_n$.
\end{enumerate}
\end{example}

\begin{theorem}[Kolmogorov 0-1 律]\label{thm:kolmogorov-01}
设 $\mathcal{F}_1,\mathcal{F}_2,\dots$ 是相互独立的 $\sigma$-代数, $\mathcal{T}$ 是它们的尾 $\sigma$-代数. 则对任意 $A\in\mathcal{T}$, $\P(A)\in\{0,1\}$.
\end{theorem}

\begin{proof}
\emph{第一步.} 对每个 $n$, 令 $\mathcal{H}_n=\sigma(\mathcal{F}_1\cup\cdots\cup\mathcal{F}_n)$. 把指标集分成 $\{1,\dots,n\}$ 与 $\{n+1,n+2,\dots\}$ 两组, 由分组原理 (推论~\ref{cor:grouping}), $\mathcal{H}_n$ 与 $\mathcal{T}_n$ 独立. 由于 $\mathcal{T}\subseteq\mathcal{T}_n$, $\mathcal{H}_n$ 与 $\mathcal{T}$ 独立.

\emph{第二步.} 令 $\mathcal{P}=\bigcup_{n=1}^\infty\mathcal{H}_n$. 由于 $\mathcal{H}_1\subseteq\mathcal{H}_2\subseteq\cdots$, 任意 $B,C\in\mathcal{P}$ 都属于同一个 $\mathcal{H}_n$, 从而 $B\cap C\in\mathcal{H}_n\subseteq\mathcal{P}$. 所以 $\mathcal{P}$ 是包含 $\Omega$ 的 $\pi$-系. 由第一步, 对任意 $B\in\mathcal{P}$, $A\in\mathcal{T}$, 有 $\P(B\cap A)=\P(B)\P(A)$. 由独立性的 $\pi$-系判别准则 (定理~\ref{thm:indep-pi}, 取 $\mathcal{P}_1=\mathcal{P}$, $\mathcal{P}_2=\mathcal{T}$), $\sigma(\mathcal{P})$ 与 $\mathcal{T}$ 独立.

\emph{第三步.} 每个 $\mathcal{F}_k$ 都包含于 $\mathcal{P}$, 所以对每个 $n$, $\mathcal{T}_n\subseteq\sigma(\mathcal{P})$, 进而 $\mathcal{T}\subseteq\sigma(\mathcal{P})$. 于是对 $A\in\mathcal{T}$, $A$ 既属于 $\mathcal{T}$ 又属于 $\sigma(\mathcal{P})$, 由第二步, $A$ 与自身独立:
\[
\P(A)=\P(A\cap A)=\P(A)^2.
\]
解得 $\P(A)\in\{0,1\}$.
\end{proof}

\begin{intuition}
尾事件既 ``与任何有限多个 $\mathcal{F}_k$ 无关'', 又 ``完全由全体 $\mathcal{F}_k$ 决定''. 前者使它与全部信息独立, 后者使它属于全部信息, 于是它只能与自己独立, 概率非 $0$ 即 $1$. 注意 0-1 律只告诉我们概率是 $0$ 或 $1$, 并不告诉我们是哪一个; 确定具体是哪一个, 需要额外的工具 (如 Borel-Cantelli 引理或大数定律). 第~3~章将用鞅收敛定理重新证明这一定律.
\end{intuition}

% =====================================================================
\section{随机变量与分布}
\label{sec:1-2}

在实际问题中, 我们通常直接观察由样本点 $\omega$ 决定的数值, 例如骰子的点数、一天后的股价或粒子的位置. 这些数值就是随机变量. 例如在无穷次抛硬币中, 前 $n$ 次正面的次数 $S_n(\omega)=\omega_1+\cdots+\omega_n$ 就是一个随机变量. 对随机变量 $X$, 最基本的问题是 ``$X\le x$ 的概率是多少'', 为此必须保证 $\{\omega: X(\omega)\le x\}$ 是事件. 这一要求引出了可测性.

% ---------------------------------------------------------------------
\subsection{可测映射与随机变量}

\begin{definition}[可测映射与随机变量]\label{def:rv}
设 $(\Omega,\mathcal{F})$ 与 $(E,\mathcal{E})$ 是可测空间. 映射 $X:\Omega\to E$ 称为 \textbf{$\mathcal{F}/\mathcal{E}$-可测}的, 如果对每个 $B\in\mathcal{E}$, 原像
\[
X^{-1}(B)=\{\omega\in\Omega: X(\omega)\in B\}\in\mathcal{F}.
\]
当 $(E,\mathcal{E})=(\R,\mathcal{B}(\R))$ 时, 称 $X$ 为 (实值) \textbf{随机变量}. 原像 $X^{-1}(B)$ 常简记为 $\{X\in B\}$, 例如 $\{X\le x\}=X^{-1}((-\infty,x])$.
\end{definition}

为了处理 $\sup_n X_n$ 这样可能取 $\pm\infty$ 的量, 我们也考虑\textbf{广义随机变量}, 即取值于 $\overline{\R}=[-\infty,+\infty]$ 的可测映射, 其中 $\mathcal{B}(\overline{\R})=\sigma(\{[-\infty,x]: x\in\R\})$.

\begin{figure}[htbp]
\centering
\begin{tikzpicture}
  % Omega
  \draw[SADark,thick,fill=SABlue!5] plot[smooth cycle,tension=0.8] coordinates {(0,0) (1.2,-0.9) (3.1,-0.6) (3.6,0.9) (2.4,2.1) (0.6,1.8)};
  \node[SADark] at (0.25,2.05) {$\Omega$};
  \fill[SABlue!35] plot[smooth cycle,tension=0.9] coordinates {(0.7,0.3) (1.3,-0.2) (1.8,0.4) (1.2,0.9)};
  \fill[SABlue!35] plot[smooth cycle,tension=0.9] coordinates {(2.3,1.0) (2.9,0.7) (3.1,1.3) (2.6,1.6)};
  \node[font=\small] at (1.25,0.35) {$X^{-1}(B)$};
  % real line
  \draw[->,SADark] (5.6,0.5) -- (10.4,0.5) node[right] {$\R$};
  \draw[SARed,line width=2.4pt] (7.3,0.5) -- (8.6,0.5);
  \node[SARed,below] at (7.95,0.45) {$B$};
  % arrows
  \draw[->,SADark!70,bend left=22] (1.5,0.55) to (7.5,0.62);
  \draw[->,SADark!70,bend left=12] (2.75,1.25) to (8.3,0.62);
  \fill (3.2,-0.2) circle (1.3pt) node[below left,font=\small] {$\omega$};
  \draw[->,SADark!50,bend right=15] (3.2,-0.2) to (9.6,0.42);
  \node[font=\small,below] at (9.6,0.4) {$X(\omega)$};
  \node[SADark] at (5.0,1.6) {$X$};
\end{tikzpicture}
\caption{随机变量是从 $\Omega$ 到 $\R$ 的映射. 可测性要求: 对每个 Borel 集 $B$, ``$X$ 落在 $B$ 中'' 的那些样本点构成一个事件}
\label{fig:rv-map}
\end{figure}

逐个检查全体 Borel 集是不可能的. 好在原像运算与集合运算可交换, 这使得只需检查生成元.

\begin{lemma}\label{lem:measurable-generator}
设 $X:\Omega\to E$, $\mathcal{C}\subseteq2^E$. 则
\begin{enumerate}
\item $\{X^{-1}(B): B\in\sigma(\mathcal{C})\}=\sigma\big(\{X^{-1}(C): C\in\mathcal{C}\}\big)$;
\item 若 $\mathcal{E}=\sigma(\mathcal{C})$, 则 $X$ 是 $\mathcal{F}/\mathcal{E}$-可测的, 当且仅当对每个 $C\in\mathcal{C}$, $X^{-1}(C)\in\mathcal{F}$.
\end{enumerate}
\end{lemma}

\begin{proof}
原像运算保持一切集合运算:
\[
X^{-1}(E\setminus B)=\Omega\setminus X^{-1}(B),\qquad X^{-1}\Big(\bigcup_n B_n\Big)=\bigcup_n X^{-1}(B_n).
\]

(1) 记左端为 $\mathcal{G}_1$, 右端为 $\mathcal{G}_2$. 由上式, $\mathcal{G}_1$ 是 $\Omega$ 上的 $\sigma$-代数, 且包含每个 $X^{-1}(C)$, 由最小性 $\mathcal{G}_2\subseteq\mathcal{G}_1$. 反过来, 令 $\mathcal{D}=\{B\subseteq E: X^{-1}(B)\in\mathcal{G}_2\}$. 仍由上式, $\mathcal{D}$ 是 $E$ 上的 $\sigma$-代数, 且 $\mathcal{C}\subseteq\mathcal{D}$, 故 $\sigma(\mathcal{C})\subseteq\mathcal{D}$, 即 $\mathcal{G}_1\subseteq\mathcal{G}_2$.

(2) 必要性显然. 充分性: 由 (1), $\{X^{-1}(B): B\in\mathcal{E}\}=\sigma(\{X^{-1}(C):C\in\mathcal{C}\})$, 而右端的生成元都在 $\mathcal{F}$ 中, 故右端包含于 $\mathcal{F}$.
\end{proof}

\begin{corollary}\label{cor:rv-criterion}
$X:\Omega\to\R$ 是随机变量, 当且仅当对每个 $x\in\R$, $\{X\le x\}\in\mathcal{F}$.
\end{corollary}

\begin{proof}
由命题~\ref{prop:borel-generator}, $\mathcal{B}(\R)=\sigma(\{(-\infty,x]: x\in\R\})$, 再由引理~\ref{lem:measurable-generator} (2) 即得.
\end{proof}

由同样的理由, 判别条件中的 $\{X\le x\}$ 也可以换成 $\{X<x\}$; 对广义随机变量, 判别条件同样是对每个 $x\in\R$ 有 $\{X\le x\}\in\mathcal{F}$.

\begin{proposition}[随机变量的运算]\label{prop:rv-operations}
\begin{enumerate}
\item 若 $X$ 是随机变量, $g:\R\to\R$ 是 Borel 可测函数 (即 $\mathcal{B}(\R)/\mathcal{B}(\R)$-可测), 则 $g(X)$ 是随机变量. 特别地, 连续函数是 Borel 可测的.
\item 若 $X,Y$ 是随机变量, $a\in\R$, 则 $aX$, $X+Y$, $XY$, $X\vee Y$, $X\wedge Y$ 都是随机变量.
\item 若 $X_1,X_2,\dots$ 是随机变量, 则 $\sup_n X_n$, $\inf_n X_n$, $\limsup_n X_n$, $\liminf_n X_n$ 都是广义随机变量, 且 $\{\omega: \lim_n X_n(\omega)\text{ 存在}\}\in\mathcal{F}$.
\end{enumerate}
\end{proposition}

\begin{proof}
(1) 对 $B\in\mathcal{B}(\R)$, $g^{-1}(B)\in\mathcal{B}(\R)$, 于是 $\{g(X)\in B\}=X^{-1}(g^{-1}(B))\in\mathcal{F}$. 若 $g$ 连续, 则开集的原像是开集; 由引理~\ref{lem:measurable-generator} (2) (以开集族为生成元), $g$ 是 Borel 可测的.

(2) 由 (1), $aX$ 是随机变量. 对 $X+Y$, 我们断言
\[
\{X+Y<x\}=\bigcup_{q\in\Q}\big(\{X<q\}\cap\{Y<x-q\}\big).
\]
若 $X(\omega)+Y(\omega)<x$, 则 $X(\omega)<x-Y(\omega)$, 取有理数 $q$ 满足 $X(\omega)<q<x-Y(\omega)$, 便有 $Y(\omega)<x-q$; 反之若 $X(\omega)<q$ 且 $Y(\omega)<x-q$, 相加即得 $X(\omega)+Y(\omega)<x$. 右端是可列个事件之并, 属于 $\mathcal{F}$. 由上面关于 $\{X<x\}$ 形式的判别法, $X+Y$ 是随机变量. 其次, 由恒等式 $XY=\frac{1}{4}\big((X+Y)^2-(X-Y)^2\big)$ 与 (1) ($x\mapsto x^2$ 连续), $XY$ 是随机变量. 最后 $\{X\vee Y\le x\}=\{X\le x\}\cap\{Y\le x\}$, $\{X\wedge Y\le x\}=\{X\le x\}\cup\{Y\le x\}$.

(3) $\{\sup_n X_n\le x\}=\bigcap_n\{X_n\le x\}\in\mathcal{F}$. 由于 $\inf_n X_n=-\sup_n(-X_n)$, 它也是广义随机变量. 再由
\[
\limsup_{n\to\infty}X_n=\inf_{m\ge1}\sup_{n\ge m}X_n,\qquad \liminf_{n\to\infty}X_n=\sup_{m\ge1}\inf_{n\ge m}X_n
\]
得到上、下极限的可测性. 最后, 极限存在 (允许为 $\pm\infty$) 当且仅当 $Y=\liminf_n X_n$ 与 $Z=\limsup_n X_n$ 相等. 而
\[
\{Y<Z\}=\bigcup_{q\in\Q}\big(\{Y<q\}\cap\{q<Z\}\big)\in\mathcal{F},
\]
由于总有 $Y\le Z$, $\{\lim_n X_n\text{ 存在}\}=\{Y<Z\}^c\in\mathcal{F}$.
\end{proof}

最简单的随机变量只取有限个值.

\begin{definition}[简单随机变量]\label{def:simple}
形如 $X=\sum_{i=1}^k a_i\1_{A_i}$ 的随机变量称为\textbf{简单随机变量}, 其中 $a_i\in\R$, $A_i\in\mathcal{F}$.
\end{definition}

下面的引理是 1.3 节构造期望、本节证明 Doob-Dynkin 引理的共同基础: 每个非负随机变量都是简单随机变量的单调增极限.

\begin{lemma}[简单函数逼近]\label{lem:simple-approx}
设 $X$ 是非负广义随机变量. 对 $n\ge1$ 令
\[
X_n=\sum_{k=0}^{n2^n-1}\frac{k}{2^n}\,\1_{\{k2^{-n}\le X<(k+1)2^{-n}\}}+n\,\1_{\{X\ge n\}}.
\]
则每个 $X_n$ 是非负简单随机变量, $0\le X_1\le X_2\le\cdots$, 且对每个 $\omega$, $X_n(\omega)\uparrow X(\omega)$. 若 $X$ 有界, 则收敛是一致的.
\end{lemma}

\begin{proof}
等价地, $X_n=\min\big\{n,\ 2^{-n}\lfloor 2^nX\rfloor\big\}$ (约定 $\lfloor+\infty\rfloor=+\infty$): 当 $X<n$ 时, $X_n$ 是把 $X$ 向下取整到 $2^{-n}$ 的整数倍; 当 $X\ge n$ 时 $X_n=n$. 各示性函数对应的集合属于 $\mathcal{F}$, 所以 $X_n$ 是简单随机变量, 且 $0\le X_n\le X$.

\emph{单调性.} 若 $X(\omega)\ge n+1$, 则 $X_{n+1}(\omega)=n+1>n=X_n(\omega)$. 若 $n\le X(\omega)<n+1$, 则 $X_{n+1}(\omega)=2^{-(n+1)}\lfloor2^{n+1}X(\omega)\rfloor\ge 2^{-(n+1)}\cdot 2^{n+1}n=n=X_n(\omega)$. 若 $X(\omega)<n$, 则 $X_n(\omega)=\frac{k}{2^n}=\frac{2k}{2^{n+1}}$, 其中 $\frac{k}{2^n}\le X(\omega)$, 而 $X_{n+1}(\omega)$ 是不超过 $X(\omega)$ 的最大的 $2^{-(n+1)}$ 的整数倍, 故 $X_{n+1}(\omega)\ge\frac{2k}{2^{n+1}}=X_n(\omega)$.

\emph{收敛性.} 若 $X(\omega)<\infty$, 则当 $n>X(\omega)$ 时, $0\le X(\omega)-X_n(\omega)<2^{-n}\to0$. 若 $X(\omega)=\infty$, 则 $X_n(\omega)=n\to\infty$. 若 $X\le M$, 则当 $n>M$ 时, 对所有 $\omega$ 同时有 $0\le X-X_n<2^{-n}$, 即一致收敛.
\end{proof}

\begin{figure}[htbp]
\centering
\begin{tikzpicture}
\begin{axis}[SAaxis, width=10cm, height=7cm, xmin=0, xmax=1.05, ymin=0, ymax=2.05,
  xtick={0,1}, ytick={0.5,1,1.5,2}, yticklabels={$\frac{1}{2}$,$1$,$\frac{3}{2}$,$2$},
  xlabel={$\omega$}, legend style={draw=none,fill=none,font=\small,at={(0.98,0.98)},anchor=north east,legend columns=1}]
  \foreach \k in {1,...,7} { \addplot[SADark!15,domain=0:1,forget plot] {\k/4}; }
  \addplot[SABlue,very thick,domain=0:1,samples=200] {1.05+0.8*sin(deg(5*x))};
  \addlegendentry{$X(\omega)$}
  \addplot[SARed,thick,domain=0:1,samples=800,const plot] {min(2,floor(4*(1.05+0.8*sin(deg(5*x))))/4)};
  \addlegendentry{$X_2(\omega)$}
  \addplot[SAGreenText,thick,domain=0:1,samples=800,const plot,dashed] {min(1,floor(2*(1.05+0.8*sin(deg(5*x))))/2)};
  \addlegendentry{$X_1(\omega)$}
\end{axis}
\end{tikzpicture}
\caption{简单函数逼近 (取 $\Omega=[0,1]$): $X_n$ 把 $X$ 的\emph{值域}按 $2^{-n}$ 分层, 向下取整. 层越细, 阶梯越贴近 $X$}
\label{fig:simple-approx}
\end{figure}

% ---------------------------------------------------------------------
\subsection{随机变量生成的 \texorpdfstring{$\sigma$}{sigma}-代数}

\begin{definition}[随机变量生成的 $\sigma$-代数]\label{def:sigma-X}
设 $X$ 是随机变量. 称
\[
\sigma(X)=\{X^{-1}(B): B\in\mathcal{B}(\R)\}
\]
为 $X$ \textbf{生成的 $\sigma$-代数}. 对一族随机变量 $\{X_i\}_{i\in I}$, 令 $\sigma(X_i: i\in I)=\sigma\big(\bigcup_{i\in I}\sigma(X_i)\big)$.
\end{definition}

由引理~\ref{lem:measurable-generator} 的证明, $\sigma(X)$ 是 $\sigma$-代数, 且 $\sigma(X)=\sigma(\{\{X\le x\}: x\in\R\})$. 它是使 $X$ 可测的最小 $\sigma$-代数: 若 $X$ 是 $\mathcal{G}/\mathcal{B}(\R)$-可测的, 则 $\sigma(X)\subseteq\mathcal{G}$.

\begin{intuition}
$\sigma(X)$ 是 ``观察到 $X$ 的值'' 所获得的信息: 知道 $X(\omega)$ 后, 能够判断是否发生的事件恰好是 $\{X\in B\}$ 这样的事件. 随机过程理论中的信息流 $\mathcal{F}_t=\sigma(X_s: s\le t)$ 正是 ``观察到时刻 $t$ 为止的全部历史'' 所得的信息.
\end{intuition}

\begin{example}\label{ex:sigma-X}
\begin{enumerate}
\item 若 $X\equiv c$ 是常数, 则 $X^{-1}(B)$ 只能是 $\varnothing$ 或 $\Omega$, 故 $\sigma(X)=\{\varnothing,\Omega\}$: 常数不提供任何信息.
\item 若 $X=\1_A$, 则 $\sigma(X)=\{\varnothing,A,A^c,\Omega\}$.
\item 若 $X$ 只取互不相同的值 $x_1,\dots,x_k$, 令 $D_i=\{X=x_i\}$, 则 $\{D_1,\dots,D_k\}$ 是 $\Omega$ 的分割, $\sigma(X)$ 就是例~\ref{ex:sigma-examples} (4) 中由该分割生成的 $\sigma$-代数. 例如掷一颗骰子, $X$ 为点数的奇偶性 ($1$ 表示奇数, $0$ 表示偶数), 则 $\sigma(X)=\{\varnothing,\{1,3,5\},\{2,4,6\},\Omega\}$.
\item 设 $\Omega=[-1,1]$, $\P$ 为均匀分布, $X(\omega)=\omega$. 则 $\sigma(X^2)\subsetneq\sigma(X)$: $\sigma(X^2)$ 中的每个集合 $\{X^2\in B\}$ 都关于原点对称, 而 $\{X>0\}\in\sigma(X)$ 不对称. 直观地, 由 $X^2$ 无法得知 $X$ 的符号.
\end{enumerate}
\end{example}

``$Y$ 由 $X$ 的信息决定'' 有两种表述方式: 一是 $Y$ 关于 $\sigma(X)$ 可测, 二是 $Y$ 是 $X$ 的函数. 下面的引理说明它们等价. 它在第~2~章用条件期望 $\E[X\,|\,Y]$ 定义条件期望函数时起关键作用.

\begin{lemma}[Doob-Dynkin 引理]\label{lem:doob-dynkin}
设 $X,Y$ 是随机变量. 则 $Y$ 是 $\sigma(X)$-可测的, 当且仅当存在 Borel 可测函数 $g:\R\to\R$, 使得 $Y=g(X)$.
\end{lemma}

\begin{proof}
``当'': 对 $B\in\mathcal{B}(\R)$, $\{g(X)\in B\}=\{X\in g^{-1}(B)\}\in\sigma(X)$.

``仅当'': 按 $Y$ 的复杂程度分四步.

\emph{第一步: $Y=\1_A$, $A\in\sigma(X)$.} 由 $\sigma(X)$ 的定义, 存在 $B\in\mathcal{B}(\R)$ 使 $A=\{X\in B\}$. 取 $g=\1_B$, 则 $g(X)=\1_{\{X\in B\}}=\1_A=Y$.

\emph{第二步: $Y=\sum_{i=1}^k a_i\1_{A_i}$, $A_i\in\sigma(X)$.} 由第一步, $\1_{A_i}=\1_{B_i}(X)$, 取 $g=\sum_i a_i\1_{B_i}$ 即可.

\emph{第三步: $Y\ge0$.} 由简单函数逼近 (引理~\ref{lem:simple-approx}), 存在 $\sigma(X)$-可测的非负简单随机变量 $Y_n\uparrow Y$ (每个 $Y_n$ 中出现的集合都形如 $\{a\le Y<b\}$ 或 $\{Y\ge n\}$, 属于 $\sigma(X)$). 由第二步, $Y_n=g_n(X)$, $g_n$ 为 Borel 函数. 令
\[
g(x)=\begin{cases}\limsup_{n\to\infty}g_n(x), & \text{若此上极限有限},\\ 0, & \text{否则}.\end{cases}
\]
由命题~\ref{prop:rv-operations} (3) (取 $\Omega=\R$, $\mathcal{F}=\mathcal{B}(\R)$), $\limsup_n g_n$ 是 Borel 可测的, 进而 $g$ 是 Borel 可测的. 对每个 $\omega$, 取 $x=X(\omega)$, 则 $g_n(x)=Y_n(\omega)\uparrow Y(\omega)<\infty$, 所以 $\limsup_n g_n(x)=Y(\omega)$ 有限, $g(X(\omega))=Y(\omega)$.

\emph{第四步: 一般情形.} 记 $Y^+=Y\vee0$, $Y^-=(-Y)\vee0$, 它们非负且 $\sigma(X)$-可测, $Y=Y^+-Y^-$. 由第三步, $Y^\pm=g_\pm(X)$, 取 $g=g_+-g_-$ 即可.
\end{proof}

\begin{remark}
对多个随机变量, 结论同样成立: $Y$ 是 $\sigma(X_1,\dots,X_n)$-可测的, 当且仅当 $Y=g(X_1,\dots,X_n)$, 其中 $g:\R^n\to\R$ 是 Borel 可测函数 (见 1.5 节), 证明完全相同. 注意第三步中的技巧: 函数 $g_n$ 只在 $X$ 的值域上被确定, 在值域以外它们可能毫无规律, 因此要用上极限并截断来定义 $g$.
\end{remark}

利用 $\sigma(X)$, 随机变量的独立性可以直接借用 $\sigma$-代数的独立性来定义.

\begin{definition}[随机变量的独立性]\label{def:indep-rv}
称随机变量族 $\{X_i\}_{i\in I}$ 相互\textbf{独立}, 如果 $\{\sigma(X_i)\}_{i\in I}$ 相互独立. 称随机变量 $X$ 与 $\sigma$-代数 $\mathcal{G}$ 独立, 如果 $\sigma(X)$ 与 $\mathcal{G}$ 独立. 独立且具有相同分布 (见下一小节) 的随机变量序列称为\textbf{独立同分布} ($\iid$) 序列.
\end{definition}

\begin{theorem}[独立性的分布函数判别法]\label{thm:indep-cdf}
随机变量 $X_1,\dots,X_n$ 相互独立, 当且仅当对任意 $x_1,\dots,x_n\in\R$,
\begin{equation}\label{eq:indep-cdf}
\P(X_1\le x_1,\dots,X_n\le x_n)=\P(X_1\le x_1)\cdots\P(X_n\le x_n).
\end{equation}
\end{theorem}

\begin{proof}
必要性由定义直接得到. 充分性: 令 $\mathcal{P}_i=\{\{X_i\le x\}: x\in\R\}\cup\{\Omega\}$. 由于 $\{X_i\le x\}\cap\{X_i\le y\}=\{X_i\le x\wedge y\}$, $\mathcal{P}_i$ 是包含 $\Omega$ 的 $\pi$-系, 且由引理~\ref{lem:measurable-generator} (1) 与命题~\ref{prop:borel-generator}, $\sigma(\mathcal{P}_i)=\sigma(X_i)$. 要用独立性的 $\pi$-系判别准则 (定理~\ref{thm:indep-pi}), 还需验证当某些 $A_i$ 取为 $\Omega$ 时乘积公式也成立. 对这些指标 $i$, 令 $x_i=m\to\infty$: 此时 $\{X_i\le m\}\uparrow\Omega$, 由概率的连续性 (定理~\ref{thm:continuity}), \eqref{eq:indep-cdf} 两端分别收敛到对应指标取 $\Omega$ 的形式. 于是乘积公式对所有 $A_i\in\mathcal{P}_i$ 成立, 由定理~\ref{thm:indep-pi}, $\sigma(X_1),\dots,\sigma(X_n)$ 独立.
\end{proof}

\begin{proposition}\label{prop:indep-functions}
若 $X_1,\dots,X_n$ 相互独立, $g_1,\dots,g_n$ 是 Borel 函数, 则 $g_1(X_1),\dots,g_n(X_n)$ 相互独立. 更一般地, 由不交的几组独立随机变量分别构造出的随机变量相互独立 (例如 $X_1+X_2$ 与 $X_3X_4$).
\end{proposition}

\begin{proof}
由命题~\ref{prop:rv-operations} 的证明 (1), $\sigma(g_i(X_i))\subseteq\sigma(X_i)$. 独立 $\sigma$-代数的子 $\sigma$-代数仍然独立 (定义~\ref{def:indep-sigma} 中的条件只会变少). 第二个论断由分组原理 (推论~\ref{cor:grouping}) 与 Doob-Dynkin 引理之后的注 (多变量函数 $g(X_1,X_2)$ 关于 $\sigma(X_1,X_2)$ 可测, 见 1.5 节) 得到.
\end{proof}

% ---------------------------------------------------------------------
\subsection{分布与分布函数}

随机变量把 $\Omega$ 上的概率 ``搬运'' 到了实直线上.

\begin{definition}[分布]\label{def:distribution}
设 $X$ 是随机变量. 称 $\R$ 上的集合函数
\[
\mu_X(B)=\P(X\in B)=\P\big(X^{-1}(B)\big),\qquad B\in\mathcal{B}(\R)
\]
为 $X$ 的\textbf{分布} (或概率律, 也称为 $\P$ 在 $X$ 下的\textbf{推前测度}, 记作 $\P\circ X^{-1}$). 若 $\mu_X=\mu_Y$, 称 $X$ 与 $Y$ \textbf{同分布}, 记作 $X\overset{d}{=}Y$; 若 $\mu_X=\mu$, 记作 $X\sim\mu$.
\end{definition}

\begin{proposition}
$\mu_X$ 是 $(\R,\mathcal{B}(\R))$ 上的概率测度.
\end{proposition}

\begin{proof}
$\mu_X(\R)=\P(\Omega)=1$. 若 $B_1,B_2,\dots$ 是两两不交的 Borel 集, 则原像 $X^{-1}(B_n)$ 两两不交 (一个 $\omega$ 不能同时使 $X(\omega)$ 落在两个不交的集合中), 且 $X^{-1}(\bigcup_n B_n)=\bigcup_n X^{-1}(B_n)$. 由 $\P$ 的可列可加性,
$\mu_X\big(\bigcup_n B_n\big)=\sum_n\P(X^{-1}(B_n))=\sum_n\mu_X(B_n)$.
\end{proof}

\begin{intuition}
分布只记录 ``$X$ 落在各处的可能性有多大'', 忘记了 $X$ 是定义在哪个概率空间上、如何依赖于 $\omega$ 的. 两个随机变量可以同分布却毫无关系 (比如两次独立抛硬币的结果). 计算单个随机变量的概率和期望时, 只需知道分布; 但涉及多个随机变量之间的关系时, 必须回到 $\Omega$ 上.
\end{intuition}

分布是定义在全体 Borel 集上的集合函数, 用起来不方便. 由 $\pi$-$\lambda$ 定理, 它其实由一个普通的实函数决定.

\begin{definition}[分布函数]\label{def:cdf}
随机变量 $X$ 的 (累积) \textbf{分布函数}定义为 $F_X(x)=\P(X\le x)=\mu_X((-\infty,x])$, $x\in\R$.
\end{definition}

\begin{theorem}[分布函数的性质]\label{thm:cdf-properties}
设 $F=F_X$. 则
\begin{enumerate}
\item $F$ 单调不减;
\item $F$ 右连续: $\lim_{y\downarrow x}F(y)=F(x)$;
\item $\lim_{x\to-\infty}F(x)=0$, $\lim_{x\to+\infty}F(x)=1$;
\item 左极限 $F(x-)=\lim_{y\uparrow x}F(y)=\P(X<x)$, 从而 $\P(X=x)=F(x)-F(x-)$.
\end{enumerate}
\end{theorem}

\begin{proof}
(1) 若 $x\le y$, 则 $\{X\le x\}\subseteq\{X\le y\}$, 由单调性 $F(x)\le F(y)$.

由于 $F$ 单调, 其单侧极限存在, 只需沿一个特定的单调数列计算.

(2) 取 $y_n\downarrow x$, 则 $\{X\le y_n\}\downarrow\{X\le x\}$, 由上连续性 (定理~\ref{thm:continuity}) $F(y_n)\to F(x)$.

(3) 当 $n\to\infty$ 时, $\{X\le -n\}\downarrow\varnothing$, $\{X\le n\}\uparrow\Omega$ (因为 $X$ 取有限实值). 由连续性, $F(-n)\to0$, $F(n)\to1$.

(4) 取 $y_n\uparrow x$ 且 $y_n<x$, 则 $\{X\le y_n\}\uparrow\{X<x\}$, 由下连续性 $F(y_n)\to\P(X<x)$. 再由 $\{X=x\}=\{X\le x\}\setminus\{X<x\}$ 和命题~\ref{prop:prob-basic} (4) 得到最后的等式.
\end{proof}

\begin{theorem}[分布函数决定分布]\label{thm:cdf-determines}
若 $F_X=F_Y$, 则 $\mu_X=\mu_Y$, 即 $X\overset{d}{=}Y$.
\end{theorem}

\begin{proof}
$\mu_X$ 与 $\mu_Y$ 在 $\pi$-系 $\{(-\infty,x]: x\in\R\}$ 上相等 (例~\ref{ex:pi-lambda} (1)), 而该 $\pi$-系生成 $\mathcal{B}(\R)$ (命题~\ref{prop:borel-generator}). 由概率测度的唯一性 (推论~\ref{cor:uniqueness}), $\mu_X=\mu_Y$.
\end{proof}

反过来, 满足定理~\ref{thm:cdf-properties} (1)-(3) 的函数都是某个随机变量的分布函数. 下面给出一个构造性的证明, 它同时是计算机模拟中 ``逆变换抽样'' 的理论基础.

\begin{theorem}[分位数变换]\label{thm:quantile}
设 $F:\R\to[0,1]$ 满足定理~\ref{thm:cdf-properties} 中的 (1)(2)(3). 对 $u\in(0,1)$ 定义广义逆
\[
F^{-1}(u)=\inf\{x\in\R: F(x)\ge u\}.
\]
设 $U$ 服从 $(0,1)$ 上的均匀分布 (例~\ref{ex:basic-spaces} (1)), 则 $X=F^{-1}(U)$ 是随机变量, 且其分布函数为 $F$.
\end{theorem}

\begin{proof}
固定 $u\in(0,1)$, 记 $S_u=\{x: F(x)\ge u\}$.
\begin{itemize}
\item $S_u$ 非空且有下界: 由 $F(x)\to1>u$ ($x\to+\infty$), 存在 $x$ 使 $F(x)\ge u$; 由 $F(x)\to0<u$ ($x\to-\infty$), 当 $x$ 充分小时 $x\notin S_u$, 而由单调性, 比这个 $x$ 更小的点也都不属于 $S_u$. 故 $F^{-1}(u)$ 是有限实数.
\item $S_u=[F^{-1}(u),\infty)$: 由单调性, 若 $x\in S_u$ 且 $y>x$, 则 $y\in S_u$, 所以 $S_u$ 是以 $F^{-1}(u)$ 为左端点的区间. 取 $x_n\in S_u$, $x_n\downarrow F^{-1}(u)$, 由右连续性 $F(F^{-1}(u))=\lim_n F(x_n)\ge u$, 所以左端点属于 $S_u$.
\end{itemize}
由此得到关键的等价关系: 对 $u\in(0,1)$, $x\in\R$,
\begin{equation}\label{eq:galois}
F^{-1}(u)\le x\iff x\in S_u\iff u\le F(x).
\end{equation}
于是 $\{X\le x\}=\{U\le F(x)\}$ 是事件, $X$ 是随机变量, 且
$\P(X\le x)=\P(U\le F(x))=F(x)$ (均匀分布满足 $\P(U\le t)=t$, $t\in[0,1]$).
\end{proof}

\begin{figure}[htbp]
\centering
\begin{tikzpicture}[xscale=1.6,yscale=3.2]
  \draw[->,SADark!80] (-0.4,0) -- (3.6,0) node[right] {$x$};
  \draw[->,SADark!80] (0,-0.05) -- (0,1.12) node[above] {$F(x)$};
  \draw[SADark!30,dashed] (0,1) -- (3.5,1);
  \node[left,font=\small] at (0,1) {$1$};
  % F pieces
  \draw[SABlue,very thick] (-0.4,0) -- (0,0) -- (1,0.3);
  \draw[SABlue,very thick] (1,0.55) -- (2,0.55) -- (3,1) -- (3.5,1);
  \draw[SABlue,fill=white] (1,0.3) circle (0.9pt);
  \fill[SABlue] (1,0.55) circle (0.9pt);
  % u1 in jump gap
  \draw[SARed,dashed] (0,0.42) -- (1,0.42) -- (1,0);
  \node[left,SARed,font=\small] at (0,0.42) {$u_1$};
  % u2 flat level
  \node[left,SARed,font=\small] at (0,0.56) {$u_2$};
  \draw[SARed,dashed] (0,0.55) -- (1,0.55);
  \node[below,SARed,font=\small,align=center] at (1,0) {$F^{-1}(u_1)$\\$=F^{-1}(u_2)$};
  % u3
  \draw[SAGreenText,dashed] (0,0.8) -- (2.556,0.8) -- (2.556,0);
  \node[left,SAGreenText,font=\small] at (0,0.8) {$u_3$};
  \node[below,SAGreenText,font=\small] at (2.556,0) {$F^{-1}(u_3)$};
\end{tikzpicture}
\caption{广义逆 $F^{-1}(u)=\inf\{x: F(x)\ge u\}$. 落在跳跃缺口中的 $u_1$ 和等于平台高度的 $u_2=0.55$ 都被送到跳跃点; $F$ 的平台区间 $(1,2)$ 不在 $F^{-1}$ 的值域中 (该区间概率为 $0$)}
\label{fig:quantile}
\end{figure}

\begin{remark}
结合定理~\ref{thm:cdf-determines} 与定理~\ref{thm:quantile}: $\R$ 上的概率测度与满足 (1)(2)(3) 的函数一一对应. 进一步, 在乘积概率空间上可以构造任意给定分布的独立随机变量序列 (见~\cite{durrett} 2.1 节). 今后我们常说 ``设 $X_1,X_2,\dots$ 是 $\iid$ 的 $N(0,1)$ 随机变量'', 其存在性即由此保证.
\end{remark}

最常见的分布有两类.

\begin{definition}[离散型与绝对连续型]\label{def:discrete-continuous}
\begin{enumerate}
\item 若存在有限或可列集 $S\subseteq\R$ 使 $\P(X\in S)=1$, 称 $X$ 为\textbf{离散型}随机变量, $p(x)=\P(X=x)$ ($x\in S$) 称为其\textbf{概率质量函数}. 此时 $\mu_X(B)=\sum_{x\in B\cap S}p(x)$, $F_X$ 是阶梯函数.
\item 若存在非负 Borel 函数 $f$ 使得对所有 $x\in\R$,
\[
F_X(x)=\int_{-\infty}^x f(t)\dd t,
\]
称 $X$ 为\textbf{绝对连续型}随机变量, $f=f_X$ 称为其\textbf{概率密度函数}. 此时 $\int_{\R}f(t)\dd t=1$.
\end{enumerate}
\end{definition}

这里的积分是关于 Lebesgue 测度的积分 (1.3 节给出定义); $f$ 分段连续时, 该积分与 Riemann (反常) 积分一致. 绝对连续型随机变量具有两条基本性质.
\begin{itemize}
\item 对每个 Borel 集 $B$, $\P(X\in B)=\int_B f(t)\dd t$. 事实上, $\nu(B)=\int_B f(t)\dd t$ 是 $\mathcal{B}(\R)$ 上的概率测度 (可列可加性由 1.3 节的单调收敛定理~\ref{thm:mct} 得到), 它与 $\mu_X$ 在 $\{(-\infty,x]\}$ 上一致, 由推论~\ref{cor:uniqueness} 二者相等.
\item $F_X$ 连续, 从而对每个 $x$, $\P(X=x)=0$; 在 $f$ 的连续点处 $F_X'(x)=f(x)$ (微积分基本定理).
\end{itemize}

\begin{remark}
条件 ``$\mu_X$ 有密度'' 可以用测度的语言刻画: \textbf{Radon-Nikodym 定理} 断言, $\mu_X$ 存在密度当且仅当 $\mu_X$ 关于 Lebesgue 测度 $\lambda$ \emph{绝对连续}, 即 $\lambda(B)=0\Rightarrow\mu_X(B)=0$; 密度 $f$ 记作 $\frac{\dd\mu_X}{\dd\lambda}$ (证明见~\cite{folland} 3.2 节). 一般的分布既不是离散型也不是绝对连续型: 它可以是两者的混合 (见下例), 甚至可以是连续但没有密度的 (如 Cantor 分布). 第~2~章与第~5~章会再次用到 Radon-Nikodym 定理.
\end{remark}

\begin{example}\label{ex:cdf-types}
\begin{enumerate}
\item \textbf{离散型.} 抛三次均匀硬币, $X$ 为正面次数, 则 $X\sim\mathrm{B}(3,\frac{1}{2})$, $p(k)=\binom{3}{k}\frac{1}{8}$, $F_X$ 在 $0,1,2,3$ 处分别跳跃 $\frac{1}{8},\frac{3}{8},\frac{3}{8},\frac{1}{8}$.
\item \textbf{绝对连续型.} 参数为 $1$ 的指数分布: $f(x)=e^{-x}\1_{\{x\ge0\}}$, $F(x)=(1-e^{-x})\1_{\{x\ge0\}}$.
\item \textbf{混合型.} 设 $T$ 服从参数为 $1$ 的指数分布 (例如一个元件的寿命), 而我们只观察到时刻 $1$ 为止, 令 $X=T\wedge1$. 则对 $0\le x<1$, $F_X(x)=1-e^{-x}$; 对 $x\ge1$, $F_X(x)=1$. $F_X$ 在 $x=1$ 处有大小为 $e^{-1}$ 的跳跃, 在 $(0,1)$ 上连续增长. $X$ 既不是离散型也不是绝对连续型.
\end{enumerate}
\end{example}

\begin{figure}[htbp]
\centering
\begin{tikzpicture}
\begin{scope}
  \draw[->,SADark!80] (-0.5,0) -- (3.9,0) node[right,font=\small] {$x$};
  \draw[->,SADark!80] (0,0) -- (0,2.4);
  \draw[SADark!30,dashed] (0,2) -- (3.7,2); \node[left,font=\small] at (0,2) {$1$};
  \foreach \a/\b/\h in {-0.5/0/0, 0/0.8/0.25, 0.8/1.6/1, 1.6/2.4/1.75, 2.4/3.7/2} { \draw[SABlue,very thick] (\a,\h) -- (\b,\h); }
  \foreach \x/\h in {0/0, 0.8/0.25, 1.6/1, 2.4/1.75} { \draw[SABlue,fill=white] (\x,\h) circle (1.6pt); }
  \foreach \x/\h in {0/0.25, 0.8/1, 1.6/1.75, 2.4/2} { \fill[SABlue] (\x,\h) circle (1.6pt); }
  \foreach \x/\l in {0/0,0.8/1,1.6/2,2.4/3} { \node[below,font=\small] at (\x,0) {$\l$}; }
  \node[font=\small] at (1.6,-0.75) {(a) 离散型};
\end{scope}
\begin{scope}[xshift=5cm]
  \draw[->,SADark!80] (-0.5,0) -- (3.9,0) node[right,font=\small] {$x$};
  \draw[->,SADark!80] (0,0) -- (0,2.4);
  \draw[SADark!30,dashed] (0,2) -- (3.7,2); \node[left,font=\small] at (0,2) {$1$};
  \draw[SABlue,very thick] (-0.5,0) -- (0,0);
  \draw[SABlue,very thick,domain=0:3.7,samples=60] plot (\x,{2*(1-exp(-1.1*\x))});
  \node[below,font=\small] at (0,0) {$0$};
  \node[font=\small] at (1.6,-0.75) {(b) 绝对连续型};
\end{scope}
\begin{scope}[xshift=10cm]
  \draw[->,SADark!80] (-0.5,0) -- (3.9,0) node[right,font=\small] {$x$};
  \draw[->,SADark!80] (0,0) -- (0,2.4);
  \draw[SADark!30,dashed] (0,2) -- (3.7,2); \node[left,font=\small] at (0,2) {$1$};
  \draw[SABlue,very thick] (-0.5,0) -- (0,0);
  \draw[SABlue,very thick,domain=0:1.6,samples=40] plot (\x,{2*(1-exp(-\x/1.6))});
  \draw[SABlue,very thick] (1.6,2) -- (3.7,2);
  \draw[SABlue,fill=white] (1.6,{2*(1-exp(-1))}) circle (1.6pt);
  \fill[SABlue] (1.6,2) circle (1.6pt);
  \draw[SARed,decorate,decoration={brace,amplitude=3pt,mirror}] (1.75,{2*(1-exp(-1))}) -- (1.75,2) node[midway,right,font=\small,xshift=2pt] {$e^{-1}$};
  \node[below,font=\small] at (0,0) {$0$};
  \node[below,font=\small] at (1.6,0) {$1$};
  \node[font=\small] at (1.6,-0.75) {(c) 混合型};
\end{scope}
\end{tikzpicture}
\caption{例~\ref{ex:cdf-types} 中三类分布函数. 实心点表示函数值, 空心点表示左极限; 跳跃的高度等于该点的概率}
\label{fig:cdf-types}
\end{figure}

表~\ref{tab:distributions} 列出了本书常用的分布, 其期望与方差将在 1.3 节计算.

\begin{table}[htbp]
\centering
\caption{常用分布}
\label{tab:distributions}
\renewcommand{\arraystretch}{1.5}
\begin{tabular}{lll}
\toprule
分布 & 概率质量函数 / 密度 & 典型背景 \\
\midrule
Bernoulli $\mathrm{B}(1,p)$ & $\P(X=1)=p$, $\P(X=0)=1-p$ & 一次试验成功与否 \\
二项 $\mathrm{B}(n,p)$ & $\binom{n}{k}p^k(1-p)^{n-k}$, $k=0,\dots,n$ & $n$ 次独立试验的成功次数 \\
几何 $\mathrm{Geo}(p)$ & $(1-p)^{k-1}p$, $k=1,2,\dots$ & 首次成功所需的试验次数 \\
Poisson $\mathrm{Poi}(\lambda)$ & $e^{-\lambda}\frac{\lambda^k}{k!}$, $k=0,1,\dots$ & 稀有事件的发生次数 \\
\midrule
均匀 $\mathrm{U}(a,b)$ & $\frac{1}{b-a}\1_{(a,b)}(x)$ & 随机取点 \\
指数 $\mathrm{Exp}(\lambda)$ & $\lambda e^{-\lambda x}\1_{\{x\ge0\}}$ & 无记忆的等待时间 \\
正态 $N(\mu,\sigma^2)$ & $\frac{1}{\sqrt{2\pi}\sigma}e^{-\frac{(x-\mu)^2}{2\sigma^2}}$ & 大量微小独立因素之和 \\
\bottomrule
\end{tabular}
\end{table}

% ---------------------------------------------------------------------
\subsection{正态分布}

正态分布在随机分析中处于中心地位: 布朗运动的每个增量都服从正态分布. 本小节给出正态分布的基本性质和精确的尾概率估计; 正态分布的矩和特征函数将在 1.3 节计算, 多元正态分布在 1.5 节讨论.

\begin{definition}[正态分布]\label{def:normal}
密度为
\[
\varphi(x)=\frac{1}{\sqrt{2\pi}}e^{-x^2/2},\qquad x\in\R
\]
的分布称为\textbf{标准正态分布}, 记为 $N(0,1)$, 其分布函数记为 $\Phi(x)=\int_{-\infty}^x\varphi(t)\dd t$. 对 $\mu\in\R$, $\sigma>0$, 若 $X=\mu+\sigma Z$, 其中 $Z\sim N(0,1)$, 称 $X$ 服从\textbf{正态分布} $N(\mu,\sigma^2)$. 约定 $N(\mu,0)$ 表示集中在点 $\mu$ 的分布 (即 $X\equiv\mu$), 称为退化正态分布.
\end{definition}

首先验证 $\varphi$ 确实是概率密度. 记 $I=\int_{\R}e^{-x^2/2}\dd x$. 利用二重积分的极坐标变换 ($x=r\cos\theta$, $y=r\sin\theta$, $\dd x\dd y=r\dd r\dd\theta$; 被积函数非负, 累次积分的交换由 1.5 节的 Fubini 定理保证),
\begin{align*}
I^2&=\int_{\R}\int_{\R}e^{-(x^2+y^2)/2}\dd x\dd y=\int_0^{2\pi}\int_0^\infty e^{-r^2/2}\,r\dd r\dd\theta\\
&=2\pi\Big[-e^{-r^2/2}\Big]_0^\infty=2\pi.
\end{align*}
所以 $I=\sqrt{2\pi}$, $\int_{\R}\varphi=1$.

\begin{proposition}[正态分布的基本性质]\label{prop:normal-basic}
\begin{enumerate}
\item 对称性: $\varphi(-x)=\varphi(x)$, $\Phi(-x)=1-\Phi(x)$; 若 $Z\sim N(0,1)$, 则 $-Z\sim N(0,1)$.
\item 若 $X\sim N(\mu,\sigma^2)$, $\sigma>0$, 则
\[
F_X(x)=\Phi\Big(\frac{x-\mu}{\sigma}\Big),\qquad f_X(x)=\frac{1}{\sqrt{2\pi}\sigma}\exp\Big(-\frac{(x-\mu)^2}{2\sigma^2}\Big).
\]
反之, 若 $X\sim N(\mu,\sigma^2)$, 则标准化变量 $\frac{X-\mu}{\sigma}\sim N(0,1)$.
\item 线性变换: 若 $X\sim N(\mu,\sigma^2)$, $a,b\in\R$, 则 $aX+b\sim N(a\mu+b,a^2\sigma^2)$.
\end{enumerate}
\end{proposition}

\begin{proof}
(1) 第一式显然. 作变量替换 $t=-s$,
\[
\Phi(-x)=\int_{-\infty}^{-x}\varphi(t)\dd t=\int_x^{\infty}\varphi(s)\dd s=1-\Phi(x).
\]
于是 $\P(-Z\le x)=\P(Z\ge-x)=1-\Phi(-x)=\Phi(x)$ (这里用到 $\P(Z=-x)=0$).

(2) $F_X(x)=\P(\mu+\sigma Z\le x)=\P\big(Z\le\frac{x-\mu}{\sigma}\big)=\Phi\big(\frac{x-\mu}{\sigma}\big)$. 由于 $\Phi$ 可微, 由链式法则得到密度. 标准化变量 $\frac{X-\mu}{\sigma}$ 的分布函数为 $\P(X\le\mu+\sigma x)=\Phi(x)$, 由定理~\ref{thm:cdf-determines} 即得.

(3) 写 $X=\mu+\sigma Z$. 若 $a\ge0$, 则 $aX+b=(a\mu+b)+(a\sigma)Z$, 由定义 (当 $a=0$ 时为退化情形) 即得. 若 $a<0$, 则 $aX+b=(a\mu+b)+\abs{a}\sigma(-Z)$, 由 (1), $-Z\sim N(0,1)$, 同样得证.
\end{proof}

参数 $\mu$ 与 $\sigma^2$ 分别是 $X$ 的期望与方差, 这一点将在 1.3 节验证. 正态分布的尾部衰减极快, 下面的估计在后文 (例如布朗运动的轨道性质、重对数律) 中反复使用.

\begin{theorem}[正态尾概率估计]\label{thm:mills}
对任意 $x>0$,
\[
\frac{x}{1+x^2}\,\varphi(x)\ \le\ 1-\Phi(x)\ \le\ \frac{\varphi(x)}{x}.
\]
特别地, 当 $x\to\infty$ 时, $1-\Phi(x)\sim\frac{\varphi(x)}{x}$, 即二者之比趋于 $1$.
\end{theorem}

\begin{proof}
关键是 $\varphi'(t)=-t\varphi(t)$.

\emph{上界.} 当 $t\ge x>0$ 时 $\frac{t}{x}\ge1$, 所以
\[
1-\Phi(x)=\int_x^\infty\varphi(t)\dd t\le\int_x^\infty\frac{t}{x}\varphi(t)\dd t=\frac{1}{x}\Big[-\varphi(t)\Big]_x^\infty=\frac{\varphi(x)}{x}.
\]

\emph{下界.} 直接求导可得
$\frac{\dd}{\dd t}\Big(-\frac{\varphi(t)}{t}\Big)=\frac{t\varphi(t)\cdot t+\varphi(t)}{t^2}=\Big(1+\frac{1}{t^2}\Big)\varphi(t)$, 因此
\[
\frac{\varphi(x)}{x}=\int_x^\infty\Big(1+\frac{1}{t^2}\Big)\varphi(t)\dd t\le\Big(1+\frac{1}{x^2}\Big)\int_x^\infty\varphi(t)\dd t=\Big(1+\frac{1}{x^2}\Big)\big(1-\Phi(x)\big).
\]
两边乘以 $\frac{x^2}{1+x^2}$ 即得下界. 最后, 上下界之比为 $\frac{x^2}{1+x^2}\to1$, 由夹逼得到渐近等价.
\end{proof}

\begin{intuition}
尾概率 $\P(Z>x)$ 大约是 ``密度在 $x$ 处的高度'' 除以 $x$, 量级为 $e^{-x^2/2}$. 表~\ref{tab:normal-tail} 显示, 当 $x\ge2$ 时上下界已经相当接近. 常用的 ``$3\sigma$ 原则'' 即 $\P(\abs{Z}>3)\approx0.27\%$.
\end{intuition}

\begin{table}[htbp]
\centering
\caption{标准正态尾概率与定理~\ref{thm:mills} 的上下界}
\label{tab:normal-tail}
\renewcommand{\arraystretch}{1.3}
\begin{tabular}{ccccc}
\toprule
$x$ & 下界 $\frac{x\varphi(x)}{1+x^2}$ & $1-\Phi(x)$ & 上界 $\frac{\varphi(x)}{x}$ & $\P(\abs{Z}>x)$\\
\midrule
$1$ & $1.210\times10^{-1}$ & $1.587\times10^{-1}$ & $2.420\times10^{-1}$ & $0.3173$\\
$2$ & $2.160\times10^{-2}$ & $2.275\times10^{-2}$ & $2.700\times10^{-2}$ & $0.0455$\\
$3$ & $1.330\times10^{-3}$ & $1.350\times10^{-3}$ & $1.477\times10^{-3}$ & $0.0027$\\
$4$ & $3.149\times10^{-5}$ & $3.167\times10^{-5}$ & $3.346\times10^{-5}$ & $6.3\times10^{-5}$\\
\bottomrule
\end{tabular}
\end{table}

\begin{figure}[htbp]
\centering
\begin{tikzpicture}
\begin{axis}[SAaxis, width=11cm, height=5cm, xmin=-4.2, xmax=4.4, ymin=0, ymax=0.46,
  xtick={-3,-2,-1,1,2,3}, xticklabels={$-3$,$-2$,$-1$,$1$,$2$,$3$}, ytick=\empty, xlabel={$x$}]
  \addplot[fill=SABlue!12,draw=none,domain=-1:1,samples=50] {exp(-x^2/2)/sqrt(2*pi)} \closedcycle;
  \addplot[fill=SABlue!25,draw=none,domain=1:2,samples=30] {exp(-x^2/2)/sqrt(2*pi)} \closedcycle;
  \addplot[fill=SABlue!25,draw=none,domain=-2:-1,samples=30] {exp(-x^2/2)/sqrt(2*pi)} \closedcycle;
  \addplot[fill=SARed!45,draw=none,domain=2:4.2,samples=40] {exp(-x^2/2)/sqrt(2*pi)} \closedcycle;
  \addplot[SADark,thick,domain=-4.2:4.2,samples=150] {exp(-x^2/2)/sqrt(2*pi)};
  \node[font=\small] at (axis cs:0,0.12) {$68.3\%$};
  \node[font=\small,SARed] at (axis cs:3.0,0.07) {$2.3\%$};
  \draw[->,SARed] (axis cs:2.9,0.055) -- (axis cs:2.35,0.012);
  \node[font=\small] at (axis cs:-1.47,0.085) {$13.6\%$};
\end{axis}
\end{tikzpicture}
\caption{标准正态密度. $\P(\abs{Z}\le1)\approx68.3\%$, $\P(\abs{Z}\le2)\approx95.4\%$, 单侧尾部 $\P(Z>2)\approx2.3\%$}
\label{fig:normal-density}
\end{figure}

下面的例子把尾概率估计与两个 Borel-Cantelli 引理结合起来, 精确地刻画了 $\iid$ 正态序列的 ``最大波动幅度''.

\begin{example}[正态序列的增长速度]\label{ex:normal-max}
设 $Z_1,Z_2,\dots$ 是 $\iid$ 的 $N(0,1)$ 随机变量. 则
\[
\limsup_{n\to\infty}\frac{Z_n}{\sqrt{2\log n}}=1\qquad\as
\]

\emph{证明.} 固定 $\varepsilon\in(0,1)$, 对 $n\ge3$ 记 $a_n=(1+\varepsilon)\sqrt{2\log n}$, $b_n=(1-\varepsilon)\sqrt{2\log n}$.

\emph{上界.} 由定理~\ref{thm:mills} 的上界, 当 $a_n\ge1$ 时,
\[
\P(Z_n>a_n)\le\frac{\varphi(a_n)}{a_n}\le\varphi(a_n)=\frac{1}{\sqrt{2\pi}}e^{-(1+\varepsilon)^2\log n}=\frac{1}{\sqrt{2\pi}}\,n^{-(1+\varepsilon)^2}.
\]
由于 $(1+\varepsilon)^2>1$, $\sum_n\P(Z_n>a_n)<\infty$. 由 Borel-Cantelli 第一引理 (定理~\ref{thm:bc1}), 几乎必然地, 只有有限个 $n$ 使 $Z_n>a_n$, 从而 $\limsup_n\frac{Z_n}{\sqrt{2\log n}}\le1+\varepsilon$ $\as$

\emph{下界.} 记 $s=(1-\varepsilon)^2<1$. 当 $b_n\ge1$ 时 $\frac{b_n}{1+b_n^2}\ge\frac{1}{2b_n}$, 由定理~\ref{thm:mills} 的下界,
\[
\P(Z_n>b_n)\ge\frac{\varphi(b_n)}{2b_n}=\frac{1}{2\sqrt{2\pi}}\cdot\frac{n^{-s}}{(1-\varepsilon)\sqrt{2\log n}}\ge\frac{1}{2\sqrt{2\pi}}\cdot\frac{n^{-s}}{\sqrt{2\log n}}.
\]
由于 $\log n$ 比 $n$ 的任何正幂次增长得慢, 当 $n$ 充分大时 $\sqrt{2\log n}\le n^{1-s}$, 于是右端不小于 $\frac{1}{2\sqrt{2\pi}}\cdot\frac{1}{n}$, 所以 $\sum_n\P(Z_n>b_n)=\infty$. 事件 $\{Z_n>b_n\}\in\sigma(Z_n)$ 相互独立, 由 Borel-Cantelli 第二引理 (定理~\ref{thm:bc2}), 几乎必然有无穷多个 $n$ 使 $Z_n>b_n$, 从而 $\limsup_n\frac{Z_n}{\sqrt{2\log n}}\ge1-\varepsilon$ $\as$

\emph{合并.} 取 $\varepsilon=\frac{1}{k}$ ($k=2,3,\dots$), 上面得到可列个概率为 $1$ 的事件, 它们的交的概率仍为 $1$ (其补集是可列个零概率集之并). 在这个交集上, 对所有 $k$ 都有 $1-\frac{1}{k}\le\limsup_n\frac{Z_n}{\sqrt{2\log n}}\le1+\frac{1}{k}$, 令 $k\to\infty$ 即得结论.
\end{example}

\begin{intuition}
$n$ 个独立标准正态变量中最大的那个大约是 $\sqrt{2\log n}$: 增长极其缓慢, 这反映了正态尾部 $e^{-x^2/2}$ 的快速衰减. 第~4~章中布朗运动的重对数律 $\limsup_{t\to\infty}\frac{B_t}{\sqrt{2t\log\log t}}=1$ 是同一现象在连续时间中的体现.
\end{intuition}

\begin{remark}[中心极限定理的背景]
正态分布之所以无处不在, 是因为大量独立微小因素之和近似服从正态分布. 例如抛 $n$ 次硬币, 正面次数 $S_n$ 经标准化后, $\frac{S_n-n/2}{\sqrt{n}/2}$ 的分布当 $n\to\infty$ 时趋于 $N(0,1)$ (de Moivre-Laplace 定理). 这一现象的一般形式是中心极限定理, 将在 1.4 节叙述; 它在第~4~章发展为 Donsker 不变原理, 由此得到布朗运动.
\end{remark}

% =====================================================================
\section{数学期望与积分收敛定理}
\label{sec:1-3}

对只取有限个值的随机变量, 期望就是加权平均: 若 $X$ 以概率 $p_i$ 取值 $x_i$, 则 $\E[X]=\sum_i x_ip_i$. 对一般的随机变量, 我们需要一种 ``积分''. Riemann 积分把\emph{定义域}切成小段, 但 $\Omega$ 只是一个抽象集合, 没有 ``小段'' 可言. Lebesgue 的想法是改为切分\emph{值域}: 把 $X$ 的取值分成若干层, 每层的 ``宽度'' 用该层原像的概率来度量 (图~\ref{fig:riemann-lebesgue}). 这正是 1.2 节简单函数逼近 (引理~\ref{lem:simple-approx}) 所做的事情.

\begin{figure}[htbp]
\centering
\begin{tikzpicture}[scale=1.05]
\begin{scope}
  \draw[->,SADark!80] (0,0) -- (4.3,0) node[right,font=\small] {$\omega$};
  \draw[->,SADark!80] (0,0) -- (0,2.7);
  \foreach \i in {0,...,7} {
    \pgfmathsetmacro{\x}{0.5*\i}
    \pgfmathsetmacro{\h}{1.3+0.9*sin(deg(1.6*(\x+0.25)))}
    \fill[SABlue!22] (\x,0) rectangle ({\x+0.5},\h);
    \draw[SABlue!70] (\x,0) rectangle ({\x+0.5},\h);
  }
  \draw[SADark,thick,domain=0:4,samples=60] plot (\x,{1.3+0.9*sin(deg(1.6*\x))});
  \node[font=\small] at (2,-0.5) {(a) Riemann: 切分定义域};
\end{scope}
\begin{scope}[xshift=6cm]
  \draw[->,SADark!80] (0,0) -- (4.3,0) node[right,font=\small] {$\omega$};
  \draw[->,SADark!80] (0,0) -- (0,2.7);
  \foreach \y in {0.6,1.2,1.8} { \draw[SADark!25,dashed] (0,\y) -- (4.1,\y); }
  % level set {X >= 1.8}: sin(1.6x) >= 0.5556 -> 1.6x in [0.5891, 2.5525] -> x in [0.368,1.595]
  \fill[SARed!30] (0.368,1.2) rectangle (1.595,1.8);
  \draw[SARed,line width=2pt] (0.368,0) -- (1.595,0);
  \draw[SARed!70,dashed] (0.368,0) -- (0.368,1.8);
  \draw[SARed!70,dashed] (1.595,0) -- (1.595,1.8);
  \draw[SADark,thick,domain=0:4,samples=60] plot (\x,{1.3+0.9*sin(deg(1.6*\x))});
  \node[left,font=\small] at (0,1.8) {$y_{k+1}$};
  \node[left,font=\small] at (0,1.2) {$y_k$};
  \node[SARed,font=\small,below] at (0.98,0) {$\{X\ge y_{k+1}\}$};
  \node[font=\small] at (2,-0.85) {(b) Lebesgue: 切分值域};
\end{scope}
\end{tikzpicture}
\caption{两种积分思想. Lebesgue 积分把值域分层, 用 ``层高 $\times$ 原像的概率'' 求和, 因而不依赖于 $\Omega$ 的任何几何结构}
\label{fig:riemann-lebesgue}
\end{figure}

% ---------------------------------------------------------------------
\subsection{期望的构造}

构造分四步: 示性函数 $\to$ 非负简单随机变量 $\to$ 非负随机变量 $\to$ 可积随机变量. 每一步都由前一步通过极限或线性组合得到.

\paragraph{第一、二步: 简单随机变量.}
设 $X\ge0$ 是简单随机变量, 其互不相同的取值为 $a_1,\dots,a_k$, 令 $A_i=\{X=a_i\}$, 则 $\{A_1,\dots,A_k\}$ 是 $\Omega$ 的分割且 $X=\sum_{i=1}^k a_i\1_{A_i}$. 定义
\begin{equation}\label{eq:simple-expectation}
\E[X]=\sum_{i=1}^k a_i\P(A_i).
\end{equation}
特别地, 对事件 $A$ 有 $\E[\1_A]=\P(A)$.

\begin{lemma}\label{lem:simple-properties}
\begin{enumerate}
\item 若 $X=\sum_{j=1}^m b_j\1_{B_j}$, 其中 $\{B_1,\dots,B_m\}$ 是 $\Omega$ 的分割 ($b_j\ge0$ 不必互不相同), 则 $\E[X]=\sum_{j=1}^m b_j\P(B_j)$.
\item 设 $X,Y$ 是非负简单随机变量, $a,b\ge0$, 则 $\E[aX+bY]=a\E[X]+b\E[Y]$.
\item 若非负简单随机变量满足 $X\le Y$, 则 $\E[X]\le\E[Y]$.
\end{enumerate}
\end{lemma}

\begin{proof}
(1) 对每个取值 $a_i$, $A_i=\{X=a_i\}$ 恰是满足 $b_j=a_i$ 的那些 $B_j$ 之不交并. 由有限可加性,
\[
\sum_{j=1}^m b_j\P(B_j)=\sum_{i=1}^k\sum_{j:\,b_j=a_i}a_i\P(B_j)=\sum_{i=1}^k a_i\P(A_i)=\E[X].
\]

(2)(3) 设 $X=\sum_i a_i\1_{A_i}$, $Y=\sum_j b_j\1_{B_j}$ 是按取值的分割表示. 令 $C_{ij}=A_i\cap B_j$, 则 $\{C_{ij}\}$ 也是 $\Omega$ 的分割, 且在 $C_{ij}$ 上 $X=a_i$, $Y=b_j$, $aX+bY=aa_i+bb_j$. 由 (1) 和有限可加性,
\begin{align*}
\E[aX+bY]&=\sum_{i,j}(aa_i+bb_j)\P(C_{ij})=a\sum_i a_i\sum_j\P(C_{ij})+b\sum_j b_j\sum_i\P(C_{ij})\\
&=a\sum_i a_i\P(A_i)+b\sum_j b_j\P(B_j)=a\E[X]+b\E[Y].
\end{align*}
若 $X\le Y$, 则对每个非空的 $C_{ij}$ 有 $a_i\le b_j$ (空集的概率为 $0$, 不影响求和), 所以 $\E[X]=\sum_{i,j}a_i\P(C_{ij})\le\sum_{i,j}b_j\P(C_{ij})=\E[Y]$.
\end{proof}

由 (2), 对任意 (不必不交的) 事件 $A_i$ 和 $a_i\ge0$, 都有 $\E\big[\sum_i a_i\1_{A_i}\big]=\sum_i a_i\P(A_i)$.

\paragraph{第三步: 非负随机变量.}

\begin{definition}[非负随机变量的期望]\label{def:expectation-nonneg}
设 $X$ 是非负广义随机变量. 定义
\[
\E[X]=\sup\big\{\E[Y]:\ Y\text{ 是简单随机变量},\ 0\le Y\le X\big\}\in[0,+\infty].
\]
\end{definition}

由引理~\ref{lem:simple-properties} (3), 当 $X$ 本身是简单随机变量时, 上确界在 $Y=X$ 处取到, 与 \eqref{eq:simple-expectation} 一致. 定义中的上确界不便计算. 下面的引理说明, 沿任何一列单调逼近 $X$ 的简单随机变量取极限即可, 它是本节一切收敛定理的基础.

\begin{lemma}\label{lem:simple-increasing}
设 $X\ge0$, $Y_n$ 是非负简单随机变量且 $Y_n\uparrow X$ (逐点). 则 $\E[Y_n]\uparrow\E[X]$.
\end{lemma}

\begin{proof}
由引理~\ref{lem:simple-properties} (3), $\E[Y_n]$ 单调不减; 由定义, $\E[Y_n]\le\E[X]$. 所以 $\lim_n\E[Y_n]$ 存在且 $\le\E[X]$. 下证反向不等式, 只需证明: 对任意满足 $0\le Y\le X$ 的简单随机变量 $Y$, 有 $\lim_n\E[Y_n]\ge\E[Y]$.

设 $Y=\sum_{i=1}^k a_i\1_{A_i}$ (按取值的分割表示), 固定 $c\in(0,1)$, 令 $B_n=\{Y_n\ge cY\}$. 我们断言 $B_n\uparrow\Omega$. 由于 $Y_n$ 单调增加, $B_n$ 单调增加. 任取 $\omega$: 若 $Y(\omega)=0$, 则对所有 $n$ 都有 $\omega\in B_n$; 若 $Y(\omega)>0$, 则 $cY(\omega)<Y(\omega)\le X(\omega)$, 而 $Y_n(\omega)\uparrow X(\omega)$, 故当 $n$ 充分大时 $Y_n(\omega)\ge cY(\omega)$, 即 $\omega\in B_n$.

由于 $Y_n\ge Y_n\1_{B_n}\ge cY\1_{B_n}$, 而 $cY\1_{B_n}=\sum_i ca_i\1_{A_i\cap B_n}$ 是简单随机变量, 由引理~\ref{lem:simple-properties},
\[
\E[Y_n]\ge\E[cY\1_{B_n}]=c\sum_{i=1}^k a_i\P(A_i\cap B_n).
\]
令 $n\to\infty$, 由于 $A_i\cap B_n\uparrow A_i$, 由概率的连续性 (定理~\ref{thm:continuity}), 右端趋于 $c\sum_i a_i\P(A_i)=c\E[Y]$. 于是 $\lim_n\E[Y_n]\ge c\E[Y]$, 再令 $c\uparrow1$ 即得.
\end{proof}

\begin{corollary}\label{cor:nonneg-properties}
设 $X,Y$ 是非负广义随机变量, $a,b\ge0$. 则
\begin{enumerate}
\item 线性: $\E[aX+bY]=a\E[X]+b\E[Y]$;
\item 单调性: 若 $X\le Y$, 则 $\E[X]\le\E[Y]$;
\item 若 $\E[X]=0$, 则 $X=0$ $\as$;
\item 若 $\E[X]<\infty$, 则 $X<\infty$ $\as$;
\item 若 $X=Y$ $\as$, 则 $\E[X]=\E[Y]$.
\end{enumerate}
\end{corollary}

\begin{proof}
(1) 由引理~\ref{lem:simple-approx}, 取简单随机变量 $X_n\uparrow X$, $Y_n\uparrow Y$. 则 $aX_n+bY_n$ 是简单随机变量且 $\uparrow aX+bY$. 由引理~\ref{lem:simple-increasing} 和引理~\ref{lem:simple-properties} (2),
\[
\E[aX+bY]=\lim_n\E[aX_n+bY_n]=\lim_n\big(a\E[X_n]+b\E[Y_n]\big)=a\E[X]+b\E[Y].
\]

(2) 满足 $0\le Z\le X$ 的简单随机变量 $Z$ 也满足 $0\le Z\le Y$, 所以定义中前者的上确界不超过后者.

(3) 对每个 $n$, $\frac{1}{n}\1_{\{X\ge1/n\}}\le X$, 由 (2), $\frac{1}{n}\P(X\ge\frac{1}{n})\le\E[X]=0$. 于是 $\P(X>0)=\P\big(\bigcup_n\{X\ge\frac{1}{n}\}\big)\le\sum_n\P(X\ge\frac{1}{n})=0$.

(4) 对每个 $n$, $n\1_{\{X=\infty\}}\le X$, 故 $n\P(X=\infty)\le\E[X]$. 令 $n\to\infty$, 由 $\E[X]<\infty$ 得 $\P(X=\infty)=0$.

(5) 设 $N=\{X\ne Y\}$, $\P(N)=0$. 对满足 $0\le Z\le X\1_N$ 的简单随机变量 $Z$, $Z$ 在 $N^c$ 上为 $0$, 故 $\E[Z]\le(\max Z)\P(N)=0$, 从而 $\E[X\1_N]=0$; 同理 $\E[Y\1_N]=0$. 由 (1),
$\E[X]=\E[X\1_{N^c}]+\E[X\1_N]=\E[Y\1_{N^c}]=\E[Y]$.
\end{proof}

\paragraph{第四步: 一般随机变量.}
对一般的 $X$, 令 $X^+=X\vee0$, $X^-=(-X)\vee0$, 则 $X^\pm\ge0$, $X=X^+-X^-$, $\abs{X}=X^++X^-$.

\begin{definition}[期望]\label{def:expectation}
若 $\E[\abs{X}]<\infty$, 称 $X$ \textbf{可积}, 并定义
\[
\E[X]=\E[X^+]-\E[X^-].
\]
可积随机变量全体记为 $L^1(\Omega,\mathcal{F},\P)$ 或简记为 $L^1$. 期望也写作积分 $\int_\Omega X\dd\P$. 对事件 $A$, 记 $\E[X;A]=\E[X\1_A]=\int_A X\dd\P$.
\end{definition}

\begin{remark}
若 $\E[X^+]$ 与 $\E[X^-]$ 中只有一个为 $+\infty$, 仍可按上式定义 $\E[X]=\pm\infty$, 称 $X$ 的期望\emph{存在}. 若二者都为 $+\infty$, 期望无定义.
\end{remark}

\begin{theorem}[期望的基本性质]\label{thm:expectation-properties}
设 $X,Y\in L^1$, $a,b\in\R$. 则
\begin{enumerate}
\item 线性: $aX+bY\in L^1$, 且 $\E[aX+bY]=a\E[X]+b\E[Y]$;
\item 单调性: 若 $X\le Y$ $\as$, 则 $\E[X]\le\E[Y]$;
\item $\abs{\E[X]}\le\E[\abs{X}]$;
\item 若 $X=Y$ $\as$, 则 $\E[X]=\E[Y]$.
\end{enumerate}
\end{theorem}

\begin{proof}
(1) 先证数乘. 若 $a\ge0$, 则 $(aX)^\pm=aX^\pm$; 若 $a<0$, 则 $(aX)^+=\abs{a}X^-$, $(aX)^-=\abs{a}X^+$. 两种情况下由推论~\ref{cor:nonneg-properties} (1) 都有 $\E[aX]=a\E[X]$.

再证加法. 记 $Z=X+Y$. 由 $\abs{Z}\le\abs{X}+\abs{Y}$ 及推论~\ref{cor:nonneg-properties} (1)(2), $\E\abs{Z}\le\E\abs{X}+\E\abs{Y}<\infty$. 由
$Z^+-Z^-=Z=X^+-X^-+Y^+-Y^-$, 移项得到非负随机变量之间的等式
\[
Z^++X^-+Y^-=Z^-+X^++Y^+.
\]
两边取期望, 由推论~\ref{cor:nonneg-properties} (1),
$\E[Z^+]+\E[X^-]+\E[Y^-]=\E[Z^-]+\E[X^+]+\E[Y^+]$. 各项都有限, 移项即得 $\E[Z]=\E[X]+\E[Y]$.

(2) 由 (1), $Y-X\in L^1$, 且 $(Y-X)^-=0$ $\as$, 由推论~\ref{cor:nonneg-properties} (5), $\E[(Y-X)^-]=0$. 于是 $\E[Y]-\E[X]=\E[Y-X]=\E[(Y-X)^+]\ge0$.

(3) 由 $-\abs{X}\le X\le\abs{X}$ 及 (2), $-\E\abs{X}\le\E[X]\le\E\abs{X}$.

(4) 由 (2), $\E[X]\le\E[Y]$ 且 $\E[Y]\le\E[X]$.
\end{proof}

\begin{intuition}
性质 (4) 说明, 期望 ``看不见'' 零概率集. 因此在概率论中, 两个几乎必然相等的随机变量常被视为同一个. 后文中 ``$X_n\to X$''、``$X\ge0$'' 之类的条件, 若不特别说明, 都只需几乎必然成立.
\end{intuition}

\begin{remark}[一般测度下的积分]
上述四步构造对任意测度空间 $(E,\mathcal{E},\mu)$ 都适用, 只需把 $\P$ 换成 $\mu$, 得到积分 $\int_E g\dd\mu$. 特别地, 对 $\R$ 上的 Lebesgue 测度 $\lambda$, 积分 $\int g\dd\lambda$ 记作 $\int g(x)\dd x$. 对 Riemann 可积函数 (以及绝对收敛的反常积分), Lebesgue 积分与 Riemann 积分相等 (见~\cite{folland} 2.5 节). 1.2 节中的密度积分 $\int_B f(t)\dd t$ 即按此理解.
\end{remark}

实际计算期望时, 我们通常把 $\Omega$ 上的积分转移到 $\R$ 上计算.

\begin{theorem}[变量替换公式]\label{thm:change-of-variables}
设 $X$ 是随机变量, $g:\R\to\R$ 是 Borel 函数, 且 $g\ge0$ 或 $g(X)\in L^1$. 则
\[
\E[g(X)]=\int_{\R}g(x)\,\mu_X(\dd x).
\]
特别地, 若 $X$ 有密度 $f$, 则 $\E[g(X)]=\int_{\R}g(x)f(x)\dd x$; 若 $X$ 是离散型, 则 $\E[g(X)]=\sum_{x\in S}g(x)p(x)$.
\end{theorem}

\begin{proof}
按四步构造的顺序证明.

\emph{示性函数.} 若 $g=\1_B$, $B\in\mathcal{B}(\R)$, 则 $\E[\1_B(X)]=\P(X\in B)=\mu_X(B)=\int\1_B\dd\mu_X$.

\emph{非负简单函数.} 若 $g=\sum_i a_i\1_{B_i}$, $a_i\ge0$, 由两边的线性即得.

\emph{非负函数.} 由引理~\ref{lem:simple-approx} (取 $\Omega=\R$), 存在非负简单 Borel 函数 $g_n\uparrow g$. 于是 $g_n(X)$ 是非负简单随机变量且 $g_n(X)\uparrow g(X)$. 在 $\Omega$ 上和在 $(\R,\mathcal{B}(\R),\mu_X)$ 上分别使用引理~\ref{lem:simple-increasing},
\[
\E[g(X)]=\lim_n\E[g_n(X)]=\lim_n\int g_n\dd\mu_X=\int g\dd\mu_X.
\]

\emph{一般情形.} 对 $g^+$ 和 $g^-$ 分别使用上一步 ($g(X)^\pm=g^\pm(X)$), 再相减.

对有密度的情形, 只需对示性函数验证
\[
\int\1_B\dd\mu_X=\mu_X(B)=\int_B f(x)\dd x=\int\1_B(x)f(x)\dd x
\]
(见 1.2 节绝对连续型的讨论), 其余三步完全相同 (第三步中 $g_nf\uparrow gf$). 离散情形同理, 此时 $\mu_X(B)=\sum_{x\in B\cap S}p(x)$.
\end{proof}

\begin{intuition}
期望只依赖于分布: 要计算 $\E[g(X)]$, 只需知道 $X$ 落在各处的概率, 而无需知道 $X$ 如何依赖于 $\omega$. 这在英文文献中被戏称为 ``无意识统计学家法则'' (law of the unconscious statistician): 人们常常不假思索地这样计算, 而上面的定理说明这样做是正确的.
\end{intuition}

\begin{example}[常用分布的期望]\label{ex:means}
\begin{enumerate}
\item \textbf{二项分布.} 设 $S_n$ 为 $n$ 次独立试验中的成功次数, 每次成功概率为 $p$, 则 $S_n=\sum_{k=1}^n\1_{A_k}$, 其中 $A_k$ 为第 $k$ 次成功. 由线性, $\E[S_n]=\sum_k\P(A_k)=np$. 用示性函数之和表示计数, 可以避开复杂的组合求和.
\item \textbf{Poisson 分布.} $\E[X]=\sum_{k\ge1}k\,e^{-\lambda}\frac{\lambda^k}{k!}=\lambda\sum_{k\ge1}e^{-\lambda}\frac{\lambda^{k-1}}{(k-1)!}=\lambda$.
\item \textbf{指数分布.} 分部积分,
$\E[X]=\int_0^\infty x\lambda e^{-\lambda x}\dd x=\Big[-xe^{-\lambda x}\Big]_0^\infty+\int_0^\infty e^{-\lambda x}\dd x=\frac{1}{\lambda}$.
\item \textbf{正态分布.} 设 $Z\sim N(0,1)$. 由 $\varphi'(x)=-x\varphi(x)$,
$\E\abs{Z}=2\int_0^\infty x\varphi(x)\dd x=2\varphi(0)=\sqrt{2/\pi}<\infty$, 所以 $Z$ 可积. 又 $x\varphi(x)$ 是奇函数, 故 $\E[Z]=0$. 对 $X=\mu+\sigma Z\sim N(\mu,\sigma^2)$, 由线性 $\E[X]=\mu$.
\end{enumerate}
\end{example}

% ---------------------------------------------------------------------
\subsection{积分收敛定理}

随机分析中处处需要交换极限与期望, 例如计算随机积分的期望时要对黎曼和取极限. 但交换并不总是允许的.

\begin{example}[质量逃逸]\label{ex:mass-escape}
在 $([0,1],\mathcal{B}([0,1]),\lambda)$ 上, 令 $X_n=n\1_{(0,1/n)}$. 对每个 $\omega\in[0,1]$, 当 $n>1/\omega$ (或 $\omega=0$) 时 $X_n(\omega)=0$, 所以 $X_n\to0$ 处处成立; 但 $\E[X_n]=n\cdot\frac{1}{n}=1$ 对所有 $n$ 成立. 因此
\[
\lim_{n\to\infty}\E[X_n]=1\ne0=\E\Big[\lim_{n\to\infty}X_n\Big].
\]
如图~\ref{fig:mass-escape}, $X_n$ 的 ``质量'' 被挤压到越来越窄的区间上, 最终从积分中 ``逃逸''.
\end{example}

\begin{figure}[htbp]
\centering
\begin{tikzpicture}[xscale=5,yscale=0.42]
  \draw[->,SADark!80] (0,0) -- (1.1,0) node[right,font=\small] {$\omega$};
  \draw[->,SADark!80] (0,0) -- (0,8.8);
  \foreach \n/\c in {2/SAGreenText,4/SABlue,8/SARed} {
    \fill[\c!18] (0,0) rectangle ({1/\n},\n);
    \draw[\c,thick] (0,0) -- (0,\n) -- ({1/\n},\n) -- ({1/\n},0);
    \node[\c,font=\small,anchor=west] at ({1/\n+0.01},\n) {$X_{\n}$};
  }
  \node[below,font=\small] at (1,0) {$1$};
  \foreach \n in {2,4,8} { \node[left,font=\small] at (0,\n) {$\n$}; }
  \node[font=\small,align=left,anchor=west] at (0.55,5) {每个矩形面积都是 $1$,\\ 但 $X_n(\omega)\to0$ 对每个 $\omega$ 成立};
\end{tikzpicture}
\caption{例~\ref{ex:mass-escape}: $X_n=n\1_{(0,1/n)}$ 逐点趋于 $0$, 期望却恒为 $1$}
\label{fig:mass-escape}
\end{figure}

下面三个定理给出允许交换的条件: 单调 (MCT)、非负时的单向不等式 (Fatou)、存在可积控制 (DCT).

\begin{theorem}[单调收敛定理, MCT]\label{thm:mct}
设 $0\le X_1\le X_2\le\cdots$ 且 $X_n\uparrow X$ $\as$ 则 $\E[X_n]\uparrow\E[X]$.
\end{theorem}

\begin{proof}
由推论~\ref{cor:nonneg-properties} (5), 在零概率集上修改各随机变量不影响期望, 故不妨设单调性与收敛对每个 $\omega$ 成立.

对每个 $n$, 由引理~\ref{lem:simple-approx}, 取非负简单随机变量 $Y_{n,k}\uparrow X_n$ ($k\to\infty$). 令
\[
Z_k=\max\{Y_{1,k},Y_{2,k},\dots,Y_{k,k}\}.
\]
$Z_k$ 是非负简单随机变量, 且有以下三条性质.
\begin{itemize}
\item $Z_k$ 单调不减: $Z_{k+1}\ge\max_{n\le k}Y_{n,k+1}\ge\max_{n\le k}Y_{n,k}=Z_k$.
\item $Z_k\le X_k$: 对 $n\le k$, $Y_{n,k}\le X_n\le X_k$.
\item $Z_k\uparrow X$: 由上一条, $\lim_k Z_k\le\lim_k X_k=X$. 反过来, 固定 $n$, 当 $k\ge n$ 时 $Z_k\ge Y_{n,k}$, 令 $k\to\infty$ 得 $\lim_k Z_k\ge X_n$, 再令 $n\to\infty$ 得 $\lim_k Z_k\ge X$.
\end{itemize}
由引理~\ref{lem:simple-increasing}, $\E[Z_k]\uparrow\E[X]$. 又由单调性 (推论~\ref{cor:nonneg-properties} (2)), $\E[Z_k]\le\E[X_k]\le\E[X]$. 令 $k\to\infty$, 夹逼得 $\E[X_k]\to\E[X]$.
\end{proof}

\begin{corollary}[逐项求期望]\label{cor:series}
若 $X_n\ge0$, 则 $\E\big[\sum_{n=1}^\infty X_n\big]=\sum_{n=1}^\infty\E[X_n]$.
\end{corollary}

\begin{proof}
对部分和 $\sum_{n\le N}X_n\uparrow\sum_{n}X_n$ 使用单调收敛定理~\ref{thm:mct} 与有限线性 (推论~\ref{cor:nonneg-properties} (1)).
\end{proof}

\begin{example}[Borel-Cantelli 第一引理的另一证明]
令 $N=\sum_{n=1}^\infty\1_{A_n}$ 为事件 $A_n$ 的发生次数. 由推论~\ref{cor:series}, $\E[N]=\sum_n\P(A_n)$. 若此级数收敛, 则由推论~\ref{cor:nonneg-properties} (4), $N<\infty$ $\as$, 这正是定理~\ref{thm:bc1}. 同理, 1.2 节中 $\nu(B)=\int_B f(t)\dd t$ 的可列可加性也是推论~\ref{cor:series} 的直接结果: 对两两不交的 $B_n$, $f\1_{\bigcup_n B_n}=\sum_n f\1_{B_n}$.
\end{example}

\begin{theorem}[Fatou 引理]\label{thm:fatou}
设 $X_n\ge0$. 则
\[
\E\Big[\liminf_{n\to\infty}X_n\Big]\le\liminf_{n\to\infty}\E[X_n].
\]
\end{theorem}

\begin{proof}
令 $Y_m=\inf_{n\ge m}X_n$. 则 $0\le Y_m\uparrow\liminf_n X_n$, 且对每个 $n\ge m$ 有 $Y_m\le X_n$, 从而 $\E[Y_m]\le\inf_{n\ge m}\E[X_n]$. 由单调收敛定理~\ref{thm:mct},
\[
\E\Big[\liminf_{n\to\infty}X_n\Big]=\lim_{m\to\infty}\E[Y_m]\le\lim_{m\to\infty}\inf_{n\ge m}\E[X_n]=\liminf_{n\to\infty}\E[X_n].\qedhere
\]
\end{proof}

\begin{intuition}
取极限时, 期望只可能 ``丢失'' 质量, 不会 ``凭空产生'' 质量. 例~\ref{ex:mass-escape} 说明不等号可以严格成立: 左边为 $0$, 右边为 $1$.
\end{intuition}

\begin{theorem}[控制收敛定理, DCT]\label{thm:dct}
设 $X_n\to X$ $\as$, 且存在 $Y\in L^1$ 使对所有 $n$ 有 $\abs{X_n}\le Y$ $\as$ 则 $X\in L^1$, 且
\[
\lim_{n\to\infty}\E\abs{X_n-X}=0,\qquad \lim_{n\to\infty}\E[X_n]=\E[X].
\]
\end{theorem}

\begin{proof}
仍不妨设所有条件对每个 $\omega$ 成立. 取极限得 $\abs{X}\le Y$, 故 $X\in L^1$, 且 $\abs{X_n-X}\le2Y$. 对非负随机变量 $2Y-\abs{X_n-X}$ 使用 Fatou 引理~\ref{thm:fatou}. 由于 $\abs{X_n-X}\to0$, 左端的下极限为 $2Y$:
\begin{align*}
\E[2Y]&=\E\Big[\liminf_{n\to\infty}\big(2Y-\abs{X_n-X}\big)\Big]\le\liminf_{n\to\infty}\E\big[2Y-\abs{X_n-X}\big]\\
&=\E[2Y]-\limsup_{n\to\infty}\E\abs{X_n-X}.
\end{align*}
由于 $\E[2Y]<\infty$, 两边消去得 $\limsup_n\E\abs{X_n-X}\le0$, 即 $\E\abs{X_n-X}\to0$. 最后由定理~\ref{thm:expectation-properties} (1)(3), $\abs{\E[X_n]-\E[X]}\le\E\abs{X_n-X}\to0$.
\end{proof}

\begin{intuition}
质量逃逸的根源是 $X_n$ 在小概率集合上取极大的值. 可积控制 $\abs{X_n}\le Y$ 把所有 $X_n$ 关进同一个 ``笼子'', 使它们在小概率集上的贡献一致地小, 质量无处可逃. 例~\ref{ex:mass-escape} 中 $\sup_n X_n$ 在 $(\frac{1}{k+1},\frac{1}{k})$ 上等于 $k$, 其期望 $\sum_k k\cdot\frac{1}{k(k+1)}=\sum_k\frac{1}{k+1}=\infty$, 因而不存在可积控制. 1.4 节的一致可积性将给出交换极限与期望的 ``最佳'' 条件.
\end{intuition}

\begin{corollary}[有界收敛定理]\label{cor:bct}
若 $X_n\to X$ $\as$, 且存在常数 $M$ 使 $\abs{X_n}\le M$, 则 $\E[X_n]\to\E[X]$.
\end{corollary}

\begin{proof}
常数 $M$ 可积 ($\E[M]=M$), 由控制收敛定理~\ref{thm:dct} 即得.
\end{proof}

控制收敛定理最常用的推论之一是 ``积分号下求导'', 后面计算矩母函数和特征函数的导数时要用到它.

\begin{theorem}[期望号下求导]\label{thm:diff-under-E}
设 $I\subseteq\R$ 是开区间, 对每个 $t\in I$, $h(t,\cdot)$ 是可积随机变量, 且
\begin{enumerate}
\item 对每个 $\omega$, $t\mapsto h(t,\omega)$ 在 $I$ 上可微;
\item 存在 $Y\in L^1$, 使得对所有 $t\in I$ 和 $\omega$, $\abs{\partial_t h(t,\omega)}\le Y(\omega)$.
\end{enumerate}
则 $t\mapsto\E[h(t,\cdot)]$ 在 $I$ 上可微, 且 $\frac{\dd}{\dd t}\E[h(t,\cdot)]=\E[\partial_t h(t,\cdot)]$.
\end{theorem}

\begin{proof}
固定 $t\in I$, 任取数列 $\delta_n\to0$, $\delta_n\ne0$, $t+\delta_n\in I$. 令
\[
D_n(\omega)=\frac{h(t+\delta_n,\omega)-h(t,\omega)}{\delta_n}.
\]
由条件 (1), $D_n(\omega)\to\partial_t h(t,\omega)$. 由微分中值定理, 存在 $\theta=\theta(n,\omega)\in(0,1)$ 使 $D_n(\omega)=\partial_t h(t+\theta\delta_n,\omega)$, 再由条件 (2), $\abs{D_n}\le Y$. ($\partial_t h(t,\cdot)$ 作为可测函数 $D_n$ 的极限是随机变量.) 由控制收敛定理~\ref{thm:dct} 与期望的线性,
\[
\frac{\E[h(t+\delta_n,\cdot)]-\E[h(t,\cdot)]}{\delta_n}=\E[D_n]\to\E[\partial_t h(t,\cdot)].
\]
由于数列 $\delta_n$ 任意, 导数存在且等于 $\E[\partial_t h(t,\cdot)]$.
\end{proof}

% ---------------------------------------------------------------------
\subsection{方差、矩与常用不等式}

\begin{definition}[矩与方差]\label{def:moments}
设 $p>0$. 若 $\E\abs{X}^p<\infty$, 称 $X$ 的 $p$ 阶矩有限, 记 $X\in L^p$, 并记 $\norm{X}_p=(\E\abs{X}^p)^{1/p}$. 对正整数 $k$, $\E[X^k]$ 称为 $X$ 的 $k$ 阶\textbf{矩}. 对 $X\in L^2$, 定义\textbf{方差}与\textbf{标准差}
\[
\Var(X)=\E\big[(X-\E X)^2\big],\qquad \sigma_X=\sqrt{\Var(X)};
\]
对 $X,Y\in L^2$, 定义\textbf{协方差} $\Cov(X,Y)=\E\big[(X-\E X)(Y-\E Y)\big]$.
\end{definition}

定义中的期望都是有意义的. 首先, 若 $0<p<q$ 且 $X\in L^q$, 则 $X\in L^p$: 因为 $\abs{x}^p\le1+\abs{x}^q$ (分 $\abs{x}\le1$ 与 $\abs{x}>1$ 讨论), 所以 $\E\abs{X}^p\le1+\E\abs{X}^q<\infty$. 特别地, $L^2\subseteq L^1$. 其次, 由 $\abs{xy}\le\frac{1}{2}(x^2+y^2)$, $X,Y\in L^2$ 蕴含 $XY\in L^1$.

\begin{proposition}[方差的性质]\label{prop:variance}
设 $X,Y\in L^2$, $a,b\in\R$. 则
\begin{enumerate}
\item $\Var(X)=\E[X^2]-(\E X)^2$;
\item 平移不变与伸缩: $\Var(aX+b)=a^2\Var(X)$;
\item $\Var(X)=0$ 当且仅当 $X=\E X$ $\as$;
\item $\Var(X+Y)=\Var(X)+\Var(Y)+2\Cov(X,Y)$, 且 $\Cov(X,Y)=\E[XY]-\E[X]\E[Y]$.
\end{enumerate}
\end{proposition}

\begin{proof}
记 $m=\E X$. (1) 由线性, $\E[(X-m)^2]=\E[X^2]-2m\E[X]+m^2=\E[X^2]-m^2$.
(2) $\E[aX+b]=am+b$, 故 $\Var(aX+b)=\E[(aX-am)^2]=a^2\Var(X)$.
(3) 由推论~\ref{cor:nonneg-properties} (3), $\E[(X-m)^2]=0$ 当且仅当 $(X-m)^2=0$ $\as$
(4) 记 $\tilde X=X-\E X$, $\tilde Y=Y-\E Y$, 展开 $\E[(\tilde X+\tilde Y)^2]=\E[\tilde X^2]+\E[\tilde Y^2]+2\E[\tilde X\tilde Y]$ 即得第一式; 第二式与 (1) 的证明相同.
\end{proof}

\begin{example}[常用分布的方差]\label{ex:variances}
\begin{enumerate}
\item \textbf{Poisson 分布.} 与例~\ref{ex:means} 类似, $\E[X(X-1)]=\sum_{k\ge2}k(k-1)e^{-\lambda}\frac{\lambda^k}{k!}=\lambda^2$, 故 $\Var(X)=\E[X(X-1)]+\E[X]-(\E X)^2=\lambda$.
\item \textbf{指数分布.} 分部积分两次得 $\E[X^2]=\frac{2}{\lambda^2}$, 故 $\Var(X)=\frac{1}{\lambda^2}$.
\item \textbf{正态分布.} 设 $Z\sim N(0,1)$. 由 $\varphi'(x)=-x\varphi(x)$ 分部积分,
\[
\E[Z^2]=\int_{\R}x\cdot x\varphi(x)\dd x=\Big[-x\varphi(x)\Big]_{-\infty}^{\infty}+\int_{\R}\varphi(x)\dd x=1.
\]
故 $\Var(Z)=1$, 由命题~\ref{prop:variance} (2), $X=\mu+\sigma Z$ 满足 $\Var(X)=\sigma^2$. 这验证了 $N(\mu,\sigma^2)$ 中两个参数的含义.
\item \textbf{二项分布.} $\Var(S_n)=np(1-p)$. 这需要用到独立随机变量的协方差为零 (定理~\ref{thm:indep-product}), 将在 1.5 节说明.
\end{enumerate}
\end{example}

下面是一组基本不等式, 它们分别用矩控制概率、用期望控制函数值、用范数控制乘积.

\begin{theorem}[Markov 不等式]\label{thm:markov}
设 $\phi:[0,\infty)\to[0,\infty)$ 单调不减, $a>0$ 且 $\phi(a)>0$. 则
\[
\P(\abs{X}\ge a)\le\frac{\E[\phi(\abs{X})]}{\phi(a)}.
\]
特别地, $\P(\abs{X}\ge a)\le\frac{\E\abs{X}}{a}$, 以及 $\P(\abs{X}\ge a)\le\frac{\E\abs{X}^p}{a^p}$ ($p>0$).
\end{theorem}

\begin{proof}
在 $\{\abs{X}\ge a\}$ 上, 由单调性 $\phi(\abs{X})\ge\phi(a)$; 在其余位置 $\phi(\abs{X})\ge0$. 所以 $\phi(a)\1_{\{\abs{X}\ge a\}}\le\phi(\abs{X})$. 两边取期望即得.
\end{proof}

\begin{figure}[htbp]
\centering
\begin{tikzpicture}[scale=0.95]
  \draw[->,SADark!80] (0,0) -- (5.4,0) node[right,font=\small] {$x$};
  \draw[->,SADark!80] (0,0) -- (0,3.3);
  \draw[SABlue,very thick] (0,0) -- (3.2,3.2) node[above right,font=\small] {$y=x$};
  \draw[SARed,very thick] (0,0) -- (2,0);
  \draw[SARed,very thick] (2,2) -- (5.2,2) node[above,font=\small] {$y=a\1_{\{x\ge a\}}$};
  \fill[SARed] (2,2) circle (1.6pt); \draw[SARed,fill=white] (2,0) circle (1.6pt);
  \draw[SADark!40,dashed] (2,0) -- (2,2);
  \node[below,font=\small] at (2,0) {$a$};
  \node[left,font=\small] at (0,2) {$a$};
\end{tikzpicture}
\caption{Markov 不等式的证明只用了一个逐点不等式 $a\1_{\{x\ge a\}}\le x$ ($x\ge0$), 再取期望}
\label{fig:markov}
\end{figure}

\begin{corollary}[Chebyshev 不等式]\label{cor:chebyshev}
若 $X\in L^2$, $a>0$, 则 $\P(\abs{X-\E X}\ge a)\le\frac{\Var(X)}{a^2}$.
\end{corollary}

\begin{proof}
对 $X-\E X$ 用 Markov 不等式~\ref{thm:markov}, 取 $\phi(x)=x^2$.
\end{proof}

\begin{example}[抛硬币的弱大数定律]\label{ex:wlln}
回到 1.1 节引言的无穷次抛硬币, 令 $\xi_k=\omega_k-\frac{1}{2}$, 它以概率 $\frac{1}{2}$ 取 $\pm\frac{1}{2}$, 且由例~\ref{ex:basic-spaces} (2) 中柱集的概率, 对 $j\ne k$, $\xi_j\xi_k$ 以概率 $\frac{1}{2}$ 取 $\pm\frac{1}{4}$, 从而 $\E[\xi_j\xi_k]=0$, 而 $\E[\xi_k^2]=\frac{1}{4}$. 于是
\[
\E\Big[\Big(S_n-\frac{n}{2}\Big)^2\Big]=\E\Big[\Big(\sum_{k=1}^n\xi_k\Big)^2\Big]=\sum_{k=1}^n\E[\xi_k^2]+\sum_{j\ne k}\E[\xi_j\xi_k]=\frac{n}{4}.
\]
由 Chebyshev 不等式, 对任意 $\varepsilon>0$,
\[
\P\Big(\Big|\frac{S_n}{n}-\frac{1}{2}\Big|\ge\varepsilon\Big)\le\frac{1}{n^2\varepsilon^2}\cdot\frac{n}{4}=\frac{1}{4n\varepsilon^2}\to0.
\]
这就是说, 频率 $\frac{S_n}{n}$ ``依概率收敛'' 到 $\frac{1}{2}$ (见 1.4 节).
\end{example}

\begin{theorem}[Jensen 不等式]\label{thm:jensen}
设 $I\subseteq\R$ 是区间, $\phi:I\to\R$ 是凸函数, $X\in L^1$ 且 $X\in I$ $\as$, $\phi(X)\in L^1$. 则
\[
\phi(\E X)\le\E[\phi(X)].
\]
\end{theorem}

\begin{proof}
记 $m=\E X$. 由 $X\in I$ $\as$ 与单调性, $m\in\overline{I}$; 若 $m$ 是 $I$ 的端点 (比如 $m=\inf I$), 则 $X-m\ge0$ $\as$ 且 $\E[X-m]=0$, 由推论~\ref{cor:nonneg-properties} (3), $X=m$ $\as$, 此时不等式取等号. 下设 $m$ 是 $I$ 的内点.

凸函数在内点 $m$ 处有\textbf{支撑直线}: 存在 $c\in\R$, 使得
\begin{equation}\label{eq:support-line}
\phi(x)\ge\phi(m)+c(x-m),\qquad x\in I.
\end{equation}
事实上, 由凸性, 差商 $\frac{\phi(x)-\phi(m)}{x-m}$ 关于 $x$ 单调不减, 所以右导数 $c=\lim_{x\downarrow m}\frac{\phi(x)-\phi(m)}{x-m}$ 存在且有限, 并且当 $x>m$ 时差商 $\ge c$, 当 $x<m$ 时差商 $\le c$, 两种情况都给出 \eqref{eq:support-line}. 在 \eqref{eq:support-line} 中代入 $x=X$, 取期望, 由线性与单调性,
\[
\E[\phi(X)]\ge\phi(m)+c\,\E[X-m]=\phi(m).\qedhere
\]
\end{proof}

\begin{figure}[htbp]
\centering
\begin{tikzpicture}[scale=1.0]
  \draw[->,SADark!80] (-0.3,0) -- (5.6,0) node[right,font=\small] {$x$};
  \draw[->,SADark!80] (0,-0.3) -- (0,2.6);
  \draw[SABlue,very thick,domain=0.3:4.9,samples=80] plot (\x,{0.28*(\x-2.4)^2+0.4});
  \node[SABlue,font=\small,anchor=west] at (4.9,{0.28*2.5^2+0.4}) {$\phi$};
  \draw[SARed,thick,domain=0.3:5.2] plot (\x,{0.4112+0.112*(\x-2.6)});
  \node[SARed,font=\small,anchor=north west] at (3.9,0.5) {支撑直线};
  \foreach \x in {0.8,4.4} {
    \fill[SABlue] (\x,{0.28*(\x-2.4)^2+0.4}) circle (1.6pt);
    \draw[SADark!30,dashed] (\x,0) -- (\x,{0.28*(\x-2.4)^2+0.4});
  }
  \node[below,font=\small] at (0.8,0) {$x_1$};
  \node[below,font=\small] at (4.4,0) {$x_2$};
  \draw[SAGreenText,thick] (0.8,1.1168) -- (4.4,1.52);
  \draw[SADark!40,dashed] (2.6,0) -- (2.6,1.3184);
  \fill[SAGreenText] (2.6,1.3184) circle (1.6pt);
  \node[SAGreenText,font=\small,anchor=south east] at (2.55,1.35) {$\E[\phi(X)]$};
  \fill[SARed] (2.6,0.4112) circle (1.6pt);
  \node[SARed,font=\small,anchor=north east] at (2.55,0.40) {$\phi(\E X)$};
  \node[below,font=\small] at (2.6,0) {$m=\E X$};
\end{tikzpicture}
\caption{Jensen 不等式: $X$ 以等概率取两点时, $\E[\phi(X)]$ 是弦的中点高度, $\phi(\E X)$ 在曲线上, 位于弦的下方. 一般情形用支撑直线代替弦}
\label{fig:jensen}
\end{figure}

\begin{example}\label{ex:jensen}
取 $\phi(x)=x^2$ 得 $(\E X)^2\le\E[X^2]$ (即 $\Var(X)\ge0$); 取 $\phi(x)=e^x$ 得 $e^{\E X}\le\E[e^X]$; 取 $\phi(x)=\abs{x}$ 得 $\abs{\E X}\le\E\abs{X}$. 取 $I=(0,\infty)$, $\phi(x)=-\log x$, 对以概率 $\frac{1}{n}$ 取值 $x_1,\dots,x_n>0$ 的 $X$, 得到算术-几何平均不等式 $\sqrt[n]{x_1\cdots x_n}\le\frac{x_1+\cdots+x_n}{n}$.
\end{example}

\begin{lemma}[Young 不等式]\label{lem:young}
设 $p,q>1$, $\frac{1}{p}+\frac{1}{q}=1$, $a,b\ge0$. 则 $ab\le\frac{a^p}{p}+\frac{b^q}{q}$.
\end{lemma}

\begin{proof}
若 $a=0$ 或 $b=0$, 结论显然. 否则, 由 $\log$ 的凹性 (即 $-\log$ 的凸性),
\[
\log\Big(\frac{1}{p}a^p+\frac{1}{q}b^q\Big)\ge\frac{1}{p}\log a^p+\frac{1}{q}\log b^q=\log(ab).
\]
由 $\log$ 的单调性即得.
\end{proof}

\begin{theorem}[H\"older 不等式]\label{thm:holder}
设 $p,q>1$, $\frac{1}{p}+\frac{1}{q}=1$. 则 $\E\abs{XY}\le\norm{X}_p\norm{Y}_q$. 特别地, 取 $p=q=2$ 得 \textbf{Cauchy-Schwarz 不等式}
\[
\E\abs{XY}\le\sqrt{\E[X^2]}\sqrt{\E[Y^2]}.
\]
\end{theorem}

\begin{proof}
若 $\norm{X}_p=0$, 则由推论~\ref{cor:nonneg-properties} (3), $X=0$ $\as$, 两边都是 $0$; $\norm{Y}_q=0$ 同理. 若 $\norm{X}_p$ 或 $\norm{Y}_q$ 为 $+\infty$, 结论显然. 下设二者都是正有限数. 令 $U=\frac{\abs{X}}{\norm{X}_p}$, $V=\frac{\abs{Y}}{\norm{Y}_q}$, 则 $\E[U^p]=\E[V^q]=1$. 由 Young 不等式~\ref{lem:young}, $UV\le\frac{U^p}{p}+\frac{V^q}{q}$, 取期望,
\[
\frac{\E\abs{XY}}{\norm{X}_p\norm{Y}_q}=\E[UV]\le\frac{1}{p}+\frac{1}{q}=1.\qedhere
\]
\end{proof}

\begin{theorem}[Minkowski 不等式]\label{thm:minkowski}
设 $p\ge1$, $X,Y\in L^p$. 则 $X+Y\in L^p$, 且 $\norm{X+Y}_p\le\norm{X}_p+\norm{Y}_p$.
\end{theorem}

\begin{proof}
$p=1$ 时由 $\abs{X+Y}\le\abs{X}+\abs{Y}$ 直接得到. 设 $p>1$. 由 $x\mapsto\abs{x}^p$ 的凸性, $\big|\frac{x+y}{2}\big|^p\le\frac{\abs{x}^p+\abs{y}^p}{2}$, 即 $\abs{x+y}^p\le2^{p-1}(\abs{x}^p+\abs{y}^p)$, 所以 $X+Y\in L^p$. 令 $q=\frac{p}{p-1}$, 则 $\frac{1}{p}+\frac{1}{q}=1$ 且 $(p-1)q=p$. 由三角不等式与 H\"older 不等式~\ref{thm:holder},
\begin{align*}
\E\abs{X+Y}^p&\le\E\big[\abs{X}\abs{X+Y}^{p-1}\big]+\E\big[\abs{Y}\abs{X+Y}^{p-1}\big]\\
&\le\big(\norm{X}_p+\norm{Y}_p\big)\big(\E\abs{X+Y}^{(p-1)q}\big)^{1/q}\\
&=\big(\norm{X}_p+\norm{Y}_p\big)\big(\E\abs{X+Y}^p\big)^{1-1/p}.
\end{align*}
若 $\E\abs{X+Y}^p=0$, 结论显然; 否则两边除以 $(\E\abs{X+Y}^p)^{1-1/p}$ (有限正数) 即得.
\end{proof}

\begin{corollary}[Lyapunov 不等式]\label{cor:lyapunov}
若 $0<p<q$, 则 $\norm{X}_p\le\norm{X}_q$.
\end{corollary}

\begin{proof}
若 $\norm{X}_q=\infty$, 结论显然. 否则由前面的讨论 $X\in L^p$. 对非负随机变量 $\abs{X}^p$ 与凸函数 $\phi(x)=x^{q/p}$ ($x\ge0$) 使用 Jensen 不等式~\ref{thm:jensen},
$(\E\abs{X}^p)^{q/p}\le\E\abs{X}^q$, 两边开 $q$ 次方即得.
\end{proof}

\begin{intuition}
Minkowski 不等式说 $\norm{\cdot}_p$ ($p\ge1$) 满足三角不等式, 因而 $L^p$ 是赋范线性空间 (模去几乎必然相等, 见 1.6 节); Lyapunov 不等式说在概率空间上 $L^q\subseteq L^p$ ($p<q$), 高阶矩有限蕴含低阶矩有限. 注意后者依赖于 $\P(\Omega)=1$: 在 $(\R,\lambda)$ 上包含关系并不成立.
\end{intuition}

% ---------------------------------------------------------------------
\subsection{矩母函数与特征函数}

矩把分布压缩成一列数. 把它们 ``打包'' 成一个函数, 就得到矩母函数; 为了让这个函数总是存在, 把实指数换成虚指数, 就得到特征函数. 特征函数是研究独立随机变量之和 (1.5 节) 与分布收敛 (1.4 节) 的主要工具.

\begin{definition}[矩母函数]\label{def:mgf}
随机变量 $X$ 的\textbf{矩母函数}定义为 $M_X(t)=\E[e^{tX}]\in(0,+\infty]$, $t\in\R$.
\end{definition}

\begin{theorem}\label{thm:mgf-moments}
若存在 $t_0>0$, 使得 $M_X(t)<\infty$ 对所有 $\abs{t}<t_0$ 成立, 则 $X$ 的各阶矩都有限, 且
\[
M_X(t)=\sum_{k=0}^\infty\frac{t^k}{k!}\E[X^k],\qquad \abs{t}<t_0.
\]
特别地, $\E[X^k]=M_X^{(k)}(0)$.
\end{theorem}

\begin{proof}
固定 $t$ 满足 $0<\abs{t}<t_0$. 由于 $e^{\abs{tX}}\le e^{tX}+e^{-tX}$, 而 $\pm t$ 都在 $(-t_0,t_0)$ 中, 所以 $e^{\abs{tX}}\in L^1$. 对每个 $k$, $\frac{\abs{tX}^k}{k!}\le e^{\abs{tX}}$, 故 $\E\abs{X}^k\le k!\,\abs{t}^{-k}\E[e^{\abs{tX}}]<\infty$. 部分和满足
\[
\Big|\sum_{k=0}^n\frac{(tX)^k}{k!}\Big|\le\sum_{k=0}^n\frac{\abs{tX}^k}{k!}\le e^{\abs{tX}},
\]
且当 $n\to\infty$ 时逐点收敛到 $e^{tX}$. 由控制收敛定理~\ref{thm:dct} 与有限线性,
\[
M_X(t)=\lim_{n\to\infty}\E\Big[\sum_{k=0}^n\frac{(tX)^k}{k!}\Big]=\lim_{n\to\infty}\sum_{k=0}^n\frac{t^k}{k!}\E[X^k].
\]
因此幂级数 $\sum_k\frac{t^k}{k!}\E[X^k]$ 在 $(-t_0,t_0)$ 上收敛于 $M_X(t)$; 由幂级数可在收敛区间内逐项求导, $M_X^{(k)}(0)=\E[X^k]$.
\end{proof}

\begin{example}[正态分布的矩母函数与矩]\label{ex:normal-mgf}
设 $Z\sim N(0,1)$. 配方 $tx-\frac{x^2}{2}=\frac{t^2}{2}-\frac{(x-t)^2}{2}$, 由变量替换公式~\ref{thm:change-of-variables},
\[
M_Z(t)=\int_{\R}e^{tx}\varphi(x)\dd x=e^{t^2/2}\int_{\R}\varphi(x-t)\dd x=e^{t^2/2},\qquad t\in\R.
\]
对 $X=\mu+\sigma Z$, $M_X(t)=e^{\mu t}M_Z(\sigma t)=\exp\big(\mu t+\frac{1}{2}\sigma^2t^2\big)$. 把 $e^{t^2/2}=\sum_{k\ge0}\frac{t^{2k}}{2^kk!}$ 与定理~\ref{thm:mgf-moments} 中的级数比较系数:
\[
\E[Z^{2k+1}]=0,\qquad \E[Z^{2k}]=\frac{(2k)!}{2^kk!}=(2k-1)!!=1\cdot3\cdot5\cdots(2k-1).
\]
例如 $\E[Z^2]=1$, $\E[Z^4]=3$, $\E[Z^6]=15$. 第~4~章计算布朗运动的二次变差时将用到 $\E[Z^4]=3$.
\end{example}

\begin{example}[Chernoff 界]\label{ex:chernoff}
对任意 $t>0$, 与 Markov 不等式~\ref{thm:markov} 的证明相同, 由逐点不等式 $e^{ta}\1_{\{X\ge a\}}\le e^{tX}$ 取期望,
$\P(X\ge a)\le e^{-ta}M_X(t)$. 对 $Z\sim N(0,1)$, 得 $\P(Z\ge a)\le\exp\big(-ta+\frac{t^2}{2}\big)$, 取 $t=a$ 使右端最小,
\[
\P(Z\ge a)\le e^{-a^2/2},\qquad a>0.
\]
这个界比定理~\ref{thm:mills} 粗糙 (少了因子 $\frac{1}{\sqrt{2\pi}a}$), 但它的推导只用到矩母函数, 因而适用于任何矩母函数满足 $M_X(t)\le e^{t^2/2}$ 的随机变量 (例如以概率 $\frac{1}{2}$ 取 $\pm1$ 的随机变量).
\end{example}

矩母函数的缺点是它可能不存在: 例如参数为 $\lambda$ 的指数分布, 当 $t\ge\lambda$ 时 $M_X(t)=\int_0^\infty e^{tx}\lambda e^{-\lambda x}\dd x=+\infty$. 改用虚指数可以完全避免这个问题. 为此先约定复值随机变量的期望: 若 $W=U+\mathrm{i}V$, $U,V\in L^1$, 定义 $\E[W]=\E[U]+\mathrm{i}\,\E[V]$. 期望的线性对复系数仍然成立; 控制收敛定理与期望号下求导对实部、虚部分别使用即可推广到复值情形.

\begin{lemma}\label{lem:complex-abs}
对复值可积随机变量 $W$, $\abs{\E W}\le\E\abs{W}$.
\end{lemma}

\begin{proof}
若 $\E W=0$, 结论显然. 否则写 $\E W=re^{\mathrm{i}\theta}$, $r>0$. 则
\[
\abs{\E W}=r=e^{-\mathrm{i}\theta}\E W=\E\big[e^{-\mathrm{i}\theta}W\big]=\E\big[\mathrm{Re}(e^{-\mathrm{i}\theta}W)\big]\le\E\abs{W},
\]
其中第三个等号是因为左边为实数, 故右边的虚部期望为零; 最后一步用到 $\mathrm{Re}\,w\le\abs{w}$.
\end{proof}

\begin{definition}[特征函数]\label{def:cf}
随机变量 $X$ 的\textbf{特征函数}定义为
\[
\varphi_X(t)=\E\big[e^{\mathrm{i}tX}\big]=\E[\cos tX]+\mathrm{i}\,\E[\sin tX],\qquad t\in\R.
\]
\end{definition}

由于 $\abs{e^{\mathrm{i}tX}}=1$, 特征函数对任何随机变量和任何 $t$ 都有定义. 若 $X$ 有密度 $f$, 则 $\varphi_X(t)=\int e^{\mathrm{i}tx}f(x)\dd x$ 就是 $f$ 的 Fourier 变换.

\begin{theorem}[特征函数的基本性质]\label{thm:cf-properties}
\begin{enumerate}
\item $\varphi_X(0)=1$, $\abs{\varphi_X(t)}\le1$, $\varphi_X(-t)=\overline{\varphi_X(t)}$;
\item $\varphi_X$ 在 $\R$ 上一致连续;
\item $\varphi_{aX+b}(t)=e^{\mathrm{i}bt}\varphi_X(at)$;
\item 若 $\E\abs{X}^n<\infty$, 则 $\varphi_X$ 有 $n$ 阶连续导数, 且 $\varphi_X^{(k)}(t)=\E\big[(\mathrm{i}X)^ke^{\mathrm{i}tX}\big]$ ($k\le n$), 特别地 $\varphi_X^{(k)}(0)=\mathrm{i}^k\E[X^k]$.
\end{enumerate}
\end{theorem}

\begin{proof}
(1) $\varphi_X(0)=\E[1]=1$; 由引理~\ref{lem:complex-abs}, $\abs{\varphi_X(t)}\le\E\abs{e^{\mathrm{i}tX}}=1$; $\cos$ 为偶函数、$\sin$ 为奇函数给出第三式.

(2) 对任意 $t,h$,
\[
\abs{\varphi_X(t+h)-\varphi_X(t)}=\big|\E\big[e^{\mathrm{i}tX}(e^{\mathrm{i}hX}-1)\big]\big|\le\E\abs{e^{\mathrm{i}hX}-1}.
\]
右端与 $t$ 无关. 当 $h\to0$ 时, $\abs{e^{\mathrm{i}hX}-1}\to0$ 逐点成立且被常数 $2$ 控制, 由有界收敛定理~\ref{cor:bct}, 右端趋于 $0$.

(3) $\E[e^{\mathrm{i}t(aX+b)}]=e^{\mathrm{i}bt}\E[e^{\mathrm{i}(at)X}]$.

(4) 对 $k$ 归纳. $k=0$ 时即定义. 设 $k<n$ 且 $\varphi_X^{(k)}(t)=\E[h(t,\cdot)]$, 其中 $h(t,\omega)=(\mathrm{i}X)^ke^{\mathrm{i}tX}$. 则 $\partial_t h=(\mathrm{i}X)^{k+1}e^{\mathrm{i}tX}$, $\abs{\partial_t h}\le\abs{X}^{k+1}$, 而由 Lyapunov 不等式~\ref{cor:lyapunov}, $\E\abs{X}^{k+1}<\infty$. 由期望号下求导定理~\ref{thm:diff-under-E} (对实部、虚部分别使用), $\varphi_X^{(k+1)}(t)=\E[(\mathrm{i}X)^{k+1}e^{\mathrm{i}tX}]$. 导数的连续性与 (2) 的证明相同 (用控制收敛定理, 控制函数为 $2\abs{X}^{k}$).
\end{proof}

\begin{example}[正态分布的特征函数]\label{ex:normal-cf}
设 $Z\sim N(0,1)$, 我们证明 $\varphi_Z(t)=e^{-t^2/2}$.

由于 $Z\overset{d}{=}-Z$ (命题~\ref{prop:normal-basic} (1)) 而 $\sin$ 是奇函数, $\E[\sin tZ]=\E[\sin(-tZ)]=-\E[\sin tZ]$, 故 $\E[\sin tZ]=0$, $\varphi_Z(t)=\E[\cos tZ]$. 由于 $\E\abs{Z}<\infty$, 由定理~\ref{thm:cf-properties} (4) 和分部积分 (利用 $\varphi'(x)=-x\varphi(x)$, 注意这里 $\varphi$ 是正态密度),
\begin{align*}
\varphi_Z'(t)&=-\E[Z\sin tZ]=-\int_{\R}\sin(tx)\,x\varphi(x)\dd x\\
&=-\Big[-\sin(tx)\varphi(x)\Big]_{-\infty}^{\infty}-\int_{\R}t\cos(tx)\varphi(x)\dd x=-t\,\varphi_Z(t).
\end{align*}
于是 $\frac{\dd}{\dd t}\big(e^{t^2/2}\varphi_Z(t)\big)=e^{t^2/2}\big(t\varphi_Z(t)+\varphi_Z'(t)\big)=0$, 结合 $\varphi_Z(0)=1$ 得 $\varphi_Z(t)=e^{-t^2/2}$. 由定理~\ref{thm:cf-properties} (3), 对 $X\sim N(\mu,\sigma^2)$ ($\sigma\ge0$, 包括退化情形),
\begin{equation}\label{eq:normal-cf}
\varphi_X(t)=\exp\Big(\mathrm{i}\mu t-\frac{1}{2}\sigma^2t^2\Big).
\end{equation}
\end{example}

\begin{table}[htbp]
\centering
\caption{常用分布的期望、方差与特征函数}
\label{tab:cf}
\renewcommand{\arraystretch}{1.6}
\begin{tabular}{llll}
\toprule
分布 & 期望 & 方差 & 特征函数 $\varphi_X(t)$ \\
\midrule
$\mathrm{B}(1,p)$ & $p$ & $p(1-p)$ & $1-p+pe^{\mathrm{i}t}$ \\
$\mathrm{B}(n,p)$ & $np$ & $np(1-p)$ & $(1-p+pe^{\mathrm{i}t})^n$ \\
$\mathrm{Poi}(\lambda)$ & $\lambda$ & $\lambda$ & $\exp\big(\lambda(e^{\mathrm{i}t}-1)\big)$ \\
$\mathrm{U}(-a,a)$ & $0$ & $\frac{a^2}{3}$ & $\frac{\sin at}{at}$ \\
$\mathrm{Exp}(\lambda)$ & $\frac{1}{\lambda}$ & $\frac{1}{\lambda^2}$ & $\frac{\lambda}{\lambda-\mathrm{i}t}$ \\
$N(\mu,\sigma^2)$ & $\mu$ & $\sigma^2$ & $\exp\big(\mathrm{i}\mu t-\frac{1}{2}\sigma^2t^2\big)$ \\
\bottomrule
\end{tabular}
\end{table}

特征函数的价值在于它完全刻画了分布. 下面两个定理的证明属于 Fourier 分析, 从略.

\begin{theorem}[唯一性定理]\label{thm:cf-uniqueness}
若 $\varphi_X(t)=\varphi_Y(t)$ 对所有 $t\in\R$ 成立, 则 $X\overset{d}{=}Y$. (证明见~\cite{durrett} 定理 3.3.4.)
\end{theorem}

\begin{theorem}[反演公式]\label{thm:inversion}
若 $\int_{\R}\abs{\varphi_X(t)}\dd t<\infty$, 则 $X$ 有有界连续的密度
\[
f(x)=\frac{1}{2\pi}\int_{\R}e^{-\mathrm{i}tx}\varphi_X(t)\dd t.
\]
(证明见~\cite{durrett} 定理 3.3.5.)
\end{theorem}

\begin{intuition}
特征函数把分布一对一地 ``编码'' 成连续函数. 独立随机变量之和的特征函数等于特征函数之积 (1.5 节), 从而把卷积化为乘法; 分布的收敛则等价于特征函数的逐点收敛 (1.4 节的 L\'evy 连续性定理). 例如由 \eqref{eq:normal-cf} 与唯一性定理, ``$X$ 服从正态分布'' 等价于 ``$\varphi_X$ 是二次多项式的指数'', 这将是 1.5 节研究多元正态分布的出发点.
\end{intuition}

% =====================================================================
\section{随机变量序列的收敛}
\label{sec:1-4}

在例~\ref{ex:wlln} 中, 我们证明了抛硬币的频率 $\frac{S_n}{n}$ 偏离 $\frac{1}{2}$ 超过 $\varepsilon$ 的概率趋于零; 强大数定律则断言, 对几乎每个样本点 $\omega$, 数列 $\frac{S_n(\omega)}{n}$ 收敛于 $\frac{1}{2}$; 而中心极限定理说的是 $\frac{S_n-n/2}{\sqrt{n}/2}$ 的\emph{分布}趋于正态分布. 这是三种不同强度的 ``收敛''. 随机分析要同时用到它们: Itô 积分在 $L^2$ 意义下定义 (第~5~章), 布朗运动的轨道性质是几乎必然的 (第~4~章), 二次变差的收敛常在依概率意义下讨论, 而由随机游动得到布朗运动则依赖分布收敛. 本节理清它们之间的关系.

% ---------------------------------------------------------------------
\subsection{四种收敛模式}

\begin{definition}[收敛模式]\label{def:convergence-modes}
设 $X,X_1,X_2,\dots$ 是定义在同一概率空间上的随机变量.
\begin{enumerate}
\item 若 $\P\big(\lim_{n\to\infty}X_n=X\big)=1$, 称 $X_n$ \textbf{几乎必然收敛}到 $X$, 记为 $X_n\xrightarrow{\as}X$.
\item 若对任意 $\varepsilon>0$, $\lim_{n\to\infty}\P(\abs{X_n-X}>\varepsilon)=0$, 称 $X_n$ \textbf{依概率收敛}到 $X$, 记为 $X_n\xrightarrow{\P}X$.
\item 设 $p>0$, $X_n,X\in L^p$. 若 $\lim_{n\to\infty}\E\abs{X_n-X}^p=0$, 称 $X_n$ \textbf{$L^p$ 收敛}到 $X$, 记为 $X_n\xrightarrow{L^p}X$. $p=2$ 时也称\textbf{均方收敛}.
\item 若对 $F_X$ 的每个连续点 $x$, $\lim_{n\to\infty}F_{X_n}(x)=F_X(x)$, 称 $X_n$ \textbf{依分布收敛}到 $X$, 记为 $X_n\xrightarrow{d}X$. 这一定义只涉及分布, 因此 $X_n$ 与 $X$ 可以定义在不同的概率空间上.
\end{enumerate}
\end{definition}

由命题~\ref{prop:rv-operations} (3), $\{\lim_n X_n=X\}=\{\lim_n(X_n-X)\text{ 存在且等于 }0\}$ 是事件, 所以 (1) 有意义.

\begin{remark}
依分布收敛的定义中排除 $F_X$ 的不连续点是必要的. 设 $X_n\equiv\frac{1}{n}$, $X\equiv0$. 无论按哪种意义, 都应该有 $X_n\to X$. 但 $F_{X_n}(0)=\P(\frac{1}{n}\le0)=0$, 而 $F_X(0)=1$: 在 $F_X$ 的跳跃点 $0$ 处, 收敛不成立. 在其余每一点 $x\ne0$ 处, $F_{X_n}(x)\to F_X(x)$ 都成立.
\end{remark}

几乎必然极限与依概率极限在几乎必然相等的意义下是唯一的. 几乎必然收敛的情形由数列极限的唯一性直接得到. 对依概率收敛, 若 $X_n\xrightarrow{\P}X$ 且 $X_n\xrightarrow{\P}Y$, 则对任意 $\varepsilon>0$, 由 $\abs{X-Y}\le\abs{X-X_n}+\abs{X_n-Y}$,
\[
\P(\abs{X-Y}>2\varepsilon)\le\P(\abs{X_n-X}>\varepsilon)+\P(\abs{X_n-Y}>\varepsilon)\to0,
\]
所以 $\P(\abs{X-Y}>2\varepsilon)=0$; 令 $\varepsilon=\frac{1}{m}$, $m\to\infty$, 由概率的连续性得 $\P(X\ne Y)=0$.

下面的引理把几乎必然收敛转化为关于概率的条件, 这是比较几种收敛模式的出发点.

\begin{lemma}[几乎必然收敛的刻画]\label{lem:as-criterion}
以下三条等价:
\begin{enumerate}
\item $X_n\xrightarrow{\as}X$;
\item 对任意 $\varepsilon>0$, $\P(\abs{X_n-X}>\varepsilon\ \text{i.o.})=0$;
\item 对任意 $\varepsilon>0$, $\lim_{n\to\infty}\P\big(\sup_{k\ge n}\abs{X_k-X}>\varepsilon\big)=0$.
\end{enumerate}
\end{lemma}

\begin{proof}
记 $A_n(\varepsilon)=\{\abs{X_n-X}>\varepsilon\}$.

(1)$\Leftrightarrow$(2): 数列 $X_n(\omega)$ 不收敛到 $X(\omega)$, 当且仅当存在 $m$, 使得 $\abs{X_n(\omega)-X(\omega)}>\frac{1}{m}$ 对无穷多个 $n$ 成立. 因此
\[
\{X_n\not\to X\}=\bigcup_{m=1}^\infty\big\{A_n(\tfrac{1}{m})\ \text{i.o.}\big\}.
\]
若 (2) 成立, 右端是可列个零概率集之并, 概率为 $0$, 即 (1) 成立. 反之, 若 (1) 成立, 则对每个 $m$, $\P(A_n(\frac{1}{m})\ \text{i.o.})\le\P(X_n\not\to X)=0$; 对任意 $\varepsilon>0$, 取 $\frac{1}{m}<\varepsilon$, 由 $A_n(\varepsilon)\subseteq A_n(\frac{1}{m})$ 得 (2).

(2)$\Leftrightarrow$(3): $\{\sup_{k\ge n}\abs{X_k-X}>\varepsilon\}=\bigcup_{k\ge n}A_k(\varepsilon)$, 当 $n\to\infty$ 时它单调下降到 $\limsup_k A_k(\varepsilon)=\{A_k(\varepsilon)\ \text{i.o.}\}$. 由上连续性 (定理~\ref{thm:continuity}), (3) 中的极限等于 $\P(A_k(\varepsilon)\ \text{i.o.})$.
\end{proof}

\begin{corollary}[快速收敛判别法]\label{cor:fast-convergence}
若对任意 $\varepsilon>0$, $\sum_{n=1}^\infty\P(\abs{X_n-X}>\varepsilon)<\infty$, 则 $X_n\xrightarrow{\as}X$.
\end{corollary}

\begin{proof}
由 Borel-Cantelli 第一引理 (定理~\ref{thm:bc1}), 引理~\ref{lem:as-criterion} 的条件 (2) 成立.
\end{proof}

% ---------------------------------------------------------------------
\subsection{蕴含关系与反例}

\begin{theorem}[收敛模式的蕴含关系]\label{thm:convergence-implications}
\begin{enumerate}
\item $X_n\xrightarrow{\as}X\ \Longrightarrow\ X_n\xrightarrow{\P}X$;
\item 对 $p>0$, $X_n\xrightarrow{L^p}X\ \Longrightarrow\ X_n\xrightarrow{\P}X$;
\item 对 $0<p<q$, $X_n\xrightarrow{L^q}X\ \Longrightarrow\ X_n\xrightarrow{L^p}X$;
\item $X_n\xrightarrow{\P}X\ \Longrightarrow\ X_n\xrightarrow{d}X$.
\end{enumerate}
\end{theorem}

\begin{proof}
(1) $\P(\abs{X_n-X}>\varepsilon)\le\P(\sup_{k\ge n}\abs{X_k-X}>\varepsilon)$, 由引理~\ref{lem:as-criterion} (3), 右端趋于 $0$.

(2) 由 Markov 不等式~\ref{thm:markov}, $\P(\abs{X_n-X}>\varepsilon)\le\frac{\E\abs{X_n-X}^p}{\varepsilon^p}\to0$.

(3) 由 Lyapunov 不等式~\ref{cor:lyapunov}, $\norm{X_n-X}_p\le\norm{X_n-X}_q\to0$.

(4) 记 $F_n=F_{X_n}$, $F=F_X$, 设 $x$ 是 $F$ 的连续点, $\varepsilon>0$. 若 $X_n\le x$, 则或者 $X\le x+\varepsilon$, 或者 $\abs{X_n-X}>\varepsilon$; 若 $X\le x-\varepsilon$, 则或者 $X_n\le x$, 或者 $\abs{X_n-X}>\varepsilon$. 即
\[
\{X_n\le x\}\subseteq\{X\le x+\varepsilon\}\cup\{\abs{X_n-X}>\varepsilon\},\qquad
\{X\le x-\varepsilon\}\subseteq\{X_n\le x\}\cup\{\abs{X_n-X}>\varepsilon\}.
\]
由单调性与次可加性,
\[
F(x-\varepsilon)-\P(\abs{X_n-X}>\varepsilon)\le F_n(x)\le F(x+\varepsilon)+\P(\abs{X_n-X}>\varepsilon).
\]
令 $n\to\infty$ 得 $F(x-\varepsilon)\le\liminf_n F_n(x)\le\limsup_n F_n(x)\le F(x+\varepsilon)$. 再令 $\varepsilon\downarrow0$, 由 $F$ 在 $x$ 处连续得 $F_n(x)\to F(x)$.
\end{proof}

反方向的蕴含一般都不成立.

\begin{example}[反例]\label{ex:convergence-counterexamples}
以下 (1)(2)(3) 都取 $\Omega=[0,1]$, $\P$ 为均匀分布 (Lebesgue 测度).
\begin{enumerate}
\item \textbf{依概率收敛 (甚至 $L^p$ 收敛) 不蕴含几乎必然收敛: ``打字机序列''.} 每个正整数 $n$ 可唯一地写成 $n=2^k+j$, $k\ge0$, $0\le j<2^k$. 令
\[
X_n=\1_{[j2^{-k},\,(j+1)2^{-k}]}.
\]
当 $n$ 从 $2^k$ 增加到 $2^{k+1}-1$ 时, 长为 $2^{-k}$ 的区间从左到右 ``扫过'' 整个 $[0,1]$ (图~\ref{fig:typewriter}). 由于 $n\to\infty$ 时 $k\to\infty$, 对任意 $\varepsilon\in(0,1)$ 和 $p>0$,
\[
\P(\abs{X_n}>\varepsilon)=\E\abs{X_n}^p=2^{-k}\to0,
\]
所以 $X_n\xrightarrow{\P}0$ 且 $X_n\xrightarrow{L^p}0$. 但对每个 $\omega\in[0,1]$, 每一层 $k$ 中都至少有一个区间包含 $\omega$, 所以 $X_n(\omega)=1$ 对无穷多个 $n$ 成立; 而当 $k\ge1$ 时每层又至少有一个区间不含 $\omega$, 所以 $X_n(\omega)=0$ 也对无穷多个 $n$ 成立. 于是 $X_n(\omega)$ 对\emph{每个} $\omega$ 都不收敛.
\item \textbf{几乎必然收敛不蕴含 $L^1$ 收敛.} 例~\ref{ex:mass-escape} 中的 $X_n=n\1_{(0,1/n)}$ 处处收敛到 $0$, 但 $\E\abs{X_n-0}=1$.
\item \textbf{$L^p$ 收敛不蕴含 $L^q$ 收敛 ($p<q$).} 令 $X_n=n^{1/q}\1_{(0,1/n)}$, 则 $\E\abs{X_n}^p=n^{p/q-1}\to0$, 但 $\E\abs{X_n}^q=1$.
\item \textbf{依分布收敛不蕴含依概率收敛.} 设 $X\sim N(0,1)$, 令 $X_n=-X$. 由对称性, $X_n\overset{d}{=}X$, 所以 $X_n\xrightarrow{d}X$. 但 $\P(\abs{X_n-X}>\varepsilon)=\P(\abs{X}>\frac{\varepsilon}{2})$ 是与 $n$ 无关的正数. 依分布收敛只比较分布, 完全不关心 $X_n$ 与 $X$ 在同一个 $\omega$ 上的取值是否接近.
\end{enumerate}
\end{example}

\begin{figure}[htbp]
\centering
\begin{tikzpicture}[xscale=7,yscale=0.42]
  \draw[->,SADark!80] (0,0) -- (1.08,0) node[right,font=\small] {$\omega$};
  \node[below,font=\small] at (0,0) {$0$};
  \node[below,font=\small] at (1,0) {$1$};
  \foreach \n/\a/\b/\c in {1/0/1/SADark, 2/0/0.5/SABlue, 3/0.5/1/SABlue, 4/0/0.25/SARed, 5/0.25/0.5/SARed, 6/0.5/0.75/SARed, 7/0.75/1/SARed, 8/0/0.125/SAGreenText, 9/0.125/0.25/SAGreenText, 10/0.25/0.375/SAGreenText, 11/0.375/0.5/SAGreenText, 12/0.5/0.625/SAGreenText, 13/0.625/0.75/SAGreenText, 14/0.75/0.875/SAGreenText, 15/0.875/1/SAGreenText} {
    \pgfmathsetmacro{\y}{16-\n}
    \draw[SADark!12] (0,\y) -- (1,\y);
    \draw[\c,line width=2.6pt] (\a,\y) -- (\b,\y);
    \node[left,font=\scriptsize] at (-0.01,\y) {$X_{\n}$};
  }
  \draw[SAYellow!90!black,thick,dashed] (0.3,0) -- (0.3,15.6);
  \node[above,font=\small] at (0.3,15.6) {$\omega=0.3$};
  \node[font=\small,align=left,anchor=west] at (1.05,10) {$k=0$: \ $X_1$\\ $k=1$: \ $X_2,X_3$\\ $k=2$: \ $X_4,\dots,X_7$\\ $k=3$: \ $X_8,\dots,X_{15}$};
\end{tikzpicture}
\caption{打字机序列: 粗线段标出 $X_n=1$ 的区间. 区间越来越短, 所以 $X_n\xrightarrow{\P}0$; 但每一层都会扫过任意给定的 $\omega$ (虚线), 所以 $X_n(\omega)$ 无穷多次等于 $1$}
\label{fig:typewriter}
\end{figure}

尽管反向蕴含不成立, 但在附加条件下可以部分地恢复. 下面的定理中, (1) 是处理依概率收敛最常用的技巧.

\begin{theorem}[部分逆命题]\label{thm:partial-converses}
\begin{enumerate}
\item 若 $X_n\xrightarrow{\P}X$, 则存在子列 $\{X_{n_k}\}$ 使 $X_{n_k}\xrightarrow{\as}X$.
\item $X_n\xrightarrow{\P}X$, 当且仅当 $\{X_n\}$ 的每个子列都有进一步的子列几乎必然收敛到 $X$.
\item 若 $X_n\xrightarrow{d}c$, 其中 $c$ 是常数, 则 $X_n\xrightarrow{\P}c$.
\item 若 $X_n\xrightarrow{\P}X$, 且存在 $Y\in L^1$ 使 $\abs{X_n}\le Y$, 则 $X_n\xrightarrow{L^1}X$.
\end{enumerate}
\end{theorem}

\begin{proof}
(1) 由依概率收敛, 对每个 $k$ 可以取 $n_k$ (并使 $n_1<n_2<\cdots$) 满足 $\P(\abs{X_{n_k}-X}>2^{-k})\le2^{-k}$. 由于 $\sum_k2^{-k}<\infty$, 由 Borel-Cantelli 第一引理 (定理~\ref{thm:bc1}), 几乎必然地, 只有有限个 $k$ 使 $\abs{X_{n_k}-X}>2^{-k}$, 即当 $k$ 充分大时 $\abs{X_{n_k}-X}\le2^{-k}\to0$.

(2) ``仅当'': 子列仍依概率收敛到 $X$, 对它用 (1). ``当'': 反设 $X_n$ 不依概率收敛到 $X$, 则存在 $\varepsilon>0$、$\delta>0$ 及子列 $\{n_k\}$ 使 $\P(\abs{X_{n_k}-X}>\varepsilon)\ge\delta$ 对所有 $k$ 成立. 由假设, 它有进一步的子列几乎必然收敛到 $X$, 由定理~\ref{thm:convergence-implications} (1), 该子列也依概率收敛到 $X$, 与 ``概率始终 $\ge\delta$'' 矛盾.

(3) 常数 $c$ 的分布函数 $F=\1_{[c,\infty)}$ 只在 $c$ 处不连续. 对 $\varepsilon>0$, $c-\varepsilon$ 与 $c+\varepsilon$ 都是连续点, 故
\[
\P(\abs{X_n-c}>\varepsilon)\le F_{X_n}(c-\varepsilon)+1-F_{X_n}(c+\varepsilon)\to F(c-\varepsilon)+1-F(c+\varepsilon)=0+1-1=0.
\]

(4) 先证 $\abs{X}\le Y$ $\as$: 由 (1), 取子列 $X_{n_k}\xrightarrow{\as}X$, 在 $\abs{X_{n_k}}\le Y$ 中取极限即得. 于是 $X\in L^1$. 反设 $\E\abs{X_n-X}\not\to0$, 则存在 $\delta>0$ 及子列使 $\E\abs{X_{n_k}-X}\ge\delta$. 由 (2), 它有进一步的子列几乎必然收敛到 $X$; 对这个子列用控制收敛定理~\ref{thm:dct} (控制函数为 $Y$), 其 $\E\abs{\cdot-X}\to0$, 矛盾.
\end{proof}

\begin{intuition}
(2) 体现了一个一般原则: 实数列 $a_n\to a$ 当且仅当每个子列都有进一步的子列收敛到 $a$. 借助它, 关于几乎必然收敛的定理 (如控制收敛定理) 常常可以 ``免费'' 推广到依概率收敛. 例如: 若 $X_n\xrightarrow{\P}X$, $g$ 是连续函数, 则 $g(X_n)\xrightarrow{\P}g(X)$ (对几乎必然收敛的子列, 由 $g$ 的连续性 $g(X_{n_k})\to g(X)$ $\as$, 再由 (2)); 若 $X_n\xrightarrow{\P}X$, $Y_n\xrightarrow{\P}Y$, 则 $X_n+Y_n\xrightarrow{\P}X+Y$, $X_nY_n\xrightarrow{\P}XY$ (同理, 先取公共的几乎必然收敛子列).
\end{intuition}

\begin{figure}[htbp]
\centering
\begin{tikzpicture}[
  box/.style={draw=SADark,rounded corners=3pt,fill=SABlue!8,minimum width=1.9cm,minimum height=0.8cm,font=\small},
  imp/.style={->,thick,SADark},
  cond/.style={->,thick,dashed,SARed},
  lab/.style={font=\scriptsize,SARed,inner sep=2pt}]
  \node[box] (as) at (0,0) {$\xrightarrow{\as}$};
  \node[box] (P) at (4,0) {$\xrightarrow{\P}$};
  \node[box] (d) at (8,0) {$\xrightarrow{d}$};
  \node[box] (Lp) at (4,-2.4) {$\xrightarrow{L^p}$};
  \node[box] (Lq) at (0,-2.4) {$\xrightarrow{L^q}$};
  \draw[imp] (as) -- (P);
  \draw[imp] (P) -- (d);
  \draw[imp] (Lq) -- node[below,font=\scriptsize,SADark] {$p<q$} (Lp);
  \draw[imp] ([xshift=-4mm]Lp.north) -- ([xshift=-4mm]P.south);
  \draw[cond] ([xshift=4mm]P.south) -- node[lab,right] {一致可积 ($p=1$)} ([xshift=4mm]Lp.north);
  \draw[cond] (P.north) to[bend right=40] node[lab,above] {子列} (as.north);
  \draw[cond] (d.north) to[bend right=40] node[lab,above] {极限为常数} (P.north east);
\end{tikzpicture}
\caption{收敛模式之间的关系. 实线: 总成立; 虚线: 在所标条件下成立}
\label{fig:convergence-diagram}
\end{figure}

% ---------------------------------------------------------------------
\subsection{一致可积性与 Vitali 收敛定理}

例~\ref{ex:mass-escape} 说明, 依概率收敛 (甚至处处收敛) 加上 $L^1$ 有界性仍不足以推出期望收敛, 因为质量可以逃逸到小概率集合上. 控制收敛定理用可积控制排除了逃逸, 但这个条件过强: 许多重要的序列 (例如第~3~章的鞅) 并没有可积控制. 一致可积性恰好刻画了 ``没有质量逃逸''.

\begin{definition}[一致可积]\label{def:ui}
随机变量族 $\{X_i\}_{i\in I}$ 称为\textbf{一致可积} (uniformly integrable, UI) 的, 如果
\[
\lim_{K\to\infty}\sup_{i\in I}\E\big[\abs{X_i};\abs{X_i}>K\big]=0.
\]
\end{definition}

\begin{lemma}\label{lem:single-integrable}
设 $X\in L^1$. 则
\begin{enumerate}
\item $\lim_{K\to\infty}\E[\abs{X};\abs{X}>K]=0$;
\item 对任意 $\varepsilon>0$, 存在 $\delta>0$, 使得只要 $\P(A)<\delta$, 就有 $\E[\abs{X};A]<\varepsilon$.
\end{enumerate}
\end{lemma}

\begin{proof}
(1) 由推论~\ref{cor:nonneg-properties} (4), $\abs{X}<\infty$ $\as$, 所以当 $K\to\infty$ 时 $\abs{X}\1_{\{\abs{X}>K\}}\to0$ $\as$, 且被 $\abs{X}$ 控制. 由控制收敛定理~\ref{thm:dct} 即得.

(2) 对任意 $K>0$ 和事件 $A$, 把 $A$ 分成 $A\cap\{\abs{X}\le K\}$ 与 $A\cap\{\abs{X}>K\}$,
\[
\E[\abs{X};A]\le K\P(A)+\E[\abs{X};\abs{X}>K].
\]
由 (1), 取 $K$ 使第二项小于 $\frac{\varepsilon}{2}$, 再取 $\delta=\frac{\varepsilon}{2K}$ 即可.
\end{proof}

于是, 单个可积随机变量 (以及有限个可积随机变量组成的族) 总是一致可积的. 一致可积性的要点在于 ``$\sup_i$'': 尾部积分要对整个族\emph{一致地}小.

\begin{proposition}[一致可积的等价刻画]\label{prop:ui-equivalent}
$\{X_i\}_{i\in I}$ 一致可积, 当且仅当以下两条同时成立:
\begin{enumerate}
\item $L^1$ 有界: $\sup_i\E\abs{X_i}<\infty$;
\item 一致绝对连续: 对任意 $\varepsilon>0$, 存在 $\delta>0$, 使得只要 $\P(A)<\delta$, 就有 $\sup_i\E[\abs{X_i};A]<\varepsilon$.
\end{enumerate}
\end{proposition}

\begin{proof}
``仅当'': 取 $K$ 使 $\sup_i\E[\abs{X_i};\abs{X_i}>K]\le1$, 则 $\E\abs{X_i}\le K+1$ 对所有 $i$ 成立, 得 (1). 与引理~\ref{lem:single-integrable} (2) 的证明相同, $\E[\abs{X_i};A]\le K\P(A)+\sup_j\E[\abs{X_j};\abs{X_j}>K]$, 先取 $K$ 再取 $\delta$ 得 (2).

``当'': 记 $C=\sup_i\E\abs{X_i}$. 对给定的 $\varepsilon$, 取 (2) 中的 $\delta$. 由 Markov 不等式~\ref{thm:markov}, 当 $K>\frac{C}{\delta}$ 时, 对所有 $i$, $\P(\abs{X_i}>K)\le\frac{C}{K}<\delta$, 从而 $\E[\abs{X_i};\abs{X_i}>K]<\varepsilon$.
\end{proof}

\begin{proposition}[一致可积的充分条件]\label{prop:ui-sufficient}
以下任一条件都蕴含 $\{X_i\}_{i\in I}$ 一致可积:
\begin{enumerate}
\item 存在 $Y\in L^1$, 使 $\abs{X_i}\le Y$ 对所有 $i$ 成立;
\item 存在 $p>1$, 使 $\sup_i\E\abs{X_i}^p<\infty$;
\item (de la Vall\'ee-Poussin 判别法) 存在非负函数 $G$ 满足 $\lim_{x\to\infty}\frac{G(x)}{x}=+\infty$, 且 $\sup_i\E[G(\abs{X_i})]<\infty$.
\end{enumerate}
\end{proposition}

\begin{proof}
(1) $\{\abs{X_i}>K\}\subseteq\{Y>K\}$, 所以 $\E[\abs{X_i};\abs{X_i}>K]\le\E[Y;Y>K]$, 由引理~\ref{lem:single-integrable} (1), 右端趋于 $0$ 且与 $i$ 无关.

(2) 是 (3) 取 $G(x)=x^p$ 的特例.

(3) 记 $C=\sup_i\E[G(\abs{X_i})]$, $\eta(K)=\sup_{x>K}\frac{x}{G(x)}$, 由假设 $\eta(K)\to0$ ($K\to\infty$). 在 $\{\abs{X_i}>K\}$ 上, $\abs{X_i}\le\eta(K)G(\abs{X_i})$, 所以
$\E[\abs{X_i};\abs{X_i}>K]\le\eta(K)\,\E[G(\abs{X_i})]\le C\eta(K)\to0$.
\end{proof}

\begin{example}
例~\ref{ex:mass-escape} 中 $X_n=n\1_{(0,1/n)}$ 满足 $\E\abs{X_n}=1$, 是 $L^1$ 有界的, 但对任意 $K$, 当 $n>K$ 时 $\E[X_n;X_n>K]=1$, 所以 $\sup_n\E[X_n;X_n>K]=1$, 不是一致可积的. 这说明 $L^1$ 有界性不能代替一致可积性; 而由命题~\ref{prop:ui-sufficient} (2), 只要稍稍加强为 $L^p$ 有界 ($p>1$) 就足够了.
\end{example}

\begin{theorem}[Vitali 收敛定理]\label{thm:vitali}
设 $X_n\xrightarrow{\P}X$. 则以下两条等价:
\begin{enumerate}
\item $\{X_n\}$ 一致可积;
\item $X\in L^1$ 且 $X_n\xrightarrow{L^1}X$.
\end{enumerate}
此时 $\E[X_n]\to\E[X]$.
\end{theorem}

\begin{proof}
(1)$\Rightarrow$(2). 由命题~\ref{prop:ui-equivalent}, $C=\sup_n\E\abs{X_n}<\infty$. 由定理~\ref{thm:partial-converses} (1), 取子列 $X_{n_k}\xrightarrow{\as}X$, 由 Fatou 引理~\ref{thm:fatou},
\[
\E\abs{X}=\E\Big[\liminf_{k\to\infty}\abs{X_{n_k}}\Big]\le\liminf_{k\to\infty}\E\abs{X_{n_k}}\le C<\infty.
\]
固定 $\varepsilon>0$, 令 $A_n=\{\abs{X_n-X}>\varepsilon\}$, 则 $\P(A_n)\to0$. 由三角不等式,
\begin{align*}
\E\abs{X_n-X}&=\E\big[\abs{X_n-X};A_n^c\big]+\E\big[\abs{X_n-X};A_n\big]\\
&\le\varepsilon+\E[\abs{X_n};A_n]+\E[\abs{X};A_n].
\end{align*}
由命题~\ref{prop:ui-equivalent} (2) (对 $\{X_n\}$) 与引理~\ref{lem:single-integrable} (2) (对 $X$), 存在 $\delta>0$, 使得 $\P(A)<\delta$ 时 $\sup_n\E[\abs{X_n};A]<\varepsilon$ 且 $\E[\abs{X};A]<\varepsilon$. 当 $n$ 充分大时 $\P(A_n)<\delta$, 从而 $\E\abs{X_n-X}<3\varepsilon$. 由 $\varepsilon$ 的任意性, $X_n\xrightarrow{L^1}X$.

(2)$\Rightarrow$(1). 我们验证命题~\ref{prop:ui-equivalent} 中的两个条件. (1) 由 $\E\abs{X_n}\le\E\abs{X_n-X}+\E\abs{X}$, 且收敛数列有界, 得 $\sup_n\E\abs{X_n}<\infty$. (2) 给定 $\varepsilon>0$, 取 $N$ 使 $n>N$ 时 $\E\abs{X_n-X}<\frac{\varepsilon}{2}$. 对有限个可积随机变量 $X,X_1,\dots,X_N$ 分别使用引理~\ref{lem:single-integrable} (2), 取其中最小的 $\delta$, 使得 $\P(A)<\delta$ 时 $\E[\abs{X};A]<\frac{\varepsilon}{2}$ 且 $\E[\abs{X_n};A]<\varepsilon$ ($n\le N$). 于是当 $\P(A)<\delta$ 时, 对 $n>N$,
\[
\E[\abs{X_n};A]\le\E\abs{X_n-X}+\E[\abs{X};A]<\varepsilon,
\]
对 $n\le N$ 也有 $\E[\abs{X_n};A]<\varepsilon$.

最后, $\abs{\E X_n-\E X}\le\E\abs{X_n-X}\to0$.
\end{proof}

\begin{intuition}
Vitali 定理把 ``何时可以交换极限与期望'' 的问题彻底解决了: 在依概率收敛的前提下, 一致可积是 $L^1$ 收敛的\emph{充分必要}条件. 控制收敛定理只是它的特例 (命题~\ref{prop:ui-sufficient} (1)). 在第~3~章中, 一致可积鞅的收敛定理与可选停时定理都建立在这一结果之上.
\end{intuition}

% ---------------------------------------------------------------------
\subsection{依分布收敛与中心极限定理}

依分布收敛的定义使用分布函数, 只适用于实值随机变量, 也不便于推广. 下面给出一个更好用的等价刻画. 为此先证明一个有独立意义的表示定理: 依分布收敛的序列, 可以在某个概率空间上 ``实现'' 为几乎必然收敛的序列.

\begin{theorem}[Skorokhod 表示定理]\label{thm:skorokhod}
设 $X_n\xrightarrow{d}X$. 则在概率空间 $((0,1),\mathcal{B}((0,1)),\lambda)$ 上存在随机变量 $Y_n,Y$, 使得 $Y_n\overset{d}{=}X_n$, $Y\overset{d}{=}X$, 且 $Y_n\xrightarrow{\as}Y$.
\end{theorem}

\begin{proof}
记 $F_n=F_{X_n}$, $F=F_X$. 在 $\Omega=(0,1)$ 上令 $Y_n(u)=F_n^{-1}(u)$, $Y(u)=F^{-1}(u)$, 其中 $F^{-1}$ 是定理~\ref{thm:quantile} 中的广义逆. 由该定理, $Y_n\overset{d}{=}X_n$, $Y\overset{d}{=}X$. 下面利用单调函数与广义逆函数的性质证明 $Y_n\to Y$ 几乎必然成立:
\begin{itemize}
\item 单调函数至多有可列个不连续点 (每个跳跃区间 $(G(u-),G(u+))$ 中可以选出一个有理数, 不同跳跃点对应的区间不交). 因此 $F$ 的连续点在 $\R$ 中稠密, 而单调不减函数 $G=F^{-1}$ 在 $(0,1)$ 中的不连续点构成 Lebesgue 零测集.
\item 由 \eqref{eq:galois}: $F^{-1}(u)\le y\iff u\le F(y)$, 等价地 $y<F^{-1}(u)\iff F(y)<u$; 对 $F_n$ 同理.
\end{itemize}
设 $u\in(0,1)$ 是 $G$ 的连续点.

\emph{下极限.} 任取 $F$ 的连续点 $y<G(u)$. 则 $F(y)<u$, 而 $F_n(y)\to F(y)$, 所以当 $n$ 充分大时 $F_n(y)<u$, 即 $y<F_n^{-1}(u)$. 故 $\liminf_n F_n^{-1}(u)\ge y$. 由于这样的 $y$ 可以任意接近 $G(u)$, $\liminf_n F_n^{-1}(u)\ge G(u)$.

\emph{上极限.} 任取 $u'\in(u,1)$ 及 $F$ 的连续点 $y>G(u')$. 则 $G(u')\le y$, 即 $F(y)\ge u'>u$. 由 $F_n(y)\to F(y)$, 当 $n$ 充分大时 $F_n(y)\ge u$, 即 $F_n^{-1}(u)\le y$. 故 $\limsup_n F_n^{-1}(u)\le y$. 令 $y\downarrow G(u')$ 得 $\limsup_n F_n^{-1}(u)\le G(u')$, 再令 $u'\downarrow u$, 由 $G$ 在 $u$ 处连续得 $\limsup_n F_n^{-1}(u)\le G(u)$.

所以 $Y_n(u)\to Y(u)$ 在 $G$ 的每个连续点 $u$ 处成立, 不连续点构成零测集.
\end{proof}

\begin{theorem}[依分布收敛的等价刻画]\label{thm:portmanteau}
$X_n\xrightarrow{d}X$, 当且仅当对每个有界连续函数 $g:\R\to\R$, $\E[g(X_n)]\to\E[g(X)]$.
\end{theorem}

\begin{proof}
``仅当'': 取 Skorokhod 表示定理~\ref{thm:skorokhod} 中的 $Y_n,Y$. 由 $g$ 的连续性, $g(Y_n)\to g(Y)$ $\as$; 由有界收敛定理~\ref{cor:bct}, $\E[g(Y_n)]\to\E[g(Y)]$. 由于期望只依赖于分布 (定理~\ref{thm:change-of-variables}), $\E[g(X_n)]=\E[g(Y_n)]$, $\E[g(X)]=\E[g(Y)]$.

``当'': 设 $x$ 是 $F_X$ 的连续点, $\varepsilon>0$. 令连续函数 $g_\varepsilon$ 在 $(-\infty,x]$ 上取 $1$, 在 $[x+\varepsilon,\infty)$ 上取 $0$, 在 $[x,x+\varepsilon]$ 上线性. 则 $\1_{(-\infty,x]}\le g_\varepsilon\le\1_{(-\infty,x+\varepsilon]}$, 于是
\[
\limsup_{n\to\infty}F_{X_n}(x)\le\lim_{n\to\infty}\E[g_\varepsilon(X_n)]=\E[g_\varepsilon(X)]\le F_X(x+\varepsilon).
\]
同理, 用在 $(-\infty,x-\varepsilon]$ 上取 $1$、在 $[x,\infty)$ 上取 $0$ 的连续函数, 得 $\liminf_n F_{X_n}(x)\ge F_X(x-\varepsilon)$. 令 $\varepsilon\downarrow0$, 由 $F_X$ 在 $x$ 处连续即得.
\end{proof}

\begin{remark}
定理~\ref{thm:portmanteau} 给出的刻画不涉及分布函数, 因而可以直接推广到随机向量, 乃至取值于函数空间的随机元 (例如布朗运动的整条轨道), 第~4~章的 Donsker 不变原理就是在这种意义下理解的.
\end{remark}

由于 $\cos(tx)$ 与 $\sin(tx)$ 都是有界连续函数, 定理~\ref{thm:portmanteau} 立即给出: 若 $X_n\xrightarrow{d}X$, 则对每个 $t$, $\varphi_{X_n}(t)\to\varphi_X(t)$. 其逆命题也成立, 这是证明依分布收敛最有力的工具.

\begin{theorem}[L\'evy 连续性定理]\label{thm:levy-continuity}
$X_n\xrightarrow{d}X$, 当且仅当对每个 $t\in\R$, $\varphi_{X_n}(t)\to\varphi_X(t)$. 更进一步, 若 $\varphi_{X_n}(t)$ 对每个 $t$ 收敛到某个函数 $\varphi(t)$, 且 $\varphi$ 在 $0$ 处连续, 则 $\varphi$ 是某个随机变量 $X$ 的特征函数, 且 $X_n\xrightarrow{d}X$. (``当'' 的方向的证明见~\cite{durrett} 定理 3.3.17.)
\end{theorem}

\begin{example}\label{ex:conv-dist}
\begin{enumerate}
\item \textbf{极值的极限分布.} 设 $U_1,U_2,\dots$ 是 $\iid$ 的 $\mathrm{U}(0,1)$ 随机变量, $M_n=\max_{k\le n}U_k$. 对 $x\in[0,n]$, 由独立性 (定理~\ref{thm:indep-cdf}),
\[
\P\big(n(1-M_n)>x\big)=\P\Big(M_n<1-\frac{x}{n}\Big)=\Big(1-\frac{x}{n}\Big)^n\to e^{-x}.
\]
所以 $n(1-M_n)\xrightarrow{d}\mathrm{Exp}(1)$: 最大值与 $1$ 的距离是 $\frac{1}{n}$ 量级的, 经过放大后服从指数分布.
\item \textbf{正态分布的极限.} 设 $X_n\sim N(\mu_n,\sigma_n^2)$, $\mu_n\to\mu$, $\sigma_n^2\to\sigma^2$. 由 \eqref{eq:normal-cf}, $\varphi_{X_n}(t)\to\exp(\mathrm{i}\mu t-\frac{1}{2}\sigma^2t^2)$, 由 L\'evy 连续性定理, $X_n\xrightarrow{d}N(\mu,\sigma^2)$. 1.6 节将证明一个相关的结论: 正态随机变量的 $L^2$ 极限仍是正态的.
\end{enumerate}
\end{example}

最后叙述概率论中最重要的极限定理. 它解释了为什么正态分布在自然界中无处不在.

\begin{theorem}[中心极限定理, Lindeberg-L\'evy]\label{thm:clt}
设 $X_1,X_2,\dots$ 是 $\iid$ 随机变量, $\E[X_1]=\mu$, $\Var(X_1)=\sigma^2\in(0,\infty)$, $S_n=X_1+\cdots+X_n$. 则
\[
\frac{S_n-n\mu}{\sigma\sqrt{n}}\xrightarrow{d}N(0,1).
\]
(证明见~\cite{durrett} 定理 3.4.1.)
\end{theorem}

\begin{intuition}
标准化 $\frac{S_n-n\mu}{\sigma\sqrt{n}}$ 的期望为 $0$、方差为 $1$ (方差的计算见 1.5 节), 定理断言其分布的极限与 $X_1$ 的具体分布无关. 记 $\psi$ 为 $\frac{X_1-\mu}{\sigma}$ 的特征函数, 由定理~\ref{thm:cf-properties} (4), $\psi(s)=1-\frac{s^2}{2}+o(s^2)$; 而独立和的特征函数是乘积 (1.5 节), 所以标准化和的特征函数为
\[
\psi\Big(\frac{t}{\sqrt{n}}\Big)^n=\Big(1-\frac{t^2}{2n}+o\Big(\frac{1}{n}\Big)\Big)^n\to e^{-t^2/2}.
\]
$X_1$ 的分布只通过前两阶矩进入极限, 更高阶的信息都在 $o(\frac{1}{n})$ 中消失了. 图~\ref{fig:clt} 显示, 即使对 $n=16$ 次抛硬币, 标准化后的分布函数已经与 $\Phi$ 相当接近.
\end{intuition}

\begin{figure}[htbp]
\centering
\begin{tikzpicture}
\begin{axis}[SAaxis, width=11cm, height=6cm, xmin=-4.2, xmax=4.4, ymin=0, ymax=1.1,
  xtick={-3,-2,-1,0,1,2,3}, ytick={0.5,1}, xlabel={$x$},
  legend style={draw=none,fill=white,font=\small,at={(0.97,0.35)},anchor=east}]
  \addplot[SARed,thick,const plot] coordinates {(-4.2,0) (-4.000,0.0000) (-3.500,0.0003) (-3.000,0.0021) (-2.500,0.0106) (-2.000,0.0384) (-1.500,0.1051) (-1.000,0.2272) (-0.500,0.4018) (0.000,0.5982) (0.500,0.7728) (1.000,0.8949) (1.500,0.9616) (2.000,0.9894) (2.500,0.9979) (3.000,0.9997) (3.500,1.0000) (4.000,1.0000) (4.3,1.0)};
  \addlegendentry{$\P\big(\frac{S_{16}-8}{2}\le x\big)$}
  \addplot[SABlue,very thick,smooth] coordinates {(-4.2,0.0000) (-3.50,0.0002) (-3.30,0.0005) (-3.10,0.0010) (-2.90,0.0019) (-2.70,0.0035) (-2.50,0.0062) (-2.30,0.0107) (-2.10,0.0179) (-1.90,0.0287) (-1.70,0.0446) (-1.50,0.0668) (-1.30,0.0968) (-1.10,0.1357) (-0.90,0.1841) (-0.70,0.2420) (-0.50,0.3085) (-0.30,0.3821) (-0.10,0.4602) (0.10,0.5398) (0.30,0.6179) (0.50,0.6915) (0.70,0.7580) (0.90,0.8159) (1.10,0.8643) (1.30,0.9032) (1.50,0.9332) (1.70,0.9554) (1.90,0.9713) (2.10,0.9821) (2.30,0.9893) (2.50,0.9938) (2.70,0.9965) (2.90,0.9981) (3.10,0.9990) (3.30,0.9995) (3.50,0.9998) (4.3,1.0000)};
  \addlegendentry{$\Phi(x)$}
\end{axis}
\end{tikzpicture}
\caption{中心极限定理: 抛 $16$ 次均匀硬币, 正面次数标准化后的分布函数 (阶梯) 与标准正态分布函数 $\Phi$ (光滑曲线)}
\label{fig:clt}
\end{figure}

% =====================================================================
\section{随机向量}
\label{sec:1-5}

研究随机过程 $\{X_t\}$ 时, 我们总要同时考察它在若干时刻的取值 $(X_{t_1},\dots,X_{t_n})$: 它们之间是否独立? 相关程度如何? 联合起来服从什么分布? 例如, 布朗运动的定义要求 $(B_{t_1},\dots,B_{t_n})$ 服从多元正态分布. 本节把 1.2、1.3 节的概念推广到多个随机变量, 重点是独立性的刻画与多元正态分布.

% ---------------------------------------------------------------------
\subsection{联合分布}

$\R^n$ 上由全体开集生成的 $\sigma$-代数称为 Borel $\sigma$-代数, 记为 $\mathcal{B}(\R^n)$. 与命题~\ref{prop:borel-generator} 类似, 它可以由 ``左下象限'' 生成.

\begin{proposition}\label{prop:borel-rn}
记 $(-\infty,\rvec{x}]=(-\infty,x_1]\times\cdots\times(-\infty,x_n]$, 其中 $\rvec{x}=(x_1,\dots,x_n)^\T$. 则
$\mathcal{B}(\R^n)=\sigma\big(\{(-\infty,\rvec{x}]: \rvec{x}\in\R^n\}\big)$, 并且全体象限构成 $\pi$-系.
\end{proposition}

\begin{proof}
记右端为 $\mathcal{S}$. 每个象限是闭集, 故 $\mathcal{S}\subseteq\mathcal{B}(\R^n)$. 反过来, 令 $\pi_i(\rvec{x})=x_i$ 为第 $i$ 个坐标. 对 $c\in\R$,
\[
\{\rvec{x}: x_i\le c\}=\bigcup_{m=1}^\infty(-\infty,m]\times\cdots\times\underbrace{(-\infty,c]}_{\text{第 }i\text{ 个}}\times\cdots\times(-\infty,m]\in\mathcal{S},
\]
即 $\pi_i^{-1}((-\infty,c])\in\mathcal{S}$. 由引理~\ref{lem:measurable-generator} (2) 和命题~\ref{prop:borel-generator}, 对每个 Borel 集 $B\subseteq\R$, $\pi_i^{-1}(B)\in\mathcal{S}$. 于是开长方体 $(a_1,b_1)\times\cdots\times(a_n,b_n)=\bigcap_i\pi_i^{-1}((a_i,b_i))\in\mathcal{S}$. 与命题~\ref{prop:borel-generator} 的证明相同, $\R^n$ 中每个开集是可列个有理端点开长方体之并, 所以属于 $\mathcal{S}$, 进而 $\mathcal{B}(\R^n)\subseteq\mathcal{S}$. 最后, 两个象限之交是以分量取最小值为顶点的象限, 故象限全体是 $\pi$-系.
\end{proof}

\begin{definition}[随机向量]\label{def:random-vector}
$\mathcal{F}/\mathcal{B}(\R^n)$-可测映射 $\rvec{X}=(X_1,\dots,X_n)^\T:\Omega\to\R^n$ 称为 $n$ 维\textbf{随机向量}. 称
\[
\mu_{\rvec{X}}(B)=\P(\rvec{X}\in B)\quad(B\in\mathcal{B}(\R^n)),\qquad F_{\rvec{X}}(\rvec{x})=\P(X_1\le x_1,\dots,X_n\le x_n)
\]
分别为 $\rvec{X}$ 的\textbf{联合分布}与\textbf{联合分布函数}. 各分量 $X_i$ 的分布称为\textbf{边缘分布}.
\end{definition}

\begin{proposition}\label{prop:random-vector}
\begin{enumerate}
\item $\rvec{X}$ 是随机向量, 当且仅当每个分量 $X_i$ 都是随机变量.
\item 联合分布函数决定联合分布.
\item 边缘分布由联合分布决定: $\mu_{X_i}(B)=\mu_{\rvec{X}}(\pi_i^{-1}(B))$, $F_{X_i}(x_i)=\lim F_{\rvec{X}}(\rvec{x})$, 其中极限是令 $x_j\to+\infty$ ($j\ne i$).
\item 若 $g:\R^n\to\R^k$ 是 Borel 可测的 (特别地, 连续的), 则 $g(\rvec{X})$ 是随机向量.
\end{enumerate}
\end{proposition}

\begin{proof}
(1) 若 $\rvec{X}$ 可测, 由于 $\pi_i$ 连续 (从而 Borel 可测), $X_i=\pi_i\circ\rvec{X}$ 可测. 反之, 若每个 $X_i$ 可测, 则 $\{\rvec{X}\in(-\infty,\rvec{x}]\}=\bigcap_i\{X_i\le x_i\}\in\mathcal{F}$, 由命题~\ref{prop:borel-rn} 和引理~\ref{lem:measurable-generator} (2), $\rvec{X}$ 可测.
(2) 由命题~\ref{prop:borel-rn} 与推论~\ref{cor:uniqueness}.
(3) 第一式即 $\{X_i\in B\}=\{\rvec{X}\in\pi_i^{-1}(B)\}$. 第二式: 令 $x_j=m\to\infty$ ($j\ne i$), 事件 $\bigcap_j\{X_j\le x_j\}$ 单调上升到 $\{X_i\le x_i\}$, 由概率的连续性即得.
(4) 与命题~\ref{prop:rv-operations} (1) 的证明相同.
\end{proof}

若存在非负 Borel 函数 $f$, 使对一切 $B\in\mathcal{B}(\R^n)$ 都有 $\P(\rvec{X}\in B)=\int_B f(\rvec{x})\dd\rvec{x}$ (积分关于 $\R^n$ 上的 Lebesgue 测度), 称 $f$ 为 $\rvec{X}$ 的\textbf{联合密度}. 计算这类积分时, 我们需要把 $n$ 重积分化为累次积分. 这由下面的测度论定理保证, 它也是第~5~章交换 $\E$ 与 $\int\dd t$ 的依据.

\begin{theorem}[乘积测度与 Fubini-Tonelli 定理]\label{thm:fubini}
设 $(E_1,\mathcal{E}_1,\mu_1)$ 与 $(E_2,\mathcal{E}_2,\mu_2)$ 是 $\sigma$-有限测度空间 (如概率空间、带 Lebesgue 测度的 $\R^n$). 则在 $\mathcal{E}_1\otimes\mathcal{E}_2=\sigma(\{B_1\times B_2: B_i\in\mathcal{E}_i\})$ 上存在唯一的测度 $\mu_1\otimes\mu_2$, 使得 $(\mu_1\otimes\mu_2)(B_1\times B_2)=\mu_1(B_1)\mu_2(B_2)$. 设 $h:E_1\times E_2\to\R$ 是 $\mathcal{E}_1\otimes\mathcal{E}_2$-可测的.
\begin{enumerate}
\item (Tonelli) 若 $h\ge0$, 则
\begin{align*}
\int_{E_1\times E_2}h\dd(\mu_1\otimes\mu_2)&=\int_{E_1}\Big(\int_{E_2}h(x,y)\,\mu_2(\dd y)\Big)\mu_1(\dd x)\\
&=\int_{E_2}\Big(\int_{E_1}h(x,y)\,\mu_1(\dd x)\Big)\mu_2(\dd y),
\end{align*}
其中内层积分作为外层变量的函数是可测的.
\item (Fubini) 若 $\int\abs{h}\dd(\mu_1\otimes\mu_2)<\infty$, 则上式对 $h$ 同样成立 (内层积分对几乎每个外层变量有定义).
\end{enumerate}
(证明见~\cite{folland} 2.5 节.) $\R^n$ 上的 Lebesgue 测度就是 $n$ 个一维 Lebesgue 测度的乘积.
\end{theorem}

\begin{example}[期望的尾积分公式]\label{ex:tail-formula}
设 $X\ge0$. 则 $\E[X]=\int_0^\infty\P(X>t)\dd t$.

\emph{证明.} 对每个 $\omega$, $X(\omega)=\int_0^\infty\1_{\{t<X(\omega)\}}\dd t$. 函数 $h(\omega,t)=\1_{\{t<X(\omega)\}}$ 在 $\Omega\times[0,\infty)$ 上关于 $\mathcal{F}\otimes\mathcal{B}([0,\infty))$ 可测, 因为
\[
\{(\omega,t): t<X(\omega)\}=\bigcup_{q\in\Q_{\ge0}}\{X>q\}\times[0,q)
\]
(若 $t<X(\omega)$, 取有理数 $q\in(t,X(\omega))$ 即可). 对 $\P\otimes\lambda$ 用 Tonelli 定理交换积分次序,
\[
\E[X]=\int_\Omega\int_0^\infty h(\omega,t)\dd t\,\P(\dd\omega)=\int_0^\infty\int_\Omega h(\omega,t)\,\P(\dd\omega)\dd t=\int_0^\infty\P(X>t)\dd t.
\]
例如, 对 $X\sim\mathrm{Exp}(\lambda)$, $\E[X]=\int_0^\infty e^{-\lambda t}\dd t=\frac{1}{\lambda}$, 与例~\ref{ex:means} 一致.
\end{example}

% ---------------------------------------------------------------------
\subsection{独立性的刻画}

独立性可以完全用联合分布刻画: 独立当且仅当联合分布是边缘分布的乘积. 以下把随机变量的独立性 (定义~\ref{def:indep-rv}) 推广到随机向量: 随机向量 $\rvec{X}_1,\dots,\rvec{X}_n$ (分别取值于 $\R^{d_1},\dots,\R^{d_n}$) 称为独立的, 如果 $\sigma(\rvec{X}_1),\dots,\sigma(\rvec{X}_n)$ 独立, 其中 $\sigma(\rvec{X}_i)=\{\rvec{X}_i^{-1}(B): B\in\mathcal{B}(\R^{d_i})\}$.

\begin{theorem}[独立性与乘积测度]\label{thm:indep-product-measure}
随机向量 $\rvec{X}_1,\dots,\rvec{X}_n$ 独立, 当且仅当 $\rvec{X}=(\rvec{X}_1,\dots,\rvec{X}_n)$ 的联合分布等于乘积测度:
\[
\mu_{\rvec{X}}=\mu_{\rvec{X}_1}\otimes\cdots\otimes\mu_{\rvec{X}_n}.
\]
若各 $\rvec{X}_i$ 有密度 $f_i$, 则这又等价于 $\rvec{X}$ 有联合密度 $f(\rvec{x}_1,\dots,\rvec{x}_n)=f_1(\rvec{x}_1)\cdots f_n(\rvec{x}_n)$.
\end{theorem}

\begin{proof}
对 Borel 集 $B_i\subseteq\R^{d_i}$, 有 $\{\rvec{X}\in B_1\times\cdots\times B_n\}=\bigcap_i\{\rvec{X}_i\in B_i\}$. 因此独立性等价于对一切这样的 $B_i$ 都有
\[
\mu_{\rvec{X}}(B_1\times\cdots\times B_n)=\mu_{\rvec{X}_1}(B_1)\cdots\mu_{\rvec{X}_n}(B_n).
\]
这恰好说明 $\mu_{\rvec{X}}$ 与乘积测度在 ``可测矩形'' 上一致. 可测矩形构成 $\pi$-系 (两个矩形之交仍是矩形), 并生成乘积 $\sigma$-代数 (它包含全体象限, 由命题~\ref{prop:borel-rn}, 生成的就是 Borel $\sigma$-代数). 由推论~\ref{cor:uniqueness}, ``在矩形上一致'' 等价于 ``测度相等''.

对密度的论断: 由 Tonelli 定理, 以 $\prod_i f_i(\rvec{x}_i)$ 为密度的测度在矩形上的值为 $\prod_i\int_{B_i}f_i=\prod_i\mu_{\rvec{X}_i}(B_i)$, 所以它就是乘积测度. 因此 $\rvec{X}$ 独立当且仅当 $\mu_{\rvec{X}}$ 等于这个测度, 即 $\prod_i f_i$ 是 $\rvec{X}$ 的联合密度.
\end{proof}

独立性最重要的计算后果是期望对乘积可以分解.

\begin{theorem}[独立随机变量乘积的期望]\label{thm:indep-product}
设 $X,Y$ 独立且可积. 则 $XY$ 可积, 且 $\E[XY]=\E[X]\E[Y]$.
\end{theorem}

\begin{proof}
按期望构造的四步证明.

\emph{示性函数.} 若 $A\in\sigma(X)$, $B\in\sigma(Y)$, 则由独立性 $\E[\1_A\1_B]=\P(A\cap B)=\P(A)\P(B)=\E[\1_A]\E[\1_B]$.

\emph{简单随机变量.} 若 $U=\sum_i a_i\1_{A_i}$, $V=\sum_j b_j\1_{B_j}$, $A_i\in\sigma(X)$, $B_j\in\sigma(Y)$, 由线性与上一步,
$\E[UV]=\sum_{i,j}a_ib_j\E[\1_{A_i}\1_{B_j}]=\sum_{i,j}a_ib_j\P(A_i)\P(B_j)=\E[U]\E[V]$.

\emph{非负随机变量.} 设 $X,Y\ge0$. 由引理~\ref{lem:simple-approx}, 取简单随机变量 $U_n\uparrow X$, $V_n\uparrow Y$; 由于 $U_n$ 中出现的集合形如 $\{a\le X<b\}$, $U_n$ 是 $\sigma(X)$-可测的, 同理 $V_n$ 是 $\sigma(Y)$-可测的. 于是 $U_nV_n\uparrow XY$, 由单调收敛定理~\ref{thm:mct} 与上一步,
\[
\E[XY]=\lim_n\E[U_nV_n]=\lim_n\E[U_n]\E[V_n]=\E[X]\E[Y].
\]

\emph{一般情形.} 对 $\abs{X},\abs{Y}$ (它们分别是 $\sigma(X)$、$\sigma(Y)$-可测的) 用上一步, $\E\abs{XY}=\E\abs{X}\E\abs{Y}<\infty$, 故 $XY$ 可积. 再把 $XY=(X^+-X^-)(Y^+-Y^-)$ 展开成四项, 每一项都是 $\sigma(X)$-可测与 $\sigma(Y)$-可测的非负随机变量之积, 由上一步与线性,
\[
\E[XY]=\E[X^+]\E[Y^+]-\E[X^+]\E[Y^-]-\E[X^-]\E[Y^+]+\E[X^-]\E[Y^-]=\E[X]\E[Y].\qedhere
\]
\end{proof}

\begin{corollary}\label{cor:indep-consequences}
\begin{enumerate}
\item 若 $X,Y$ 独立, $g,h$ 是 Borel 函数且 $g(X),h(Y)$ 可积, 则 $\E[g(X)h(Y)]=\E[g(X)]\E[h(Y)]$.
\item 若 $X,Y\in L^2$ 独立, 则 $\Cov(X,Y)=0$. 若 $X_1,\dots,X_n\in L^2$ 两两不相关 (例如相互独立), 则 $\Var\big(\sum_iX_i\big)=\sum_i\Var(X_i)$.
\item 若 $X,Y$ 独立, 则 $\varphi_{X+Y}(t)=\varphi_X(t)\varphi_Y(t)$.
\end{enumerate}
\end{corollary}

\begin{proof}
(1) 由命题~\ref{prop:indep-functions}, $g(X)$ 与 $h(Y)$ 独立, 再用定理~\ref{thm:indep-product}.
(2) $\Cov(X,Y)=\E[XY]-\E X\E Y=0$. 对和式展开 $\Var(\sum_iX_i)=\sum_i\Var(X_i)+\sum_{i\ne j}\Cov(X_i,X_j)$ (命题~\ref{prop:variance} (4) 的推广).
(3) $e^{\mathrm{i}t(X+Y)}=e^{\mathrm{i}tX}e^{\mathrm{i}tY}$. 把两个因子写成实部加虚部并展开, 得到四项, 每项都是 $X$ 的有界函数与 $Y$ 的有界函数之积, 由 (1) 分别分解后再合并即得.
\end{proof}

\begin{example}\label{ex:indep-sums}
\begin{enumerate}
\item \textbf{二项分布的方差.} $S_n=\sum_{k=1}^n\1_{A_k}$, 其中 $A_k$ 独立且 $\P(A_k)=p$. 由推论~\ref{cor:indep-consequences} (2), $\Var(S_n)=\sum_k\Var(\1_{A_k})=np(1-p)$.
\item \textbf{独立正态之和.} 若 $X\sim N(\mu_1,\sigma_1^2)$ 与 $Y\sim N(\mu_2,\sigma_2^2)$ 独立, 则由 \eqref{eq:normal-cf},
\[
\varphi_{X+Y}(t)=e^{\mathrm{i}\mu_1t-\frac{1}{2}\sigma_1^2t^2}e^{\mathrm{i}\mu_2t-\frac{1}{2}\sigma_2^2t^2}=\exp\Big(\mathrm{i}(\mu_1+\mu_2)t-\frac{1}{2}(\sigma_1^2+\sigma_2^2)t^2\Big),
\]
由唯一性定理~\ref{thm:cf-uniqueness}, $X+Y\sim N(\mu_1+\mu_2,\sigma_1^2+\sigma_2^2)$. 这是布朗运动增量 ``可加'' 的基础: $B_t-B_0=(B_t-B_s)+(B_s-B_0)$.
\item \textbf{独立 Poisson 之和.} 同理, 若 $X\sim\mathrm{Poi}(\lambda_1)$ 与 $Y\sim\mathrm{Poi}(\lambda_2)$ 独立, 则 $X+Y\sim\mathrm{Poi}(\lambda_1+\lambda_2)$.
\end{enumerate}
\end{example}

\begin{proposition}[卷积公式]\label{prop:convolution}
若 $X,Y$ 独立, 分别有密度 $f,g$, 则 $X+Y$ 有密度 $(f*g)(z)=\int_{\R}f(z-y)g(y)\dd y$.
\end{proposition}

\begin{proof}
由定理~\ref{thm:indep-product-measure}, $(X,Y)$ 有联合密度 $f(x)g(y)$. 对 $z\in\R$, 由 Tonelli 定理与内层积分的平移 $x=u-y$,
\begin{align*}
\P(X+Y\le z)&=\int_{\R}\Big(\int_{-\infty}^{z-y}f(x)\dd x\Big)g(y)\dd y=\int_{\R}\Big(\int_{-\infty}^{z}f(u-y)\dd u\Big)g(y)\dd y\\
&=\int_{-\infty}^z\Big(\int_{\R}f(u-y)g(y)\dd y\Big)\dd u.
\end{align*}
由定义~\ref{def:discrete-continuous} (2), $f*g$ 是 $X+Y$ 的密度.
\end{proof}

对随机向量 $\rvec{X}\in\R^n$, 其特征函数定义为 $\varphi_{\rvec{X}}(\rvec{t})=\E\big[e^{\mathrm{i}\rvec{t}^\T\rvec{X}}\big]$, $\rvec{t}\in\R^n$. 定理~\ref{thm:cf-uniqueness} 对随机向量同样成立: 特征函数相同的随机向量同分布 (见~\cite{durrett} 3.10 节). 由此得到一个非常方便的独立性判别法.

\begin{theorem}[独立性的特征函数判别法]\label{thm:indep-cf}
设 $\rvec{X}\in\R^k$, $\rvec{Y}\in\R^m$ 是随机向量. 则 $\rvec{X}$ 与 $\rvec{Y}$ 独立, 当且仅当对一切 $\rvec{s}\in\R^k$, $\rvec{t}\in\R^m$,
\[
\varphi_{(\rvec{X},\rvec{Y})}(\rvec{s},\rvec{t})=\varphi_{\rvec{X}}(\rvec{s})\,\varphi_{\rvec{Y}}(\rvec{t}).
\]
\end{theorem}

\begin{proof}
``仅当'': $e^{\mathrm{i}(\rvec{s}^\T\rvec{X}+\rvec{t}^\T\rvec{Y})}=e^{\mathrm{i}\rvec{s}^\T\rvec{X}}e^{\mathrm{i}\rvec{t}^\T\rvec{Y}}$, 与推论~\ref{cor:indep-consequences} (3) 的证明相同.

``当'': 在乘积概率空间上构造独立的随机向量 $\rvec{X}'\overset{d}{=}\rvec{X}$ 与 $\rvec{Y}'\overset{d}{=}\rvec{Y}$ (取 $(\R^k\times\R^m,\mathcal{B}(\R^{k+m}),\mu_{\rvec{X}}\otimes\mu_{\rvec{Y}})$ 上的坐标映射即可). 由 ``仅当'' 部分与假设,
\[
\varphi_{(\rvec{X}',\rvec{Y}')}(\rvec{s},\rvec{t})=\varphi_{\rvec{X}}(\rvec{s})\varphi_{\rvec{Y}}(\rvec{t})=\varphi_{(\rvec{X},\rvec{Y})}(\rvec{s},\rvec{t}).
\]
由唯一性定理, $\mu_{(\rvec{X},\rvec{Y})}=\mu_{(\rvec{X}',\rvec{Y}')}=\mu_{\rvec{X}}\otimes\mu_{\rvec{Y}}$, 由定理~\ref{thm:indep-product-measure}, $\rvec{X}$ 与 $\rvec{Y}$ 独立.
\end{proof}

% ---------------------------------------------------------------------
\subsection{密度变换公式}

\begin{theorem}[密度变换公式]\label{thm:jacobian}
设随机向量 $\rvec{X}$ 有密度 $f_{\rvec{X}}$, 且 $\P(\rvec{X}\in U)=1$, 其中 $U\subseteq\R^n$ 是开集. 设 $g:U\to V$ 是开集之间的双射, $g$ 与其逆 $h=g^{-1}$ 都连续可微. 则 $\rvec{Y}=g(\rvec{X})$ 有密度
\[
f_{\rvec{Y}}(\rvec{y})=f_{\rvec{X}}(h(\rvec{y}))\,\abs{\det J_h(\rvec{y})}\,\1_V(\rvec{y}),
\]
其中 $J_h$ 是 $h$ 的 Jacobi 矩阵.
\end{theorem}

\begin{proof}
对 Borel 集 $B\subseteq V$, $\{\rvec{Y}\in B\}=\{\rvec{X}\in h(B)\}$. 由多元积分的变量替换公式 ($\rvec{x}=h(\rvec{y})$, 见~\cite{folland} 定理 2.47),
\[
\P(\rvec{Y}\in B)=\int_{h(B)}f_{\rvec{X}}(\rvec{x})\dd\rvec{x}=\int_B f_{\rvec{X}}(h(\rvec{y}))\abs{\det J_h(\rvec{y})}\dd\rvec{y}.
\]
又 $\P(\rvec{Y}\notin V)=0$, 故对一般的 Borel 集 $B$, $\P(\rvec{Y}\in B)=\P(\rvec{Y}\in B\cap V)$, 即得所述密度.
\end{proof}

\begin{example}[正态向量的极坐标与 Box-Muller 方法]\label{ex:box-muller}
设 $Z_1,Z_2$ 是独立的 $N(0,1)$ 随机变量. 由定理~\ref{thm:indep-product-measure}, $(Z_1,Z_2)$ 有联合密度 $\frac{1}{2\pi}e^{-(z_1^2+z_2^2)/2}$. 去掉零概率集 $\{(z_1,0): z_1\ge0\}$ 后, 极坐标 $h(r,\theta)=(r\cos\theta,r\sin\theta)$ 是从 $(0,\infty)\times(0,2\pi)$ 到其余部分的连续可微双射, $\abs{\det J_h}=r$. 由定理~\ref{thm:jacobian}, $(R,\Theta)=h^{-1}(Z_1,Z_2)$ 的密度为
\[
f_{R,\Theta}(r,\theta)=\frac{1}{2\pi}e^{-r^2/2}\cdot r=\Big(re^{-r^2/2}\Big)\cdot\Big(\frac{1}{2\pi}\Big),\qquad r>0,\ 0<\theta<2\pi.
\]
密度分解为 $r$ 的函数与 $\theta$ 的函数之积, 由定理~\ref{thm:indep-product-measure}, $R$ 与 $\Theta$ 独立, $\Theta\sim\mathrm{U}(0,2\pi)$, 且 $\P(R>r)=e^{-r^2/2}$, 即 $R^2\sim\mathrm{Exp}(\frac{1}{2})$. 反过来, 若 $U_1,U_2$ 是独立的 $\mathrm{U}(0,1)$ 随机变量, 令 $R=\sqrt{-2\log U_1}$, $\Theta=2\pi U_2$, 则 $(R\cos\Theta,R\sin\Theta)$ 是一对独立的标准正态随机变量. 这就是计算机中生成正态随机数的 Box-Muller 方法.
\end{example}

% ---------------------------------------------------------------------
\subsection{期望向量与协方差矩阵}

\begin{definition}[期望向量与协方差矩阵]\label{def:covariance-matrix}
设 $\rvec{X}=(X_1,\dots,X_n)^\T$ 的各分量属于 $L^2$. 称 $\E\rvec{X}=(\E X_1,\dots,\E X_n)^\T$ 为\textbf{期望向量}, 称 $n\times n$ 矩阵
\[
\mat{\Sigma}=\Cov(\rvec{X})=\E\big[(\rvec{X}-\E\rvec{X})(\rvec{X}-\E\rvec{X})^\T\big]=\big(\Cov(X_i,X_j)\big)_{i,j=1}^n
\]
为\textbf{协方差矩阵} (随机矩阵的期望按元素逐个定义).
\end{definition}

\begin{proposition}[协方差矩阵的性质]\label{prop:covariance-matrix}
设 $\rvec{X}$ 的期望向量为 $\rvec{\mu}$, 协方差矩阵为 $\mat{\Sigma}$.
\begin{enumerate}
\item $\mat{\Sigma}$ 对称半正定: 对任意 $\rvec{a}\in\R^n$, $\rvec{a}^\T\mat{\Sigma}\rvec{a}=\Var(\rvec{a}^\T\rvec{X})\ge0$.
\item 线性变换: 对 $k\times n$ 矩阵 $\mat{A}$ 与 $\rvec{b}\in\R^k$, $\E[\mat{A}\rvec{X}+\rvec{b}]=\mat{A}\rvec{\mu}+\rvec{b}$, $\Cov(\mat{A}\rvec{X}+\rvec{b})=\mat{A}\mat{\Sigma}\mat{A}^\T$.
\item $\mat{\Sigma}$ 奇异, 当且仅当存在 $\rvec{a}\ne\rvec{0}$, 使得 $\rvec{a}^\T\rvec{X}=\rvec{a}^\T\rvec{\mu}$ $\as$, 即 $\rvec{X}$ 几乎必然落在某个超平面上.
\end{enumerate}
\end{proposition}

\begin{proof}
(1) 对称性来自 $\Cov(X_i,X_j)=\Cov(X_j,X_i)$. 由期望的线性,
\[
\rvec{a}^\T\mat{\Sigma}\rvec{a}=\E\big[\rvec{a}^\T(\rvec{X}-\rvec{\mu})(\rvec{X}-\rvec{\mu})^\T\rvec{a}\big]=\E\big[(\rvec{a}^\T\rvec{X}-\rvec{a}^\T\rvec{\mu})^2\big]=\Var(\rvec{a}^\T\rvec{X})\ge0.
\]
(2) 由线性, $\E[\mat{A}\rvec{X}+\rvec{b}]=\mat{A}\rvec{\mu}+\rvec{b}$, 且
$\Cov(\mat{A}\rvec{X}+\rvec{b})=\E\big[\mat{A}(\rvec{X}-\rvec{\mu})(\rvec{X}-\rvec{\mu})^\T\mat{A}^\T\big]=\mat{A}\mat{\Sigma}\mat{A}^\T$.

(3) 若 $\mat{\Sigma}\rvec{a}=\rvec{0}$, $\rvec{a}\ne\rvec{0}$, 则由 (1), $\Var(\rvec{a}^\T\rvec{X})=\rvec{a}^\T\mat{\Sigma}\rvec{a}=0$, 由命题~\ref{prop:variance} (3), $\rvec{a}^\T\rvec{X}=\rvec{a}^\T\rvec{\mu}$ $\as$ 反之, 若这样的 $\rvec{a}$ 存在, 则 $\rvec{a}^\T\mat{\Sigma}\rvec{a}=0$. 由于 $\mat{\Sigma}$ 对称半正定, 可写 $\mat{\Sigma}=\mat{\Sigma}^{1/2}\mat{\Sigma}^{1/2}$ ($\mat{\Sigma}^{1/2}$ 对称), 于是 $\norm{\mat{\Sigma}^{1/2}\rvec{a}}^2=0$, 从而 $\mat{\Sigma}\rvec{a}=\mat{\Sigma}^{1/2}(\mat{\Sigma}^{1/2}\rvec{a})=\rvec{0}$, $\mat{\Sigma}$ 奇异.
\end{proof}

由推论~\ref{cor:indep-consequences} (2), 独立蕴含不相关. 反之不然.

\begin{example}[不相关但不独立]\label{ex:uncorrelated-dependent}
设 $X\sim N(0,1)$, $Y=X^2$. 由例~\ref{ex:normal-mgf}, $\Cov(X,Y)=\E[X^3]-\E[X]\E[X^2]=0$. 但 $X$ 与 $Y$ 显然不独立: 令 $p=\P(\abs{X}\le1)\in(0,1)$, 则
\[
\P(\abs{X}\le1,\,Y\le1)=\P(\abs{X}\le1)=p\ne p^2=\P(\abs{X}\le1)\P(Y\le1).
\]
协方差只度量\emph{线性}相关性; $Y$ 由 $X$ 完全决定, 但这种依赖是非线性的, 协方差察觉不到.
\end{example}

% ---------------------------------------------------------------------
\subsection{多元正态分布}

一维正态分布由均值和方差决定. 多元正态分布是随机分析中最重要的分布, 布朗运动、Wiener 积分以及 Ornstein-Uhlenbeck 过程的有限维分布都是多元正态的. 它有几种形式不同但彼此等价的定义.

\begin{definition}[正态随机向量]\label{def:gaussian-vector}
随机向量 $\rvec{X}\in\R^n$ 称为\textbf{正态随机向量} (或称 $X_1,\dots,X_n$ \textbf{联合正态}), 如果对每个 $\rvec{a}\in\R^n$, 线性组合 $\rvec{a}^\T\rvec{X}=a_1X_1+\cdots+a_nX_n$ 都服从 (可能退化的) 一维正态分布.
\end{definition}

取 $\rvec{a}$ 为第 $i$ 个坐标向量可知, 正态随机向量的每个分量都是正态的, 因而属于 $L^2$, 期望向量与协方差矩阵都有定义.

\begin{theorem}[多元正态分布的等价刻画]\label{thm:gaussian-equivalent}
设 $\rvec{\mu}\in\R^n$, $\mat{\Sigma}$ 是 $n\times n$ 对称半正定矩阵. 对随机向量 $\rvec{X}\in\R^n$, 以下三条等价:
\begin{enumerate}[(A)]
\item $\rvec{X}$ 是正态随机向量, 且 $\E\rvec{X}=\rvec{\mu}$, $\Cov(\rvec{X})=\mat{\Sigma}$;
\item $\rvec{X}$ 的特征函数为 $\varphi_{\rvec{X}}(\rvec{t})=\exp\big(\mathrm{i}\,\rvec{t}^\T\rvec{\mu}-\frac{1}{2}\rvec{t}^\T\mat{\Sigma}\rvec{t}\big)$;
\item $\rvec{X}\overset{d}{=}\rvec{\mu}+\mat{A}\rvec{Z}$, 其中 $\mat{A}$ 是满足 $\mat{A}\mat{A}^\T=\mat{\Sigma}$ 的 $n\times m$ 矩阵, $\rvec{Z}=(Z_1,\dots,Z_m)^\T$ 的分量是独立的 $N(0,1)$ 随机变量.
\end{enumerate}
若进一步 $\mat{\Sigma}$ 正定, 则它们还等价于
\begin{enumerate}[(A),start=4]
\item $\rvec{X}$ 有密度
$\displaystyle f(\rvec{x})=\frac{1}{(2\pi)^{n/2}(\det\mat{\Sigma})^{1/2}}\exp\Big(-\frac{1}{2}(\rvec{x}-\rvec{\mu})^\T\mat{\Sigma}^{-1}(\rvec{x}-\rvec{\mu})\Big)$.
\end{enumerate}
满足上述条件的分布记为 $N(\rvec{\mu},\mat{\Sigma})$. 特别地, (C) 说明对任意 $\rvec{\mu}$ 与对称半正定的 $\mat{\Sigma}$, $N(\rvec{\mu},\mat{\Sigma})$ 存在.
\end{theorem}

\begin{proof}
(A)$\Rightarrow$(B). 对 $\rvec{t}\in\R^n$, 由 (A), $\rvec{t}^\T\rvec{X}$ 是正态随机变量, 由命题~\ref{prop:covariance-matrix}, 其均值为 $\rvec{t}^\T\rvec{\mu}$, 方差为 $\rvec{t}^\T\mat{\Sigma}\rvec{t}$. 由 \eqref{eq:normal-cf} (取 $t=1$),
\[
\varphi_{\rvec{X}}(\rvec{t})=\E\big[e^{\mathrm{i}\cdot1\cdot(\rvec{t}^\T\rvec{X})}\big]=\varphi_{\rvec{t}^\T\rvec{X}}(1)=\exp\Big(\mathrm{i}\,\rvec{t}^\T\rvec{\mu}-\frac{1}{2}\rvec{t}^\T\mat{\Sigma}\rvec{t}\Big).
\]

(B)$\Rightarrow$(A). 对 $\rvec{a}\in\R^n$ 与 $s\in\R$, 由 (B),
\[
\varphi_{\rvec{a}^\T\rvec{X}}(s)=\E\big[e^{\mathrm{i}(s\rvec{a})^\T\rvec{X}}\big]=\varphi_{\rvec{X}}(s\rvec{a})=\exp\Big(\mathrm{i}s\,\rvec{a}^\T\rvec{\mu}-\frac{1}{2}s^2\,\rvec{a}^\T\mat{\Sigma}\rvec{a}\Big).
\]
这是 $N(\rvec{a}^\T\rvec{\mu},\rvec{a}^\T\mat{\Sigma}\rvec{a})$ 的特征函数, 由唯一性定理~\ref{thm:cf-uniqueness}, $\rvec{a}^\T\rvec{X}\sim N(\rvec{a}^\T\rvec{\mu},\rvec{a}^\T\mat{\Sigma}\rvec{a})$, 故 $\rvec{X}$ 是正态随机向量. 取 $\rvec{a}=\rvec{e}_i$ 得 $\E X_i=\mu_i$, $\Var(X_i)=\Sigma_{ii}$; 取 $\rvec{a}=\rvec{e}_i+\rvec{e}_j$ 得 $\Var(X_i+X_j)=\Sigma_{ii}+\Sigma_{jj}+2\Sigma_{ij}$, 与命题~\ref{prop:variance} (4) 比较得 $\Cov(X_i,X_j)=\Sigma_{ij}$.

(C)$\Rightarrow$(B). 记 $\rvec{Y}=\rvec{\mu}+\mat{A}\rvec{Z}$. 由于 $Z_1,\dots,Z_m$ 独立, 由推论~\ref{cor:indep-consequences} (1) (对 $m$ 个因子归纳), 对 $\rvec{u}\in\R^m$,
$\varphi_{\rvec{Z}}(\rvec{u})=\prod_{k=1}^m\E[e^{\mathrm{i}u_kZ_k}]=\prod_k e^{-u_k^2/2}=e^{-\frac{1}{2}\rvec{u}^\T\rvec{u}}$. 于是
\[
\varphi_{\rvec{Y}}(\rvec{t})=e^{\mathrm{i}\rvec{t}^\T\rvec{\mu}}\,\E\big[e^{\mathrm{i}(\mat{A}^\T\rvec{t})^\T\rvec{Z}}\big]=e^{\mathrm{i}\rvec{t}^\T\rvec{\mu}}\varphi_{\rvec{Z}}(\mat{A}^\T\rvec{t})=\exp\Big(\mathrm{i}\,\rvec{t}^\T\rvec{\mu}-\frac{1}{2}\rvec{t}^\T\mat{A}\mat{A}^\T\rvec{t}\Big).
\]
由 $\mat{A}\mat{A}^\T=\mat{\Sigma}$ 及 $\rvec{X}\overset{d}{=}\rvec{Y}$ (特征函数只依赖于分布) 得 (B).

(B)$\Rightarrow$(C). 由谱分解, $\mat{\Sigma}=\mat{Q}\mat{\Lambda}\mat{Q}^\T$, 其中 $\mat{Q}$ 正交, $\mat{\Lambda}=\mathrm{diag}(\lambda_1,\dots,\lambda_n)$, $\lambda_i\ge0$. 令 $\mat{A}=\mat{Q}\mat{\Lambda}^{1/2}$, 则 $\mat{A}\mat{A}^\T=\mat{\Sigma}$. 由上一步, $\rvec{\mu}+\mat{A}\rvec{Z}$ 的特征函数与 $\rvec{X}$ 的相同, 由唯一性定理, $\rvec{X}\overset{d}{=}\rvec{\mu}+\mat{A}\rvec{Z}$.

(C)$\Rightarrow$(D), 设 $\mat{\Sigma}$ 正定. 取 $\mat{A}=\mat{\Sigma}^{1/2}=\mat{Q}\mat{\Lambda}^{1/2}\mat{Q}^\T$, 它对称可逆, 且 (由 (B)$\Rightarrow$(C)) $\rvec{X}\overset{d}{=}\rvec{Y}=\rvec{\mu}+\mat{A}\rvec{Z}$, $m=n$. 由定理~\ref{thm:indep-product-measure}, $\rvec{Z}$ 有密度 $\prod_k\varphi(z_k)=(2\pi)^{-n/2}e^{-\rvec{z}^\T\rvec{z}/2}$. 映射 $\rvec{z}\mapsto\rvec{\mu}+\mat{A}\rvec{z}$ 是 $\R^n$ 上的双射, 逆映射 $h(\rvec{y})=\mat{A}^{-1}(\rvec{y}-\rvec{\mu})$ 的 Jacobi 行列式为 $\det\mat{A}^{-1}=(\det\mat{\Sigma})^{-1/2}$. 由密度变换公式~\ref{thm:jacobian}, 并利用 $(\mat{A}^{-1})^\T\mat{A}^{-1}=(\mat{A}\mat{A}^\T)^{-1}=\mat{\Sigma}^{-1}$,
\[
f_{\rvec{Y}}(\rvec{y})=\frac{1}{(2\pi)^{n/2}}\exp\Big(-\frac{1}{2}(\rvec{y}-\rvec{\mu})^\T(\mat{A}^{-1})^\T\mat{A}^{-1}(\rvec{y}-\rvec{\mu})\Big)(\det\mat{\Sigma})^{-1/2},
\]
正是 (D) 中的 $f$. 由于 $\rvec{X}\overset{d}{=}\rvec{Y}$, $\rvec{X}$ 也以 $f$ 为密度.

(D)$\Rightarrow$(C). 若 $\rvec{X}$ 以 $f$ 为密度, 由上一步, $\rvec{Y}=\rvec{\mu}+\mat{\Sigma}^{1/2}\rvec{Z}$ 也以 $f$ 为密度, 所以 $\rvec{X}\overset{d}{=}\rvec{Y}$.
\end{proof}

\begin{intuition}
三种刻画各有用处. (A) 最便于验证 (例如证明布朗运动是高斯过程); (B) 最便于计算 (证明独立性、求极限); (C) 给出构造方法: 正态随机向量就是独立标准正态变量经过线性变换再平移. 当 $\mat{\Sigma}$ 退化时 (例如 $\rvec{X}=(Z,Z)^\T$, $\mat{\Sigma}=\begin{psmallmatrix}1&1\\1&1\end{psmallmatrix}$), $\rvec{X}$ 落在一个低维子空间上, 没有密度, 但 (A)(B)(C) 依然有效. 这正是我们把退化正态分布也纳入定义的原因.
\end{intuition}

\begin{figure}[htbp]
\centering
\begin{tikzpicture}[scale=1.0]
  \begin{scope}
    \clip (-3.2,-2.5) rectangle (3.2,2.5);
    \foreach \c/\o in {1.45/10, 1/22, 0.55/40} {
      \fill[SABlue!\o,rotate=37.98] (0,0) ellipse ({1.6676*\c} and {0.8481*\c});
    }
    \foreach \c in {0.55,1,1.45} {
      \draw[SABlue!80!black,rotate=37.98] (0,0) ellipse ({1.6676*\c} and {0.8481*\c});
    }
  \end{scope}
  \draw[->,SADark!70] (-3.3,0) -- (3.4,0) node[right,font=\small] {$x_1$};
  \draw[->,SADark!70] (0,-2.6) -- (0,2.7) node[above,font=\small] {$x_2$};
  \draw[->,SARed,very thick] (0,0) -- ({1.6676*cos(37.98)},{1.6676*sin(37.98)});
  \node[SARed,font=\small,anchor=north west] at ({1.6676*cos(37.98)},{1.6676*sin(37.98)}) {$\sqrt{\lambda_1}\rvec{q}_1$};
  \draw[->,SARed,very thick] (0,0) -- ({-0.8481*sin(37.98)},{0.8481*cos(37.98)});
  \node[SARed,font=\small,anchor=south east] at ({-0.8481*sin(37.98)},{0.8481*cos(37.98)}) {$\sqrt{\lambda_2}\rvec{q}_2$};
  \node[font=\small,align=left,anchor=west] at (3.6,1.2) {$\mat{\Sigma}=\begin{pmatrix}2&1\\1&1.5\end{pmatrix}$};
  \node[font=\small,align=left,anchor=west] at (3.6,-0.2) {$\lambda_1\approx2.78$, $\lambda_2\approx0.72$};
  \node[font=\small,align=left,anchor=west] at (3.6,-1.2) {$\rho=\frac{1}{\sqrt{3}}\approx0.58$};
\end{tikzpicture}
\caption{二维正态分布 $N(\rvec{0},\mat{\Sigma})$ 密度的等高线是椭圆 $\rvec{x}^\T\mat{\Sigma}^{-1}\rvec{x}=c$. 椭圆主轴沿 $\mat{\Sigma}$ 的特征向量 $\rvec{q}_1,\rvec{q}_2$ 方向, 半轴长正比于 $\sqrt{\lambda_1},\sqrt{\lambda_2}$. 正相关使椭圆向第一、三象限倾斜}
\label{fig:gaussian-contours}
\end{figure}

\begin{proposition}[正态随机向量的运算]\label{prop:gaussian-operations}
设 $\rvec{X}\sim N(\rvec{\mu},\mat{\Sigma})$.
\begin{enumerate}
\item 对 $k\times n$ 矩阵 $\mat{B}$ 与 $\rvec{c}\in\R^k$, $\mat{B}\rvec{X}+\rvec{c}\sim N(\mat{B}\rvec{\mu}+\rvec{c},\mat{B}\mat{\Sigma}\mat{B}^\T)$. 特别地, $\rvec{X}$ 的任意子向量都是正态随机向量.
\item 若 $\rvec{X}$ 与 $\rvec{Y}$ 是独立的正态随机向量, 则 $(\rvec{X},\rvec{Y})$ 是正态随机向量.
\item 若 $\rvec{Z}\sim N(\rvec{0},\mat{I}_n)$, $\mat{Q}$ 是正交矩阵, 则 $\mat{Q}\rvec{Z}\sim N(\rvec{0},\mat{I}_n)$.
\end{enumerate}
\end{proposition}

\begin{proof}
(1) 对 $\rvec{a}\in\R^k$, $\rvec{a}^\T(\mat{B}\rvec{X}+\rvec{c})=(\mat{B}^\T\rvec{a})^\T\rvec{X}+\rvec{a}^\T\rvec{c}$ 是正态随机变量加常数, 仍是正态的. 均值与协方差由命题~\ref{prop:covariance-matrix} (2) 给出. 取 $\mat{B}$ 为选取部分坐标的矩阵即得子向量的结论.
(2) 由定理~\ref{thm:indep-cf} 与 (B), $\varphi_{(\rvec{X},\rvec{Y})}(\rvec{s},\rvec{t})=\varphi_{\rvec{X}}(\rvec{s})\varphi_{\rvec{Y}}(\rvec{t})$ 是二次多项式的指数, 形如 (B), 其协方差矩阵为分块对角阵 $\mathrm{diag}(\mat{\Sigma}_{\rvec{X}},\mat{\Sigma}_{\rvec{Y}})$.
(3) 由 (1), $\mat{Q}\rvec{Z}\sim N(\rvec{0},\mat{Q}\mat{Q}^\T)=N(\rvec{0},\mat{I}_n)$.
\end{proof}

(3) 说明标准正态向量的分布是旋转不变的, 这与例~\ref{ex:box-muller} 中 $\Theta$ 服从均匀分布是同一件事. 下面是正态随机向量最重要的性质.

\begin{theorem}[正态向量: 不相关等价于独立]\label{thm:gaussian-uncorrelated}
设 $(\rvec{X},\rvec{Y})=(X_1,\dots,X_k,Y_1,\dots,Y_m)^\T$ 是正态随机向量. 则 $\rvec{X}$ 与 $\rvec{Y}$ 独立, 当且仅当 $\Cov(X_i,Y_j)=0$ 对一切 $i,j$ 成立.
\end{theorem}

\begin{proof}
``仅当'': 由推论~\ref{cor:indep-consequences} (2).

``当'': 由假设, 协方差矩阵是分块对角的: $\mat{\Sigma}=\mathrm{diag}(\mat{\Sigma}_{\rvec{X}},\mat{\Sigma}_{\rvec{Y}})$, 其中 $\mat{\Sigma}_{\rvec{X}}=\Cov(\rvec{X})$, $\mat{\Sigma}_{\rvec{Y}}=\Cov(\rvec{Y})$. 记期望向量为 $(\rvec{\mu}_{\rvec{X}},\rvec{\mu}_{\rvec{Y}})$. 由定理~\ref{thm:gaussian-equivalent} (B), 对 $\rvec{s}\in\R^k$, $\rvec{t}\in\R^m$,
\begin{align*}
\varphi_{(\rvec{X},\rvec{Y})}(\rvec{s},\rvec{t})&=\exp\Big(\mathrm{i}(\rvec{s}^\T\rvec{\mu}_{\rvec{X}}+\rvec{t}^\T\rvec{\mu}_{\rvec{Y}})-\frac{1}{2}\big(\rvec{s}^\T\mat{\Sigma}_{\rvec{X}}\rvec{s}+\rvec{t}^\T\mat{\Sigma}_{\rvec{Y}}\rvec{t}\big)\Big)\\
&=\exp\Big(\mathrm{i}\rvec{s}^\T\rvec{\mu}_{\rvec{X}}-\frac{1}{2}\rvec{s}^\T\mat{\Sigma}_{\rvec{X}}\rvec{s}\Big)\exp\Big(\mathrm{i}\rvec{t}^\T\rvec{\mu}_{\rvec{Y}}-\frac{1}{2}\rvec{t}^\T\mat{\Sigma}_{\rvec{Y}}\rvec{t}\Big)\\
&=\varphi_{\rvec{X}}(\rvec{s})\,\varphi_{\rvec{Y}}(\rvec{t}),
\end{align*}
最后一步是因为 $\rvec{X}$、$\rvec{Y}$ 作为子向量也是正态的 (命题~\ref{prop:gaussian-operations} (1)), 其特征函数由 (B) 给出. 由独立性的特征函数判别法 (定理~\ref{thm:indep-cf}), $\rvec{X}$ 与 $\rvec{Y}$ 独立.
\end{proof}

\begin{intuition}
对一般随机变量, 协方差为零只排除了线性依赖 (例~\ref{ex:uncorrelated-dependent}); 但对联合正态的随机变量, 一切依赖关系都是线性的, 所以 ``没有线性依赖'' 就是 ``没有依赖''. 这一性质使得正态随机向量的计算可以完全化归为线性代数.
\end{intuition}

定理中 ``联合正态'' 的假设不能减弱为 ``各分量分别是正态的''.

\begin{example}[两个正态变量不相关却不独立]\label{ex:normal-not-joint}
设 $X\sim N(0,1)$, $S$ 与 $X$ 独立, 且 $\P(S=1)=\P(S=-1)=\frac{1}{2}$. 令 $Y=SX$.
\begin{itemize}
\item \emph{$Y\sim N(0,1)$}: 由全概率与独立性, $\P(Y\le y)=\frac{1}{2}\P(X\le y)+\frac{1}{2}\P(-X\le y)=\Phi(y)$.
\item \emph{不相关}: 由定理~\ref{thm:indep-product}, $\Cov(X,Y)=\E[SX^2]=\E[S]\E[X^2]=0$.
\item \emph{不独立}: $\abs{Y}=\abs{X}$, 于是 $\P(\abs{X}\le1,\abs{Y}\le1)=\P(\abs{X}\le1)\ne\P(\abs{X}\le1)^2$.
\item \emph{$(X,Y)$ 不是正态随机向量}: $X+Y=(1+S)X$ 以概率 $\frac{1}{2}$ 等于 $0$, 但不恒为 $0$. 非退化正态分布没有原子, 退化正态分布是常数, 所以 $X+Y$ 不服从正态分布.
\end{itemize}
如图~\ref{fig:normal-not-joint}, $(X,Y)$ 的全部概率集中在两条对角线上, 完全不像图~\ref{fig:gaussian-contours} 中的椭圆.
\end{example}

\begin{figure}[htbp]
\centering
\begin{tikzpicture}[scale=0.85]
  \draw[->,SADark!70] (-3,0) -- (3.2,0) node[right,font=\small] {$x$};
  \draw[->,SADark!70] (0,-3) -- (0,3.2) node[above,font=\small] {$y$};
  \foreach \t in {0,0.1,...,2.8} {
    \pgfmathsetmacro{\o}{90*exp(-\t*\t/2)+5}
    \fill[SARed!\o] (\t,\t) circle (2.4pt);
    \fill[SARed!\o] (-\t,-\t) circle (2.4pt);
    \fill[SABlue!\o] (\t,-\t) circle (2.4pt);
    \fill[SABlue!\o] (-\t,\t) circle (2.4pt);
  }
  \node[SARed,font=\small,anchor=west] at (2.9,2.7) {$S=1$: $y=x$};
  \node[SABlue,font=\small,anchor=west] at (2.9,-2.7) {$S=-1$: $y=-x$};
\end{tikzpicture}
\caption{例~\ref{ex:normal-not-joint} 中 $(X,Y)$ 的分布: 概率以各 $\frac{1}{2}$ 的比例分布在两条对角线上 (颜色深浅表示密度大小). 两个边缘分布都是 $N(0,1)$, 但联合分布不是二维正态}
\label{fig:normal-not-joint}
\end{figure}

定理~\ref{thm:gaussian-uncorrelated} 的一个直接应用是把一个正态变量分解为 ``可由另一个变量线性解释的部分'' 与 ``独立的噪声''. 它预示了第~2~章的条件期望, 也是第~4~章研究布朗桥的基本工具.

\begin{example}[正态变量的正交分解]\label{ex:gaussian-regression}
设 $(X,Y)$ 是二维正态随机向量, $\Var(Y)>0$. 令
\[
b=\frac{\Cov(X,Y)}{\Var(Y)},\qquad W=X-bY.
\]
由命题~\ref{prop:gaussian-operations} (1), $(W,Y)$ 是正态随机向量 (它是 $(X,Y)$ 的线性变换). 又
$\Cov(W,Y)=\Cov(X,Y)-b\Var(Y)=0$, 由定理~\ref{thm:gaussian-uncorrelated}, $W$ 与 $Y$ 独立. 于是
\[
X=bY+W,\qquad W\ \text{与}\ Y\ \text{独立},\quad W\sim N\big(\E X-b\E Y,\ \Var(X)-b^2\Var(Y)\big).
\]
例如对布朗运动, 取 $X=B_s$, $Y=B_t$ ($0<s<t$), 则 $\Cov(B_s,B_t)=s$, $b=\frac{s}{t}$, 从而 $B_s=\frac{s}{t}B_t+W$, $W$ 与 $B_t$ 独立: 已知终点 $B_t$ 时, $B_s$ 的 ``最佳猜测'' 是线性插值 $\frac{s}{t}B_t$. 方差的计算: $\Var(W)=\Var(X)-2b\Cov(X,Y)+b^2\Var(Y)=\Var(X)-b^2\Var(Y)$.
\end{example}

% =====================================================================
\section{平方可积随机变量空间 \texorpdfstring{$L^2(\Omega,\mathcal{F},\P)$}{L2}}
\label{sec:1-6}

前面几节分别讨论了方差、协方差与正态分布. 本节换一个视角: 把平方可积的随机变量看作一个无穷维空间中的\emph{向量}. 在这个空间里, 标准差是向量的长度, 协方差是内积, 相关系数是夹角的余弦, 而 ``用已知信息对未知量作最佳预测'' 就是向子空间作正交投影. 这一几何观点有两个直接的后续: 第~2~章把条件期望定义为正交投影; 第~5~章把 Itô 积分构造为一个保持长度的映射 (Itô 等距). 两者都依赖于本节证明的完备性.

% ---------------------------------------------------------------------
\subsection{\texorpdfstring{$L^p$}{Lp} 空间与 \texorpdfstring{$L^2$}{L2} 的内积}

几乎必然相等的随机变量具有相同的期望 (定理~\ref{thm:expectation-properties} (4)), 在讨论矩和范数时无法区分. 因此我们把它们视为同一个元素.

\begin{definition}[$L^p$ 空间]\label{def:Lp}
设 $p\ge1$. 在满足 $\E\abs{X}^p<\infty$ 的随机变量全体上, 定义等价关系 $X\sim Y\iff X=Y$ $\as$, 其等价类全体记为 $L^p(\Omega,\mathcal{F},\P)$, 简记为 $L^p$. 在 $L^p$ 上定义 $\norm{X}_p=(\E\abs{X}^p)^{1/p}$.
\end{definition}

按照惯例, 我们仍把 $L^p$ 的元素称为随机变量, 并用其中任意一个代表元进行运算; 由于期望不依赖于零概率集上的取值, 加法、数乘、范数都与代表元的选取无关.

\begin{proposition}\label{prop:Lp-normed}
对 $p\ge1$, $(L^p,\norm{\cdot}_p)$ 是赋范线性空间.
\end{proposition}

\begin{proof}
由 Minkowski 不等式~\ref{thm:minkowski}, $L^p$ 对加法封闭且三角不等式成立; 齐次性 $\norm{aX}_p=\abs{a}\norm{X}_p$ 是显然的; 最后, $\norm{X}_p=0$ 当且仅当 $\E\abs{X}^p=0$, 由推论~\ref{cor:nonneg-properties} (3), 当且仅当 $X=0$ $\as$, 即 $X$ 是 $L^p$ 中的零元.
\end{proof}

$p=2$ 的情形有额外的结构: 内积.

\begin{proposition}[$L^2$ 的内积]\label{prop:L2-inner}
对 $X,Y\in L^2$ 定义 $\ip{X,Y}=\E[XY]$. 则
\begin{enumerate}
\item $\ip{\cdot,\cdot}$ 是 $L^2$ 上的内积 (对称、双线性、正定), 且 $\ip{X,X}=\norm{X}_2^2$;
\item Cauchy-Schwarz 不等式: $\abs{\ip{X,Y}}\le\norm{X}_2\norm{Y}_2$, 等号成立当且仅当 $X,Y$ 线性相关 (即存在不全为零的 $a,b$ 使 $aX+bY=0$ $\as$);
\item 平行四边形法则: $\norm{X+Y}_2^2+\norm{X-Y}_2^2=2\norm{X}_2^2+2\norm{Y}_2^2$;
\item 勾股定理: 若 $\ip{X,Y}=0$ (称 $X$ 与 $Y$ \textbf{正交}, 记作 $X\perp Y$), 则 $\norm{X+Y}_2^2=\norm{X}_2^2+\norm{Y}_2^2$.
\end{enumerate}
\end{proposition}

\begin{proof}
以下记 $\norm{\cdot}=\norm{\cdot}_2$. (1) 由 H\"older 不等式~\ref{thm:holder}, $XY\in L^1$, 内积有定义. 对称与双线性来自期望的线性; 正定性即命题~\ref{prop:Lp-normed} 中 $\norm{X}=0\iff X=0$.

(2) 若 $\norm{Y}=0$, 则 $Y=0$, 两边为 $0$, 且 $X,Y$ 线性相关. 设 $\norm{Y}>0$. 对任意 $t\in\R$, 由双线性,
\[
0\le\norm{X-tY}^2=\norm{X}^2-2t\ip{X,Y}+t^2\norm{Y}^2.
\]
这个关于 $t$ 的二次函数非负, 判别式 $4\ip{X,Y}^2-4\norm{X}^2\norm{Y}^2\le0$, 即得不等式. 等号成立当且仅当判别式为零, 当且仅当二次函数有实根 $t_0$, 即 $\norm{X-t_0Y}=0$, $X=t_0Y$ $\as$; 而当 $\norm{Y}>0$ 时, 线性相关恰好意味着 $X$ 是 $Y$ 的倍数.

(3)(4) 展开 $\norm{X\pm Y}^2=\norm{X}^2\pm2\ip{X,Y}+\norm{Y}^2$, 两式相加得 (3); 当 $\ip{X,Y}=0$ 时得 (4).
\end{proof}

现在可以用几何语言重新解读 1.3 节的数字特征. 记中心化的随机变量为 $\tilde X=X-\E X$. 则
\[
\norm{\tilde X}_2^2=\Var(X),\qquad \ip{\tilde X,\tilde Y}=\Cov(X,Y),\qquad \E[X]=\ip{X,1}.
\]
标准差是中心化向量的长度, 协方差是中心化向量的内积. 于是自然地定义夹角.

\begin{definition}[相关系数]\label{def:correlation}
设 $X,Y\in L^2$ 且 $\Var(X),\Var(Y)>0$. 称
\[
\rho(X,Y)=\frac{\Cov(X,Y)}{\sqrt{\Var(X)}\sqrt{\Var(Y)}}=\frac{\ip{\tilde X,\tilde Y}}{\norm{\tilde X}_2\norm{\tilde Y}_2}
\]
为 $X$ 与 $Y$ 的\textbf{相关系数}. 它是 $\tilde X$ 与 $\tilde Y$ 夹角 $\theta$ 的余弦. $\rho=0$ 时称 $X$ 与 $Y$ \textbf{不相关}, 即 $\tilde X\perp\tilde Y$.
\end{definition}

\begin{proposition}\label{prop:correlation}
$\abs{\rho(X,Y)}\le1$; 且 $\abs{\rho(X,Y)}=1$ 当且仅当存在常数 $a\ne0$, $b$ 使 $Y=aX+b$ $\as$, 此时 $a$ 与 $\rho$ 同号.
\end{proposition}

\begin{proof}
对 $\tilde X,\tilde Y$ 用 Cauchy-Schwarz 不等式 (命题~\ref{prop:L2-inner} (2)) 得 $\abs{\rho}\le1$. 等号成立当且仅当 $\tilde X,\tilde Y$ 线性相关; 由于二者都非零, 这等价于 $\tilde Y=a\tilde X$ ($a\ne0$), 即 $Y=aX+(\E Y-a\E X)$. 此时 $\Cov(X,Y)=a\Var(X)$ 与 $a$ 同号.
\end{proof}

\begin{figure}[htbp]
\centering
\begin{tikzpicture}[scale=1.2]
  \coordinate (O) at (0,0);
  \coordinate (X) at (3.2,0.4);
  \coordinate (Y) at (1.6,2.2);
  \draw[->,SABlue,very thick] (O) -- (X) node[below right,font=\small] {$\tilde X=X-\E X$};
  \draw[->,SARed,very thick] (O) -- (Y) node[above,font=\small] {$\tilde Y=Y-\E Y$};
  \draw[SADark] ($(O)!0.55cm!(X)$) arc[start angle={atan(0.4/3.2)}, end angle={atan(2.2/1.6)}, radius=0.55cm];
  \node[font=\small] at (0.62,0.42) {$\theta$};
  \node[font=\small,align=left,anchor=west] at (5.6,1.8) {长度: $\norm{\tilde X}_2=\sigma_X$};
  \node[font=\small,align=left,anchor=west] at (5.6,1.1) {内积: $\ip{\tilde X,\tilde Y}=\Cov(X,Y)$};
  \node[font=\small,align=left,anchor=west] at (5.6,0.4) {夹角: $\cos\theta=\rho(X,Y)$};
\end{tikzpicture}
\caption{$L^2$ 中的几何: 中心化后, 标准差是长度, 协方差是内积, 相关系数是夹角余弦. 不相关即正交; 方差可加性 $\Var(X+Y)=\Var(X)+\Var(Y)$ ($\rho=0$ 时) 就是勾股定理}
\label{fig:correlation-angle}
\end{figure}

% ---------------------------------------------------------------------
\subsection{完备性}

有限维欧氏空间的一个基本性质是 Cauchy 序列必收敛. 无穷维空间未必如此 (例如区间上的连续函数全体在 $L^2$ 范数下不完备). $L^p$ 空间的完备性是 Lebesgue 积分理论最重要的成果之一; 在随机分析中, 它保证了 ``用简单过程逼近, 再取极限'' 的构造方法 (例如 Itô 积分) 能够得到一个真正存在的对象.

\begin{definition}[Cauchy 序列]
称 $\{X_n\}\subseteq L^p$ 为 \textbf{Cauchy 序列}, 如果对任意 $\varepsilon>0$, 存在 $N$, 使得 $m,n\ge N$ 时 $\norm{X_n-X_m}_p<\varepsilon$.
\end{definition}

\begin{theorem}[Riesz-Fischer 定理]\label{thm:riesz-fischer}
设 $p\ge1$. 若 $\{X_n\}$ 是 $L^p$ 中的 Cauchy 序列, 则存在 $X\in L^p$ 使 $X_n\xrightarrow{L^p}X$, 并且存在子列 $X_{n_k}\xrightarrow{\as}X$. 即 $L^p$ 是完备的赋范线性空间 (Banach 空间); 特别地, $L^2$ 是 Hilbert 空间.
\end{theorem}

\begin{proof}
\emph{第一步: 选取快速收敛的子列.} 由 Cauchy 性, 可取 $n_1<n_2<\cdots$, 使得 $\norm{X_n-X_m}_p\le2^{-k}$ 对一切 $m,n\ge n_k$ 成立. 特别地 $\norm{X_{n_{k+1}}-X_{n_k}}_p\le2^{-k}$.

\emph{第二步: 子列几乎必然收敛.} 令 $G_K=\sum_{k=1}^K\abs{X_{n_{k+1}}-X_{n_k}}$, $G=\lim_KG_K=\sum_{k=1}^\infty\abs{X_{n_{k+1}}-X_{n_k}}$. 由 Minkowski 不等式~\ref{thm:minkowski},
$\norm{G_K}_p\le\sum_{k=1}^K\norm{X_{n_{k+1}}-X_{n_k}}_p\le\sum_{k\ge1}2^{-k}=1$. 由于 $G_K^p\uparrow G^p$, 由单调收敛定理~\ref{thm:mct}, $\E[G^p]=\lim_K\E[G_K^p]\le1$, 从而由推论~\ref{cor:nonneg-properties} (4), $G<\infty$ $\as$ 在 $\{G<\infty\}$ 上, 级数
\[
X_{n_1}+\sum_{k=1}^\infty\big(X_{n_{k+1}}-X_{n_k}\big)
\]
绝对收敛, 而其部分和恰为 $X_{n_{K+1}}$. 所以 $X_{n_k}$ 几乎必然收敛. 令 $X=\lim_kX_{n_k}$ (在极限不存在的零概率集上令 $X=0$), 由命题~\ref{prop:rv-operations} (3), $X$ 是随机变量.

\emph{第三步: $L^p$ 收敛.} 给定 $\varepsilon>0$, 取 $N$ 使 $m,n\ge N$ 时 $\norm{X_n-X_m}_p<\varepsilon$. 固定 $n\ge N$, 对 $\abs{X_n-X_{n_k}}^p\to\abs{X_n-X}^p$ ($k\to\infty$, $\as$) 使用 Fatou 引理~\ref{thm:fatou},
\[
\E\abs{X_n-X}^p\le\liminf_{k\to\infty}\E\abs{X_n-X_{n_k}}^p\le\varepsilon^p,
\]
这里用到当 $k$ 充分大时 $n_k\ge N$. 于是 $X_n-X\in L^p$, 从而 $X=X_n-(X_n-X)\in L^p$, 且 $\norm{X_n-X}_p\le\varepsilon$ 对一切 $n\ge N$ 成立.
\end{proof}

\begin{proposition}[内积的连续性]\label{prop:inner-continuity}
若 $X_n\xrightarrow{L^2}X$, $Y_n\xrightarrow{L^2}Y$, 则 $\ip{X_n,Y_n}\to\ip{X,Y}$. 特别地, $\E X_n\to\E X$, $\Var(X_n)\to\Var(X)$, $\Cov(X_n,Y_n)\to\Cov(X,Y)$.
\end{proposition}

\begin{proof}
由三角不等式 $\norm{Y_n}_2\le\norm{Y}_2+\norm{Y_n-Y}_2$, 序列 $\norm{Y_n}_2$ 有界. 由 Cauchy-Schwarz 不等式,
\begin{align*}
\abs{\ip{X_n,Y_n}-\ip{X,Y}}&\le\abs{\ip{X_n-X,Y_n}}+\abs{\ip{X,Y_n-Y}}\\
&\le\norm{X_n-X}_2\norm{Y_n}_2+\norm{X}_2\norm{Y_n-Y}_2\to0.
\end{align*}
取 $Y_n=Y=1$ 得 $\E X_n\to\E X$; 再由 $\Var(X_n)=\ip{X_n,X_n}-(\E X_n)^2$ 及协方差的类似表达式得其余结论.
\end{proof}

\begin{example}[独立随机变量级数]\label{ex:L2-series}
设 $\xi_1,\xi_2,\dots$ 独立, $\E\xi_k=0$, $\Var(\xi_k)=1$, 实数列 $\{a_k\}$ 满足 $\sum_ka_k^2<\infty$. 令 $S_n=\sum_{k=1}^na_k\xi_k$. 由推论~\ref{cor:indep-consequences} (2), $\xi_k$ 两两正交, 于是对 $m<n$,
\[
\norm{S_n-S_m}_2^2=\Big\|\sum_{k=m+1}^na_k\xi_k\Big\|_2^2=\sum_{k=m+1}^na_k^2\to0\quad(m,n\to\infty).
\]
所以 $\{S_n\}$ 是 $L^2$ 中的 Cauchy 序列. 由 Riesz-Fischer 定理~\ref{thm:riesz-fischer}, 级数 $S=\sum_{k=1}^\infty a_k\xi_k$ 在 $L^2$ 中收敛, 且由内积的连续性, $\E S=0$, $\Var(S)=\sum_ka_k^2$. 例如 $\sum_k\frac{\xi_k}{k}$ 收敛, 而 $\xi_k$ 取 $\pm1$ 时这是 ``随机符号的调和级数''. 第~4~章将用这种级数 (取 $\xi_k$ 为标准正态变量) 构造布朗运动.
\end{example}

% ---------------------------------------------------------------------
\subsection{正交投影}

设我们观察到了随机变量 $Y$, 想用 $Y$ 的某个函数去预测未知的 $X$, 并以均方误差 $\E[(X-\hat X)^2]$ 衡量预测的好坏. 若只允许常数预测, 最佳选择是什么? 若允许 $Y$ 的线性函数呢? 这些问题都是 ``在 $L^2$ 的某个子空间中, 求离 $X$ 最近的点''.

\begin{definition}[闭子空间]
$L^2$ 的线性子空间 $\mathcal{K}$ 称为\textbf{闭}的, 如果它对 $L^2$ 极限封闭: 若 $X_n\in\mathcal{K}$ 且 $X_n\xrightarrow{L^2}X$, 则 $X\in\mathcal{K}$.
\end{definition}

\begin{example}\label{ex:closed-subspaces}
\begin{enumerate}
\item 有限维子空间总是闭的, 例如常数全体 $\{c\cdot1: c\in\R\}$, 以及 $\mathrm{span}\{1,Y_1,\dots,Y_n\}$ (这是线性代数中的事实: 有限维赋范空间完备).
\item \textbf{子 $\sigma$-代数对应的子空间.} 设 $\mathcal{G}\subseteq\mathcal{F}$ 是子 $\sigma$-代数. 令 $L^2(\mathcal{G})$ 为含有 $\mathcal{G}$-可测代表元的 $L^2$ 元素全体, 即 ``只依赖于信息 $\mathcal{G}$'' 的平方可积随机变量. 它显然是线性子空间. 它是闭的: 若 $X_n\in L^2(\mathcal{G})$ ($X_n$ 取为 $\mathcal{G}$-可测的代表元) 且 $X_n\xrightarrow{L^2}X$, 由 Riesz-Fischer 定理~\ref{thm:riesz-fischer}, 存在子列 $X_{n_k}\to X$ $\as$ 令 $X'=\limsup_kX_{n_k}$ (在其不是有限实数处令 $X'=0$), 由命题~\ref{prop:rv-operations} (3) (在可测空间 $(\Omega,\mathcal{G})$ 上使用), $X'$ 是 $\mathcal{G}$-可测的, 且 $X'=X$ $\as$ 所以 $X\in L^2(\mathcal{G})$. 这个子空间将是第~2~章的主角.
\end{enumerate}
\end{example}

\begin{theorem}[正交投影定理]\label{thm:projection}
设 $\mathcal{K}$ 是 $L^2$ 的闭子空间, $X\in L^2$. 则存在唯一的 $Y\in\mathcal{K}$, 使得
\[
\norm{X-Y}_2=\inf_{W\in\mathcal{K}}\norm{X-W}_2.
\]
并且 $Y$ 由以下性质刻画: $Y\in\mathcal{K}$, 且 $X-Y\perp\mathcal{K}$ (即对一切 $W\in\mathcal{K}$, $\ip{X-Y,W}=0$). 称 $Y$ 为 $X$ 在 $\mathcal{K}$ 上的\textbf{正交投影}, 记作 $\Pi_{\mathcal{K}}X$.
\end{theorem}

\begin{proof}
记 $\norm{\cdot}=\norm{\cdot}_2$, $d=\inf_{W\in\mathcal{K}}\norm{X-W}$.

\emph{存在性.} 取 $Y_n\in\mathcal{K}$ 使 $\norm{X-Y_n}^2<d^2+\frac{1}{n}$. 对 $X-Y_n$ 与 $X-Y_m$ 用平行四边形法则 (命题~\ref{prop:L2-inner} (3)),
\[
\norm{Y_m-Y_n}^2=2\norm{X-Y_n}^2+2\norm{X-Y_m}^2-4\Big\|X-\frac{Y_n+Y_m}{2}\Big\|^2.
\]
由于 $\mathcal{K}$ 是线性子空间, $\frac{Y_n+Y_m}{2}\in\mathcal{K}$, 所以最后一项中的范数平方 $\ge d^2$. 于是
\[
\norm{Y_m-Y_n}^2\le2\Big(d^2+\frac{1}{n}\Big)+2\Big(d^2+\frac{1}{m}\Big)-4d^2=\frac{2}{n}+\frac{2}{m}.
\]
因此 $\{Y_n\}$ 是 Cauchy 序列. 由 Riesz-Fischer 定理~\ref{thm:riesz-fischer}, $Y_n\xrightarrow{L^2}Y$; 由 $\mathcal{K}$ 闭, $Y\in\mathcal{K}$. 由范数的连续性 ($\abs{\norm{A}-\norm{B}}\le\norm{A-B}$), $\norm{X-Y}=\lim_n\norm{X-Y_n}=d$.

\emph{唯一性.} 若 $Y,Y'\in\mathcal{K}$ 都满足 $\norm{X-Y}=\norm{X-Y'}=d$, 与上面同样用平行四边形法则, $\norm{Y-Y'}^2\le2d^2+2d^2-4d^2=0$, 故 $Y=Y'$.

\emph{刻画: 最近点 $\Rightarrow$ 正交.} 设 $Y$ 是最近点, $W\in\mathcal{K}$. 对 $t\in\R$, $Y+tW\in\mathcal{K}$, 故
\[
f(t)=\norm{X-Y-tW}^2=d^2-2t\ip{X-Y,W}+t^2\norm{W}^2\ge d^2=f(0).
\]
二次函数 $f$ 在 $t=0$ 取最小值, 故 $f'(0)=-2\ip{X-Y,W}=0$.

\emph{刻画: 正交 $\Rightarrow$ 最近点.} 设 $Y\in\mathcal{K}$ 且 $X-Y\perp\mathcal{K}$. 对任意 $W\in\mathcal{K}$, $Y-W\in\mathcal{K}$, 由勾股定理 (命题~\ref{prop:L2-inner} (4)),
\[
\norm{X-W}^2=\norm{(X-Y)+(Y-W)}^2=\norm{X-Y}^2+\norm{Y-W}^2\ge\norm{X-Y}^2.\qedhere
\]
\end{proof}

\begin{figure}[htbp]
\centering
\begin{tikzpicture}[scale=1.1]
  \fill[SABlue!10] (-0.5,-0.6) -- (5.2,-0.6) -- (6.6,1.4) -- (0.9,1.4) -- cycle;
  \draw[SABlue!60] (-0.5,-0.6) -- (5.2,-0.6) -- (6.6,1.4) -- (0.9,1.4) -- cycle;
  \node[SABlue!80!black,font=\small] at (5.9,0.3) {$\mathcal{K}$};
  \coordinate (O) at (0.6,0);
  \coordinate (Y) at (3.6,0.5);
  \coordinate (X) at (3.6,3.1);
  \coordinate (W) at (1.9,0.9);
  \fill (O) circle (1.2pt) node[below left,font=\small] {$0$};
  \draw[->,SADark,very thick] (O) -- (X) node[above,font=\small] {$X$};
  \draw[->,SABlue,very thick] (O) -- (Y) node[below right,font=\small] {$\Pi_{\mathcal{K}}X$};
  \draw[SARed,very thick,dashed] (Y) -- (X) node[midway,right,font=\small] {$X-\Pi_{\mathcal{K}}X\perp\mathcal{K}$};
  \draw[SARed] ($(Y)+(0,0.22)$) -- ($(Y)+(-0.24,0.18)$) -- ($(Y)+(-0.24,-0.04)$);
  \fill[SAGreenText] (W) circle (1.4pt) node[below,font=\small] {$W$};
  \draw[SAGreenText,dashed] (W) -- (X);
  \node[SAGreenText,font=\small,anchor=east] at (2.6,2.2) {$\norm{X-W}_2$ 更大};
\end{tikzpicture}
\caption{正交投影: 子空间 $\mathcal{K}$ 中离 $X$ 最近的点 $\Pi_{\mathcal{K}}X$, 恰是使残差 $X-\Pi_{\mathcal{K}}X$ 与 $\mathcal{K}$ 垂直的那个点}
\label{fig:projection}
\end{figure}

\begin{proposition}[投影算子的性质]\label{prop:projection-properties}
设 $\mathcal{K},\mathcal{K}_1\subseteq\mathcal{K}_2$ 是 $L^2$ 的闭子空间, $\Pi=\Pi_{\mathcal{K}}$, $X,Z\in L^2$, $a,b\in\R$. 则
\begin{enumerate}
\item 线性: $\Pi(aX+bZ)=a\Pi X+b\Pi Z$;
\item $\Pi X=X$ 当且仅当 $X\in\mathcal{K}$; 特别地 $\Pi(\Pi X)=\Pi X$;
\item 勾股分解: $\norm{X}_2^2=\norm{\Pi X}_2^2+\norm{X-\Pi X}_2^2$, 从而 $\norm{\Pi X}_2\le\norm{X}_2$;
\item 自伴性: $\ip{\Pi X,Z}=\ip{X,\Pi Z}$;
\item 塔性质: $\Pi_{\mathcal{K}_1}\big(\Pi_{\mathcal{K}_2}X\big)=\Pi_{\mathcal{K}_1}X$.
\end{enumerate}
\end{proposition}

\begin{proof}
各条都利用定理~\ref{thm:projection} 中的刻画: 要证 $Y=\Pi X$, 只需验证 $Y\in\mathcal{K}$ 且 $X-Y\perp\mathcal{K}$.

(1) $a\Pi X+b\Pi Z\in\mathcal{K}$, 且 $(aX+bZ)-(a\Pi X+b\Pi Z)=a(X-\Pi X)+b(Z-\Pi Z)$ 与 $\mathcal{K}$ 中每个元素正交.

(2) 若 $X\in\mathcal{K}$, 则 $Y=X$ 满足刻画 ($X-Y=0$). 反之 $\Pi X\in\mathcal{K}$. 由于 $\Pi X\in\mathcal{K}$, 再用一次得 $\Pi(\Pi X)=\Pi X$.

(3) $X=\Pi X+(X-\Pi X)$, 两项正交, 由勾股定理即得.

(4) 由于 $\Pi X\in\mathcal{K}$ 而 $Z-\Pi Z\perp\mathcal{K}$, $\ip{\Pi X,Z}=\ip{\Pi X,\Pi Z}+\ip{\Pi X,Z-\Pi Z}=\ip{\Pi X,\Pi Z}$. 同理 $\ip{X,\Pi Z}=\ip{\Pi X,\Pi Z}$.

(5) 记 $Y=\Pi_{\mathcal{K}_1}(\Pi_{\mathcal{K}_2}X)\in\mathcal{K}_1$. 写
$X-Y=(X-\Pi_{\mathcal{K}_2}X)+(\Pi_{\mathcal{K}_2}X-Y)$. 对 $W\in\mathcal{K}_1$: 由于 $W\in\mathcal{K}_2$, 第一项与 $W$ 正交; 由 $Y$ 的定义, 第二项与 $\mathcal{K}_1$ 正交, 也与 $W$ 正交. 所以 $X-Y\perp\mathcal{K}_1$, $Y=\Pi_{\mathcal{K}_1}X$.
\end{proof}

\begin{intuition}
(5) 说: 先投影到大子空间, 再投影到小子空间, 等于直接投影到小子空间. 把子空间理解为 ``信息'', 这就是说 ``先用较多信息作预测, 再用较少信息对这个预测作预测, 结果等于直接用较少信息作预测''. 第~2~章中, 它将成为条件期望的塔性质 $\E[\E[X\,|\,\mathcal{G}]\,|\,\mathcal{H}]=\E[X\,|\,\mathcal{H}]$ ($\mathcal{H}\subseteq\mathcal{G}$).
\end{intuition}

\begin{example}[最佳常数预测]\label{ex:best-constant}
取 $\mathcal{K}=\{c\cdot1: c\in\R\}$. 由于 $\ip{X-\E X,c}=c\,\E[X-\E X]=0$, 由投影的刻画, $\Pi_{\mathcal{K}}X=\E X$. 所以 $\min_c\E[(X-c)^2]=\E[(X-\E X)^2]=\Var(X)$: 期望是均方意义下最好的常数预测, 方差是相应的最小误差.
\end{example}

\begin{example}[最佳线性预测]\label{ex:best-linear}
设 $X,Y\in L^2$, $\Var(Y)>0$. 求 $a,b$ 使 $\E[(X-a-bY)^2]$ 最小, 即求 $X$ 在 $\mathcal{K}=\mathrm{span}\{1,Y\}$ 上的投影. 由刻画, 只需 $X-a-bY$ 与 $1$ 和 $Y$ 都正交:
\[
\E[X-a-bY]=0,\qquad \E[(X-a-bY)Y]=0.
\]
由第一式 $a=\E X-b\E Y$, 代入第二式得 $\E[(\tilde X-b\tilde Y)Y]=\Cov(X,Y)-b\Var(Y)=0$, 所以
\[
b=\frac{\Cov(X,Y)}{\Var(Y)},\qquad \Pi_{\mathcal{K}}X=\E X+\frac{\Cov(X,Y)}{\Var(Y)}(Y-\E Y).
\]
由勾股分解, 最小均方误差为
\[
\E\big[(X-\Pi_{\mathcal{K}}X)^2\big]=\Var(X)-b^2\Var(Y)=\Var(X)\big(1-\rho(X,Y)^2\big).
\]
所以 $\rho^2$ 是 $X$ 的方差中能被 $Y$ \emph{线性}解释的比例. 若 $(X,Y)$ 是正态随机向量, 例~\ref{ex:gaussian-regression} 表明残差 $X-\Pi_{\mathcal{K}}X$ 与 $Y$ 独立, 这时线性预测已经用尽了 $Y$ 中关于 $X$ 的全部信息. 在非正态情形下, 线性预测未必能提取全部信息. 若 $X\sim N(0,1)$, $Y=X^2$, 则 $\rho(X,Y)=0$ (例~\ref{ex:uncorrelated-dependent}), 用 $X$ 对 $Y$ 作最佳线性预测只能得到常数 $\E Y=1$, 而 $Y$ 其实是 $X$ 的函数, 可以被完全预测. 要利用 $X$ 的全部信息, 应当在更大的闭子空间 $L^2(\sigma(X))$ (例~\ref{ex:closed-subspaces} (2)) 上投影, 这正是第~2~章条件期望 $\E[Y\,|\,X]$ 的定义.
\end{example}

% ---------------------------------------------------------------------
\subsection{正态随机变量的 \texorpdfstring{$L^2$}{L2} 极限}

最后证明一个对第~5~章至关重要的基本事实, 即正态性在 $L^2$ 极限下保持. Wiener 积分 $\int_0^tf(s)\dd B_s$ 定义为黎曼和 $\sum_if(t_i)(B_{t_{i+1}}-B_{t_i})$ 的 $L^2$ 极限, 每个黎曼和都是正态的, 由下面的定理, 极限也是正态的.

\begin{theorem}[正态性在 $L^2$ 极限下保持]\label{thm:gaussian-L2-limit}
设 $X_n\sim N(\mu_n,\sigma_n^2)$, 且 $X_n\xrightarrow{L^2}X$. 则极限 $\mu=\lim_n\mu_n$ 与 $\sigma^2=\lim_n\sigma_n^2$ 存在, 且 $X\sim N(\mu,\sigma^2)$.
\end{theorem}

\begin{proof}
由内积的连续性 (命题~\ref{prop:inner-continuity}), $\mu_n=\E X_n\to\E X$, $\sigma_n^2=\Var(X_n)\to\Var(X)$. 记 $\mu=\E X$, $\sigma^2=\Var(X)$.

对实数 $a,b$, $\abs{e^{\mathrm{i}a}-e^{\mathrm{i}b}}=\Big|\int_b^a\mathrm{i}e^{\mathrm{i}s}\dd s\Big|\le\abs{a-b}$. 于是由引理~\ref{lem:complex-abs} 与 Lyapunov 不等式~\ref{cor:lyapunov}, 对每个 $t$,
\[
\abs{\varphi_{X_n}(t)-\varphi_X(t)}\le\E\abs{e^{\mathrm{i}tX_n}-e^{\mathrm{i}tX}}\le\abs{t}\,\E\abs{X_n-X}\le\abs{t}\,\norm{X_n-X}_2\to0.
\]
另一方面, 由 \eqref{eq:normal-cf},
\[
\varphi_{X_n}(t)=\exp\Big(\mathrm{i}\mu_nt-\frac{1}{2}\sigma_n^2t^2\Big)\to\exp\Big(\mathrm{i}\mu t-\frac{1}{2}\sigma^2t^2\Big).
\]
因此 $\varphi_X(t)=\exp(\mathrm{i}\mu t-\frac{1}{2}\sigma^2t^2)$, 由唯一性定理~\ref{thm:cf-uniqueness}, $X\sim N(\mu,\sigma^2)$.
\end{proof}

\begin{corollary}[高斯族的闭线性张成]\label{cor:gaussian-space}
设 $\{G_i\}_{i\in I}$ 是一族随机变量, 其中任意有限个都构成正态随机向量. 令 $\mathcal{H}$ 为 $\{G_i\}_{i\in I}$ 的线性张成在 $L^2$ 中的闭包. 则 $\mathcal{H}$ 中任意有限个元素 $H_1,\dots,H_n$ 构成正态随机向量.
\end{corollary}

\begin{proof}
$\mathcal{H}$ 是线性子空间 (线性子空间的闭包仍是线性子空间: 若 $A_k\to A$, $B_k\to B$, 则 $aA_k+bB_k\to aA+bB$). 对 $\rvec{a}\in\R^n$, $\sum_ja_jH_j\in\mathcal{H}$, 所以只需证明 $\mathcal{H}$ 中每个元素服从正态分布. 设 $H\in\mathcal{H}$, 则存在 $\{G_i\}$ 的有限线性组合 $L_k$ 使 $L_k\xrightarrow{L^2}H$. 每个 $L_k$ 是有限个 $G_i$ 组成的正态随机向量的线性组合, 由定义~\ref{def:gaussian-vector}, $L_k$ 服从正态分布. 由定理~\ref{thm:gaussian-L2-limit}, $H$ 服从正态分布.
\end{proof}

\begin{example}
在例~\ref{ex:L2-series} 中取 $\xi_k=Z_k$ 为 $\iid$ 的 $N(0,1)$ 随机变量. 由定理~\ref{thm:gaussian-L2-limit}, $\sum_{k=1}^\infty a_kZ_k\sim N\big(0,\sum_ka_k^2\big)$. 更进一步, 由推论~\ref{cor:gaussian-space}, 对任意平方可和数列 $\{a_k\},\{b_k\}$, $\big(\sum_ka_kZ_k,\sum_kb_kZ_k\big)$ 是二维正态随机向量, 其协方差为 $\sum_ka_kb_k$ (内积的连续性). 第~4~章正是用这种方式, 取适当的函数列 $a_k=a_k(t)$, 构造出一族联合正态、协方差为 $s\wedge t$ 的随机变量 $\{B_t\}$, 即布朗运动.
\end{example}

% ---------------------------------------------------------------------
\subsection*{本章小结}
\addcontentsline{toc}{subsection}{本章小结}

本章建立了全书使用的概率论框架. 表~\ref{tab:ch1-summary} 列出了本章的主要工具及其在后续各章中的作用.

\begin{center}
\captionof{table}{第~1~章的主要工具及其后续应用}
\label{tab:ch1-summary}
\renewcommand{\arraystretch}{1.4}
\begin{tabular}{>{\raggedright\arraybackslash}p{0.36\textwidth} >{\raggedright\arraybackslash}p{0.54\textwidth}}
\toprule
本章工具 & 后续应用 \\
\midrule
$\sigma$-代数、$\sigma(X)$、独立性 & 信息流 $\{\mathcal{F}_t\}$、适应过程、独立增量 (第~2、4~章) \\
$\pi$-$\lambda$ 定理 & 验证条件期望、马尔可夫性对整个 $\sigma$-代数成立 (第~2~章) \\
Borel-Cantelli 引理、0-1 律 & 布朗运动轨道性质、重对数律 (第~4~章) \\
MCT、Fatou、DCT & 条件期望的极限定理、随机积分的构造 (第~2、5~章) \\
一致可积与 Vitali 定理 & 鞅收敛定理、可选停时定理 (第~3~章) \\
特征函数、多元正态分布 & 布朗运动作为高斯过程、Girsanov 定理 (第~4、5~章) \\
$L^2$ 完备性与正交投影 & 条件期望的定义 (第~2~章)、Itô 等距与 Itô 积分 (第~5~章) \\
\bottomrule
\end{tabular}
\end{center}

% =====================================================================
%  第 2 章  条件数学期望
% =====================================================================
\chapter{条件数学期望}
\label{ch:2}

随机分析处处在问同一个问题: 已知到目前为止的信息, 对未来的随机量能作出怎样的估计? 赌徒根据已有的输赢决定下一次的赌注; 交易者为衍生品定价时需要依据历史价格, 物理学家则由粒子的当前位置预测其未来分布. 把 ``已知信息'' 表示为一个子 $\sigma$-代数 $\mathcal{G}$, 把 ``估计'' 理解为均方意义下的最佳预测, 就得到条件数学期望 $\E[X\,|\,\mathcal{G}]$. 它是全书最核心的工具: 第~3~章的鞅、第~4~章布朗运动的马尔可夫性、第~5~章 Itô 积分的鞅性质, 都用条件期望来表述.

本章先从初等的条件期望出发, 说明为什么条件期望应当是一个\emph{随机变量} (2.1 节); 然后给出一般定义并证明其存在唯一性 (2.2 节); 2.3 节系统推导条件期望的运算法则; 2.4 节对正态随机向量给出条件分布的显式公式; 最后 2.5 节引入随机过程、信息流与马尔可夫性.

% =====================================================================
\section{从初等条件期望到正交投影}
\label{sec:2-1}

% ---------------------------------------------------------------------
\subsection{对分割取条件}

在初等概率论中, 若 $\P(B)>0$, 事件 $A$ 在 $B$ 发生条件下的概率是 $\P(A\,|\,B)=\frac{\P(A\cap B)}{\P(B)}$. 同样地, 对可积随机变量 $X$, 定义 $X$ 在 $B$ 发生条件下的期望为
\begin{equation}\label{eq:elementary-ce}
\E[X\,|\,B]=\frac{\E[X;B]}{\P(B)},
\end{equation}
即 $X$ 在 $B$ 上的平均值. 例如掷一颗均匀骰子, $X$ 为点数, $B=\{2,4,6\}$, 则 $\E[X\,|\,B]=\frac{(2+4+6)/6}{1/2}=4$.

现在设观察者掌握的信息是 ``结果落在分割 $\{D_1,D_2,\dots\}$ 的哪一块中'' (例~\ref{ex:sigma-examples} (4)), 即 $\mathcal{G}=\sigma(\{D_1,D_2,\dots\})$. 观察者事先并不知道会落在哪一块, 所以他的预测应当是一个\emph{规则}: ``若结果落在 $D_i$ 中, 就报告 $X$ 在 $D_i$ 上的平均值''. 这个规则就是一个随机变量.

\begin{definition}[对可数分割的条件期望]\label{def:ce-partition}
设 $\{D_i\}_{i\ge1}$ 是 $\Omega$ 的有限或可列分割, $D_i\in\mathcal{F}$ 且 $\P(D_i)>0$, $\mathcal{G}=\sigma(\{D_i\}_{i\ge1})$, $X\in L^1$. 定义
\[
\E[X\,|\,\mathcal{G}](\omega)=\sum_{i\ge1}\E[X\,|\,D_i]\,\1_{D_i}(\omega)=\sum_{i\ge1}\frac{\E[X;D_i]}{\P(D_i)}\1_{D_i}(\omega).
\]
\end{definition}

\begin{figure}[htbp]
\centering
\begin{tikzpicture}
\begin{axis}[SAaxis, width=11cm, height=5.6cm, xmin=0, xmax=1.06, ymin=0, ymax=2.2,
  xtick={0.3,0.55,1}, xticklabels={$0.3$,$0.55$,$1$}, ytick=\empty, xlabel={$\omega$},
  legend style={draw=none,fill=none,font=\small,at={(0.5,-0.2)},anchor=north,legend columns=2,/tikz/every even column/.append style={column sep=1em}}]
  \addplot[SABlue,very thick,domain=0:1,samples=200] {1.0+0.6*sin(deg(9*x))+0.5*x};
  \addlegendentry{$X(\omega)$}
  % block averages computed: see caption
  \addplot[SARed,very thick] coordinates {(0,1.4981) (0.3,1.4981)};
  \addlegendentry{$\E[X\,|\,\mathcal{G}](\omega)$}
  \addplot[SARed,very thick,forget plot] coordinates {(0.3,0.9086) (0.55,0.9086)};
  \addplot[SARed,very thick,forget plot] coordinates {(0.55,1.5574) (1,1.5574)};
  \addplot[SADark!30,dashed,forget plot] coordinates {(0.3,0) (0.3,2.1)};
  \addplot[SADark!30,dashed,forget plot] coordinates {(0.55,0) (0.55,2.1)};
  \node[font=\small] at (axis cs:0.15,0.2) {$D_1$};
  \node[font=\small] at (axis cs:0.425,0.2) {$D_2$};
  \node[font=\small] at (axis cs:0.775,0.2) {$D_3$};
\end{axis}
\end{tikzpicture}
\caption{对分割取条件期望 (取 $\Omega=[0,1]$, $\P$ 为均匀分布): 在每一块 $D_i$ 上, $\E[X\,|\,\mathcal{G}]$ 取常值, 等于 $X$ 在该块上的平均. 它 ``抹平'' 了 $X$ 在块内的变化, 只保留观察者能区分的信息}
\label{fig:ce-partition}
\end{figure}

\begin{example}\label{ex:ce-partition}
\begin{enumerate}
\item 掷一颗骰子, $X$ 为点数, $\mathcal{G}=\sigma(\{1,3,5\},\{2,4,6\})$ 表示只知道奇偶性. 则 $\E[X\,|\,\mathcal{G}]$ 在奇数点上取 $3$, 在偶数点上取 $4$.
\item 掷两颗骰子, 点数分别为 $X_1,X_2$, $S=X_1+X_2$. 取 $\mathcal{G}=\sigma(S)$, 它由分割 $\{S=s\}$ ($s=2,\dots,12$) 生成. 由于 $(X_1,X_2)$ 与 $(X_2,X_1)$ 同分布, $\E[X_1;S=s]=\E[X_2;S=s]$, 而二者之和为 $\E[S;S=s]=s\,\P(S=s)$. 所以 $\E[X_1\,|\,S=s]=\frac{s}{2}$, 即
\[
\E[X_1\,|\,\sigma(S)]=\frac{S}{2}.
\]
知道两颗骰子的点数之和, 对第一颗骰子的最佳估计是和的一半.
\end{enumerate}
\end{example}

对分割的条件期望有两个基本性质, 它们将分别引出一般条件期望的两种定义.

\begin{proposition}\label{prop:ce-partition-properties}
在定义~\ref{def:ce-partition} 的记号下, 记 $Y=\E[X\,|\,\mathcal{G}]$. 则
\begin{enumerate}
\item $Y$ 是 $\mathcal{G}$-可测的, 且对每个 $G\in\mathcal{G}$, $\E[Y;G]=\E[X;G]$;
\item 若 $X\in L^2$, 则对每个 $\mathcal{G}$-可测的 $Z\in L^2$, $\E[(X-Y)^2]\le\E[(X-Z)^2]$.
\end{enumerate}
\end{proposition}

\begin{proof}
先说明: $\mathcal{G}$ 中的集合恰是若干块 $D_i$ 之并, 且 $\mathcal{G}$-可测的随机变量在每块 $D_i$ 上取常值. 前者是因为 ``若干块之并'' 的全体构成包含每个 $D_i$ 的 $\sigma$-代数 (与例~\ref{ex:sigma-examples} (4) 相同), 且它包含于 $\mathcal{G}$. 后者: 若 $Z$ 是 $\mathcal{G}$-可测的, $\omega\in D_i$, 则 $\{Z=Z(\omega)\}\in\mathcal{G}$ 是若干块之并且含有 $\omega$, 从而包含整个 $D_i$.

(1) $Y$ 在每块上取常值, $\{Y\in B\}$ 是若干块之并, 故 $Y$ 是 $\mathcal{G}$-可测的. 设 $G=\bigcup_{i\in I}D_i$. 由于 $\abs{X\1_{\bigcup_{i\in I,i\le n}D_i}}\le\abs{X}$, 由控制收敛定理~\ref{thm:dct}, $\E[X;G]=\sum_{i\in I}\E[X;D_i]$, 于是
\[
\E[Y;G]=\sum_{i\in I}\frac{\E[X;D_i]}{\P(D_i)}\P(D_i)=\sum_{i\in I}\E[X;D_i]=\E[X;G].
\]
(对 $Y$ 交换求和与期望同样由控制收敛定理保证, 控制函数为 $\abs{Y}$, 而 $\E\abs{Y}\le\sum_i\E[\abs{X};D_i]=\E\abs{X}$.)

(2) $Z$ 在 $D_i$ 上取常值 $z_i$. 于是 $\E[(X-Z)^2]=\sum_i\E[(X-z_i)^2;D_i]$, 而每一项
\[
\E[(X-z_i)^2;D_i]=\E[X^2;D_i]-2z_i\E[X;D_i]+z_i^2\P(D_i)
\]
是 $z_i$ 的二次函数, 在 $z_i=\frac{\E[X;D_i]}{\P(D_i)}$, 即 $Z=Y$ 时取得最小值.
\end{proof}

\begin{intuition}
(1) 表明, $Y$ 只依赖于信息 $\mathcal{G}$, 并且在 $\mathcal{G}$ 能分辨的每个事件上, $Y$ 与 $X$ 的 ``总量'' 相同. (2) 表明, 在只能利用信息 $\mathcal{G}$ 的预测中, $Y$ 的均方误差最小. 这两条都不再涉及分割, 因此可以推广到任意子 $\sigma$-代数. 之所以需要推广, 是因为若 $\mathcal{G}=\sigma(Y)$ 而 $Y$ 是连续型随机变量, 则每个 $\{Y=y\}$ 的概率都是 $0$, 公式 \eqref{eq:elementary-ce} 已经失效, 但 ``已知 $Y$ 的值时对 $X$ 的估计'' 显然是有意义的问题.
\end{intuition}

% ---------------------------------------------------------------------
\subsection{\texorpdfstring{$L^2$}{L2} 中的条件期望: 正交投影}

设 $\mathcal{G}\subseteq\mathcal{F}$ 是子 $\sigma$-代数. 由例~\ref{ex:closed-subspaces} (2), 只依赖于信息 $\mathcal{G}$ 的平方可积随机变量全体
\[
L^2(\mathcal{G})=L^2(\Omega,\mathcal{G},\P)
\]
是 $L^2=L^2(\Omega,\mathcal{F},\P)$ 的闭子空间. 命题~\ref{prop:ce-partition-properties} (2) 提示我们: 条件期望应当是 $L^2(\mathcal{G})$ 中离 $X$ 最近的元素, 而正交投影定理~\ref{thm:projection} 保证了这样的元素存在且唯一.

\begin{definition}[$L^2$ 中的条件期望]\label{def:ce-L2}
设 $X\in L^2$. 称 $X$ 在闭子空间 $L^2(\mathcal{G})$ 上的正交投影为 $X$ 关于 $\mathcal{G}$ 的\textbf{条件期望}, 记为 $\E[X\,|\,\mathcal{G}]$. 它是满足以下两条的唯一元素 $Y$ (定理~\ref{thm:projection}):
\begin{equation}\label{eq:ce-orthogonality}
Y\in L^2(\mathcal{G}),\qquad \E[(X-Y)Z]=0\quad\text{对一切 }Z\in L^2(\mathcal{G}).
\end{equation}
\end{definition}

\begin{figure}[htbp]
\centering
\begin{tikzpicture}[scale=1.1]
  \fill[SABlue!10] (-0.5,-0.6) -- (5.2,-0.6) -- (6.6,1.4) -- (0.9,1.4) -- cycle;
  \draw[SABlue!60] (-0.5,-0.6) -- (5.2,-0.6) -- (6.6,1.4) -- (0.9,1.4) -- cycle;
  \node[SABlue!80!black,font=\small] at (5.55,0.3) {$L^2(\mathcal{G})$};
  \coordinate (O) at (0.6,0);
  \coordinate (Y) at (3.6,0.5);
  \coordinate (X) at (3.6,3.1);
  \fill (O) circle (1.2pt) node[below left,font=\small] {$0$};
  \draw[->,SADark,very thick] (O) -- (X) node[above,font=\small] {$X$};
  \draw[->,SABlue,very thick] (O) -- (Y) node[below right,font=\small] {$\E[X\,|\,\mathcal{G}]$};
  \draw[SARed,very thick,dashed] (Y) -- (X) node[midway,right,font=\small,align=left] {预测误差 $X-\E[X\,|\,\mathcal{G}]$\\ 与一切 $\mathcal{G}$-可测变量正交};
  \draw[SARed] ($(Y)+(0,0.22)$) -- ($(Y)+(-0.24,0.18)$) -- ($(Y)+(-0.24,-0.04)$);
\end{tikzpicture}
\caption{$L^2$ 中的条件期望: 向 ``只依赖于信息 $\mathcal{G}$'' 的随机变量构成的闭子空间作正交投影}
\label{fig:ce-projection}
\end{figure}

在 \eqref{eq:ce-orthogonality} 中取 $Z=\1_G$ ($G\in\mathcal{G}$), 得到 $\E[Y;G]=\E[X;G]$, 这正是命题~\ref{prop:ce-partition-properties} (1) 中的积分恒等式. 由命题~\ref{prop:projection-properties}, $L^2$ 中的条件期望直接继承了投影的性质: 它是线性的; $\norm{\E[X\,|\,\mathcal{G}]}_2\le\norm{X}_2$; 若 $\mathcal{H}\subseteq\mathcal{G}$, 则 $L^2(\mathcal{H})\subseteq L^2(\mathcal{G})$, 从而由投影的塔性质 $\E\big[\E[X\,|\,\mathcal{G}]\,\big|\,\mathcal{H}\big]=\E[X\,|\,\mathcal{H}]$.

投影定义的缺点是只适用于 $X\in L^2$, 而许多重要的随机变量 (例如鞅) 只是可积的. 下一节以积分恒等式为定义, 把条件期望推广到 $L^1$, 并证明在 $L^2$ 中两种定义一致.

% =====================================================================
\section{一般条件期望}
\label{sec:2-2}

% ---------------------------------------------------------------------
\subsection{Kolmogorov 定义}

\begin{definition}[条件期望, Kolmogorov]\label{def:ce}
设 $X\in L^1(\Omega,\mathcal{F},\P)$, $\mathcal{G}\subseteq\mathcal{F}$ 是子 $\sigma$-代数. 若随机变量 $Y$ 满足
\begin{enumerate}
\item $Y$ 是 $\mathcal{G}$-可测的, 且 $Y\in L^1$;
\item 对每个 $G\in\mathcal{G}$, $\E[Y;G]=\E[X;G]$,
\end{enumerate}
则称 $Y$ 为 $X$ 关于 $\mathcal{G}$ 的\textbf{条件期望}, 记为 $\E[X\,|\,\mathcal{G}]$.
\end{definition}

条件期望只在几乎必然相等的意义下确定 (见下面的推论~\ref{cor:ce-uniqueness}): 满足定义的任何一个 $Y$ 称为条件期望的一个\textbf{版本}. 涉及条件期望的等式与不等式因此都理解为几乎必然成立, 后文不再逐一注明 ``$\as$''.

下面的比较引理是证明条件期望各种性质的基本工具.

\begin{lemma}[比较引理]\label{lem:ce-compare}
设 $Y_1,Y_2$ 是 $\mathcal{G}$-可测的可积随机变量, 且对每个 $G\in\mathcal{G}$, $\E[Y_1;G]\le\E[Y_2;G]$. 则 $Y_1\le Y_2$ $\as$
\end{lemma}

\begin{proof}
对 $n\ge1$, 令 $G_n=\{Y_1-Y_2>\frac{1}{n}\}$. 由于 $Y_1-Y_2$ 是 $\mathcal{G}$-可测的, $G_n\in\mathcal{G}$. 一方面由假设 $\E[Y_1-Y_2;G_n]\le0$; 另一方面在 $G_n$ 上 $Y_1-Y_2>\frac{1}{n}$, 所以 $\E[Y_1-Y_2;G_n]\ge\frac{1}{n}\P(G_n)$. 于是 $\P(G_n)=0$, 进而
$\P(Y_1>Y_2)=\P\big(\bigcup_nG_n\big)\le\sum_n\P(G_n)=0$.
\end{proof}

\begin{corollary}[唯一性]\label{cor:ce-uniqueness}
若 $Y$ 与 $Y'$ 都满足定义~\ref{def:ce}, 则 $Y=Y'$ $\as$
\end{corollary}

\begin{proof}
二者在每个 $G\in\mathcal{G}$ 上的积分都等于 $\E[X;G]$. 由比较引理~\ref{lem:ce-compare}, $Y\le Y'$ 且 $Y'\le Y$ $\as$
\end{proof}

验证一个随机变量是条件期望时, 往往不必检查全部 $G\in\mathcal{G}$, 检查一个生成 $\pi$-系就够了.

\begin{lemma}\label{lem:ce-pi-system}
设 $X\in L^1$, $Y$ 是 $\mathcal{G}$-可测的可积随机变量, $\mathcal{P}$ 是包含 $\Omega$ 的 $\pi$-系, 且 $\sigma(\mathcal{P})=\mathcal{G}$. 若 $\E[Y;G]=\E[X;G]$ 对一切 $G\in\mathcal{P}$ 成立, 则 $Y=\E[X\,|\,\mathcal{G}]$.
\end{lemma}

\begin{proof}
令 $\mathcal{L}=\{G\in\mathcal{G}: \E[Y;G]=\E[X;G]\}$, 则 $\mathcal{P}\subseteq\mathcal{L}$. 验证 $\mathcal{L}$ 是 $\lambda$-系: (a) $\Omega\in\mathcal{P}\subseteq\mathcal{L}$. (b) 若 $A,B\in\mathcal{L}$, $A\subseteq B$, 则由 $\1_{B\setminus A}=\1_B-\1_A$ 与期望的线性, $\E[Y;B\setminus A]=\E[Y;B]-\E[Y;A]=\E[X;B]-\E[X;A]=\E[X;B\setminus A]$. (c) 若 $A_n\in\mathcal{L}$, $A_n\uparrow A$, 则 $X\1_{A_n}\to X\1_A$, 且 $\abs{X\1_{A_n}}\le\abs{X}$, 由控制收敛定理~\ref{thm:dct}, $\E[X;A_n]\to\E[X;A]$; 对 $Y$ 同理. 由 $\pi$-$\lambda$ 定理~\ref{thm:pi-lambda}, $\mathcal{G}=\sigma(\mathcal{P})\subseteq\mathcal{L}$.
\end{proof}

% ---------------------------------------------------------------------
\subsection{存在性}

\begin{theorem}[条件期望的存在唯一性]\label{thm:ce-existence}
对任意 $X\in L^1$ 与子 $\sigma$-代数 $\mathcal{G}\subseteq\mathcal{F}$, 条件期望 $\E[X\,|\,\mathcal{G}]$ 存在, 且在几乎必然相等的意义下唯一. 若 $X\in L^2$, 则它与定义~\ref{def:ce-L2} 中的正交投影一致.
\end{theorem}

\begin{proof}
唯一性即推论~\ref{cor:ce-uniqueness}. 存在性分四步证明.

\emph{第一步: $X\in L^2$.} 令 $Y$ 为 $X$ 在 $L^2(\mathcal{G})$ 上的正交投影, 取其 $\mathcal{G}$-可测的代表元. 由于 $L^2\subseteq L^1$, $Y$ 可积. 对 $G\in\mathcal{G}$, $\1_G\in L^2(\mathcal{G})$, 由 \eqref{eq:ce-orthogonality}, $\E[(X-Y)\1_G]=0$, 即 $\E[Y;G]=\E[X;G]$. 所以 $Y$ 满足定义~\ref{def:ce}. 反过来, 由唯一性, 满足定义~\ref{def:ce} 的任何版本都几乎必然等于这个投影. 这证明了 $L^2$ 中两种定义一致.

\emph{第二步: 保序性.} 设 $X_1,X_2\in L^2$, $X_1\le X_2$ $\as$, 则 $\E[X_1\,|\,\mathcal{G}]\le\E[X_2\,|\,\mathcal{G}]$. 事实上, 对 $G\in\mathcal{G}$,
\[
\E\big[\E[X_1\,|\,\mathcal{G}];G\big]=\E[X_1;G]\le\E[X_2;G]=\E\big[\E[X_2\,|\,\mathcal{G}];G\big],
\]
由比较引理~\ref{lem:ce-compare} 即得.

\emph{第三步: $X\ge0$, $X\in L^1$.} 令 $X_n=X\wedge n$, 它有界, 故属于 $L^2$, 且 $0\le X_n\uparrow X$. 令 $Y_n=\E[X_n\,|\,\mathcal{G}]$ (由第一步). 由第二步, 对每个 $n$, $0\le Y_n\le Y_{n+1}$ 在 $\mathcal{G}$-可测集 $\Omega\setminus N_n$ 上成立, 其中 $\P(N_n)=0$. 令 $N=\bigcup_nN_n\in\mathcal{G}$, 定义
\[
Y=\begin{cases}\lim_{n\to\infty}Y_n, & \text{在 } N^c \text{ 上},\\ 0, & \text{在 } N \text{ 上}.\end{cases}
\]
$Y$ 是 $\mathcal{G}$-可测的非负广义随机变量 (命题~\ref{prop:rv-operations} (3)). 对 $G\in\mathcal{G}$, 在 $\Omega$ 上 $Y_n\1_{G\cap N^c}\uparrow Y\1_{G\cap N^c}$, $X_n\1_G\uparrow X\1_G$, 由单调收敛定理~\ref{thm:mct} 与 $\P(N)=0$,
\[
\E[Y;G]=\lim_{n\to\infty}\E[Y_n;G]=\lim_{n\to\infty}\E[X_n;G]=\E[X;G].
\]
取 $G=\Omega$ 得 $\E[Y]=\E[X]<\infty$, 所以 $Y$ 可积, 且 $Y<\infty$ $\as$; 把 $Y$ 在 $\mathcal{G}$-可测零概率集 $\{Y=\infty\}$ 上改为 $0$, 积分不变. 于是 $Y$ 满足定义~\ref{def:ce}.

\emph{第四步: 一般的 $X\in L^1$.} 由第三步, $Y^+=\E[X^+\,|\,\mathcal{G}]$ 与 $Y^-=\E[X^-\,|\,\mathcal{G}]$ 存在. 令 $Y=Y^+-Y^-$, 它是 $\mathcal{G}$-可测且可积的, 并且对 $G\in\mathcal{G}$, 由期望的线性,
$\E[Y;G]=\E[X^+;G]-\E[X^-;G]=\E[X;G]$.
\end{proof}

\begin{intuition}
$\E[X\,|\,\mathcal{G}]$ 是 $X$ 的一个 ``粗化'' 版本: 它只能利用 $\mathcal{G}$ 中的信息, 却在 $\mathcal{G}$ 能分辨的每个事件上保持与 $X$ 相同的总量. 粗化的程度由 $\mathcal{G}$ 决定: $\mathcal{G}$ 越小, 能分辨的事件越少, 条件期望越 ``平滑''. 两个极端是 $\E[X\,|\,\{\varnothing,\Omega\}]=\E[X]$ (毫无信息, 只能报告常数) 与 $\E[X\,|\,\mathcal{F}]=X$ (完全信息, 无需估计). 这两个等式可以直接由定义验证.
\end{intuition}

\begin{remark}[Radon-Nikodym 定理的观点]
条件期望的存在性也可以由 Radon-Nikodym 定理得到: 对 $X\ge0$, $\nu(G)=\E[X;G]$ 是 $(\Omega,\mathcal{G})$ 上的有限测度, 且关于 $\P$ 在 $\mathcal{G}$ 上的限制绝对连续 ($\P(G)=0\Rightarrow\nu(G)=0$). 定义~\ref{def:ce} (2) 恰好说明 $\E[X\,|\,\mathcal{G}]$ 是 Radon-Nikodym 导数 $\frac{\dd\nu}{\dd\P|_{\mathcal{G}}}$. 上面的证明利用 1.6 节的 Hilbert 空间结构, 避开了 Radon-Nikodym 定理; 第~5~章讨论测度变换时, 我们将从另一个方向用到 Radon-Nikodym 导数.
\end{remark}

由定义~\ref{def:ce-partition} 与命题~\ref{prop:ce-partition-properties} (1), 在可数分割的情形, 一般定义与初等定义一致. 若某块 $D_i$ 满足 $\P(D_i)=0$, 条件期望在该块上的取值还可以任意指定, 这正体现了 ``版本'' 的不唯一性.

% ---------------------------------------------------------------------
\subsection{关于随机变量的条件期望}

最常见的情形是以另一个随机变量 (或随机向量) 的信息为条件.

\begin{definition}\label{def:ce-rv}
设 $X\in L^1$, $\rvec{Y}$ 是随机向量. 定义 $\E[X\,|\,\rvec{Y}]=\E[X\,|\,\sigma(\rvec{Y})]$. 对随机变量族, $\E[X\,|\,Y_1,\dots,Y_n]=\E[X\,|\,\sigma(Y_1,\dots,Y_n)]$.
\end{definition}

由于 $\E[X\,|\,\rvec{Y}]$ 是 $\sigma(\rvec{Y})$-可测的, 由 Doob-Dynkin 引理~\ref{lem:doob-dynkin} (及其后的注), 存在 Borel 函数 $g$ 使
\[
\E[X\,|\,\rvec{Y}]=g(\rvec{Y}).
\]
记 $\E[X\,|\,\rvec{Y}=\rvec{y}]=g(\rvec{y})$, 称为\textbf{条件期望函数} (或回归函数). 它由以下性质刻画: $g$ 是 Borel 函数, 且对一切 Borel 集 $B$,
\begin{equation}\label{eq:regression-function}
\E[X;\rvec{Y}\in B]=\E[g(\rvec{Y});\rvec{Y}\in B].
\end{equation}
$g$ 在 $\mu_{\rvec{Y}}$-几乎处处的意义下唯一: 若 $g_1,g_2$ 都满足 \eqref{eq:regression-function}, 则由唯一性 $g_1(\rvec{Y})=g_2(\rvec{Y})$ $\as$, 即 $\mu_{\rvec{Y}}(\{g_1\ne g_2\})=\P(g_1(\rvec{Y})\ne g_2(\rvec{Y}))=0$. 这就是 ``条件于零概率事件 $\{\rvec{Y}=\rvec{y}\}$'' 的严格含义: 单个点上的值没有意义, 但作为 $\rvec{y}$ 的函数, $g$ 几乎处处是确定的.

若 $(X,Y)$ 有联合密度, 条件期望函数有显式公式.

\begin{proposition}[条件密度]\label{prop:conditional-density}
设 $(X,Y)$ 有联合密度 $f(x,y)$, $Y$ 的边缘密度为 $f_Y(y)=\int_{\R}f(x,y)\dd x$. 定义\textbf{条件密度}
\[
f_{X|Y}(x\,|\,y)=\begin{cases}\dfrac{f(x,y)}{f_Y(y)}, & f_Y(y)>0,\\[2mm] 0, & f_Y(y)=0.\end{cases}
\]
若 $h$ 是 Borel 函数且 $h(X)\in L^1$, 则 $\E[h(X)\,|\,Y]=g(Y)$, 其中
\[
g(y)=\int_{\R}h(x)f_{X|Y}(x\,|\,y)\dd x\qquad(\text{当此积分绝对收敛时; 否则令 }g(y)=0).
\]
\end{proposition}

\begin{proof}
由 Tonelli 定理~\ref{thm:fubini}, $y\mapsto\int\abs{h(x)}f(x,y)\dd x$ 是 Borel 函数, 且
\[
\int_{\R}\Big(\int_{\R}\abs{h(x)}f(x,y)\dd x\Big)\dd y=\E\abs{h(X)}<\infty,
\]
所以对几乎每个 $y$ (关于 Lebesgue 测度), 内层积分有限; 由 Fubini 定理, $g$ 是 Borel 函数. 对 Borel 集 $B$, 由 Fubini 定理与 $f(x,y)=f_{X|Y}(x\,|\,y)f_Y(y)$ (当 $f_Y(y)=0$ 时, 对几乎每个 $x$ 有 $f(x,y)=0$, 两边都为 $0$),
\begin{align*}
\E[h(X);Y\in B]&=\int_B\Big(\int_{\R}h(x)f(x,y)\dd x\Big)\dd y=\int_B\Big(\int_{\R}h(x)f_{X|Y}(x\,|\,y)\dd x\Big)f_Y(y)\dd y\\
&=\int_Bg(y)f_Y(y)\dd y=\E[g(Y);Y\in B].
\end{align*}
$g(Y)$ 的可积性: $\E\abs{g(Y)}\le\int\int\abs{h(x)}f(x,y)\dd x\dd y=\E\abs{h(X)}<\infty$. 由 \eqref{eq:regression-function} 即得.
\end{proof}

\begin{example}[三角形上的均匀分布]\label{ex:triangle}
设 $(X,Y)$ 在三角形 $\{(x,y): 0<x<y<1\}$ 上均匀分布, 联合密度为 $f=2$. 则 $f_Y(y)=\int_0^y2\dd x=2y$ ($0<y<1$), 条件密度 $f_{X|Y}(x\,|\,y)=\frac{1}{y}\1_{(0,y)}(x)$: 已知 $Y=y$ 时, $X$ 在 $(0,y)$ 上均匀分布. 于是
\[
\E[X\,|\,Y]=\frac{Y}{2},\qquad \E[X^2\,|\,Y]=\frac{Y^2}{3}.
\]
反过来, $f_X(x)=2(1-x)$, $f_{Y|X}(y\,|\,x)=\frac{1}{1-x}\1_{(x,1)}(y)$, 所以 $\E[Y\,|\,X]=\frac{1+X}{2}$.
\end{example}

\begin{figure}[htbp]
\centering
\begin{tikzpicture}[scale=3.6]
  \fill[SABlue!12] (0,0) -- (1,1) -- (0,1) -- cycle;
  \draw[SABlue!70] (0,0) -- (1,1) -- (0,1) -- cycle;
  \draw[->,SADark!70] (0,0) -- (1.15,0) node[right,font=\small] {$x$};
  \draw[->,SADark!70] (0,0) -- (0,1.15) node[above,font=\small] {$y$};
  \draw[SARed,very thick] (0,0.7) -- (0.7,0.7);
  \fill[SARed] (0.35,0.7) circle (0.5pt);
  \node[SARed,font=\small,anchor=west] at (0.74,0.7) {$\E[X\,|\,Y=y]=\frac{y}{2}$};
  \draw[SARed,->,thin] (0.73,0.66) to[bend left=20] (0.37,0.68);
  \node[left,font=\small] at (0,0.7) {$y$};
  \draw[SAGreenText,thick,dashed] (0,0) -- (0.5,1);
  \node[SAGreenText,font=\small,anchor=west] at (0.52,0.98) {$x=\frac{y}{2}$};
  \node[font=\small] at (0.25,0.38) {$f=2$};
  \node[below,font=\small] at (1,0) {$1$};
\end{tikzpicture}
\caption{例~\ref{ex:triangle}: 已知 $Y=y$, $X$ 在水平线段 $(0,y)$ 上均匀分布, 条件期望是线段中点. 所有中点连成回归直线 $x=\frac{y}{2}$}
\label{fig:triangle}
\end{figure}

\begin{example}[信息的不对称]\label{ex:info-asymmetry}
设 $X\sim N(0,1)$. 则 $\E[X^2\,|\,X]=X^2$ (它是 $\sigma(X)$-可测的), 但 $\E[X\,|\,X^2]=0$. 事实上, 对 Borel 集 $B$, 函数 $x\mapsto x\1_B(x^2)$ 是奇函数, 而 $X\overset{d}{=}-X$, 所以 $\E[X;X^2\in B]=\E[-X;X^2\in B]$, 从而为 $0=\E[0;X^2\in B]$. 直观地, 知道 $X$ 就知道 $X^2$; 但知道 $X^2$ 只知道 $\abs{X}$, 而 $X$ 的两个可能取值 $\pm\abs{X}$ 等可能, 平均为 $0$.
\end{example}

% =====================================================================
\section{条件期望的性质}
\label{sec:2-3}

本节推导条件期望的运算法则. 证明几乎都遵循同一个模式: 要证 $Y=\E[X\,|\,\mathcal{G}]$, 就验证 $Y$ 是 $\mathcal{G}$-可测且可积的, 并且 $\E[Y;G]=\E[X;G]$ 对一切 $G\in\mathcal{G}$ (或一个生成 $\pi$-系中的 $G$) 成立, 再由唯一性得出结论. 以下设 $\mathcal{G},\mathcal{H}\subseteq\mathcal{F}$ 是子 $\sigma$-代数, 随机变量都可积, 等式与不等式都在几乎必然的意义下成立.

% ---------------------------------------------------------------------
\subsection{基本性质}

\begin{theorem}[条件期望的基本性质]\label{thm:ce-basic}
设 $X,Y\in L^1$, $a,b\in\R$. 则
\begin{enumerate}
\item 线性: $\E[aX+bY\,|\,\mathcal{G}]=a\E[X\,|\,\mathcal{G}]+b\E[Y\,|\,\mathcal{G}]$;
\item 单调性: 若 $X\le Y$, 则 $\E[X\,|\,\mathcal{G}]\le\E[Y\,|\,\mathcal{G}]$; 特别地, 若 $X\ge0$, 则 $\E[X\,|\,\mathcal{G}]\ge0$;
\item 三角不等式: $\abs{\E[X\,|\,\mathcal{G}]}\le\E[\abs{X}\,|\,\mathcal{G}]$, 从而 $\E\abs{\E[X\,|\,\mathcal{G}]}\le\E\abs{X}$;
\item 全期望公式: $\E\big[\E[X\,|\,\mathcal{G}]\big]=\E[X]$;
\item 若 $X$ 是 $\mathcal{G}$-可测的, 则 $\E[X\,|\,\mathcal{G}]=X$; 特别地, 常数的条件期望是它自身;
\item 独立性化简: 若 $X$ 与 $\mathcal{G}$ 独立, 则 $\E[X\,|\,\mathcal{G}]=\E[X]$.
\end{enumerate}
\end{theorem}

\begin{proof}
(1) 右端是 $\mathcal{G}$-可测且可积的. 对 $G\in\mathcal{G}$, 由期望的线性与定义,
\[
\E\big[a\E[X\,|\,\mathcal{G}]+b\E[Y\,|\,\mathcal{G}];G\big]=a\E[X;G]+b\E[Y;G]=\E[aX+bY;G].
\]

(2) 对 $G\in\mathcal{G}$, $\E\big[\E[X\,|\,\mathcal{G}];G\big]=\E[X;G]\le\E[Y;G]=\E\big[\E[Y\,|\,\mathcal{G}];G\big]$, 由比较引理~\ref{lem:ce-compare} 即得. 取 $X=0$ 得第二个论断.

(3) 由 $-\abs{X}\le X\le\abs{X}$ 与 (1)(2), $-\E[\abs{X}\,|\,\mathcal{G}]\le\E[X\,|\,\mathcal{G}]\le\E[\abs{X}\,|\,\mathcal{G}]$. 两边取期望并用 (4) 得第二式.

(4) 在定义~\ref{def:ce} (2) 中取 $G=\Omega$.

(5) $X$ 自身满足定义~\ref{def:ce} 的两个条件.

(6) 常数 $\E[X]$ 是 $\mathcal{G}$-可测的. 对 $G\in\mathcal{G}$, $X$ 与 $\1_G$ 独立, 由定理~\ref{thm:indep-product}, $\E[X;G]=\E[X]\P(G)=\E\big[\E[X];G\big]$.
\end{proof}

\begin{intuition}
(4) 是 ``先分情况求平均, 再对各情况加权平均, 等于直接求平均'', 即初等概率中的全概率公式. (6) 说: 与 $X$ 无关的信息对预测 $X$ 毫无帮助, 最佳预测仍是无条件期望. (3) 的第二式说明条件期望是 $L^1$ 上的压缩映射: 若 $X_n\xrightarrow{L^1}X$, 则 $\E[X_n\,|\,\mathcal{G}]\xrightarrow{L^1}\E[X\,|\,\mathcal{G}]$.
\end{intuition}

% ---------------------------------------------------------------------
\subsection{塔性质与 ``提出已知因子''}

下面两个性质刻画了条件期望与信息结构的关系, 在鞅论与随机积分中几乎每一步都要用到.

\begin{theorem}[塔性质]\label{thm:tower}
若 $\mathcal{H}\subseteq\mathcal{G}$, 则
\[
\E\big[\E[X\,|\,\mathcal{G}]\,\big|\,\mathcal{H}\big]=\E[X\,|\,\mathcal{H}]=\E\big[\E[X\,|\,\mathcal{H}]\,\big|\,\mathcal{G}\big].
\]
\end{theorem}

\begin{proof}
第二个等号: $\E[X\,|\,\mathcal{H}]$ 是 $\mathcal{H}$-可测的, 从而是 $\mathcal{G}$-可测的, 由定理~\ref{thm:ce-basic} (5) 即得.

第一个等号: $\E[X\,|\,\mathcal{H}]$ 是 $\mathcal{H}$-可测且可积的. 对 $H\in\mathcal{H}$, 由于 $H\in\mathcal{G}$, 两次使用定义~\ref{def:ce} (2),
\[
\E\big[\E[X\,|\,\mathcal{G}];H\big]=\E[X;H]=\E\big[\E[X\,|\,\mathcal{H}];H\big].
\]
所以 $\E[X\,|\,\mathcal{H}]$ 满足 $\E[X\,|\,\mathcal{G}]$ 关于 $\mathcal{H}$ 的条件期望的定义.
\end{proof}

\begin{figure}[htbp]
\centering
\begin{tikzpicture}
\begin{axis}[SAaxis, width=11cm, height=5.8cm, xmin=0, xmax=1.06, ymin=0, ymax=2.3,
  xtick={0.25,0.5,0.75,1}, ytick=\empty, xlabel={$\omega$},
  legend style={draw=none,fill=none,font=\small,at={(0.5,-0.2)},anchor=north,legend columns=3,/tikz/every even column/.append style={column sep=1em}}]
  \addplot[SABlue!60,thick,domain=0:1,samples=200] {1.0+0.6*sin(deg(9*x))+0.5*x};
  \addlegendentry{$X$}
  \addplot[SAGreenText,very thick] coordinates {(0,1.4967) (0.25,1.4967)};
  \addlegendentry{$\E[X\,|\,\mathcal{G}]$}
  \addplot[SARed,very thick,dashed] coordinates {(0,1.2864) (0.5,1.2864)};
  \addlegendentry{$\E[X\,|\,\mathcal{H}]$}
  \addplot[SAGreenText,very thick,forget plot] coordinates {(0.25,1.0762) (0.5,1.0762)};
  \addplot[SAGreenText,very thick,forget plot] coordinates {(0.5,1.0182) (0.75,1.0182)};
  \addplot[SAGreenText,very thick,forget plot] coordinates {(0.75,1.9186) (1,1.9186)};
  \addplot[SARed,very thick,dashed,forget plot] coordinates {(0.5,1.4684) (1,1.4684)};
  \addplot[SADark!25,forget plot] coordinates {(0.25,0) (0.25,2.2)};
  \addplot[SADark!50,thick,forget plot] coordinates {(0.5,0) (0.5,2.2)};
  \addplot[SADark!25,forget plot] coordinates {(0.75,0) (0.75,2.2)};
\end{axis}
\end{tikzpicture}
\caption{塔性质 ($\Omega=[0,1]$, $\mathcal{G}$ 由四个等长区间生成, $\mathcal{H}\subseteq\mathcal{G}$ 由两个等长区间生成): 先在细分割上平均, 再把相邻两块的平均值再平均, 结果等于直接在粗分割上平均}
\label{fig:tower}
\end{figure}

\begin{intuition}
塔性质可以概括为 ``信息少的一方说了算'': 先用较多的信息 $\mathcal{G}$ 预测, 再用较少的信息 $\mathcal{H}$ 预测这个预测, 等于直接用 $\mathcal{H}$ 预测; 而用较少信息得到的预测, 在较多信息下已经完全已知, 不会再改变. 在 $L^2$ 中, 这就是投影的塔性质 (命题~\ref{prop:projection-properties} (5)).
\end{intuition}

\begin{theorem}[提出已知因子]\label{thm:take-out}
设 $Z$ 是 $\mathcal{G}$-可测的随机变量, 且 $X$ 与 $ZX$ 都可积. 则
\[
\E[ZX\,|\,\mathcal{G}]=Z\,\E[X\,|\,\mathcal{G}].
\]
\end{theorem}

\begin{proof}
按 $Z$ 的复杂程度分步.

\emph{第一步: $Z=\1_{G_0}$, $G_0\in\mathcal{G}$.} 右端 $\1_{G_0}\E[X\,|\,\mathcal{G}]$ 是 $\mathcal{G}$-可测且可积的. 对 $G\in\mathcal{G}$, 由于 $G\cap G_0\in\mathcal{G}$,
\[
\E\big[\1_{G_0}\E[X\,|\,\mathcal{G}];G\big]=\E\big[\E[X\,|\,\mathcal{G}];G\cap G_0\big]=\E[X;G\cap G_0]=\E[\1_{G_0}X;G].
\]

\emph{第二步: $Z$ 是 $\mathcal{G}$-可测的简单随机变量.} 由第一步与条件期望的线性 (定理~\ref{thm:ce-basic} (1)).

\emph{第三步: $Z\ge0$, $X\ge0$.} 由引理~\ref{lem:simple-approx}, 取 $\mathcal{G}$-可测的非负简单随机变量 $Z_n\uparrow Z$. 由定理~\ref{thm:ce-basic} (2), $\E[X\,|\,\mathcal{G}]\ge0$, 所以 $Z_n\E[X\,|\,\mathcal{G}]\uparrow Z\E[X\,|\,\mathcal{G}]$, 同时 $Z_nX\uparrow ZX$. 对 $G\in\mathcal{G}$, 由单调收敛定理~\ref{thm:mct} 与第二步,
\[
\E\big[Z\E[X\,|\,\mathcal{G}];G\big]=\lim_{n\to\infty}\E\big[Z_n\E[X\,|\,\mathcal{G}];G\big]=\lim_{n\to\infty}\E[Z_nX;G]=\E[ZX;G].
\]
取 $G=\Omega$ 知 $\E\big[Z\E[X\,|\,\mathcal{G}]\big]=\E[ZX]<\infty$, 故 $Z\E[X\,|\,\mathcal{G}]$ 可积, 它满足 $\E[ZX\,|\,\mathcal{G}]$ 的定义.

\emph{第四步: 一般情形.} 记 $Z=Z^+-Z^-$, $X=X^+-X^-$. 对 $a,b\in\{+,-\}$, $Z^a$ 是 $\mathcal{G}$-可测的非负随机变量, $Z^aX^b$ 非负且 $\abs{Z^aX^b}\le\abs{ZX}$, 所以可积. 由第三步, $\E[Z^aX^b\,|\,\mathcal{G}]=Z^a\E[X^b\,|\,\mathcal{G}]$. 展开 $ZX=Z^+X^+-Z^+X^--Z^-X^++Z^-X^-$, 由线性,
\[
\E[ZX\,|\,\mathcal{G}]=\big(Z^+-Z^-\big)\big(\E[X^+\,|\,\mathcal{G}]-\E[X^-\,|\,\mathcal{G}]\big)=Z\,\E[X\,|\,\mathcal{G}].\qedhere
\]
\end{proof}

\begin{intuition}
在信息 $\mathcal{G}$ 下, $Z$ 的值已经知道, 它像常数一样可以提到条件期望外面. 例如在第~3~章中, 时刻 $n-1$ 决定的赌注 $H_n$ 乘以下一局的输赢 $\xi_n$, 其条件期望为 $\E[H_n\xi_n\,|\,\mathcal{F}_{n-1}]=H_n\E[\xi_n\,|\,\mathcal{F}_{n-1}]$; 若每一局都是公平的, 这一项为零, 这就是 ``公平赌博中任何策略都无法获利'' 的根源.
\end{intuition}

\begin{corollary}[条件期望的自伴性]\label{cor:ce-selfadjoint}
若 $X,Z\in L^2$, 或 $X\in L^1$ 且 $Z$ 有界, 则
\[
\E\big[Z\,\E[X\,|\,\mathcal{G}]\big]=\E\big[\E[Z\,|\,\mathcal{G}]\,\E[X\,|\,\mathcal{G}]\big]=\E\big[\E[Z\,|\,\mathcal{G}]\,X\big].
\]
\end{corollary}

\begin{proof}
所涉及的乘积都可积: 在 $L^2$ 情形, 由 Cauchy-Schwarz 不等式及 $\norm{\E[\cdot\,|\,\mathcal{G}]}_2\le\norm{\cdot}_2$ (定理~\ref{thm:ce-existence} 与命题~\ref{prop:projection-properties} (3)); 在 $Z$ 有界情形, 由 $\abs{\E[Z\,|\,\mathcal{G}]}\le\sup\abs{Z}$ (定理~\ref{thm:ce-basic} (2)). 由全期望公式与提出已知因子 (提出 $\mathcal{G}$-可测的 $\E[X\,|\,\mathcal{G}]$),
\[
\E\big[Z\,\E[X\,|\,\mathcal{G}]\big]=\E\Big[\E\big[Z\,\E[X\,|\,\mathcal{G}]\,\big|\,\mathcal{G}\big]\Big]=\E\big[\E[Z\,|\,\mathcal{G}]\,\E[X\,|\,\mathcal{G}]\big].
\]
交换 $X$ 与 $Z$ 的角色即得第二个等号.
\end{proof}

% ---------------------------------------------------------------------
\subsection{条件极限定理与条件 Jensen 不等式}

1.3 节的三个积分收敛定理都有条件版本. 与无条件情形不同, 结论是随机变量之间的几乎必然收敛.

\begin{theorem}[条件单调收敛定理]\label{thm:cmct}
若 $0\le X_n\uparrow X$, 且 $X\in L^1$, 则 $\E[X_n\,|\,\mathcal{G}]\uparrow\E[X\,|\,\mathcal{G}]$.
\end{theorem}

\begin{proof}
由于 $0\le X_n\le X$, 每个 $X_n$ 可积. 记 $Y_n=\E[X_n\,|\,\mathcal{G}]$. 由单调性 (定理~\ref{thm:ce-basic} (2)), $0\le Y_n\le Y_{n+1}\le\E[X\,|\,\mathcal{G}]$ (除去一个 $\mathcal{G}$-可测零概率集, 与定理~\ref{thm:ce-existence} 第三步相同处理). 令 $Y=\lim_nY_n$, 它是 $\mathcal{G}$-可测的, 且 $Y\le\E[X\,|\,\mathcal{G}]$, 所以可积. 对 $G\in\mathcal{G}$, 两次使用单调收敛定理~\ref{thm:mct},
\[
\E[Y;G]=\lim_{n\to\infty}\E[Y_n;G]=\lim_{n\to\infty}\E[X_n;G]=\E[X;G].
\]
所以 $Y=\E[X\,|\,\mathcal{G}]$.
\end{proof}

\begin{theorem}[条件 Fatou 引理]\label{thm:cfatou}
若 $X_n\ge0$ 且 $X_n$ 与 $\liminf_nX_n$ 都可积, 则
\[
\E\Big[\liminf_{n\to\infty}X_n\,\Big|\,\mathcal{G}\Big]\le\liminf_{n\to\infty}\E[X_n\,|\,\mathcal{G}].
\]
\end{theorem}

\begin{proof}
令 $Z_m=\inf_{n\ge m}X_n$, 则 $0\le Z_m\uparrow\liminf_nX_n$, 且 $Z_m\le\liminf_nX_n$, 故 $Z_m$ 可积. 对 $n\ge m$, $Z_m\le X_n$, 由单调性 $\E[Z_m\,|\,\mathcal{G}]\le\E[X_n\,|\,\mathcal{G}]$; 对可列个 $n$ 取下确界 (例外零概率集的可列并仍为零概率集), $\E[Z_m\,|\,\mathcal{G}]\le\inf_{n\ge m}\E[X_n\,|\,\mathcal{G}]$. 令 $m\to\infty$, 由条件单调收敛定理~\ref{thm:cmct},
\[
\E\Big[\liminf_{n\to\infty}X_n\,\Big|\,\mathcal{G}\Big]=\lim_{m\to\infty}\E[Z_m\,|\,\mathcal{G}]\le\lim_{m\to\infty}\inf_{n\ge m}\E[X_n\,|\,\mathcal{G}]=\liminf_{n\to\infty}\E[X_n\,|\,\mathcal{G}].\qedhere
\]
\end{proof}

\begin{theorem}[条件控制收敛定理]\label{thm:cdct}
若 $X_n\to X$ $\as$, 且存在 $Y\in L^1$ 使 $\abs{X_n}\le Y$, 则 $\E[X_n\,|\,\mathcal{G}]\to\E[X\,|\,\mathcal{G}]$ $\as$
\end{theorem}

\begin{proof}
令 $W_n=\sup_{k\ge n}\abs{X_k-X}$. 则 $0\le W_n\le2Y$, 且 $W_n\downarrow0$ $\as$ 由条件期望的线性、三角不等式与单调性 (定理~\ref{thm:ce-basic}),
\[
\abs{\E[X_n\,|\,\mathcal{G}]-\E[X\,|\,\mathcal{G}]}\le\E[\abs{X_n-X}\,|\,\mathcal{G}]\le\E[W_n\,|\,\mathcal{G}].
\]
由单调性, $\E[W_n\,|\,\mathcal{G}]$ 关于 $n$ 单调不增且非负, 设其极限为 $W\ge0$. 由全期望公式与 $W\le\E[W_n\,|\,\mathcal{G}]$, $\E[W]\le\E[W_n]$ 对一切 $n$ 成立, 而由控制收敛定理~\ref{thm:dct} (控制函数 $2Y$), $\E[W_n]\to0$. 所以 $\E[W]=0$, 由推论~\ref{cor:nonneg-properties} (3), $W=0$ $\as$, 即得结论.
\end{proof}

为了证明条件 Jensen 不等式, 需要把凸函数写成\emph{可列个}仿射函数的上确界: 每个仿射函数对应一个例外零概率集, 只有可列个零概率集之并才仍是零概率集. 回忆实分析中的事实: 若 $\phi:\R\to\R$ 是凸函数, 则对 $s<t<u$ 有三弦不等式
\[
\frac{\phi(t)-\phi(s)}{t-s}\le\frac{\phi(u)-\phi(s)}{u-s}\le\frac{\phi(u)-\phi(t)}{u-t};
\]
因此右导数 $\phi'_+(x)$ 处处存在且单调不减, $\phi$ 连续, 并且 $\phi(y)\ge\phi(x)+\phi'_+(x)(y-x)$ 对一切 $x,y$ 成立 (这就是 \eqref{eq:support-line}).

\begin{lemma}\label{lem:convex-countable}
设 $\phi:\R\to\R$ 是凸函数. 对 $q\in\Q$ 令 $\ell_q(x)=\phi(q)+\phi'_+(q)(x-q)$. 则 $\phi(x)=\sup_{q\in\Q}\ell_q(x)$ 对一切 $x\in\R$ 成立.
\end{lemma}

\begin{proof}
由支撑直线不等式, $\ell_q\le\phi$, 故 $\sup_q\ell_q\le\phi$. 反过来, 固定 $x$, 取有理数 $q_k\to x$ 且 $q_k\in[x-1,x+1]$. 由 $\phi'_+$ 单调不减, $\abs{\phi'_+(q_k)}\le M=\max\{\abs{\phi'_+(x-1)},\abs{\phi'_+(x+1)}\}$. 于是
\[
\sup_{q\in\Q}\ell_q(x)\ge\ell_{q_k}(x)\ge\phi(q_k)-M\abs{x-q_k}\xrightarrow{k\to\infty}\phi(x),
\]
最后一步用到 $\phi$ 的连续性.
\end{proof}

\begin{theorem}[条件 Jensen 不等式]\label{thm:cjensen}
设 $\phi:\R\to\R$ 是凸函数, $X$ 与 $\phi(X)$ 都可积. 则
\[
\phi\big(\E[X\,|\,\mathcal{G}]\big)\le\E[\phi(X)\,|\,\mathcal{G}].
\]
\end{theorem}

\begin{proof}
对每个 $q\in\Q$, 由 $\ell_q(X)\le\phi(X)$ 及条件期望的线性与单调性,
\[
\ell_q\big(\E[X\,|\,\mathcal{G}]\big)=\phi(q)+\phi'_+(q)\big(\E[X\,|\,\mathcal{G}]-q\big)=\E[\ell_q(X)\,|\,\mathcal{G}]\le\E[\phi(X)\,|\,\mathcal{G}]
\]
在某个零概率集 $N_q$ 之外成立. 在 $\big(\bigcup_{q\in\Q}N_q\big)^c$ 上对 $q$ 取上确界, 由引理~\ref{lem:convex-countable} 即得.
\end{proof}

\begin{corollary}[$L^p$ 压缩性]\label{cor:ce-Lp}
设 $p\ge1$, $X\in L^p$. 则 $\abs{\E[X\,|\,\mathcal{G}]}^p\le\E[\abs{X}^p\,|\,\mathcal{G}]$, 从而 $\norm{\E[X\,|\,\mathcal{G}]}_p\le\norm{X}_p$. 特别地, $\big(\E[X\,|\,\mathcal{G}]\big)^2\le\E[X^2\,|\,\mathcal{G}]$.
\end{corollary}

\begin{proof}
$\phi(x)=\abs{x}^p$ 是凸函数, 由条件 Jensen 不等式~\ref{thm:cjensen} 得第一式, 取期望并用全期望公式得第二式.
\end{proof}

% ---------------------------------------------------------------------
\subsection{独立性与代入引理}

定理~\ref{thm:ce-basic} (6) 处理了 ``$X$ 与 $\mathcal{G}$ 独立'' 的情形. 实际中更常见的是混合情形, 例如随机游动 $S_{n+1}=S_n+\xi_{n+1}$: $S_n$ 在时刻 $n$ 已知, 而 $\xi_{n+1}$ 与过去独立. 计算 $\E[f(S_n+\xi_{n+1})\,|\,\mathcal{F}_n]$ 的直观做法是把已知的 $S_n$ ``冻结'' 为常数 $x$, 对独立的 $\xi_{n+1}$ 求期望, 再把 $x$ 换回 $S_n$. 本小节证明这一做法是正确的. 证明需要一个函数形式的单调类定理.

\begin{theorem}[函数形式的单调类定理]\label{thm:functional-mct}
设 $\mathcal{H}$ 是 $\Omega$ 上一族有界实函数, 满足
\begin{enumerate}
\item $\mathcal{H}$ 是线性空间, 且常数函数 $1\in\mathcal{H}$;
\item 若 $f_n\in\mathcal{H}$, $0\le f_n\uparrow f$, 且 $f$ 有界, 则 $f\in\mathcal{H}$.
\end{enumerate}
若 $\mathcal{P}$ 是 $\pi$-系, 且 $\1_A\in\mathcal{H}$ 对一切 $A\in\mathcal{P}$ 成立, 则 $\mathcal{H}$ 包含全体有界 $\sigma(\mathcal{P})$-可测函数.
\end{theorem}

\begin{proof}
令 $\mathcal{L}=\{A\subseteq\Omega: \1_A\in\mathcal{H}\}$. 由假设 $\mathcal{P}\subseteq\mathcal{L}$. 验证 $\mathcal{L}$ 是 $\lambda$-系: (a) $\1_\Omega=1\in\mathcal{H}$; (b) 若 $A,B\in\mathcal{L}$, $A\subseteq B$, 则 $\1_{B\setminus A}=\1_B-\1_A\in\mathcal{H}$ (线性空间); (c) 若 $A_n\in\mathcal{L}$, $A_n\uparrow A$, 则 $0\le\1_{A_n}\uparrow\1_A$, 由条件 (2), $\1_A\in\mathcal{H}$. 由 $\pi$-$\lambda$ 定理~\ref{thm:pi-lambda}, $\sigma(\mathcal{P})\subseteq\mathcal{L}$.

于是 $\sigma(\mathcal{P})$-可测集的示性函数都属于 $\mathcal{H}$, 由线性性, $\sigma(\mathcal{P})$-可测的简单函数都属于 $\mathcal{H}$. 若 $f\ge0$ 是有界 $\sigma(\mathcal{P})$-可测函数, 由引理~\ref{lem:simple-approx}, 存在 $\sigma(\mathcal{P})$-可测的非负简单函数 $f_n\uparrow f$, 由条件 (2), $f\in\mathcal{H}$. 一般的有界可测函数 $f=f^+-f^-$ 由线性性也属于 $\mathcal{H}$.
\end{proof}

\begin{intuition}
这是 ``好函数原理'': 要证明某个关于函数的线性性质对全体有界可测函数成立, 只需对一个生成 $\pi$-系中集合的示性函数验证, 再验证该性质在单调极限下保持. 它与 $\pi$-$\lambda$ 定理的关系, 就像积分的四步构造与示性函数的关系.
\end{intuition}

\begin{lemma}[代入引理]\label{lem:freezing}
设 $\rvec{X}$ 是取值于 $\R^k$ 的随机向量, 与 $\mathcal{G}$ 独立; $\rvec{Y}$ 是取值于 $\R^m$ 的 $\mathcal{G}$-可测随机向量; $\phi:\R^k\times\R^m\to\R$ 是 Borel 函数.
\begin{enumerate}
\item 若 $\phi$ 有界, 令 $\psi(\rvec{y})=\E[\phi(\rvec{X},\rvec{y})]$, 则 $\psi$ 是有界 Borel 函数, 且
\[
\E[\phi(\rvec{X},\rvec{Y})\,|\,\mathcal{G}]=\psi(\rvec{Y}).
\]
\item 若 $\phi(\rvec{X},\rvec{Y})\in L^1$, 则对 $\mu_{\rvec{Y}}$-几乎每个 $\rvec{y}$, $\phi(\rvec{X},\rvec{y})\in L^1$; 令 $\psi(\rvec{y})=\E[\phi(\rvec{X},\rvec{y})]$ (在其余的 $\rvec{y}$ 处令 $\psi=0$), 则上式仍成立.
\end{enumerate}
\end{lemma}

\begin{proof}
(1) 令 $\mathcal{H}$ 为满足以下条件的有界 Borel 函数 $\phi$ 全体: $\psi_\phi(\rvec{y})=\E[\phi(\rvec{X},\rvec{y})]$ 是 Borel 函数, 且对一切 $G\in\mathcal{G}$,
\begin{equation}\label{eq:freezing-H}
\E[\phi(\rvec{X},\rvec{Y});G]=\E[\psi_\phi(\rvec{Y});G].
\end{equation}
我们用函数形式的单调类定理~\ref{thm:functional-mct} 证明 $\mathcal{H}$ 包含全体有界 Borel 函数.
\begin{itemize}
\item \emph{矩形的示性函数.} 设 $\phi=\1_{A\times B}$, $A\in\mathcal{B}(\R^k)$, $B\in\mathcal{B}(\R^m)$. 则 $\phi(\rvec{x},\rvec{y})=\1_A(\rvec{x})\1_B(\rvec{y})$, $\psi_\phi(\rvec{y})=\P(\rvec{X}\in A)\1_B(\rvec{y})$ 是 Borel 函数. 由于 $\1_A(\rvec{X})$ 是 $\sigma(\rvec{X})$-可测的, $\1_B(\rvec{Y})\1_G$ 是 $\mathcal{G}$-可测的, 而 $\sigma(\rvec{X})$ 与 $\mathcal{G}$ 独立, 由定理~\ref{thm:indep-product},
\[
\E[\1_A(\rvec{X})\1_B(\rvec{Y})\1_G]=\P(\rvec{X}\in A)\,\E[\1_B(\rvec{Y})\1_G]=\E[\psi_\phi(\rvec{Y});G].
\]
矩形 $A\times B$ 全体构成 $\pi$-系, 且生成 $\mathcal{B}(\R^{k+m})$ (定理~\ref{thm:indep-product-measure} 的证明).
\item \emph{线性空间且含常数.} $\phi\mapsto\psi_\phi$ 是线性的, \eqref{eq:freezing-H} 两边关于 $\phi$ 线性; 常数 $1$ 是 $\Omega$ 的示性函数, 已包含在上一条中.
\item \emph{单调极限.} 若 $\phi_n\in\mathcal{H}$, $0\le\phi_n\uparrow\phi$, $\phi$ 有界, 则对每个 $\rvec{y}$, 由单调收敛定理 $\psi_{\phi_n}(\rvec{y})\uparrow\psi_\phi(\rvec{y})$, 所以 $\psi_\phi$ 是 Borel 函数; 再对 \eqref{eq:freezing-H} 两边用单调收敛定理, 得到 $\phi\in\mathcal{H}$.
\end{itemize}
由定理~\ref{thm:functional-mct}, $\mathcal{H}$ 包含全体有界 Borel 函数. 对有界 Borel 函数 $\phi$, $\psi_\phi(\rvec{Y})$ 是 $\mathcal{G}$-可测的有界随机变量, 而 \eqref{eq:freezing-H} 正是条件期望的定义.

(2) 先设 $\phi\ge0$. 对 $\phi\wedge n$ 用 (1): $\E[\phi(\rvec{X},\rvec{Y})\wedge n\,|\,\mathcal{G}]=\psi_n(\rvec{Y})$, 其中 $\psi_n(\rvec{y})=\E[\phi(\rvec{X},\rvec{y})\wedge n]$. 由单调收敛定理, $\psi_n(\rvec{y})\uparrow\psi(\rvec{y})=\E[\phi(\rvec{X},\rvec{y})]\in[0,\infty]$, $\psi$ 是 Borel 函数. 由条件单调收敛定理~\ref{thm:cmct}, $\E[\phi(\rvec{X},\rvec{Y})\,|\,\mathcal{G}]=\lim_n\psi_n(\rvec{Y})=\psi(\rvec{Y})$. 左端可积, 所以 $\psi(\rvec{Y})<\infty$ $\as$, 即 $\mu_{\rvec{Y}}(\{\psi=\infty\})=0$. 对一般的 $\phi$, 分别对 $\phi^+$ 与 $\phi^-$ 使用上述结论再相减.
\end{proof}

\begin{example}\label{ex:freezing}
设 $\xi$ 与 $\mathcal{G}$ 独立, $Y$ 是 $\mathcal{G}$-可测的.
\begin{enumerate}
\item 若 $f$ 有界 Borel, 则 $\E[f(Y+\xi)\,|\,\mathcal{G}]=\psi(Y)$, 其中 $\psi(y)=\E[f(y+\xi)]=\int f(y+z)\,\mu_\xi(\dd z)$. 特别地, 若 $\xi\sim N(0,\sigma^2)$, 则
\[
\E[f(Y+\xi)\,|\,\mathcal{G}]=\int_{\R}f(Y+z)\frac{1}{\sqrt{2\pi}\sigma}e^{-z^2/(2\sigma^2)}\dd z.
\]
第~4~章中取 $Y=B_s$, $\xi=B_t-B_s$, 这就是布朗运动马尔可夫性的计算公式.
\item 若 $\E e^{\lambda\xi}<\infty$, 则 $\E[e^{\lambda(Y+\xi)}\,|\,\mathcal{G}]=e^{\lambda Y}\E[e^{\lambda\xi}]$ (当左端可积时). 这也可以由提出已知因子与独立性化简直接得到.
\item \textbf{独立性的条件不能减弱.} 引理要求 $\xi$ 与整个 $\mathcal{G}$ 独立, 仅与 $Y$ 独立是不够的. 例如设 $\xi,Y$ 是独立的 $N(0,1)$ 随机变量, $\mathcal{G}=\sigma(Y,\xi)$, 则 $\E[(Y+\xi)^2\,|\,\mathcal{G}]=(Y+\xi)^2$, 而 ``代入'' 会给出 $\E[(y+\xi)^2]|_{y=Y}=Y^2+1$.
\end{enumerate}
\end{example}

下面的推论在验证马尔可夫性时很有用: 与所关心的量独立的额外信息可以去掉.

\begin{proposition}[去掉独立的信息]\label{prop:drop-independent}
设 $X\in L^1$. 若 $\sigma\big(\sigma(X)\cup\mathcal{H}\big)$ 与 $\mathcal{G}$ 独立, 则
\[
\E\big[X\,\big|\,\sigma(\mathcal{G}\cup\mathcal{H})\big]=\E[X\,|\,\mathcal{H}].
\]
\end{proposition}

\begin{proof}
$\E[X\,|\,\mathcal{H}]$ 是 $\sigma(\mathcal{G}\cup\mathcal{H})$-可测且可积的. 集族 $\mathcal{P}=\{G\cap H: G\in\mathcal{G},H\in\mathcal{H}\}$ 是包含 $\Omega$ 的 $\pi$-系, 且生成 $\sigma(\mathcal{G}\cup\mathcal{H})$. 由引理~\ref{lem:ce-pi-system}, 只需对 $G\cap H\in\mathcal{P}$ 验证积分恒等式. 由于 $X\1_H$ 与 $\E[X\,|\,\mathcal{H}]\1_H$ 都是 $\sigma(\sigma(X)\cup\mathcal{H})$-可测的, 它们与 $\1_G$ 独立, 由定理~\ref{thm:indep-product},
\[
\E[X;G\cap H]=\E[X\1_H]\P(G)=\E\big[\E[X\,|\,\mathcal{H}]\1_H\big]\P(G)=\E\big[\E[X\,|\,\mathcal{H}];G\cap H\big],
\]
其中第二个等号是条件期望的定义 ($H\in\mathcal{H}$).
\end{proof}

% ---------------------------------------------------------------------
\subsection{条件方差}

\begin{definition}[条件方差]\label{def:conditional-variance}
设 $X\in L^2$. 定义 $X$ 关于 $\mathcal{G}$ 的\textbf{条件方差}
\[
\Var(X\,|\,\mathcal{G})=\E\big[(X-\E[X\,|\,\mathcal{G}])^2\,\big|\,\mathcal{G}\big]=\E[X^2\,|\,\mathcal{G}]-\big(\E[X\,|\,\mathcal{G}]\big)^2.
\]
\end{definition}

第二个等号: 展开平方, 并对交叉项 $\E\big[X\E[X\,|\,\mathcal{G}]\,\big|\,\mathcal{G}\big]$ 用提出已知因子 (定理~\ref{thm:take-out}) 得 $(\E[X\,|\,\mathcal{G}])^2$. 由条件 Jensen 不等式 (推论~\ref{cor:ce-Lp}), $\Var(X\,|\,\mathcal{G})\ge0$.

\begin{theorem}[全方差公式]\label{thm:total-variance}
设 $X\in L^2$. 则
\[
\Var(X)=\E\big[\Var(X\,|\,\mathcal{G})\big]+\Var\big(\E[X\,|\,\mathcal{G}]\big).
\]
\end{theorem}

\begin{proof}
记 $Y=\E[X\,|\,\mathcal{G}]$, 则 $\E Y=\E X$. 写 $X-\E X=(X-Y)+(Y-\E Y)$. 由于 $Y-\E Y\in L^2(\mathcal{G})$, 而 $X-Y$ 与 $L^2(\mathcal{G})$ 正交 (定义~\ref{def:ce-L2}), 由勾股定理,
\[
\Var(X)=\E[(X-Y)^2]+\E[(Y-\E Y)^2]=\E\Big[\E\big[(X-Y)^2\,\big|\,\mathcal{G}\big]\Big]+\Var(Y),
\]
最后一步用到全期望公式. 由定义, $\E[(X-Y)^2\,|\,\mathcal{G}]=\Var(X\,|\,\mathcal{G})$.
\end{proof}

\begin{intuition}
总的不确定性 = ``已知信息之后仍剩下的不确定性的平均'' + ``信息本身带来的变化''. 在 $L^2$ 的几何中, 这就是图~\ref{fig:ce-projection} 中直角三角形的勾股定理.
\end{intuition}

\begin{example}[随机和]\label{ex:random-sum}
设 $\xi_1,\xi_2,\dots$ 是 $\iid$ 随机变量, $\E\xi_1=\mu$, $\Var(\xi_1)=\sigma^2$; $N$ 取非负整数值, $N\in L^2$, 且 $N$ 与 $\{\xi_k\}$ 独立. 令 $S=\sum_{k=1}^N\xi_k$ (约定 $N=0$ 时 $S=0$), 例如 $N$ 是一天中的理赔次数, $\xi_k$ 是每次理赔的金额. 记 $S_n=\sum_{k=1}^n\xi_k$, 则在 $\{N=n\}$ 上 $S=S_n$.

$\sigma(N)$ 由可数分割 $\{N=n\}$ 生成. 对 $\P(N=n)>0$, 由于 $S_n$ 与 $\{N=n\}$ 独立,
\begin{align*}
\E[S\,|\,N=n]&=\frac{\E[S_n;N=n]}{\P(N=n)}=\E[S_n]=n\mu,\\
\E[S^2\,|\,N=n]&=\frac{\E[S_n^2;N=n]}{\P(N=n)}=\E[S_n^2]=n\sigma^2+n^2\mu^2.
\end{align*}
所以 $\E[S\,|\,N]=N\mu$, $\Var(S\,|\,N)=N\sigma^2+N^2\mu^2-N^2\mu^2=N\sigma^2$. 由全期望公式与全方差公式~\ref{thm:total-variance},
\[
\E[S]=\mu\,\E[N],\qquad \Var(S)=\E[N\sigma^2]+\Var(N\mu)=\sigma^2\E[N]+\mu^2\Var(N).
\]
方差分成两部分: 每次理赔金额的波动, 以及理赔次数本身的波动. 第~3~章研究分支过程时将用到同样的计算.
\end{example}

% ---------------------------------------------------------------------
\subsection{条件概率与正则条件分布}

\begin{definition}[条件概率]\label{def:conditional-probability}
对 $A\in\mathcal{F}$, 称 $\P(A\,|\,\mathcal{G})=\E[\1_A\,|\,\mathcal{G}]$ 为 $A$ 关于 $\mathcal{G}$ 的\textbf{条件概率}.
\end{definition}

由定理~\ref{thm:ce-basic} 与条件单调收敛定理~\ref{thm:cmct}, 条件概率具有概率的各条性质, 但每条都只在一个零概率集之外成立: $0\le\P(A\,|\,\mathcal{G})\le1$; $\P(\Omega\,|\,\mathcal{G})=1$; 对两两不交的 $A_1,A_2,\dots$, $\P\big(\bigcup_nA_n\,\big|\,\mathcal{G}\big)=\sum_n\P(A_n\,|\,\mathcal{G})$. 但最后一条的例外零概率集依赖于序列 $\{A_n\}$, 而这样的序列有不可数多个, 其例外集之并未必是零概率集. 所以对固定的 $\omega$, $A\mapsto\P(A\,|\,\mathcal{G})(\omega)$ 未必是概率测度. 为了能把条件概率当作测度来积分, 需要选取一个 ``好的版本''.

\begin{definition}[正则条件分布]\label{def:regular-conditional}
设 $\rvec{X}$ 是取值于 $\R^k$ 的随机向量. 称 $\mu:\Omega\times\mathcal{B}(\R^k)\to[0,1]$ 为 $\rvec{X}$ 关于 $\mathcal{G}$ 的\textbf{正则条件分布}, 如果
\begin{enumerate}
\item 对每个 $\omega$, $\mu(\omega,\cdot)$ 是 $(\R^k,\mathcal{B}(\R^k))$ 上的概率测度;
\item 对每个 $B\in\mathcal{B}(\R^k)$, $\mu(\cdot,B)$ 是 $\P(\rvec{X}\in B\,|\,\mathcal{G})$ 的一个版本 (特别地, 它是 $\mathcal{G}$-可测的).
\end{enumerate}
\end{definition}

\begin{theorem}[正则条件分布的存在性]\label{thm:regular-existence}
取值于 $\R^k$ 的随机向量关于任意子 $\sigma$-代数的正则条件分布总存在. (证明见~\cite{durrett} 定理 4.1.18; 其要点是利用 $\R^k$ 有可数稠密子集, 先在有理点上确定条件分布函数, 再扩张为测度.)
\end{theorem}

\begin{proposition}\label{prop:regular-integral}
设 $\mu$ 是 $\rvec{X}$ 关于 $\mathcal{G}$ 的正则条件分布, $h$ 是 Borel 函数且 $h(\rvec{X})\in L^1$. 则
\[
\E[h(\rvec{X})\,|\,\mathcal{G}](\omega)=\int_{\R^k}h(\rvec{x})\,\mu(\omega,\dd\rvec{x})\qquad\as
\]
\end{proposition}

\begin{proof}
按 $h$ 的复杂程度分步. 若 $h=\1_B$, 这就是定义~\ref{def:regular-conditional} (2). 若 $h$ 是非负简单函数, 由两边的线性得到. 若 $h\ge0$, 取非负简单函数 $h_n\uparrow h$: 对每个 $\omega$, 由单调收敛定理 $\int h_n\dd\mu(\omega,\cdot)\uparrow\int h\dd\mu(\omega,\cdot)$, 所以右端作为 $\mathcal{G}$-可测函数的极限是 $\mathcal{G}$-可测的; 左端由条件单调收敛定理~\ref{thm:cmct}, $\E[h_n(\rvec{X})\,|\,\mathcal{G}]\uparrow\E[h(\rvec{X})\,|\,\mathcal{G}]$. 两端的极限相等. 一般情形对 $h^+$、$h^-$ 分别处理 (二者积分在几乎每个 $\omega$ 处有限, 因为其期望有限).
\end{proof}

\begin{example}\label{ex:regular-examples}
\begin{enumerate}
\item 若 $\rvec{X}$ 与 $\mathcal{G}$ 独立, 则 $\mu(\omega,\cdot)=\mu_{\rvec{X}}$ (与 $\omega$ 无关) 是正则条件分布 (定理~\ref{thm:ce-basic} (6)).
\item 若 $\rvec{X}$ 是 $\mathcal{G}$-可测的, 则 $\mu(\omega,\cdot)=\delta_{\rvec{X}(\omega)}$ (集中在点 $\rvec{X}(\omega)$ 的测度) 是正则条件分布 (定理~\ref{thm:ce-basic} (5)).
\item 若 $(X,Y)$ 有联合密度, 则 $\mu(\omega,B)=\int_Bf_{X|Y}(x\,|\,Y(\omega))\dd x$ 是 $X$ 关于 $\sigma(Y)$ 的正则条件分布: 对每个 $\omega$, 若 $f_Y(Y(\omega))>0$, 它是以 $f_{X|Y}(\cdot\,|\,Y(\omega))$ 为密度的概率测度 (在 $f_Y(Y(\omega))=0$ 这一零概率事件上可改为任意固定的概率测度); 条件 (2) 由命题~\ref{prop:conditional-density} (取 $h=\1_B$) 给出.
\item 在代入引理~\ref{lem:freezing} 的条件下, $\mu(\omega,B)=\P\big(\phi(\rvec{X},\rvec{y})\in B\big)\big|_{\rvec{y}=\rvec{Y}(\omega)}$ 是 $\phi(\rvec{X},\rvec{Y})$ 关于 $\mathcal{G}$ 的正则条件分布 (对 $\1_B\circ\phi$ 用代入引理).
\end{enumerate}
\end{example}

% =====================================================================
\section{正态随机向量的条件分布}
\label{sec:2-4}

通常条件期望 $\E[X\,|\,Y]$ 是 $Y$ 的非线性函数, 条件分布也可以与无条件分布完全不同. 正态随机向量则是例外: 条件期望是线性的, 条件分布仍是正态的, 而且条件协方差不依赖于观测值. 这使得正态模型下的预测、滤波与插值都可以显式计算. 例~\ref{ex:gaussian-regression} 已经给出了这一结论的雏形, 本节用条件期望的语言完整地叙述它.

\begin{theorem}[正态向量的条件分布]\label{thm:gaussian-conditional}
设 $(\rvec{X},\rvec{Y})$ 是正态随机向量, $\rvec{X}\in\R^k$, $\rvec{Y}\in\R^m$, 期望向量为 $(\rvec{\mu}_{\rvec{X}},\rvec{\mu}_{\rvec{Y}})$, 协方差矩阵分块为
\[
\begin{pmatrix}\mat{\Sigma}_{\rvec{X}\rvec{X}}&\mat{\Sigma}_{\rvec{X}\rvec{Y}}\\ \mat{\Sigma}_{\rvec{Y}\rvec{X}}&\mat{\Sigma}_{\rvec{Y}\rvec{Y}}\end{pmatrix},
\]
其中 $\mat{\Sigma}_{\rvec{X}\rvec{Y}}=\Cov(\rvec{X},\rvec{Y})=\E[(\rvec{X}-\rvec{\mu}_{\rvec{X}})(\rvec{Y}-\rvec{\mu}_{\rvec{Y}})^\T]$, $\mat{\Sigma}_{\rvec{Y}\rvec{X}}=\mat{\Sigma}_{\rvec{X}\rvec{Y}}^\T$. 设 $\mat{\Sigma}_{\rvec{Y}\rvec{Y}}$ 可逆, 令
\[
\mat{K}=\mat{\Sigma}_{\rvec{X}\rvec{Y}}\mat{\Sigma}_{\rvec{Y}\rvec{Y}}^{-1},\qquad \rvec{m}(\rvec{y})=\rvec{\mu}_{\rvec{X}}+\mat{K}(\rvec{y}-\rvec{\mu}_{\rvec{Y}}),\qquad \mat{\Sigma}_{\rvec{X}|\rvec{Y}}=\mat{\Sigma}_{\rvec{X}\rvec{X}}-\mat{\Sigma}_{\rvec{X}\rvec{Y}}\mat{\Sigma}_{\rvec{Y}\rvec{Y}}^{-1}\mat{\Sigma}_{\rvec{Y}\rvec{X}}.
\]
则
\begin{enumerate}
\item $\E[\rvec{X}\,|\,\rvec{Y}]=\rvec{m}(\rvec{Y})$ (按分量理解);
\item $\mu(\omega,\cdot)=N\big(\rvec{m}(\rvec{Y}(\omega)),\mat{\Sigma}_{\rvec{X}|\rvec{Y}}\big)$ 是 $\rvec{X}$ 关于 $\sigma(\rvec{Y})$ 的正则条件分布; 特别地, 条件协方差矩阵 $\E\big[(\rvec{X}-\rvec{m}(\rvec{Y}))(\rvec{X}-\rvec{m}(\rvec{Y}))^\T\,\big|\,\rvec{Y}\big]=\mat{\Sigma}_{\rvec{X}|\rvec{Y}}$ 是确定性的矩阵.
\end{enumerate}
\end{theorem}

\begin{proof}
令 $\rvec{W}=\rvec{X}-\rvec{\mu}_{\rvec{X}}-\mat{K}(\rvec{Y}-\rvec{\mu}_{\rvec{Y}})$.

\emph{$\rvec{W}$ 与 $\rvec{Y}$ 独立.} $(\rvec{W},\rvec{Y})$ 是 $(\rvec{X},\rvec{Y})$ 的仿射变换, 由命题~\ref{prop:gaussian-operations} (1), 它是正态随机向量. 由协方差的双线性,
\[
\Cov(\rvec{W},\rvec{Y})=\mat{\Sigma}_{\rvec{X}\rvec{Y}}-\mat{K}\mat{\Sigma}_{\rvec{Y}\rvec{Y}}=\mat{\Sigma}_{\rvec{X}\rvec{Y}}-\mat{\Sigma}_{\rvec{X}\rvec{Y}}\mat{\Sigma}_{\rvec{Y}\rvec{Y}}^{-1}\mat{\Sigma}_{\rvec{Y}\rvec{Y}}=\mat{0}.
\]
由定理~\ref{thm:gaussian-uncorrelated}, $\rvec{W}$ 与 $\rvec{Y}$ 独立.

\emph{$\rvec{W}$ 的分布.} $\E\rvec{W}=\rvec{0}$; 由命题~\ref{prop:covariance-matrix} (2) (把 $\rvec{W}$ 写成 $(\mat{I},-\mat{K})$ 乘以 $(\rvec{X},\rvec{Y})$ 再平移),
\begin{align*}
\Cov(\rvec{W})&=\mat{\Sigma}_{\rvec{X}\rvec{X}}-\mat{K}\mat{\Sigma}_{\rvec{Y}\rvec{X}}-\mat{\Sigma}_{\rvec{X}\rvec{Y}}\mat{K}^\T+\mat{K}\mat{\Sigma}_{\rvec{Y}\rvec{Y}}\mat{K}^\T\\
&=\mat{\Sigma}_{\rvec{X}\rvec{X}}-2\mat{\Sigma}_{\rvec{X}\rvec{Y}}\mat{\Sigma}_{\rvec{Y}\rvec{Y}}^{-1}\mat{\Sigma}_{\rvec{Y}\rvec{X}}+\mat{\Sigma}_{\rvec{X}\rvec{Y}}\mat{\Sigma}_{\rvec{Y}\rvec{Y}}^{-1}\mat{\Sigma}_{\rvec{Y}\rvec{Y}}\mat{\Sigma}_{\rvec{Y}\rvec{Y}}^{-1}\mat{\Sigma}_{\rvec{Y}\rvec{X}}=\mat{\Sigma}_{\rvec{X}|\rvec{Y}},
\end{align*}
这里用到 $\mat{\Sigma}_{\rvec{Y}\rvec{Y}}^{-1}$ 对称. 所以 $\rvec{W}\sim N(\rvec{0},\mat{\Sigma}_{\rvec{X}|\rvec{Y}})$.

\emph{代入.} $\rvec{X}=\rvec{m}(\rvec{Y})+\rvec{W}$, 其中 $\rvec{W}$ 与 $\sigma(\rvec{Y})$ 独立, $\rvec{Y}$ 是 $\sigma(\rvec{Y})$-可测的. 对 Borel 集 $B\subseteq\R^k$, 对有界函数 $\phi(\rvec{w},\rvec{y})=\1_B(\rvec{m}(\rvec{y})+\rvec{w})$ 使用代入引理~\ref{lem:freezing} (1),
\[
\P(\rvec{X}\in B\,|\,\rvec{Y})=\P\big(\rvec{m}(\rvec{y})+\rvec{W}\in B\big)\big|_{\rvec{y}=\rvec{Y}}=N\big(\rvec{m}(\rvec{Y}),\mat{\Sigma}_{\rvec{X}|\rvec{Y}}\big)(B),
\]
最后一步用到命题~\ref{prop:gaussian-operations} (1): $\rvec{m}(\rvec{y})+\rvec{W}\sim N(\rvec{m}(\rvec{y}),\mat{\Sigma}_{\rvec{X}|\rvec{Y}})$. 对每个 $\omega$, 右端是概率测度, 这证明了 (2) 的第一部分. 对 $\rvec{X}$ 的每个分量 $X_i=m_i(\rvec{Y})+W_i$ 用 $\mathcal{G}$-可测变量的条件期望 (定理~\ref{thm:ce-basic} (5)) 与独立性化简 (定理~\ref{thm:ce-basic} (6)), $\E[X_i\,|\,\rvec{Y}]=m_i(\rvec{Y})+\E[W_i]=m_i(\rvec{Y})$, 这是 (1). 最后, $(\rvec{X}-\rvec{m}(\rvec{Y}))(\rvec{X}-\rvec{m}(\rvec{Y}))^\T=\rvec{W}\rvec{W}^\T$ 与 $\sigma(\rvec{Y})$ 独立, 其条件期望等于 $\E[\rvec{W}\rvec{W}^\T]=\mat{\Sigma}_{\rvec{X}|\rvec{Y}}$.
\end{proof}

\begin{intuition}
定理把 $\rvec{X}$ 分解为两部分: $\rvec{m}(\rvec{Y})$ 是 $\rvec{Y}$ 能解释的部分, $\rvec{W}$ 是与 $\rvec{Y}$ 完全无关的噪声. 对正态向量, 这个 ``线性解释 + 独立噪声'' 的分解是精确的, 所以第一章的最佳\emph{线性}预测 (例~\ref{ex:best-linear}) 就是最佳预测, 不存在任何非线性改进的余地. 矩阵 $\mat{\Sigma}_{\rvec{X}|\rvec{Y}}$ 是协方差矩阵的 Schur 补, 它半正定, 且 $\mat{\Sigma}_{\rvec{X}|\rvec{Y}}\le\mat{\Sigma}_{\rvec{X}\rvec{X}}$ (差为半正定矩阵 $\mat{\Sigma}_{\rvec{X}\rvec{Y}}\mat{\Sigma}_{\rvec{Y}\rvec{Y}}^{-1}\mat{\Sigma}_{\rvec{Y}\rvec{X}}$): 观测只会减少不确定性, 而且减少的幅度与观测到的具体数值无关.
\end{intuition}

\begin{remark}
若 $\mat{\Sigma}_{\rvec{Y}\rvec{Y}}$ 奇异, 把 $\mat{\Sigma}_{\rvec{Y}\rvec{Y}}^{-1}$ 换成 Moore-Penrose 广义逆, 定理仍然成立. 这是因为此时 $\rvec{Y}$ 几乎必然落在 $\mat{\Sigma}_{\rvec{Y}\rvec{Y}}$ 的值域 (平移后) 上 (命题~\ref{prop:covariance-matrix} (3)), 证明中只需在该子空间上进行计算.
\end{remark}

二维情形最常用. 设 $(X,Y)$ 是二维正态随机向量, $\E X=\mu_X$, $\E Y=\mu_Y$, $\Var(X)=\sigma_X^2$, $\Var(Y)=\sigma_Y^2>0$, 相关系数为 $\rho$. 由定理~\ref{thm:gaussian-conditional}, 已知 $Y$ 时 $X$ 的条件分布为
\begin{equation}\label{eq:bivariate-conditional}
N\Big(\mu_X+\rho\frac{\sigma_X}{\sigma_Y}(Y-\mu_Y),\ \sigma_X^2(1-\rho^2)\Big).
\end{equation}

\begin{figure}[htbp]
\centering
\begin{tikzpicture}[scale=1.0]
  \begin{scope}
    \clip (-3.3,-2.4) rectangle (3.3,2.4);
    \foreach \c/\o in {1.45/10, 1/22, 0.55/40} {
      \fill[SABlue!\o,rotate=37.98] (0,0) ellipse ({1.6676*\c} and {0.8481*\c});
    }
    \foreach \c in {0.55,1,1.45} {
      \draw[SABlue!80!black,rotate=37.98] (0,0) ellipse ({1.6676*\c} and {0.8481*\c});
    }
    \draw[SADark!60,dashed,thick] (-3.3,{-3.3*0.7808}) -- (3.3,{3.3*0.7808});
    \draw[SARed,very thick] (-3.3,-1.65) -- (3.3,1.65);
  \end{scope}
  \draw[->,SADark!70] (-3.4,0) -- (3.5,0) node[right,font=\small] {$y$};
  \draw[->,SADark!70] (0,-2.5) -- (0,2.6) node[above,font=\small] {$x$};
  \foreach \s in {1,-1} {
    \fill[SARed] ({\s*2.0506},{\s*1.0253}) circle (1.8pt);
    \draw[SARed!60,thin] ({\s*2.0506},-2.4) -- ({\s*2.0506},2.4);
  }
  \node[SARed,font=\small,anchor=west] at (3.35,1.65) {$x=\E[X\,|\,Y=y]$};
  \node[SADark!80,font=\small,anchor=west] at (3.35,2.45) {主轴};
\end{tikzpicture}
\caption{二维正态分布的等高椭圆 ($\Var(Y)=2$, $\Var(X)=1.5$, $\Cov(X,Y)=1$). 回归直线 $x=\E[X\,|\,Y=y]=\frac{1}{2}y$ (实线) 经过椭圆上竖直切线的切点, 斜率小于椭圆主轴 (虚线) 的斜率: 这就是 ``向均值回归''}
\label{fig:regression}
\end{figure}

\begin{intuition}
图~\ref{fig:regression} 解释了 Galton 发现的 ``回归'' 现象: 父亲身高 ($Y$) 比平均高出 $d$ 时, 儿子身高 ($X$) 的条件期望只比平均高出 $\rho\frac{\sigma_X}{\sigma_Y}d$, 当 $\sigma_X=\sigma_Y$ 时就是 $\rho d<d$. 这并不意味着下一代的差异会缩小: 条件方差 $\sigma_X^2(1-\rho^2)$ 补上了被 ``回归'' 掉的那部分方差, 总方差由全方差公式~\ref{thm:total-variance} 保持为 $\rho^2\sigma_X^2+\sigma_X^2(1-\rho^2)=\sigma_X^2$.
\end{intuition}

\begin{example}[带噪声的观测]\label{ex:signal-noise}
设未知量 $X\sim N(m,\sigma^2)$, 我们对它作 $n$ 次独立的带噪声观测 $Y_i=X+\varepsilon_i$, 其中 $\varepsilon_1,\dots,\varepsilon_n$ 是 $\iid$ 的 $N(0,\tau^2)$, 且与 $X$ 独立. 由命题~\ref{prop:gaussian-operations} (2), $(X,\varepsilon_1,\dots,\varepsilon_n)$ 是正态随机向量, 从而 $(X,Y_1,\dots,Y_n)$ 也是. 记 $\rvec{1}=(1,\dots,1)^\T\in\R^n$, $\rvec{Y}=(Y_1,\dots,Y_n)^\T$, 则
\[
\mat{\Sigma}_{X\rvec{Y}}=\sigma^2\rvec{1}^\T,\qquad \mat{\Sigma}_{\rvec{Y}\rvec{Y}}=\tau^2\mat{I}+\sigma^2\rvec{1}\rvec{1}^\T.
\]
我们断言
\[
\mat{\Sigma}_{\rvec{Y}\rvec{Y}}^{-1}=\frac{1}{\tau^2}\Big(\mat{I}-\frac{\sigma^2}{\tau^2+n\sigma^2}\rvec{1}\rvec{1}^\T\Big).
\]
事实上, 记 $c=\frac{\sigma^2}{\tau^2+n\sigma^2}$, 利用 $\rvec{1}\rvec{1}^\T\rvec{1}\rvec{1}^\T=n\,\rvec{1}\rvec{1}^\T$ 直接相乘,
\begin{align*}
\big(\tau^2\mat{I}+\sigma^2\rvec{1}\rvec{1}^\T\big)\big(\mat{I}-c\,\rvec{1}\rvec{1}^\T\big)
&=\tau^2\mat{I}+\big(\sigma^2-c\tau^2-cn\sigma^2\big)\rvec{1}\rvec{1}^\T\\
&=\tau^2\mat{I}+\big(\sigma^2-c(\tau^2+n\sigma^2)\big)\rvec{1}\rvec{1}^\T=\tau^2\mat{I}.
\end{align*}
于是 (再次利用 $\rvec{1}^\T\rvec{1}\rvec{1}^\T=n\rvec{1}^\T$)
\[
\mat{K}=\sigma^2\rvec{1}^\T\mat{\Sigma}_{\rvec{Y}\rvec{Y}}^{-1}=\frac{\sigma^2}{\tau^2}\Big(1-\frac{n\sigma^2}{\tau^2+n\sigma^2}\Big)\rvec{1}^\T=\frac{\sigma^2}{\tau^2+n\sigma^2}\rvec{1}^\T.
\] 记 $\bar Y=\frac{1}{n}\sum_iY_i$, 由定理~\ref{thm:gaussian-conditional},
\begin{align*}
\E[X\,|\,Y_1,\dots,Y_n]&=m+\frac{n\sigma^2}{\tau^2+n\sigma^2}(\bar Y-m),\\
\Var(X\,|\,Y_1,\dots,Y_n)&=\sigma^2-\frac{n\sigma^4}{\tau^2+n\sigma^2}=\frac{\sigma^2\tau^2}{\tau^2+n\sigma^2}.
\end{align*}
\begin{itemize}
\item 估计值是先验均值 $m$ 与样本均值 $\bar Y$ 的加权平均, 权重 $\frac{n\sigma^2}{\tau^2+n\sigma^2}$ 随观测次数增加而趋于 $1$, 即数据越多, 越相信数据.
\item 条件方差满足 $\frac{1}{\Var(X\,|\,\rvec{Y})}=\frac{1}{\sigma^2}+\frac{n}{\tau^2}$, 即 ``精度'' (方差的倒数) 是可加的, 每次观测贡献 $\frac{1}{\tau^2}$.
\item 条件期望只通过 $\bar Y$ 依赖于观测, 估计 $X$ 所需的全部有用信息都已包含在样本均值中.
\end{itemize}
这一计算是 Kalman 滤波的单步版本.
\end{example}

\begin{example}[布朗桥的雏形]\label{ex:bridge-preview}
设 $(U,V)$ 是均值为零的二维正态随机向量, $\Var(U)=s$, $\Var(V)=t$, $\Cov(U,V)=s$, 其中 $0<s<t$. (第~4~章将看到, 这正是布朗运动 $(B_s,B_t)$ 的联合分布.) 由 \eqref{eq:bivariate-conditional} 及 $\rho\frac{\sigma_U}{\sigma_V}=\frac{\Cov(U,V)}{\Var(V)}$, $\rho^2=\frac{s^2}{st}$,
\[
U\,|\,V\ \sim\ N\Big(\frac{s}{t}V,\ \frac{s(t-s)}{t}\Big),\qquad V\,|\,U\ \sim\ N\big(U,\ t-s\big).
\]
第二式说明: 已知 $U$ 时, $V-U$ 与 $U$ 无关且服从 $N(0,t-s)$, 这就是布朗运动的独立增量性. 第一式说明: 已知终点 $V$ 时, 中间时刻的最佳估计是线性插值 $\frac{s}{t}V$, 条件方差 $\frac{s(t-s)}{t}$ 在两端为零、在中点最大. 取 $t=1$ 并以 $V=0$ 为条件, 得到方差 $s(1-s)$, 这就是第~4~章布朗桥 $B_s-sB_1$ 的方差.
\end{example}

% =====================================================================
\section{随机过程与马尔可夫性}
\label{sec:2-5}

到目前为止, 我们研究的是有限个随机变量. 随机分析的对象则是随时间演化的一族随机变量: 随机过程. 描述随机过程需要三个层次的概念: 过程本身及其分布 (2.5.1--2.5.2 小节); 随时间积累的信息, 即信息流 (2.5.3 小节); 以及过程与信息之间最重要的一类关系, 即马尔可夫性 (2.5.4--2.5.6 小节).

% ---------------------------------------------------------------------
\subsection{随机过程及其有限维分布}

\begin{definition}[随机过程]\label{def:process}
设 $T\subseteq[0,\infty)$ 是指标集 (通常为 $\N=\{0,1,2,\dots\}$, $[0,\infty)$ 或 $[0,T_0]$). 定义在同一概率空间 $(\Omega,\mathcal{F},\P)$ 上的一族随机变量 $X=\{X_t\}_{t\in T}$ 称为\textbf{随机过程}. $T$ 可数时称为离散时间过程 (随机序列), $T$ 为区间时称为连续时间过程. 对固定的 $\omega$, 函数 $t\mapsto X_t(\omega)$ 称为过程的一条\textbf{样本轨道}.
\end{definition}

随机过程有两种看法: 固定 $t$, $X_t$ 是一个随机变量; 固定 $\omega$, $X_\cdot(\omega)$ 是一个函数. 后一种看法把随机过程视为 ``随机函数'', 是随机分析的出发点: 布朗运动的轨道连续而处处不可微 (第~4~章), 正是这一点使得经典微积分失效而需要 Itô 积分 (第~5~章).

\begin{figure}[htbp]
\centering
\begin{tikzpicture}
\begin{axis}[SAaxis, width=11.5cm, height=6cm, xmin=0, xmax=21, ymin=-8, ymax=10,
  xtick={0,5,10,15,20}, ytick={-6,-3,3,6,9}, xlabel={$n$}, ylabel={$S_n(\omega)$}]
  \addplot[SABlue,thick,mark=*,mark size=1pt] coordinates {(0,0) (1,-1) (2,0) (3,-1) (4,0) (5,1) (6,2) (7,1) (8,2) (9,3) (10,4) (11,5) (12,4) (13,3) (14,4) (15,5) (16,6) (17,5) (18,6) (19,7) (20,8)};
  \addplot[SARed,thick,mark=*,mark size=1pt] coordinates {(0,0) (1,1) (2,2) (3,1) (4,2) (5,3) (6,4) (7,5) (8,6) (9,7) (10,6) (11,7) (12,8) (13,9) (14,8) (15,7) (16,6) (17,7) (18,6) (19,5) (20,4)};
  \addplot[SAGreenText,thick,mark=*,mark size=1pt] coordinates {(0,0) (1,-1) (2,-2) (3,-3) (4,-4) (5,-5) (6,-6) (7,-5) (8,-4) (9,-3) (10,-2) (11,-3) (12,-4) (13,-5) (14,-6) (15,-7) (16,-6) (17,-5) (18,-6) (19,-5) (20,-6)};
  \addplot[SAYellowText,thick,dashed] coordinates {(12,-8) (12,10)};
  \node[font=\small,anchor=west] at (axis cs:12.2,-7.3) {固定 $n=12$: 随机变量 $S_{12}$};
\end{axis}
\end{tikzpicture}
\caption{简单对称随机游动 $S_n=\xi_1+\cdots+\xi_n$ ($\xi_k$ 以等概率取 $\pm1$) 的三条样本轨道. 固定 $\omega$ 看是一条折线; 固定时间 (虚线) 看是一个随机变量, 三条轨道在虚线上的取值是它的三个可能结果}
\label{fig:rw-paths}
\end{figure}

\begin{example}\label{ex:processes}
\begin{enumerate}
\item \textbf{独立同分布序列} $\{\xi_n\}_{n\ge1}$: 最简单的离散时间过程, 也称为 (离散) 白噪声.
\item \textbf{随机游动} $S_n=S_0+\sum_{k=1}^n\xi_k$, 其中 $\xi_k$ 是 $\iid$ 的且与 $S_0$ 独立. 图~\ref{fig:rw-paths} 是 $S_0=0$、$\xi_k=\pm1$ 时的情形.
\item \textbf{随机相位过程} $X_t=\cos(t+\Theta)$, $t\ge0$, 其中 $\Theta\sim\mathrm{U}(0,2\pi)$. 每条轨道都是一条平移后的余弦曲线, 只有一个随机参数. 由积化和差,
\[
\E[X_sX_t]=\frac{1}{2}\E\big[\cos(t-s)+\cos(s+t+2\Theta)\big]=\frac{1}{2}\cos(t-s),
\]
且 $\E X_t=0$. 协方差只依赖于时间差 $t-s$.
\item \textbf{高斯过程}: 若对任意 $t_1,\dots,t_n\in T$, $(X_{t_1},\dots,X_{t_n})$ 都是正态随机向量, 称 $X$ 为高斯过程. 它的全部概率性质由均值函数 $m(t)=\E X_t$ 与协方差函数 $K(s,t)=\Cov(X_s,X_t)$ 决定 (见下面的推论~\ref{cor:gaussian-process}).
\end{enumerate}
\end{example}

\begin{definition}[有限维分布]\label{def:fdd}
对互不相同的 $t_1,\dots,t_n\in T$, 随机向量 $(X_{t_1},\dots,X_{t_n})$ 的联合分布
\[
\mu_{t_1,\dots,t_n}(B)=\P\big((X_{t_1},\dots,X_{t_n})\in B\big),\qquad B\in\mathcal{B}(\R^n)
\]
称为 $X$ 的\textbf{有限维分布}. 全体有限维分布称为 $X$ 的有限维分布族.
\end{definition}

有限维分布族自然满足两个\textbf{相容性条件}:
\begin{enumerate}
\item[(C1)] 对 $(1,\dots,n)$ 的任意置换 $\pi$, $\mu_{t_{\pi(1)},\dots,t_{\pi(n)}}$ 是 $\mu_{t_1,\dots,t_n}$ 在坐标置换 $(x_1,\dots,x_n)\mapsto(x_{\pi(1)},\dots,x_{\pi(n)})$ 下的像;
\item[(C2)] 边缘相容: $\mu_{t_1,\dots,t_n,t_{n+1}}(B\times\R)=\mu_{t_1,\dots,t_n}(B)$.
\end{enumerate}
反过来, 相容性条件也足以保证过程的存在. 设 $\R^T$ 为一切函数 $\omega:T\to\R$ 的集合, 其上的乘积 $\sigma$-代数 $\mathcal{B}(\R)^{\otimes T}$ 由全体\textbf{柱集} $\{\omega: (\omega(t_1),\dots,\omega(t_n))\in B\}$ 生成.

\begin{theorem}[Kolmogorov 延拓定理]\label{thm:kolmogorov-extension}
设 $\{\mu_{t_1,\dots,t_n}\}$ 是满足 (C1)(C2) 的一族概率测度 ($\mu_{t_1,\dots,t_n}$ 定义在 $\mathcal{B}(\R^n)$ 上). 则 $(\R^T,\mathcal{B}(\R)^{\otimes T})$ 上存在唯一的概率测度 $\P$, 使得坐标过程 $X_t(\omega)=\omega(t)$ 的有限维分布恰为 $\{\mu_{t_1,\dots,t_n}\}$. (证明见~\cite{durrett} 附录 A.3 或~\cite{kallenberg} 定理 8.23.)
\end{theorem}

唯一性由柱集构成 $\pi$-系与推论~\ref{cor:uniqueness} 得到; 存在性的证明依赖于 Carath\'eodory 扩张定理~\ref{thm:caratheodory} 与 $\R^n$ 的拓扑性质.

\begin{corollary}[高斯过程的存在性]\label{cor:gaussian-process}
设 $m:T\to\R$, $K:T\times T\to\R$ 对称且半正定, 即对任意 $n$、$t_1,\dots,t_n\in T$ 与 $a_1,\dots,a_n\in\R$, $\sum_{i,j}a_ia_jK(t_i,t_j)\ge0$. 则存在高斯过程 $X$, 使 $\E X_t=m(t)$, $\Cov(X_s,X_t)=K(s,t)$.
\end{corollary}

\begin{proof}
令 $\mu_{t_1,\dots,t_n}=N\big((m(t_i))_i,(K(t_i,t_j))_{i,j}\big)$, 由假设协方差矩阵半正定, 由定理~\ref{thm:gaussian-equivalent} 该分布存在. 验证相容性: (C1) 正态随机向量的分量置换是线性变换, 由命题~\ref{prop:gaussian-operations} (1), 置换后服从均值与协方差相应置换后的正态分布; (C2) 去掉最后一个分量也是线性变换, 得到的子向量服从 $N\big((m(t_i))_{i\le n},(K(t_i,t_j))_{i,j\le n}\big)$. 由 Kolmogorov 延拓定理~\ref{thm:kolmogorov-extension} 即得.
\end{proof}

\begin{example}[布朗运动的协方差]\label{ex:min-kernel}
$K(s,t)=s\wedge t$ ($s,t\ge0$) 是半正定的. 事实上, $s\wedge t=\int_0^\infty\1_{[0,s]}(u)\1_{[0,t]}(u)\dd u$, 所以
\[
\sum_{i,j}a_ia_j(t_i\wedge t_j)=\int_0^\infty\Big(\sum_ia_i\1_{[0,t_i]}(u)\Big)^2\dd u\ge0.
\]
由推论~\ref{cor:gaussian-process}, 存在均值为零、协方差为 $s\wedge t$ 的高斯过程. 但这样得到的过程是定义在 $\R^{[0,\infty)}$ 上的, 我们无法断言它的轨道连续: 可以证明, $\mathcal{B}(\R)^{\otimes[0,\infty)}$ 中的每个事件都只依赖于可列个坐标 (见~\cite{billingsley} 第 36 节), 而 ``轨道连续'' 依赖于不可列个坐标, 不是这个 $\sigma$-代数中的事件. 第~4~章将通过修改过程的 ``版本'' 来得到连续轨道, 下一小节准备所需的概念.
\end{example}

% ---------------------------------------------------------------------
\subsection{修正与不可区分}

\begin{definition}[修正与不可区分]\label{def:modification}
设 $X,Y$ 是定义在同一概率空间上、指标集相同的随机过程.
\begin{enumerate}
\item 若对每个 $t\in T$, $\P(X_t=Y_t)=1$, 称 $Y$ 是 $X$ 的\textbf{修正} (或版本).
\item 若 $\P(X_t=Y_t\text{ 对一切 }t\in T\text{ 成立})=1$, 称 $X$ 与 $Y$ \textbf{不可区分}.
\end{enumerate}
\end{definition}

不可区分显然蕴含互为修正; $T$ 可数时二者等价 (可列个零概率集之并仍是零概率集); $T$ 不可数时, 二者有本质区别.

\begin{example}\label{ex:modification}
设 $\tau\sim\mathrm{U}(0,1)$, 在 $T=[0,1]$ 上令 $X_t\equiv0$, $Y_t=\1_{\{t=\tau\}}$. 对每个固定的 $t$, $\P(X_t\ne Y_t)=\P(\tau=t)=0$, 所以 $Y$ 是 $X$ 的修正, 二者有相同的有限维分布 (见下面的命题). 但每条轨道 $Y_\cdot(\omega)$ 都在 $t=\tau(\omega)$ 处跳到 $1$, 所以 $\P(X_t=Y_t\ \forall t)=0$. $X$ 的轨道连续而 $Y$ 的轨道处处不连续. 有限维分布 ``看不见'' 不可列个时刻上的行为.
\end{example}

\begin{figure}[htbp]
\centering
\begin{tikzpicture}[xscale=6,yscale=1.4]
  \draw[->,SADark!80] (0,0) -- (1.1,0) node[right,font=\small] {$t$};
  \draw[->,SADark!80] (0,0) -- (0,1.3);
  \draw[SABlue,very thick] (0,0) -- (0.62,0);
  \draw[SABlue,very thick] (0.62,0) -- (1,0);
  \draw[SABlue,fill=white] (0.62,0) circle (0.25pt);
  \fill[SARed] (0.62,1) circle (0.25pt);
  \draw[SADark!30,dashed] (0.62,0) -- (0.62,1);
  \node[below,font=\small] at (0.62,0) {$\tau(\omega)$};
  \node[left,font=\small] at (0,1) {$1$};
  \node[below,font=\small] at (1,0) {$1$};
  \node[font=\small,anchor=west,SARed] at (0.65,1) {$Y_\tau(\omega)=1$};
  \node[font=\small,anchor=south,SABlue] at (0.3,0.05) {$Y_t(\omega)=0$};
\end{tikzpicture}
\caption{例~\ref{ex:modification}: $Y$ 的一条轨道. 它只在随机时刻 $\tau(\omega)$ 处不为零. 对任一\emph{固定}时刻 $t$, 轨道恰好在 $t$ 处跳跃的概率为零}
\label{fig:modification}
\end{figure}

\begin{proposition}\label{prop:modification}
\begin{enumerate}
\item 若 $Y$ 是 $X$ 的修正, 则 $X$ 与 $Y$ 有相同的有限维分布.
\item 设 $T=[0,\infty)$ 或 $[0,T_0]$, $Y$ 是 $X$ 的修正, 且 $X$ 与 $Y$ 的轨道几乎必然右连续. 则 $X$ 与 $Y$ 不可区分.
\end{enumerate}
\end{proposition}

\begin{proof}
(1) 对 $t_1,\dots,t_n$, 事件 $N=\bigcup_i\{X_{t_i}\ne Y_{t_i}\}$ 的概率为零, 而在 $N^c$ 上两个随机向量相等, 所以它们的分布相同.

(2) 令 $Q$ 为 $T$ 中的有理数与右端点 (若有) 构成的可数集, $N_0=\bigcup_{q\in Q}\{X_q\ne Y_q\}$, $N_1$ 为 $X$ 或 $Y$ 的轨道不右连续的零概率集 (更准确地说, 是包含这些 $\omega$ 的某个零概率事件). 在 $(N_0\cup N_1)^c$ 上, 对任意 $t\in T$ (若 $t$ 是右端点则它属于 $Q$), 取 $q_k\in Q$, $q_k\downarrow t$, 由右连续性,
$X_t=\lim_kX_{q_k}=\lim_kY_{q_k}=Y_t$. 所以 $\P(X_t=Y_t\ \forall t)\ge1-\P(N_0\cup N_1)=1$.
\end{proof}

% ---------------------------------------------------------------------
\subsection{信息流与适应过程}

随着时间推移, 观察者掌握的信息越来越多. 用一族递增的 $\sigma$-代数表示这一点.

\begin{definition}[信息流]\label{def:filtration}
一族子 $\sigma$-代数 $\{\mathcal{F}_t\}_{t\in T}$, 若满足 $s\le t\Rightarrow\mathcal{F}_s\subseteq\mathcal{F}_t\subseteq\mathcal{F}$, 称为\textbf{信息流} (或 $\sigma$-代数流, filtration). $(\Omega,\mathcal{F},\{\mathcal{F}_t\},\P)$ 称为\textbf{带流的概率空间}. 记 $\mathcal{F}_\infty=\sigma\big(\bigcup_t\mathcal{F}_t\big)$.
\end{definition}

\begin{definition}[适应过程与自然信息流]\label{def:adapted}
若对每个 $t$, $X_t$ 是 $\mathcal{F}_t$-可测的, 称过程 $X$ \textbf{适应}于 $\{\mathcal{F}_t\}$. 称 $\mathcal{F}^X_t=\sigma(X_s: s\le t,\ s\in T)$ 为 $X$ 的\textbf{自然信息流}; 它是使 $X$ 适应的最小信息流.
\end{definition}

``$X$ 适应于 $\{\mathcal{F}_t\}$'' 是指在时刻 $t$, 观察者已经知道 $X_t$ 的值. 信息流可以比自然信息流更大, 例如它可以包含其他过程的信息.

\begin{example}[抛硬币的信息流]\label{ex:coin-filtration}
在无穷次抛硬币中, 令 $\mathcal{F}_0=\{\varnothing,\Omega\}$, $\mathcal{F}_n=\sigma(\omega_1,\dots,\omega_n)$. $\mathcal{F}_n$ 由 $2^n$ 个长为 $n$ 的柱集构成的分割生成, 每增加一次抛掷, 分割的每一块就一分为二 (图~\ref{fig:coin-tree}). $S_n=\omega_1+\cdots+\omega_n$ 与 $M_n=\max_{k\le n}S_k$ 都适应于 $\{\mathcal{F}_n\}$; 而 $S_{n+1}$ 不是 $\mathcal{F}_n$-可测的, 因为它依赖于尚未发生的第 $n+1$ 次抛掷.
\end{example}

\begin{figure}[htbp]
\centering
\begin{tikzpicture}[scale=0.82,
  every node/.style={font=\small},
  nd/.style={circle,draw=SADark,fill=SABlue!12,inner sep=1.4pt,minimum size=5mm},
  level distance=1.25cm, sibling distance=2.0cm]
  \node[nd] (r) at (0,0) {$\varnothing$};
  \node[nd] (a) at (-3.2,-1.3) {$0$};
  \node[nd] (b) at (3.2,-1.3) {$1$};
  \foreach \p/\x/\l in {a/-4.8/00, a/-1.6/01, b/1.6/10, b/4.8/11} { \node[nd] (\l) at (\x,-2.6) {\l}; \draw[SADark!60] (\p) -- (\l); }
  \foreach \p/\x/\l in {00/-5.6/000, 00/-4.0/001, 01/-2.4/010, 01/-0.8/011, 10/0.8/100, 10/2.4/101, 11/4.0/110, 11/5.6/111} { \node[nd,inner sep=0.6pt] (n\l) at (\x,-3.9) {\scriptsize\l}; \draw[SADark!60] (\p) -- (n\l); }
  \draw[SADark!60] (r) -- (a); \draw[SADark!60] (r) -- (b);
  \draw[SARed,very thick,->] (r) -- (b); \draw[SARed,very thick,->] (b) -- (10); \draw[SARed,very thick,->] (10) -- (n101);
  \foreach \y/\l in {0/{$\mathcal{F}_0$: 1 块}, -1.3/{$\mathcal{F}_1$: 2 块}, -2.6/{$\mathcal{F}_2$: 4 块}, -3.9/{$\mathcal{F}_3$: 8 块}} { \node[anchor=west,SADark] at (6.4,\y) {\l}; }
\end{tikzpicture}
\caption{抛硬币的信息流. 第 $n$ 层的每个结点代表 $\mathcal{F}_n$ 的分割中的一块 (前 $n$ 次的结果). 一个样本点 $\omega$ 对应树中一条无穷路径 (粗线表示 $\omega_1\omega_2\omega_3=101$); 时刻 $n$ 的信息就是 ``路径走到了第 $n$ 层的哪个结点''}
\label{fig:coin-tree}
\end{figure}

在离散时间中, 还需要区分 ``时刻 $n$ 已知'' 与 ``时刻 $n$ 之前已知''.

\begin{definition}[可料序列]\label{def:predictable}
设 $\{\mathcal{F}_n\}_{n\ge0}$ 是离散时间信息流. 若对每个 $n\ge1$, $H_n$ 是 $\mathcal{F}_{n-1}$-可测的, 称序列 $\{H_n\}_{n\ge1}$ 为\textbf{可料的}.
\end{definition}

\begin{example}[赌博策略]\label{ex:betting}
设 $\xi_1,\xi_2,\dots$ 是各局的输赢 ($\pm1$), $\mathcal{F}_n=\sigma(\xi_1,\dots,\xi_n)$. 赌徒在第 $n$ 局开始前决定赌注 $H_n$, 他只能利用前 $n-1$ 局的结果, 所以 $\{H_n\}$ 是可料的. 他在第 $n$ 局后的财富为
\[
W_n=W_0+\sum_{k=1}^nH_k\xi_k.
\]
例如 ``加倍策略'': $H_1=1$, 若前 $n-1$ 局全输则 $H_n=2^{n-1}$, 一旦赢了就停止下注 ($H_n=0$). 第~3~章将这类和式称为鞅变换, 并证明在公平赌博中, 有界期限内的任何可料策略都不能使期望财富增加. 第~5~章的 Itô 积分 $\int H_s\dd B_s$ 是它的连续时间版本: 要求被积过程在每个小区间的\emph{左端点}取值, 正是为了保持这种 ``可料'' 的结构.
\end{example}

\begin{remark}
连续时间中可料性与可测性的精确定义 (可料 $\sigma$-代数、循序可测过程) 涉及 $(t,\omega)$ 的联合可测性, 将在第~5~章构造 Itô 积分时引入.
\end{remark}

% ---------------------------------------------------------------------
\subsection{条件独立}

马尔可夫性的直观含义是 ``给定现在, 过去与未来无关''. 为了严格表述它, 需要条件独立的概念.

\begin{definition}[条件独立]\label{def:conditional-independence}
称子 $\sigma$-代数 $\mathcal{G}_1$ 与 $\mathcal{G}_2$ 关于 $\mathcal{H}$ \textbf{条件独立}, 如果对一切 $A\in\mathcal{G}_1$, $B\in\mathcal{G}_2$,
\[
\P(A\cap B\,|\,\mathcal{H})=\P(A\,|\,\mathcal{H})\,\P(B\,|\,\mathcal{H}).
\]
\end{definition}

取 $\mathcal{H}=\{\varnothing,\Omega\}$ 就回到通常的独立性. 下面的等价刻画更便于使用: 在已知 $\mathcal{H}$ 的前提下, 额外知道 $\mathcal{G}_1$ 对预测 $\mathcal{G}_2$ 中的事件毫无帮助.

\begin{proposition}\label{prop:conditional-independence}
$\mathcal{G}_1$ 与 $\mathcal{G}_2$ 关于 $\mathcal{H}$ 条件独立, 当且仅当对一切 $B\in\mathcal{G}_2$,
\begin{equation}\label{eq:ci-equivalent}
\P\big(B\,\big|\,\sigma(\mathcal{G}_1\cup\mathcal{H})\big)=\P(B\,|\,\mathcal{H}).
\end{equation}
\end{proposition}

\begin{proof}
``当'': 设 \eqref{eq:ci-equivalent} 成立, $A\in\mathcal{G}_1$, $B\in\mathcal{G}_2$. 由塔性质与提出已知因子 ($\1_A$ 是 $\sigma(\mathcal{G}_1\cup\mathcal{H})$-可测的, $\P(B\,|\,\mathcal{H})$ 是 $\mathcal{H}$-可测的),
\begin{align*}
\P(A\cap B\,|\,\mathcal{H})&=\E\Big[\E\big[\1_A\1_B\,\big|\,\sigma(\mathcal{G}_1\cup\mathcal{H})\big]\,\Big|\,\mathcal{H}\Big]=\E\Big[\1_A\,\P\big(B\,\big|\,\sigma(\mathcal{G}_1\cup\mathcal{H})\big)\,\Big|\,\mathcal{H}\Big]\\
&=\E\big[\1_A\,\P(B\,|\,\mathcal{H})\,\big|\,\mathcal{H}\big]=\P(B\,|\,\mathcal{H})\,\P(A\,|\,\mathcal{H}).
\end{align*}

``仅当'': 设条件独立, $B\in\mathcal{G}_2$. $\P(B\,|\,\mathcal{H})$ 是 $\sigma(\mathcal{G}_1\cup\mathcal{H})$-可测的. 集族 $\{A\cap H: A\in\mathcal{G}_1,H\in\mathcal{H}\}$ 是包含 $\Omega$ 的 $\pi$-系, 生成 $\sigma(\mathcal{G}_1\cup\mathcal{H})$; 由引理~\ref{lem:ce-pi-system}, 只需对 $A\cap H$ 验证积分恒等式. 由全期望公式、提出已知因子 ($\1_H\P(B\,|\,\mathcal{H})$ 是 $\mathcal{H}$-可测的) 与条件独立性,
\begin{align*}
\E\big[\P(B\,|\,\mathcal{H});A\cap H\big]&=\E\Big[\E\big[\1_A\1_H\P(B\,|\,\mathcal{H})\,\big|\,\mathcal{H}\big]\Big]=\E\big[\1_H\,\P(B\,|\,\mathcal{H})\,\P(A\,|\,\mathcal{H})\big]\\
&=\E\big[\1_H\,\P(A\cap B\,|\,\mathcal{H})\big]=\E\big[\E[\1_H\1_{A\cap B}\,|\,\mathcal{H}]\big]=\P(A\cap B\cap H).
\end{align*}
而 $\P(A\cap B\cap H)=\E[\1_B;A\cap H]$, 所以 $\P(B\,|\,\mathcal{H})$ 满足 $\P(B\,|\,\sigma(\mathcal{G}_1\cup\mathcal{H}))$ 的定义.
\end{proof}

% ---------------------------------------------------------------------
\subsection{马尔可夫性}

\begin{definition}[马尔可夫过程]\label{def:markov}
设 $X$ 适应于信息流 $\{\mathcal{F}_t\}_{t\in T}$. 若对一切 $s\le t$ ($s,t\in T$) 与一切有界 Borel 函数 $f$,
\[
\E[f(X_t)\,|\,\mathcal{F}_s]=\E[f(X_t)\,|\,X_s],
\]
称 $X$ 是关于 $\{\mathcal{F}_t\}$ 的\textbf{马尔可夫过程}.
\end{definition}

定义只涉及未来\emph{一个}时刻. 下面的定理说明, 它等价于对整个未来成立, 也等价于 ``过去与未来关于现在条件独立''.

\begin{theorem}[马尔可夫性的等价形式]\label{thm:markov-equivalent}
设 $X$ 适应于 $\{\mathcal{F}_t\}$. 对 $s\in T$, 记 $\mathcal{F}^{\ge s}=\sigma(X_u: u\ge s,\ u\in T)$ 为 ``未来'' 的信息. 以下四条等价:
\begin{enumerate}[(a)]
\item $X$ 是关于 $\{\mathcal{F}_t\}$ 的马尔可夫过程;
\item 对任意 $s\le t_1<\cdots<t_n$ 与有界 Borel 函数 $f_1,\dots,f_n$,
$\E\big[\prod_{i=1}^nf_i(X_{t_i})\,\big|\,\mathcal{F}_s\big]=\E\big[\prod_{i=1}^nf_i(X_{t_i})\,\big|\,X_s\big]$;
\item 对任意 $s$ 与有界 $\mathcal{F}^{\ge s}$-可测随机变量 $F$, $\E[F\,|\,\mathcal{F}_s]=\E[F\,|\,X_s]$;
\item 对任意 $s$, ``过去'' $\mathcal{F}_s$ 与 ``未来'' $\mathcal{F}^{\ge s}$ 关于 ``现在'' $\sigma(X_s)$ 条件独立.
\end{enumerate}
\end{theorem}

\begin{proof}
(c)$\Rightarrow$(b)$\Rightarrow$(a) 是显然的 (后者取 $n=1$).

(a)$\Rightarrow$(b): 对 $n$ 归纳, $n=1$ 即 (a). 设结论对 $n-1$ 成立. 由 (a) 与 Doob-Dynkin 引理~\ref{lem:doob-dynkin}, 存在 Borel 函数 $g$ 使 (由于左端的绝对值不超过 $\sup\abs{f_n}$, 可把 $g$ 截断为有界函数) $\E[f_n(X_{t_n})\,|\,\mathcal{F}_{t_{n-1}}]=\E[f_n(X_{t_n})\,|\,X_{t_{n-1}}]=g(X_{t_{n-1}})$. 记 $F=\prod_{i=1}^nf_i(X_{t_i})$. 由于 $\prod_{i<n}f_i(X_{t_i})$ 是 $\mathcal{F}_{t_{n-1}}$-可测的, 由塔性质与提出已知因子,
\begin{align*}
\E[F\,|\,\mathcal{F}_s]&=\E\Big[\prod_{i<n}f_i(X_{t_i})\,\E[f_n(X_{t_n})\,|\,\mathcal{F}_{t_{n-1}}]\,\Big|\,\mathcal{F}_s\Big]\\
&=\E\Big[\prod_{i<n-1}f_i(X_{t_i})\cdot(f_{n-1}g)(X_{t_{n-1}})\,\Big|\,\mathcal{F}_s\Big].
\end{align*}
右端是 $n-1$ 个有界函数之积, 由归纳假设, 它等于把 $\mathcal{F}_s$ 换成 $\sigma(X_s)$ 后的条件期望. 再由同样的计算 (塔性质对 $\sigma(X_s)\subseteq\mathcal{F}_{t_{n-1}}$ 同样适用) 倒推回去,
\[
\E\Big[\prod_{i<n-1}f_i(X_{t_i})\cdot(f_{n-1}g)(X_{t_{n-1}})\,\Big|\,X_s\Big]=\E\Big[\E\big[F\,\big|\,\mathcal{F}_{t_{n-1}}\big]\,\Big|\,X_s\Big]=\E[F\,|\,X_s].
\]

(b)$\Rightarrow$(c): 固定 $s$. 令 $\mathcal{H}$ 为满足 $\E[F\,|\,\mathcal{F}_s]=\E[F\,|\,X_s]$ 的有界随机变量 $F$ 全体. 由条件期望的线性, $\mathcal{H}$ 是含常数的线性空间; 由条件单调收敛定理~\ref{thm:cmct}, 它对有界单调增极限封闭. 令 $\mathcal{P}$ 为一切形如 $\bigcap_{i=1}^n\{X_{t_i}\in B_i\}$ ($s\le t_1<\cdots<t_n$, $B_i$ 为 Borel 集) 的事件全体. 两个这样的事件之交仍是这种形式 (把时刻合并, 同一时刻上的集合取交), 所以 $\mathcal{P}$ 是 $\pi$-系, 且 $\sigma(\mathcal{P})=\mathcal{F}^{\ge s}$. 由 (b) (取 $f_i=\1_{B_i}$), $\mathcal{P}$ 中事件的示性函数属于 $\mathcal{H}$. 由函数形式的单调类定理~\ref{thm:functional-mct} 即得 (c).

(c)$\Leftrightarrow$(d): 由于 $X_s$ 是 $\mathcal{F}_s$-可测的, $\sigma(\mathcal{F}_s\cup\sigma(X_s))=\mathcal{F}_s$. 由命题~\ref{prop:conditional-independence} (取 $\mathcal{G}_1=\mathcal{F}_s$, $\mathcal{G}_2=\mathcal{F}^{\ge s}$, $\mathcal{H}=\sigma(X_s)$), (d) 等价于: 对一切 $B\in\mathcal{F}^{\ge s}$, $\P(B\,|\,\mathcal{F}_s)=\P(B\,|\,X_s)$. 这是 (c) 对示性函数的特例; 反之, 由它出发, 与上一段同样用函数形式的单调类定理 (此时取 $\mathcal{P}=\mathcal{F}^{\ge s}$ 本身) 得到 (c).
\end{proof}

\begin{intuition}
(d) 是马尔可夫性最本质的表述: 只要知道现在的状态 $X_s$, 过去的全部历史就不再提供关于未来的任何额外信息. 过程 ``没有记忆'', 现在的状态概括了过去中与未来有关的一切.
\end{intuition}

\begin{remark}
若 $X$ 是关于某个信息流 $\{\mathcal{F}_t\}$ 的马尔可夫过程, 则它也是关于自然信息流 $\{\mathcal{F}^X_t\}$ 的马尔可夫过程: 由 $\sigma(X_s)\subseteq\mathcal{F}^X_s\subseteq\mathcal{F}_s$ 与塔性质,
$\E[f(X_t)\,|\,\mathcal{F}^X_s]=\E\big[\E[f(X_t)\,|\,\mathcal{F}_s]\,\big|\,\mathcal{F}^X_s\big]=\E\big[\E[f(X_t)\,|\,X_s]\,\big|\,\mathcal{F}^X_s\big]=\E[f(X_t)\,|\,X_s]$. 反之不然: 信息流越大, 马尔可夫性越难成立.
\end{remark}

最重要的一类马尔可夫过程是独立增量过程, 随机游动与布朗运动都属于此类.

\begin{proposition}[独立增量过程是马尔可夫过程]\label{prop:independent-increments}
设 $X$ 适应于 $\{\mathcal{F}_t\}$, 且对一切 $s\le t$, 增量 $X_t-X_s$ 与 $\mathcal{F}_s$ 独立. 则 $X$ 是关于 $\{\mathcal{F}_t\}$ 的马尔可夫过程, 并且对有界 Borel 函数 $f$,
\[
\E[f(X_t)\,|\,\mathcal{F}_s]=\psi_{s,t}(X_s),\qquad \psi_{s,t}(x)=\E\big[f(x+X_t-X_s)\big].
\]
\end{proposition}

\begin{proof}
写 $X_t=X_s+(X_t-X_s)$, 其中 $X_s$ 是 $\mathcal{F}_s$-可测的, $X_t-X_s$ 与 $\mathcal{F}_s$ 独立. 由代入引理~\ref{lem:freezing} (1) (取 $\phi(z,x)=f(x+z)$), $\E[f(X_t)\,|\,\mathcal{F}_s]=\psi_{s,t}(X_s)$. 右端是 $\sigma(X_s)$-可测的, 由塔性质, $\E[f(X_t)\,|\,X_s]=\E\big[\E[f(X_t)\,|\,\mathcal{F}_s]\,\big|\,X_s\big]=\psi_{s,t}(X_s)$. 二者相等.
\end{proof}

\begin{example}[随机游动的马尔可夫性]\label{ex:rw-markov}
设 $S_n=S_0+\sum_{k=1}^n\xi_k$ 如例~\ref{ex:processes} (2), $\mathcal{F}_n=\sigma(S_0,\xi_1,\dots,\xi_n)$. 对 $n\le m$, $S_m-S_n=\xi_{n+1}+\cdots+\xi_m$ 由 $\xi_{n+1},\dots,\xi_m$ 决定, 由分组原理 (推论~\ref{cor:grouping}), 它与 $\mathcal{F}_n$ 独立. 由命题~\ref{prop:independent-increments}, $\{S_n\}$ 是马尔可夫过程.
\end{example}

\begin{example}[不是马尔可夫过程的例子]\label{ex:non-markov}
设 $\xi_0,\xi_1,\dots$ 是 $\iid$ 的, 以等概率取 $\pm1$. 令 $X_n=\xi_n+\xi_{n+1}$ (``滑动和''), 取 $X$ 的自然信息流. $X_n$ 取值于 $\{-2,0,2\}$. 比较事件 $\{X_2=2\}$ 在不同历史信息下的条件概率:
\begin{itemize}
\item 在事件 $\{X_0=2,X_1=0\}$ 上, 必有 $\xi_1=1$, $\xi_2=-1$, 于是 $X_2=-1+\xi_3\le0$. 由于 $\mathcal{F}^X_1=\sigma(X_0,X_1)$ 由有限分割生成, 在这一块 (其概率为 $\P(\xi_0=1,\xi_1=1,\xi_2=-1)=\frac{1}{8}>0$) 上 $\P(X_2=2\,|\,\mathcal{F}^X_1)=0$.
\item 在事件 $\{X_1=0\}$ 上, $(\xi_1,\xi_2)$ 以等概率取 $(1,-1)$ 或 $(-1,1)$, 而 $X_2=2$ 当且仅当 $\xi_2=\xi_3=1$, 所以 $\P(X_2=2\,|\,X_1=0)=\frac{\P(\xi_1=-1,\xi_2=1,\xi_3=1)}{\P(X_1=0)}=\frac{1/8}{1/2}=\frac{1}{4}$.
\end{itemize}
两者在概率为 $\frac{1}{8}$ 的事件上不相等, 所以 $X$ 不是马尔可夫过程: 知道 $X_0$ 能推出关于 $\xi_2$ 的额外信息, 而 $\xi_2$ 会影响 $X_2$. 不过, 二维过程 $(\xi_n,\xi_{n+1})$ 是马尔可夫的. 适当扩大状态空间, 往往可以把非马尔可夫过程变成马尔可夫过程.
\end{example}

% ---------------------------------------------------------------------
\subsection{转移核与 Chapman-Kolmogorov 方程}

马尔可夫过程从时刻 $s$ 的状态 $x$ 出发、在时刻 $t$ 落在 $B$ 中的概率, 由一个 ``转移核'' 描述.

\begin{definition}[转移核]\label{def:kernel}
映射 $P:\R\times\mathcal{B}(\R)\to[0,1]$ 称为\textbf{转移核}, 如果对每个 $x$, $P(x,\cdot)$ 是概率测度, 且对每个 $B$, $x\mapsto P(x,B)$ 是 Borel 函数. 对有界 Borel 函数 $f$, 记 $Pf(x)=\int_{\R}f(y)\,P(x,\dd y)$.
\end{definition}

\begin{definition}[转移概率]\label{def:transition}
设 $X$ 是关于 $\{\mathcal{F}_t\}$ 的马尔可夫过程. 若存在一族转移核 $\{P_{s,t}\}_{s\le t}$, 使得对一切 $s\le t$ 与 $B\in\mathcal{B}(\R)$,
\[
\P(X_t\in B\,|\,\mathcal{F}_s)=P_{s,t}(X_s,B),
\]
称 $\{P_{s,t}\}$ 为 $X$ 的\textbf{转移概率}. 若 $P_{s,t}$ 只依赖于 $t-s$, 记 $P_{s,t}=P_{t-s}$, 称 $X$ 是\textbf{时齐}的.
\end{definition}

由定义, $P_{s,t}(X_s(\omega),\cdot)$ 是 $X_t$ 关于 $\mathcal{F}_s$ 的正则条件分布 (定义~\ref{def:regular-conditional}), 所以由命题~\ref{prop:regular-integral}, 对有界 Borel 函数 $f$,
\begin{equation}\label{eq:kernel-expectation}
\E[f(X_t)\,|\,\mathcal{F}_s]=\int_{\R}f(y)\,P_{s,t}(X_s,\dd y)=P_{s,t}f(X_s).
\end{equation}
例如在命题~\ref{prop:independent-increments} 中, $P_{s,t}(x,B)=\P(x+X_t-X_s\in B)$.

\begin{theorem}[Chapman-Kolmogorov 方程]\label{thm:chapman-kolmogorov}
设 $X$ 是具有转移概率 $\{P_{s,t}\}$ 的马尔可夫过程, $s\le u\le t$, $B\in\mathcal{B}(\R)$. 则对 $\mu_{X_s}$-几乎每个 $x$,
\[
P_{s,t}(x,B)=\int_{\R}P_{u,t}(z,B)\,P_{s,u}(x,\dd z).
\]
\end{theorem}

\begin{proof}
由塔性质 ($\mathcal{F}_s\subseteq\mathcal{F}_u$) 与转移概率的定义,
\[
P_{s,t}(X_s,B)=\P(X_t\in B\,|\,\mathcal{F}_s)=\E\big[\P(X_t\in B\,|\,\mathcal{F}_u)\,\big|\,\mathcal{F}_s\big]=\E\big[P_{u,t}(X_u,B)\,\big|\,\mathcal{F}_s\big].
\]
函数 $h(z)=P_{u,t}(z,B)$ 是有界 Borel 函数, 由 \eqref{eq:kernel-expectation} (把 $t$ 换成 $u$, $f$ 换成 $h$), 右端等于 $\int h(z)\,P_{s,u}(X_s,\dd z)$. 于是两个 Borel 函数 $x\mapsto P_{s,t}(x,B)$ 与 $x\mapsto\int P_{u,t}(z,B)P_{s,u}(x,\dd z)$ 在 $X_s$ 处的值几乎必然相等, 即它们 $\mu_{X_s}$-几乎处处相等.
\end{proof}

\begin{figure}[htbp]
\centering
\begin{tikzpicture}[xscale=1.6,yscale=1.0]
  \draw[->,SADark!80] (0,-0.3) -- (6.4,-0.3) node[right,font=\small] {时间};
  \foreach \t/\l in {0.5/s,3/u,5.5/t} { \draw[SADark!40,dashed] (\t,-0.3) -- (\t,3.6); \node[below,font=\small] at (\t,-0.3) {$\l$}; }
  \fill[SADark] (0.5,1.6) circle (1.6pt) node[left,font=\small] {$x$};
  \draw[SARed,line width=3pt] (5.5,2.3) -- (5.5,3.1); \node[right,SARed,font=\small] at (5.55,2.7) {$B$};
  \foreach \z/\c in {0.4/SABlue!50, 1.3/SABlue!80, 2.2/SABlue, 3.1/SABlue!60} {
    \fill[\c] (3,\z) circle (1.4pt);
    \draw[\c,thick] plot[smooth,tension=0.7] coordinates {(0.5,1.6) (1.4,{1.6+0.3*(\z-1.6)+0.15}) (2.2,{1.6+0.7*(\z-1.6)-0.1}) (3,\z)};
    \draw[\c,thick] plot[smooth,tension=0.7] coordinates {(3,\z) (3.9,{\z+0.35*(2.7-\z)-0.12}) (4.7,{\z+0.75*(2.7-\z)+0.1}) (5.5,{2.7+0.08*(\z-1.6)})};
  }
  \node[SABlue,font=\small,align=center] at (3,-1.35) {中间状态 $z$\\ 按 $P_{s,u}(x,\dd z)$ 分布};
\end{tikzpicture}
\caption{Chapman-Kolmogorov 方程: 从 $x$ 出发经时间 $s\to t$ 到达 $B$ 的概率, 等于对所有可能的中间状态 $z$ (时刻 $u$) 求和: 先从 $x$ 到 $z$, 再从 $z$ 到 $B$. 两段之间的 ``无记忆'' 正是马尔可夫性}
\label{fig:chapman-kolmogorov}
\end{figure}

\begin{remark}
反过来, 给定初始分布 $\nu$ 与一族对\emph{一切} $x$ 满足 Chapman-Kolmogorov 方程的转移核, 令
\begin{multline*}
\P(X_{t_0}\in B_0,\dots,X_{t_n}\in B_n)\\
=\int_{B_0}\nu(\dd x_0)\int_{B_1}P_{t_0,t_1}(x_0,\dd x_1)\cdots\int_{B_n}P_{t_{n-1},t_n}(x_{n-1},\dd x_n),
\end{multline*}
Chapman-Kolmogorov 方程恰好保证这族有限维分布满足相容性条件 (C2), 由 Kolmogorov 延拓定理即得一个以 $\{P_{s,t}\}$ 为转移概率的马尔可夫过程 (见~\cite{kallenberg} 第 11 章). 在时齐情形, Chapman-Kolmogorov 方程写成 $P_{t+s}=P_tP_s$ (算子的复合), 即 $\{P_t\}$ 构成一个半群. 第~6~章将从这一半群出发研究扩散过程的无穷小生成元.
\end{remark}

\begin{example}[有限状态马氏链]\label{ex:markov-chain}
设 $X_n$ 取值于有限集 $\{1,\dots,N\}$, 且 $\P(X_{n+1}=j\,|\,\mathcal{F}_n)=p_{X_n,j}$, 其中 $\mat{P}=(p_{ij})$ 是每行之和为 $1$ 的非负矩阵 (转移矩阵). 此时转移核就是矩阵, $k$ 步转移概率为 $\mat{P}^k$, Chapman-Kolmogorov 方程就是矩阵乘法 $\mat{P}^{m+n}=\mat{P}^m\mat{P}^n$.

以两状态 ``天气'' 链为例 (图~\ref{fig:two-state}): 状态 $1$ (晴) 以概率 $p$ 转为状态 $2$ (雨), 状态 $2$ 以概率 $q$ 转为状态 $1$, $0<p+q<2$. 令
\[
\mat{A}=\frac{1}{p+q}\begin{pmatrix}q&p\\q&p\end{pmatrix},\qquad \mat{B}=\frac{1}{p+q}\begin{pmatrix}p&-p\\-q&q\end{pmatrix}.
\]
直接计算可得 $\mat{A}^2=\mat{A}$, $\mat{B}^2=\mat{B}$, $\mat{A}\mat{B}=\mat{B}\mat{A}=\mat{0}$, 且 $\mat{P}=\begin{psmallmatrix}1-p&p\\q&1-q\end{psmallmatrix}=\mat{A}+(1-p-q)\mat{B}$. 由二项式展开 (交叉项全为零),
\[
\mat{P}^n=\mat{A}+(1-p-q)^n\mat{B}\xrightarrow{n\to\infty}\mat{A}.
\]
极限矩阵 $\mat{A}$ 的每一行都是 $\pi=\big(\frac{q}{p+q},\frac{p}{p+q}\big)$: 无论初始天气如何, 长期来看晴天的比例趋于 $\frac{q}{p+q}$. 容易验证 $\pi\mat{P}=\pi$, 即 $\pi$ 是\textbf{平稳分布}. 收敛速度由 $\abs{1-p-q}^n$ 刻画.
\end{example}

\begin{figure}[htbp]
\centering
\begin{tikzpicture}[
  st/.style={circle,draw=SADark,thick,fill=SABlue!12,minimum size=11mm,font=\small},
  ar/.style={->,thick,SADark!80,>={Stealth[length=2.4mm]}}]
  \node[st] (s1) at (0,0) {$1$ 晴};
  \node[st] (s2) at (4,0) {$2$ 雨};
  \draw[ar] (s1) to[bend left=30] node[above,font=\small] {$p$} (s2);
  \draw[ar] (s2) to[bend left=30] node[below,font=\small] {$q$} (s1);
  \draw[ar] (s1) to[out=150,in=210,looseness=6] node[left,font=\small] {$1-p$} (s1);
  \draw[ar] (s2) to[out=30,in=-30,looseness=6] node[right,font=\small] {$1-q$} (s2);
\end{tikzpicture}
\caption{两状态马氏链的转移图. 每个状态出发的箭头上的概率之和为 $1$}
\label{fig:two-state}
\end{figure}

\begin{example}[正态自回归过程]\label{ex:ar1}
设 $X_0$ 与 $\iid$ 的 $N(0,1)$ 序列 $\{\xi_n\}$ 独立, $X_{n+1}=aX_n+\sigma\xi_{n+1}$. 对 $\mathcal{F}_n=\sigma(X_0,\xi_1,\dots,\xi_n)$, 由代入引理, $X$ 是时齐马尔可夫链, 一步转移核为 $P(x,\cdot)=N(ax,\sigma^2)$. 迭代得 $X_{n+k}=a^kX_n+\sigma\sum_{j=1}^ka^{k-j}\xi_{n+j}$, 所以 $k$ 步转移核为
\[
P_k(x,\cdot)=N\Big(a^kx,\ \sigma^2\sum_{j=0}^{k-1}a^{2j}\Big).
\]
若 $\abs{a}<1$, 当 $k\to\infty$ 时它趋于 $N\big(0,\frac{\sigma^2}{1-a^2}\big)$, 与起点 $x$ 无关, 这是平稳分布. 这一过程是第~6~章 Ornstein-Uhlenbeck 过程的离散版本.
\end{example}

\begin{example}[热核]\label{ex:heat-kernel}
对 $t>0$, 令
\[
p_t(x,y)=\frac{1}{\sqrt{2\pi t}}\exp\Big(-\frac{(y-x)^2}{2t}\Big),\qquad P_t(x,B)=\int_Bp_t(x,y)\dd y,
\]
即 $P_t(x,\cdot)=N(x,t)$, 并令 $P_0(x,\cdot)=\delta_x$. 它对一切 $x$ 满足 Chapman-Kolmogorov 方程 $P_{s+t}=P_sP_t$: $\int P_t(z,B)P_s(x,\dd z)$ 是 $x+Z_1+Z_2\in B$ 的概率, 其中 $Z_1\sim N(0,s)$ 与 $Z_2\sim N(0,t)$ 独立 (先按 $Z_1$ 取条件, 再用代入引理), 由例~\ref{ex:indep-sums} (2), $x+Z_1+Z_2\sim N(x,s+t)$. 用密度写出来, 就是
\[
\int_{\R}p_s(x,z)\,p_t(z,y)\dd z=p_{s+t}(x,y).
\]
第~4~章将看到, 布朗运动正是以 $\{P_t\}$ 为转移概率的马尔可夫过程; $p_t(x,y)$ 满足热传导方程 $\partial_tp=\frac{1}{2}\partial_{yy}p$, 这是概率与偏微分方程联系的起点 (第~6~章).
\end{example}

% ---------------------------------------------------------------------
\subsection*{本章小结}
\addcontentsline{toc}{subsection}{本章小结}

本章把条件期望定义为 ``在给定信息下的最佳预测'', 并建立了它的运算法则. 表~\ref{tab:ch2-summary} 汇总了最常用的法则; 它们将在后续各章中反复使用, 通常不再注明出处.

\begin{center}
\captionof{table}{条件期望的常用法则 (均在几乎必然意义下成立)}
\label{tab:ch2-summary}
\renewcommand{\arraystretch}{1.5}
\begin{tabular}{>{\raggedright\arraybackslash}p{0.24\textwidth} >{\raggedright\arraybackslash}p{0.46\textwidth} >{\raggedright\arraybackslash}p{0.20\textwidth}}
\toprule
法则 & 内容 & 出处 \\
\midrule
全期望公式 & $\E[\E[X\,|\,\mathcal{G}]]=\E X$ & 定理~\ref{thm:ce-basic} \\
已知量不变 & $X$ 是 $\mathcal{G}$-可测的 $\Rightarrow\E[X\,|\,\mathcal{G}]=X$ & 定理~\ref{thm:ce-basic} \\
独立性化简 & $X$ 与 $\mathcal{G}$ 独立 $\Rightarrow\E[X\,|\,\mathcal{G}]=\E X$ & 定理~\ref{thm:ce-basic} \\
塔性质 & $\mathcal{H}\subseteq\mathcal{G}\Rightarrow\E[\E[X\,|\,\mathcal{G}]\,|\,\mathcal{H}]=\E[X\,|\,\mathcal{H}]$ & 定理~\ref{thm:tower} \\
提出已知因子 & $Z$ 是 $\mathcal{G}$-可测的 $\Rightarrow\E[ZX\,|\,\mathcal{G}]=Z\E[X\,|\,\mathcal{G}]$ & 定理~\ref{thm:take-out} \\
条件 MCT / Fatou / DCT & 极限与条件期望交换 & 定理~\ref{thm:cmct}--\ref{thm:cdct} \\
条件 Jensen & $\phi(\E[X\,|\,\mathcal{G}])\le\E[\phi(X)\,|\,\mathcal{G}]$ & 定理~\ref{thm:cjensen} \\
代入引理 & $\E[\phi(\rvec{X},\rvec{Y})\,|\,\mathcal{G}]=\E[\phi(\rvec{X},\rvec{y})]|_{\rvec{y}=\rvec{Y}}$ & 引理~\ref{lem:freezing} \\
正态条件分布 & 条件均值线性, 条件协方差为 Schur 补 & 定理~\ref{thm:gaussian-conditional} \\
\bottomrule
\end{tabular}
\end{center}

% =====================================================================
%  第 3 章  鞅
% =====================================================================
\chapter{鞅}
\label{ch:3}

``鞅'' (martingale) 一词原指 18 世纪法国流行的一种赌博策略: 输了就加倍下注, 直到赢一次为止. 数学上的鞅刻画的是\emph{公平}的赌博, 无论过去发生了什么, 下一局的期望输赢都是零. 这个看似朴素的条件蕴含了很强的结论. 公平赌博中, 任何下注策略、任何 (合理的) 离场时机都无法带来期望收益 (3.1--3.2 节); 鞅的最大值可以用终值控制 (3.3 节); 有界的鞅必定收敛 (3.4 节). 这些结论构成了现代概率论最有力的工具箱. 在随机分析中, 布朗运动是连续时间鞅的原型, Itô 积分是鞅变换的连续版本, 而随机微积分的全部理论都建立在连续鞅之上 (3.5 节).

本章 3.1--3.4 节讨论离散时间, 总在带流的概率空间 $(\Omega,\mathcal{F},\{\mathcal{F}_n\}_{n\ge0},\P)$ 上进行 (定义~\ref{def:filtration}), 记 $\mathcal{F}_\infty=\sigma\big(\bigcup_n\mathcal{F}_n\big)$. 对过程 $X$, 记增量 $\Delta X_n=X_n-X_{n-1}$ ($n\ge1$).

% =====================================================================
\section{鞅的定义与经典例子}
\label{sec:3-1}

% ---------------------------------------------------------------------
\subsection{定义与基本性质}

\begin{definition}[鞅、下鞅、上鞅]\label{def:martingale}
设 $X=\{X_n\}_{n\ge0}$ 适应于 $\{\mathcal{F}_n\}$, 且每个 $X_n$ 可积. 若对一切 $n\ge0$,
\begin{enumerate}
\item $\E[X_{n+1}\,|\,\mathcal{F}_n]=X_n$, 称 $X$ 为\textbf{鞅};
\item $\E[X_{n+1}\,|\,\mathcal{F}_n]\ge X_n$, 称 $X$ 为\textbf{下鞅};
\item $\E[X_{n+1}\,|\,\mathcal{F}_n]\le X_n$, 称 $X$ 为\textbf{上鞅}.
\end{enumerate}
需要强调信息流时, 称 $X$ 为 $\{\mathcal{F}_n\}$-鞅.
\end{definition}

把 $X_n$ 想象成赌徒在第 $n$ 局后的财富: 鞅是公平赌博, 下鞅是有利的赌博, 上鞅是不利的赌博 (例如赌场里的轮盘赌对赌客是上鞅). 名称中的 ``上''、``下'' 与直觉相反, 它来自位势理论: 下鞅对应下调和函数. $X$ 是上鞅当且仅当 $-X$ 是下鞅, 所以关于下鞅的结论都可以平行地转换为上鞅的结论.

\begin{proposition}\label{prop:martingale-basic}
\begin{enumerate}
\item $X$ 是鞅, 当且仅当对一切 $m\ge n$, $\E[X_m\,|\,\mathcal{F}_n]=X_n$; 下鞅 (上鞅) 的情形把等号换成 $\ge$ ($\le$).
\item 鞅的期望不变: $\E X_n=\E X_0$; 下鞅的期望单调不减.
\item 若 $X$ 是 $\{\mathcal{F}_n\}$-鞅, 则它也是关于自然信息流 $\{\mathcal{F}^X_n\}$ 的鞅.
\end{enumerate}
\end{proposition}

\begin{proof}
(1) 对 $m-n$ 归纳. 设 $\E[X_{m}\,|\,\mathcal{F}_n]=X_n$, 则由塔性质~\ref{thm:tower},
$\E[X_{m+1}\,|\,\mathcal{F}_n]=\E\big[\E[X_{m+1}\,|\,\mathcal{F}_m]\,\big|\,\mathcal{F}_n\big]=\E[X_m\,|\,\mathcal{F}_n]=X_n$. 下鞅情形在第二个等号处换成 $\ge$, 并用条件期望的单调性.
(2) 对 (1) 取期望并用全期望公式.
(3) $\mathcal{F}^X_n\subseteq\mathcal{F}_n$ (因为 $X$ 适应), 由塔性质 $\E[X_{n+1}\,|\,\mathcal{F}^X_n]=\E\big[\E[X_{n+1}\,|\,\mathcal{F}_n]\,\big|\,\mathcal{F}^X_n\big]=\E[X_n\,|\,\mathcal{F}^X_n]=X_n$.
\end{proof}

\begin{proposition}[鞅的凸函数是下鞅]\label{prop:convex-submartingale}
\begin{enumerate}
\item 若 $X$ 是鞅, $\phi$ 是凸函数, 且每个 $\phi(X_n)$ 可积, 则 $\phi(X)$ 是下鞅.
\item 若 $X$ 是下鞅, $\phi$ 是单调不减的凸函数, 且每个 $\phi(X_n)$ 可积, 则 $\phi(X)$ 是下鞅.
\end{enumerate}
特别地, 若 $X$ 是鞅, 则 $\abs{X}$、$X^+$ 以及 (当 $X_n\in L^p$ 时) $\abs{X}^p$ ($p\ge1$) 都是下鞅; 若 $X$ 是下鞅, 则 $X^+$ 与 $(X-a)^+$ 是下鞅.
\end{proposition}

\begin{proof}
由条件 Jensen 不等式~\ref{thm:cjensen}, $\E[\phi(X_{n+1})\,|\,\mathcal{F}_n]\ge\phi(\E[X_{n+1}\,|\,\mathcal{F}_n])$. 在 (1) 中右端等于 $\phi(X_n)$; 在 (2) 中 $\E[X_{n+1}\,|\,\mathcal{F}_n]\ge X_n$, 由 $\phi$ 单调不减, 右端 $\ge\phi(X_n)$.
\end{proof}

% ---------------------------------------------------------------------
\subsection{经典例子}

\begin{example}[独立随机变量的和与积]\label{ex:sum-martingales}
设 $\xi_1,\xi_2,\dots$ 独立, $\mathcal{F}_0=\{\varnothing,\Omega\}$, $\mathcal{F}_n=\sigma(\xi_1,\dots,\xi_n)$, $S_0=0$, $S_n=\sum_{k=1}^n\xi_k$. 由于 $\xi_{n+1}$ 与 $\mathcal{F}_n$ 独立 (分组原理~\ref{cor:grouping}), 而 $S_n$ 是 $\mathcal{F}_n$-可测的, 由条件期望的基本性质~\ref{thm:ce-basic} 与提出已知因子~\ref{thm:take-out}, 以下三个过程都是鞅:
\begin{enumerate}
\item \textbf{和.} 若 $\E\xi_k=0$, 则 $S_n$ 是鞅: $\E[S_{n+1}\,|\,\mathcal{F}_n]=S_n+\E[\xi_{n+1}]=S_n$. 一般地, 若 $\E\xi_k=\mu_k$, 则 $S_n-\sum_{k\le n}\mu_k$ 是鞅; 若各 $\mu_k\ge0$, 则 $S_n$ 是下鞅.
\item \textbf{平方补偿.} 若 $\E\xi_k=0$, $\Var(\xi_k)=\sigma_k^2<\infty$, 则 $S_n^2-\sum_{k\le n}\sigma_k^2$ 是鞅:
\[
\E[S_{n+1}^2\,|\,\mathcal{F}_n]=S_n^2+2S_n\E[\xi_{n+1}]+\E[\xi_{n+1}^2]=S_n^2+\sigma_{n+1}^2.
\]
\item \textbf{指数鞅.} 若 $\xi_k$ 同分布, $\varphi(\theta)=\E[e^{\theta\xi_1}]<\infty$, 则 $Z_n=e^{\theta S_n}/\varphi(\theta)^n$ 是鞅:
$\E[Z_{n+1}\,|\,\mathcal{F}_n]=Z_n\,\E[e^{\theta\xi_{n+1}}]/\varphi(\theta)=Z_n$. 更一般地, 若 $Y_k\ge0$ 独立且 $\E Y_k=1$, 则乘积 $\prod_{k\le n}Y_k$ 是鞅.
\end{enumerate}
第~4~章将看到, 对布朗运动 $B_t$, 三个对应的鞅分别是 $B_t$、$B_t^2-t$ 与 $e^{\theta B_t-\theta^2t/2}$.
\end{example}

\begin{example}[Doob 鞅]\label{ex:doob-martingale}
设 $X\in L^1$, 令 $M_n=\E[X\,|\,\mathcal{F}_n]$. 由塔性质, $\E[M_{n+1}\,|\,\mathcal{F}_n]=\E\big[\E[X\,|\,\mathcal{F}_{n+1}]\,\big|\,\mathcal{F}_n\big]=M_n$, 所以 $M$ 是鞅. $M_n$ 是在时刻 $n$ 对 $X$ 的最佳预测, 随着信息增加不断修正. 3.4 节将证明, 每个一致可积鞅都是这种形式, 并且 $M_n\to\E[X\,|\,\mathcal{F}_\infty]$.
\end{example}

\begin{example}[折现资产价格]\label{ex:binomial-model}
在二叉树模型中, 股票价格每期以因子 $u$ 或 $d$ 变化 ($0<d<1+r<u$, $r$ 为无风险利率): $S_{n+1}=S_n\eta_{n+1}$, 其中 $\eta_1,\eta_2,\dots$ 独立, 以概率 $p^*$ 取 $u$、以概率 $1-p^*$ 取 $d$. 取
\[
p^*=\frac{1+r-d}{u-d}\in(0,1),
\]
则 $\E[\eta_{n+1}]=p^*u+(1-p^*)d=1+r$, 从而折现价格 $\tilde S_n=S_n/(1+r)^n$ 满足 $\E[\tilde S_{n+1}\,|\,\mathcal{F}_n]=\tilde S_n\E[\eta_{n+1}]/(1+r)=\tilde S_n$, 是鞅. 使折现价格成为鞅的概率测度称为\textbf{风险中性测度}. 金融数学的基本定理断言: 在无套利市场中, 衍生品的价格等于其折现收益在风险中性测度下的期望. 第~5~章的 Girsanov 定理正是在连续时间中构造这种测度的工具.
\end{example}

\begin{example}[Galton-Watson 分支过程]\label{ex:galton-watson}
设 $\{\xi_{n,i}\}_{n,i\ge1}$ 是 $\iid$ 的非负整数值随机变量 (第 $n$ 代第 $i$ 个个体的后代数), $\mu=\E\xi_{1,1}\in(0,\infty)$, 母函数 $f(s)=\E[s^{\xi_{1,1}}]$ ($s\in[0,1]$). 令 $Z_0=1$,
\[
Z_{n+1}=\sum_{i=1}^{Z_n}\xi_{n+1,i}\qquad(\text{约定空和为 }0),
\]
$\mathcal{F}_n=\sigma(\xi_{k,i}: k\le n,\ i\ge1)$. $Z_n$ 是第 $n$ 代的个体数 (图~\ref{fig:galton-watson}).
\begin{enumerate}
\item $W_n=Z_n/\mu^n$ 是鞅. 事实上, $\sigma(Z_n)$ 由可数分割 $\{Z_n=k\}$ 生成, 而 $\{\xi_{n+1,i}\}_i$ 与 $\mathcal{F}_n$ 独立. 对 $A\in\mathcal{F}_n$, 由于在 $A\cap\{Z_n=k\}$ 上 $Z_{n+1}=\sum_{i=1}^k\xi_{n+1,i}$,
\begin{align*}
\E[Z_{n+1};A]&=\sum_{k\ge0}\E\Big[\sum_{i=1}^k\xi_{n+1,i};A\cap\{Z_n=k\}\Big]\\
&=\sum_{k\ge0}k\mu\,\P(A\cap\{Z_n=k\})=\E[\mu Z_n;A]
\end{align*}
(第一个等号由单调收敛定理, 第二个等号由独立性). 所以 $\E[Z_{n+1}\,|\,\mathcal{F}_n]=\mu Z_n$.
\item 若 $q\in[0,1]$ 满足 $f(q)=q$, 则 $q^{Z_n}$ 是鞅. 同理, 在 $\{Z_n=k\}$ 上, $\E[q^{Z_{n+1}}\,|\,\mathcal{F}_n]=\E[q^{\sum_{i\le k}\xi_{n+1,i}}]=f(q)^k=q^k$.
\end{enumerate}
3.4 节将用这两个鞅研究种群的灭绝: 当 $\mu\le1$ 时种群几乎必然灭绝, 当 $\mu>1$ 时灭绝概率是方程 $f(q)=q$ 在 $[0,1)$ 中的根.
\end{example}

\begin{figure}[htbp]
\centering
\begin{tikzpicture}[
  ind/.style={circle,draw=SADark,fill=SABlue!20,inner sep=0pt,minimum size=3.2mm},
  dead/.style={circle,draw=SADark!50,fill=white,inner sep=0pt,minimum size=3.2mm},
  ed/.style={SADark!60}]
  \node[ind] (a) at (0,0) {};
  \node[ind] (b1) at (-3,-1.2) {}; \node[ind] (b2) at (0,-1.2) {}; \node[ind] (b3) at (3,-1.2) {};
  \foreach \c in {b1,b2,b3} \draw[ed] (a) -- (\c);
  \node[ind] (c1) at (-3.6,-2.4) {}; \node[ind] (c2) at (-2.4,-2.4) {};
  \node[ind] (c3) at (2.2,-2.4) {}; \node[ind] (c4) at (3.0,-2.4) {}; \node[ind] (c5) at (3.8,-2.4) {};
  \draw[ed] (b1) -- (c1); \draw[ed] (b1) -- (c2);
  \draw[ed] (b3) -- (c3); \draw[ed] (b3) -- (c4); \draw[ed] (b3) -- (c5);
  \node[font=\scriptsize,SADark!60] at (0,-1.65) {无后代};
  \node[ind] (d1) at (-3.6,-3.6) {};
  \node[ind] (d2) at (2.0,-3.6) {}; \node[ind] (d3) at (2.6,-3.6) {};
  \node[ind] (d4) at (3.8,-3.6) {};
  \draw[ed] (c1) -- (d1); \draw[ed] (c3) -- (d2); \draw[ed] (c3) -- (d3); \draw[ed] (c5) -- (d4);
  \foreach \y/\l in {0/{$Z_0=1$},-1.2/{$Z_1=3$},-2.4/{$Z_2=5$},-3.6/{$Z_3=4$}} { \node[anchor=west,font=\small] at (5,\y) {\l}; }
\end{tikzpicture}
\caption{Galton-Watson 分支过程的一个实现. 每个个体独立地产生随机个后代, $Z_n$ 是第 $n$ 代的个体数}
\label{fig:galton-watson}
\end{figure}

\begin{example}[P\'olya 罐子]\label{ex:polya}
罐中起初有一个红球和一个蓝球. 每次随机取出一球, 把它连同一个同色的新球放回. 设 $R_n$ 为 $n$ 次操作后的红球数, 则罐中共有 $n+2$ 个球, 令 $X_n=\frac{R_n}{n+2}$ 为红球比例. 严格地, 设 $U_1,U_2,\dots$ 是 $\iid$ 的 $\mathrm{U}(0,1)$ 随机变量, 第 $n+1$ 次取出红球当且仅当 $U_{n+1}\le X_n$, 取 $\mathcal{F}_n=\sigma(U_1,\dots,U_n)$. 则 $R_{n+1}=R_n+\1_{\{U_{n+1}\le X_n\}}$. 由于 $U_{n+1}$ 与 $\mathcal{F}_n$ 独立而 $X_n$ 是 $\mathcal{F}_n$-可测的, 由代入引理~\ref{lem:freezing}, $\E[\1_{\{U_{n+1}\le X_n\}}\,|\,\mathcal{F}_n]=X_n$. 于是
\[
\E[X_{n+1}\,|\,\mathcal{F}_n]=\frac{R_n+X_n}{n+3}=\frac{R_n+\frac{R_n}{n+2}}{n+3}=\frac{R_n}{n+2}=X_n,
\]
红球比例是鞅. 直观上, 红球多时更可能再抽到红球 (``富者愈富''), 但比例的增量恰好被总数的增长抵消. 3.4 节将证明 $X_n$ 几乎必然收敛, 且极限服从 $\mathrm{U}(0,1)$.
\end{example}

\begin{example}[似然比鞅]\label{ex:likelihood-ratio}
设 $Y_1,Y_2,\dots$ 是 $\iid$ 的, 密度为 $f>0$; $g$ 是另一个概率密度. 令 $\mathcal{F}_n=\sigma(Y_1,\dots,Y_n)$,
\[
L_n=\prod_{k=1}^n\frac{g(Y_k)}{f(Y_k)}.
\]
由于 $\E\big[\frac{g(Y_k)}{f(Y_k)}\big]=\int\frac{g(y)}{f(y)}f(y)\dd y=\int g(y)\dd y=1$, 由例~\ref{ex:sum-martingales} (3), $L_n$ 是鞅. 统计学上, $L_n$ 是检验 ``数据来自 $g$'' 对 ``数据来自 $f$'' 的似然比. 若在 $(\Omega,\mathcal{F}_n)$ 上定义新测度 $\mathbb{Q}(A)=\E[L_n;A]$, 则在 $\mathbb{Q}$ 下 $Y_1,\dots,Y_n$ 是 $\iid$ 的、密度为 $g$ (对乘积集合直接验证). 因此 $L_n$ 就是 Radon-Nikodym 导数 $\frac{\dd\mathbb{Q}}{\dd\P}\big|_{\mathcal{F}_n}$. 这是第~5~章 Girsanov 定理中 ``测度变换'' 的离散原型. 3.4 节将证明: 若 $g\ne f$, 则在 $\P$ 下 $L_n\to0$ 几乎必然成立, 尽管 $\E L_n\equiv1$.
\end{example}

% ---------------------------------------------------------------------
\subsection{鞅变换: 公平赌博中的策略}

在例~\ref{ex:betting} 中我们引入了可料的赌博策略. 现在可以精确地回答: 策略能否改变赌博的公平性?

\begin{definition}[鞅变换]\label{def:martingale-transform}
设 $X$ 适应, $H=\{H_n\}_{n\ge1}$ 可料. 定义
\[
(H\cdot X)_0=0,\qquad (H\cdot X)_n=\sum_{k=1}^nH_k(X_k-X_{k-1})=\sum_{k=1}^nH_k\,\Delta X_k,\quad n\ge1.
\]
称 $H\cdot X$ 为 $X$ 关于 $H$ 的\textbf{鞅变换} (或离散随机积分).
\end{definition}

若 $X_n$ 是每单位赌注在第 $n$ 局后的累计输赢, $H_k$ 是第 $k$ 局的赌注, 则 $(H\cdot X)_n$ 就是按策略 $H$ 下注的累计收益. 记号 $H\cdot X$ 有意模仿积分 $\int H\dd X$: 第~5~章的 Itô 积分正是它的连续时间极限.

\begin{theorem}[鞅变换定理]\label{thm:martingale-transform}
设 $H$ 可料, 且每个 $H_n$ 有界.
\begin{enumerate}
\item 若 $X$ 是鞅, 则 $H\cdot X$ 是鞅.
\item 若 $X$ 是下鞅 (上鞅), 且 $H_n\ge0$, 则 $H\cdot X$ 是下鞅 (上鞅).
\end{enumerate}
若 $X$ 是平方可积鞅, 则把 ``$H_n$ 有界'' 换成 ``$H_n\in L^2$'' 时 (1) 仍成立.
\end{theorem}

\begin{proof}
$(H\cdot X)_n$ 是 $\mathcal{F}_n$-可测的有限和. 若 $\abs{H_k}\le c_k$, 则 $\abs{(H\cdot X)_n}\le\sum_{k\le n}c_k(\abs{X_k}+\abs{X_{k-1}})$, 故可积. 由于 $H_{n+1}$ 是 $\mathcal{F}_n$-可测的, 由提出已知因子 (定理~\ref{thm:take-out}),
\[
\E\big[(H\cdot X)_{n+1}-(H\cdot X)_n\,\big|\,\mathcal{F}_n\big]=\E[H_{n+1}\Delta X_{n+1}\,|\,\mathcal{F}_n]=H_{n+1}\E[\Delta X_{n+1}\,|\,\mathcal{F}_n].
\]
在 (1) 中 $\E[\Delta X_{n+1}\,|\,\mathcal{F}_n]=0$; 在 (2) 中它 $\ge0$ ($\le0$), 乘以 $H_{n+1}\ge0$ 后符号不变. 平方可积情形: 由 Cauchy-Schwarz 不等式, $H_k\Delta X_k\in L^1$, 其余论证相同.
\end{proof}

\begin{corollary}
若 $X$ 是鞅, $H$ 可料且有界, 则 $\E[(H\cdot X)_n]=0$ 对一切 $n$ 成立: 在公平赌博中, 任何有界的可料策略在任何固定时刻的期望收益都是零.
\end{corollary}

\begin{intuition}
定理~\ref{thm:martingale-transform} 的关键是可料性: 赌注 $H_{n+1}$ 必须在第 $n+1$ 局开始前确定. 若允许 ``偷看'' 结果再下注 ($H_{n+1}=\sgn(\Delta X_{n+1})$), 显然稳赚不赔. 第~5~章 Itô 积分的定义要求被积过程在每个小区间的左端点取值, 原因完全相同.
\end{intuition}

\begin{example}[加倍策略]\label{ex:doubling}
设 $\xi_k$ 是 $\iid$ 的, 以等概率取 $\pm1$, $X_n=\sum_{k\le n}\xi_k$ 是鞅. 加倍策略为: $H_1=1$; 若前 $n-1$ 局全输, 则 $H_n=2^{n-1}$; 一旦赢了一局, 此后 $H_n=0$. $H$ 是可料的, 每个 $H_n$ 有界 (至多 $2^{n-1}$). 若前 $k-1$ 局全输而第 $k$ 局赢, 则累计收益为 $-(1+2+\cdots+2^{k-2})+2^{k-1}=1$. 由于 ``永远不赢'' 的概率为 $\lim_n2^{-n}=0$, 收益几乎必然最终等于 $1$. 但由鞅变换定理, $\E[(H\cdot X)_n]=0$ 对每个 $n$ 成立. 前 $n$ 局全输的概率是 $2^{-n}$, 此时损失为 $2^n-1$, 恰好抵消了其余情形下的收益 $1$, 即
\[
1\cdot(1-2^{-n})+(-(2^n-1))\cdot2^{-n}=0.
\]
加倍策略在任何固定的时刻都不占便宜; 它 ``几乎必然赢'' 的代价是需要无限的本金和无限的时间. 3.2 节的可选停时定理将精确地说明这一点.
\end{example}

\begin{figure}[htbp]
\centering
\begin{tikzpicture}
\begin{axis}[SAaxis, width=10cm, height=5.6cm, xmin=0, xmax=7.5, ymin=-16, ymax=4,
  xtick={1,2,3,4,5,6,7}, ytick={-15,-7,-3,-1,1}, xlabel={$n$}, ylabel={$(H\cdot X)_n$}]
  \addplot[SABlue,very thick,mark=*,mark size=1.5pt] coordinates {(0,0) (1,-1) (2,-3) (3,-7) (4,1) (5,1) (6,1) (7,1)};
  \addplot[SARed,thick,dashed,mark=o,mark size=1.5pt] coordinates {(0,0) (1,-1) (2,-3) (3,-7) (4,-15)};
  \node[SABlue,font=\small,anchor=south west] at (axis cs:4.1,1.2) {第 4 局赢: 收益为 $1$};
  \node[SARed,font=\small,anchor=north west] at (axis cs:4.1,-14) {若再输: 损失 $2^n-1$};
\end{axis}
\end{tikzpicture}
\caption{加倍策略的收益轨道. 每次赢后收益恰为 $1$, 但在赢之前, 损失以指数速度增长}
\label{fig:doubling}
\end{figure}

% =====================================================================
\section{停时与可选停时定理}
\label{sec:3-2}

鞅变换定理说: 公平赌博中, 调整赌注无法获利. 另一个自然的想法是选择\emph{何时离场}: 比如 ``领先时就收手''. 离场时刻依赖于赌局的进展, 是一个随机时刻; 但它只能依据已经发生的事情来决定, 而不能预知未来. 这样的随机时刻称为停时.

% ---------------------------------------------------------------------
\subsection{停时}

\begin{definition}[停时]\label{def:stopping-time}
取值于 $\{0,1,2,\dots\}\cup\{\infty\}$ 的随机变量 $\tau$ 称为 (关于 $\{\mathcal{F}_n\}$ 的) \textbf{停时}, 如果对每个 $n\ge0$, $\{\tau\le n\}\in\mathcal{F}_n$.
\end{definition}

由于 $\{\tau=n\}=\{\tau\le n\}\setminus\{\tau\le n-1\}$ 且 $\{\tau\le n\}=\bigcup_{k\le n}\{\tau=k\}$, 条件 $\{\tau\le n\}\in\mathcal{F}_n$ 等价于对每个 $n$, $\{\tau=n\}\in\mathcal{F}_n$: 在时刻 $n$, 根据已有的信息就能判断 ``是否现在停止''.

\begin{example}[首达时]\label{ex:hitting-time}
设 $X$ 适应, $B$ 是 Borel 集. 令 $\tau_B=\inf\{n\ge0: X_n\in B\}$ (约定 $\inf\varnothing=\infty$) 为 $X$ 首次进入 $B$ 的时刻. 由于
\[
\{\tau_B\le n\}=\bigcup_{k=0}^n\{X_k\in B\}\in\mathcal{F}_n,
\]
$\tau_B$ 是停时 (图~\ref{fig:hitting-time}). 相反, $X$ 在时刻 $N$ 之前\emph{最后一次}位于 $B$ 中的时刻 $\lambda=\sup\{n\le N: X_n\in B\}$ 一般不是停时: 判断 $\{\lambda=n\}$ 需要知道 $X_{n+1},\dots,X_N$ 都不在 $B$ 中, 这依赖于未来.
\end{example}

\begin{figure}[htbp]
\centering
\begin{tikzpicture}
\begin{axis}[SAaxis, axis x line=bottom, axis y line=left, width=11.5cm, height=5.4cm, xmin=0, xmax=41, ymin=-3, ymax=15,
  xtick={10,20,30,40}, ytick={5,10}, xlabel={$n$}, ylabel={$X_n$}]
  \addplot[SADark!20,forget plot] coordinates {(0,5) (41,5)};
  \addplot[SABlue!50,thick,mark=*,mark size=0.8pt,forget plot] coordinates {(0,0) (1,-1) (2,-2) (3,-1) (4,-2) (5,-1) (6,0) (7,1) (8,2) (9,1) (10,0) (11,1) (12,0) (13,1) (14,2) (15,3) (16,4) (17,3) (18,2) (19,3) (20,2) (21,3) (22,4) (23,5) (24,4) (25,5) (26,6) (27,7) (28,8) (29,9) (30,10) (31,11) (32,12) (33,11) (34,10) (35,11) (36,12) (37,13) (38,14) (39,13) (40,12)};
  \addplot[SABlue,very thick,mark=*,mark size=1.1pt] coordinates {(0,0) (1,-1) (2,-2) (3,-1) (4,-2) (5,-1) (6,0) (7,1) (8,2) (9,1) (10,0) (11,1) (12,0) (13,1) (14,2) (15,3) (16,4) (17,3) (18,2) (19,3) (20,2) (21,3) (22,4) (23,5)};
  \addplot[SARed,dashed] coordinates {(23,-3) (23,5)};
  \fill[SARed] (axis cs:23,5) circle (2.2pt);
  \node[SARed,font=\small,anchor=south west] at (axis cs:23.5,-2.8) {$\tau=\inf\{n: X_n\ge5\}=23$};
  \node[font=\small,anchor=south west] at (axis cs:0.5,5.2) {$5$};
\end{axis}
\end{tikzpicture}
\caption{首达时: 在时刻 $n$, 只需查看 $X_0,\dots,X_n$ 就能判断 $\tau\le n$ 是否发生. 粗线是停止过程 $X_{\tau\wedge n}$ 在 $\tau$ 之前的部分, 此后它保持在 $X_\tau=5$}
\label{fig:hitting-time}
\end{figure}

\begin{proposition}[停时的运算]\label{prop:stopping-operations}
设 $\sigma,\tau$ 是停时, $m\ge0$ 是整数. 则 $\sigma\wedge\tau$、$\sigma\vee\tau$、$\sigma+\tau$、$\tau+m$ 都是停时.
\end{proposition}

\begin{proof}
$\{\sigma\wedge\tau\le n\}=\{\sigma\le n\}\cup\{\tau\le n\}$, $\{\sigma\vee\tau\le n\}=\{\sigma\le n\}\cap\{\tau\le n\}$. 对和,
$\{\sigma+\tau=n\}=\bigcup_{k=0}^n\{\sigma=k\}\cap\{\tau=n-k\}$, 其中 $\{\sigma=k\}\in\mathcal{F}_k\subseteq\mathcal{F}_n$, $\{\tau=n-k\}\in\mathcal{F}_{n-k}\subseteq\mathcal{F}_n$. 常数 $m$ 显然是停时.
\end{proof}

注意 $\tau-1$ 一般不是停时: $\{\tau-1=n\}=\{\tau=n+1\}$ 需要时刻 $n+1$ 的信息.

\begin{definition}[停时前的信息]\label{def:F-tau}
设 $\tau$ 是停时. 定义
\[
\mathcal{F}_\tau=\big\{A\in\mathcal{F}_\infty:\ A\cap\{\tau\le n\}\in\mathcal{F}_n\text{ 对一切 }n\ge0\text{ 成立}\big\}.
\]
\end{definition}

$A\in\mathcal{F}_\tau$ 是指在 $\tau$ 发生的那一刻, 已经能够判断 $A$ 是否发生. 若 $\tau\equiv n$, 则 $\mathcal{F}_\tau=\mathcal{F}_n$.

\begin{proposition}[$\mathcal{F}_\tau$ 的性质]\label{prop:F-tau}
设 $\sigma,\tau$ 是停时.
\begin{enumerate}
\item $\mathcal{F}_\tau$ 是 $\sigma$-代数, 且 $\tau$ 是 $\mathcal{F}_\tau$-可测的.
\item 若 $\sigma\le\tau$, 则 $\mathcal{F}_\sigma\subseteq\mathcal{F}_\tau$.
\item 若 $X$ 适应, 则 $X_\tau\1_{\{\tau<\infty\}}$ 是 $\mathcal{F}_\tau$-可测的.
\end{enumerate}
\end{proposition}

\begin{proof}
(1) $\Omega\cap\{\tau\le n\}\in\mathcal{F}_n$. 若 $A\in\mathcal{F}_\tau$, 则 $A^c\cap\{\tau\le n\}=\{\tau\le n\}\setminus(A\cap\{\tau\le n\})\in\mathcal{F}_n$. 若 $A_k\in\mathcal{F}_\tau$, 则 $\big(\bigcup_kA_k\big)\cap\{\tau\le n\}=\bigcup_k(A_k\cap\{\tau\le n\})\in\mathcal{F}_n$. 对 $\tau$ 的可测性: $\{\tau\le m\}\cap\{\tau\le n\}=\{\tau\le m\wedge n\}\in\mathcal{F}_n$.

(2) 若 $A\in\mathcal{F}_\sigma$, 由 $\sigma\le\tau$ 得 $\{\tau\le n\}\subseteq\{\sigma\le n\}$, 于是 $A\cap\{\tau\le n\}=\big(A\cap\{\sigma\le n\}\big)\cap\{\tau\le n\}\in\mathcal{F}_n$.

(3) 对 Borel 集 $B$, $\{X_\tau\1_{\{\tau<\infty\}}\in B\}$ 在 $\{\tau<\infty\}$ 上的部分与 $\{\tau\le n\}$ 之交为 $\bigcup_{k\le n}\big(\{X_k\in B\}\cap\{\tau=k\}\big)\in\mathcal{F}_n$; 在 $\{\tau=\infty\}$ 上取值为 $0$, 而 $\{\tau=\infty\}\in\mathcal{F}_\tau$ (由 (1)). 两部分合起来得到结论.
\end{proof}

\begin{definition}[停止过程]\label{def:stopped-process}
设 $X$ 适应, $\tau$ 是停时. 称 $X^\tau_n=X_{\tau\wedge n}$ 为 $X$ 在 $\tau$ 处的\textbf{停止过程}.
\end{definition}

\begin{theorem}[停止过程仍是鞅]\label{thm:stopped-martingale}
若 $X$ 是鞅 (下鞅、上鞅), $\tau$ 是停时, 则 $X^\tau$ 也是鞅 (下鞅、上鞅).
\end{theorem}

\begin{proof}
令 $H_k=\1_{\{\tau\ge k\}}=1-\1_{\{\tau\le k-1\}}$, 它是 $\mathcal{F}_{k-1}$-可测的, 所以 $H$ 是可料、非负且有界的. 由于 $\sum_{k=1}^nH_k\Delta X_k=\sum_{k=1}^{\tau\wedge n}\Delta X_k=X_{\tau\wedge n}-X_0$,
\[
X^\tau_n=X_0+(H\cdot X)_n.
\]
由鞅变换定理~\ref{thm:martingale-transform}, $X^\tau$ 是鞅 (下鞅、上鞅).
\end{proof}

% ---------------------------------------------------------------------
\subsection{可选停时定理}

\begin{theorem}[可选停时定理: 有界停时]\label{thm:ost-bounded}
设 $X$ 是鞅 (下鞅), $\sigma\le\tau\le N$ 是有界停时 ($N$ 为常数). 则 $X_\tau$, $X_\sigma$ 可积, 且
\[
\E[X_\tau\,|\,\mathcal{F}_\sigma]=X_\sigma\qquad(\text{下鞅时为 }\ge).
\]
特别地, $\E X_\tau=\E X_\sigma=\E X_0$ (下鞅时 $\E X_0\le\E X_\sigma\le\E X_\tau\le\E X_N$).
\end{theorem}

\begin{proof}
可积性: $\abs{X_\tau}\le\sum_{k=0}^N\abs{X_k}$.

\emph{第一步: $\E X_\tau=\E X_\sigma$ (下鞅时为 $\ge$).} 令 $H_k=\1_{\{\sigma<k\le\tau\}}=\1_{\{\sigma\le k-1\}}-\1_{\{\tau\le k-1\}}$, 它可料、非负、有界, 且
$(H\cdot X)_N=\sum_{k=\sigma+1}^{\tau}\Delta X_k=X_\tau-X_\sigma$. 由鞅变换定理, $\E[X_\tau-X_\sigma]=\E[(H\cdot X)_N]=\E[(H\cdot X)_0]=0$ (下鞅时 $\ge0$).

\emph{第二步: 条件形式.} 设 $A\in\mathcal{F}_\sigma$. 定义
\[
\sigma_A=\begin{cases}\sigma,&\text{在 }A\text{ 上},\\ N,&\text{在 }A^c\text{ 上},\end{cases}\qquad
\tau_A=\begin{cases}\tau,&\text{在 }A\text{ 上},\\ N,&\text{在 }A^c\text{ 上}.\end{cases}
\]
它们是停时: 对 $n<N$, $\{\sigma_A\le n\}=A\cap\{\sigma\le n\}\in\mathcal{F}_n$; 对 $n\ge N$, $\{\sigma_A\le n\}=\Omega$. 由于 $A\in\mathcal{F}_\sigma\subseteq\mathcal{F}_\tau$ (命题~\ref{prop:F-tau} (2)), 对 $\tau_A$ 同理. 又 $\sigma_A\le\tau_A\le N$. 由第一步, $\E X_{\tau_A}=\E X_{\sigma_A}$, 即
\[
\E[X_\tau;A]+\E[X_N;A^c]=\E[X_\sigma;A]+\E[X_N;A^c],
\]
所以 $\E[X_\tau;A]=\E[X_\sigma;A]$ (下鞅时为 $\ge$). 由命题~\ref{prop:F-tau} (3), $X_\sigma$ 是 $\mathcal{F}_\sigma$-可测的. 鞅情形由条件期望的定义即得; 下鞅情形, 由于 $\E\big[\E[X_\tau\,|\,\mathcal{F}_\sigma];A\big]=\E[X_\tau;A]\ge\E[X_\sigma;A]$ 对一切 $A\in\mathcal{F}_\sigma$ 成立, 由比较引理~\ref{lem:ce-compare} 即得.
\end{proof}

\begin{intuition}
在有限期限内, 任何不偷看未来的离场规则都无法使公平赌博变得有利. ``领先时收手'' 看似聪明, 但它同时意味着在落后时继续赌下去, 而落后时的损失恰好抵消了领先时的收益.
\end{intuition}

有界性不能简单地去掉.

\begin{example}\label{ex:ost-fails}
设 $X_n$ 是从 $0$ 出发的简单对称随机游动, $\tau=\inf\{n: X_n=1\}$. 下面的例~\ref{ex:gambler-ruin} 将证明 $\tau<\infty$ 几乎必然成立, 从而 $X_\tau=1$, $\E X_\tau=1\ne0=\E X_0$. 加倍策略 (例~\ref{ex:doubling}) 是同一现象: 离场时刻 ``首次赢的时刻'' 几乎必然有限, 但无界. 问题出在离场之前可能经历巨大的亏损: 由例~\ref{ex:gambler-ruin} 中的计算, $\E\tau=\infty$.
\end{example}

对无界的停时, 需要附加条件. 下面的充分条件不依赖于收敛理论; 3.4 节将给出基于一致可积性的最一般的版本.

\begin{theorem}[可选停时定理: 有界增量]\label{thm:ost-bounded-increments}
设 $X$ 是鞅, $\tau$ 是停时, $\E\tau<\infty$, 且存在常数 $c$ 使对一切 $n\ge1$,
$\E\big[\abs{\Delta X_n}\,\big|\,\mathcal{F}_{n-1}\big]\le c$. 则 $X_\tau$ 可积, 且 $\E X_\tau=\E X_0$.
\end{theorem}

\begin{proof}
由 $\E\tau<\infty$, $\tau<\infty$ $\as$, 所以 $X_{\tau\wedge n}\to X_\tau$ $\as$ 令 $Y=\sum_{k=1}^\tau\abs{\Delta X_k}=\sum_{k=1}^\infty\abs{\Delta X_k}\1_{\{\tau\ge k\}}$, 则 $\abs{X_{\tau\wedge n}}\le\abs{X_0}+Y$ 对一切 $n$ 成立. 由于 $\{\tau\ge k\}\in\mathcal{F}_{k-1}$, 由逐项求期望 (推论~\ref{cor:series})、全期望公式与提出已知因子,
\[
\E Y=\sum_{k=1}^\infty\E\Big[\1_{\{\tau\ge k\}}\E\big[\abs{\Delta X_k}\,\big|\,\mathcal{F}_{k-1}\big]\Big]\le c\sum_{k=1}^\infty\P(\tau\ge k)=c\,\E\tau<\infty,
\]
最后一步是尾积分公式的离散形式 ($\tau=\sum_{k\ge1}\1_{\{\tau\ge k\}}$, 再逐项求期望). 由控制收敛定理~\ref{thm:dct} (控制函数 $\abs{X_0}+Y$) 与定理~\ref{thm:ost-bounded} (用于有界停时 $\tau\wedge n$),
\[
\E X_\tau=\lim_{n\to\infty}\E X_{\tau\wedge n}=\E X_0.\qedhere
\]
\end{proof}

\begin{corollary}[Wald 恒等式]\label{cor:wald}
设 $\xi_1,\xi_2,\dots$ 是 $\iid$ 可积随机变量, $\E\xi_1=\mu$, $S_n=\sum_{k\le n}\xi_k$, $\mathcal{F}_n=\sigma(\xi_1,\dots,\xi_n)$, $\tau$ 是停时且 $\E\tau<\infty$.
\begin{enumerate}
\item $\E S_\tau=\mu\,\E\tau$.
\item 若进一步 $\mu=0$, $\Var(\xi_1)=\sigma^2<\infty$, 则 $\E S_\tau^2=\sigma^2\E\tau$.
\end{enumerate}
\end{corollary}

\begin{proof}
(1) $M_n=S_n-n\mu$ 是鞅 (例~\ref{ex:sum-martingales} (1)), 且 $\E[\abs{\Delta M_n}\,|\,\mathcal{F}_{n-1}]=\E\abs{\xi_n-\mu}=\E\abs{\xi_1-\mu}$ 是常数 (独立性化简). 由定理~\ref{thm:ost-bounded-increments}, $\E[S_\tau-\tau\mu]=0$.

(2) 记 $T_n=S_{\tau\wedge n}=\sum_{k=1}^n\xi_k\1_{\{\tau\ge k\}}$. 对 $j<k$, $\xi_j\1_{\{\tau\ge j\}}\1_{\{\tau\ge k\}}$ 是 $\mathcal{F}_{k-1}$-可测的, 而 $\xi_k$ 与 $\mathcal{F}_{k-1}$ 独立, 所以交叉项 $\E[\xi_j\1_{\{\tau\ge j\}}\xi_k\1_{\{\tau\ge k\}}]=0$; 同理 $\E[\xi_k^2\1_{\{\tau\ge k\}}]=\sigma^2\P(\tau\ge k)$. 于是对 $m<n$,
\[
\E\big[(T_n-T_m)^2\big]=\sum_{k=m+1}^n\sigma^2\P(\tau\ge k)\xrightarrow{m,n\to\infty}0,
\]
因为 $\sum_k\P(\tau\ge k)=\E\tau<\infty$. 由 Riesz-Fischer 定理~\ref{thm:riesz-fischer}, $T_n$ 在 $L^2$ 中收敛, 而它几乎必然收敛到 $S_\tau$, 故 $L^2$ 极限就是 $S_\tau$. 由内积的连续性 (命题~\ref{prop:inner-continuity}), $\E S_\tau^2=\lim_n\E T_n^2=\lim_n\sigma^2\sum_{k\le n}\P(\tau\ge k)=\sigma^2\E\tau$.
\end{proof}

% ---------------------------------------------------------------------
\subsection{应用: 赌徒破产与模式等待时间}

\begin{example}[赌徒破产问题]\label{ex:gambler-ruin}
赌徒的初始财富为整数 $x$, $0<x<N$. 每局以概率 $p$ 赢 $1$、以概率 $q=1-p$ 输 $1$, 各局独立. 他在财富达到 $N$ 或 $0$ (破产) 时停止. 设 $S_n=x+\sum_{k\le n}\xi_k$, $\tau=\inf\{n: S_n\in\{0,N\}\}$ (图~\ref{fig:gambler-ruin}).

\emph{$\E\tau<\infty$.} 把时间分成长为 $N$ 的块. 若某一块中 $N$ 局全赢, 则无论起点在哪里, 赌徒在这一块结束之前必然已经停止. 各块独立, 每块全赢的概率为 $p^N$, 所以 $\P(\tau>kN)\le(1-p^N)^k$. 于是
\[
\E\tau=\sum_{n\ge0}\P(\tau>n)\le\sum_{k\ge0}N\,\P(\tau>kN)\le N\sum_{k\ge0}(1-p^N)^k=\frac{N}{p^N}<\infty.
\]

\emph{对称情形 $p=q=\frac{1}{2}$.} $S_n$ 是鞅, 增量有界, 由定理~\ref{thm:ost-bounded-increments}, $\E S_\tau=x$. 而 $S_\tau$ 只取 $0$ 与 $N$, 所以
\[
\P(S_\tau=N)=\frac{x}{N},\qquad\P(S_\tau=0)=\frac{N-x}{N}.
\]
再由 Wald 第二恒等式 (推论~\ref{cor:wald} (2), $\sigma^2=1$),
\[
\E\tau=\E(S_\tau-x)^2=\frac{x}{N}(N-x)^2+\frac{N-x}{N}x^2=x(N-x).
\]

\emph{非对称情形 $p\ne q$.} 令 $r=\frac{q}{p}$. 由于 $\E[r^{\xi_k}]=p\cdot\frac{q}{p}+q\cdot\frac{p}{q}=1$, $r^{S_n}$ 是鞅 (例~\ref{ex:sum-martingales} (3)). 停止过程 $r^{S_{\tau\wedge n}}$ 是有界鞅 (取值于 $r^0,\dots,r^N$ 之间), 由有界收敛定理~\ref{cor:bct} 与定理~\ref{thm:ost-bounded},
$\E r^{S_\tau}=\lim_n\E r^{S_{\tau\wedge n}}=r^x$. 解 $r^N\P(S_\tau=N)+(1-\P(S_\tau=N))=r^x$ 得
\[
\P(S_\tau=N)=\frac{r^x-1}{r^N-1}.
\]
由 Wald 第一恒等式 (推论~\ref{cor:wald} (1), $\mu=p-q$), $\E\tau=\frac{\E S_\tau-x}{p-q}=\frac{1}{p-q}\Big(N\frac{r^x-1}{r^N-1}-x\Big)$.

\emph{对称随机游动的常返性.} 设 $X$ 是从 $0$ 出发的对称随机游动. 把它看作 $x=N_0$、目标为 $N_0+1$、下界为 $0$ 的赌徒破产问题 (平移 $N_0$), 则 $X$ 在跌到 $-N_0$ 之前到达 $1$ 的概率为 $\frac{N_0}{N_0+1}$. 令 $N_0\to\infty$, 由概率的连续性, $\P(X\text{ 终将到达 }1)=1$. 但若 $\E\tau_1<\infty$ ($\tau_1$ 为首达 $1$ 的时刻), 由 Wald 第一恒等式将得到 $1=\E X_{\tau_1}=0\cdot\E\tau_1=0$, 矛盾. 所以 $\E\tau_1=\infty$: 对称随机游动必然到达 $1$, 但平均需要无穷长的时间. 这正是例~\ref{ex:ost-fails} 中可选停时定理失效的原因.
\end{example}

\begin{figure}[htbp]
\centering
\begin{tikzpicture}
\begin{axis}[SAaxis, width=11cm, height=5.4cm, xmin=0, xmax=16, ymin=-0.8, ymax=10.5,
  xtick={5,10,15}, ytick={3,8}, yticklabels={$x$,$N$}, xlabel={$n$}, ylabel={$S_n$}]
  \addplot[SAGreenText,thick,forget plot] coordinates {(0,8) (16,8)};
  \addplot[SARed,thick,forget plot] coordinates {(0,0) (16,0)};
  \addplot[SABlue,very thick,mark=*,mark size=1.2pt] coordinates {(0,3) (1,4) (2,3) (3,2) (4,3) (5,4) (6,5) (7,4) (8,5) (9,6) (10,7) (11,6) (12,7) (13,8)};
  \addplot[SADark!50,dashed] coordinates {(13,-0.8) (13,8)};
  \node[SAGreenText,font=\small,anchor=south west] at (axis cs:0.3,8.1) {达到目标 (概率 $\frac{x}{N}$)};
  \node[SARed,font=\small,anchor=south east] at (axis cs:15.8,0.1) {破产 (概率 $\frac{N-x}{N}$)};
  \node[font=\small,anchor=south] at (axis cs:13,8.3) {$\tau$};
\end{axis}
\end{tikzpicture}
\caption{对称赌徒破产问题: 从 $x$ 出发的随机游动在两条吸收边界之间运动, 直到首次碰到其中一条. 由可选停时定理, 碰到上边界的概率为 $\frac{x}{N}$, 平均持续时间为 $x(N-x)$}
\label{fig:gambler-ruin}
\end{figure}

\begin{example}[ABRACADABRA]\label{ex:abracadabra}
一只猴子在打字机上独立地、均匀随机地从 26 个字母中敲击. 问: 平均需要敲多少次才首次出现单词 ABRACADABRA?

一般地, 设字母表有 $m$ 个字母, $X_1,X_2,\dots$ 是 $\iid$ 的均匀随机字母, 目标单词为 $w=w_1w_2\cdots w_L$, $\tau=\inf\{n\ge L: X_{n-L+1}\cdots X_n=w\}$, $\mathcal{F}_n=\sigma(X_1,\dots,X_n)$. 我们用一群 ``公平赌徒'' 构造鞅.

\emph{构造.} 在第 $j$ 次敲击前 ($j=1,2,\dots$), 第 $j$ 个赌徒带着本金 $1$ 入场, 押注 ``$X_j=w_1$'': 若押中, 财富变为 $m$ (赔率公平, 因为押中的概率为 $\frac{1}{m}$), 否则变为 $0$ 并离场. 若押中, 他把全部财富押在 ``$X_{j+1}=w_2$'' 上, 依此类推, 直到押错或押完整个单词为止. 记第 $j$ 个赌徒在时刻 $n$ 的财富为 $G^{(j)}_n$ (入场前约定为 $1$). 具体地, 对 $n\ge j$, 令 $\ell=\min\{n-j+1,L\}$, 则
\[
G^{(j)}_n=m^{\ell}\,\1_{\{X_j=w_1,\dots,X_{j+\ell-1}=w_\ell\}}.
\]
\begin{itemize}
\item \emph{每个 $G^{(j)}$ 是鞅.} 若 $j\le n+1\le j+L-1$, 则 $G^{(j)}_{n+1}=G^{(j)}_n\cdot m\1_{\{X_{n+1}=w_{n+2-j}\}}$, 其中 $G^{(j)}_n$ 是 $\mathcal{F}_n$-可测的, 第二个因子与 $\mathcal{F}_n$ 独立且期望为 $1$; 否则 $G^{(j)}_{n+1}=G^{(j)}_n$. 由提出已知因子与独立性化简, $\E[G^{(j)}_{n+1}\,|\,\mathcal{F}_n]=G^{(j)}_n$.
\item \emph{总净收益是鞅.} 令 $M_n=\sum_{j=1}^n(G^{(j)}_n-1)$, 即已入场的 $n$ 个赌徒的总财富减去总本金. 由于 $j>n$ 时 $G^{(j)}_n-1=0$, 也可写成 $M_n=\sum_{j\ge1}(G^{(j)}_n-1)$, 其中对每个 $n$ 只有有限项非零. 于是 $M_{n+1}-M_n=\sum_{j=n-L+2}^{n+1}(G^{(j)}_{n+1}-G^{(j)}_n)$ (只有仍在场的赌徒的财富会变化), 每项的条件期望为零, 所以 $M$ 是鞅, $M_0=0$; 且 $\abs{\Delta M_n}\le L(m^L+1)$, 增量有界.
\item \emph{$\E\tau<\infty$.} 与例~\ref{ex:gambler-ruin} 相同, 把时间分成长为 $L$ 的不交块, 每块恰为 $w$ 的概率为 $m^{-L}$, 故 $\E\tau\le Lm^L<\infty$.
\end{itemize}
由定理~\ref{thm:ost-bounded-increments}, $\E M_\tau=0$.

\emph{计算 $M_\tau$.} 在时刻 $\tau$, 最后 $L$ 个字母恰为 $w$. 第 $j$ 个赌徒的财富分两种情形:
\begin{itemize}
\item 若 $j\le\tau-L$, 他已经离场. 事实上, 若他押完了整个单词, 则单词在 $\tau$ 之前已经出现, 与 $\tau$ 的定义矛盾; 所以他必然押错过, 财富为 $0$.
\item 若 $j=\tau-k+1$ ($1\le k\le L$), 他押了 $k$ 次, 需要 $X_{\tau-k+1}\cdots X_\tau=w_1\cdots w_k$. 而已知 $X_{\tau-k+1}\cdots X_\tau=w_{L-k+1}\cdots w_L$ (单词的最后 $k$ 个字母). 所以他的财富为 $m^k$ 当且仅当 $w$ 的长为 $k$ 的前缀等于长为 $k$ 的后缀, 否则为 $0$.
\end{itemize}
于是 $M_\tau=c(w)-\tau$, 其中 $c(w)=\sum_{k\in K(w)}m^k$, $K(w)=\{k\le L: w_1\cdots w_k=w_{L-k+1}\cdots w_L\}$ 是确定的数. 由 $\E M_\tau=0$ 得
\[
\E\tau=c(w)=\sum_{k\in K(w)}m^k.
\]
对 ABRACADABRA, 前缀与后缀相同的长度为 $k=1$ (A)、$k=4$ (ABRA)、$k=11$ (整个单词), 所以平均等待时间为 $26^{11}+26^4+26$. 再如抛硬币 ($m=2$): 等待 ``正正'' 平均需要 $2^2+2=6$ 次, 而等待 ``正反'' 只需 $2^2=4$ 次. 自身重叠的模式反而更 ``难等'': 一旦部分匹配失败, 重叠模式浪费的进度更多.
\end{example}

% =====================================================================
\section{Doob 分解与 Doob 鞅不等式}
\label{sec:3-3}

% ---------------------------------------------------------------------
\subsection{Doob 分解}

下鞅是 ``有利的赌博'', 能否把有利的部分与公平的部分分开? 答案是肯定的, 每个适应过程都可以唯一地写成 ``鞅 + 可料过程''.

\begin{theorem}[Doob 分解]\label{thm:doob-decomposition}
设 $X$ 适应且每个 $X_n$ 可积. 则存在鞅 $M$ (满足 $M_0=0$) 与可料过程 $A$ (满足 $A_0=0$, 即 $A_n$ 是 $\mathcal{F}_{n-1}$-可测的, $n\ge1$), 使得
\[
X_n=X_0+M_n+A_n,\qquad n\ge0,
\]
并且这种分解在几乎必然的意义下唯一. $X$ 是下鞅当且仅当 $A$ 单调不减; $X$ 是鞅当且仅当 $A\equiv0$.
\end{theorem}

\begin{proof}
\emph{存在性.} 令
\[
A_n=\sum_{k=1}^n\E[\Delta X_k\,|\,\mathcal{F}_{k-1}],\qquad M_n=X_n-X_0-A_n=\sum_{k=1}^n\big(\Delta X_k-\E[\Delta X_k\,|\,\mathcal{F}_{k-1}]\big).
\]
$A_n$ 是 $\mathcal{F}_{n-1}$-可测的, 所以 $A$ 可料. $M$ 适应且可积, 并且
$\E[\Delta M_n\,|\,\mathcal{F}_{n-1}]=\E[\Delta X_n\,|\,\mathcal{F}_{n-1}]-\E[\Delta X_n\,|\,\mathcal{F}_{n-1}]=0$, 所以 $M$ 是鞅.

\emph{唯一性.} 若 $X_n=X_0+M_n+A_n=X_0+M'_n+A'_n$ 是两个这样的分解, 则 $N=M-M'=A'-A$ 既是鞅又是可料的, 且 $N_0=0$. 于是 $N_n=\E[N_n\,|\,\mathcal{F}_{n-1}]=N_{n-1}$ (前一个等号因为 $N_n$ 是 $\mathcal{F}_{n-1}$-可测的), 递推得 $N_n=N_0=0$.

\emph{刻画.} $\Delta A_n=\E[\Delta X_n\,|\,\mathcal{F}_{n-1}]$, 所以 $A$ 单调不减当且仅当对一切 $n$, $\E[\Delta X_n\,|\,\mathcal{F}_{n-1}]\ge0$, 即 $X$ 是下鞅. 鞅的情形同理.
\end{proof}

\begin{intuition}
$\Delta A_n=\E[\Delta X_n\,|\,\mathcal{F}_{n-1}]$ 是 ``在时刻 $n-1$ 就能预见的那部分变化'' (漂移), $\Delta M_n$ 是 ``不可预见的意外'' (噪声). 这是随机微分方程 $\dd X_t=b\dd t+\sigma\dd B_t$ 中 ``漂移 + 扩散'' 分解的离散原型. 可料性保证了分解的唯一性: 若允许 $A$ 只是适应的, 则 $X$ 自己就是一个平凡的分解.
\end{intuition}

对平方可积鞅 $M$, $M^2$ 是下鞅 (命题~\ref{prop:convex-submartingale}), 它的 Doob 分解中的可料部分具有特别重要的意义.

\begin{definition}[尖括号过程与二次变差]\label{def:bracket}
设 $M$ 是平方可积鞅 (即每个 $M_n\in L^2$). 定义
\[
\ip{M}_n=\sum_{k=1}^n\E\big[(\Delta M_k)^2\,\big|\,\mathcal{F}_{k-1}\big],\qquad [M]_n=\sum_{k=1}^n(\Delta M_k)^2,
\]
$\ip{M}_0=[M]_0=0$. 称 $\ip{M}$ 为 $M$ 的\textbf{尖括号过程} (条件二次变差, 或可料二次变差), 称 $[M]$ 为 $M$ 的\textbf{二次变差}.
\end{definition}

\begin{proposition}\label{prop:bracket}
设 $M$ 是平方可积鞅. 则
\begin{enumerate}
\item 鞅增量正交: 对 $j<k$, $\E[\Delta M_j\Delta M_k]=0$; 且 $\E[(\Delta M_k)^2\,|\,\mathcal{F}_{k-1}]=\E[M_k^2-M_{k-1}^2\,|\,\mathcal{F}_{k-1}]$.
\item $M^2-\ip{M}$ 是鞅, 即 $M_n^2=M_0^2+(\text{鞅})+\ip{M}_n$ 是 $M^2$ 的 Doob 分解.
\item 若 $[M]_n$ 可积, 则 $M^2-[M]$ 也是鞅.
\item $\E[M_n^2]=\E[M_0^2]+\E\ip{M}_n=\E[M_0^2]+\E[M]_n$.
\end{enumerate}
\end{proposition}

\begin{proof}
(1) $\Delta M_j$ 是 $\mathcal{F}_{k-1}$-可测的, 由提出已知因子 (乘积可积由 Cauchy-Schwarz 不等式保证),
\[
\E[\Delta M_j\Delta M_k]=\E\big[\Delta M_j\E[\Delta M_k\,|\,\mathcal{F}_{k-1}]\big]=0.
\]
其次, $M_k^2-M_{k-1}^2=(\Delta M_k)^2+2M_{k-1}\Delta M_k$, 而
\[
\E[M_{k-1}\Delta M_k\,|\,\mathcal{F}_{k-1}]=M_{k-1}\E[\Delta M_k\,|\,\mathcal{F}_{k-1}]=0.
\]

(2) 由 (1), $\E[\Delta(M^2)_k\,|\,\mathcal{F}_{k-1}]=\E[(\Delta M_k)^2\,|\,\mathcal{F}_{k-1}]=\Delta\ip{M}_k$, 这正是 Doob 分解中的可料部分 (定理~\ref{thm:doob-decomposition} 的证明).

(3) $M_n^2-[M]_n=M_0^2+\sum_{k=1}^n\big(M_k^2-M_{k-1}^2-(\Delta M_k)^2\big)=M_0^2+2\sum_{k=1}^nM_{k-1}\Delta M_k$, 由 (1) 的证明, 每一项的条件期望为零.

(4) 对 (2)(3) 取期望; 或直接由 (1): $\E[M_n^2]=\E[M_0^2]+\sum_{k\le n}\E[(\Delta M_k)^2]$.
\end{proof}

\begin{example}\label{ex:bracket-rw}
设 $S_n=\sum_{k\le n}\xi_k$, $\xi_k$ 独立, $\E\xi_k=0$, $\Var(\xi_k)=\sigma_k^2$. 则 $\ip{S}_n=\sum_{k\le n}\sigma_k^2$ 是确定性的, 而 $[S]_n=\sum_{k\le n}\xi_k^2$ 是随机的. 对简单对称随机游动 ($\xi_k=\pm1$), 二者都等于 $n$: Doob 分解 $S_n^2=(S_n^2-n)+n$ 把下鞅 $S_n^2$ 分成了鞅 $S_n^2-n$ 与确定性的增过程 $n$ (图~\ref{fig:doob-decomposition}). 第~4~章将证明布朗运动的二次变差也是确定性的: $[B]_t=t$.
\end{example}

\begin{figure}[htbp]
\centering
\begin{tikzpicture}
\begin{axis}[SAaxis, width=11.5cm, height=6cm, xmin=0, xmax=25, ymin=-18, ymax=28,
  xtick={5,10,15,20}, ytick={-10,10,20}, xlabel={$n$},
  legend style={draw=none,fill=none,font=\small,at={(0.5,-0.2)},anchor=north,legend columns=3,/tikz/every even column/.append style={column sep=1em}}]
  \addplot[SABlue,very thick,mark=*,mark size=0.9pt] coordinates {(0,0) (1,1) (2,4) (3,1) (4,4) (5,1) (6,0) (7,1) (8,4) (9,1) (10,0) (11,1) (12,0) (13,1) (14,4) (15,9) (16,16) (17,9) (18,4) (19,9) (20,4) (21,9) (22,16) (23,25) (24,16)};
  \addlegendentry{下鞅 $S_n^2$}
  \addplot[SARed,thick,dashed] coordinates {(0,0) (24,24)};
  \addlegendentry{可料增过程 $\ip{S}_n=n$}
  \addplot[SAGreenText,thick,mark=o,mark size=0.9pt] coordinates {(0,0) (1,0) (2,2) (3,-2) (4,0) (5,-4) (6,-6) (7,-6) (8,-4) (9,-8) (10,-10) (11,-10) (12,-12) (13,-12) (14,-10) (15,-6) (16,0) (17,-8) (18,-14) (19,-10) (20,-16) (21,-12) (22,-6) (23,2) (24,-8)};
  \addlegendentry{鞅 $S_n^2-n$}
\end{axis}
\end{tikzpicture}
\caption{Doob 分解: 简单对称随机游动的平方 $S_n^2$ (下鞅) 等于鞅 $S_n^2-n$ 加上确定性的增过程 $n$}
\label{fig:doob-decomposition}
\end{figure}

命题~\ref{prop:bracket} (4) 对鞅变换给出一个非常有用的等距公式.

\begin{corollary}[离散 Itô 等距]\label{cor:discrete-ito-isometry}
设 $M$ 是平方可积鞅, $H$ 可料且有界. 则
\[
\E\big[(H\cdot M)_n^2\big]=\E\Big[\sum_{k=1}^nH_k^2\,\Delta\ip{M}_k\Big].
\]
特别地, 对简单对称随机游动 $S$, $\E[(H\cdot S)_n^2]=\E\big[\sum_{k=1}^nH_k^2\big]$.
\end{corollary}

\begin{proof}
由鞅变换定理, $H\cdot M$ 是平方可积鞅, 其增量为 $H_k\Delta M_k$. 由命题~\ref{prop:bracket} (4) 与提出已知因子 ($H_k^2$ 是 $\mathcal{F}_{k-1}$-可测的),
\[
\E\big[(H\cdot M)_n^2\big]=\sum_{k=1}^n\E\big[H_k^2(\Delta M_k)^2\big]=\sum_{k=1}^n\E\Big[H_k^2\,\E\big[(\Delta M_k)^2\,\big|\,\mathcal{F}_{k-1}\big]\Big]=\E\Big[\sum_{k=1}^nH_k^2\Delta\ip{M}_k\Big].\qedhere
\]
\end{proof}

\begin{intuition}
左边是 ``随机积分'' 的均方, 右边是被积函数对 ``时钟'' $\ip{M}$ 的普通积分. 第~5~章的 Itô 等距 $\E\big[(\int_0^tH_s\dd B_s)^2\big]=\E\int_0^tH_s^2\dd s$ 是它的连续极限, 而它正是构造 Itô 积分的基石: 等距映射可以从简单过程唯一地延拓到它们的 $L^2$ 闭包.
\end{intuition}

\begin{remark}[Doob-Meyer 分解]
在连续时间中, Doob 分解的推广称为 Doob-Meyer 分解: 满足适当可积条件 (``(D) 类'') 的右连续下鞅可以唯一地写成鞅与可料增过程之和 (见~\cite{karatzas} 定理 1.4.10). 其证明远比离散情形困难, 因为连续时间中没有 ``前一个时刻''. 对平方可积的连续鞅 $M$, 这给出连续增过程 $\ip{M}$, 使 $M^2-\ip{M}$ 是鞅 (3.5 节).
\end{remark}

% ---------------------------------------------------------------------
\subsection{Doob 鞅不等式}

Markov 不等式用 $\E\abs{X_n}$ 控制单个时刻 $\P(\abs{X_n}\ge\lambda)$; 鞅的同一个期望还可以控制\emph{整条轨道}在 $[0,n]$ 上的最大值: 这是鞅与一般随机序列最显著的区别之一.

\begin{theorem}[Doob 最大值不等式]\label{thm:doob-maximal}
设 $X$ 是下鞅, $n\ge0$, $\lambda>0$. 记 $X^*_n=\max_{k\le n}X_k$. 则
\[
\lambda\,\P(X^*_n\ge\lambda)\le\E\big[X_n;X^*_n\ge\lambda\big]\le\E[X_n^+].
\]
特别地, 若 $M$ 是鞅, 则 $\lambda\,\P\big(\max_{k\le n}\abs{M_k}\ge\lambda\big)\le\E\abs{M_n}$.
\end{theorem}

\begin{proof}
令 $\tau=\inf\{k: X_k\ge\lambda\}\wedge n$, 它是有界停时, 且 $\tau\le n$. 记 $A=\{X^*_n\ge\lambda\}$. 在 $A$ 上, $X$ 在时刻 $n$ 之前达到 $\lambda$, 所以 $X_\tau\ge\lambda$; 在 $A^c$ 上, $\tau=n$. 由可选停时定理~\ref{thm:ost-bounded} (下鞅情形, $\tau\le n$), $\E X_\tau\le\E X_n$, 即
\[
\E[X_\tau;A]+\E[X_n;A^c]\le\E[X_n;A]+\E[X_n;A^c].
\]
消去 $\E[X_n;A^c]$ 得 $\E[X_\tau;A]\le\E[X_n;A]$. 于是
\[
\lambda\P(A)\le\E[X_\tau;A]\le\E[X_n;A]\le\E[X_n^+;A]\le\E[X_n^+].
\]
对鞅 $M$, $\abs{M}$ 是下鞅 (命题~\ref{prop:convex-submartingale}), 应用上式即得.
\end{proof}

\begin{figure}[htbp]
\centering
\begin{tikzpicture}
\begin{axis}[SAaxis, width=11cm, height=5.4cm, xmin=0, xmax=21, ymin=-3, ymax=6,
  xtick={5,10,15,20}, xticklabels={5,10,15,$n=20$}, ytick={3}, yticklabels={$\lambda$}, xlabel={$k$}, ylabel={$X_k$}]
  \addplot[SARed!60,forget plot] coordinates {(0,3) (21,3)};
  \addplot[SABlue,very thick,mark=*,mark size=0.9pt] coordinates {(0,0) (1,-1) (2,-2) (3,-1) (4,-2) (5,-1) (6,0) (7,1) (8,2) (9,1) (10,0) (11,1) (12,0) (13,1) (14,2) (15,3) (16,4) (17,3) (18,2) (19,3) (20,2)};
  \addplot[SADark!50,dashed] coordinates {(15,-3) (15,3)};
  \fill[SARed] (axis cs:15,3) circle (2.2pt);
  \node[font=\small,anchor=north west] at (axis cs:15.2,-1.3) {$\tau=15$: 首次 $X_k\ge\lambda$};
  \node[font=\small,SADark!80,anchor=south west] at (axis cs:16.2,4.3) {$X^*_{20}=4\ge\lambda$};
\end{axis}
\end{tikzpicture}
\caption{Doob 最大值不等式的证明思路: 在首次越过 $\lambda$ 的时刻 $\tau$ 停下. 在事件 $\{X^*_n\ge\lambda\}$ 上 $X_\tau\ge\lambda$, 而下鞅性质保证从 $\tau$ 到 $n$ 平均不会下降, 所以 $X_n$ 在该事件上的平均也至少为 $\lambda$}
\label{fig:doob-maximal}
\end{figure}

\begin{corollary}[Kolmogorov 不等式]\label{cor:kolmogorov-inequality}
设 $\xi_1,\dots,\xi_n$ 独立, $\E\xi_k=0$, $\Var(\xi_k)<\infty$, $S_k=\sum_{j\le k}\xi_j$. 则对 $\lambda>0$,
\[
\P\Big(\max_{k\le n}\abs{S_k}\ge\lambda\Big)\le\frac{\Var(S_n)}{\lambda^2}.
\]
\end{corollary}

\begin{proof}
$S$ 是鞅, $S^2$ 是下鞅. 由 Doob 最大值不等式 (对 $S^2$, 水平为 $\lambda^2$), $\P(\max_k S_k^2\ge\lambda^2)\le\E[S_n^2]/\lambda^2$.
\end{proof}

与 Chebyshev 不等式 $\P(\abs{S_n}\ge\lambda)\le\Var(S_n)/\lambda^2$ 相比, Kolmogorov 不等式用\emph{同一个}上界控制了整个部分和序列的最大值.

\begin{theorem}[Doob $L^p$ 不等式]\label{thm:doob-Lp}
设 $X$ 是非负下鞅, $p>1$, $q=\frac{p}{p-1}$. 则
\[
\Big\|\max_{k\le n}X_k\Big\|_p\le q\,\norm{X_n}_p.
\]
特别地, 若 $M$ 是鞅, 则 $\E\big[\max_{k\le n}\abs{M_k}^p\big]\le q^p\,\E\abs{M_n}^p$; $p=2$ 时,
\[
\E\Big[\max_{k\le n}M_k^2\Big]\le4\,\E[M_n^2].
\]
\end{theorem}

\begin{proof}
若 $\norm{X_n}_p=\infty$, 结论显然, 下设 $X_n\in L^p$. 记 $X^*=\max_{k\le n}X_k$, 并取截断 $K>0$. 由尾积分公式 (例~\ref{ex:tail-formula}, 用于 $(X^*\wedge K)^p$, 并作变量替换 $t=\lambda^p$),
\[
\E\big[(X^*\wedge K)^p\big]=\int_0^\infty p\lambda^{p-1}\P(X^*\wedge K>\lambda)\dd\lambda\le\int_0^Kp\lambda^{p-1}\P(X^*\ge\lambda)\dd\lambda.
\]
由 Doob 最大值不等式~\ref{thm:doob-maximal}, $\P(X^*\ge\lambda)\le\frac{1}{\lambda}\E[X_n;X^*\ge\lambda]$. 代入并用 Tonelli 定理~\ref{thm:fubini} 交换积分次序,
\begin{align*}
\E\big[(X^*\wedge K)^p\big]&\le\int_0^Kp\lambda^{p-2}\E\big[X_n\1_{\{X^*\ge\lambda\}}\big]\dd\lambda=\E\Big[X_n\int_0^{X^*\wedge K}p\lambda^{p-2}\dd\lambda\Big]\\
&=\frac{p}{p-1}\E\big[X_n(X^*\wedge K)^{p-1}\big]\le q\,\norm{X_n}_p\,\big\|(X^*\wedge K)^{p-1}\big\|_q,
\end{align*}
最后一步是 H\"older 不等式~\ref{thm:holder}. 由于 $(p-1)q=p$, $\big\|(X^*\wedge K)^{p-1}\big\|_q=\big(\E[(X^*\wedge K)^p]\big)^{1-1/p}$. 若 $\E[(X^*\wedge K)^p]=0$ 则无需证明; 否则它不超过 $K^p$, 是有限正数, 两边除以 $\big(\E[(X^*\wedge K)^p]\big)^{1-1/p}$ 得
$\norm{X^*\wedge K}_p\le q\norm{X_n}_p$. 令 $K\to\infty$, 由单调收敛定理即得. 对鞅 $M$, 取 $X=\abs{M}$.
\end{proof}

\begin{remark}
$p=1$ 时类似的不等式不成立: $\E[\max_{k\le n}\abs{M_k}]$ 不能被 $C\,\E\abs{M_n}$ 控制 (需要 $\E[\abs{M_n}\log^+\abs{M_n}]$, 见~\cite{durrett} 4.4 节). 这是 $L^1$ 与 $L^p$ ($p>1$) 在鞅论中表现不同的又一体现.
\end{remark}

% ---------------------------------------------------------------------
\subsection{Azuma-Hoeffding 不等式}

Doob 不等式只用到矩. 若鞅的增量有界, 可以得到指数型的尾概率估计, 其形式与正态尾概率 $e^{-\lambda^2/2}$ 相同.

\begin{lemma}[Hoeffding 引理]\label{lem:hoeffding}
设 $Y$ 是随机变量, $\abs{Y}\le c$, $\E[Y\,|\,\mathcal{G}]=0$. 则对任意 $\theta\in\R$, $\E[e^{\theta Y}\,|\,\mathcal{G}]\le e^{\theta^2c^2/2}$.
\end{lemma}

\begin{proof}
若 $c=0$, 结论显然, 下设 $c>0$. 由 $y\mapsto e^{\theta y}$ 的凸性, 对 $y\in[-c,c]$, 把 $y$ 写成端点 $\pm c$ 的凸组合 $y=\frac{c-y}{2c}(-c)+\frac{c+y}{2c}c$,
\[
e^{\theta y}\le\frac{c-y}{2c}e^{-\theta c}+\frac{c+y}{2c}e^{\theta c}.
\]
代入 $y=Y$, 取条件期望并用 $\E[Y\,|\,\mathcal{G}]=0$, 得 $\E[e^{\theta Y}\,|\,\mathcal{G}]\le\frac{1}{2}(e^{-\theta c}+e^{\theta c})=\cosh(\theta c)$. 最后, 比较幂级数的系数并利用 $(2k)!\ge2^kk!$,
\[
\cosh x=\sum_{k\ge0}\frac{x^{2k}}{(2k)!}\le\sum_{k\ge0}\frac{x^{2k}}{2^kk!}=e^{x^2/2}.\qedhere
\]
\end{proof}

\begin{theorem}[Azuma-Hoeffding 不等式]\label{thm:azuma}
设 $M$ 是鞅, $\abs{\Delta M_k}\le c_k$ ($c_k$ 为常数). 则对 $\lambda>0$,
\[
\P\Big(\max_{k\le n}(M_k-M_0)\ge\lambda\Big)\le\exp\Big(-\frac{\lambda^2}{2\sum_{k=1}^nc_k^2}\Big).
\]
\end{theorem}

\begin{proof}
记 $v_n=\sum_{k\le n}c_k^2$, 设 $v_n>0$ (否则 $M_k=M_0$, 结论显然). 固定 $\theta>0$. 先证明 $\E[e^{\theta(M_n-M_0)}]\le e^{\theta^2v_n/2}$. 对 $n$ 归纳: 由于 $e^{\theta(M_{n-1}-M_0)}$ 是 $\mathcal{F}_{n-1}$-可测的, 由提出已知因子与 Hoeffding 引理~\ref{lem:hoeffding} (取 $Y=\Delta M_n$, $\mathcal{G}=\mathcal{F}_{n-1}$),
\[
\E[e^{\theta(M_n-M_0)}]=\E\Big[e^{\theta(M_{n-1}-M_0)}\E\big[e^{\theta\Delta M_n}\,\big|\,\mathcal{F}_{n-1}\big]\Big]\le e^{\theta^2c_n^2/2}\E[e^{\theta(M_{n-1}-M_0)}].
\]
其次, $e^{\theta(M_k-M_0)}$ 是鞅的凸函数, 因而是非负下鞅 (命题~\ref{prop:convex-submartingale}), 由 Doob 最大值不等式~\ref{thm:doob-maximal},
\begin{align*}
\P\Big(\max_{k\le n}(M_k-M_0)\ge\lambda\Big)&=\P\Big(\max_{k\le n}e^{\theta(M_k-M_0)}\ge e^{\theta\lambda}\Big)\\
&\le e^{-\theta\lambda}\E[e^{\theta(M_n-M_0)}]\le e^{-\theta\lambda+\theta^2v_n/2}.
\end{align*}
取 $\theta=\lambda/v_n$ 使右端最小, 即得.
\end{proof}

\begin{example}\label{ex:azuma}
对简单对称随机游动 ($c_k=1$), $\P(\max_{k\le n}S_k\ge\lambda)\le e^{-\lambda^2/(2n)}$. 取 $\lambda=x\sqrt{n}$, 得 $\P(\max_{k\le n}S_k\ge x\sqrt{n})\le e^{-x^2/2}$: 随机游动在 $n$ 步内的最大偏离以 $\sqrt{n}$ 为尺度, 并且有正态型的尾部. 这正是第~4~章把随机游动按 $\sqrt{n}$ 缩放得到布朗运动的原因. Azuma-Hoeffding 不等式也广泛用于随机算法与组合概率: 只要一个量对每个独立输入的依赖都有界, 它的 Doob 鞅就有有界增量, 从而它高度集中在其均值附近.
\end{example}

% =====================================================================
\section{鞅收敛定理}
\label{sec:3-4}

图~\ref{fig:polya-paths} 画出了 P\'olya 罐子中红球比例的几条轨道. 每条轨道最终都稳定在某个值附近, 这是有界鞅必定收敛的结果; 不同轨道的极限则不同. 本节证明这一事实及其各种变形. 核心思想是 Doob 的上穿计数. 若一条轨道不收敛, 它必须在某两个水平之间来回穿越无穷多次; 而 (下) 鞅穿越的期望次数是有限的.

\begin{figure}[htbp]
\centering
\begin{tikzpicture}
\begin{axis}[SAaxis, width=11.5cm, height=5.6cm, xmin=0, xmax=205, ymin=0, ymax=1.02,
  xtick={50,100,150,200}, ytick={0.25,0.5,0.75,1}, xlabel={$n$}, ylabel={$X_n$}]
  \addplot[SABlue,thick] coordinates {(0,0.500) (5,0.857) (10,0.917) (15,0.882) (20,0.909) (25,0.926) (30,0.938) (35,0.946) (40,0.952) (45,0.936) (50,0.942) (55,0.930) (60,0.919) (65,0.925) (70,0.917) (75,0.922) (80,0.927) (85,0.908) (90,0.913) (95,0.918) (100,0.912) (105,0.907) (110,0.911) (115,0.906) (120,0.910) (125,0.913) (130,0.909) (135,0.912) (140,0.915) (145,0.918) (150,0.921) (155,0.924) (160,0.926) (165,0.928) (170,0.930) (175,0.932) (180,0.934) (185,0.936) (190,0.922) (195,0.919) (200,0.921)};
  \addplot[SARed,thick] coordinates {(0,0.500) (5,0.429) (10,0.333) (15,0.353) (20,0.318) (25,0.296) (30,0.281) (35,0.270) (40,0.310) (45,0.340) (50,0.365) (55,0.386) (60,0.419) (65,0.433) (70,0.431) (75,0.416) (80,0.402) (85,0.391) (90,0.391) (95,0.381) (100,0.373) (105,0.383) (110,0.402) (115,0.393) (120,0.385) (125,0.378) (130,0.386) (135,0.387) (140,0.394) (145,0.388) (150,0.375) (155,0.376) (160,0.370) (165,0.383) (170,0.378) (175,0.384) (180,0.390) (185,0.385) (190,0.375) (195,0.365) (200,0.366)};
  \addplot[SAGreenText,thick] coordinates {(0,0.500) (5,0.714) (10,0.583) (15,0.647) (20,0.636) (25,0.704) (30,0.719) (35,0.703) (40,0.667) (45,0.660) (50,0.635) (55,0.632) (60,0.613) (65,0.612) (70,0.611) (75,0.610) (80,0.622) (85,0.621) (90,0.630) (95,0.639) (100,0.637) (105,0.626) (110,0.643) (115,0.641) (120,0.639) (125,0.630) (130,0.629) (135,0.635) (140,0.641) (145,0.633) (150,0.645) (155,0.637) (160,0.630) (165,0.617) (170,0.622) (175,0.627) (180,0.632) (185,0.626) (190,0.635) (195,0.640) (200,0.649)};
  \addplot[SAYellowText,thick] coordinates {(0,0.500) (5,0.143) (10,0.083) (15,0.059) (20,0.045) (25,0.037) (30,0.031) (35,0.027) (40,0.024) (45,0.021) (50,0.019) (55,0.018) (60,0.016) (65,0.015) (70,0.014) (75,0.026) (80,0.037) (85,0.046) (90,0.054) (95,0.052) (100,0.049) (105,0.047) (110,0.045) (115,0.043) (120,0.041) (125,0.039) (130,0.038) (135,0.044) (140,0.042) (145,0.041) (150,0.046) (155,0.045) (160,0.043) (165,0.042) (170,0.041) (175,0.040) (180,0.038) (185,0.037) (190,0.036) (195,0.041) (200,0.040)};
\end{axis}
\end{tikzpicture}
\caption{P\'olya 罐子中红球比例 $X_n$ 的四条模拟轨道. 每条轨道都收敛, 但极限是随机的; 定理~\ref{thm:mart-conv} 与例~\ref{ex:polya-limit} 将证明极限服从 $\mathrm{U}(0,1)$}
\label{fig:polya-paths}
\end{figure}

% ---------------------------------------------------------------------
\subsection{上穿不等式与几乎必然收敛}

设 $a<b$. 对适应过程 $X$, 递归地定义停时
\[
\sigma_1=\inf\{k\ge0: X_k\le a\},\quad \tau_j=\inf\{k>\sigma_j: X_k\ge b\},\quad \sigma_{j+1}=\inf\{k>\tau_j: X_k\le a\},
\]
并令 $U_n[a,b]=\sup\{j: \tau_j\le n\}$ (约定 $\sup\varnothing=0$) 为 $X$ 在时刻 $n$ 之前完成的\textbf{上穿} $[a,b]$ 的次数: 每次上穿是指从 $a$ 以下出发、到达 $b$ 以上 (图~\ref{fig:upcrossing}). 这些 $\sigma_j,\tau_j$ 都是停时: 例如 $\{\tau_j\le n\}=\bigcup_{k\le n}\big(\{\sigma_j\le k-1\}\cap\{X_k\ge b\}\big)$, 由归纳法可知它属于 $\mathcal{F}_n$.

\begin{figure}[htbp]
\centering
\begin{tikzpicture}
\begin{axis}[SAaxis, width=11.5cm, height=5.6cm, xmin=0, xmax=21, ymin=0, ymax=5.6,
  xtick={2,5,8,11,15}, xticklabels={$\sigma_1$,$\tau_1$,$\sigma_2$,$\tau_2$,$\sigma_3$}, ytick={1,4}, yticklabels={$a$,$b$}, xlabel={$k$}]
  \addplot[SADark!25,forget plot] coordinates {(0,1) (21,1)};
  \addplot[SADark!25,forget plot] coordinates {(0,4) (21,4)};
  \addplot[SABlue!50,thick,mark=*,mark size=0.9pt] coordinates {(0,3) (1,2) (2,0.5) (3,1.5) (4,3) (5,4.5) (6,3.5) (7,2) (8,0.8) (9,2) (10,3.2) (11,4.2) (12,5) (13,3) (14,1.5) (15,0.7) (16,1.8) (17,3) (18,3.8) (19,2.5) (20,3.5)};
  \addplot[SARed,very thick] coordinates {(2,0.5) (3,1.5) (4,3) (5,4.5)};
  \addplot[SARed,very thick] coordinates {(8,0.8) (9,2) (10,3.2) (11,4.2)};
  \addplot[SARed,very thick,dashed] coordinates {(15,0.7) (16,1.8) (17,3) (18,3.8) (19,2.5) (20,3.5)};
  \node[font=\small,anchor=south east] at (axis cs:20.8,4.7) {$U_{20}[a,b]=2$};
\end{axis}
\end{tikzpicture}
\caption{上穿计数. 粗线是两次完成的上穿; 从 $\sigma_3$ 开始的最后一段 (虚线) 尚未到达 $b$, 不计入 $U_{20}[a,b]$. 证明中的策略 $H$ 恰在粗线与虚线所在的时段持有一单位 ``股票''}
\label{fig:upcrossing}
\end{figure}

\begin{theorem}[Doob 上穿不等式]\label{thm:upcrossing}
设 $X$ 是下鞅, $a<b$. 则
\[
(b-a)\,\E\,U_n[a,b]\le\E(X_n-a)^+-\E(X_0-a)^+\le\E X_n^++\abs{a}.
\]
\end{theorem}

\begin{proof}
令 $Y_k=(X_k-a)^+$. 由命题~\ref{prop:convex-submartingale} (2), $Y$ 是非负下鞅. $X$ 上穿 $[a,b]$ 恰好是 $Y$ 上穿 $[0,b-a]$, 且 $Y_{\sigma_j}=0$, $Y_{\tau_j}\ge b-a$.

考虑如下策略: 在每段 $(\sigma_j,\tau_j]$ 中持有一单位, 其余时间不持有, 即
\[
H_k=\sum_{j\ge1}\1_{\{\sigma_j<k\le\tau_j\}}=\sum_{j\ge1}\big(\1_{\{\sigma_j\le k-1\}}-\1_{\{\tau_j\le k-1\}}\big).
\]
由于这些时段互不相交, $H_k\in\{0,1\}$; 由第二个表达式, $H_k$ 是 $\mathcal{F}_{k-1}$-可测的, 所以 $H$ 可料. 收益 $(H\cdot Y)_n$ 是各段中 $Y$ 的增量之和: 每段完成的上穿贡献 $Y_{\tau_j}-Y_{\sigma_j}\ge b-a$; 若在时刻 $n$ 有一段尚未完成, 它贡献 $Y_n-Y_{\sigma_j}=Y_n\ge0$. 所以
\[
(H\cdot Y)_n\ge(b-a)\,U_n[a,b].
\]
另一方面, $1-H$ 可料、非负、有界, 由鞅变换定理~\ref{thm:martingale-transform} (2), $((1-H)\cdot Y)$ 是下鞅, $\E[((1-H)\cdot Y)_n]\ge0$. 由于 $(H\cdot Y)_n+((1-H)\cdot Y)_n=Y_n-Y_0$,
\[
(b-a)\,\E\,U_n[a,b]\le\E[(H\cdot Y)_n]\le\E[(H\cdot Y)_n]+\E[((1-H)\cdot Y)_n]=\E Y_n-\E Y_0.
\]
最后 $Y_n\le X_n^++\abs{a}$, $Y_0\ge0$.
\end{proof}

\begin{intuition}
策略 $H$ 是 ``低买高卖'': 价格跌到 $a$ 以下时买入, 涨到 $b$ 以上时卖出. 每完成一次上穿就获利至少 $b-a$. 但在下鞅 (至多 ``略有利'' 的赌博) 中, 任何策略的期望收益都不超过 ``一直持有'' 的收益 $\E Y_n-\E Y_0$. 所以上穿次数的期望受到控制.
\end{intuition}

\begin{theorem}[鞅收敛定理]\label{thm:mart-conv}
设 $X$ 是下鞅, 且 $\sup_n\E X_n^+<\infty$. 则存在 $X_\infty\in L^1$, 使得 $X_n\to X_\infty$ $\as$
\end{theorem}

\begin{proof}
对 $a<b$, 令 $U_\infty[a,b]=\lim_nU_n[a,b]$ (单调不减). 由单调收敛定理与上穿不等式~\ref{thm:upcrossing},
\[
\E\,U_\infty[a,b]=\lim_n\E\,U_n[a,b]\le\frac{\sup_n\E X_n^++\abs{a}}{b-a}<\infty,
\]
所以 $U_\infty[a,b]<\infty$ $\as$ 若 $\liminf_nX_n<a<b<\limsup_nX_n$, 则 $X_n$ 无穷多次低于 $a$ 又无穷多次高于 $b$, 从而 $U_\infty[a,b]=\infty$. 于是
\begin{align*}
\Big\{\liminf_nX_n<\limsup_nX_n\Big\}&=\bigcup_{a<b,\ a,b\in\Q}\Big\{\liminf_nX_n<a<b<\limsup_nX_n\Big\}\\
&\subseteq\bigcup_{a<b,\ a,b\in\Q}\{U_\infty[a,b]=\infty\}
\end{align*}
是可列个零概率集之并, 概率为零. 所以 $X_\infty=\lim_nX_n$ 几乎必然存在 (可能为 $\pm\infty$). 由 Fatou 引理~\ref{thm:fatou}, 并利用 $\abs{X_n}=2X_n^+-X_n$ 与 $\E X_n\ge\E X_0$ (下鞅期望单调不减),
\[
\E\abs{X_\infty}\le\liminf_n\E\abs{X_n}=\liminf_n\big(2\E X_n^+-\E X_n\big)\le2\sup_n\E X_n^+-\E X_0<\infty.
\]
所以 $X_\infty$ 可积, 特别地几乎必然有限.
\end{proof}

由证明中的等式 $\E\abs{X_n}=2\E X_n^+-\E X_n$ 可知, 下鞅的条件 $\sup_n\E X_n^+<\infty$ 等价于 $L^1$ 有界 $\sup_n\E\abs{X_n}<\infty$; 上述定理对鞅的条件也就是 $L^1$ 有界.

\begin{corollary}[非负上鞅收敛定理]\label{cor:nonneg-supermart}
若 $X$ 是非负上鞅, 则 $X_n\to X_\infty$ $\as$, 且 $\E X_\infty\le\E X_0$.
\end{corollary}

\begin{proof}
$-X$ 是下鞅, 且 $(-X_n)^+=0$. 由定理~\ref{thm:mart-conv}, $X_n$ 几乎必然收敛. 由 Fatou 引理与上鞅期望单调不增, $\E X_\infty\le\liminf_n\E X_n\le\E X_0$.
\end{proof}

\begin{remark}
定理只断言几乎必然收敛, 不断言 $L^1$ 收敛. 加倍策略的财富 $W_n=1+(H\cdot X)_n$ (例~\ref{ex:doubling}, 起始本金 $1$) 是非负鞅, $W_n\to2$ $\as$, 但 $\E W_n=1\ne2$. 收敛何时是 $L^1$ 的, 由一致可积性回答 (3.4.3 小节).
\end{remark}

% ---------------------------------------------------------------------
\subsection{应用: 回到经典例子}

\begin{example}[P\'olya 罐子的极限]\label{ex:polya-limit}
红球比例 $X_n$ 是取值于 $[0,1]$ 的鞅, 由定理~\ref{thm:mart-conv}, $X_n\to X_\infty$ $\as$ 下面证明 $X_\infty\sim\mathrm{U}(0,1)$.

先计算 $R_n$ 的分布. 一个含 $k$ 个红、$n-k$ 个蓝的特定抽取序列, 其概率的分母依次为 $2,3,\dots,n+1$; 分子中, 第 $i$ 次抽到红球时罐中有 $i$ 个红球, 第 $i$ 次抽到蓝球时有 $i$ 个蓝球, 与红蓝出现的次序无关. 所以这个序列的概率为 $\frac{k!\,(n-k)!}{(n+1)!}$, 乘以序列个数 $\binom{n}{k}$ 得
\[
\P(R_n=k+1)=\frac{1}{n+1},\qquad k=0,1,\dots,n.
\]
即 $X_n$ 在 $\big\{\frac{1}{n+2},\dots,\frac{n+1}{n+2}\big\}$ 上均匀分布. 于是对 $x\in[0,1]$, $\P(X_n\le x)=\frac{\#\{k: (k+1)/(n+2)\le x\}}{n+1}\to x$. 几乎必然收敛蕴含依分布收敛 (定理~\ref{thm:convergence-implications}), 所以在 $F_{X_\infty}$ 的每个连续点处 $F_{X_\infty}(x)=x$; 连续点稠密, 由右连续性 $F_{X_\infty}(x)=x$ 对一切 $x\in[0,1]$ 成立. 罐子的长期组成是完全随机的: 最初几次抽取的偶然结果被 ``富者愈富'' 的机制放大并永久保留.
\end{example}

\begin{example}[分支过程的灭绝]\label{ex:gw-extinction}
沿用例~\ref{ex:galton-watson} 的记号, 设 $\P(\xi_{1,1}=1)<1$.

\emph{第一步: 种群不会稳定在正数.} 对整数 $k\ge1$, 令 $r_k=\P\big(\sum_{i=1}^k\xi_{1,i}=k\big)$. 则 $r_k<1$: 若 $\P(\xi_{1,1}\ge2)>0$, 则 ``$\xi_{1,1}\ge2$ 且其余 $\xi_{1,i}\ge1$'' 有正概率, 此时和大于 $k$; 否则 $\xi_{1,1}\in\{0,1\}$ 且 $\P(\xi_{1,1}=0)>0$, ``全部为 $0$'' 有正概率. 与例~\ref{ex:galton-watson} 的计算相同, 在 $\{Z_n=k\}$ 上 $\P(Z_{n+1}=k\,|\,\mathcal{F}_n)=r_k$, 所以由全期望公式递推,
$\P(Z_N=Z_{N+1}=\cdots=Z_{N+m}=k)=r_k^m\,\P(Z_N=k)\to0$ ($m\to\infty$). 对可列个 $(N,k)$ 取并, $\P(\text{存在 }k\ge1\text{ 使 }Z_n\text{ 最终恒等于 }k)=0$.

\emph{第二步: $\mu\le1$ 时必然灭绝.} $W_n=Z_n/\mu^n$ 是非负鞅, 由推论~\ref{cor:nonneg-supermart} 几乎必然收敛到有限极限. 当 $\mu\le1$ 时, $Z_n=\mu^nW_n$ 也几乎必然收敛到有限极限. 整数值序列收敛意味着最终恒定, 由第一步, 极限只能是 $0$. 所以 $\P(\text{灭绝})=1$.

\emph{第三步: $\mu>1$ 时的灭绝概率.} 母函数 $f$ 在 $[0,1]$ 上连续, $f(1)=1$, 且左导数 $f'(1-)=\mu>1$ (由单调收敛定理求导), 所以 $s$ 略小于 $1$ 时 $f(s)<s$; 又 $f(0)=\P(\xi_{1,1}=0)\ge0$. 由介值定理, 存在 $q\in[0,1)$ 使 $f(q)=q$. 若 $q=0$, 则每个个体至少有一个后代, 种群不会灭绝. 设 $q\in(0,1)$. 由例~\ref{ex:galton-watson} (2), $q^{Z_n}$ 是取值于 $[0,1]$ 的鞅, 几乎必然收敛; 由于 $s\mapsto q^s$ 严格单调, $Z_n$ 在 $[0,\infty]$ 中几乎必然收敛, 由第一步, 极限 $Z_\infty\in\{0,\infty\}$. 由有界收敛定理 (约定 $q^\infty=0$),
\[
q=\E\,q^{Z_0}=\E\,q^{Z_n}\to\E\,q^{Z_\infty}=\P(Z_\infty=0)=\P(\text{灭绝}).
\]
所以灭绝概率等于 $q$ (这也说明 $f$ 在 $(0,1)$ 中至多有一个不动点), 且在不灭绝的事件上 $Z_n\to\infty$. 至于增长速度, $W_n\to W_\infty$, 而 $W_\infty>0$ 在 ``不灭绝'' 上几乎必然成立当且仅当 $\E[\xi_{1,1}\log^+\xi_{1,1}]<\infty$ (Kesten-Stigum 定理, 见~\cite{durrett} 5.4 节): 此时种群以 $\mu^n$ 的速度增长.
\end{example}

\begin{example}[似然比趋于零]\label{ex:likelihood-limit}
在例~\ref{ex:likelihood-ratio} 中, 设 $g\ne f$ (即二者在正 Lebesgue 测度集上不同). $L_n$ 是非负鞅, 由推论~\ref{cor:nonneg-supermart}, $L_n\to L_\infty$ $\as$ 我们证明 $L_\infty=0$. 令 $\rho=\E\sqrt{g(Y_1)/f(Y_1)}=\int\sqrt{f(y)g(y)}\dd y$. 由 Cauchy-Schwarz 不等式, $\rho\le\big(\int f\big)^{1/2}\big(\int g\big)^{1/2}=1$, 且等号成立当且仅当 $\sqrt{g}$ 与 $\sqrt{f}$ 成比例, 即 $g=f$ 几乎处处; 所以 $\rho<1$. 由独立性, $\E\sqrt{L_n}=\rho^n\to0$. 由 Fatou 引理,
\[
\E\sqrt{L_\infty}\le\liminf_n\E\sqrt{L_n}=0,
\]
故 $L_\infty=0$ $\as$, 即数据真的来自 $f$ 时, 似然比终将压倒性地支持 $f$. 但 $\E L_n=1$ 对一切 $n$ 成立: 期望的 ``质量'' 集中在概率越来越小的事件上 (那些看起来像是来自 $g$ 的数据), 这正是 $L_n$ 不一致可积的表现. 图~\ref{fig:likelihood} 中的 $\log L_n$ 大致以斜率 $\E\log\frac{g(Y_1)}{f(Y_1)}<0$ 下降.
\end{example}

\begin{figure}[htbp]
\centering
\begin{tikzpicture}
\begin{axis}[SAaxis, axis x line=bottom, axis y line=left, clip=false, width=11.5cm, height=6cm, xmin=0, xmax=124, ymin=-22, ymax=8,
  xtick={30,60,90,120}, ytick={-20,-10,0}, xlabel={$n$}]
  \node[font=\small,anchor=south west] at (rel axis cs:0,1.0) {$\log L_n$};
  \addplot[SABlue,thick] coordinates {(0,0.00) (3,-0.30) (6,-1.63) (9,-1.82) (12,-2.26) (15,-1.42) (18,-0.79) (21,-2.08) (24,-2.59) (27,-1.93) (30,-2.75) (33,-4.66) (36,-5.31) (39,-3.64) (42,-4.49) (45,-4.67) (48,-5.68) (51,-5.02) (54,-5.22) (57,-4.53) (60,-4.91) (63,-3.85) (66,-3.88) (69,-3.92) (72,-5.30) (75,-4.13) (78,-3.24) (81,-2.25) (84,-2.16) (87,-1.49) (90,-0.33) (93,-0.17) (96,-0.86) (99,-0.29) (102,0.02) (105,-0.84) (108,-1.28) (111,-1.38) (114,-2.70) (117,-1.01) (120,-1.94)};
  \addplot[SARed,thick] coordinates {(0,0.00) (3,1.24) (6,2.53) (9,1.93) (12,2.75) (15,0.80) (18,0.29) (21,0.18) (24,0.51) (27,-0.91) (30,0.16) (33,-0.99) (36,-0.05) (39,-1.33) (42,-2.32) (45,-1.20) (48,-2.45) (51,-3.37) (54,-4.66) (57,-5.57) (60,-5.57) (63,-7.23) (66,-6.98) (69,-6.43) (72,-8.02) (75,-8.21) (78,-7.38) (81,-8.26) (84,-10.77) (87,-11.31) (90,-10.59) (93,-11.75) (96,-11.44) (99,-10.76) (102,-12.16) (105,-12.74) (108,-13.00) (111,-14.91) (114,-15.24) (117,-15.41) (120,-15.70)};
  \addplot[SAGreenText,thick] coordinates {(0,0.00) (3,0.22) (6,-0.10) (9,-1.59) (12,-0.84) (15,-1.11) (18,-1.59) (21,-3.24) (24,-4.88) (27,-6.00) (30,-5.73) (33,-6.54) (36,-7.23) (39,-9.49) (42,-9.25) (45,-8.85) (48,-10.38) (51,-10.12) (54,-10.77) (57,-11.78) (60,-14.18) (63,-11.90) (66,-13.51) (69,-13.08) (72,-13.08) (75,-14.65) (78,-15.12) (81,-16.14) (84,-15.36) (87,-16.21) (90,-16.87) (93,-17.01) (96,-16.20) (99,-17.07) (102,-17.71) (105,-18.12) (108,-17.07) (111,-18.47) (114,-19.15) (117,-19.78) (120,-20.39)};
  \addplot[SAYellowText,thick] coordinates {(0,0.00) (3,-1.12) (6,-0.83) (9,-0.34) (12,-0.88) (15,-1.37) (18,-0.99) (21,-1.32) (24,-1.94) (27,-1.55) (30,-0.59) (33,-1.16) (36,-1.64) (39,-2.16) (42,-2.72) (45,-3.01) (48,-4.64) (51,-5.35) (54,-2.58) (57,-2.87) (60,-2.23) (63,-2.74) (66,-3.55) (69,-4.01) (72,-5.19) (75,-7.20) (78,-7.34) (81,-8.09) (84,-9.03) (87,-8.14) (90,-9.70) (93,-9.30) (96,-9.24) (99,-10.26) (102,-11.89) (105,-11.57) (108,-12.07) (111,-11.83) (114,-14.00) (117,-15.58) (120,-15.50)};
  \addplot[SADark,dashed,thick,domain=0:120] {-x/8};
  \node[font=\small,anchor=south] (sl) at (axis cs:105,2) {斜率 $-\frac{1}{8}$}; \draw[->,SADark!70] (sl.south) -- (axis cs:105,-12.9);
\end{axis}
\end{tikzpicture}
\caption{似然比的对数 $\log L_n$ 的四条模拟轨道 ($f$ 为 $N(0,1)$ 密度, $g$ 为 $N(\frac{1}{2},1)$ 密度, 数据来自 $f$). 此时 $\log\frac{g}{f}(y)=\frac{y}{2}-\frac{1}{8}$, 期望为 $-\frac{1}{8}$ (虚线). $L_n\to0$ 几乎必然成立, 而 $\E L_n\equiv1$}
\label{fig:likelihood}
\end{figure}

% ---------------------------------------------------------------------
\subsection{一致可积鞅}

几乎必然收敛何时可以加强为 $L^1$ 收敛? 由 Vitali 定理~\ref{thm:vitali}, 答案是一致可积性. 鞅的一致可积性还与 ``鞅由其极限决定'' 等价.

\begin{lemma}\label{lem:ce-family-ui}
设 $X\in L^1$. 则族 $\{\E[X\,|\,\mathcal{G}]: \mathcal{G}\text{ 是 }\mathcal{F}\text{ 的子 }\sigma\text{-代数}\}$ 一致可积.
\end{lemma}

\begin{proof}
给定 $\varepsilon>0$, 由引理~\ref{lem:single-integrable} (2), 取 $\delta>0$ 使 $\P(A)<\delta$ 时 $\E[\abs{X};A]<\varepsilon$. 取 $K>\E\abs{X}/\delta$. 对任意子 $\sigma$-代数 $\mathcal{G}$, 记 $Y=\E[X\,|\,\mathcal{G}]$. 由条件期望的三角不等式~\ref{thm:ce-basic} (3), $\abs{Y}\le\E[\abs{X}\,|\,\mathcal{G}]$, 所以由 Markov 不等式 $\P(\abs{Y}>K)\le\frac{\E\abs{Y}}{K}\le\frac{\E\abs{X}}{K}<\delta$. 由于 $\{\abs{Y}>K\}\in\mathcal{G}$,
\[
\E\big[\abs{Y};\abs{Y}>K\big]\le\E\big[\E[\abs{X}\,|\,\mathcal{G}];\abs{Y}>K\big]=\E\big[\abs{X};\abs{Y}>K\big]<\varepsilon.\qedhere
\]
\end{proof}

\begin{theorem}[一致可积鞅]\label{thm:ui-mart}
设 $M$ 是鞅. 以下四条等价:
\begin{enumerate}
\item $\{M_n\}$ 一致可积;
\item $M_n$ 几乎必然且在 $L^1$ 中收敛;
\item $M_n$ 在 $L^1$ 中收敛;
\item 存在 $X\in L^1$, 使 $M_n=\E[X\,|\,\mathcal{F}_n]$ 对一切 $n$ 成立 (称 $M$ 为\textbf{闭鞅}).
\end{enumerate}
此时, 记 $M_\infty$ 为极限, 则 $M_n=\E[M_\infty\,|\,\mathcal{F}_n]$, 且在 (4) 中 $M_\infty=\E[X\,|\,\mathcal{F}_\infty]$.
\end{theorem}

\begin{proof}
(1)$\Rightarrow$(2): 一致可积蕴含 $L^1$ 有界 (命题~\ref{prop:ui-equivalent}), 由鞅收敛定理~\ref{thm:mart-conv}, $M_n\to M_\infty$ $\as$; 几乎必然收敛蕴含依概率收敛, 由 Vitali 定理~\ref{thm:vitali}, 收敛也是 $L^1$ 的.

(2)$\Rightarrow$(3): 显然.

(3)$\Rightarrow$(4): 设 $M_n\xrightarrow{L^1}M_\infty$. 由定理~\ref{thm:partial-converses} (1), 有子列几乎必然收敛到 $M_\infty$, 故可取 $M_\infty$ 为 $\mathcal{F}_\infty$-可测的. 固定 $n$ 与 $A\in\mathcal{F}_n$. 对 $m\ge n$, 由鞅性质 $\E[M_m;A]=\E[M_n;A]$; 而 $\abs{\E[M_m;A]-\E[M_\infty;A]}\le\E\abs{M_m-M_\infty}\to0$. 所以 $\E[M_\infty;A]=\E[M_n;A]$, 即 $M_n=\E[M_\infty\,|\,\mathcal{F}_n]$. 取 $X=M_\infty$.

(4)$\Rightarrow$(1): 由引理~\ref{lem:ce-family-ui}.

最后, 在 (4) 中, 由 (4)$\Rightarrow$(1)$\Rightarrow$(2), $M_n\to M_\infty$. 要证 $M_\infty=\E[X\,|\,\mathcal{F}_\infty]$. 集族 $\bigcup_n\mathcal{F}_n$ 是包含 $\Omega$ 的 $\pi$-系 (它是递增 $\sigma$-代数之并), 生成 $\mathcal{F}_\infty$. 对 $A\in\mathcal{F}_n$, 由 (3)$\Rightarrow$(4) 的证明, $\E[M_\infty;A]=\E[M_n;A]=\E[X;A]$. 由引理~\ref{lem:ce-pi-system}, $M_\infty=\E[X\,|\,\mathcal{F}_\infty]$.
\end{proof}

\begin{corollary}[L\'evy 向上定理]\label{cor:levy-up}
设 $X\in L^1$. 则 $\E[X\,|\,\mathcal{F}_n]\to\E[X\,|\,\mathcal{F}_\infty]$ 几乎必然且在 $L^1$ 中成立. 特别地, 对 $A\in\mathcal{F}_\infty$, $\P(A\,|\,\mathcal{F}_n)\to\1_A$ $\as$ (L\'evy 0-1 律).
\end{corollary}

\begin{proof}
Doob 鞅 $\E[X\,|\,\mathcal{F}_n]$ 满足定理~\ref{thm:ui-mart} (4). 取 $X=\1_A$, $\E[\1_A\,|\,\mathcal{F}_\infty]=\1_A$.
\end{proof}

\begin{intuition}
随着信息不断积累, 对 $X$ 的预测收敛到 ``用全部信息'' 所作的预测. L\'evy 0-1 律说: 由全部信息决定的事件, 其条件概率最终会趋于 $0$ 或 $1$, 即我们最终会 ``知道'' 它是否发生.
\end{intuition}

\begin{corollary}[Kolmogorov 0-1 律的鞅证明]\label{cor:kolmogorov-01-martingale}
设 $\mathcal{G}_1,\mathcal{G}_2,\dots$ 是相互独立的 $\sigma$-代数, $A$ 属于它们的尾 $\sigma$-代数. 则 $\P(A)\in\{0,1\}$.
\end{corollary}

\begin{proof}
令 $\mathcal{F}_n=\sigma(\mathcal{G}_1\cup\cdots\cup\mathcal{G}_n)$. 由于 $A\in\sigma\big(\bigcup_{k>n}\mathcal{G}_k\big)$, 由分组原理~\ref{cor:grouping}, $A$ 与 $\mathcal{F}_n$ 独立, 故 $\P(A\,|\,\mathcal{F}_n)=\P(A)$ (定理~\ref{thm:ce-basic} (6)). 又 $A\in\mathcal{F}_\infty$, 由 L\'evy 0-1 律, $\P(A)=\P(A\,|\,\mathcal{F}_n)\to\1_A$ $\as$ 常数 $\P(A)$ 几乎必然等于 $\1_A$, 只能是 $0$ 或 $1$.
\end{proof}

对一致可积鞅, 可选停时定理对任意停时成立 (包括可能取 $\infty$ 的停时, 此时约定 $M_\tau=M_\infty$ 在 $\{\tau=\infty\}$ 上).

\begin{theorem}[可选停时定理: 一致可积鞅]\label{thm:ost-ui}
设 $M$ 是一致可积鞅, $\sigma\le\tau$ 是停时. 则 $M_\tau$ 可积, $M_\tau=\E[M_\infty\,|\,\mathcal{F}_\tau]$, 且
\[
\E[M_\tau\,|\,\mathcal{F}_\sigma]=M_\sigma,\qquad\E M_\tau=\E M_0.
\]
\end{theorem}

\begin{proof}
由定理~\ref{thm:ui-mart}, $M_n=\E[M_\infty\,|\,\mathcal{F}_n]$, 且可取 $M_\infty$ 为 $\mathcal{F}_\infty$-可测的.

\emph{可积性与可测性.} 在 $\{\tau=n\}$ 上 $\abs{M_\tau}=\abs{M_n}\le\E[\abs{M_\infty}\,|\,\mathcal{F}_n]$, 而 $\{\tau=n\}\in\mathcal{F}_n$, 所以 $\E[\abs{M_\tau};\tau=n]\le\E[\abs{M_\infty};\tau=n]$ (对 $n=\infty$ 显然). 对 $n\in\{0,1,\dots,\infty\}$ 求和得 $\E\abs{M_\tau}\le\E\abs{M_\infty}<\infty$. 由命题~\ref{prop:F-tau} (3), $M_\tau\1_{\{\tau<\infty\}}$ 是 $\mathcal{F}_\tau$-可测的; 而对 $A\in\mathcal{F}_\infty$, $A\cap\{\tau=\infty\}\in\mathcal{F}_\tau$ (它与 $\{\tau\le n\}$ 之交为空集), 所以 $M_\infty\1_{\{\tau=\infty\}}$ 也是 $\mathcal{F}_\tau$-可测的.

\emph{$M_\tau=\E[M_\infty\,|\,\mathcal{F}_\tau]$.} 设 $A\in\mathcal{F}_\tau$. 对有限的 $n$, $A\cap\{\tau=n\}\in\mathcal{F}_n$, 所以 $\E[M_\infty;A\cap\{\tau=n\}]=\E[M_n;A\cap\{\tau=n\}]=\E[M_\tau;A\cap\{\tau=n\}]$; 对 $n=\infty$ 两边相同. 对 $n$ 求和 (由控制收敛定理, 控制函数分别为 $\abs{M_\infty}$ 与 $\abs{M_\tau}$), $\E[M_\infty;A]=\E[M_\tau;A]$.

\emph{结论.} 由 $\mathcal{F}_\sigma\subseteq\mathcal{F}_\tau$ 与塔性质, $\E[M_\tau\,|\,\mathcal{F}_\sigma]=\E\big[\E[M_\infty\,|\,\mathcal{F}_\tau]\,\big|\,\mathcal{F}_\sigma\big]=\E[M_\infty\,|\,\mathcal{F}_\sigma]=M_\sigma$. 取 $\sigma=0$ 并取期望得 $\E M_\tau=\E M_0$.
\end{proof}

\begin{corollary}\label{cor:ost-bounded-stopped}
设 $M$ 是鞅, $\tau$ 是停时, 停止过程 $M^\tau$ 一致可积 (例如 $\abs{M_{\tau\wedge n}}\le K$). 则 $M_\tau=\lim_nM_{\tau\wedge n}$ 几乎必然存在, 且 $\E M_\tau=\E M_0$.
\end{corollary}

\begin{proof}
由定理~\ref{thm:stopped-martingale}, $M^\tau$ 是一致可积鞅, 由定理~\ref{thm:ui-mart} 它收敛, 极限在 $\{\tau<\infty\}$ 上等于 $M_\tau$. 由定理~\ref{thm:ost-ui} (对鞅 $M^\tau$ 与停时 $\tau$, 或直接由 $L^1$ 收敛), $\E M_\tau=\lim_n\E M_{\tau\wedge n}=\E M_0$.
\end{proof}

例如在例~\ref{ex:gambler-ruin} 的非对称情形中, 停止鞅 $r^{S_{\tau\wedge n}}$ 有界, 推论~\ref{cor:ost-bounded-stopped} 直接给出 $\E r^{S_\tau}=r^x$, 无需事先证明 $\tau<\infty$ (这一点可以由此反推出来).

% ---------------------------------------------------------------------
\subsection{\texorpdfstring{$L^2$}{L2} 有界鞅}

\begin{theorem}[$L^2$ 有界鞅收敛定理]\label{thm:L2-mart}
设 $M$ 是鞅, $\sup_n\E[M_n^2]<\infty$ (由命题~\ref{prop:bracket} (4), 这等价于 $\sum_k\E[(\Delta M_k)^2]<\infty$). 则 $M_n\to M_\infty$ 几乎必然且在 $L^2$ 中成立, 且
\[
\E[M_\infty^2]=\E[M_0^2]+\sum_{k=1}^\infty\E[(\Delta M_k)^2]=\E[M_0^2]+\E\ip{M}_\infty,
\]
\[
\E\Big[\sup_nM_n^2\Big]\le4\sup_n\E[M_n^2].
\]
\end{theorem}

\begin{proof}
记 $C=\sup_n\E[M_n^2]$. 由命题~\ref{prop:ui-sufficient} (2), $M$ 一致可积, 由定理~\ref{thm:ui-mart}, $M_n\to M_\infty$ $\as$ 由 Doob $L^2$ 不等式~\ref{thm:doob-Lp}, $\E[\max_{k\le n}M_k^2]\le4C$; 令 $n\to\infty$, 由单调收敛定理, $\E[\sup_nM_n^2]\le4C$. 于是 $\abs{M_n-M_\infty}^2\le4\sup_kM_k^2$ 是可积控制, 由控制收敛定理, $\E\abs{M_n-M_\infty}^2\to0$. 最后由命题~\ref{prop:bracket} (4), 内积的连续性与单调收敛定理, $\E[M_\infty^2]=\lim_n\E[M_n^2]=\E[M_0^2]+\sum_k\E[(\Delta M_k)^2]$.
\end{proof}

\begin{example}[随机级数]\label{ex:random-series}
设 $\xi_1,\xi_2,\dots$ 独立, $\E\xi_k=0$, $\Var(\xi_k)=1$, 且 $\sum_ka_k^2<\infty$. 则 $\sum_ka_k\xi_k$ \emph{几乎必然}收敛 (例~\ref{ex:L2-series} 只得到了 $L^2$ 收敛). 例如取 $\xi_k=\pm1$ (等概率), 则 ``随机符号的调和级数'' $\sum_k\pm\frac{1}{k}$ 几乎必然收敛, 尽管 $\sum_k\frac{1}{k}=\infty$: 随机符号带来的相消足以使级数收敛.
\end{example}

% ---------------------------------------------------------------------
\subsection{逆向鞅与强大数定律}

最后讨论信息\emph{递减}的情形. 设 $\mathcal{G}_0\supseteq\mathcal{G}_1\supseteq\mathcal{G}_2\supseteq\cdots$ 是递减的子 $\sigma$-代数, $\mathcal{G}_\infty=\bigcap_n\mathcal{G}_n$.

\begin{theorem}[L\'evy 向下定理]\label{thm:levy-down}
设 $X\in L^1$. 则 $\E[X\,|\,\mathcal{G}_n]\to\E[X\,|\,\mathcal{G}_\infty]$ 几乎必然且在 $L^1$ 中成立.
\end{theorem}

\begin{proof}
记 $Y_n=\E[X\,|\,\mathcal{G}_n]$.

\emph{几乎必然收敛.} 固定 $N$, 把时间反过来: 令 $Z_k=Y_{N-k}$, $\mathcal{H}_k=\mathcal{G}_{N-k}$ ($k=0,\dots,N$). 则 $\{\mathcal{H}_k\}$ 递增, 由塔性质 $\E[Z_{k+1}\,|\,\mathcal{H}_k]=\E\big[\E[X\,|\,\mathcal{G}_{N-k-1}]\,\big|\,\mathcal{G}_{N-k}\big]=Z_k$, 所以 $(Z_0,\dots,Z_N)$ 是鞅. 由上穿不等式~\ref{thm:upcrossing}, $Z$ 在 $[0,N]$ 上上穿 $[a,b]$ 的期望次数不超过 $\frac{\E\abs{Z_N}+\abs{a}}{b-a}\le\frac{\E\abs{X}+\abs{a}}{b-a}$, 与 $N$ 无关. $Z$ 的上穿对应于序列 $Y_N,Y_{N-1},\dots,Y_0$ 的上穿, 即序列 $Y_0,Y_1,\dots,Y_N$ 的 ``下穿'' ($Y$ 从 $b$ 以上下降到 $a$ 以下). 令 $N\to\infty$, 由单调收敛定理, 序列 $\{Y_n\}$ 下穿 $[a,b]$ 的总次数的期望有限, 故几乎必然有限. 与定理~\ref{thm:mart-conv} 的证明相同 (一个序列若不收敛, 它必然无穷多次下穿某个有理区间), $Y_n$ 几乎必然收敛到 $[-\infty,\infty]$ 中的极限 $Y_\infty$.

\emph{$L^1$ 收敛与极限的识别.} 由引理~\ref{lem:ce-family-ui}, $\{Y_n\}$ 一致可积. 由 Fatou 引理 $Y_\infty$ 可积, 由 Vitali 定理, $Y_n\xrightarrow{L^1}Y_\infty$. 对每个 $m$, 当 $n\ge m$ 时 $Y_n$ 是 $\mathcal{G}_m$-可测的, 所以 $Y_\infty=\limsup_nY_n$ 是 $\mathcal{G}_m$-可测的; 对一切 $m$ 成立, 故 $Y_\infty$ 是 $\mathcal{G}_\infty$-可测的. 对 $A\in\mathcal{G}_\infty\subseteq\mathcal{G}_n$, $\E[Y_n;A]=\E[X;A]$, 令 $n\to\infty$ 由 $L^1$ 收敛得 $\E[Y_\infty;A]=\E[X;A]$. 所以 $Y_\infty=\E[X\,|\,\mathcal{G}_\infty]$.
\end{proof}

这样的序列 $Y_n=\E[X\,|\,\mathcal{G}_n]$ 称为\textbf{逆向鞅}. 与普通的鞅不同, 逆向鞅\emph{总是}一致可积的, 因而收敛总是 $L^1$ 的. 逆向鞅最著名的应用是强大数定律.

\begin{theorem}[强大数定律]\label{thm:slln}
设 $\xi_1,\xi_2,\dots$ 是 $\iid$ 可积随机变量, $\E\xi_1=\mu$, $S_n=\sum_{k\le n}\xi_k$. 则
\[
\frac{S_n}{n}\to\mu\qquad\text{几乎必然且在 }L^1\text{ 中成立}.
\]
\end{theorem}

\begin{proof}
令 $\mathcal{G}_n=\sigma(S_n,\xi_{n+1},\xi_{n+2},\dots)$. 由于 $S_{n+1}=S_n+\xi_{n+1}$, $\mathcal{G}_{n+1}\subseteq\mathcal{G}_n$, 即 $\{\mathcal{G}_n\}$ 递减.

\emph{第一步: $\E[\xi_1\,|\,\mathcal{G}_n]=\frac{S_n}{n}$.} 由于 $\sigma(\xi_{n+1},\xi_{n+2},\dots)$ 与 $\sigma(\xi_1,S_n)$ 独立 (分组原理), 由命题~\ref{prop:drop-independent}, $\E[\xi_1\,|\,\mathcal{G}_n]=\E[\xi_1\,|\,S_n]$. 对 $k\le n$, $(\xi_k,S_n)$ 与 $(\xi_1,S_n)$ 同分布 ($S_n$ 关于 $\xi_1,\dots,\xi_n$ 对称), 所以对 Borel 集 $B$, $\E[\xi_k;S_n\in B]=\E[\xi_1;S_n\in B]$. 对 $k=1,\dots,n$ 求和, $\E[S_n;S_n\in B]=n\E[\xi_1;S_n\in B]$. 所以 $\E[\xi_1\,|\,S_n]=\frac{S_n}{n}$.

\emph{第二步: 收敛.} 由 L\'evy 向下定理~\ref{thm:levy-down}, $\frac{S_n}{n}\to L=\E[\xi_1\,|\,\mathcal{G}_\infty]$ 几乎必然且在 $L^1$ 中成立.

\emph{第三步: 极限是常数.} 对任意固定的 $k$, $\frac{S_n}{n}-\frac{S_n-S_k}{n}=\frac{S_k}{n}\to0$, 所以 $L=\limsup_n\frac{\xi_{k+1}+\cdots+\xi_n}{n}$ 是 $\sigma(\xi_{k+1},\xi_{k+2},\dots)$-可测的. 对一切 $k$ 成立, 所以 $L$ 关于 $\{\sigma(\xi_k)\}$ 的尾 $\sigma$-代数可测. 由 Kolmogorov 0-1 律~\ref{thm:kolmogorov-01}, 对每个 $c$, $\P(L\le c)\in\{0,1\}$, 所以 $L$ 几乎必然等于常数 $c^*=\inf\{c: \P(L\le c)=1\}$. 最后由 $L^1$ 收敛, $c^*=\E L=\lim_n\E\frac{S_n}{n}=\mu$.
\end{proof}

\begin{remark}
特别地, 在 1.1 节引言的无穷次抛硬币中, 频率 $\frac{S_n}{n}\to\frac{1}{2}$ 几乎必然成立: 事件 \eqref{eq:1-freq-event} 的概率为 $1$. 这个证明的精妙之处在于, 它只假设了一阶矩有限, 而所需的全部 ``估计'' 都由上穿不等式与一致可积性提供.
\end{remark}

% =====================================================================
\section{连续时间鞅}
\label{sec:3-5}

布朗运动及其随机积分都是连续时间的过程. 本节把离散时间鞅论的主要结果推广到连续时间. 推广的基本方法是\textbf{离散化}: 用二进有理时刻上的离散鞅逼近连续时间的鞅, 再利用轨道的右连续性取极限. 这要求对信息流和轨道作一些正则性假设.

以下指标集为 $[0,\infty)$, $\{\mathcal{F}_t\}_{t\ge0}$ 是信息流, $\mathcal{F}_\infty=\sigma\big(\bigcup_t\mathcal{F}_t\big)$.

% ---------------------------------------------------------------------
\subsection{通常条件与轨道正则性}

\begin{definition}[通常条件]\label{def:usual-conditions}
称信息流 $\{\mathcal{F}_t\}$ 满足\textbf{通常条件}, 如果
\begin{enumerate}
\item \textbf{完备}: $\mathcal{F}_0$ 包含 $\mathcal{F}$ 中的全部零概率集 (从而每个 $\mathcal{F}_t$ 都包含);
\item \textbf{右连续}: 对每个 $t\ge0$, $\mathcal{F}_t=\mathcal{F}_{t+}$, 其中 $\mathcal{F}_{t+}=\bigcap_{s>t}\mathcal{F}_s$.
\end{enumerate}
\end{definition}

完备性保证: 若 $X$ 适应, 而 $Y$ 与 $X$ 不可区分, 则 $Y$ 也适应. 右连续性的含义是 ``在时刻 $t$ 之后无穷小的时间内不会突然获得新信息'', 它保证了 ``首次进入开集的时刻'' 是停时 (命题~\ref{prop:hitting-continuous}). 任何信息流都可以通过加入零概率集并取右极限而扩充为满足通常条件的信息流; 对布朗运动, 这样的扩充不改变其鞅性质与马尔可夫性 (第~4~章). \emph{今后在连续时间中, 总假设信息流满足通常条件.}

\begin{definition}[连续时间鞅]\label{def:ct-martingale}
适应于 $\{\mathcal{F}_t\}$ 的可积过程 $X$ 称为鞅 (下鞅、上鞅), 如果对一切 $0\le s\le t$, $\E[X_t\,|\,\mathcal{F}_s]=X_s$ ($\ge$, $\le$).
\end{definition}

\begin{example}\label{ex:ct-martingales}
\begin{enumerate}
\item 若 $X\in L^1$, 则 $M_t=\E[X\,|\,\mathcal{F}_t]$ 是鞅 (与例~\ref{ex:doob-martingale} 相同).
\item 若 $X$ 有独立增量 (即 $X_t-X_s$ 与 $\mathcal{F}_s$ 独立, $s\le t$), 且 $\E X_t=0$, 则 $X$ 是鞅: $\E[X_t-X_s\,|\,\mathcal{F}_s]=\E[X_t-X_s]=0$. 第~4~章将证明标准布朗运动 $B_t$、$B_t^2-t$ 与 $e^{\theta B_t-\theta^2t/2}$ 都是连续鞅. 另一个例子是补偿 Poisson 过程 $N_t-\lambda t$, 它的轨道有跳跃.
\end{enumerate}
\end{example}

对连续时间过程, 我们需要轨道至少是右连续的, 否则可数个时刻上的信息无法决定整条轨道. 若轨道右连续且处处有左极限, 称为 \textbf{c\`adl\`ag} (法语 continue \`a droite, limite \`a gauche) 轨道. 幸运的是, 在通常条件下这总可以做到.

\begin{theorem}[鞅的正则化]\label{thm:regularization}
设信息流满足通常条件. 则每个鞅都有一个 c\`adl\`ag 修正; 更一般地, 下鞅 $X$ 有 c\`adl\`ag 修正当且仅当 $t\mapsto\E X_t$ 右连续. (证明见~\cite{karatzas} 定理 1.3.13.)
\end{theorem}

由于修正有相同的有限维分布 (命题~\ref{prop:modification}), 且在通常条件下修正仍然适应, \emph{今后总取鞅的 c\`adl\`ag 版本}. 由命题~\ref{prop:modification} (2), 两个 c\`adl\`ag 修正是不可区分的, 所以这个版本本质上是唯一的.

% ---------------------------------------------------------------------
\subsection{停时与离散化}

\begin{definition}[连续时间停时]\label{def:ct-stopping-time}
取值于 $[0,\infty]$ 的随机变量 $\tau$ 称为\textbf{停时}, 如果对每个 $t\ge0$, $\{\tau\le t\}\in\mathcal{F}_t$. 定义 $\mathcal{F}_\tau=\{A\in\mathcal{F}_\infty: A\cap\{\tau\le t\}\in\mathcal{F}_t\text{ 对一切 }t\ge0\}$.
\end{definition}

命题~\ref{prop:stopping-operations} 与命题~\ref{prop:F-tau} (1)(2) 对连续时间停时成立, 证明完全相同 (把 $n$ 换成 $t$; 对 $\sigma+\tau$ 的证明需稍作修改, 见~\cite{karatzas} 引理 1.2.9). 若 $X$ 适应且轨道右连续, 则 $X_\tau\1_{\{\tau<\infty\}}$ 也是 $\mathcal{F}_\tau$-可测的 (见~\cite{karatzas} 命题 1.2.18).

\begin{lemma}\label{lem:optional-time}
在右连续的信息流下, $\tau$ 是停时当且仅当对每个 $t>0$, $\{\tau<t\}\in\mathcal{F}_t$.
\end{lemma}

\begin{proof}
若 $\tau$ 是停时, 则 $\{\tau<t\}=\bigcup_n\{\tau\le t-\frac{1}{n}\}\in\mathcal{F}_t$. 反之, 对每个 $t\ge0$ 与 $m$, $\{\tau\le t\}=\bigcap_{n\ge m}\{\tau<t+\frac{1}{n}\}\in\mathcal{F}_{t+1/m}$; 对一切 $m$ 成立, 所以 $\{\tau\le t\}\in\mathcal{F}_{t+}=\mathcal{F}_t$.
\end{proof}

\begin{proposition}[首达时]\label{prop:hitting-continuous}
设 $X$ 是适应过程, 取值于 $\R^d$.
\begin{enumerate}
\item 若 $X$ 的轨道连续, $F$ 是非空闭集, 则 $\tau_F=\inf\{t\ge0: X_t\in F\}$ 是停时.
\item 若 $X$ 的轨道右连续, $G$ 是开集, 则 $\tau_G=\inf\{t\ge0: X_t\in G\}$ 是停时.
\end{enumerate}
\end{proposition}

\begin{proof}
(1) 令 $d(x,F)=\inf_{y\in F}\abs{x-y}$, 它是 $x$ 的连续函数, 且 $d(x,F)=0\iff x\in F$. 记 $Q_t=(\Q\cap[0,t])\cup\{t\}$ (可数集). 我们断言
\[
\{\tau_F\le t\}=\Big\{\inf_{s\in Q_t}d(X_s,F)=0\Big\}.
\]
``$\subseteq$'': 若 $\tau_F\le t$, 由下确界的定义, 存在 $s_k\downarrow\tau_F$ ($s_k\ge\tau_F$) 使 $X_{s_k}\in F$ (或 $X_{\tau_F}\in F$); 由轨道连续与 $F$ 闭, $X_{\tau_F}\in F$. 取 $r_k\in Q_t$, $r_k\to\tau_F$, 则 $d(X_{r_k},F)\to d(X_{\tau_F},F)=0$. ``$\supseteq$'': 若下确界为零, 取 $s_k\in Q_t$ 使 $d(X_{s_k},F)\to0$; 由 $[0,t]$ 的紧性, 有子列 $s_k\to s\in[0,t]$, 由连续性 $d(X_s,F)=0$, 即 $X_s\in F$, 所以 $\tau_F\le s\le t$. 右端是可数个 $\mathcal{F}_t$-可测随机变量的下确界等于零的事件, 属于 $\mathcal{F}_t$.

(2) 我们断言 $\{\tau_G<t\}=\bigcup_{s\in\Q\cap[0,t)}\{X_s\in G\}$. ``$\supseteq$'' 是显然的. ``$\subseteq$'': 若 $\tau_G<t$, 存在 $u\in[\tau_G,t)$ 使 $X_u\in G$; 由于 $G$ 开且轨道右连续, 存在 $\delta>0$ 使 $X_s\in G$ 对 $s\in[u,u+\delta)$ 成立, 在其中取有理数 $s<t$ 即可. 于是 $\{\tau_G<t\}\in\mathcal{F}_t$, 由引理~\ref{lem:optional-time}, $\tau_G$ 是停时.
\end{proof}

连续时间结果的证明都依赖于下面的离散化.

\begin{lemma}[停时的二进离散化]\label{lem:dyadic-approx}
设 $\tau$ 是停时. 对 $n\ge1$, 令
\[
\tau_n=\frac{\lfloor2^n\tau\rfloor+1}{2^n}\qquad(\tau=\infty\text{ 时 }\tau_n=\infty).
\]
则 $\tau_n$ 是取值于 $D_n=\{k2^{-n}: k\ge1\}\cup\{\infty\}$ 的停时, $\tau_n>\tau$, 且 $\tau_n\downarrow\tau$ ($n\to\infty$).
\end{lemma}

\begin{proof}
$\tau_n$ 是严格大于 $\tau$ 的最小的 $2^{-n}$ 整数倍, 所以 $\tau<\tau_n\le\tau+2^{-n}$, 且 $D_n\subseteq D_{n+1}$ 蕴含 $\tau_{n+1}\le\tau_n$. 对 $k\ge1$, $\{\tau_n=k2^{-n}\}=\{(k-1)2^{-n}\le\tau<k2^{-n}\}\in\mathcal{F}_{k2^{-n}}$ (由引理~\ref{lem:optional-time} 的证明, $\{\tau<t\}\in\mathcal{F}_t$). 由于 $\tau_n$ 只取 $D_n$ 中的值, 对任意 $t$, $\{\tau_n\le t\}=\bigcup_{k2^{-n}\le t}\{\tau_n=k2^{-n}\}\in\mathcal{F}_t$.
\end{proof}

\begin{figure}[htbp]
\centering
\begin{tikzpicture}
  \draw[->,SADark!80] (-0.6,0) -- (12.8,0) node[right,font=\small] {$t$};
  \draw[SADark!45] (0.000,-0.1) -- (0.000,0.1);
  \draw[SADark!45] (1.500,-0.1) -- (1.500,0.1);
  \draw[SADark!45] (3.000,-0.1) -- (3.000,0.1);
  \draw[SADark!45] (4.500,-0.1) -- (4.500,0.1);
  \draw[SADark!45] (6.000,-0.1) -- (6.000,0.1);
  \draw[SADark!45] (7.500,-0.1) -- (7.500,0.1);
  \draw[SADark!45] (9.000,-0.1) -- (9.000,0.1);
  \draw[SADark!45] (10.500,-0.1) -- (10.500,0.1);
  \draw[SADark!45] (12.000,-0.1) -- (12.000,0.1);
  \draw[SADark,thick] (0.000,-0.2) -- (0.000,0.2);
  \draw[SADark,thick] (6.000,-0.2) -- (6.000,0.2);
  \draw[SADark,thick] (12.000,-0.2) -- (12.000,0.2);
  \node[below,font=\small] at (0.000,-0.22) {$\tfrac{1}{4}$};
  \node[below,font=\small] at (6.000,-0.22) {$\tfrac{1}{2}$};
  \node[below,font=\small] at (12.000,-0.22) {$\tfrac{3}{4}$};
  \node[below,font=\small,SADark!70] at (4.500,-0.22) {$\tfrac{7}{16}$};
  \fill[SARed] (3.600,0) circle (2pt);
  \draw[SARed,dashed] (3.600,0) -- (3.600,2.0) node[above,font=\small] {$\tau$};
  \draw[SAGreenText,dashed] (4.500,0) -- (4.500,0.7);
  \draw[->,SAGreenText,thick] (3.600,0.7) -- (4.500,0.7) node[right,font=\small] {$\tau_4=\tfrac{7}{16}$};
  \draw[SABlue,dashed] (6.000,0) -- (6.000,1.4);
  \draw[->,SABlue,thick] (3.600,1.4) -- (6.000,1.4) node[right,font=\small] {$\tau_2=\tau_3=\tfrac{1}{2}$};
\end{tikzpicture}
\caption{停时的二进离散化: $\tau_n$ 是严格大于 $\tau$ 的第一个 $2^{-n}$ 的整数倍. 图中 $\tau=0.4$, 粗刻度为 $\frac{1}{4}$ 的倍数, 细刻度为 $\frac{1}{16}$ 的倍数. 随着 $n$ 增大, $\tau_n$ 从右侧单调下降到 $\tau$. 由于 $\tau_n$ 只用到 ``$\tau$ 是否已经发生'', 它仍是停时; 而向下取整则会用到未来的信息}
\label{fig:dyadic-approx}
\end{figure}

% ---------------------------------------------------------------------
\subsection{Doob 不等式、收敛定理与可选停时定理}

\begin{theorem}[连续时间 Doob 不等式]\label{thm:ct-doob}
设 $X$ 是轨道右连续的下鞅, $t>0$, $\lambda>0$. 则
\[
\lambda\,\P\Big(\sup_{s\le t}X_s\ge\lambda\Big)\le\E X_t^+.
\]
若 $X$ 非负, $p>1$, 则 $\E\big[\sup_{s\le t}X_s^p\big]\le\big(\frac{p}{p-1}\big)^p\E X_t^p$. 特别地, 对轨道右连续的鞅 $M$, $\E\big[\sup_{s\le t}M_s^2\big]\le4\E M_t^2$.
\end{theorem}

\begin{proof}
令 $E_n=\{kt2^{-n}: 0\le k\le2^n\}$, 它们是递增的有限集, 其并 $E$ 在 $[0,t]$ 中稠密且包含 $t$. 把 $(X_s)_{s\in E_n}$ 看作 (按时间顺序排列的) 离散时间下鞅, 由 Doob 最大值不等式~\ref{thm:doob-maximal}, 对任意 $\mu>0$, $\mu\P(\max_{s\in E_n}X_s>\mu)\le\E X_t^+$. 由轨道的右连续性, 对每个 $s\in[0,t)$, 存在 $E$ 中的点从右侧逼近 $s$, 所以 $\sup_{s\le t}X_s=\sup_{s\in E}X_s$. 事件 $\{\max_{s\in E_n}X_s>\mu\}$ 关于 $n$ 递增, 其并为 $\{\sup_{s\le t}X_s>\mu\}$, 由概率的连续性, $\mu\P(\sup_{s\le t}X_s>\mu)\le\E X_t^+$. 取 $\mu=\lambda-\varepsilon$ 并令 $\varepsilon\downarrow0$, 由于 $\{\sup_{s\le t}X_s>\lambda-\varepsilon\}\downarrow\{\sup_{s\le t}X_s\ge\lambda\}$, 即得第一个不等式. $L^p$ 不等式同理: 对 $E_n$ 用定理~\ref{thm:doob-Lp}, 再由单调收敛定理令 $n\to\infty$.
\end{proof}

\begin{theorem}[连续时间鞅收敛定理]\label{thm:ct-convergence}
设 $X$ 是轨道右连续的下鞅, 且 $\sup_t\E X_t^+<\infty$. 则 $X_t\to X_\infty$ 几乎必然成立 ($t\to\infty$), 且 $X_\infty\in L^1$.
\end{theorem}

\begin{proof}
记 $C=\sup_t\E X_t^+$, $D=\{k2^{-m}: k,m\ge0\}$ 为非负二进有理数全体. 对 $D$ 的有限子集 $F$ 与 $a<b$, 记 $U_F[a,b]$ 为 $(X_s)_{s\in F}$ (按时间排序) 上穿 $[a,b]$ 的次数. 由上穿不等式~\ref{thm:upcrossing} (应用于有限个时刻构成的离散下鞅), $(b-a)\E U_F[a,b]\le C+\abs{a}$. 取有限集 $F_m=\{k2^{-m}: 0\le k\le m2^m\}\uparrow D$, 则 $U_{F_m}[a,b]$ 单调不减, 其极限 $U_D[a,b]$ 满足 (由单调收敛定理) $\E U_D[a,b]\le\frac{C+\abs{a}}{b-a}$, 故几乎必然有限. 与定理~\ref{thm:mart-conv} 的证明相同, 在一个零概率集之外, 对一切有理数 $a<b$, $U_D[a,b]<\infty$, 从而极限 $L=\lim_{s\to\infty,\,s\in D}X_s$ 存在于 $[-\infty,\infty]$.

现在固定这个零概率集之外的 $\omega$. 由轨道在每个 $t$ 处右连续, 可取 $s(t)\in D\cap(t,t+1)$ 使 $\abs{X_{s(t)}(\omega)-X_t(\omega)}<\frac{1}{t}$. 当 $t\to\infty$ 时 $s(t)\to\infty$, 所以 $X_{s(t)}(\omega)\to L(\omega)$, 从而 $X_t(\omega)\to L(\omega)$. 最后与定理~\ref{thm:mart-conv} 相同, 由 Fatou 引理 (沿整数时刻) $\E\abs{L}\le2C-\E X_0<\infty$.
\end{proof}

与离散时间相同 (证明完全一样), 对轨道右连续的鞅 $M$, 一致可积等价于 $L^1$ 收敛, 也等价于存在 $X\in L^1$ 使 $M_t=\E[X\,|\,\mathcal{F}_t]$; 此时 $M_t=\E[M_\infty\,|\,\mathcal{F}_t]$.

\begin{theorem}[连续时间可选停时定理]\label{thm:ct-ost}
设 $M$ 是轨道右连续的一致可积鞅, $\sigma\le\tau$ 是停时 (在 $\{\tau=\infty\}$ 上令 $M_\tau=M_\infty$). 则 $M_\tau$ 可积, 且
\[
\E[M_\tau\,|\,\mathcal{F}_\sigma]=M_\sigma,\qquad\E M_\tau=\E M_0.
\]
\end{theorem}

\begin{proof}
设 $\sigma_n,\tau_n$ 是引理~\ref{lem:dyadic-approx} 中的离散化. 由于离散化是单调的, $\sigma_n\le\tau_n$.

\emph{离散情形.} 固定 $n$, 考虑离散时间鞅 $(M_{k2^{-n}})_{k\ge0}$ 与信息流 $(\mathcal{F}_{k2^{-n}})_{k\ge0}$. 它是一致可积的 ($M_{k2^{-n}}=\E[M_\infty\,|\,\mathcal{F}_{k2^{-n}}]$), 而 $2^n\sigma_n$, $2^n\tau_n$ 是关于这个离散信息流的停时. 设 $A\in\mathcal{F}_\sigma$. 则 $A$ 也属于离散意义下的 $\mathcal{F}_{\sigma_n}$: 对每个 $k$,
$A\cap\{\sigma_n\le k2^{-n}\}=A\cap\{\sigma<k2^{-n}\}=\bigcup_m\big(A\cap\{\sigma\le k2^{-n}-\tfrac{1}{m}\}\big)\in\mathcal{F}_{k2^{-n}}$. 由离散的可选停时定理~\ref{thm:ost-ui},
\[
\E[M_{\tau_n};A]=\E[M_{\sigma_n};A],\qquad M_{\tau_n}=\E[M_\infty\,|\,\mathcal{F}_{\tau_n}].
\]

\emph{取极限.} 由引理~\ref{lem:ce-family-ui}, $\{M_{\tau_n}\}_n$ 与 $\{M_{\sigma_n}\}_n$ 都一致可积. 由于 $\tau_n\downarrow\tau$ 且轨道右连续, $M_{\tau_n}\to M_\tau$ (在 $\{\tau=\infty\}$ 上两边都是 $M_\infty$); 同理 $M_{\sigma_n}\to M_\sigma$. 由 Vitali 定理~\ref{thm:vitali}, 收敛也是 $L^1$ 的, 所以 $M_\tau,M_\sigma$ 可积, 且 $\E[M_\tau;A]=\E[M_\sigma;A]$. 由于 $M_\sigma$ 是 $\mathcal{F}_\sigma$-可测的, 这正是 $\E[M_\tau\,|\,\mathcal{F}_\sigma]=M_\sigma$.
\end{proof}

\begin{corollary}\label{cor:ct-ost-bounded}
设 $M$ 是轨道右连续的鞅, $\tau$ 是停时.
\begin{enumerate}
\item 若 $\sigma\le\tau\le T$ 是有界停时, 则 $\E[M_\tau\,|\,\mathcal{F}_\sigma]=M_\sigma$, $\E M_\tau=\E M_0$.
\item 停止过程 $M^\tau_t=M_{\tau\wedge t}$ 是鞅.
\end{enumerate}
\end{corollary}

\begin{proof}
(1) 过程 $N_t=M_{t\wedge T}$ 是一致可积鞅: $N_t=\E[M_T\,|\,\mathcal{F}_{t\wedge T}]$ (由鞅性质), 由引理~\ref{lem:ce-family-ui} 一致可积. 对 $N$ 用定理~\ref{thm:ct-ost}, 并注意 $N_\tau=M_\tau$, $N_\sigma=M_\sigma$.

(2) 设 $s\le t$. 把 $M_{\tau\wedge t}$ 拆成两部分. 在 $\{\tau\le s\}$ 上 $M_{\tau\wedge t}=M_{\tau\wedge s}$. 对另一部分, 考虑有界停时 $s\le(\tau\vee s)\wedge t$; 由 (1),
$\E\big[M_{(\tau\vee s)\wedge t}\,\big|\,\mathcal{F}_s\big]=M_s$. 而 $M_{(\tau\vee s)\wedge t}=M_{\tau\wedge t}\1_{\{\tau>s\}}+M_s\1_{\{\tau\le s\}}$, 其中第二项是 $\mathcal{F}_s$-可测的, 所以
$\E[M_{\tau\wedge t}\1_{\{\tau>s\}}\,|\,\mathcal{F}_s]=M_s-M_s\1_{\{\tau\le s\}}=M_s\1_{\{\tau>s\}}$. 由于 $M_{\tau\wedge s}\1_{\{\tau\le s\}}$ 是 $\mathcal{F}_s$-可测的 ($\mathcal{F}_{\tau\wedge s}\subseteq\mathcal{F}_s$),
\[
\E[M_{\tau\wedge t}\,|\,\mathcal{F}_s]=M_{\tau\wedge s}\1_{\{\tau\le s\}}+M_s\1_{\{\tau>s\}}=M_{\tau\wedge s}.\qedhere
\]
\end{proof}

% ---------------------------------------------------------------------
\subsection{局部鞅}

随机积分 $\int_0^tH_s\dd B_s$ 只在被积过程满足较强的可积条件时才是鞅; 一般情形下它只是 ``局部地'' 是鞅. 这引出了局部鞅的概念, 它是 Itô 积分理论的自然框架 (第~5~章).

\begin{definition}[局部鞅]\label{def:local-martingale}
设 $M$ 是适应的连续过程. 若存在停时序列 $\tau_n\uparrow\infty$ $\as$, 使得每个停止过程 $M^{\tau_n}-M_0$ 都是鞅, 称 $M$ 为\textbf{连续局部鞅}, 称 $\{\tau_n\}$ 为它的\textbf{局部化序列}.
\end{definition}

\begin{proposition}\label{prop:local-martingale}
设 $M$ 是连续局部鞅.
\begin{enumerate}
\item 每个连续鞅都是连续局部鞅.
\item 若 $M_0=0$, 则 $\sigma_n=\inf\{t: \abs{M_t}\ge n\}$ 是局部化序列, 且每个 $M^{\sigma_n}$ 是有界鞅.
\item 若存在常数 $K$ 使 $\abs{M_t}\le K$ 对一切 $t$ 成立, 则 $M$ 是鞅.
\item 若 $M\ge0$ 且 $M_0$ 可积, 则 $M$ 是上鞅.
\end{enumerate}
\end{proposition}

\begin{proof}
(1) 取 $\tau_n=n$, 由推论~\ref{cor:ct-ost-bounded} (2).

(2) 由命题~\ref{prop:hitting-continuous} (1), $\sigma_n$ 是停时 (闭集 $\{\abs{x}\ge n\}$ 的首达时); 由轨道连续, 在任意有限区间上 $M$ 有界, 所以 $\sigma_n\uparrow\infty$. 设 $\{\tau_k\}$ 是一个局部化序列. 由推论~\ref{cor:ct-ost-bounded} (2), 鞅 $M^{\tau_k}$ 在 $\sigma_n$ 处停止后仍是鞅, 即 $M^{\sigma_n\wedge\tau_k}$ 是鞅. 由连续性与 $M_0=0$, $\abs{M^{\sigma_n\wedge\tau_k}_t}\le n$. 令 $k\to\infty$, $M^{\sigma_n\wedge\tau_k}_t\to M^{\sigma_n}_t$, 由条件控制收敛定理~\ref{thm:cdct}, 对 $s\le t$, $\E[M^{\sigma_n}_t\,|\,\mathcal{F}_s]=\lim_k\E[M^{\sigma_n\wedge\tau_k}_t\,|\,\mathcal{F}_s]=\lim_kM^{\sigma_n\wedge\tau_k}_s=M^{\sigma_n}_s$.

(3) 与 (2) 的最后一步相同: 对 $s\le t$, $\E[M_{\tau_k\wedge t}\,|\,\mathcal{F}_s]=M_{\tau_k\wedge s}$, 令 $k\to\infty$, 由条件控制收敛定理 (控制函数为常数 $2K$) 得 $\E[M_t\,|\,\mathcal{F}_s]=M_s$.

(4) 设 $\{\tau_k\}$ 是局部化序列, 则 $M^{\tau_k}$ 是鞅 (因为 $M^{\tau_k}-M_0$ 是鞅而 $M_0$ 是 $\mathcal{F}_0$-可测且可积的), 特别地 $M_{\tau_k\wedge t}$ 可积, 且 $M_{\tau_k\wedge t}\to M_t$ ($k\to\infty$).

\emph{第一步: $M_t$ 可积.} 取 $K>0$. 序列 $M_{\tau_k\wedge t}\wedge K$ 非负、有界, 其极限 $M_t\wedge K$ 也有界. 由条件 Fatou 引理~\ref{thm:cfatou} 与鞅性质,
\[
\E[M_t\wedge K\,|\,\mathcal{F}_0]\le\liminf_k\E[M_{\tau_k\wedge t}\wedge K\,|\,\mathcal{F}_0]\le\liminf_k\E[M_{\tau_k\wedge t}\,|\,\mathcal{F}_0]=M_0.
\]
取期望得 $\E[M_t\wedge K]\le\E M_0$; 令 $K\to\infty$, 由单调收敛定理 $\E M_t\le\E M_0<\infty$.

\emph{第二步: 上鞅性质.} 现在 $M_t=\liminf_kM_{\tau_k\wedge t}$ 可积, 对 $s\le t$ 直接用条件 Fatou 引理,
\[
\E[M_t\,|\,\mathcal{F}_s]\le\liminf_k\E[M_{\tau_k\wedge t}\,|\,\mathcal{F}_s]=\liminf_kM_{\tau_k\wedge s}=M_s.
\]
\end{proof}

\begin{remark}
局部鞅未必是鞅, 即使它可积. 经典的例子是: 设 $\rvec{B}$ 是从 $(1,0,0)$ 出发的三维布朗运动, 则 $M_t=1/\abs{\rvec{B}_t}$ 是连续局部鞅, 且 $\sup_t\E M_t^2<\infty$, 但 $\E M_t$ 严格递减, 所以它不是鞅 (见~\cite{revuz} 第 V 章). 这类 ``严格局部鞅'' 在金融数学中对应于资产价格泡沫. 判断一个局部鞅是否为真鞅是第~5~章 Novikov 条件要解决的问题.
\end{remark}

% ---------------------------------------------------------------------
\subsection{二次变差}

离散时间中, 平方可积鞅 $M$ 的二次变差 $[M]_n=\sum_k(\Delta M_k)^2$ 是使 $M^2-[M]$ 成为鞅的过程 (命题~\ref{prop:bracket}). 连续时间中, 我们用越来越细的分割上的平方和来定义它. 先看一个关键事实: 非常数的连续鞅不可能有有限的变差.

\begin{definition}[变差]
函数 $f:[0,t]\to\R$ 的 (全) \textbf{变差}定义为 $V_t(f)=\sup_\Pi\sum_i\abs{f(t_i)-f(t_{i-1})}$, 上确界取遍 $[0,t]$ 的一切分割 $\Pi: 0=t_0<t_1<\cdots<t_m=t$. 若 $V_t(f)<\infty$ 对一切 $t$ 成立, 称 $f$ 有有限变差.
\end{definition}

单调函数与连续可微函数都有有限变差. 有限变差函数正是可以作 Riemann-Stieltjes 积分 $\int g\dd f$ 的积分器.

\begin{theorem}[连续鞅没有有限变差]\label{thm:no-finite-variation}
设 $M$ 是连续鞅, $M_0=0$, 且存在常数 $K$ 使 $\abs{M_t}\le K$ 与 $V_t(M)\le K$ 对一切 $t\le T$ 几乎必然成立. 则 $M_t=0$ 对一切 $t\le T$ 几乎必然成立.
\end{theorem}

\begin{proof}
固定 $t\le T$ 与分割 $\Pi: 0=t_0<\cdots<t_m=t$, 记 $\Delta_i=M_{t_i}-M_{t_{i-1}}$. 与命题~\ref{prop:bracket} (1) 相同, 鞅增量两两正交, 所以
\[
\E[M_t^2]=\E\Big[\Big(\sum_i\Delta_i\Big)^2\Big]=\E\Big[\sum_i\Delta_i^2\Big].
\]
而 $\sum_i\Delta_i^2\le\max_i\abs{\Delta_i}\cdot\sum_i\abs{\Delta_i}\le K\max_i\abs{\Delta_i}$. 当分割的步长 $\max_i(t_i-t_{i-1})\to0$ 时, 由轨道在 $[0,t]$ 上一致连续, $\max_i\abs{\Delta_i}\to0$; 又 $K\max_i\abs{\Delta_i}\le2K^2$, 由有界收敛定理, $\E[M_t^2]=\lim\E\big[\sum_i\Delta_i^2\big]=0$. 所以对每个 $t\le T$, $M_t=0$ $\as$; 取 $t$ 为有理数并利用轨道连续, 得 $M_t=0$ 对一切 $t\le T$ 几乎必然成立.
\end{proof}

通过局部化 (用 $\inf\{t: \abs{M_t}+V_t(M)\ge n\}$ 截断), 这一结论推广到一切连续局部鞅: \emph{有限变差的连续局部鞅是常数} (见~\cite{karatzas} 引理 1.5.10).

\begin{intuition}
这个定理是随机微积分存在的理由. 布朗运动是连续鞅, 所以它的轨道没有有限变差, Riemann-Stieltjes 积分 $\int H\dd B$ 无法逐条轨道定义. 证明还指出了替代的方向: 变差虽然无穷, 平方和 $\sum_i\Delta_i^2$ 却有有限极限. 下面的定理说这个极限总是存在的.
\end{intuition}

\begin{theorem}[连续局部鞅的二次变差]\label{thm:quadratic-variation}
设 $M$ 是连续局部鞅. 则存在唯一的 (在不可区分的意义下) 连续、适应、单调不减的过程 $\ip{M}$, $\ip{M}_0=0$, 使得 $M^2-\ip{M}$ 是连续局部鞅. 并且对任意 $t$ 与 $[0,t]$ 的分割序列 $\Pi_n$, 只要步长趋于零,
\[
\sum_{i}\big(M_{t_i}-M_{t_{i-1}}\big)^2\xrightarrow{\P}\ip{M}_t.
\]
若 $M$ 是平方可积的连续鞅, 则 $M^2-\ip{M}$ 是鞅, 从而 $\E[M_t^2]=\E[M_0^2]+\E\ip{M}_t$. (证明见~\cite{karatzas} 定理 1.5.8 与定理 1.5.13, 或~\cite{revuz} 第 IV 章.)
\end{theorem}

$\ip{M}$ 称为 $M$ 的\textbf{二次变差过程}. 它是离散时间中 $\ip{M}$ 与 $[M]$ 的共同推广: 对连续鞅, 可料二次变差与 ``实现的'' 二次变差合而为一. 唯一性部分是定理~\ref{thm:no-finite-variation} 的直接推论: 若 $A,A'$ 都满足要求, 则 $A-A'=(M^2-A')-(M^2-A)$ 是有限变差的连续局部鞅, 因而为常数零.

\begin{intuition}
$\ip{M}_t$ 是鞅 $M$ 的 ``内在时钟'': 它度量到时刻 $t$ 为止 $M$ 累积了多少随机性. 第~4~章将直接计算出布朗运动的二次变差 $\ip{B}_t=t$ (不依赖于本定理): 布朗运动的内在时钟就是普通的时间. 第~5~章的 L\'evy 刻画定理将说明这一性质完全刻画了布朗运动; 更一般地, 每个连续局部鞅都是一个按其二次变差进行时间变换的布朗运动 (Dambis-Dubins-Schwarz 定理). 在这个意义上, 布朗运动是连续鞅的 ``原子''.
\end{intuition}

% ---------------------------------------------------------------------
\subsection*{本章小结}
\addcontentsline{toc}{subsection}{本章小结}

\begin{center}
\captionof{table}{鞅论的主要结论}
\label{tab:ch3-summary}
\renewcommand{\arraystretch}{1.5}
\begin{tabular}{>{\raggedright\arraybackslash}p{0.26\textwidth} >{\raggedright\arraybackslash}p{0.44\textwidth} >{\raggedright\arraybackslash}p{0.20\textwidth}}
\toprule
结论 & 内容 & 后续应用 \\
\midrule
鞅变换定理 & 可料有界策略不改变公平性 & Itô 积分是鞅 (第~5~章) \\
可选停时定理 & 有界停时或一致可积时 $\E M_\tau=\E M_0$ & 布朗运动首达时 (第~4~章) \\
Doob 分解, $\ip{M}$ & 下鞅 = 鞅 + 可料增过程 & Itô 等距, 二次变差 \\
Doob 不等式 & $\E\sup_{s\le t}M_s^2\le4\E M_t^2$ & 构造 Itô 积分, SDE 解的存在性 \\
鞅收敛定理 & $L^1$ 有界 $\Rightarrow$ 几乎必然收敛 & 鞅表示, 测度变换 \\
一致可积鞅 & $L^1$ 收敛 $\iff$ 闭鞅 & Girsanov 定理 (第~5~章) \\
连续鞅无有限变差 & 需要二次变差 $\ip{M}$ 代替 & Itô 公式的修正项 \\
\bottomrule
\end{tabular}
\end{center}

% =====================================================================
%  第 4 章  布朗运动
% =====================================================================
\chapter{布朗运动}
\label{ch:4}

1827 年, 植物学家 Brown 在显微镜下观察到悬浮在水中的花粉微粒在做永不停息的无规则运动. 1905 年, Einstein 解释说这是微粒受到周围水分子大量微小且相互独立的随机碰撞的结果, 并由此推出微粒位移的方差与时间成正比; 稍早的 1900 年, Bachelier 在博士论文中用同样的模型描述巴黎交易所的股价. 1923 年, Wiener 第一次在数学上严格构造了这个过程, 所以它也称为 \textbf{Wiener 过程}.

布朗运动在随机分析中的地位, 相当于正态分布在概率论中的地位. 正态分布是独立小量之和的普适极限 (中心极限定理), 布朗运动则是随机游动在时间与空间同时缩放下的普适极限 (4.1 节). 它同时属于前三章讨论过的每一类对象: 它是高斯过程 (第~2~章), 是马尔可夫过程 (第~2~章), 也是连续鞅 (第~3~章). 正因为三者兼备, 布朗运动成为构造一切连续随机模型的基本 ``原料''. 第~5~章用它作积分器来定义随机积分, 第~6~章的随机微分方程则由它提供驱动噪声.

本章的结构如下. 4.1 节从随机游动出发, 说明布朗运动应该是什么样子; 4.2 节给出定义并用 Kolmogorov 连续性定理证明它的存在; 4.3 节讨论它的不变性与轨道性质, 核心结论是二次变差等于 $t$ 而轨道处处不可微; 4.4 节研究鞅性质与 (强) 马尔可夫性; 4.5 节用这些工具计算首达时、最大值、最后零点的分布, 并介绍由布朗运动派生的几类重要过程, 最后证明重对数律.

% =====================================================================
\section{从随机游动到布朗运动}
\label{sec:4-1}

% ---------------------------------------------------------------------
\subsection{随机游动的缩放}

设 $\xi_1,\xi_2,\dots$ 是 $\iid$ 随机变量, $\E\xi_1=0$, $\Var(\xi_1)=1$ (例如 $\P(\xi_1=\pm1)=\frac{1}{2}$), $S_0=0$, $S_n=\xi_1+\cdots+\xi_n$. 我们想在时间与空间上同时 ``拉远镜头'', 使得在一个单位时间内走 $n$ 步. 时间的缩放是自然的: 第 $k$ 步发生在时刻 $k/n$. 空间应该如何缩放? 若令
\[
X^{(n)}(t)=\frac{S_{\lfloor nt\rfloor}}{n^{\alpha}},\qquad t\ge0,
\]
则 $\Var\big(X^{(n)}(t)\big)=\lfloor nt\rfloor/n^{2\alpha}$. 当 $n\to\infty$ 时, 若 $\alpha>\frac{1}{2}$, 方差趋于 $0$, 极限是平凡的零过程 (这正是弱大数定律); 若 $\alpha<\frac{1}{2}$, 方差趋于无穷, 没有极限. 只有 $\alpha=\frac{1}{2}$ 时方差收敛到 $t$, 极限既不平凡也不发散. 于是令
\begin{equation}\label{eq:scaled-rw}
W^{(n)}(t)=\frac{S_{\lfloor nt\rfloor}}{\sqrt{n}},\qquad t\ge0.
\end{equation}
$W^{(n)}$ 是阶梯函数: 在每个区间 $[\frac{k}{n},\frac{k+1}{n})$ 上取常值 $S_k/\sqrt{n}$, 每步跳跃的大小为 $\abs{\xi_k}/\sqrt{n}$. 图~\ref{fig:scaled-rw} 画出了 $n=10,100,1000$ 时的一条轨道 (为便于观察, 相邻的点用线段相连). 随着 $n$ 增大, 跳跃越来越小, 轨道看起来越来越像一条连续而极不规则的曲线.

\begin{figure}[htbp]
\centering
\begin{tikzpicture}
\begin{axis}[SAaxis, axis x line=bottom, axis y line=left, width=5.6cm, height=4.6cm, xmin=0, xmax=1.02, ymin=-1.2, ymax=1.0,
  xtick={0,0.5,1}, ytick={-1,0,1}, title={\small $n=10$}, xlabel={$t$}, clip=false]
  \addplot[SADark!20] coordinates {(0,0) (1,0)};
  \addplot[SABlue,thick,mark=*,mark size=1.2pt] coordinates {(0.000,0.000) (0.100,0.316) (0.200,0.632) (0.300,0.316) (0.400,0.632) (0.500,0.316) (0.600,0.000) (0.700,-0.316) (0.800,-0.632) (0.900,-0.316) (1.000,0.000)};
\end{axis}
\end{tikzpicture}\hfill
\begin{tikzpicture}
\begin{axis}[SAaxis, axis x line=bottom, axis y line=left, width=5.6cm, height=4.6cm, xmin=0, xmax=1.02, ymin=-1.2, ymax=1.0,
  xtick={0,0.5,1}, ytick={-1,0,1}, title={\small $n=100$}, xlabel={$t$}]
  \addplot[SADark!20] coordinates {(0,0) (1,0)};
  \addplot[SABlue,semithick] coordinates {(0.000,0.000) (0.010,-0.100) (0.020,0.000) (0.030,-0.100) (0.040,-0.200) (0.050,-0.100) (0.060,-0.200) (0.070,-0.300) (0.080,-0.200) (0.090,-0.100) (0.100,-0.200) (0.110,-0.100) (0.120,0.000) (0.130,0.100) (0.140,0.200) (0.150,0.100) (0.160,0.200) (0.170,0.100) (0.180,0.200) (0.190,0.300) (0.200,0.200) (0.210,0.100) (0.220,0.200) (0.230,0.100) (0.240,0.200) (0.250,0.300) (0.260,0.200) (0.270,0.100) (0.280,0.200) (0.290,0.100) (0.300,0.200) (0.310,0.300) (0.320,0.200) (0.330,0.300) (0.340,0.200) (0.350,0.100) (0.360,0.200) (0.370,0.100) (0.380,0.000) (0.390,-0.100) (0.400,-0.200) (0.410,-0.100) (0.420,-0.200) (0.430,-0.100) (0.440,-0.200) (0.450,-0.300) (0.460,-0.200) (0.470,-0.300) (0.480,-0.400) (0.490,-0.300) (0.500,-0.400) (0.510,-0.300) (0.520,-0.200) (0.530,-0.300) (0.540,-0.400) (0.550,-0.500) (0.560,-0.400) (0.570,-0.500) (0.580,-0.400) (0.590,-0.500) (0.600,-0.600) (0.610,-0.500) (0.620,-0.400) (0.630,-0.300) (0.640,-0.200) (0.650,-0.100) (0.660,0.000) (0.670,-0.100) (0.680,-0.200) (0.690,-0.300) (0.700,-0.400) (0.710,-0.500) (0.720,-0.400) (0.730,-0.500) (0.740,-0.400) (0.750,-0.300) (0.760,-0.400) (0.770,-0.300) (0.780,-0.400) (0.790,-0.500) (0.800,-0.400) (0.810,-0.500) (0.820,-0.600) (0.830,-0.700) (0.840,-0.800) (0.850,-0.900) (0.860,-0.800) (0.870,-0.900) (0.880,-1.000) (0.890,-0.900) (0.900,-0.800) (0.910,-0.700) (0.920,-0.600) (0.930,-0.500) (0.940,-0.600) (0.950,-0.500) (0.960,-0.400) (0.970,-0.300) (0.980,-0.200) (0.990,-0.300) (1.000,-0.200)};
\end{axis}
\end{tikzpicture}\hfill
\begin{tikzpicture}
\begin{axis}[SAaxis, axis x line=bottom, axis y line=left, width=5.6cm, height=4.6cm, xmin=0, xmax=1.02, ymin=-1.2, ymax=1.0,
  xtick={0,0.5,1}, ytick={-1,0,1}, title={\small $n=1000$}, xlabel={$t$}]
  \addplot[SADark!20] coordinates {(0,0) (1,0)};
  \addplot[SABlue,thin] coordinates {(0.000,0.000) (0.002,0.000) (0.004,0.000) (0.006,0.000) (0.008,0.000) (0.010,0.063) (0.012,0.063) (0.014,0.000) (0.016,-0.063) (0.018,-0.126) (0.020,-0.063) (0.022,-0.126) (0.024,-0.190) (0.026,-0.190) (0.028,-0.190) (0.030,-0.190) (0.032,-0.126) (0.034,-0.126) (0.036,-0.063) (0.038,-0.063) (0.040,-0.063) (0.042,-0.063) (0.044,-0.063) (0.046,-0.126) (0.048,-0.126) (0.050,-0.063) (0.052,-0.063) (0.054,0.000) (0.056,-0.063) (0.058,-0.063) (0.060,-0.063) (0.062,0.000) (0.064,0.000) (0.066,0.063) (0.068,0.063) (0.070,0.063) (0.072,0.063) (0.074,0.063) (0.076,0.063) (0.078,0.000) (0.080,0.000) (0.082,-0.063) (0.084,0.000) (0.086,0.000) (0.088,0.000) (0.090,0.000) (0.092,0.000) (0.094,-0.063) (0.096,-0.063) (0.098,0.000) (0.100,0.063) (0.102,0.126) (0.104,0.126) (0.106,0.063) (0.108,0.000) (0.110,0.000) (0.112,0.000) (0.114,0.000) (0.116,0.000) (0.118,0.000) (0.120,0.000) (0.122,0.063) (0.124,0.063) (0.126,0.063) (0.128,0.126) (0.130,0.063) (0.132,0.063) (0.134,0.000) (0.136,0.063) (0.138,0.063) (0.140,0.126) (0.142,0.126) (0.144,0.190) (0.146,0.190) (0.148,0.253) (0.150,0.253) (0.152,0.316) (0.154,0.316) (0.156,0.253) (0.158,0.190) (0.160,0.190) (0.162,0.190) (0.164,0.253) (0.166,0.253) (0.168,0.190) (0.170,0.126) (0.172,0.126) (0.174,0.126) (0.176,0.126) (0.178,0.126) (0.180,0.063) (0.182,0.063) (0.184,0.063) (0.186,0.126) (0.188,0.126) (0.190,0.126) (0.192,0.126) (0.194,0.126) (0.196,0.063) (0.198,0.063) (0.200,0.063) (0.202,0.063) (0.204,0.063) (0.206,0.126) (0.208,0.190) (0.210,0.190) (0.212,0.190) (0.214,0.253) (0.216,0.316) (0.218,0.253) (0.220,0.190) (0.222,0.253) (0.224,0.253) (0.226,0.316) (0.228,0.379) (0.230,0.316) (0.232,0.316) (0.234,0.379) (0.236,0.443) (0.238,0.443) (0.240,0.443) (0.242,0.443) (0.244,0.506) (0.246,0.506) (0.248,0.443) (0.250,0.379) (0.252,0.379) (0.254,0.316) (0.256,0.316) (0.258,0.316) (0.260,0.316) (0.262,0.316) (0.264,0.379) (0.266,0.379) (0.268,0.316) (0.270,0.316) (0.272,0.379) (0.274,0.379) (0.276,0.379) (0.278,0.379) (0.280,0.379) (0.282,0.316) (0.284,0.316) (0.286,0.253) (0.288,0.316) (0.290,0.316) (0.292,0.379) (0.294,0.316) (0.296,0.253) (0.298,0.253) (0.300,0.316) (0.302,0.316) (0.304,0.316) (0.306,0.316) (0.308,0.316) (0.310,0.253) (0.312,0.253) (0.314,0.316) (0.316,0.316) (0.318,0.316) (0.320,0.379) (0.322,0.379) (0.324,0.316) (0.326,0.316) (0.328,0.253) (0.330,0.253) (0.332,0.316) (0.334,0.253) (0.336,0.253) (0.338,0.253) (0.340,0.253) (0.342,0.253) (0.344,0.316) (0.346,0.316) (0.348,0.379) (0.350,0.379) (0.352,0.379) (0.354,0.316) (0.356,0.316) (0.358,0.253) (0.360,0.253) (0.362,0.253) (0.364,0.253) (0.366,0.316) (0.368,0.379) (0.370,0.379) (0.372,0.379) (0.374,0.316) (0.376,0.253) (0.378,0.316) (0.380,0.253) (0.382,0.190) (0.384,0.126) (0.386,0.126) (0.388,0.063) (0.390,0.126) (0.392,0.126) (0.394,0.063) (0.396,0.063) (0.398,0.063) (0.400,0.063) (0.402,0.000) (0.404,-0.063) (0.406,-0.063) (0.408,-0.063) (0.410,-0.063) (0.412,-0.063) (0.414,-0.063) (0.416,0.000) (0.418,-0.063) (0.420,-0.126) (0.422,-0.126) (0.424,-0.190) (0.426,-0.190) (0.428,-0.190) (0.430,-0.190) (0.432,-0.190) (0.434,-0.190) (0.436,-0.126) (0.438,-0.063) (0.440,-0.126) (0.442,-0.126) (0.444,-0.063) (0.446,-0.126) (0.448,-0.063) (0.450,-0.126) (0.452,-0.063) (0.454,-0.126) (0.456,-0.126) (0.458,-0.190) (0.460,-0.190) (0.462,-0.253) (0.464,-0.253) (0.466,-0.190) (0.468,-0.253) (0.470,-0.253) (0.472,-0.253) (0.474,-0.190) (0.476,-0.190) (0.478,-0.253) (0.480,-0.190) (0.482,-0.190) (0.484,-0.253) (0.486,-0.253) (0.488,-0.190) (0.490,-0.190) (0.492,-0.126) (0.494,-0.190) (0.496,-0.190) (0.498,-0.126) (0.500,-0.063) (0.502,0.000) (0.504,-0.063) (0.506,-0.126) (0.508,-0.063) (0.510,0.000) (0.512,0.000) (0.514,-0.063) (0.516,-0.063) (0.518,0.000) (0.520,-0.063) (0.522,-0.063) (0.524,-0.126) (0.526,-0.190) (0.528,-0.126) (0.530,-0.063) (0.532,0.000) (0.534,-0.063) (0.536,-0.126) (0.538,-0.126) (0.540,-0.063) (0.542,0.000) (0.544,0.000) (0.546,0.000) (0.548,0.000) (0.550,0.000) (0.552,0.000) (0.554,0.063) (0.556,0.126) (0.558,0.126) (0.560,0.126) (0.562,0.190) (0.564,0.190) (0.566,0.190) (0.568,0.126) (0.570,0.063) (0.572,0.063) (0.574,0.126) (0.576,0.126) (0.578,0.126) (0.580,0.063) (0.582,0.063) (0.584,0.063) (0.586,0.063) (0.588,0.126) (0.590,0.190) (0.592,0.190) (0.594,0.253) (0.596,0.253) (0.598,0.253) (0.600,0.253) (0.602,0.316) (0.604,0.316) (0.606,0.253) (0.608,0.253) (0.610,0.253) (0.612,0.253) (0.614,0.253) (0.616,0.253) (0.618,0.190) (0.620,0.190) (0.622,0.190) (0.624,0.190) (0.626,0.190) (0.628,0.190) (0.630,0.126) (0.632,0.126) (0.634,0.126) (0.636,0.126) (0.638,0.190) (0.640,0.126) (0.642,0.126) (0.644,0.063) (0.646,0.126) (0.648,0.126) (0.650,0.063) (0.652,0.000) (0.654,-0.063) (0.656,-0.063) (0.658,0.000) (0.660,-0.063) (0.662,-0.063) (0.664,-0.126) (0.666,-0.190) (0.668,-0.253) (0.670,-0.190) (0.672,-0.126) (0.674,-0.190) (0.676,-0.190) (0.678,-0.253) (0.680,-0.316) (0.682,-0.379) (0.684,-0.316) (0.686,-0.316) (0.688,-0.253) (0.690,-0.190) (0.692,-0.190) (0.694,-0.190) (0.696,-0.126) (0.698,-0.126) (0.700,-0.063) (0.702,0.000) (0.704,0.000) (0.706,-0.063) (0.708,-0.063) (0.710,-0.063) (0.712,-0.063) (0.714,-0.126) (0.716,-0.126) (0.718,-0.126) (0.720,-0.126) (0.722,-0.126) (0.724,-0.126) (0.726,-0.190) (0.728,-0.190) (0.730,-0.190) (0.732,-0.190) (0.734,-0.126) (0.736,-0.190) (0.738,-0.253) (0.740,-0.316) (0.742,-0.379) (0.744,-0.379) (0.746,-0.316) (0.748,-0.316) (0.750,-0.316) (0.752,-0.379) (0.754,-0.443) (0.756,-0.379) (0.758,-0.316) (0.760,-0.379) (0.762,-0.443) (0.764,-0.506) (0.766,-0.569) (0.768,-0.569) (0.770,-0.632) (0.772,-0.632) (0.774,-0.569) (0.776,-0.569) (0.778,-0.569) (0.780,-0.632) (0.782,-0.696) (0.784,-0.759) (0.786,-0.822) (0.788,-0.822) (0.790,-0.885) (0.792,-0.885) (0.794,-0.822) (0.796,-0.759) (0.798,-0.696) (0.800,-0.696) (0.802,-0.696) (0.804,-0.696) (0.806,-0.632) (0.808,-0.632) (0.810,-0.569) (0.812,-0.632) (0.814,-0.569) (0.816,-0.569) (0.818,-0.632) (0.820,-0.632) (0.822,-0.569) (0.824,-0.569) (0.826,-0.569) (0.828,-0.632) (0.830,-0.632) (0.832,-0.632) (0.834,-0.632) (0.836,-0.569) (0.838,-0.632) (0.840,-0.696) (0.842,-0.696) (0.844,-0.696) (0.846,-0.696) (0.848,-0.696) (0.850,-0.759) (0.852,-0.696) (0.854,-0.759) (0.856,-0.759) (0.858,-0.696) (0.860,-0.696) (0.862,-0.696) (0.864,-0.696) (0.866,-0.759) (0.868,-0.822) (0.870,-0.885) (0.872,-0.885) (0.874,-0.949) (0.876,-0.949) (0.878,-0.885) (0.880,-0.885) (0.882,-0.885) (0.884,-0.822) (0.886,-0.885) (0.888,-0.949) (0.890,-0.885) (0.892,-0.885) (0.894,-0.822) (0.896,-0.759) (0.898,-0.696) (0.900,-0.696) (0.902,-0.759) (0.904,-0.696) (0.906,-0.696) (0.908,-0.632) (0.910,-0.696) (0.912,-0.759) (0.914,-0.696) (0.916,-0.759) (0.918,-0.696) (0.920,-0.696) (0.922,-0.696) (0.924,-0.632) (0.926,-0.632) (0.928,-0.632) (0.930,-0.632) (0.932,-0.632) (0.934,-0.569) (0.936,-0.632) (0.938,-0.569) (0.940,-0.632) (0.942,-0.632) (0.944,-0.632) (0.946,-0.632) (0.948,-0.632) (0.950,-0.696) (0.952,-0.759) (0.954,-0.759) (0.956,-0.696) (0.958,-0.632) (0.960,-0.569) (0.962,-0.632) (0.964,-0.632) (0.966,-0.696) (0.968,-0.696) (0.970,-0.696) (0.972,-0.696) (0.974,-0.696) (0.976,-0.759) (0.978,-0.696) (0.980,-0.632) (0.982,-0.632) (0.984,-0.696) (0.986,-0.696) (0.988,-0.696) (0.990,-0.759) (0.992,-0.696) (0.994,-0.759) (0.996,-0.696) (0.998,-0.696) (1.000,-0.696)};
\end{axis}
\end{tikzpicture}
\caption{缩放随机游动 $W^{(n)}(t)=S_{\lfloor nt\rfloor}/\sqrt{n}$ ($\xi_k=\pm1$) 的轨道. 时间每缩小 $n$ 倍, 空间只缩小 $\sqrt{n}$ 倍, 所以单步跳跃 $1/\sqrt{n}$ 趋于零, 而在单位时间内的总位移保持 $O(1)$}
\label{fig:scaled-rw}
\end{figure}

% ---------------------------------------------------------------------
\subsection{有限维分布的收敛}

第~1~章只讨论了随机变量的依分布收敛. 对随机向量, 定义完全类似: 称 $\R^k$ 值随机向量 $\rvec{X}_n\xrightarrow{d}\rvec{X}$, 如果对一切有界连续函数 $g:\R^k\to\R$, $\E[g(\rvec{X}_n)]\to\E[g(\rvec{X})]$ (一维时这与定理~\ref{thm:portmanteau} 一致). L\'evy 连续性定理~\ref{thm:levy-continuity} 对随机向量同样成立: $\rvec{X}_n\xrightarrow{d}\rvec{X}$ 当且仅当 $\varphi_{\rvec{X}_n}(\rvec{u})\to\varphi_{\rvec{X}}(\rvec{u})$ 对一切 $\rvec{u}\in\R^k$ 成立 (见~\cite{durrett} 定理 3.10.5).

\begin{proposition}[缩放随机游动的有限维分布收敛]\label{prop:rw-fdd}
对任意 $0\le t_1<t_2<\cdots<t_k$,
\[
\big(W^{(n)}(t_1),\dots,W^{(n)}(t_k)\big)\xrightarrow{d}\big(G_1,\dots,G_k\big),
\]
其中 $(G_1,\dots,G_k)$ 是均值为零的正态随机向量, $\Cov(G_i,G_j)=t_i\wedge t_j$. 等价地, 极限向量的增量 $G_1,G_2-G_1,\dots,G_k-G_{k-1}$ 相互独立, 且 $G_i-G_{i-1}\sim N(0,t_i-t_{i-1})$ (约定 $t_0=0$, $G_0=0$).
\end{proposition}

\begin{proof}
\textbf{第一步: 单个增量.} 固定 $s<t$, 令 $m_n=\lfloor nt\rfloor-\lfloor ns\rfloor$. 则 $W^{(n)}(t)-W^{(n)}(s)=\frac{1}{\sqrt{n}}\big(\xi_{\lfloor ns\rfloor+1}+\cdots+\xi_{\lfloor nt\rfloor}\big)$ 与 $S_{m_n}/\sqrt{n}$ 同分布. 由于 $nt-1<\lfloor nt\rfloor\le nt$, 有 $\abs{m_n-n(t-s)}<1$, 从而 $m_n\to\infty$ 且 $c_n:=\sqrt{m_n/n}\to\sqrt{t-s}$. 记 $X_n=S_{m_n}/\sqrt{m_n}$. 由中心极限定理~\ref{thm:clt} (沿子列 $m_n\to\infty$), $X_n\xrightarrow{d}Z\sim N(0,1)$. 由 Skorokhod 表示定理~\ref{thm:skorokhod}, 存在 $Y_n\overset{d}{=}X_n$, $Y\overset{d}{=}Z$, 且 $Y_n\to Y$ $\as$, 从而 $c_nY_n\to\sqrt{t-s}\,Y$ $\as$, 由定理~\ref{thm:convergence-implications} 也依分布收敛. 由于 $c_nY_n\overset{d}{=}c_nX_n=S_{m_n}/\sqrt{n}$, 得
\[
W^{(n)}(t)-W^{(n)}(s)\xrightarrow{d}N(0,t-s).
\]

\textbf{第二步: 增量向量.} 记 $\Delta^{(n)}_i=W^{(n)}(t_i)-W^{(n)}(t_{i-1})$, $i=1,\dots,k$. 它们分别由不相交的几组 $\xi$ 决定, 由分组原理 (推论~\ref{cor:grouping}) 相互独立. 由第一步与 L\'evy 连续性定理~\ref{thm:levy-continuity}, 对每个 $v\in\R$, $\varphi_{\Delta^{(n)}_i}(v)\to e^{-(t_i-t_{i-1})v^2/2}$ ($t_1=0$ 时 $\Delta^{(n)}_1=0$, 结论平凡). 由独立性 (定理~\ref{thm:indep-cf}), 对 $\rvec{v}=(v_1,\dots,v_k)$,
\[
\varphi_{\rvec{\Delta}^{(n)}}(\rvec{v})=\prod_{i=1}^k\varphi_{\Delta^{(n)}_i}(v_i)\ \longrightarrow\ \prod_{i=1}^ke^{-(t_i-t_{i-1})v_i^2/2}=\varphi_{\rvec{\Delta}}(\rvec{v}),
\]
其中 $\rvec{\Delta}=(\Delta_1,\dots,\Delta_k)$ 的分量独立, $\Delta_i\sim N(0,t_i-t_{i-1})$.

\textbf{第三步: 线性变换.} 令 $\mat{A}$ 为 $k\times k$ 下三角全 $1$ 矩阵, 则 $\big(W^{(n)}(t_1),\dots,W^{(n)}(t_k)\big)^\T=\mat{A}\rvec{\Delta}^{(n)}$. 对 $\rvec{u}\in\R^k$,
\[
\varphi_{\mat{A}\rvec{\Delta}^{(n)}}(\rvec{u})=\E\big[e^{\mathrm{i}\rvec{u}^\T\mat{A}\rvec{\Delta}^{(n)}}\big]=\varphi_{\rvec{\Delta}^{(n)}}(\mat{A}^\T\rvec{u})\ \longrightarrow\ \varphi_{\rvec{\Delta}}(\mat{A}^\T\rvec{u})=\varphi_{\mat{A}\rvec{\Delta}}(\rvec{u}).
\]
由多维 L\'evy 连续性定理, $\mat{A}\rvec{\Delta}^{(n)}\xrightarrow{d}\rvec{G}:=\mat{A}\rvec{\Delta}$. 由命题~\ref{prop:gaussian-operations}, $\rvec{G}$ 是均值为零的正态向量; 对 $i\le j$, 由增量的独立性,
\[
\Cov(G_i,G_j)=\Cov\Big(\sum_{l\le i}\Delta_l,\ \sum_{l\le j}\Delta_l\Big)=\sum_{l\le i}\Var(\Delta_l)=t_i=t_i\wedge t_j.
\]
等价性的说法就是 $\rvec{G}=\mat{A}\rvec{\Delta}$ 的构造本身.
\end{proof}

\begin{remark}
极限与 $\xi_1$ 的具体分布无关, 只用到了 $\E\xi_1=0$ 与 $\Var(\xi_1)=1$. 这种普适性称为\textbf{不变原理}. 另外, 若把 $W^{(n)}$ 换成相邻格点之间线性插值得到的连续折线 $\widetilde{W}^{(n)}$, 则 $\abs{\widetilde{W}^{(n)}(t)-W^{(n)}(t)}\le\abs{\xi_{\lfloor nt\rfloor+1}}/\sqrt{n}\xrightarrow{\P}0$, 二者的有限维极限相同.
\end{remark}

% ---------------------------------------------------------------------
\subsection{Donsker 不变原理}

有限维分布收敛只涉及有限个时刻. 要研究 ``随机游动在 $n$ 步内的最大值'' 这类依赖于整条轨道的量, 我们需要在轨道空间中讨论收敛. 记 $C[0,1]$ 为 $[0,1]$ 上连续函数全体, 赋予一致范数 $\norm{w}_\infty=\sup_{t\le1}\abs{w(t)}$.

\begin{theorem}[Donsker 不变原理]\label{thm:donsker}
设 $\widetilde{W}^{(n)}$ 是缩放随机游动的线性插值, $B$ 是标准布朗运动 (定义~\ref{def:bm}). 则对每个有界函数 $F:C[0,1]\to\R$, 只要 $F$ 关于一致范数连续,
\[
\E\big[F(\widetilde{W}^{(n)})\big]\to\E\big[F(B)\big].
\]
(证明见~\cite{billingsley} 或~\cite{karatzas} 定理 2.4.20.)
\end{theorem}

例如, $F(w)=g\big(\max_{t\le1}w(t)\big)$ 对有界连续的 $g$ 满足条件, 所以 $\max_{k\le n}S_k/\sqrt{n}\xrightarrow{d}\max_{t\le1}B_t$. 4.5 节的反射原理将给出右端的分布: 它与 $\abs{B_1}$ 同分布. 于是对 $a>0$,
\[
\P\Big(\max_{k\le n}S_k\ge a\sqrt{n}\Big)\to2\big(1-\Phi(a)\big).
\]
这个关于随机游动的结论, 无论 $\xi_1$ 是什么分布 (只要均值为零、方差为一) 都成立, 而它的计算只需要在布朗运动上完成一次. 这正是布朗运动作为 ``普适极限'' 的价值.

\begin{intuition}
中心极限定理说: 无论单个小量服从什么分布, 大量独立小量之和都近似正态. Donsker 定理是它的 ``过程版本'': 无论单步服从什么分布, 缩放后的整条随机游动轨道都近似于布朗运动. 正如正态分布由均值与方差两个参数完全决定, 布朗运动作为极限也只 ``记得'' 单步的均值 (为零) 与方差 (为一), 其余细节在极限中全部消失.
\end{intuition}

% ---------------------------------------------------------------------
\subsection{极限过程应有的性质}

假设缩放随机游动有一个极限过程 $B$. 从 $W^{(n)}$ 的构造与命题~\ref{prop:rw-fdd} 可以读出 $B$ 应具备的四条性质.
\begin{enumerate}
\item $B_0=0$;
\item \textbf{独立增量}: 不相交时间段上的增量由不相交的几组 $\xi$ 决定, 因而相互独立;
\item \textbf{平稳正态增量}: $B_t-B_s\sim N(0,t-s)$, 只依赖于时间段的长度;
\item \textbf{连续轨道}: 单步跳跃 $\abs{\xi_k}/\sqrt{n}$ 趋于零, 极限轨道不应有跳跃.
\end{enumerate}
物理上, (2) 对应于 ``不同时间段内的分子碰撞相互独立'', (3) 中方差与时间成正比正是 Einstein 的扩散定律. 下一节以这四条为定义, 并证明这样的过程确实存在. 前三条只涉及有限维分布, 由 Kolmogorov 延拓定理即可保证; 真正的困难在第 (4) 条.

% =====================================================================
\section{布朗运动的定义与存在性}
\label{sec:4-2}

% ---------------------------------------------------------------------
\subsection{定义与等价刻画}

\begin{definition}[标准布朗运动]\label{def:bm}
概率空间 $(\Omega,\mathcal{F},\P)$ 上的实值过程 $B=\{B_t\}_{t\ge0}$ 称为\textbf{标准布朗运动}, 如果
\begin{enumerate}
\item $B_0=0$;
\item 对任意 $0\le t_0<t_1<\cdots<t_n$, 增量 $B_{t_1}-B_{t_0},\dots,B_{t_n}-B_{t_{n-1}}$ 相互独立;
\item 对任意 $0\le s<t$, $B_t-B_s\sim N(0,t-s)$;
\item 对每个 $\omega\in\Omega$, 轨道 $t\mapsto B_t(\omega)$ 连续.
\end{enumerate}
若 (1)(4) 只在一个概率为 $1$ 的事件上成立, 只需在其补集上把 $B$ 重新定义为 $0$, 就得到满足全部四条的过程, 所以通常不加区分.
\end{definition}

\begin{figure}[htbp]
\centering
\begin{tikzpicture}
\begin{axis}[SAaxis, axis x line=bottom, axis y line=left, width=12cm, height=5.6cm, xmin=0, xmax=1.03, ymin=-0.9, ymax=1.3,
  xtick={0,0.25,0.5,0.75,1}, ytick={-0.5,0,0.5,1}, xlabel={$t$}, ylabel={$B_t$}, clip=false]
  \addplot[SADark!20] coordinates {(0,0) (1,0)};
  \addplot[SABlue!25,fill=SABlue!8,draw=none,domain=0:1,samples=60] {sqrt(x)} \closedcycle;
  \addplot[SABlue!25,fill=SABlue!8,draw=none,domain=0:1,samples=60] {-sqrt(x)} \closedcycle;
  \addplot[SABlue!50,dashed,domain=0:1,samples=60] {sqrt(x)};
  \addplot[SABlue!50,dashed,domain=0:1,samples=60] {-sqrt(x)};
  \addplot[SADark,thin] coordinates {(0.000,0.000) (0.002,0.002) (0.004,0.034) (0.006,-0.039) (0.008,-0.005) (0.010,-0.010) (0.012,-0.050) (0.014,-0.043) (0.016,-0.060) (0.018,-0.082) (0.020,-0.140) (0.021,-0.082) (0.023,-0.102) (0.025,-0.056) (0.027,-0.039) (0.029,-0.064) (0.031,-0.016) (0.033,0.052) (0.035,0.043) (0.037,0.039) (0.039,0.017) (0.041,0.033) (0.043,-0.006) (0.045,0.054) (0.047,0.114) (0.049,0.035) (0.051,0.094) (0.053,0.097) (0.055,0.178) (0.057,0.163) (0.059,0.177) (0.061,0.197) (0.062,0.155) (0.064,0.193) (0.066,0.183) (0.068,0.102) (0.070,0.024) (0.072,0.013) (0.074,0.041) (0.076,0.113) (0.078,0.134) (0.080,0.161) (0.082,0.188) (0.084,0.184) (0.086,0.180) (0.088,0.208) (0.090,0.228) (0.092,0.214) (0.094,0.222) (0.096,0.124) (0.098,0.112) (0.100,0.195) (0.102,0.277) (0.104,0.358) (0.105,0.401) (0.107,0.363) (0.109,0.305) (0.111,0.340) (0.113,0.215) (0.115,0.254) (0.117,0.182) (0.119,0.166) (0.121,0.066) (0.123,0.057) (0.125,0.004) (0.127,0.050) (0.129,0.094) (0.131,0.120) (0.133,0.164) (0.135,0.223) (0.137,0.258) (0.139,0.333) (0.141,0.338) (0.143,0.275) (0.145,0.249) (0.146,0.238) (0.148,0.296) (0.150,0.314) (0.152,0.247) (0.154,0.247) (0.156,0.273) (0.158,0.300) (0.160,0.282) (0.162,0.279) (0.164,0.282) (0.166,0.262) (0.168,0.303) (0.170,0.295) (0.172,0.242) (0.174,0.220) (0.176,0.286) (0.178,0.313) (0.180,0.288) (0.182,0.311) (0.184,0.338) (0.186,0.358) (0.188,0.364) (0.189,0.346) (0.191,0.417) (0.193,0.474) (0.195,0.452) (0.197,0.423) (0.199,0.504) (0.201,0.518) (0.203,0.461) (0.205,0.438) (0.207,0.444) (0.209,0.403) (0.211,0.457) (0.213,0.456) (0.215,0.494) (0.217,0.488) (0.219,0.517) (0.221,0.549) (0.223,0.568) (0.225,0.562) (0.227,0.562) (0.229,0.533) (0.230,0.566) (0.232,0.542) (0.234,0.522) (0.236,0.512) (0.238,0.465) (0.240,0.447) (0.242,0.483) (0.244,0.423) (0.246,0.445) (0.248,0.404) (0.250,0.358) (0.252,0.421) (0.254,0.420) (0.256,0.423) (0.258,0.432) (0.260,0.326) (0.262,0.278) (0.264,0.325) (0.266,0.361) (0.268,0.397) (0.270,0.392) (0.271,0.482) (0.273,0.512) (0.275,0.536) (0.277,0.540) (0.279,0.500) (0.281,0.493) (0.283,0.451) (0.285,0.394) (0.287,0.300) (0.289,0.429) (0.291,0.375) (0.293,0.431) (0.295,0.458) (0.297,0.398) (0.299,0.440) (0.301,0.454) (0.303,0.442) (0.305,0.405) (0.307,0.479) (0.309,0.490) (0.311,0.467) (0.312,0.504) (0.314,0.496) (0.316,0.491) (0.318,0.484) (0.320,0.513) (0.322,0.597) (0.324,0.634) (0.326,0.609) (0.328,0.537) (0.330,0.499) (0.332,0.479) (0.334,0.425) (0.336,0.458) (0.338,0.424) (0.340,0.501) (0.342,0.523) (0.344,0.599) (0.346,0.607) (0.348,0.633) (0.350,0.615) (0.352,0.656) (0.354,0.677) (0.355,0.764) (0.357,0.791) (0.359,0.760) (0.361,0.720) (0.363,0.702) (0.365,0.662) (0.367,0.595) (0.369,0.570) (0.371,0.556) (0.373,0.532) (0.375,0.430) (0.377,0.441) (0.379,0.478) (0.381,0.389) (0.383,0.350) (0.385,0.272) (0.387,0.265) (0.389,0.287) (0.391,0.234) (0.393,0.363) (0.395,0.376) (0.396,0.271) (0.398,0.329) (0.400,0.382) (0.402,0.387) (0.404,0.439) (0.406,0.417) (0.408,0.347) (0.410,0.370) (0.412,0.285) (0.414,0.289) (0.416,0.261) (0.418,0.247) (0.420,0.277) (0.422,0.325) (0.424,0.351) (0.426,0.337) (0.428,0.341) (0.430,0.253) (0.432,0.252) (0.434,0.302) (0.436,0.227) (0.438,0.148) (0.439,0.176) (0.441,0.246) (0.443,0.312) (0.445,0.288) (0.447,0.340) (0.449,0.438) (0.451,0.465) (0.453,0.524) (0.455,0.505) (0.457,0.512) (0.459,0.518) (0.461,0.488) (0.463,0.448) (0.465,0.489) (0.467,0.498) (0.469,0.528) (0.471,0.557) (0.473,0.612) (0.475,0.587) (0.477,0.585) (0.479,0.559) (0.480,0.528) (0.482,0.487) (0.484,0.538) (0.486,0.510) (0.488,0.525) (0.490,0.431) (0.492,0.464) (0.494,0.493) (0.496,0.533) (0.498,0.495) (0.500,0.537) (0.502,0.566) (0.504,0.579) (0.506,0.638) (0.508,0.542) (0.510,0.507) (0.512,0.400) (0.514,0.370) (0.516,0.376) (0.518,0.372) (0.520,0.348) (0.521,0.243) (0.523,0.242) (0.525,0.214) (0.527,0.216) (0.529,0.217) (0.531,0.190) (0.533,0.202) (0.535,0.183) (0.537,0.104) (0.539,0.048) (0.541,0.082) (0.543,0.041) (0.545,0.066) (0.547,-0.017) (0.549,-0.027) (0.551,-0.044) (0.553,0.015) (0.555,0.051) (0.557,0.030) (0.559,0.059) (0.561,0.035) (0.562,0.035) (0.564,0.067) (0.566,0.108) (0.568,0.166) (0.570,0.183) (0.572,0.199) (0.574,0.202) (0.576,0.184) (0.578,0.165) (0.580,0.257) (0.582,0.206) (0.584,0.172) (0.586,0.193) (0.588,0.179) (0.590,0.171) (0.592,0.194) (0.594,0.231) (0.596,0.200) (0.598,0.250) (0.600,0.313) (0.602,0.348) (0.604,0.297) (0.605,0.345) (0.607,0.370) (0.609,0.348) (0.611,0.347) (0.613,0.308) (0.615,0.280) (0.617,0.245) (0.619,0.248) (0.621,0.296) (0.623,0.290) (0.625,0.269) (0.627,0.274) (0.629,0.341) (0.631,0.358) (0.633,0.299) (0.635,0.355) (0.637,0.374) (0.639,0.395) (0.641,0.450) (0.643,0.386) (0.645,0.374) (0.646,0.434) (0.648,0.477) (0.650,0.412) (0.652,0.419) (0.654,0.472) (0.656,0.469) (0.658,0.469) (0.660,0.441) (0.662,0.378) (0.664,0.387) (0.666,0.345) (0.668,0.367) (0.670,0.379) (0.672,0.381) (0.674,0.321) (0.676,0.368) (0.678,0.294) (0.680,0.225) (0.682,0.266) (0.684,0.273) (0.686,0.270) (0.688,0.254) (0.689,0.185) (0.691,0.109) (0.693,0.025) (0.695,-0.002) (0.697,0.027) (0.699,-0.002) (0.701,-0.002) (0.703,0.083) (0.705,0.114) (0.707,0.161) (0.709,0.186) (0.711,0.121) (0.713,0.158) (0.715,0.292) (0.717,0.269) (0.719,0.285) (0.721,0.377) (0.723,0.460) (0.725,0.520) (0.727,0.551) (0.729,0.439) (0.730,0.418) (0.732,0.388) (0.734,0.391) (0.736,0.393) (0.738,0.328) (0.740,0.238) (0.742,0.253) (0.744,0.244) (0.746,0.253) (0.748,0.285) (0.750,0.237) (0.752,0.172) (0.754,0.210) (0.756,0.155) (0.758,0.151) (0.760,0.119) (0.762,0.146) (0.764,0.141) (0.766,0.086) (0.768,-0.008) (0.770,0.033) (0.771,0.052) (0.773,0.079) (0.775,0.142) (0.777,0.137) (0.779,0.091) (0.781,0.091) (0.783,0.100) (0.785,0.071) (0.787,0.053) (0.789,0.030) (0.791,0.035) (0.793,0.098) (0.795,0.069) (0.797,0.062) (0.799,0.054) (0.801,0.079) (0.803,0.039) (0.805,0.009) (0.807,0.045) (0.809,0.041) (0.811,0.100) (0.812,0.152) (0.814,0.143) (0.816,0.199) (0.818,0.250) (0.820,0.239) (0.822,0.228) (0.824,0.209) (0.826,0.159) (0.828,0.181) (0.830,0.241) (0.832,0.203) (0.834,0.250) (0.836,0.223) (0.838,0.201) (0.840,0.269) (0.842,0.192) (0.844,0.149) (0.846,0.173) (0.848,0.149) (0.850,0.088) (0.852,0.184) (0.854,0.184) (0.855,0.240) (0.857,0.261) (0.859,0.270) (0.861,0.230) (0.863,0.230) (0.865,0.255) (0.867,0.256) (0.869,0.229) (0.871,0.166) (0.873,0.173) (0.875,0.234) (0.877,0.263) (0.879,0.257) (0.881,0.202) (0.883,0.170) (0.885,0.205) (0.887,0.134) (0.889,0.106) (0.891,0.071) (0.893,0.054) (0.895,0.085) (0.896,-0.002) (0.898,-0.030) (0.900,-0.079) (0.902,-0.051) (0.904,-0.014) (0.906,0.043) (0.908,0.083) (0.910,0.039) (0.912,0.075) (0.914,0.038) (0.916,0.099) (0.918,0.082) (0.920,0.117) (0.922,0.178) (0.924,0.177) (0.926,0.148) (0.928,0.170) (0.930,0.165) (0.932,0.153) (0.934,0.120) (0.936,0.153) (0.938,0.097) (0.939,0.113) (0.941,0.150) (0.943,0.152) (0.945,0.233) (0.947,0.288) (0.949,0.387) (0.951,0.296) (0.953,0.265) (0.955,0.257) (0.957,0.224) (0.959,0.164) (0.961,0.062) (0.963,-0.003) (0.965,0.030) (0.967,0.089) (0.969,0.089) (0.971,0.063) (0.973,0.124) (0.975,-0.010) (0.977,-0.013) (0.979,-0.037) (0.980,-0.009) (0.982,-0.034) (0.984,0.042) (0.986,0.029) (0.988,0.001) (0.990,-0.028) (0.992,-0.026) (0.994,0.105) (0.996,0.151) (0.998,0.129) (1.000,0.088)};
  \node[SABlue,font=\small,anchor=west] at (axis cs:1.0,1.0) {$\sqrt{t}$};
  \node[SABlue,font=\small,anchor=west] at (axis cs:1.0,-1.0) {$-\sqrt{t}$};
\end{axis}
\end{tikzpicture}
\caption{标准布朗运动的一条轨道. 阴影区域 $\abs{x}\le\sqrt{t}$ 标出一倍标准差的范围: 对每个固定的 $t$, $\P(\abs{B_t}\le\sqrt{t})\approx0.683$}
\label{fig:bm-path}
\end{figure}

许多问题中, 除了布朗运动本身, 还有其他随机量 (例如初始条件或另一个独立的过程) 需要一同考虑, 这时需要相对于一个更大的信息流定义布朗运动.

\begin{definition}[$\{\mathcal{F}_t\}$-布朗运动]\label{def:Ft-bm}
设 $\{\mathcal{F}_t\}_{t\ge0}$ 是信息流. 称 $B$ 为 \textbf{$\{\mathcal{F}_t\}$-布朗运动}, 如果 $B$ 满足定义~\ref{def:bm} 的 (1)(3)(4), 适应于 $\{\mathcal{F}_t\}$, 且对一切 $0\le s<t$, $B_t-B_s$ 与 $\mathcal{F}_s$ 独立.
\end{definition}

由于 $B_{t_1}-B_{t_0},\dots,B_{t_{n-1}}-B_{t_{n-2}}$ 都是 $\mathcal{F}_{t_{n-1}}$-可测的, 而 $B_{t_n}-B_{t_{n-1}}$ 与 $\mathcal{F}_{t_{n-1}}$ 独立, 对 $n$ 归纳可知 $\{\mathcal{F}_t\}$-布朗运动满足定义~\ref{def:bm} (2). 反过来, 标准布朗运动是关于其自然信息流 $\mathcal{F}^B_t=\sigma(B_s: s\le t)$ 的布朗运动: 对 $s<t$, 事件族 $\mathcal{P}=\{\{(B_{s_1},\dots,B_{s_m})\in A\}: m\ge1,\ s_1<\cdots<s_m\le s,\ A\in\mathcal{B}(\R^m)\}$ 是生成 $\mathcal{F}^B_s$ 的 $\pi$-系, 而 $(B_{s_1},\dots,B_{s_m})$ 是增量 $B_{s_1}-B_0,\dots,B_{s_m}-B_{s_{m-1}}$ 的函数, 后者与 $B_t-B_s$ 独立 (由定义~\ref{def:bm} (2), 必要时把 $s$ 加入分点), 由定理~\ref{thm:indep-pi}, $B_t-B_s$ 与 $\mathcal{F}^B_s$ 独立.

下面的刻画把布朗运动归结为高斯过程, 这是验证一个过程是布朗运动的最常用的方法.

\begin{theorem}[布朗运动的高斯刻画]\label{thm:bm-gaussian}
设 $B=\{B_t\}_{t\ge0}$ 的轨道连续且 $B_0=0$. 则 $B$ 是标准布朗运动, 当且仅当 $B$ 是均值为零的高斯过程, 且
\[
\Cov(B_s,B_t)=s\wedge t,\qquad s,t\ge0.
\]
\end{theorem}

\begin{proof}
``仅当'': 任取 $0<t_1<\cdots<t_n$. 由定义~\ref{def:bm} (2)(3), 增量 $\Delta_i=B_{t_i}-B_{t_{i-1}}$ ($t_0=0$) 是独立的正态随机变量, 所以 $(\Delta_1,\dots,\Delta_n)$ 是正态随机向量 (定理~\ref{thm:gaussian-equivalent} (C)). $(B_{t_1},\dots,B_{t_n})$ 是它的线性变换, 由命题~\ref{prop:gaussian-operations} 仍是正态的, 且均值为零. 对 $s\le t$, 由于 $B_t-B_s$ 与 $B_s=B_s-B_0$ 独立,
\[
\Cov(B_s,B_t)=\Cov(B_s,B_s)+\Cov(B_s,B_t-B_s)=s+0=s\wedge t.
\]

``当'': 任取 $0=t_0<t_1<\cdots<t_n$. 增量向量 $(\Delta_1,\dots,\Delta_n)$ 是正态向量 $(B_{t_1},\dots,B_{t_n})$ 的线性变换, 因而是均值为零的正态向量. 对 $i<j$, 由于 $t_{i-1}<t_i\le t_{j-1}<t_j$,
\begin{align*}
\Cov(\Delta_i,\Delta_j)&=\Cov(B_{t_i},B_{t_j})-\Cov(B_{t_i},B_{t_{j-1}})-\Cov(B_{t_{i-1}},B_{t_j})+\Cov(B_{t_{i-1}},B_{t_{j-1}})\\
&=t_i-t_i-t_{i-1}+t_{i-1}=0,
\end{align*}
而 $\Var(\Delta_i)=t_i-2t_{i-1}+t_{i-1}=t_i-t_{i-1}$. 由定理~\ref{thm:gaussian-uncorrelated}, 不相关的正态分量相互独立, 于是 (2)(3) 成立 (若 $t_0>0$, 把 $0$ 加入分点即可).
\end{proof}

% ---------------------------------------------------------------------
\subsection{Kolmogorov 连续性定理}

由推论~\ref{cor:gaussian-process} 与例~\ref{ex:min-kernel}, 存在均值为零、协方差为 $s\wedge t$ 的高斯过程 $X$. 它满足定义~\ref{def:bm} 的 (1)--(3) (定理~\ref{thm:bm-gaussian} 的证明中, ``当'' 的部分没有用到连续性), 但如例~\ref{ex:min-kernel} 所说, 我们无法断言它的轨道连续. 出路是寻找 $X$ 的一个修正 (定义~\ref{def:modification}): 修正有相同的有限维分布, 因此仍满足 (1)--(3), 而我们希望修正的轨道是连续的. 下面的定理给出了一个只依赖于二维分布的充分条件.

\begin{theorem}[Kolmogorov-Chentsov 连续性定理]\label{thm:kolmogorov-chentsov}
设 $X=\{X_t\}_{t\in[0,T]}$ 是实值过程, 存在常数 $\alpha,\beta,C>0$, 使得
\begin{equation}\label{eq:kc-condition}
\E\abs{X_t-X_s}^\alpha\le C\abs{t-s}^{1+\beta},\qquad s,t\in[0,T].
\end{equation}
则 $X$ 有一个连续修正 $\widetilde{X}$. 进一步, 对每个 $\gamma\in(0,\beta/\alpha)$, $\widetilde{X}$ 的轨道几乎必然是 $\gamma$ 阶 H\"older 连续的: 存在几乎必然有限的随机变量 $K_\gamma$, 使
\[
\abs{\widetilde{X}_t-\widetilde{X}_s}\le K_\gamma\abs{t-s}^\gamma,\qquad s,t\in[0,T].
\]
\end{theorem}

证明的思路是先在二进有理点上控制增量, 再用 ``链式'' 论证把二进点上的增量拼接起来, 得到二进点上的一致连续性, 最后连续延拓到整个区间 (图~\ref{fig:chaining}).

\begin{proof}
用 $\widetilde{X}_t=X_{Tt}$ 代替 $X$, 条件 \eqref{eq:kc-condition} 只改变常数 $C$, 所以不妨设 $T=1$. 记 $D_n=\{k2^{-n}: 0\le k\le2^n\}$, $D=\bigcup_nD_n$ 为 $[0,1]$ 中的二进有理数全体. 固定 $\gamma\in(0,\beta/\alpha)$.

\textbf{第一步: 相邻二进点上的增量.} 由 Markov 不等式~\ref{thm:markov} 与 \eqref{eq:kc-condition}, 对每个 $1\le k\le2^n$,
\[
\P\big(\abs{X_{k2^{-n}}-X_{(k-1)2^{-n}}}\ge2^{-\gamma n}\big)\le2^{\gamma\alpha n}\,\E\abs{X_{k2^{-n}}-X_{(k-1)2^{-n}}}^\alpha\le C\,2^{\gamma\alpha n}2^{-n(1+\beta)}.
\]
对 $k$ 取并,
\[
\P\Big(\max_{1\le k\le2^n}\abs{X_{k2^{-n}}-X_{(k-1)2^{-n}}}\ge2^{-\gamma n}\Big)\le2^n\cdot C\,2^{\gamma\alpha n}2^{-n(1+\beta)}=C\,2^{-n(\beta-\alpha\gamma)}.
\]
由于 $\beta-\alpha\gamma>0$, 右端关于 $n$ 可和. 由 Borel-Cantelli 引理~\ref{thm:bc1}, 存在概率为 $1$ 的事件 $\Omega_0$, 使得对每个 $\omega\in\Omega_0$, 存在 $n^*(\omega)<\infty$, 当 $n\ge n^*(\omega)$ 时
\begin{equation}\label{eq:kc-adjacent}
\abs{X_{k2^{-n}}(\omega)-X_{(k-1)2^{-n}}(\omega)}<2^{-\gamma n},\qquad 1\le k\le2^n.
\end{equation}

\textbf{第二步: 链式论证.} 固定 $\omega\in\Omega_0$ 与 $n\ge n^*(\omega)$. 我们对 $m>n$ 归纳证明: 若 $s,t\in D_m$, $0<t-s<2^{-n}$, 则
\begin{equation}\label{eq:kc-chain}
\abs{X_t-X_s}\le2\sum_{j=n+1}^m2^{-\gamma j}.
\end{equation}
$m=n+1$ 时, $t-s$ 是 $2^{-(n+1)}$ 的正整数倍且小于 $2^{-n}$, 只能等于 $2^{-(n+1)}$, 由 \eqref{eq:kc-adjacent} (把 $n$ 换成 $n+1$) 得 $\abs{X_t-X_s}<2^{-\gamma(n+1)}$. 设 \eqref{eq:kc-chain} 对 $m$ 成立, 取 $s,t\in D_{m+1}$, $0<t-s<2^{-n}$. 令 $s'=\min\{u\in D_m: u\ge s\}$, $t'=\max\{u\in D_m: u\le t\}$. 区间 $[s,t]$ 中必含 $D_m$ 的点: 否则 $s,t$ 落在 $D_m$ 的两个相邻点之间的开区间内, 而这个开区间中只有一个 $D_{m+1}$ 的点 (中点), 与 $s<t$ 矛盾. 所以 $s\le s'\le t'\le t$, 且 $s'-s\le2^{-(m+1)}$, $t-t'\le2^{-(m+1)}$. 由于 $s,s'\in D_{m+1}$ 相等或相邻, 由 \eqref{eq:kc-adjacent} 得 $\abs{X_{s'}-X_s}\le2^{-\gamma(m+1)}$, 同理 $\abs{X_t-X_{t'}}\le2^{-\gamma(m+1)}$. 又 $0\le t'-s'\le t-s<2^{-n}$, 由归纳假设 (或 $s'=t'$), $\abs{X_{t'}-X_{s'}}\le2\sum_{j=n+1}^m2^{-\gamma j}$. 于是
\[
\abs{X_t-X_s}\le\abs{X_{s'}-X_s}+\abs{X_{t'}-X_{s'}}+\abs{X_t-X_{t'}}\le2\sum_{j=n+1}^m2^{-\gamma j}+2\cdot2^{-\gamma(m+1)}=2\sum_{j=n+1}^{m+1}2^{-\gamma j}.
\]
归纳完成. 由于 $D=\bigcup_mD_m$, 对任意 $s,t\in D$ 且 $0<t-s<2^{-n}$,
\[
\abs{X_t-X_s}\le2\sum_{j=n+1}^\infty2^{-\gamma j}=\frac{2\cdot2^{-\gamma(n+1)}}{1-2^{-\gamma}}.
\]

\textbf{第三步: 二进点上的 H\"older 连续性.} 记 $c_\gamma=\frac{2}{1-2^{-\gamma}}$. 设 $s,t\in D$, $0<t-s<2^{-n^*}$. 取 $n\ge n^*$ 使 $2^{-(n+1)}\le t-s<2^{-n}$. 由第二步,
\begin{equation}\label{eq:kc-holder}
\abs{X_t-X_s}\le c_\gamma2^{-\gamma(n+1)}\le c_\gamma\abs{t-s}^\gamma.
\end{equation}
特别地, 对 $\omega\in\Omega_0$, $t\mapsto X_t(\omega)$ 在 $D$ 上一致连续.

\textbf{第四步: 连续延拓.} 对 $\omega\in\Omega_0$ 与 $t\in[0,1]$, 令 $\widetilde{X}_t(\omega)=\lim_{s\to t,\,s\in D}X_s(\omega)$; 对 $\omega\notin\Omega_0$, 令 $\widetilde{X}_t(\omega)=0$. 一致连续函数在稠密子集上的极限存在 (由 Cauchy 准则), 且延拓后的函数连续, 在 $D$ 上与 $X$ 一致, 并且 \eqref{eq:kc-holder} 通过取极限对一切 $s,t\in[0,1]$, $\abs{t-s}<2^{-n^*}$ 成立. 当 $\abs{t-s}\ge2^{-n^*}$ 时, $\abs{\widetilde{X}_t-\widetilde{X}_s}\le2\sup_{u\le1}\abs{\widetilde{X}_u}\le2\sup_{u\le1}\abs{\widetilde{X}_u}\,2^{\gamma n^*}\abs{t-s}^\gamma$, 其中上确界有限 (连续函数在紧集上有界). 取 $K_\gamma=\max\big(c_\gamma,\,2^{1+\gamma n^*}\sup_{u\le1}\abs{\widetilde{X}_u}\big)$ 即得 H\"older 估计. 每个 $\widetilde{X}_t$ 是随机变量: 取 $s_k\in D$, $s_k\to t$, 则 $\widetilde{X}_t=\lim_kX_{s_k}\1_{\Omega_0}$.

\textbf{第五步: $\widetilde{X}$ 是修正.} 固定 $t\in[0,1]$, 取 $s_k\in D$, $s_k\to t$. 由 Markov 不等式与 \eqref{eq:kc-condition}, 对 $\varepsilon>0$,
\[
\P\big(\abs{X_{s_k}-X_t}>\varepsilon\big)\le\varepsilon^{-\alpha}C\abs{s_k-t}^{1+\beta}\to0,
\]
即 $X_{s_k}\xrightarrow{\P}X_t$. 另一方面 $X_{s_k}=\widetilde{X}_{s_k}\to\widetilde{X}_t$ 在 $\Omega_0$ 上成立. 由定理~\ref{thm:partial-converses} (1), 存在子列 $X_{s_{k_l}}\to X_t$ $\as$, 而它同时几乎必然收敛到 $\widetilde{X}_t$, 所以 $\widetilde{X}_t=X_t$ $\as$
\end{proof}

\begin{figure}[htbp]
\centering
\begin{tikzpicture}[x=12cm,y=1cm]
  \draw[SADark!60,->] (-0.02,0) -- (1.04,0);
  \foreach \k in {0,...,16}{\draw[SADark!40] (\k/16,0.08) -- (\k/16,-0.08);}
  \foreach \k in {0,...,8}{\draw[SADark!60] (\k/8,0.13) -- (\k/8,-0.13);}
  \foreach \k in {0,...,4}{\draw[SADark] (\k/4,0.2) -- (\k/4,-0.2);}
  \node[font=\small,below] at (0,-0.22) {$0$};
  \node[font=\small,below] at (1,-0.22) {$1$};
  \node[font=\scriptsize,SADark!70,anchor=west] at (1.05,0.2) {$D_2$};
  \node[font=\scriptsize,SADark!70,anchor=west] at (1.05,0.0) {$D_3$};
  \node[font=\scriptsize,SADark!70,anchor=west] at (1.05,-0.2) {$D_4$};
  \fill[SARed] (5/16,0) circle (2pt) node[below=4pt,font=\small] {$s$};
  \fill[SARed] (13/16,0) circle (2pt) node[below=4pt,font=\small] {$t$};
  \fill[SABlue] (3/8,0) circle (1.6pt) node[below=4pt,font=\small] {$s'$};
  \fill[SABlue] (3/4,0) circle (1.6pt) node[below=4pt,font=\small] {$t'$};
  \draw[SARed,thick] (5/16,0.1) to[out=70,in=110] node[above,font=\scriptsize]{$D_4$} (3/8,0.1);
  \draw[SABlue,thick] (3/8,0.1) to[out=60,in=120] node[above,font=\scriptsize]{$D_3$} (1/2,0.1);
  \draw[SAGreenText,thick] (1/2,0.1) to[out=45,in=135] node[above,font=\scriptsize]{$D_2$} (3/4,0.1);
  \draw[SARed,thick] (3/4,0.1) to[out=70,in=110] node[above,font=\scriptsize]{$D_4$} (13/16,0.1);
\end{tikzpicture}
\caption{链式论证 ($s,t\in D_4$): 先走到 $D_3$ 中的点 $s'$, $t'$, 再逐层变粗. 每一层至多用两段相邻点之间的增量, 第 $j$ 层的每段不超过 $2^{-\gamma j}$, 总和被几何级数 $2\sum_{j>n}2^{-\gamma j}$ 控制}
\label{fig:chaining}
\end{figure}

\begin{remark}
条件中的 ``$1+\beta$'' 不能换成 ``$1$''. 例如强度为 $\lambda$ 的 Poisson 过程 $N$ 满足 $\E\abs{N_t-N_s}=\lambda\abs{t-s}$, 而且对任意 $\alpha>0$, 当 $\abs{t-s}\to0$ 时 $\E\abs{N_t-N_s}^\alpha\sim\lambda\abs{t-s}$ (因为 $\abs{N_t-N_s}\ge1$ 的概率约为 $\lambda\abs{t-s}$). 它的轨道有跳跃, 没有连续修正. 指数 $1+\beta$ 中多出来的 $\beta$ 正好抵消第一步中对 $2^n$ 个区间取并的代价.
\end{remark}

% ---------------------------------------------------------------------
\subsection{布朗运动的存在性}

\begin{theorem}[布朗运动的存在性]\label{thm:bm-existence}
存在标准布朗运动. 进一步, 对任意 $\gamma<\frac{1}{2}$ 与 $T>0$, 布朗运动的轨道在 $[0,T]$ 上几乎必然是 $\gamma$ 阶 H\"older 连续的.
\end{theorem}

\begin{proof}
由推论~\ref{cor:gaussian-process} 与例~\ref{ex:min-kernel}, 存在均值为零、协方差为 $s\wedge t$ 的高斯过程 $X=\{X_t\}_{t\ge0}$. 对 $s<t$, $X_t-X_s$ 是正态的, 均值为零, 方差为 $t-2s+s=t-s$, 所以 $X_t-X_s\overset{d}{=}\sqrt{t-s}\,Z$, $Z\sim N(0,1)$. 对任意正整数 $p$,
\[
\E\abs{X_t-X_s}^{2p}=\E\abs{Z}^{2p}\cdot\abs{t-s}^p=(2p-1)!!\,\abs{t-s}^p,
\]
其中 $\E Z^{2p}=(2p-1)!!$ 来自例~\ref{ex:normal-mgf}. 取 $p\ge2$, 条件 \eqref{eq:kc-condition} 对 $\alpha=2p$, $\beta=p-1$ 成立.

对每个正整数 $N$, 由定理~\ref{thm:kolmogorov-chentsov}, $\{X_t\}_{t\in[0,N]}$ 有连续修正 $Y^{(N)}$. $Y^{(N)}$ 与 $Y^{(N+1)}$ 在 $[0,N]$ 上都是 $X$ 的连续修正, 从而互为修正, 由命题~\ref{prop:modification} (2) 不可区分. 于是存在概率为 $1$ 的事件 $\Omega_1$, 在其上对一切 $N$ 与 $t\le N$, $Y^{(N)}_t=Y^{(N+1)}_t$, 并且 $Y^{(1)}_0=X_0=0$ (注意 $\Var(X_0)=0$). 在 $\Omega_1$ 上令 $B_t=Y^{(N)}_t$ ($N\ge t$ 任取, 结果与 $N$ 无关), 在 $\Omega_1^c$ 上令 $B_t\equiv0$. 则 $B$ 的每条轨道连续, $B_0=0$, 且 $B$ 是 $X$ 的修正, 由命题~\ref{prop:modification} (1) 与 $X$ 有相同的有限维分布, 因而是均值为零、协方差为 $s\wedge t$ 的高斯过程. 由定理~\ref{thm:bm-gaussian}, $B$ 是标准布朗运动.

H\"older 连续性: 对 $\alpha=2p$, $\beta=p-1$, 定理~\ref{thm:kolmogorov-chentsov} 允许的阶数为 $\gamma<\frac{p-1}{2p}=\frac{1}{2}-\frac{1}{2p}$. 对给定的 $\gamma<\frac{1}{2}$, 取 $p$ 充分大即可. (定理中的 H\"older 性对 $Y^{(N)}$ 成立, 而 $B$ 与 $Y^{(N)}$ 在 $[0,N]$ 上不可区分.)
\end{proof}

4.3 节将证明, $\gamma<\frac{1}{2}$ 的限制是本质的: 布朗运动的轨道在任何区间上都不是 $\frac{1}{2}$ 阶 H\"older 连续的 (更不可能可微).

\begin{remark}[Wiener 测度与典则布朗运动]
记 $C[0,\infty)$ 为 $[0,\infty)$ 上连续函数全体, $\mathcal{C}$ 为由坐标映射 $w\mapsto w(t)$ ($t\ge0$) 生成的 $\sigma$-代数. 布朗运动可以看成映射 $B:\Omega\to C[0,\infty)$, $\omega\mapsto B_\cdot(\omega)$; 由引理~\ref{lem:measurable-generator}, 它是 $\mathcal{F}/\mathcal{C}$ 可测的 (只需验证每个坐标 $B_t$ 可测). 其像测度 $\mathbb{W}=\P\circ B^{-1}$ 称为 \textbf{Wiener 测度}. 由于 $\mathcal{C}$ 由柱集 $\pi$-系生成, $\mathbb{W}$ 由有限维分布唯一决定 (推论~\ref{cor:uniqueness}), 与 $B$ 的具体构造无关. 在 $(C[0,\infty),\mathcal{C},\mathbb{W})$ 上, 坐标过程 $B_t(w)=w(t)$ 是标准布朗运动, 称为\textbf{典则布朗运动}. 对 $x\in\R$, 记 $\P_x$ 为 $x+B$ 在 $C[0,\infty)$ 上的分布, $\E_x$ 为相应的期望, 称为\textbf{从 $x$ 出发的布朗运动}.
\end{remark}

% ---------------------------------------------------------------------
\subsection{L\'evy-Ciesielski 构造}

Kolmogorov 连续性定理给出的是存在性. 另有一种直接的构造: 把布朗运动展开成随机系数的函数级数. 设 $\{h_n\}$ 是 $L^2[0,1]$ 的一组标准正交基, 令 $B_t=\sum_nZ_n\int_0^th_n(u)\dd u$, 其中 $Z_n$ 是 $\iid$ $N(0,1)$. 形式上, 由 Parseval 等式,
\[
\E[B_sB_t]=\sum_n\int_0^sh_n\dd u\int_0^th_n\dd u=\sum_n\ip{\1_{[0,s]},h_n}\ip{\1_{[0,t]},h_n}=\ip{\1_{[0,s]},\1_{[0,t]}}=s\wedge t,
\]
正好是布朗运动的协方差. 用 Haar 基 $h_0\equiv1$, $h_{n,k}=2^{n/2}\big(\1_{[k2^{-n},(k+\frac{1}{2})2^{-n})}-\1_{[(k+\frac{1}{2})2^{-n},(k+1)2^{-n})}\big)$ ($n\ge0$, $0\le k<2^n$), 其积分是一族 ``帐篷函数'' (Schauder 函数, 图~\ref{fig:schauder} 左): $\int_0^th_{n,k}\dd u$ 支撑在 $[k2^{-n},(k+1)2^{-n}]$ 上, 在中点取最大值 $2^{-n/2-1}$.

\begin{theorem}[L\'evy-Ciesielski]\label{thm:levy-ciesielski}
设 $Z_0$, $Z_{n,k}$ ($n\ge0$, $0\le k<2^n$) 是 $\iid$ 的 $N(0,1)$ 随机变量. 则级数
\[
B_t=Z_0\,t+\sum_{n=0}^\infty\sum_{k=0}^{2^n-1}Z_{n,k}\int_0^th_{n,k}(u)\dd u
\]
几乎必然在 $[0,1]$ 上一致收敛, 其和是 $[0,1]$ 上的标准布朗运动. (证明见~\cite{karatzas} 2.3 节.)
\end{theorem}

一致收敛的关键在于第 $n$ 层的 $2^n$ 个帐篷函数支撑不交, 所以该层之和的一致范数为 $2^{-n/2-1}\max_k\abs{Z_{n,k}}$; 而由正态尾估计~\ref{thm:mills}, $\max_k\abs{Z_{n,k}}$ 几乎必然只以 $\sqrt{n}$ 的速度增长, 于是 $\sum_n2^{-n/2}\sqrt{n}<\infty$. 前 $n$ 层的部分和恰是布朗运动在 $D_n$ 上取值的线性插值 (图~\ref{fig:schauder} 右). 这个构造在概率上有清晰的含义: 已知 $B$ 在区间两端的值时, 中点的值等于两端平均值加上独立的正态扰动, 其方差为区间长度的四分之一 (这正是例~\ref{ex:bridge-preview} 中的条件方差 $\frac{s(t-s)}{t}$ 取 $s=\frac{t}{2}$). 逐层加细, 就从直线 $Z_0t$ 长出了整条布朗轨道.

\begin{figure}[htbp]
\centering
\begin{tikzpicture}
\begin{axis}[SAaxis, axis x line=bottom, axis y line=left, width=6.4cm, height=5.2cm, xmin=0, xmax=1.05, ymin=0, ymax=1.08,
  xtick={0,0.25,0.5,0.75,1}, xticklabels={$0$,$\frac14$,$\frac12$,$\frac34$,$1$}, ytick={0.25,0.5,1}, yticklabels={$\frac14$,$\frac12$,$1$}, xlabel={$t$}, clip=false]
  \addplot[SADark,thick] coordinates {(0,0) (1,1)};
  \addplot[SARed,thick] coordinates {(0,0) (0.5,0.5) (1,0)};
  \addplot[SABlue,thick] coordinates {(0,0) (0.25,0.354) (0.5,0) (0.75,0.354) (1,0)};
  \addplot[SAGreenText,thick] coordinates {(0,0) (0.125,0.25) (0.25,0) (0.375,0.25) (0.5,0) (0.625,0.25) (0.75,0) (0.875,0.25) (1,0)};
  \node[font=\scriptsize,SADark,anchor=south east] at (axis cs:0.8,0.8) {$t$};
  \node[font=\scriptsize,SARed] at (axis cs:0.5,0.57) {$n=0$};
  \node[font=\scriptsize,SABlue] at (axis cs:0.75,0.42) {$n=1$};
  \node[font=\scriptsize,SAGreenText] at (axis cs:0.94,0.3) {$n=2$};
\end{axis}
\end{tikzpicture}\hfill
\begin{tikzpicture}
\begin{axis}[SAaxis, axis x line=bottom, axis y line=left, width=8cm, height=5.2cm, xmin=0, xmax=1.03, ymin=-0.2, ymax=0.9,
  xtick={0,0.25,0.5,0.75,1}, ytick={0,0.5}, xlabel={$t$}, legend style={font=\scriptsize,draw=none,fill=none,at={(0.99,0.99)},anchor=north east}, legend cell align=left]
  \addplot[SADark!25,thin,forget plot] coordinates {(0.000,0.000) (0.002,0.002) (0.004,0.034) (0.006,-0.039) (0.008,-0.005) (0.010,-0.010) (0.012,-0.050) (0.014,-0.043) (0.016,-0.060) (0.018,-0.082) (0.020,-0.140) (0.021,-0.082) (0.023,-0.102) (0.025,-0.056) (0.027,-0.039) (0.029,-0.064) (0.031,-0.016) (0.033,0.052) (0.035,0.043) (0.037,0.039) (0.039,0.017) (0.041,0.033) (0.043,-0.006) (0.045,0.054) (0.047,0.114) (0.049,0.035) (0.051,0.094) (0.053,0.097) (0.055,0.178) (0.057,0.163) (0.059,0.177) (0.061,0.197) (0.062,0.155) (0.064,0.193) (0.066,0.183) (0.068,0.102) (0.070,0.024) (0.072,0.013) (0.074,0.041) (0.076,0.113) (0.078,0.134) (0.080,0.161) (0.082,0.188) (0.084,0.184) (0.086,0.180) (0.088,0.208) (0.090,0.228) (0.092,0.214) (0.094,0.222) (0.096,0.124) (0.098,0.112) (0.100,0.195) (0.102,0.277) (0.104,0.358) (0.105,0.401) (0.107,0.363) (0.109,0.305) (0.111,0.340) (0.113,0.215) (0.115,0.254) (0.117,0.182) (0.119,0.166) (0.121,0.066) (0.123,0.057) (0.125,0.004) (0.127,0.050) (0.129,0.094) (0.131,0.120) (0.133,0.164) (0.135,0.223) (0.137,0.258) (0.139,0.333) (0.141,0.338) (0.143,0.275) (0.145,0.249) (0.146,0.238) (0.148,0.296) (0.150,0.314) (0.152,0.247) (0.154,0.247) (0.156,0.273) (0.158,0.300) (0.160,0.282) (0.162,0.279) (0.164,0.282) (0.166,0.262) (0.168,0.303) (0.170,0.295) (0.172,0.242) (0.174,0.220) (0.176,0.286) (0.178,0.313) (0.180,0.288) (0.182,0.311) (0.184,0.338) (0.186,0.358) (0.188,0.364) (0.189,0.346) (0.191,0.417) (0.193,0.474) (0.195,0.452) (0.197,0.423) (0.199,0.504) (0.201,0.518) (0.203,0.461) (0.205,0.438) (0.207,0.444) (0.209,0.403) (0.211,0.457) (0.213,0.456) (0.215,0.494) (0.217,0.488) (0.219,0.517) (0.221,0.549) (0.223,0.568) (0.225,0.562) (0.227,0.562) (0.229,0.533) (0.230,0.566) (0.232,0.542) (0.234,0.522) (0.236,0.512) (0.238,0.465) (0.240,0.447) (0.242,0.483) (0.244,0.423) (0.246,0.445) (0.248,0.404) (0.250,0.358) (0.252,0.421) (0.254,0.420) (0.256,0.423) (0.258,0.432) (0.260,0.326) (0.262,0.278) (0.264,0.325) (0.266,0.361) (0.268,0.397) (0.270,0.392) (0.271,0.482) (0.273,0.512) (0.275,0.536) (0.277,0.540) (0.279,0.500) (0.281,0.493) (0.283,0.451) (0.285,0.394) (0.287,0.300) (0.289,0.429) (0.291,0.375) (0.293,0.431) (0.295,0.458) (0.297,0.398) (0.299,0.440) (0.301,0.454) (0.303,0.442) (0.305,0.405) (0.307,0.479) (0.309,0.490) (0.311,0.467) (0.312,0.504) (0.314,0.496) (0.316,0.491) (0.318,0.484) (0.320,0.513) (0.322,0.597) (0.324,0.634) (0.326,0.609) (0.328,0.537) (0.330,0.499) (0.332,0.479) (0.334,0.425) (0.336,0.458) (0.338,0.424) (0.340,0.501) (0.342,0.523) (0.344,0.599) (0.346,0.607) (0.348,0.633) (0.350,0.615) (0.352,0.656) (0.354,0.677) (0.355,0.764) (0.357,0.791) (0.359,0.760) (0.361,0.720) (0.363,0.702) (0.365,0.662) (0.367,0.595) (0.369,0.570) (0.371,0.556) (0.373,0.532) (0.375,0.430) (0.377,0.441) (0.379,0.478) (0.381,0.389) (0.383,0.350) (0.385,0.272) (0.387,0.265) (0.389,0.287) (0.391,0.234) (0.393,0.363) (0.395,0.376) (0.396,0.271) (0.398,0.329) (0.400,0.382) (0.402,0.387) (0.404,0.439) (0.406,0.417) (0.408,0.347) (0.410,0.370) (0.412,0.285) (0.414,0.289) (0.416,0.261) (0.418,0.247) (0.420,0.277) (0.422,0.325) (0.424,0.351) (0.426,0.337) (0.428,0.341) (0.430,0.253) (0.432,0.252) (0.434,0.302) (0.436,0.227) (0.438,0.148) (0.439,0.176) (0.441,0.246) (0.443,0.312) (0.445,0.288) (0.447,0.340) (0.449,0.438) (0.451,0.465) (0.453,0.524) (0.455,0.505) (0.457,0.512) (0.459,0.518) (0.461,0.488) (0.463,0.448) (0.465,0.489) (0.467,0.498) (0.469,0.528) (0.471,0.557) (0.473,0.612) (0.475,0.587) (0.477,0.585) (0.479,0.559) (0.480,0.528) (0.482,0.487) (0.484,0.538) (0.486,0.510) (0.488,0.525) (0.490,0.431) (0.492,0.464) (0.494,0.493) (0.496,0.533) (0.498,0.495) (0.500,0.537) (0.502,0.566) (0.504,0.579) (0.506,0.638) (0.508,0.542) (0.510,0.507) (0.512,0.400) (0.514,0.370) (0.516,0.376) (0.518,0.372) (0.520,0.348) (0.521,0.243) (0.523,0.242) (0.525,0.214) (0.527,0.216) (0.529,0.217) (0.531,0.190) (0.533,0.202) (0.535,0.183) (0.537,0.104) (0.539,0.048) (0.541,0.082) (0.543,0.041) (0.545,0.066) (0.547,-0.017) (0.549,-0.027) (0.551,-0.044) (0.553,0.015) (0.555,0.051) (0.557,0.030) (0.559,0.059) (0.561,0.035) (0.562,0.035) (0.564,0.067) (0.566,0.108) (0.568,0.166) (0.570,0.183) (0.572,0.199) (0.574,0.202) (0.576,0.184) (0.578,0.165) (0.580,0.257) (0.582,0.206) (0.584,0.172) (0.586,0.193) (0.588,0.179) (0.590,0.171) (0.592,0.194) (0.594,0.231) (0.596,0.200) (0.598,0.250) (0.600,0.313) (0.602,0.348) (0.604,0.297) (0.605,0.345) (0.607,0.370) (0.609,0.348) (0.611,0.347) (0.613,0.308) (0.615,0.280) (0.617,0.245) (0.619,0.248) (0.621,0.296) (0.623,0.290) (0.625,0.269) (0.627,0.274) (0.629,0.341) (0.631,0.358) (0.633,0.299) (0.635,0.355) (0.637,0.374) (0.639,0.395) (0.641,0.450) (0.643,0.386) (0.645,0.374) (0.646,0.434) (0.648,0.477) (0.650,0.412) (0.652,0.419) (0.654,0.472) (0.656,0.469) (0.658,0.469) (0.660,0.441) (0.662,0.378) (0.664,0.387) (0.666,0.345) (0.668,0.367) (0.670,0.379) (0.672,0.381) (0.674,0.321) (0.676,0.368) (0.678,0.294) (0.680,0.225) (0.682,0.266) (0.684,0.273) (0.686,0.270) (0.688,0.254) (0.689,0.185) (0.691,0.109) (0.693,0.025) (0.695,-0.002) (0.697,0.027) (0.699,-0.002) (0.701,-0.002) (0.703,0.083) (0.705,0.114) (0.707,0.161) (0.709,0.186) (0.711,0.121) (0.713,0.158) (0.715,0.292) (0.717,0.269) (0.719,0.285) (0.721,0.377) (0.723,0.460) (0.725,0.520) (0.727,0.551) (0.729,0.439) (0.730,0.418) (0.732,0.388) (0.734,0.391) (0.736,0.393) (0.738,0.328) (0.740,0.238) (0.742,0.253) (0.744,0.244) (0.746,0.253) (0.748,0.285) (0.750,0.237) (0.752,0.172) (0.754,0.210) (0.756,0.155) (0.758,0.151) (0.760,0.119) (0.762,0.146) (0.764,0.141) (0.766,0.086) (0.768,-0.008) (0.770,0.033) (0.771,0.052) (0.773,0.079) (0.775,0.142) (0.777,0.137) (0.779,0.091) (0.781,0.091) (0.783,0.100) (0.785,0.071) (0.787,0.053) (0.789,0.030) (0.791,0.035) (0.793,0.098) (0.795,0.069) (0.797,0.062) (0.799,0.054) (0.801,0.079) (0.803,0.039) (0.805,0.009) (0.807,0.045) (0.809,0.041) (0.811,0.100) (0.812,0.152) (0.814,0.143) (0.816,0.199) (0.818,0.250) (0.820,0.239) (0.822,0.228) (0.824,0.209) (0.826,0.159) (0.828,0.181) (0.830,0.241) (0.832,0.203) (0.834,0.250) (0.836,0.223) (0.838,0.201) (0.840,0.269) (0.842,0.192) (0.844,0.149) (0.846,0.173) (0.848,0.149) (0.850,0.088) (0.852,0.184) (0.854,0.184) (0.855,0.240) (0.857,0.261) (0.859,0.270) (0.861,0.230) (0.863,0.230) (0.865,0.255) (0.867,0.256) (0.869,0.229) (0.871,0.166) (0.873,0.173) (0.875,0.234) (0.877,0.263) (0.879,0.257) (0.881,0.202) (0.883,0.170) (0.885,0.205) (0.887,0.134) (0.889,0.106) (0.891,0.071) (0.893,0.054) (0.895,0.085) (0.896,-0.002) (0.898,-0.030) (0.900,-0.079) (0.902,-0.051) (0.904,-0.014) (0.906,0.043) (0.908,0.083) (0.910,0.039) (0.912,0.075) (0.914,0.038) (0.916,0.099) (0.918,0.082) (0.920,0.117) (0.922,0.178) (0.924,0.177) (0.926,0.148) (0.928,0.170) (0.930,0.165) (0.932,0.153) (0.934,0.120) (0.936,0.153) (0.938,0.097) (0.939,0.113) (0.941,0.150) (0.943,0.152) (0.945,0.233) (0.947,0.288) (0.949,0.387) (0.951,0.296) (0.953,0.265) (0.955,0.257) (0.957,0.224) (0.959,0.164) (0.961,0.062) (0.963,-0.003) (0.965,0.030) (0.967,0.089) (0.969,0.089) (0.971,0.063) (0.973,0.124) (0.975,-0.010) (0.977,-0.013) (0.979,-0.037) (0.980,-0.009) (0.982,-0.034) (0.984,0.042) (0.986,0.029) (0.988,0.001) (0.990,-0.028) (0.992,-0.026) (0.994,0.105) (0.996,0.151) (0.998,0.129) (1.000,0.088)};
  \addplot[SARed,thick] coordinates {(0.000,0.000) (0.500,0.537) (1.000,0.088)};
  \addplot[SABlue,thick] coordinates {(0.000,0.000) (0.250,0.358) (0.500,0.537) (0.750,0.237) (1.000,0.088)};
  \addplot[SAGreenText,thick] coordinates {(0.000,0.000) (0.062,0.155) (0.125,0.004) (0.188,0.364) (0.250,0.358) (0.312,0.504) (0.375,0.430) (0.438,0.148) (0.500,0.537) (0.562,0.035) (0.625,0.269) (0.688,0.254) (0.750,0.237) (0.812,0.152) (0.875,0.234) (0.938,0.097) (1.000,0.088)};
  \legend{$D_1$, $D_2$, $D_4$}
\end{axis}
\end{tikzpicture}
\caption{左: Schauder 函数 $t$ 与第 $n=0,1,2$ 层的帐篷函数 $\int_0^th_{n,k}\dd u$, 第 $n$ 层的高度为 $2^{-n/2-1}$. 右: L\'evy-Ciesielski 级数的部分和, 即布朗轨道 (灰色) 在 $D_1,D_2,D_4$ 上的线性插值}
\label{fig:schauder}
\end{figure}

% ---------------------------------------------------------------------
\subsection{多维布朗运动}

\begin{definition}[$d$ 维布朗运动]\label{def:bm-d}
设 $B^1,\dots,B^d$ 是相互独立的标准布朗运动. 称 $\rvec{B}_t=(B^1_t,\dots,B^d_t)^\T$ 为 \textbf{$d$ 维标准布朗运动}. 对 $\rvec{x}\in\R^d$, 称 $\rvec{x}+\rvec{B}$ 为从 $\rvec{x}$ 出发的 $d$ 维布朗运动.
\end{definition}

存在性: 在乘积空间上取 $d$ 个独立的布朗运动即可 (或者由一个布朗运动在 $[0,1]$ 上的 L\'evy-Ciesielski 构造, 把可列个 $\iid$ 正态变量分成 $d$ 组). 与定理~\ref{thm:bm-gaussian} 的证明相同, $\rvec{B}$ 由以下性质刻画: $\rvec{B}_0=\rvec{0}$, 轨道连续, 增量独立, 且 $\rvec{B}_t-\rvec{B}_s\sim N(\rvec{0},(t-s)\mat{I}_d)$.

\begin{proposition}[旋转不变性]\label{prop:bm-rotation}
设 $\rvec{B}$ 是 $d$ 维标准布朗运动, $\mat{O}$ 是 $d\times d$ 正交矩阵. 则 $\mat{O}\rvec{B}$ 也是 $d$ 维标准布朗运动. 更一般地, 若 $\mat{A}$ 是 $m\times d$ 矩阵, 则 $\rvec{X}=\mat{A}\rvec{B}$ 的分量 $X^i$ 是 (未必独立的) 布朗运动的常数倍, 且 $\Cov(X^i_s,X^j_t)=(s\wedge t)(\mat{A}\mat{A}^\T)_{ij}$.
\end{proposition}

\begin{proof}
$\mat{A}\rvec{B}$ 的轨道连续, 从 $\rvec{0}$ 出发, 有限个时刻上的取值是正态向量 $(\rvec{B}_{t_1},\dots,\rvec{B}_{t_n})$ 的线性变换, 因而是联合正态的. 由各分量的独立性, $\Cov(B^k_s,B^l_t)=(s\wedge t)\delta_{kl}$, 所以
\[
\Cov(X^i_s,X^j_t)=\sum_{k,l}A_{ik}A_{jl}\Cov(B^k_s,B^l_t)=(s\wedge t)\sum_kA_{ik}A_{jk}=(s\wedge t)(\mat{A}\mat{A}^\T)_{ij}.
\]
当 $\mat{A}=\mat{O}$ 正交时, $\mat{O}\mat{O}^\T=\mat{I}_d$. 于是 $\mat{O}\rvec{B}$ 的分量都是标准布朗运动 (定理~\ref{thm:bm-gaussian}), 且不同分量的 (所有时刻的) 取值不相关; 由于它们是联合正态的, 由定理~\ref{thm:gaussian-uncorrelated}, 不同分量的有限维取值相互独立, 再由 $\pi$-系判别准则~\ref{thm:indep-pi}, 各分量作为过程相互独立.
\end{proof}

例如, 若 $W^1,W^2$ 是独立的标准布朗运动, $\rho\in[-1,1]$, 则 $W^1$ 与 $\rho W^1+\sqrt{1-\rho^2}\,W^2$ 是两个相关系数为 $\rho$ 的标准布朗运动. 这是金融中构造相关资产价格的标准方法. 图~\ref{fig:bm-2d} 画出了平面布朗运动的一条轨道.

\begin{figure}[htbp]
\centering
\begin{tikzpicture}
\begin{axis}[SAaxis, axis lines=box, axis line style={-,SADark!60}, width=9cm, height=7cm, xmin=-1.9, xmax=0.3, ymin=-0.7, ymax=1.1, axis equal image,
  xtick={-1.5,-1,-0.5,0}, ytick={-0.5,0,0.5,1}, xlabel={$B^1_t$}, ylabel={$B^2_t$}]
  \addplot[SABlue,very thin] coordinates {(0.000,0.000) (-0.016,0.001) (-0.033,0.009) (-0.018,0.019) (-0.034,0.002) (-0.079,0.053) (-0.092,0.041) (-0.041,0.009) (-0.045,0.033) (-0.067,0.037) (-0.088,0.013) (-0.066,0.018) (-0.079,-0.016) (-0.078,-0.032) (-0.094,-0.013) (-0.074,-0.022) (-0.089,-0.024) (-0.125,-0.009) (-0.143,-0.033) (-0.147,0.010) (-0.169,0.009) (-0.205,-0.022) (-0.182,-0.020) (-0.169,-0.010) (-0.157,0.007) (-0.164,0.032) (-0.186,0.045) (-0.138,0.078) (-0.170,0.072) (-0.206,0.063) (-0.207,0.047) (-0.228,0.050) (-0.219,0.054) (-0.197,0.031) (-0.170,0.043) (-0.168,0.033) (-0.163,0.013) (-0.194,0.015) (-0.175,-0.003) (-0.173,-0.022) (-0.168,-0.024) (-0.134,0.027) (-0.179,-0.004) (-0.170,-0.013) (-0.172,0.017) (-0.173,0.018) (-0.197,0.023) (-0.168,-0.015) (-0.165,-0.000) (-0.173,-0.004) (-0.165,-0.027) (-0.192,-0.002) (-0.193,-0.015) (-0.174,-0.029) (-0.177,0.001) (-0.181,0.023) (-0.156,0.059) (-0.173,0.087) (-0.166,0.117) (-0.182,0.108) (-0.202,0.123) (-0.190,0.157) (-0.153,0.159) (-0.147,0.128) (-0.146,0.152) (-0.165,0.132) (-0.176,0.183) (-0.141,0.264) (-0.151,0.262) (-0.166,0.255) (-0.168,0.249) (-0.151,0.239) (-0.160,0.213) (-0.183,0.195) (-0.179,0.194) (-0.159,0.218) (-0.124,0.193) (-0.117,0.203) (-0.135,0.190) (-0.177,0.186) (-0.160,0.177) (-0.175,0.172) (-0.161,0.154) (-0.191,0.183) (-0.180,0.177) (-0.231,0.190) (-0.193,0.135) (-0.209,0.131) (-0.210,0.147) (-0.198,0.177) (-0.169,0.160) (-0.210,0.202) (-0.193,0.228) (-0.197,0.264) (-0.182,0.250) (-0.182,0.255) (-0.216,0.259) (-0.216,0.257) (-0.179,0.268) (-0.199,0.225) (-0.213,0.207) (-0.223,0.203) (-0.223,0.218) (-0.240,0.223) (-0.236,0.216) (-0.240,0.179) (-0.234,0.202) (-0.212,0.231) (-0.186,0.241) (-0.218,0.261) (-0.200,0.259) (-0.218,0.264) (-0.213,0.252) (-0.213,0.232) (-0.201,0.197) (-0.227,0.172) (-0.221,0.187) (-0.211,0.156) (-0.177,0.177) (-0.187,0.204) (-0.195,0.228) (-0.211,0.225) (-0.179,0.232) (-0.206,0.255) (-0.211,0.262) (-0.147,0.262) (-0.130,0.244) (-0.113,0.234) (-0.102,0.188) (-0.135,0.203) (-0.152,0.245) (-0.151,0.199) (-0.178,0.236) (-0.169,0.251) (-0.183,0.280) (-0.217,0.269) (-0.251,0.242) (-0.241,0.246) (-0.228,0.265) (-0.228,0.286) (-0.214,0.304) (-0.220,0.308) (-0.218,0.325) (-0.243,0.315) (-0.236,0.305) (-0.210,0.283) (-0.192,0.263) (-0.200,0.276) (-0.218,0.243) (-0.232,0.203) (-0.201,0.203) (-0.228,0.192) (-0.206,0.211) (-0.218,0.215) (-0.209,0.182) (-0.214,0.213) (-0.228,0.251) (-0.253,0.262) (-0.236,0.253) (-0.240,0.264) (-0.245,0.241) (-0.294,0.234) (-0.261,0.229) (-0.260,0.205) (-0.271,0.211) (-0.248,0.231) (-0.251,0.206) (-0.221,0.179) (-0.216,0.201) (-0.233,0.190) (-0.257,0.183) (-0.224,0.152) (-0.212,0.146) (-0.238,0.132) (-0.241,0.141) (-0.235,0.161) (-0.214,0.200) (-0.215,0.166) (-0.201,0.173) (-0.202,0.180) (-0.197,0.143) (-0.170,0.125) (-0.183,0.155) (-0.198,0.152) (-0.174,0.088) (-0.178,0.134) (-0.167,0.158) (-0.168,0.146) (-0.115,0.168) (-0.119,0.135) (-0.094,0.092) (-0.061,0.083) (-0.040,0.119) (-0.057,0.069) (-0.078,0.078) (-0.068,0.077) (-0.121,0.091) (-0.155,0.062) (-0.136,0.049) (-0.125,0.089) (-0.144,0.096) (-0.126,0.103) (-0.125,0.120) (-0.128,0.080) (-0.122,0.078) (-0.146,0.083) (-0.156,0.070) (-0.150,0.061) (-0.122,0.034) (-0.141,0.062) (-0.147,0.054) (-0.152,0.069) (-0.141,0.072) (-0.171,0.106) (-0.170,0.090) (-0.175,0.081) (-0.166,0.020) (-0.194,-0.006) (-0.195,-0.030) (-0.232,-0.034) (-0.208,-0.031) (-0.216,-0.024) (-0.228,-0.029) (-0.217,-0.051) (-0.234,-0.107) (-0.231,-0.136) (-0.246,-0.149) (-0.251,-0.149) (-0.245,-0.181) (-0.244,-0.197) (-0.236,-0.171) (-0.245,-0.171) (-0.258,-0.182) (-0.243,-0.183) (-0.260,-0.194) (-0.259,-0.189) (-0.224,-0.177) (-0.237,-0.145) (-0.289,-0.160) (-0.333,-0.193) (-0.296,-0.210) (-0.246,-0.215) (-0.257,-0.225) (-0.277,-0.216) (-0.258,-0.224) (-0.247,-0.239) (-0.276,-0.265) (-0.256,-0.253) (-0.257,-0.230) (-0.237,-0.200) (-0.246,-0.181) (-0.277,-0.167) (-0.299,-0.160) (-0.281,-0.133) (-0.276,-0.127) (-0.323,-0.117) (-0.276,-0.103) (-0.334,-0.119) (-0.370,-0.129) (-0.391,-0.124) (-0.376,-0.111) (-0.365,-0.153) (-0.348,-0.156) (-0.371,-0.138) (-0.370,-0.169) (-0.375,-0.213) (-0.367,-0.205) (-0.366,-0.192) (-0.380,-0.215) (-0.372,-0.209) (-0.384,-0.192) (-0.377,-0.214) (-0.355,-0.199) (-0.381,-0.243) (-0.371,-0.255) (-0.350,-0.225) (-0.368,-0.224) (-0.363,-0.208) (-0.384,-0.213) (-0.352,-0.199) (-0.379,-0.201) (-0.407,-0.243) (-0.378,-0.263) (-0.412,-0.239) (-0.411,-0.214) (-0.395,-0.208) (-0.381,-0.195) (-0.369,-0.171) (-0.348,-0.161) (-0.321,-0.135) (-0.292,-0.139) (-0.303,-0.125) (-0.278,-0.125) (-0.241,-0.087) (-0.291,-0.096) (-0.281,-0.083) (-0.338,-0.139) (-0.358,-0.127) (-0.374,-0.117) (-0.414,-0.084) (-0.378,-0.109) (-0.390,-0.127) (-0.422,-0.124) (-0.418,-0.143) (-0.396,-0.154) (-0.404,-0.144) (-0.460,-0.137) (-0.433,-0.110) (-0.404,-0.097) (-0.341,-0.106) (-0.376,-0.142) (-0.384,-0.124) (-0.356,-0.156) (-0.349,-0.134) (-0.364,-0.148) (-0.373,-0.154) (-0.375,-0.132) (-0.368,-0.146) (-0.392,-0.125) (-0.396,-0.090) (-0.414,-0.071) (-0.411,-0.102) (-0.426,-0.126) (-0.438,-0.115) (-0.442,-0.127) (-0.451,-0.161) (-0.443,-0.153) (-0.464,-0.157) (-0.430,-0.140) (-0.417,-0.158) (-0.457,-0.175) (-0.445,-0.171) (-0.447,-0.178) (-0.398,-0.209) (-0.386,-0.195) (-0.394,-0.172) (-0.361,-0.167) (-0.347,-0.204) (-0.337,-0.182) (-0.348,-0.172) (-0.315,-0.138) (-0.316,-0.160) (-0.331,-0.144) (-0.327,-0.125) (-0.331,-0.152) (-0.300,-0.115) (-0.293,-0.115) (-0.288,-0.147) (-0.248,-0.172) (-0.249,-0.190) (-0.291,-0.195) (-0.302,-0.184) (-0.328,-0.154) (-0.312,-0.140) (-0.321,-0.159) (-0.336,-0.143) (-0.329,-0.147) (-0.300,-0.146) (-0.324,-0.149) (-0.287,-0.136) (-0.278,-0.133) (-0.296,-0.118) (-0.303,-0.153) (-0.294,-0.123) (-0.284,-0.154) (-0.305,-0.165) (-0.325,-0.134) (-0.337,-0.183) (-0.322,-0.163) (-0.327,-0.141) (-0.307,-0.165) (-0.326,-0.159) (-0.347,-0.204) (-0.337,-0.228) (-0.342,-0.181) (-0.364,-0.170) (-0.344,-0.176) (-0.374,-0.177) (-0.371,-0.184) (-0.392,-0.170) (-0.434,-0.169) (-0.440,-0.172) (-0.446,-0.175) (-0.437,-0.194) (-0.438,-0.214) (-0.464,-0.202) (-0.494,-0.177) (-0.506,-0.171) (-0.501,-0.178) (-0.518,-0.199) (-0.543,-0.196) (-0.555,-0.158) (-0.583,-0.194) (-0.586,-0.195) (-0.557,-0.244) (-0.540,-0.243) (-0.512,-0.232) (-0.470,-0.188) (-0.464,-0.181) (-0.483,-0.206) (-0.511,-0.225) (-0.517,-0.218) (-0.515,-0.182) (-0.509,-0.192) (-0.481,-0.193) (-0.494,-0.213) (-0.481,-0.182) (-0.528,-0.179) (-0.562,-0.200) (-0.518,-0.174) (-0.543,-0.172) (-0.532,-0.177) (-0.487,-0.161) (-0.481,-0.170) (-0.480,-0.175) (-0.463,-0.201) (-0.447,-0.163) (-0.468,-0.154) (-0.477,-0.143) (-0.493,-0.151) (-0.503,-0.121) (-0.530,-0.143) (-0.506,-0.137) (-0.510,-0.115) (-0.514,-0.124) (-0.486,-0.141) (-0.476,-0.189) (-0.465,-0.165) (-0.490,-0.149) (-0.501,-0.161) (-0.506,-0.152) (-0.477,-0.177) (-0.451,-0.203) (-0.434,-0.218) (-0.439,-0.220) (-0.469,-0.239) (-0.470,-0.257) (-0.493,-0.246) (-0.505,-0.227) (-0.565,-0.215) (-0.591,-0.189) (-0.573,-0.199) (-0.538,-0.227) (-0.536,-0.211) (-0.509,-0.235) (-0.495,-0.254) (-0.517,-0.238) (-0.460,-0.259) (-0.412,-0.218) (-0.477,-0.209) (-0.486,-0.225) (-0.518,-0.237) (-0.510,-0.265) (-0.529,-0.246) (-0.534,-0.239) (-0.524,-0.265) (-0.591,-0.275) (-0.599,-0.291) (-0.567,-0.292) (-0.604,-0.319) (-0.631,-0.340) (-0.630,-0.330) (-0.616,-0.321) (-0.590,-0.313) (-0.570,-0.308) (-0.557,-0.319) (-0.588,-0.363) (-0.549,-0.324) (-0.577,-0.340) (-0.621,-0.331) (-0.636,-0.345) (-0.614,-0.357) (-0.581,-0.371) (-0.619,-0.357) (-0.639,-0.383) (-0.635,-0.352) (-0.618,-0.329) (-0.603,-0.319) (-0.602,-0.313) (-0.597,-0.340) (-0.597,-0.338) (-0.606,-0.367) (-0.597,-0.361) (-0.622,-0.386) (-0.637,-0.412) (-0.632,-0.429) (-0.655,-0.418) (-0.639,-0.387) (-0.621,-0.410) (-0.607,-0.388) (-0.653,-0.355) (-0.650,-0.331) (-0.646,-0.298) (-0.664,-0.271) (-0.672,-0.235) (-0.668,-0.249) (-0.662,-0.278) (-0.646,-0.302) (-0.648,-0.253) (-0.649,-0.221) (-0.626,-0.180) (-0.613,-0.200) (-0.638,-0.202) (-0.646,-0.199) (-0.630,-0.244) (-0.640,-0.226) (-0.616,-0.230) (-0.620,-0.230) (-0.634,-0.256) (-0.613,-0.197) (-0.614,-0.190) (-0.600,-0.199) (-0.609,-0.199) (-0.618,-0.212) (-0.616,-0.212) (-0.577,-0.236) (-0.557,-0.250) (-0.551,-0.247) (-0.541,-0.235) (-0.541,-0.224) (-0.587,-0.222) (-0.609,-0.210) (-0.621,-0.208) (-0.644,-0.239) (-0.631,-0.239) (-0.632,-0.234) (-0.637,-0.240) (-0.669,-0.220) (-0.640,-0.245) (-0.652,-0.243) (-0.661,-0.245) (-0.661,-0.272) (-0.678,-0.292) (-0.696,-0.303) (-0.732,-0.359) (-0.744,-0.352) (-0.778,-0.356) (-0.759,-0.338) (-0.752,-0.353) (-0.790,-0.368) (-0.787,-0.342) (-0.789,-0.338) (-0.814,-0.316) (-0.770,-0.292) (-0.782,-0.261) (-0.784,-0.276) (-0.800,-0.296) (-0.817,-0.316) (-0.838,-0.366) (-0.888,-0.376) (-0.883,-0.388) (-0.870,-0.406) (-0.875,-0.413) (-0.872,-0.419) (-0.905,-0.418) (-0.870,-0.462) (-0.908,-0.452) (-0.915,-0.471) (-0.900,-0.438) (-0.920,-0.448) (-0.897,-0.480) (-0.903,-0.490) (-0.936,-0.463) (-0.925,-0.514) (-0.914,-0.476) (-0.868,-0.474) (-0.866,-0.465) (-0.891,-0.459) (-0.910,-0.462) (-0.932,-0.474) (-0.889,-0.461) (-0.885,-0.467) (-0.820,-0.474) (-0.834,-0.451) (-0.833,-0.444) (-0.868,-0.475) (-0.867,-0.478) (-0.865,-0.468) (-0.816,-0.439) (-0.828,-0.420) (-0.818,-0.429) (-0.841,-0.444) (-0.873,-0.424) (-0.861,-0.402) (-0.901,-0.431) (-0.849,-0.414) (-0.838,-0.373) (-0.825,-0.375) (-0.785,-0.353) (-0.790,-0.347) (-0.842,-0.331) (-0.830,-0.346) (-0.810,-0.319) (-0.807,-0.336) (-0.802,-0.341) (-0.827,-0.348) (-0.804,-0.349) (-0.826,-0.352) (-0.837,-0.358) (-0.852,-0.321) (-0.849,-0.299) (-0.835,-0.307) (-0.852,-0.288) (-0.846,-0.296) (-0.832,-0.322) (-0.819,-0.282) (-0.836,-0.325) (-0.864,-0.339) (-0.870,-0.350) (-0.859,-0.332) (-0.836,-0.343) (-0.792,-0.337) (-0.829,-0.363) (-0.803,-0.361) (-0.836,-0.365) (-0.818,-0.367) (-0.810,-0.377) (-0.797,-0.367) (-0.787,-0.353) (-0.770,-0.323) (-0.775,-0.297) (-0.755,-0.301) (-0.740,-0.320) (-0.706,-0.303) (-0.709,-0.303) (-0.670,-0.325) (-0.676,-0.288) (-0.665,-0.301) (-0.669,-0.294) (-0.669,-0.309) (-0.665,-0.343) (-0.660,-0.381) (-0.687,-0.385) (-0.674,-0.369) (-0.685,-0.389) (-0.701,-0.403) (-0.724,-0.431) (-0.705,-0.458) (-0.721,-0.417) (-0.714,-0.415) (-0.726,-0.405) (-0.731,-0.373) (-0.726,-0.341) (-0.700,-0.317) (-0.721,-0.329) (-0.677,-0.305) (-0.650,-0.309) (-0.620,-0.317) (-0.583,-0.282) (-0.599,-0.309) (-0.609,-0.333) (-0.580,-0.317) (-0.579,-0.319) (-0.589,-0.298) (-0.562,-0.281) (-0.563,-0.267) (-0.597,-0.272) (-0.570,-0.298) (-0.633,-0.281) (-0.674,-0.290) (-0.719,-0.292) (-0.747,-0.247) (-0.782,-0.280) (-0.764,-0.308) (-0.798,-0.297) (-0.814,-0.337) (-0.796,-0.292) (-0.796,-0.295) (-0.796,-0.314) (-0.825,-0.297) (-0.794,-0.315) (-0.802,-0.312) (-0.800,-0.334) (-0.796,-0.324) (-0.792,-0.328) (-0.785,-0.325) (-0.780,-0.331) (-0.773,-0.341) (-0.775,-0.363) (-0.776,-0.376) (-0.738,-0.371) (-0.761,-0.369) (-0.773,-0.358) (-0.774,-0.364) (-0.776,-0.408) (-0.786,-0.359) (-0.752,-0.281) (-0.750,-0.265) (-0.741,-0.298) (-0.727,-0.299) (-0.763,-0.315) (-0.760,-0.289) (-0.715,-0.293) (-0.718,-0.312) (-0.709,-0.297) (-0.737,-0.307) (-0.723,-0.311) (-0.758,-0.310) (-0.746,-0.315) (-0.723,-0.301) (-0.701,-0.288) (-0.661,-0.305) (-0.692,-0.318) (-0.710,-0.331) (-0.720,-0.369) (-0.720,-0.375) (-0.702,-0.390) (-0.686,-0.385) (-0.630,-0.373) (-0.632,-0.345) (-0.645,-0.284) (-0.688,-0.298) (-0.668,-0.316) (-0.670,-0.305) (-0.651,-0.316) (-0.659,-0.331) (-0.684,-0.306) (-0.665,-0.327) (-0.657,-0.345) (-0.675,-0.339) (-0.693,-0.340) (-0.703,-0.354) (-0.711,-0.362) (-0.701,-0.366) (-0.734,-0.402) (-0.690,-0.412) (-0.691,-0.392) (-0.671,-0.403) (-0.645,-0.370) (-0.656,-0.334) (-0.637,-0.370) (-0.664,-0.387) (-0.677,-0.412) (-0.671,-0.418) (-0.685,-0.416) (-0.704,-0.403) (-0.683,-0.425) (-0.724,-0.401) (-0.730,-0.398) (-0.716,-0.396) (-0.702,-0.419) (-0.679,-0.433) (-0.675,-0.406) (-0.630,-0.428) (-0.634,-0.438) (-0.625,-0.444) (-0.603,-0.470) (-0.603,-0.460) (-0.590,-0.466) (-0.570,-0.499) (-0.565,-0.504) (-0.526,-0.501) (-0.497,-0.524) (-0.522,-0.502) (-0.509,-0.518) (-0.533,-0.524) (-0.518,-0.511) (-0.499,-0.525) (-0.486,-0.532) (-0.458,-0.490) (-0.491,-0.483) (-0.508,-0.432) (-0.478,-0.408) (-0.473,-0.438) (-0.462,-0.425) (-0.487,-0.399) (-0.443,-0.401) (-0.459,-0.393) (-0.494,-0.365) (-0.481,-0.384) (-0.494,-0.389) (-0.524,-0.396) (-0.509,-0.423) (-0.535,-0.403) (-0.522,-0.389) (-0.492,-0.396) (-0.512,-0.368) (-0.528,-0.370) (-0.495,-0.379) (-0.508,-0.405) (-0.479,-0.436) (-0.475,-0.408) (-0.497,-0.432) (-0.492,-0.402) (-0.479,-0.437) (-0.471,-0.392) (-0.468,-0.372) (-0.484,-0.380) (-0.497,-0.370) (-0.481,-0.373) (-0.431,-0.348) (-0.454,-0.345) (-0.471,-0.336) (-0.487,-0.387) (-0.507,-0.400) (-0.511,-0.410) (-0.562,-0.436) (-0.561,-0.461) (-0.573,-0.423) (-0.589,-0.417) (-0.587,-0.403) (-0.601,-0.420) (-0.611,-0.365) (-0.623,-0.363) (-0.639,-0.367) (-0.650,-0.354) (-0.655,-0.322) (-0.643,-0.299) (-0.637,-0.298) (-0.659,-0.317) (-0.706,-0.314) (-0.693,-0.300) (-0.691,-0.275) (-0.702,-0.254) (-0.698,-0.273) (-0.670,-0.247) (-0.653,-0.209) (-0.627,-0.229) (-0.656,-0.209) (-0.655,-0.148) (-0.647,-0.137) (-0.658,-0.129) (-0.633,-0.145) (-0.633,-0.169) (-0.621,-0.137) (-0.587,-0.110) (-0.590,-0.100) (-0.609,-0.125) (-0.620,-0.092) (-0.617,-0.081) (-0.641,-0.078) (-0.628,-0.107) (-0.642,-0.090) (-0.630,-0.105) (-0.620,-0.098) (-0.607,-0.098) (-0.651,-0.062) (-0.711,-0.052) (-0.716,-0.038) (-0.736,-0.044) (-0.733,-0.041) (-0.757,-0.015) (-0.751,-0.034) (-0.784,-0.060) (-0.818,-0.089) (-0.816,-0.101) (-0.811,-0.100) (-0.806,-0.083) (-0.834,-0.078) (-0.809,-0.042) (-0.802,-0.025) (-0.807,-0.010) (-0.781,-0.032) (-0.788,-0.041) (-0.797,-0.044) (-0.814,-0.048) (-0.797,-0.045) (-0.754,-0.040) (-0.777,-0.061) (-0.757,-0.057) (-0.725,-0.054) (-0.720,-0.042) (-0.674,-0.015) (-0.648,-0.002) (-0.638,0.021) (-0.692,0.011) (-0.701,0.009) (-0.728,0.019) (-0.731,0.039) (-0.726,0.028) (-0.753,0.040) (-0.786,0.058) (-0.773,0.105) (-0.781,0.083) (-0.764,0.114) (-0.763,0.129) (-0.722,0.128) (-0.742,0.145) (-0.804,0.163) (-0.779,0.119) (-0.826,0.108) (-0.794,0.107) (-0.804,0.110) (-0.797,0.077) (-0.820,0.058) (-0.816,0.031) (-0.827,0.032) (-0.866,0.032) (-0.884,0.070) (-0.894,0.046) (-0.907,0.057) (-0.897,0.047) (-0.887,0.054) (-0.875,0.078) (-0.915,0.065) (-0.930,0.104) (-0.887,0.132) (-0.842,0.140) (-0.803,0.110) (-0.807,0.079) (-0.779,0.077) (-0.786,0.069) (-0.782,0.073) (-0.761,0.077) (-0.764,0.044) (-0.757,0.038) (-0.783,0.022) (-0.817,-0.001) (-0.794,-0.024) (-0.778,-0.005) (-0.793,-0.021) (-0.833,-0.058) (-0.856,-0.063) (-0.882,-0.040) (-0.897,-0.019) (-0.913,-0.025) (-0.911,-0.021) (-0.902,-0.031) (-0.918,-0.042) (-0.910,-0.084) (-0.929,-0.093) (-0.917,-0.110) (-0.906,-0.122) (-0.914,-0.118) (-0.914,-0.103) (-0.911,-0.103) (-0.904,-0.107) (-0.914,-0.095) (-0.928,-0.115) (-0.907,-0.164) (-0.889,-0.136) (-0.890,-0.166) (-0.855,-0.151) (-0.825,-0.148) (-0.870,-0.134) (-0.838,-0.110) (-0.864,-0.090) (-0.899,-0.085) (-0.901,-0.082) (-0.871,-0.115) (-0.861,-0.088) (-0.887,-0.102) (-0.872,-0.107) (-0.888,-0.103) (-0.877,-0.144) (-0.833,-0.133) (-0.814,-0.130) (-0.813,-0.099) (-0.820,-0.076) (-0.808,-0.049) (-0.773,-0.045) (-0.781,-0.041) (-0.767,-0.047) (-0.761,-0.056) (-0.738,-0.038) (-0.725,-0.039) (-0.719,-0.016) (-0.719,0.003) (-0.712,0.002) (-0.680,-0.002) (-0.697,-0.029) (-0.731,-0.034) (-0.704,-0.036) (-0.683,-0.044) (-0.643,-0.040) (-0.629,-0.065) (-0.662,-0.077) (-0.614,-0.071) (-0.627,-0.056) (-0.686,-0.075) (-0.690,-0.035) (-0.692,-0.014) (-0.705,-0.015) (-0.728,-0.002) (-0.752,0.018) (-0.784,0.039) (-0.789,0.061) (-0.792,0.044) (-0.786,0.041) (-0.797,0.057) (-0.797,0.108) (-0.858,0.088) (-0.824,0.098) (-0.830,0.117) (-0.824,0.132) (-0.870,0.183) (-0.853,0.169) (-0.871,0.170) (-0.895,0.172) (-0.854,0.189) (-0.868,0.154) (-0.864,0.156) (-0.858,0.174) (-0.858,0.204) (-0.864,0.222) (-0.875,0.164) (-0.881,0.181) (-0.859,0.187) (-0.863,0.219) (-0.850,0.221) (-0.866,0.221) (-0.872,0.183) (-0.867,0.166) (-0.925,0.171) (-0.914,0.182) (-0.919,0.204) (-0.902,0.194) (-0.914,0.215) (-0.864,0.209) (-0.888,0.199) (-0.888,0.194) (-0.874,0.205) (-0.912,0.232) (-0.897,0.237) (-0.880,0.240) (-0.857,0.262) (-0.805,0.260) (-0.808,0.283) (-0.780,0.244) (-0.796,0.237) (-0.777,0.220) (-0.751,0.220) (-0.782,0.230) (-0.741,0.255) (-0.750,0.249) (-0.710,0.217) (-0.702,0.172) (-0.705,0.181) (-0.730,0.187) (-0.750,0.189) (-0.735,0.176) (-0.745,0.150) (-0.726,0.165) (-0.738,0.181) (-0.750,0.174) (-0.760,0.178) (-0.740,0.181) (-0.764,0.208) (-0.752,0.172) (-0.690,0.160) (-0.710,0.187) (-0.734,0.206) (-0.748,0.179) (-0.735,0.173) (-0.752,0.148) (-0.769,0.158) (-0.790,0.138) (-0.775,0.176) (-0.796,0.210) (-0.773,0.244) (-0.759,0.262) (-0.774,0.253) (-0.793,0.198) (-0.802,0.224) (-0.868,0.245) (-0.836,0.228) (-0.824,0.182) (-0.794,0.184) (-0.824,0.160) (-0.819,0.168) (-0.773,0.154) (-0.768,0.159) (-0.815,0.193) (-0.818,0.207) (-0.839,0.207) (-0.840,0.215) (-0.845,0.228) (-0.866,0.248) (-0.875,0.228) (-0.924,0.224) (-0.891,0.199) (-0.878,0.173) (-0.892,0.180) (-0.897,0.188) (-0.877,0.153) (-0.902,0.155) (-0.866,0.178) (-0.867,0.201) (-0.857,0.210) (-0.911,0.211) (-0.907,0.211) (-0.918,0.170) (-0.914,0.185) (-0.911,0.190) (-0.931,0.163) (-0.952,0.108) (-0.950,0.062) (-0.960,0.043) (-0.973,0.032) (-1.000,0.022) (-1.052,0.024) (-1.068,0.003) (-1.050,0.073) (-1.051,0.055) (-1.021,0.083) (-0.986,0.035) (-0.991,0.055) (-1.024,0.089) (-1.008,0.096) (-0.980,0.089) (-0.950,0.096) (-0.917,0.115) (-0.928,0.075) (-0.943,0.069) (-0.957,0.042) (-0.976,0.044) (-0.969,0.084) (-0.995,0.044) (-0.998,0.059) (-1.005,0.038) (-1.026,0.066) (-1.019,0.066) (-1.059,0.085) (-1.130,0.129) (-1.113,0.139) (-1.111,0.142) (-1.064,0.177) (-1.070,0.191) (-1.070,0.212) (-1.068,0.194) (-1.090,0.204) (-1.039,0.223) (-1.034,0.221) (-1.055,0.226) (-1.082,0.210) (-1.046,0.212) (-1.050,0.223) (-1.065,0.265) (-1.060,0.243) (-1.066,0.259) (-1.058,0.240) (-1.033,0.251) (-1.028,0.252) (-1.040,0.232) (-1.000,0.238) (-1.021,0.226) (-1.054,0.236) (-1.064,0.214) (-1.039,0.225) (-1.081,0.241) (-1.092,0.216) (-1.088,0.240) (-1.124,0.238) (-1.102,0.236) (-1.082,0.236) (-1.043,0.231) (-1.040,0.227) (-1.070,0.234) (-1.022,0.256) (-0.972,0.271) (-0.983,0.287) (-0.983,0.295) (-0.936,0.281) (-0.914,0.286) (-0.982,0.278) (-0.970,0.273) (-0.961,0.284) (-0.937,0.289) (-0.952,0.304) (-0.987,0.319) (-0.975,0.315) (-1.000,0.363) (-0.960,0.341) (-0.945,0.317) (-0.951,0.305) (-0.976,0.300) (-0.942,0.285) (-0.923,0.318) (-0.955,0.319) (-0.943,0.336) (-0.916,0.320) (-0.946,0.365) (-0.934,0.347) (-0.939,0.360) (-0.952,0.364) (-0.965,0.341) (-0.981,0.384) (-1.011,0.357) (-1.018,0.333) (-0.992,0.343) (-0.948,0.342) (-0.946,0.347) (-0.939,0.323) (-0.965,0.316) (-0.972,0.304) (-0.977,0.320) (-1.005,0.289) (-1.011,0.310) (-1.022,0.292) (-1.013,0.296) (-1.012,0.284) (-1.027,0.290) (-1.025,0.252) (-1.066,0.272) (-1.099,0.296) (-1.068,0.277) (-1.079,0.279) (-1.094,0.324) (-1.099,0.346) (-1.069,0.325) (-1.092,0.298) (-1.123,0.311) (-1.129,0.314) (-1.156,0.290) (-1.119,0.277) (-1.116,0.260) (-1.105,0.234) (-1.103,0.230) (-1.056,0.231) (-1.044,0.273) (-1.059,0.286) (-1.071,0.286) (-1.083,0.276) (-1.128,0.309) (-1.156,0.294) (-1.197,0.319) (-1.210,0.293) (-1.235,0.307) (-1.234,0.343) (-1.228,0.336) (-1.237,0.330) (-1.228,0.360) (-1.190,0.387) (-1.199,0.353) (-1.203,0.359) (-1.270,0.366) (-1.286,0.344) (-1.286,0.355) (-1.312,0.358) (-1.323,0.312) (-1.313,0.310) (-1.340,0.322) (-1.375,0.319) (-1.347,0.328) (-1.366,0.339) (-1.337,0.328) (-1.351,0.319) (-1.329,0.301) (-1.328,0.295) (-1.359,0.312) (-1.344,0.356) (-1.338,0.376) (-1.345,0.340) (-1.322,0.353) (-1.316,0.373) (-1.327,0.391) (-1.311,0.395) (-1.332,0.389) (-1.344,0.385) (-1.320,0.355) (-1.316,0.340) (-1.291,0.355) (-1.340,0.314) (-1.350,0.315) (-1.380,0.314) (-1.384,0.259) (-1.391,0.252) (-1.378,0.224) (-1.386,0.168) (-1.374,0.145) (-1.384,0.132) (-1.376,0.126) (-1.370,0.160) (-1.361,0.156) (-1.345,0.193) (-1.346,0.233) (-1.387,0.257) (-1.335,0.293) (-1.324,0.283) (-1.296,0.298) (-1.305,0.300) (-1.318,0.316) (-1.349,0.341) (-1.351,0.364) (-1.371,0.353) (-1.408,0.318) (-1.429,0.288) (-1.425,0.294) (-1.407,0.297) (-1.419,0.323) (-1.407,0.344) (-1.371,0.326) (-1.381,0.326) (-1.391,0.355) (-1.401,0.333) (-1.376,0.305) (-1.389,0.354) (-1.385,0.319) (-1.387,0.335) (-1.391,0.329) (-1.353,0.316) (-1.337,0.335) (-1.323,0.370) (-1.316,0.343) (-1.336,0.350) (-1.340,0.337) (-1.361,0.344) (-1.349,0.373) (-1.380,0.357) (-1.374,0.358) (-1.376,0.325) (-1.371,0.317) (-1.393,0.319) (-1.399,0.299) (-1.422,0.322) (-1.475,0.282) (-1.500,0.283) (-1.526,0.272) (-1.512,0.265) (-1.523,0.292) (-1.486,0.295) (-1.499,0.297) (-1.511,0.300) (-1.493,0.306) (-1.475,0.301) (-1.442,0.282) (-1.434,0.310) (-1.429,0.315) (-1.387,0.373) (-1.420,0.320) (-1.433,0.324) (-1.441,0.373) (-1.452,0.371) (-1.454,0.387) (-1.435,0.379) (-1.432,0.360) (-1.439,0.380) (-1.419,0.398) (-1.443,0.399) (-1.460,0.411) (-1.464,0.422) (-1.437,0.435) (-1.426,0.463) (-1.408,0.471) (-1.427,0.447) (-1.419,0.437) (-1.419,0.421) (-1.462,0.417) (-1.466,0.426) (-1.500,0.442) (-1.514,0.425) (-1.515,0.414) (-1.507,0.378) (-1.510,0.425) (-1.498,0.396) (-1.523,0.407) (-1.546,0.403) (-1.555,0.445) (-1.566,0.437) (-1.532,0.456) (-1.531,0.450) (-1.539,0.468) (-1.548,0.489) (-1.536,0.545) (-1.565,0.535) (-1.575,0.541) (-1.597,0.552) (-1.595,0.583) (-1.582,0.564) (-1.588,0.575) (-1.553,0.555) (-1.563,0.521) (-1.575,0.494) (-1.546,0.503) (-1.551,0.516) (-1.530,0.553) (-1.574,0.541) (-1.617,0.506) (-1.659,0.473) (-1.670,0.465) (-1.665,0.489) (-1.660,0.483) (-1.647,0.436) (-1.618,0.433) (-1.611,0.463) (-1.582,0.427) (-1.583,0.389) (-1.576,0.407) (-1.560,0.394) (-1.549,0.405) (-1.530,0.395) (-1.487,0.397) (-1.479,0.391) (-1.501,0.393) (-1.470,0.403) (-1.469,0.380) (-1.461,0.389) (-1.475,0.356) (-1.505,0.372) (-1.514,0.406) (-1.563,0.436) (-1.587,0.456) (-1.577,0.485) (-1.571,0.476) (-1.569,0.458) (-1.635,0.459) (-1.667,0.489) (-1.667,0.458) (-1.655,0.451) (-1.643,0.489) (-1.647,0.486) (-1.658,0.475) (-1.686,0.497) (-1.695,0.503) (-1.709,0.531) (-1.674,0.542) (-1.652,0.540) (-1.655,0.535) (-1.688,0.550) (-1.656,0.577) (-1.658,0.617) (-1.662,0.637) (-1.680,0.667) (-1.686,0.662) (-1.637,0.708) (-1.663,0.710) (-1.655,0.727) (-1.628,0.695) (-1.627,0.702) (-1.637,0.755) (-1.640,0.728) (-1.649,0.721) (-1.645,0.713) (-1.623,0.667) (-1.596,0.668) (-1.622,0.700) (-1.641,0.706) (-1.626,0.723) (-1.601,0.722) (-1.613,0.732) (-1.590,0.714) (-1.600,0.721) (-1.591,0.754) (-1.573,0.727) (-1.586,0.706) (-1.583,0.753) (-1.588,0.731) (-1.572,0.715) (-1.567,0.709) (-1.558,0.705) (-1.518,0.691) (-1.466,0.662) (-1.497,0.629) (-1.518,0.623) (-1.536,0.653) (-1.555,0.699) (-1.566,0.746) (-1.582,0.771) (-1.549,0.742) (-1.543,0.739) (-1.527,0.756) (-1.533,0.744) (-1.516,0.718) (-1.521,0.701) (-1.516,0.708) (-1.536,0.715) (-1.550,0.706) (-1.556,0.671) (-1.548,0.638) (-1.522,0.648) (-1.527,0.611) (-1.543,0.598) (-1.599,0.602) (-1.601,0.617) (-1.629,0.588) (-1.635,0.635) (-1.654,0.639) (-1.622,0.637) (-1.626,0.600) (-1.651,0.624) (-1.646,0.594) (-1.647,0.597) (-1.645,0.601) (-1.617,0.595) (-1.609,0.604) (-1.617,0.615) (-1.626,0.597) (-1.626,0.620) (-1.599,0.683) (-1.586,0.679) (-1.605,0.663) (-1.582,0.664) (-1.582,0.700) (-1.596,0.707) (-1.564,0.683) (-1.581,0.632) (-1.517,0.622) (-1.496,0.634) (-1.474,0.648) (-1.495,0.614) (-1.481,0.592) (-1.467,0.602) (-1.487,0.623) (-1.537,0.624) (-1.523,0.603) (-1.496,0.577) (-1.496,0.569) (-1.503,0.588) (-1.488,0.583) (-1.498,0.580) (-1.520,0.617) (-1.512,0.627) (-1.494,0.664) (-1.549,0.634) (-1.534,0.665) (-1.538,0.700) (-1.561,0.725) (-1.550,0.687) (-1.557,0.695) (-1.592,0.683) (-1.578,0.679) (-1.591,0.717) (-1.556,0.728) (-1.589,0.687) (-1.571,0.716) (-1.578,0.706) (-1.595,0.682) (-1.592,0.635) (-1.586,0.629) (-1.560,0.639) (-1.550,0.626) (-1.515,0.633) (-1.495,0.641) (-1.464,0.622) (-1.482,0.640) (-1.459,0.683) (-1.407,0.667) (-1.394,0.695) (-1.378,0.671) (-1.384,0.680) (-1.380,0.636) (-1.381,0.627) (-1.389,0.644) (-1.383,0.680) (-1.409,0.646) (-1.435,0.631) (-1.429,0.622) (-1.435,0.614) (-1.434,0.579) (-1.423,0.539) (-1.399,0.528) (-1.424,0.554) (-1.436,0.553) (-1.448,0.538) (-1.448,0.527) (-1.456,0.538) (-1.475,0.546) (-1.493,0.536) (-1.494,0.535) (-1.475,0.541) (-1.458,0.510) (-1.432,0.544) (-1.452,0.544) (-1.468,0.544) (-1.483,0.516) (-1.464,0.529) (-1.514,0.528) (-1.497,0.562) (-1.475,0.560) (-1.466,0.553) (-1.472,0.566) (-1.469,0.583) (-1.458,0.620) (-1.452,0.608) (-1.450,0.598) (-1.456,0.619) (-1.462,0.626) (-1.429,0.653) (-1.432,0.637) (-1.424,0.638) (-1.448,0.635) (-1.447,0.641) (-1.426,0.659) (-1.421,0.657) (-1.435,0.687) (-1.449,0.680) (-1.460,0.653) (-1.468,0.633) (-1.487,0.601) (-1.471,0.580) (-1.477,0.583) (-1.524,0.604) (-1.602,0.587) (-1.608,0.527) (-1.609,0.555) (-1.622,0.571) (-1.653,0.627) (-1.628,0.632) (-1.657,0.643) (-1.626,0.618) (-1.624,0.596) (-1.588,0.590) (-1.580,0.603) (-1.580,0.598) (-1.545,0.605) (-1.553,0.613) (-1.565,0.615) (-1.541,0.609) (-1.545,0.615) (-1.536,0.591) (-1.516,0.595) (-1.527,0.581) (-1.511,0.555) (-1.539,0.565) (-1.549,0.585) (-1.527,0.585) (-1.517,0.597) (-1.519,0.605) (-1.517,0.618) (-1.544,0.631) (-1.526,0.647) (-1.509,0.611) (-1.505,0.568) (-1.496,0.603) (-1.500,0.614) (-1.479,0.622) (-1.461,0.647) (-1.440,0.675) (-1.427,0.673) (-1.407,0.693) (-1.413,0.716) (-1.391,0.712) (-1.366,0.733) (-1.386,0.736) (-1.358,0.742) (-1.391,0.746) (-1.427,0.727) (-1.422,0.688) (-1.443,0.695) (-1.436,0.702) (-1.410,0.722) (-1.434,0.739) (-1.426,0.715) (-1.461,0.691) (-1.467,0.694) (-1.432,0.662) (-1.447,0.663) (-1.432,0.668) (-1.480,0.643) (-1.412,0.627) (-1.382,0.629) (-1.394,0.600) (-1.429,0.613) (-1.433,0.636) (-1.456,0.640) (-1.426,0.658) (-1.407,0.646) (-1.389,0.662) (-1.391,0.676) (-1.359,0.683) (-1.388,0.668) (-1.362,0.679) (-1.401,0.664) (-1.373,0.661) (-1.340,0.687) (-1.353,0.681) (-1.345,0.637) (-1.314,0.649) (-1.340,0.599) (-1.334,0.573) (-1.358,0.588) (-1.380,0.589) (-1.375,0.585) (-1.373,0.567) (-1.399,0.570) (-1.373,0.576) (-1.365,0.607) (-1.380,0.618) (-1.386,0.630) (-1.362,0.649) (-1.356,0.663) (-1.345,0.675) (-1.332,0.688) (-1.338,0.673) (-1.298,0.665) (-1.282,0.697) (-1.274,0.728) (-1.256,0.713) (-1.254,0.724) (-1.249,0.699) (-1.268,0.704) (-1.315,0.708) (-1.307,0.735) (-1.307,0.723) (-1.296,0.738) (-1.297,0.736) (-1.285,0.756) (-1.278,0.797) (-1.266,0.812) (-1.256,0.838) (-1.256,0.854) (-1.229,0.811) (-1.205,0.830) (-1.214,0.842) (-1.205,0.817) (-1.168,0.851) (-1.190,0.885) (-1.162,0.872) (-1.159,0.867) (-1.156,0.856) (-1.138,0.843) (-1.095,0.885) (-1.108,0.922) (-1.084,0.903) (-1.136,0.911) (-1.138,0.892) (-1.153,0.932) (-1.157,0.937) (-1.161,0.937) (-1.160,0.902) (-1.124,0.920) (-1.131,0.911) (-1.147,0.889) (-1.127,0.884) (-1.120,0.900) (-1.124,0.864) (-1.144,0.866) (-1.166,0.882) (-1.188,0.897) (-1.176,0.920) (-1.209,0.925) (-1.193,0.911) (-1.215,0.899) (-1.206,0.920) (-1.190,0.878) (-1.189,0.862) (-1.204,0.844) (-1.237,0.809) (-1.270,0.862) (-1.244,0.885) (-1.224,0.840) (-1.234,0.839) (-1.256,0.865) (-1.225,0.857) (-1.205,0.864) (-1.198,0.878) (-1.187,0.921) (-1.183,0.920) (-1.173,0.906) (-1.164,0.905) (-1.136,0.925) (-1.159,0.945) (-1.196,0.942) (-1.232,0.918) (-1.209,0.933) (-1.173,0.962) (-1.185,0.933) (-1.194,0.902) (-1.183,0.892) (-1.176,0.909) (-1.171,0.927) (-1.140,0.942) (-1.119,0.936) (-1.130,0.949) (-1.149,0.952) (-1.181,0.940) (-1.187,0.949) (-1.138,0.980) (-1.160,0.975) (-1.205,0.978) (-1.197,0.961) (-1.189,0.983) (-1.243,0.949) (-1.245,0.928) (-1.287,0.918) (-1.347,0.921) (-1.343,0.901) (-1.345,0.882) (-1.356,0.879) (-1.339,0.902) (-1.320,0.880) (-1.312,0.864) (-1.305,0.880) (-1.280,0.838) (-1.269,0.834) (-1.260,0.824) (-1.281,0.799) (-1.285,0.793) (-1.272,0.779) (-1.258,0.820) (-1.235,0.837) (-1.208,0.883) (-1.192,0.902) (-1.232,0.920) (-1.222,0.902) (-1.228,0.878) (-1.218,0.852) (-1.231,0.836) (-1.256,0.837) (-1.252,0.840) (-1.268,0.806) (-1.297,0.842) (-1.271,0.843) (-1.275,0.839) (-1.294,0.869) (-1.304,0.864) (-1.293,0.909) (-1.277,0.900) (-1.262,0.897) (-1.263,0.875) (-1.296,0.865) (-1.306,0.821) (-1.364,0.829) (-1.390,0.827) (-1.390,0.785) (-1.396,0.774) (-1.432,0.793) (-1.426,0.791) (-1.430,0.797) (-1.427,0.801) (-1.431,0.794) (-1.424,0.837) (-1.381,0.845) (-1.398,0.857) (-1.410,0.865) (-1.421,0.891) (-1.431,0.915) (-1.441,0.919) (-1.429,0.894) (-1.439,0.882) (-1.445,0.852) (-1.436,0.833) (-1.434,0.795) (-1.447,0.816) (-1.446,0.823) (-1.455,0.827) (-1.431,0.829) (-1.443,0.830) (-1.476,0.865) (-1.457,0.823) (-1.467,0.791) (-1.466,0.815) (-1.483,0.807) (-1.489,0.789) (-1.483,0.766) (-1.476,0.769) (-1.487,0.793) (-1.502,0.792) (-1.479,0.781) (-1.440,0.753) (-1.439,0.769) (-1.462,0.778) (-1.453,0.777) (-1.523,0.765) (-1.566,0.723) (-1.576,0.707) (-1.553,0.719) (-1.584,0.697) (-1.605,0.682) (-1.577,0.677) (-1.518,0.665) (-1.486,0.665) (-1.473,0.671) (-1.465,0.657) (-1.429,0.640) (-1.445,0.661) (-1.450,0.701) (-1.424,0.728) (-1.391,0.705) (-1.366,0.751) (-1.313,0.735) (-1.320,0.712) (-1.347,0.695) (-1.351,0.692) (-1.404,0.702) (-1.398,0.679) (-1.419,0.686) (-1.403,0.679) (-1.432,0.685) (-1.407,0.680) (-1.410,0.641) (-1.425,0.643) (-1.435,0.649) (-1.451,0.616) (-1.438,0.603) (-1.422,0.625) (-1.412,0.654) (-1.390,0.667) (-1.397,0.676) (-1.385,0.654) (-1.419,0.672) (-1.441,0.651) (-1.454,0.653) (-1.435,0.620) (-1.420,0.610) (-1.417,0.624) (-1.465,0.635) (-1.455,0.633) (-1.452,0.624) (-1.451,0.650) (-1.458,0.653) (-1.441,0.658) (-1.461,0.659) (-1.468,0.715) (-1.478,0.774) (-1.461,0.757) (-1.461,0.759) (-1.484,0.783) (-1.488,0.751) (-1.521,0.730) (-1.573,0.699) (-1.586,0.720) (-1.624,0.739) (-1.649,0.765) (-1.615,0.784) (-1.597,0.781) (-1.553,0.762) (-1.510,0.744) (-1.534,0.760) (-1.520,0.723) (-1.513,0.787) (-1.444,0.713) (-1.437,0.717) (-1.429,0.730) (-1.409,0.748) (-1.408,0.741) (-1.396,0.756) (-1.430,0.771) (-1.451,0.698) (-1.456,0.683) (-1.482,0.673) (-1.455,0.682) (-1.433,0.693) (-1.407,0.675) (-1.356,0.664) (-1.343,0.662) (-1.348,0.663) (-1.352,0.667) (-1.343,0.667) (-1.407,0.666) (-1.433,0.615) (-1.441,0.629) (-1.407,0.659) (-1.426,0.672) (-1.432,0.695) (-1.419,0.644) (-1.378,0.668) (-1.344,0.691) (-1.357,0.688) (-1.406,0.717) (-1.387,0.696) (-1.358,0.668) (-1.366,0.647) (-1.352,0.662) (-1.391,0.643) (-1.384,0.661) (-1.362,0.666) (-1.338,0.633) (-1.333,0.614) (-1.274,0.589) (-1.296,0.615) (-1.278,0.618) (-1.273,0.618) (-1.292,0.586) (-1.274,0.609) (-1.252,0.599) (-1.266,0.621) (-1.267,0.616) (-1.235,0.604) (-1.244,0.627) (-1.223,0.650) (-1.199,0.632) (-1.215,0.617) (-1.221,0.627) (-1.218,0.620) (-1.218,0.631) (-1.268,0.652) (-1.281,0.635) (-1.285,0.664) (-1.291,0.655) (-1.279,0.620) (-1.266,0.593) (-1.240,0.594) (-1.258,0.583) (-1.251,0.582) (-1.252,0.613) (-1.246,0.641) (-1.236,0.644) (-1.227,0.651) (-1.225,0.670) (-1.196,0.690) (-1.181,0.727) (-1.154,0.743) (-1.152,0.769) (-1.142,0.786) (-1.129,0.799) (-1.115,0.793) (-1.125,0.811) (-1.095,0.829) (-1.106,0.816) (-1.138,0.846) (-1.160,0.839) (-1.155,0.845) (-1.191,0.858) (-1.214,0.875) (-1.197,0.874) (-1.159,0.899) (-1.161,0.837) (-1.174,0.831) (-1.196,0.845) (-1.134,0.869)};
  \fill[SARed] (axis cs:0,0) circle (2pt);
\end{axis}
\end{tikzpicture}
\caption{平面布朗运动 $(B^1_t,B^2_t)_{t\le1}$ 的轨道, 红点为起点. 轨道连续, 但极其曲折: 可以证明它的像集的 Hausdorff 维数为 $2$}
\label{fig:bm-2d}
\end{figure}

% =====================================================================
\section{不变性与轨道性质}
\label{sec:4-3}

% ---------------------------------------------------------------------
\subsection{布朗运动的不变性}

布朗运动在许多变换下保持不变. 这些不变性是计算的利器: 一个关于布朗运动的问题, 往往可以通过适当的变换化归为一个已经解决的问题.

\begin{theorem}[不变性]\label{thm:bm-invariance}
设 $B$ 是标准布朗运动. 则下列过程都是标准布朗运动:
\begin{enumerate}
\item \textbf{对称性}: $X_t=-B_t$;
\item \textbf{尺度不变性}: $X_t=\frac{1}{\sqrt{c}}B_{ct}$, 其中 $c>0$ 是常数;
\item \textbf{平移不变性}: $X_t=B_{s+t}-B_s$, 其中 $s\ge0$ 是常数; 并且 $X$ 与 $\mathcal{F}^B_s$ 独立;
\item \textbf{时间反演}: $X_0=0$, $X_t=tB_{1/t}$ ($t>0$).
\end{enumerate}
\textbf{时间逆转}同样保持布朗运动性质. 对固定的 $T>0$, $\{B_T-B_{T-t}\}_{t\in[0,T]}$ 是 $[0,T]$ 上的标准布朗运动.
\end{theorem}

\begin{proof}
所有情形都用高斯刻画 (定理~\ref{thm:bm-gaussian}): 只需验证 $X$ 的轨道连续、$X_0=0$, 且 $X$ 是均值为零、协方差为 $s\wedge t$ 的高斯过程. 在每种情形中, $X$ 在有限个时刻上的取值都是 $B$ 在有限个时刻上取值的线性变换, 因而是均值为零的正态向量 (命题~\ref{prop:gaussian-operations}). 下面对 $u\le v$ 计算协方差并验证连续性.

(1) $\Cov(-B_u,-B_v)=u$. 轨道连续性显然.

(2) $\Cov\big(c^{-1/2}B_{cu},c^{-1/2}B_{cv}\big)=c^{-1}(cu\wedge cv)=u$. 轨道连续性显然.

(3) $\Cov(B_{s+u}-B_s,B_{s+v}-B_s)=(s+u)-s-s+s=u$. 轨道连续性显然. 为验证独立性, 注意到 $X$ 在有限个时刻 $u_1<\cdots<u_n$ 上的取值是增量 $B_{s+u_1}-B_s,\dots,B_{s+u_n}-B_{s+u_{n-1}}$ 的函数, 由 4.2 节中的论证, 这些增量与 $\mathcal{F}^B_s$ 独立. 事件 $\{(X_{u_1},\dots,X_{u_n})\in A\}$ 构成生成 $\sigma(X)$ 的 $\pi$-系, 由定理~\ref{thm:indep-pi}, $\sigma(X)$ 与 $\mathcal{F}^B_s$ 独立.

(4) 对 $0<u\le v$, $\Cov(uB_{1/u},vB_{1/v})=uv\big(\frac{1}{u}\wedge\frac{1}{v}\big)=uv\cdot\frac{1}{v}=u$; 涉及 $X_0=0$ 的协方差为零, 也等于 $0\wedge v$. 所以 $X$ 是均值为零、协方差为 $s\wedge t$ 的高斯过程. $X$ 在 $(0,\infty)$ 上的轨道显然连续, 只需证明 $\lim_{t\downarrow0}X_t=0$ $\as$ 该事件可以只用可数个时刻表达. 由 $X$ 在 $(0,\infty)$ 上连续,
\[
\Big\{\lim_{t\downarrow0}X_t=0\Big\}=\bigcap_{m\ge1}\bigcup_{n\ge1}\bigcap_{q\in\Q\cap(0,1/n]}\Big\{\abs{X_q}\le\frac{1}{m}\Big\},
\]
对 $B$ 有同样的表达式. 记 $\Q_+=\Q\cap(0,\infty)$. 由于 $\{X_q\}_{q\in\Q_+}$ 与 $\{B_q\}_{q\in\Q_+}$ 有相同的有限维分布, 由推论~\ref{cor:uniqueness} 它们作为 $\R^{\Q_+}$ 值随机元有相同的分布, 而上式右端是 $\R^{\Q_+}$ 中的可测集的原像. 所以
\[
\P\Big(\lim_{t\downarrow0}X_t=0\Big)=\P\Big(\lim_{t\downarrow0}B_t=0\Big)=1.
\]
在这个零概率集的补集之外重新定义 $X\equiv0$, 即得连续轨道.

时间逆转: 对 $0\le u\le v\le T$, $\Cov(B_T-B_{T-u},B_T-B_{T-v})=T-(T-v)-(T-u)+(T-v)=u$, 轨道连续且起点为零.
\end{proof}

\begin{figure}[htbp]
\centering
\begin{tikzpicture}
\begin{axis}[SAaxis, axis x line=bottom, axis y line=left, width=7.6cm, height=5.4cm, xmin=0, xmax=1.03, ymin=-0.5, ymax=1.2,
  xtick={0,0.25,0.5,0.75,1}, ytick={0,0.5,1}, xlabel={$t$}, title={\small $B_t$, $t\in[0,1]$}]
  \fill[SAYellow!30] (axis cs:0,-0.5) rectangle (axis cs:0.25,1.2);
  \addplot[SADark!20] coordinates {(0,0) (1,0)};
  \addplot[SADark,thin] coordinates {(0.000,0.000) (0.002,0.002) (0.004,0.034) (0.006,-0.039) (0.008,-0.005) (0.010,-0.010) (0.012,-0.050) (0.014,-0.043) (0.016,-0.060) (0.018,-0.082) (0.020,-0.140) (0.021,-0.082) (0.023,-0.102) (0.025,-0.056) (0.027,-0.039) (0.029,-0.064) (0.031,-0.016) (0.033,0.052) (0.035,0.043) (0.037,0.039) (0.039,0.017) (0.041,0.033) (0.043,-0.006) (0.045,0.054) (0.047,0.114) (0.049,0.035) (0.051,0.094) (0.053,0.097) (0.055,0.178) (0.057,0.163) (0.059,0.177) (0.061,0.197) (0.062,0.155) (0.064,0.193) (0.066,0.183) (0.068,0.102) (0.070,0.024) (0.072,0.013) (0.074,0.041) (0.076,0.113) (0.078,0.134) (0.080,0.161) (0.082,0.188) (0.084,0.184) (0.086,0.180) (0.088,0.208) (0.090,0.228) (0.092,0.214) (0.094,0.222) (0.096,0.124) (0.098,0.112) (0.100,0.195) (0.102,0.277) (0.104,0.358) (0.105,0.401) (0.107,0.363) (0.109,0.305) (0.111,0.340) (0.113,0.215) (0.115,0.254) (0.117,0.182) (0.119,0.166) (0.121,0.066) (0.123,0.057) (0.125,0.004) (0.127,0.050) (0.129,0.094) (0.131,0.120) (0.133,0.164) (0.135,0.223) (0.137,0.258) (0.139,0.333) (0.141,0.338) (0.143,0.275) (0.145,0.249) (0.146,0.238) (0.148,0.296) (0.150,0.314) (0.152,0.247) (0.154,0.247) (0.156,0.273) (0.158,0.300) (0.160,0.282) (0.162,0.279) (0.164,0.282) (0.166,0.262) (0.168,0.303) (0.170,0.295) (0.172,0.242) (0.174,0.220) (0.176,0.286) (0.178,0.313) (0.180,0.288) (0.182,0.311) (0.184,0.338) (0.186,0.358) (0.188,0.364) (0.189,0.346) (0.191,0.417) (0.193,0.474) (0.195,0.452) (0.197,0.423) (0.199,0.504) (0.201,0.518) (0.203,0.461) (0.205,0.438) (0.207,0.444) (0.209,0.403) (0.211,0.457) (0.213,0.456) (0.215,0.494) (0.217,0.488) (0.219,0.517) (0.221,0.549) (0.223,0.568) (0.225,0.562) (0.227,0.562) (0.229,0.533) (0.230,0.566) (0.232,0.542) (0.234,0.522) (0.236,0.512) (0.238,0.465) (0.240,0.447) (0.242,0.483) (0.244,0.423) (0.246,0.445) (0.248,0.404) (0.250,0.358) (0.252,0.421) (0.254,0.420) (0.256,0.423) (0.258,0.432) (0.260,0.326) (0.262,0.278) (0.264,0.325) (0.266,0.361) (0.268,0.397) (0.270,0.392) (0.271,0.482) (0.273,0.512) (0.275,0.536) (0.277,0.540) (0.279,0.500) (0.281,0.493) (0.283,0.451) (0.285,0.394) (0.287,0.300) (0.289,0.429) (0.291,0.375) (0.293,0.431) (0.295,0.458) (0.297,0.398) (0.299,0.440) (0.301,0.454) (0.303,0.442) (0.305,0.405) (0.307,0.479) (0.309,0.490) (0.311,0.467) (0.312,0.504) (0.314,0.496) (0.316,0.491) (0.318,0.484) (0.320,0.513) (0.322,0.597) (0.324,0.634) (0.326,0.609) (0.328,0.537) (0.330,0.499) (0.332,0.479) (0.334,0.425) (0.336,0.458) (0.338,0.424) (0.340,0.501) (0.342,0.523) (0.344,0.599) (0.346,0.607) (0.348,0.633) (0.350,0.615) (0.352,0.656) (0.354,0.677) (0.355,0.764) (0.357,0.791) (0.359,0.760) (0.361,0.720) (0.363,0.702) (0.365,0.662) (0.367,0.595) (0.369,0.570) (0.371,0.556) (0.373,0.532) (0.375,0.430) (0.377,0.441) (0.379,0.478) (0.381,0.389) (0.383,0.350) (0.385,0.272) (0.387,0.265) (0.389,0.287) (0.391,0.234) (0.393,0.363) (0.395,0.376) (0.396,0.271) (0.398,0.329) (0.400,0.382) (0.402,0.387) (0.404,0.439) (0.406,0.417) (0.408,0.347) (0.410,0.370) (0.412,0.285) (0.414,0.289) (0.416,0.261) (0.418,0.247) (0.420,0.277) (0.422,0.325) (0.424,0.351) (0.426,0.337) (0.428,0.341) (0.430,0.253) (0.432,0.252) (0.434,0.302) (0.436,0.227) (0.438,0.148) (0.439,0.176) (0.441,0.246) (0.443,0.312) (0.445,0.288) (0.447,0.340) (0.449,0.438) (0.451,0.465) (0.453,0.524) (0.455,0.505) (0.457,0.512) (0.459,0.518) (0.461,0.488) (0.463,0.448) (0.465,0.489) (0.467,0.498) (0.469,0.528) (0.471,0.557) (0.473,0.612) (0.475,0.587) (0.477,0.585) (0.479,0.559) (0.480,0.528) (0.482,0.487) (0.484,0.538) (0.486,0.510) (0.488,0.525) (0.490,0.431) (0.492,0.464) (0.494,0.493) (0.496,0.533) (0.498,0.495) (0.500,0.537) (0.502,0.566) (0.504,0.579) (0.506,0.638) (0.508,0.542) (0.510,0.507) (0.512,0.400) (0.514,0.370) (0.516,0.376) (0.518,0.372) (0.520,0.348) (0.521,0.243) (0.523,0.242) (0.525,0.214) (0.527,0.216) (0.529,0.217) (0.531,0.190) (0.533,0.202) (0.535,0.183) (0.537,0.104) (0.539,0.048) (0.541,0.082) (0.543,0.041) (0.545,0.066) (0.547,-0.017) (0.549,-0.027) (0.551,-0.044) (0.553,0.015) (0.555,0.051) (0.557,0.030) (0.559,0.059) (0.561,0.035) (0.562,0.035) (0.564,0.067) (0.566,0.108) (0.568,0.166) (0.570,0.183) (0.572,0.199) (0.574,0.202) (0.576,0.184) (0.578,0.165) (0.580,0.257) (0.582,0.206) (0.584,0.172) (0.586,0.193) (0.588,0.179) (0.590,0.171) (0.592,0.194) (0.594,0.231) (0.596,0.200) (0.598,0.250) (0.600,0.313) (0.602,0.348) (0.604,0.297) (0.605,0.345) (0.607,0.370) (0.609,0.348) (0.611,0.347) (0.613,0.308) (0.615,0.280) (0.617,0.245) (0.619,0.248) (0.621,0.296) (0.623,0.290) (0.625,0.269) (0.627,0.274) (0.629,0.341) (0.631,0.358) (0.633,0.299) (0.635,0.355) (0.637,0.374) (0.639,0.395) (0.641,0.450) (0.643,0.386) (0.645,0.374) (0.646,0.434) (0.648,0.477) (0.650,0.412) (0.652,0.419) (0.654,0.472) (0.656,0.469) (0.658,0.469) (0.660,0.441) (0.662,0.378) (0.664,0.387) (0.666,0.345) (0.668,0.367) (0.670,0.379) (0.672,0.381) (0.674,0.321) (0.676,0.368) (0.678,0.294) (0.680,0.225) (0.682,0.266) (0.684,0.273) (0.686,0.270) (0.688,0.254) (0.689,0.185) (0.691,0.109) (0.693,0.025) (0.695,-0.002) (0.697,0.027) (0.699,-0.002) (0.701,-0.002) (0.703,0.083) (0.705,0.114) (0.707,0.161) (0.709,0.186) (0.711,0.121) (0.713,0.158) (0.715,0.292) (0.717,0.269) (0.719,0.285) (0.721,0.377) (0.723,0.460) (0.725,0.520) (0.727,0.551) (0.729,0.439) (0.730,0.418) (0.732,0.388) (0.734,0.391) (0.736,0.393) (0.738,0.328) (0.740,0.238) (0.742,0.253) (0.744,0.244) (0.746,0.253) (0.748,0.285) (0.750,0.237) (0.752,0.172) (0.754,0.210) (0.756,0.155) (0.758,0.151) (0.760,0.119) (0.762,0.146) (0.764,0.141) (0.766,0.086) (0.768,-0.008) (0.770,0.033) (0.771,0.052) (0.773,0.079) (0.775,0.142) (0.777,0.137) (0.779,0.091) (0.781,0.091) (0.783,0.100) (0.785,0.071) (0.787,0.053) (0.789,0.030) (0.791,0.035) (0.793,0.098) (0.795,0.069) (0.797,0.062) (0.799,0.054) (0.801,0.079) (0.803,0.039) (0.805,0.009) (0.807,0.045) (0.809,0.041) (0.811,0.100) (0.812,0.152) (0.814,0.143) (0.816,0.199) (0.818,0.250) (0.820,0.239) (0.822,0.228) (0.824,0.209) (0.826,0.159) (0.828,0.181) (0.830,0.241) (0.832,0.203) (0.834,0.250) (0.836,0.223) (0.838,0.201) (0.840,0.269) (0.842,0.192) (0.844,0.149) (0.846,0.173) (0.848,0.149) (0.850,0.088) (0.852,0.184) (0.854,0.184) (0.855,0.240) (0.857,0.261) (0.859,0.270) (0.861,0.230) (0.863,0.230) (0.865,0.255) (0.867,0.256) (0.869,0.229) (0.871,0.166) (0.873,0.173) (0.875,0.234) (0.877,0.263) (0.879,0.257) (0.881,0.202) (0.883,0.170) (0.885,0.205) (0.887,0.134) (0.889,0.106) (0.891,0.071) (0.893,0.054) (0.895,0.085) (0.896,-0.002) (0.898,-0.030) (0.900,-0.079) (0.902,-0.051) (0.904,-0.014) (0.906,0.043) (0.908,0.083) (0.910,0.039) (0.912,0.075) (0.914,0.038) (0.916,0.099) (0.918,0.082) (0.920,0.117) (0.922,0.178) (0.924,0.177) (0.926,0.148) (0.928,0.170) (0.930,0.165) (0.932,0.153) (0.934,0.120) (0.936,0.153) (0.938,0.097) (0.939,0.113) (0.941,0.150) (0.943,0.152) (0.945,0.233) (0.947,0.288) (0.949,0.387) (0.951,0.296) (0.953,0.265) (0.955,0.257) (0.957,0.224) (0.959,0.164) (0.961,0.062) (0.963,-0.003) (0.965,0.030) (0.967,0.089) (0.969,0.089) (0.971,0.063) (0.973,0.124) (0.975,-0.010) (0.977,-0.013) (0.979,-0.037) (0.980,-0.009) (0.982,-0.034) (0.984,0.042) (0.986,0.029) (0.988,0.001) (0.990,-0.028) (0.992,-0.026) (0.994,0.105) (0.996,0.151) (0.998,0.129) (1.000,0.088)};
\end{axis}
\end{tikzpicture}\hfill
\begin{tikzpicture}
\begin{axis}[SAaxis, axis x line=bottom, axis y line=left, width=7.6cm, height=5.4cm, xmin=0, xmax=1.03, ymin=-0.5, ymax=1.2,
  xtick={0,0.25,0.5,0.75,1}, ytick={0,0.5,1}, xlabel={$t$}, title={\small $2B_{t/4}$, $t\in[0,1]$}]
  \addplot[SADark!20] coordinates {(0,0) (1,0)};
  \addplot[SABlue,thin] coordinates {(0.000,0.000) (0.002,0.008) (0.004,-0.009) (0.006,-0.045) (0.008,0.003) (0.010,0.044) (0.012,0.062) (0.014,0.036) (0.016,0.068) (0.018,-0.040) (0.020,-0.082) (0.021,-0.074) (0.023,-0.078) (0.025,-0.056) (0.027,-0.023) (0.029,0.032) (0.031,-0.010) (0.033,-0.025) (0.035,0.003) (0.037,-0.041) (0.039,-0.019) (0.041,-0.037) (0.043,-0.070) (0.045,-0.027) (0.047,-0.100) (0.049,-0.129) (0.051,-0.116) (0.053,-0.136) (0.055,-0.085) (0.057,-0.029) (0.059,-0.066) (0.061,-0.066) (0.062,-0.119) (0.064,-0.166) (0.066,-0.189) (0.068,-0.227) (0.070,-0.164) (0.072,-0.302) (0.074,-0.314) (0.076,-0.319) (0.078,-0.279) (0.080,-0.258) (0.082,-0.189) (0.084,-0.155) (0.086,-0.164) (0.088,-0.118) (0.090,-0.200) (0.092,-0.230) (0.094,-0.204) (0.096,-0.194) (0.098,-0.182) (0.100,-0.151) (0.102,-0.112) (0.104,-0.145) (0.105,-0.112) (0.107,-0.110) (0.109,-0.078) (0.111,-0.125) (0.113,-0.139) (0.115,-0.128) (0.117,-0.128) (0.119,-0.133) (0.121,-0.071) (0.123,-0.055) (0.125,-0.033) (0.127,-0.029) (0.129,-0.011) (0.131,0.030) (0.133,0.103) (0.135,0.078) (0.137,0.107) (0.139,0.108) (0.141,0.085) (0.143,0.058) (0.145,0.078) (0.146,0.085) (0.148,0.077) (0.150,0.164) (0.152,0.144) (0.154,0.135) (0.156,0.034) (0.158,0.029) (0.160,0.057) (0.162,0.130) (0.164,0.066) (0.166,0.075) (0.168,0.025) (0.170,0.014) (0.172,-0.011) (0.174,0.032) (0.176,0.033) (0.178,0.062) (0.180,0.107) (0.182,0.131) (0.184,0.181) (0.186,0.201) (0.188,0.227) (0.189,0.254) (0.191,0.159) (0.193,0.149) (0.195,0.071) (0.197,0.133) (0.199,0.150) (0.201,0.115) (0.203,0.188) (0.205,0.209) (0.207,0.234) (0.209,0.217) (0.211,0.195) (0.213,0.220) (0.215,0.235) (0.217,0.317) (0.219,0.357) (0.221,0.365) (0.223,0.373) (0.225,0.409) (0.227,0.325) (0.229,0.293) (0.230,0.284) (0.232,0.309) (0.234,0.353) (0.236,0.384) (0.238,0.481) (0.240,0.487) (0.242,0.395) (0.244,0.367) (0.246,0.328) (0.248,0.321) (0.250,0.310) (0.252,0.367) (0.254,0.415) (0.256,0.370) (0.258,0.386) (0.260,0.316) (0.262,0.372) (0.264,0.397) (0.266,0.366) (0.268,0.297) (0.270,0.308) (0.271,0.256) (0.273,0.204) (0.275,0.179) (0.277,0.115) (0.279,0.035) (0.281,0.049) (0.283,-0.048) (0.285,-0.054) (0.287,-0.006) (0.289,0.025) (0.291,0.087) (0.293,0.036) (0.295,0.105) (0.297,0.081) (0.299,0.087) (0.301,0.177) (0.303,0.170) (0.305,0.225) (0.307,0.271) (0.309,0.240) (0.311,0.275) (0.312,0.268) (0.314,0.225) (0.316,0.256) (0.318,0.203) (0.320,0.322) (0.322,0.261) (0.324,0.295) (0.326,0.361) (0.328,0.375) (0.330,0.347) (0.332,0.335) (0.334,0.433) (0.336,0.368) (0.338,0.383) (0.340,0.340) (0.342,0.439) (0.344,0.360) (0.346,0.444) (0.348,0.526) (0.350,0.507) (0.352,0.416) (0.354,0.430) (0.355,0.404) (0.357,0.430) (0.359,0.456) (0.361,0.471) (0.363,0.447) (0.365,0.427) (0.367,0.429) (0.369,0.434) (0.371,0.390) (0.373,0.437) (0.375,0.444) (0.377,0.428) (0.379,0.371) (0.381,0.365) (0.383,0.249) (0.385,0.265) (0.387,0.211) (0.389,0.203) (0.391,0.224) (0.393,0.293) (0.395,0.311) (0.396,0.393) (0.398,0.391) (0.400,0.409) (0.402,0.433) (0.404,0.480) (0.406,0.553) (0.408,0.585) (0.410,0.669) (0.412,0.686) (0.414,0.717) (0.416,0.686) (0.418,0.732) (0.420,0.758) (0.422,0.801) (0.424,0.809) (0.426,0.823) (0.428,0.771) (0.430,0.725) (0.432,0.649) (0.434,0.645) (0.436,0.661) (0.438,0.610) (0.439,0.683) (0.441,0.689) (0.443,0.627) (0.445,0.680) (0.447,0.619) (0.449,0.584) (0.451,0.494) (0.453,0.430) (0.455,0.439) (0.457,0.481) (0.459,0.564) (0.461,0.507) (0.463,0.441) (0.465,0.438) (0.467,0.404) (0.469,0.364) (0.471,0.349) (0.473,0.323) (0.475,0.356) (0.477,0.332) (0.479,0.242) (0.480,0.212) (0.482,0.171) (0.484,0.133) (0.486,0.115) (0.488,0.149) (0.490,0.121) (0.492,0.114) (0.494,0.147) (0.496,0.082) (0.498,0.047) (0.500,0.009) (0.502,-0.004) (0.504,-0.000) (0.506,0.062) (0.508,0.101) (0.510,0.061) (0.512,0.046) (0.514,0.137) (0.516,0.188) (0.518,0.241) (0.520,0.252) (0.521,0.234) (0.523,0.239) (0.525,0.314) (0.527,0.307) (0.529,0.344) (0.531,0.327) (0.533,0.368) (0.535,0.398) (0.537,0.451) (0.539,0.446) (0.541,0.468) (0.543,0.474) (0.545,0.475) (0.547,0.516) (0.549,0.520) (0.551,0.529) (0.553,0.611) (0.555,0.667) (0.557,0.691) (0.559,0.670) (0.561,0.671) (0.562,0.675) (0.564,0.578) (0.566,0.600) (0.568,0.589) (0.570,0.551) (0.572,0.507) (0.574,0.514) (0.576,0.521) (0.578,0.498) (0.580,0.532) (0.582,0.552) (0.584,0.516) (0.586,0.475) (0.588,0.501) (0.590,0.584) (0.592,0.562) (0.594,0.591) (0.596,0.580) (0.598,0.598) (0.600,0.602) (0.602,0.628) (0.604,0.541) (0.605,0.508) (0.607,0.495) (0.609,0.495) (0.611,0.528) (0.613,0.540) (0.615,0.496) (0.617,0.493) (0.619,0.513) (0.621,0.506) (0.623,0.547) (0.625,0.547) (0.627,0.605) (0.629,0.611) (0.631,0.610) (0.633,0.600) (0.635,0.573) (0.637,0.593) (0.639,0.618) (0.641,0.565) (0.643,0.593) (0.645,0.519) (0.646,0.491) (0.648,0.557) (0.650,0.635) (0.652,0.598) (0.654,0.631) (0.656,0.563) (0.658,0.580) (0.660,0.533) (0.662,0.530) (0.664,0.524) (0.666,0.546) (0.668,0.578) (0.670,0.580) (0.672,0.606) (0.674,0.641) (0.676,0.645) (0.678,0.599) (0.680,0.589) (0.682,0.492) (0.684,0.502) (0.686,0.512) (0.688,0.485) (0.689,0.451) (0.691,0.429) (0.693,0.450) (0.695,0.440) (0.697,0.501) (0.699,0.553) (0.701,0.534) (0.703,0.571) (0.705,0.467) (0.707,0.511) (0.709,0.613) (0.711,0.627) (0.713,0.667) (0.715,0.681) (0.717,0.561) (0.719,0.575) (0.721,0.671) (0.723,0.663) (0.725,0.621) (0.727,0.622) (0.729,0.625) (0.730,0.648) (0.732,0.659) (0.734,0.675) (0.736,0.684) (0.738,0.669) (0.740,0.714) (0.742,0.715) (0.744,0.753) (0.746,0.755) (0.748,0.737) (0.750,0.729) (0.752,0.742) (0.754,0.752) (0.756,0.684) (0.758,0.692) (0.760,0.722) (0.762,0.768) (0.764,0.780) (0.766,0.834) (0.768,0.850) (0.770,0.854) (0.771,0.893) (0.773,0.947) (0.775,0.945) (0.777,0.927) (0.779,0.936) (0.781,0.905) (0.783,0.880) (0.785,0.841) (0.787,0.855) (0.789,0.845) (0.791,0.912) (0.793,0.923) (0.795,0.971) (0.797,1.008) (0.799,1.015) (0.801,1.027) (0.803,1.078) (0.805,1.037) (0.807,1.072) (0.809,1.006) (0.811,0.945) (0.812,0.923) (0.814,0.915) (0.816,0.875) (0.818,0.853) (0.820,0.877) (0.822,0.967) (0.824,0.929) (0.826,0.911) (0.828,0.888) (0.830,0.861) (0.832,0.875) (0.834,0.827) (0.836,0.806) (0.838,0.838) (0.840,0.866) (0.842,0.900) (0.844,0.915) (0.846,0.912) (0.848,0.963) (0.850,0.972) (0.852,0.911) (0.854,0.947) (0.855,0.976) (0.857,1.017) (0.859,0.988) (0.861,1.018) (0.863,0.991) (0.865,1.002) (0.867,0.976) (0.869,0.995) (0.871,0.987) (0.873,0.968) (0.875,1.034) (0.877,1.027) (0.879,1.071) (0.881,1.093) (0.883,1.099) (0.885,1.088) (0.887,1.091) (0.889,1.067) (0.891,1.136) (0.893,1.112) (0.895,1.094) (0.896,1.080) (0.898,1.124) (0.900,1.027) (0.902,1.087) (0.904,1.089) (0.906,1.124) (0.908,1.044) (0.910,1.095) (0.912,1.046) (0.914,1.067) (0.916,1.117) (0.918,1.107) (0.920,1.121) (0.922,1.133) (0.924,1.124) (0.926,1.108) (0.928,1.092) (0.930,1.084) (0.932,1.121) (0.934,1.087) (0.936,1.121) (0.938,1.044) (0.939,1.093) (0.941,1.085) (0.943,1.080) (0.945,1.025) (0.947,0.987) (0.949,0.971) (0.951,0.950) (0.953,0.930) (0.955,0.979) (0.957,0.958) (0.959,0.908) (0.961,0.894) (0.963,0.882) (0.965,0.929) (0.967,0.976) (0.969,0.966) (0.971,0.948) (0.973,0.888) (0.975,0.877) (0.977,0.846) (0.979,0.857) (0.980,0.840) (0.982,0.777) (0.984,0.890) (0.986,0.816) (0.988,0.823) (0.990,0.770) (0.992,0.808) (0.994,0.791) (0.996,0.796) (0.998,0.787) (1.000,0.715)};
\end{axis}
\end{tikzpicture}
\caption{尺度不变性. 左图阴影部分 ($t\le\frac14$) 在时间上放大 $4$ 倍、空间上放大 $2$ 倍后得到右图. 右图仍是一条标准布朗运动的轨道, 在统计意义上与左图无法区分: 布朗轨道是\emph{自相似}的}
\label{fig:bm-scaling}
\end{figure}

\begin{intuition}
尺度不变性说: 把布朗轨道在时间上放大 $c$ 倍、在空间上放大 $\sqrt{c}$ 倍, 看到的仍是布朗运动 (图~\ref{fig:bm-scaling}). 无论用多高倍的 ``显微镜'' 观察, 轨道都同样曲折, 永远不会显得 ``光滑''. 这就预示了 4.3.3 小节的结论: 轨道处处不可微. 时间反演则把 $t\to0$ 时的局部行为与 $t\to\infty$ 时的长期行为联系起来, 二者可以互相转化.
\end{intuition}

\begin{corollary}[布朗运动的强大数定律]\label{cor:bm-slln}
$\lim_{t\to\infty}B_t/t=0$ $\as$
\end{corollary}

\begin{proof}
令 $X$ 为定理~\ref{thm:bm-invariance} (4) 中的时间反演. 则 $B_t/t=X_{1/t}$, 而 $X$ 是标准布朗运动, 轨道在 $0$ 处连续, $X_0=0$, 所以 $X_{1/t}\to0$ ($t\to\infty$) 几乎必然成立.
\end{proof}

这也可以用第~3~章的强大数定律~\ref{thm:slln} 理解: $B_n=\sum_{k=1}^n(B_k-B_{k-1})$ 是 $\iid$ 标准正态变量之和, 所以 $B_n/n\to0$ $\as$; 再控制整数之间的波动, 即得连续时间的结论.

尺度不变性的一个典型用法是从一个参数的结论推出所有参数的结论. 例如, 对 $a>0$ 令 $\tau_a=\inf\{t\ge0: B_t=a\}$. 由 (2), $X_t=\frac{1}{a}B_{a^2t}$ 是布朗运动, 而 $\tau_a=\inf\{t: B_t=a\}=a^2\inf\{u: X_u=1\}$. 所以 $\tau_a\overset{d}{=}a^2\tau_1$: 走到两倍远的地方, 需要四倍长的时间. 同理, $\sup_{s\le t}B_s\overset{d}{=}\sqrt{t}\sup_{s\le1}B_s$.

% ---------------------------------------------------------------------
\subsection{二次变差与无穷变差}

设 $\Pi=\{0=t_0<t_1<\cdots<t_m=t\}$ 是 $[0,t]$ 的分割, $\abs{\Pi}=\max_i(t_i-t_{i-1})$ 为其步长. 记
\[
Q(\Pi)=\sum_{i=1}^m\big(B_{t_i}-B_{t_{i-1}}\big)^2,\qquad V(\Pi)=\sum_{i=1}^m\abs{B_{t_i}-B_{t_{i-1}}}.
\]
对光滑函数 $f$, $\sum_i\abs{f(t_i)-f(t_{i-1})}\to\int_0^t\abs{f'}\dd s<\infty$, 而平方和 $\sum_i(f(t_i)-f(t_{i-1}))^2\le\max_i\abs{f(t_i)-f(t_{i-1})}\cdot\int_0^t\abs{f'}\to0$. 布朗运动的情形正好相反.

\begin{theorem}[布朗运动的二次变差]\label{thm:bm-qv}
设 $\Pi_n$ 是 $[0,t]$ 的分割序列, $\abs{\Pi_n}\to0$. 则
\[
Q(\Pi_n)\xrightarrow{L^2}t.
\]
若进一步 $\sum_n\abs{\Pi_n}<\infty$ (例如二进分割 $t_i=it2^{-n}$), 则 $Q(\Pi_n)\to t$ 几乎必然成立.
\end{theorem}

\begin{proof}
记 $\Delta_i=B_{t_i}-B_{t_{i-1}}$, $\delta_i=t_i-t_{i-1}$. 则 $\Delta_i^2-\delta_i$ ($i=1,\dots,m$) 相互独立, 均值为零. 由于 $\Delta_i\overset{d}{=}\sqrt{\delta_i}Z$, $Z\sim N(0,1)$, 而 $\E Z^4=3$,
\[
\Var\big(\Delta_i^2\big)=\delta_i^2\Var(Z^2)=\delta_i^2(\E Z^4-1)=2\delta_i^2.
\]
所以, 由独立随机变量方差的可加性 (推论~\ref{cor:indep-consequences}),
\begin{equation}\label{eq:qv-variance}
\E\big[(Q(\Pi)-t)^2\big]=\Var\Big(\sum_{i=1}^m\Delta_i^2\Big)=\sum_{i=1}^m2\delta_i^2\le2\abs{\Pi}\sum_{i=1}^m\delta_i=2t\abs{\Pi}.
\end{equation}
当 $\abs{\Pi_n}\to0$ 时右端趋于零, 这就是 $L^2$ 收敛. 若 $\sum_n\abs{\Pi_n}<\infty$, 则由单调收敛定理 (推论~\ref{cor:series}),
\[
\E\Big[\sum_n\big(Q(\Pi_n)-t\big)^2\Big]=\sum_n\E\big[(Q(\Pi_n)-t)^2\big]\le2t\sum_n\abs{\Pi_n}<\infty,
\]
所以 $\sum_n(Q(\Pi_n)-t)^2<\infty$ $\as$, 从而一般项 $Q(\Pi_n)-t\to0$ $\as$
\end{proof}

\begin{remark}
对一般的 (不满足 $\sum_n\abs{\Pi_n}<\infty$ 的) 分割序列, 几乎必然收敛可能不成立; 但只要分割是\emph{嵌套}的 ($\Pi_n\subseteq\Pi_{n+1}$), 几乎必然收敛总成立, 证明可以用逆向鞅收敛定理 (见~\cite{revuz} 第 I 章). 另一方面, 若允许分割依赖于轨道 $\omega$ 本身, 则 $\sup_\Pi Q(\Pi)=\infty$ $\as$ 所以计算 ``二次变差'' 时, 应沿着预先给定的分割序列取极限.
\end{remark}

定理~\ref{thm:bm-qv} 与第~3~章的二次变差定理~\ref{thm:quadratic-variation} 一致: 布朗运动是平方可积的连续鞅 (4.4 节), 它的二次变差过程是
\[
\ip{B}_t=t.
\]
这里的证明是直接计算, 不依赖于那个定理. 由二次变差立即推出无穷变差.

\begin{corollary}[布朗运动的轨道是无穷变差的]\label{cor:bm-infinite-variation}
几乎必然地, 对一切 $0\le a<b$, 布朗轨道在 $[a,b]$ 上的全变差
\[
V_{[a,b]}(B)=\sup_\Pi\sum_i\abs{B_{t_i}-B_{t_{i-1}}}=\infty,
\]
其中上确界取遍 $[a,b]$ 的所有分割.
\end{corollary}

\begin{proof}
先固定有理数 $a<b$. 由平移不变性 (定理~\ref{thm:bm-invariance} (3)) 与定理~\ref{thm:bm-qv}, 沿 $[a,b]$ 的二进分割 $\Pi_n$, 几乎必然有 $Q(\Pi_n)\to b-a>0$. 另一方面, 对任一轨道,
\[
Q(\Pi_n)\le\max_i\abs{B_{t_i}-B_{t_{i-1}}}\cdot\sum_i\abs{B_{t_i}-B_{t_{i-1}}}\le\max_i\abs{B_{t_i}-B_{t_{i-1}}}\cdot V_{[a,b]}(B).
\]
由轨道在 $[a,b]$ 上一致连续, $\max_i\abs{B_{t_i}-B_{t_{i-1}}}\to0$. 若 $V_{[a,b]}(B)<\infty$, 则右端趋于零, 与 $Q(\Pi_n)\to b-a$ 矛盾. 所以 $\P\big(V_{[a,b]}(B)=\infty\big)=1$. 对可数个有理区间取交, 得到一个概率为 $1$ 的事件, 在其上对一切有理数 $a<b$, $V_{[a,b]}(B)=\infty$. 任意区间 $[a,b]$ ($a<b$) 都包含一个有理区间, 而全变差关于区间单调, 结论对一切 $a<b$ 成立.
\end{proof}

\begin{figure}[htbp]
\centering
\begin{tikzpicture}
\begin{axis}[SAaxis, axis x line=bottom, axis y line=left, width=11cm, height=5.6cm, xmin=0.5, xmax=12.8, ymin=0.15, ymax=120, ymode=log,
  xtick={1,...,12}, ytick={0.25,1,4,16,64}, yticklabels={$\frac14$,$1$,$4$,$16$,$64$}, xlabel={$n$ (分割 $[0,1]$ 为 $2^n$ 段)},
  legend style={font=\small,draw=none,fill=none,at={(0.02,0.98)},anchor=north west}, legend cell align=left]
  \addplot[SARed,thick,mark=*,mark size=1.6pt] coordinates {(1.000,0.986) (2.000,0.986) (3.000,0.986) (4.000,2.646) (5.000,4.007) (6.000,6.981) (7.000,9.690) (8.000,13.945) (9.000,19.051) (10.000,26.166) (11.000,36.031) (12.000,51.041)};
  \addplot[SARed!50,dashed,domain=1:12,samples=12] {sqrt(2/pi)*2^(x/2)};
  \addplot[SABlue,thick,mark=square*,mark size=1.6pt] coordinates {(1.000,0.490) (2.000,0.272) (3.000,0.236) (4.000,0.775) (5.000,0.829) (6.000,1.205) (7.000,1.133) (8.000,1.182) (9.000,1.089) (10.000,1.036) (11.000,0.995) (12.000,1.001)};
  \addplot[SABlue!50,dashed,domain=0.5:12.8] {1};
  \legend{$V(\Pi_n)=\sum\abs{\Delta B}$, $\E V(\Pi_n)=\sqrt{2/\pi}\,2^{n/2}$, $Q(\Pi_n)=\sum(\Delta B)^2$, $t=1$}
\end{axis}
\end{tikzpicture}
\caption{沿二进分割计算同一条布朗轨道 (图~\ref{fig:bm-path}) 在 $[0,1]$ 上的一阶变差与二次变差 (纵轴为对数刻度). 一阶变差以 $2^{n/2}$ 的速度发散, 二次变差稳定在 $1$ 附近}
\label{fig:qv-tv}
\end{figure}

\begin{intuition}
为什么一阶变差发散而二次变差收敛? 在长为 $\delta$ 的时间段内, 布朗运动的增量的典型大小是 $\abs{\Delta B}\approx\sqrt{\delta}$. 把 $[0,t]$ 分成 $m$ 段, 每段长 $\delta=t/m$, 则
\[
\sum\abs{\Delta B}\approx m\sqrt{t/m}=\sqrt{mt}\to\infty,\qquad\sum(\Delta B)^2\approx m\cdot\frac{t}{m}=t.
\]
更进一步, 由于 $\sum(\Delta B)^2$ 是大量独立小量之和, 大数定律使它的涨落相对于均值可以忽略, 极限是\emph{确定}的 $t$. 这就是随机微积分中的形式法则
\[
(\dd B_t)^2=\dd t.
\]
在 Taylor 展开 $f(B_{t+\delta})-f(B_t)=f'(B_t)\Delta B+\frac{1}{2}f''(B_t)(\Delta B)^2+\cdots$ 中, 二阶项与一阶项同阶 (都对总和有 $O(1)$ 的贡献), 不能像普通微积分那样略去. 这正是第~5~章 It\^o 公式中出现修正项 $\frac{1}{2}f''(B_t)\dd t$ 的原因. 对 $p>2$, 同样的估计给出 $\sum\abs{\Delta B}^p\approx m(t/m)^{p/2}\to0$, 所以更高阶的项都可以略去.
\end{intuition}

% ---------------------------------------------------------------------
\subsection{处处不可微}

无穷变差说明布朗轨道不是有界变差函数, 因此不可能在一个区间上处处可微且导数有界. 实际上结论要强得多.

\begin{theorem}[Paley-Wiener-Zygmund]\label{thm:nowhere-diff}
几乎必然地, 布朗轨道 $t\mapsto B_t$ 在 $[0,\infty)$ 的任何一点处都不可微. 更确切地, 存在概率为 $1$ 的事件, 在其上对每个 $s\ge0$,
\[
\limsup_{h\downarrow0}\frac{\abs{B_{s+h}-B_s}}{h}=\infty.
\]
\end{theorem}

若轨道在某点 $s$ 处有有限的 (右) 导数, 则在 $s$ 右侧附近, 轨道被一个 Lipschitz 锥控制. 由此, 在 $s$ 右侧的三段相邻的长为 $\frac{1}{n}$ 的小区间上, 增量都不超过 $O(\frac{1}{n})$. 而每段增量的典型大小是 $\frac{1}{\sqrt{n}}$, 三段同时这么小的概率约为 $(\frac{1}{\sqrt{n}})^3=n^{-3/2}$; 对 $s$ 的约 $n$ 个可能位置求和, 总概率约为 $n^{-1/2}\to0$ (图~\ref{fig:pwz}).

\begin{proof}
由平移不变性, 只需对 $s\in[0,1]$ 证明 (对 $[k,k+1]$ 应用于 $B_{k+\cdot}-B_k$, 再对 $k$ 取可列并). 对正整数 $C$ 与 $n$, 令
\[
A_{C,n}=\Big\{\omega:\ \text{存在 }s\in[0,1],\ \text{使得对一切 }t\in\big[s,s+\tfrac{4}{n}\big],\ \abs{B_t-B_s}\le C(t-s)\Big\}.
\]
若在某个 $s\in[0,1]$ 处 $\limsup_{h\downarrow0}\abs{B_{s+h}-B_s}/h<\infty$, 则由上极限的定义, 存在正整数 $C$ 与 $\delta>0$ 使 $\abs{B_t-B_s}\le C(t-s)$ 对 $t\in[s,s+\delta]$ 成立, 取 $n\ge4/\delta$ 即知 $\omega\in A_{C,n}$. 所以只需证明每个 $A_{C,n}$ 包含于一个零概率事件.

固定 $C$. 设 $\omega\in A_{C,n}$, $s$ 如定义. 令 $k=\lceil ns\rceil$, 则 $0\le k\le n$, 且 $s\le\frac{k}{n}<s+\frac{1}{n}$. 对 $j=0,1,2$, 两点 $\frac{k+j}{n}$ 与 $\frac{k+j+1}{n}$ 都落在 $[s,s+\frac{4}{n}]$ 中, 于是
\[
\Big|B_{\frac{k+j+1}{n}}-B_{\frac{k+j}{n}}\Big|\le\Big|B_{\frac{k+j+1}{n}}-B_s\Big|+\Big|B_{\frac{k+j}{n}}-B_s\Big|\le\frac{4C}{n}+\frac{3C}{n}=\frac{7C}{n}.
\]
所以 $A_{C,n}\subseteq E_n$, 其中
\[
E_n=\bigcup_{k=0}^n\bigcap_{j=0}^2\Big\{\Big|B_{\frac{k+j+1}{n}}-B_{\frac{k+j}{n}}\Big|\le\frac{7C}{n}\Big\}.
\]
三个增量相互独立, 都服从 $N(0,\frac{1}{n})$. 对 $Z\sim N(0,1)$ 与 $x>0$, $\P(\abs{Z}\le x)=\int_{-x}^x\varphi\le\frac{2x}{\sqrt{2\pi}}$. 于是
\[
\P(E_n)\le(n+1)\,\P\Big(\abs{Z}\le\frac{7C}{\sqrt{n}}\Big)^3\le(n+1)\Big(\frac{14C}{\sqrt{2\pi n}}\Big)^3\xrightarrow{n\to\infty}0.
\]
又 $A_{C,n}$ 关于 $n$ 单调不减 (区间 $[s,s+\frac{4}{n}]$ 随 $n$ 增大而缩小, 条件变弱), 所以对任意 $m\ge n$, $A_{C,n}\subseteq A_{C,m}\subseteq E_m$, 从而 $A_{C,n}\subseteq\bigcap_{m\ge n}E_m$. 而 $\P\big(\bigcap_{m\ge n}E_m\big)\le\inf_{m\ge n}\P(E_m)=0$. 因此 $\bigcup_{C,n}A_{C,n}$ 包含于可列个零概率事件之并, 在其补集上, 轨道在每个 $s\in[0,1]$ 处满足 $\limsup_{h\downarrow0}\abs{B_{s+h}-B_s}/h=\infty$.
\end{proof}

\begin{figure}[htbp]
\centering
\begin{tikzpicture}[x=1.15cm,y=1.15cm]
  \fill[SAYellow!30] (0,0) -- (5,1.6) -- (5,-1.6) -- cycle;
  \draw[SAYellow!80!SADark,thick] (0,0) -- (5,1.6) node[right,font=\small] {$B_s+C(t-s)$};
  \draw[SAYellow!80!SADark,thick] (0,0) -- (5,-1.6) node[right,font=\small] {$B_s-C(t-s)$};
  \draw[SADark!60,->] (-0.6,0) -- (5.6,0) node[right,font=\small] {$t$};
  \fill[SARed] (0,0) circle (2pt) node[below left,font=\small] {$s$};
  \foreach \x/\l in {0.7/\frac{k}{n},1.95/\frac{k+1}{n},3.2/\frac{k+2}{n},4.45/\frac{k+3}{n}}{
    \draw[SADark!50,dashed] (\x,-1.7) -- (\x,1.7);
    \node[font=\scriptsize,below] at (\x,-1.7) {$\l$};}
  \draw[SABlue,thick] (0,0) .. controls (0.3,0.12) and (0.5,-0.08) .. (0.7,0.1) .. controls (1.0,0.3) and (1.5,-0.2) .. (1.95,0.25)
    .. controls (2.3,0.55) and (2.8,0.1) .. (3.2,0.45) .. controls (3.6,0.8) and (4.1,0.3) .. (4.45,0.6);
  \foreach \x/\y in {0.7/0.1,1.95/0.25,3.2/0.45,4.45/0.6}{\fill[SABlue] (\x,\y) circle (1.6pt);}
  \node[font=\small,SABlue,anchor=west,text width=4.6cm] at (6.6,0.35) {三段增量都 $\le7C/n$\\ 概率 $\approx(C/\sqrt{n})^3$};
\end{tikzpicture}
\caption{处处不可微的证明. 若轨道在 $s$ 处被 Lipschitz 锥 (阴影) 控制, 则 $s$ 右侧三段长为 $\frac1n$ 的格点区间上的增量都很小. 对 $n+1$ 个可能的 $k$ 求和, 总概率 $\lesssim n\cdot n^{-3/2}\to0$}
\label{fig:pwz}
\end{figure}

\begin{remark}
定理~\ref{thm:bm-existence} 说布朗轨道是 $\gamma<\frac{1}{2}$ 阶 H\"older 连续的. 上面的论证稍加修改 (用 $C(t-s)^\gamma$ 代替 $C(t-s)$, 并考虑 $l$ 段而非三段, 使 $l(\gamma-\frac{1}{2})<-1$) 可以证明: 对 $\gamma>\frac{1}{2}$, 轨道在任何一点都不是 $\gamma$ 阶 H\"older 连续的. 临界情形 $\gamma=\frac{1}{2}$ 由 \textbf{L\'evy 连续模定理}精确刻画:
\[
\limsup_{h\downarrow0}\ \sup_{0\le s\le1-h}\frac{\abs{B_{s+h}-B_s}}{\sqrt{2h\log(1/h)}}=1\qquad\as
\]
(见~\cite{karatzas} 定理 2.9.25). 对数因子 $\sqrt{\log(1/h)}$ 来自于对 $1/h$ 个位置取最大值: 它与正态尾概率 $e^{-x^2/2}$ 相匹配. 4.5 节的重对数律描述的则是\emph{固定}时刻附近的振荡, 那里的因子是更小的 $\sqrt{\log\log(1/h)}$.
\end{remark}

\begin{intuition}
处处不可微的轨道曾被认为是病态的数学构造 (Weierstrass 函数), 布朗运动却说明这样的函数是 ``典型的'': 在 Wiener 测度下, 几乎所有连续函数都处处不可微. 物理上, 花粉微粒的速度在每个瞬间都被分子碰撞彻底改变, 不存在有意义的 ``瞬时速度''. 因此 $\dd B_t/\dd t$ 只能作为广义函数 (``白噪声'') 理解, 而形如 $\frac{\dd X_t}{\dd t}=b(X_t)+\sigma(X_t)\frac{\dd B_t}{\dd t}$ 的方程必须改写为积分形式 $X_t=X_0+\int_0^tb(X_s)\dd s+\int_0^t\sigma(X_s)\dd B_s$, 并为第二个积分赋予意义. 这就是第~5~章的任务.
\end{intuition}

% =====================================================================
\section{鞅性质与马尔可夫性}
\label{sec:4-4}

4.3 节研究的是单条轨道的形状. 本节转向布朗运动与信息流的关系: 它在怎样的意义下是 ``公平的'' (鞅性质), 又在怎样的意义下是 ``无记忆的'' (马尔可夫性). 前者使第~3~章的可选停时定理可以用于布朗运动, 后者使我们可以在任意停时处让布朗运动 ``重新开始''. 二者结合, 就是 4.5 节一切计算的基础. 本节中 $B$ 总是 $\{\mathcal{F}_t\}$-布朗运动 (定义~\ref{def:Ft-bm}); $\{\mathcal{F}^B_t\}$ 表示其自然信息流.

% ---------------------------------------------------------------------
\subsection{布朗运动的鞅}

\begin{theorem}[布朗运动的三个基本鞅]\label{thm:bm-martingales}
设 $B$ 是 $\{\mathcal{F}_t\}$-布朗运动. 则以下过程都是关于 $\{\mathcal{F}_t\}$ 的连续鞅:
\begin{enumerate}
\item $B_t$;
\item $B_t^2-t$;
\item $\exp\big(\theta B_t-\frac{1}{2}\theta^2t\big)$, 其中 $\theta\in\R$ 是常数.
\end{enumerate}
\end{theorem}

\begin{proof}
三个过程显然连续且适应. 可积性: $\E\abs{B_t}\le(\E B_t^2)^{1/2}=\sqrt{t}$; 由例~\ref{ex:normal-mgf}, $\E e^{\theta B_t}=\E e^{\theta\sqrt{t}Z}=e^{\theta^2t/2}<\infty$.

固定 $s<t$, 记 $\Delta=B_t-B_s$. 由定义, $\Delta\sim N(0,t-s)$ 且与 $\mathcal{F}_s$ 独立, 而 $B_s$ 是 $\mathcal{F}_s$-可测的. 由独立随机变量的条件期望等于其期望, 以及提出已知因子,

(1) $\E[B_t\,|\,\mathcal{F}_s]=B_s+\E[\Delta\,|\,\mathcal{F}_s]=B_s+\E\Delta=B_s$.

(2) 展开 $B_t^2=B_s^2+2B_s\Delta+\Delta^2$,
\[
\E[B_t^2\,|\,\mathcal{F}_s]=B_s^2+2B_s\E\Delta+\E\Delta^2=B_s^2+(t-s),
\]
即 $\E[B_t^2-t\,|\,\mathcal{F}_s]=B_s^2-s$. (这里 $B_s\Delta$ 的可积性由 Cauchy-Schwarz 不等式保证.)

(3) $e^{\theta B_t}=e^{\theta B_s}e^{\theta\Delta}$, 其中 $e^{\theta\Delta}$ 与 $\mathcal{F}_s$ 独立, 且 $\E e^{\theta\Delta}=e^{\theta^2(t-s)/2}$. 于是
\[
\E\big[e^{\theta B_t-\frac{1}{2}\theta^2t}\,\big|\,\mathcal{F}_s\big]=e^{\theta B_s-\frac{1}{2}\theta^2t}\,e^{\frac{1}{2}\theta^2(t-s)}=e^{\theta B_s-\frac{1}{2}\theta^2s}.\qedhere
\]
\end{proof}

由 (2) 与二次变差的唯一性 (定理~\ref{thm:quadratic-variation}), $\ip{B}_t=t$, 这与 4.3 节直接计算的结果 (定理~\ref{thm:bm-qv}) 一致. 第~5~章的 L\'evy 刻画定理将证明其逆: 一个连续鞅 $M$ 若满足 $M_0=0$ 且 $M_t^2-t$ 也是鞅, 则 $M$ 必是布朗运动.

(3) 中的过程称为\textbf{指数鞅}. 它是离散时间中似然比鞅 (例~\ref{ex:likelihood-ratio}) 的连续版本: 第~5~章的 Girsanov 定理将用它改变概率测度, 从而给布朗运动加上漂移. 本章则主要把它当作计算工具: 对不同的 $\theta$ 应用可选停时定理, 可以得到首达时的 Laplace 变换 (4.5 节).

指数鞅还是一族鞅的 ``生成函数''. 把 $e^{\theta x-\frac{1}{2}\theta^2t}$ 按 $\theta$ 展开为幂级数,
\begin{equation}\label{eq:hermite-generating}
\exp\Big(\theta x-\frac{1}{2}\theta^2t\Big)=\sum_{n=0}^\infty\frac{\theta^n}{n!}H_n(x,t),
\end{equation}
系数 $H_n(x,t)=\partial_\theta^n e^{\theta x-\frac{1}{2}\theta^2t}\big|_{\theta=0}$ 是 $(x,t)$ 的多项式, 称为 (带参数的) \textbf{Hermite 多项式}. 前几个是
\[
H_0=1,\quad H_1=x,\quad H_2=x^2-t,\quad H_3=x^3-3tx,\quad H_4=x^4-6tx^2+3t^2.
\]

\begin{proposition}[Hermite 鞅]\label{prop:hermite-martingales}
对每个 $n\ge0$, $H_n(B_t,t)$ 是关于 $\{\mathcal{F}_t\}$ 的连续鞅.
\end{proposition}

\begin{proof}
记 $g(\theta,x,t)=e^{\theta x-\frac{1}{2}\theta^2t}$. 对 $k$ 归纳可知 $\partial_\theta^kg=q_k(\theta,x,t)\,g$, 其中 $q_k$ 是多项式: $q_0=1$, $q_{k+1}=\partial_\theta q_k+(x-\theta t)q_k$. 于是对固定的 $t$, 存在常数 $C_k$, 使得当 $\abs{\theta}\le1$ 时
\[
\abs{\partial_\theta^kg(\theta,x,t)}\le C_k(1+\abs{x})^ke^{\abs{x}}.
\]
而 $\E\big[(1+\abs{B_t})^ke^{\abs{B_t}}\big]<\infty$ (正态分布的各阶矩与矩母函数都有限). 固定 $s<t$ 与 $A\in\mathcal{F}_s$. 由定理~\ref{thm:bm-martingales} (3), 对一切 $\theta$,
\[
\E\big[g(\theta,B_t,t);A\big]=\E\big[g(\theta,B_s,s);A\big].
\]
由上面的控制, 对两端反复应用期望号下求导定理~\ref{thm:diff-under-E} (在 $\theta\in(-1,1)$ 上), 两端都是 $\theta$ 的光滑函数, 且各阶导数可以在期望号下计算. 令两端的 $n$ 阶导数在 $\theta=0$ 处相等, 得 $\E[H_n(B_t,t);A]=\E[H_n(B_s,s);A]$. 由 $A\in\mathcal{F}_s$ 的任意性, 这就是鞅性质. 可积性已包含在上面的控制中.
\end{proof}

$H_2$ 就是定理~\ref{thm:bm-martingales} (2) 中的鞅. 高阶的 Hermite 鞅可以用来计算停时的高阶矩 (例~\ref{ex:exit-moments}).

\begin{intuition}
三个基本鞅以及全部 Hermite 鞅都形如 $h(B_t,t)$, 而直接验证可知, 它们都满足偏微分方程
\[
\partial_th+\frac{1}{2}\partial_{xx}h=0.
\]
例如对 $h=e^{\theta x-\theta^2t/2}$, $\partial_th=-\frac{\theta^2}{2}h$, $\frac{1}{2}\partial_{xx}h=\frac{\theta^2}{2}h$. 由 Taylor 展开与形式法则 $(\dd B_t)^2=\dd t$ (4.3 节),
\[
\dd h(B_t,t)\approx\partial_xh\,\dd B_t+\Big(\partial_th+\frac{1}{2}\partial_{xx}h\Big)\dd t,
\]
第一项的条件期望为零; 第二项消失, 恰好意味着 $h(B_t,t)$ 没有 ``趋势'', 即是鞅. 第~5~章的 It\^o 公式将使这个论证严格化.
\end{intuition}

% ---------------------------------------------------------------------
\subsection{马尔可夫性与热方程}

下面的技术引理把 ``对有界连续函数成立'' 推广到 ``对有界 Borel 函数成立''. 本节将多次用到它.

\begin{lemma}[从连续函数到 Borel 函数]\label{lem:cont-to-borel}
设 $\mathcal{H}$ 是 $\R^n$ 上一族有界 Borel 函数, 满足函数形式的单调类定理~\ref{thm:functional-mct} 的条件 (1)(2). 若 $\mathcal{H}$ 包含一切形如 $f_1(x_1)f_2(x_2)\cdots f_n(x_n)$ 的函数, 其中 $f_i$ 是有界连续函数, 则 $\mathcal{H}$ 包含 $\R^n$ 上全体有界 Borel 函数.
\end{lemma}

\begin{proof}
对 $\rvec{a}=(a_1,\dots,a_n)\in\R^n$, 记 $O_{\rvec{a}}=\prod_{i=1}^n(-\infty,a_i]$. 令
\[
h_k(u;a)=\min\big\{1,\max\{0,\,1-k(u-a)\}\big\},\qquad g_k(\rvec{x})=\prod_{i=1}^nh_k(x_i;a_i).
\]
$h_k(\cdot;a)$ 连续, 在 $(-\infty,a]$ 上等于 $1$, 在 $[a+\frac{1}{k},\infty)$ 上等于 $0$, 且当 $k\to\infty$ 时单调下降到 $\1_{(-\infty,a]}$. 所以 $g_k\in\mathcal{H}$, 且 $g_k\downarrow\1_{O_{\rvec{a}}}$. 于是 $0\le1-g_k\uparrow1-\1_{O_{\rvec{a}}}$, 由条件 (2) 与线性性, $\1_{O_{\rvec{a}}}=1-(1-\1_{O_{\rvec{a}}})\in\mathcal{H}$. 集族 $\{O_{\rvec{a}}:\rvec{a}\in\R^n\}$ 是生成 $\mathcal{B}(\R^n)$ 的 $\pi$-系, 由定理~\ref{thm:functional-mct}, 结论成立.
\end{proof}

\begin{theorem}[布朗运动的马尔可夫性]\label{thm:bm-markov}
设 $B$ 是 $\{\mathcal{F}_t\}$-布朗运动, $s\ge0$. 则
\begin{enumerate}
\item $W_t=B_{s+t}-B_s$ ($t\ge0$) 是标准布朗运动, 且 $\sigma(W_t:t\ge0)$ 与 $\mathcal{F}_s$ 独立;
\item 对每个有界 $\mathcal{C}$-可测函数 $F:C[0,\infty)\to\R$, 函数 $x\mapsto\E_x[F]$ 是 Borel 函数, 且
\begin{equation}\label{eq:markov-path}
\E\big[F(B_{s+\cdot})\,\big|\,\mathcal{F}_s\big]=\E_{B_s}[F]\qquad\as,
\end{equation}
其中右端表示函数 $x\mapsto\E_x[F]$ 在 $x=B_s$ 处的值.
\end{enumerate}
特别地, $B$ 是关于 $\{\mathcal{F}_t\}$ 的时齐马尔可夫过程, 转移概率为热核 $P_t(x,\cdot)=N(x,t)$ (例~\ref{ex:heat-kernel}).
\end{theorem}

\begin{proof}
(1) $W$ 的轨道连续, $W_0=0$. 任取 $0=u_0<u_1<\cdots<u_n$, 记 $\Delta_i=W_{u_i}-W_{u_{i-1}}=B_{s+u_i}-B_{s+u_{i-1}}$. 我们证明: 对任意 $A\in\mathcal{F}_s$ 与 Borel 集 $C_1,\dots,C_n$,
\begin{equation}\label{eq:markov-induction}
\P\Big(A\cap\bigcap_{i=1}^n\{\Delta_i\in C_i\}\Big)=\P(A)\prod_{i=1}^n\P(\Delta_i\in C_i).
\end{equation}
对 $n$ 归纳. 事件 $A\cap\bigcap_{i<n}\{\Delta_i\in C_i\}$ 属于 $\mathcal{F}_{s+u_{n-1}}$, 而 $\Delta_n$ 与 $\mathcal{F}_{s+u_{n-1}}$ 独立, 所以左端等于 $\P\big(A\cap\bigcap_{i<n}\{\Delta_i\in C_i\}\big)\P(\Delta_n\in C_n)$, 再由归纳假设即得 \eqref{eq:markov-induction}. 取 $A=\Omega$, 可知 $\Delta_1,\dots,\Delta_n$ 相互独立, 且 $\Delta_i\sim N(0,u_i-u_{i-1})$, 所以 $W$ 满足定义~\ref{def:bm}, 是标准布朗运动. 再看独立性: 对固定的 $A\in\mathcal{F}_s$, 满足 $\P(A\cap\{(\Delta_1,\dots,\Delta_n)\in C\})=\P(A)\P((\Delta_1,\dots,\Delta_n)\in C)$ 的 $C\in\mathcal{B}(\R^n)$ 构成 $\lambda$-系, 由 \eqref{eq:markov-induction} 它包含矩形全体构成的 $\pi$-系, 所以由 $\pi$-$\lambda$ 定理~\ref{thm:pi-lambda} 包含 $\mathcal{B}(\R^n)$. 由于 $(W_{u_1},\dots,W_{u_n})$ 是 $(\Delta_1,\dots,\Delta_n)$ 的线性函数, 柱集 $\{(W_{u_1},\dots,W_{u_n})\in C\}$ 都与 $\mathcal{F}_s$ 独立. 柱集构成生成 $\sigma(W)$ 的 $\pi$-系, 由定理~\ref{thm:indep-pi}, $\sigma(W)$ 与 $\mathcal{F}_s$ 独立.

(2) 先设 $F(w)=g\big(w(u_1),\dots,w(u_n)\big)$, $g$ 是 $\R^n$ 上的有界 Borel 函数. 则 $F(B_{s+\cdot})=g(B_s+W_{u_1},\dots,B_s+W_{u_n})$, 其中 $B_s$ 是 $\mathcal{F}_s$-可测的, 而由 (1), $(W_{u_1},\dots,W_{u_n})$ 与 $\mathcal{F}_s$ 独立. 由代入引理~\ref{lem:freezing} (1),
\[
\E\big[F(B_{s+\cdot})\,\big|\,\mathcal{F}_s\big]=\psi(B_s),\qquad\psi(x)=\E\big[g(x+W_{u_1},\dots,x+W_{u_n})\big]=\E_x[F],
\]
最后一个等号是因为 $x+W$ 是从 $x$ 出发的布朗运动, 其分布为 $\P_x$; 代入引理同时保证 $\psi$ 是 Borel 函数. 令 $\mathcal{H}$ 为满足 ``$x\mapsto\E_x[F]$ 是 Borel 函数且 \eqref{eq:markov-path} 成立'' 的有界 $\mathcal{C}$-可测函数 $F$ 全体. $\mathcal{H}$ 是包含常数的线性空间; 若 $0\le F_k\uparrow F$ 且 $F$ 有界, 由单调收敛定理, $\E_x[F_k]\uparrow\E_x[F]$, 所以 $x\mapsto\E_x[F]$ 是 Borel 函数的极限, 仍是 Borel 函数; 由条件单调收敛定理~\ref{thm:cmct}, \eqref{eq:markov-path} 对 $F$ 成立. 上面已证明 $\mathcal{H}$ 包含柱集的示性函数, 而柱集构成生成 $\mathcal{C}$ 的 $\pi$-系, 由定理~\ref{thm:functional-mct}, $\mathcal{H}$ 包含全体有界 $\mathcal{C}$-可测函数.

最后, 取 $F(w)=\1_C(w(t-s))$ ($t\ge s$), \eqref{eq:markov-path} 给出 $\P(B_t\in C\,|\,\mathcal{F}_s)=\P_{B_s}(w(t-s)\in C)=P_{t-s}(B_s,C)$, 其中 $P_{t-s}(x,\cdot)=N(x,t-s)$, 即 $B$ 是以热核为转移概率的时齐马尔可夫过程 (定义~\ref{def:transition}).
\end{proof}

\begin{intuition}
马尔可夫性表明, 在时刻 $s$, 未来的轨道 $B_{s+\cdot}$ 从当前位置 $B_s$ 出发, 是 ``全新的'' 布朗运动, 过去的历史除了通过 $B_s$ 以外不提供任何信息. \eqref{eq:markov-path} 对一切依赖于整条未来轨道的有界泛函 $F$ 成立 (例如 ``未来一年内的最大值超过某个水平''), 不限于单个时刻的取值.
\end{intuition}

热核把布朗运动与偏微分方程联系起来. 对有界 Borel 函数 $f$, 令
\begin{equation}\label{eq:heat-semigroup}
u(t,x)=\E_x\big[f(B_t)\big]=\int_{\R}p_t(x,y)f(y)\dd y,\qquad p_t(x,y)=\frac{1}{\sqrt{2\pi t}}e^{-\frac{(y-x)^2}{2t}}.
\end{equation}

\begin{theorem}[热方程]\label{thm:heat-equation}
\begin{enumerate}
\item 热核满足
\[
\partial_tp_t(x,y)=\frac{1}{2}\partial_{xx}p_t(x,y)=\frac{1}{2}\partial_{yy}p_t(x,y),\qquad t>0,\ x,y\in\R.
\]
\item 设 $f$ 有界连续, $u$ 如 \eqref{eq:heat-semigroup}. 则 $u$ 在 $(0,\infty)\times\R$ 上关于 $t$ 一阶、关于 $x$ 二阶连续可微, 满足
\[
\partial_tu=\frac{1}{2}\partial_{xx}u,\qquad \lim_{t\downarrow0}u(t,x)=f(x).
\]
\end{enumerate}
\end{theorem}

\begin{proof}
(1) 记 $z=y-x$. 直接求导,
\begin{align*}
\partial_tp_t&=p_t\cdot\Big(-\frac{1}{2t}+\frac{z^2}{2t^2}\Big),\\
\partial_xp_t&=p_t\cdot\frac{z}{t},\qquad \partial_{xx}p_t=p_t\cdot\frac{z^2}{t^2}+p_t\cdot\Big(-\frac{1}{t}\Big)=p_t\cdot\Big(\frac{z^2}{t^2}-\frac{1}{t}\Big).
\end{align*}
比较即得 $\partial_tp_t=\frac{1}{2}\partial_{xx}p_t$. 由于 $p_t(x,y)$ 只依赖于 $(y-x)^2$, 关于 $y$ 求两次导数与关于 $x$ 求两次导数的结果相同.

(2) 固定 $0<t_0<t_1$ 与 $R>0$. 由 (1) 的计算, 当 $t\in[t_0,t_1]$, $\abs{x}\le R$ 时, $\partial_tp_t$, $\partial_xp_t$, $\partial_{xx}p_t$ 的绝对值都不超过
\[
C(1+y^2)\,\frac{1}{\sqrt{2\pi t_0}}\exp\Big(-\frac{\big((\abs{y}-R)^+\big)^2}{2t_1}\Big),
\]
其中 $C$ 只依赖于 $t_0,t_1,R$; 它关于 $y$ 可积. 乘以有界函数 $f(y)$ 后仍可积. 由控制收敛定理 (与定理~\ref{thm:diff-under-E} 的证明相同), 可以在积分号下求导, 且导数连续:
\[
\partial_tu=\int\partial_tp_t(x,y)f(y)\dd y=\frac{1}{2}\int\partial_{xx}p_t(x,y)f(y)\dd y=\frac{1}{2}\partial_{xx}u.
\]
初值: $u(t,x)=\E f(x+\sqrt{t}Z)$, $Z\sim N(0,1)$. 当 $t\downarrow0$ 时, 由 $f$ 的连续性 $f(x+\sqrt{t}Z)\to f(x)$, 由有界收敛定理, $u(t,x)\to f(x)$.
\end{proof}

\begin{figure}[htbp]
\centering
\begin{tikzpicture}
\begin{axis}[SAaxis, axis x line=bottom, axis y line=left, width=11cm, height=5.4cm, xmin=-4.2, xmax=4.2, ymin=0, ymax=0.85,
  xtick={-4,-2,0,2,4}, ytick={0,0.2,0.4,0.6,0.8}, xlabel={$y$}, samples=161, domain=-4.2:4.2,
  legend style={font=\small,draw=none,fill=none,at={(0.99,0.98)},anchor=north east}, legend cell align=left]
  \addplot[SARed,thick] {exp(-x^2/(2*0.25))/sqrt(2*pi*0.25)};
  \addplot[SABlue,thick] {exp(-x^2/2)/sqrt(2*pi)};
  \addplot[SAGreenText,thick] {exp(-x^2/(2*4))/sqrt(2*pi*4)};
  \legend{$t=\frac14$, $t=1$, $t=4$}
\end{axis}
\end{tikzpicture}
\caption{热核 $p_t(0,\cdot)$, 即 $B_t$ 的密度. 随着 $t$ 增大, 概率质量像热量一样向两侧扩散: 峰高按 $t^{-1/2}$ 下降, 宽度按 $t^{1/2}$ 增长, 总质量保持为 $1$}
\label{fig:heat-kernel}
\end{figure}

\begin{remark}[向后方程与向前方程]
(1) 中关于 $x$ (起点) 的方程 $\partial_tp=\frac{1}{2}\partial_{xx}p$ 称为 \textbf{Kolmogorov 向后方程}, 它描述期望 $u(t,x)=\E_x f(B_t)$ 如何随起点变化; 关于 $y$ (终点) 的方程 $\partial_tp=\frac{1}{2}\partial_{yy}p$ 称为 \textbf{Kolmogorov 向前方程} (Fokker-Planck 方程), 它描述 $B_t$ 的密度如何随时间演化 (图~\ref{fig:heat-kernel}). 对布朗运动, 二者形式相同; 对第~6~章的一般扩散过程, 二者不同, 分别对应生成元及其伴随算子.
\end{remark}

\begin{intuition}
算子 $\frac{1}{2}\frac{\dd^2}{\dd x^2}$ 是布朗运动的 ``无穷小生成元''. 对光滑的 $f$, 由 Taylor 展开以及 $\E B_t=0$, $\E B_t^2=t$,
\[
\E_x f(B_t)-f(x)\approx f'(x)\E B_t+\frac{1}{2}f''(x)\E B_t^2=\frac{1}{2}f''(x)\,t,
\]
所以 $\frac{\dd}{\dd t}\E_xf(B_t)\big|_{t=0}=\frac{1}{2}f''(x)$. 热方程就是这个关系在所有时刻上的版本 (由马尔可夫性, 每个时刻都可以看成新的起点). 时间反转后, 对固定的 $T$, $v(t,x)=u(T-t,x)$ 满足 $\partial_tv+\frac{1}{2}\partial_{xx}v=0$, 这正是前一小节 Hermite 鞅所满足的方程; 相应地, $v(t,B_t)=\E[f(B_T)\,|\,\mathcal{F}_t]$ 是鞅.
\end{intuition}

% ---------------------------------------------------------------------
\subsection{Blumenthal 0-1 律与通常条件}

第~3~章约定连续时间的信息流满足通常条件 (定义~\ref{def:usual-conditions}). 布朗运动的自然信息流 $\{\mathcal{F}^B_t\}$ 一般并不满足: 它未必完备, 也未必右连续. 例如事件 ``$B$ 在 $0$ 处有右导数'' 只依赖于轨道在 $0$ 的任意小右邻域内的行为, 因而属于 $\mathcal{F}^B_{0+}=\bigcap_{t>0}\mathcal{F}^B_t$; 而 $\mathcal{F}^B_0=\sigma(B_0)$ 只含 $\varnothing$ 与 $\Omega$ 两个事件. 本小节证明, 这个缺陷只是零概率集的问题: $\mathcal{F}^B_{t+}$ 比 $\mathcal{F}^B_t$ 多出来的只有零概率集 (及其补集). 关键是下面的引理.

\begin{lemma}\label{lem:germ}
设 $B$ 是标准布朗运动, $s\ge0$. 则对每个有界 $\mathcal{F}^B_\infty$-可测随机变量 $Y$,
\[
\E\big[Y\,\big|\,\mathcal{F}^B_{s+}\big]=\E\big[Y\,\big|\,\mathcal{F}^B_s\big]\qquad\as
\]
\end{lemma}

\begin{proof}
\emph{第一步: $Y=\prod_{i=1}^nf_i(B_{t_i})$, $f_i$ 有界连续.} 不妨设 $t_1<\cdots<t_n$. 把 $t_i\le s$ 的因子之积记为 $Z$, 其余因子之积记为 $Y'$. $Z$ 是 $\mathcal{F}^B_s$-可测的 (从而也是 $\mathcal{F}^B_{s+}$-可测的), 提出已知因子后, 只需对 $Y'$ 证明. 因此不妨设 $s<t_1$. 对 $0<h<t_1-s$, 在时刻 $s+h$ 应用马尔可夫性~\ref{thm:bm-markov} (2) (信息流取 $\{\mathcal{F}^B_t\}$),
\[
\E\big[Y\,\big|\,\mathcal{F}^B_{s+h}\big]=\phi_h(B_{s+h}),\qquad \phi_h(x)=\E\Big[\prod_{i=1}^nf_i\big(x+B_{t_i-s-h}\big)\Big].
\]
函数 $(h,x)\mapsto\phi_h(x)$ 在 $[0,t_1-s)\times\R$ 上连续: 若 $(h_k,x_k)\to(h,x)$, 由 $f_i$ 与布朗轨道的连续性, 期望号内的被积函数逐点收敛, 而它有界, 由有界收敛定理即得. 由于 $\mathcal{F}^B_{s+}\subseteq\mathcal{F}^B_{s+h}$, 由塔性质,
\[
\E\big[Y\,\big|\,\mathcal{F}^B_{s+}\big]=\E\big[\phi_h(B_{s+h})\,\big|\,\mathcal{F}^B_{s+}\big].
\]
取 $h=h_k\downarrow0$. 由 $\phi$ 的连续性与轨道连续性, $\phi_{h_k}(B_{s+h_k})\to\phi_0(B_s)$ 逐点成立, 且 $\abs{\phi_{h_k}}\le\prod\sup\abs{f_i}$. 由条件控制收敛定理~\ref{thm:cdct}, 右端几乎必然收敛到 $\E[\phi_0(B_s)\,|\,\mathcal{F}^B_{s+}]=\phi_0(B_s)$. 而左端与 $k$ 无关, 所以
\[
\E\big[Y\,\big|\,\mathcal{F}^B_{s+}\big]=\phi_0(B_s)=\E\big[Y\,\big|\,\mathcal{F}^B_s\big],
\]
最后一个等号是马尔可夫性在时刻 $s$ 的情形.

\emph{第二步: 推广.} 固定 $t_1<\cdots<t_n$, 令 $\mathcal{H}_n$ 为使 $Y=g(B_{t_1},\dots,B_{t_n})$ 满足结论的有界 Borel 函数 $g$ 全体. 它是包含常数的线性空间, 由条件单调收敛定理~\ref{thm:cmct} 对单调有界极限封闭, 由第一步包含有界连续函数之积. 由引理~\ref{lem:cont-to-borel}, $\mathcal{H}_n$ 包含全体有界 Borel 函数. 于是对一切柱集 $A=\{(B_{t_1},\dots,B_{t_n})\in C\}$, $Y=\1_A$ 满足结论. 柱集构成生成 $\mathcal{F}^B_\infty$ 的 $\pi$-系, 再用一次函数形式的单调类定理~\ref{thm:functional-mct} (其条件同样由条件单调收敛定理保证), 结论对一切有界 $\mathcal{F}^B_\infty$-可测的 $Y$ 成立.
\end{proof}

\begin{theorem}[Blumenthal 0-1 律]\label{thm:blumenthal}
设 $B$ 是标准布朗运动. 则对每个 $A\in\mathcal{F}^B_{0+}$, $\P(A)\in\{0,1\}$.
\end{theorem}

\begin{proof}
由于 $B_0=0$, $\mathcal{F}^B_0=\{\varnothing,\Omega\}$, 对它的条件期望就是期望. 对 $A\in\mathcal{F}^B_{0+}$, 由引理~\ref{lem:germ} ($s=0$, $Y=\1_A$),
\[
\1_A=\E\big[\1_A\,\big|\,\mathcal{F}^B_{0+}\big]=\E\big[\1_A\,\big|\,\mathcal{F}^B_0\big]=\P(A)\qquad\as
\]
左端只取 $0,1$ 两个值, 所以 $\P(A)\in\{0,1\}$.
\end{proof}

$\mathcal{F}^B_{0+}$ 中的事件只依赖于轨道在 $0$ 的任意小右邻域内的行为, 即所谓 ``芽'' (germ). Blumenthal 0-1 律说这类事件是平凡的, 它是 Kolmogorov 0-1 律~\ref{thm:kolmogorov-01} 在 ``无穷小时刻'' 上的类比: 后者关心 $n\to\infty$ 的尾部, 前者关心 $t\downarrow0$ 的 ``头部''. 它的威力在于, 只要能证明某个芽事件的概率为正, 它的概率就自动等于 $1$.

\begin{corollary}[布朗运动在起点附近的振荡]\label{cor:bm-germ}
以下两条性质几乎必然成立:
\begin{enumerate}
\item 对每个 $\varepsilon>0$, $B$ 在 $(0,\varepsilon)$ 中既取到正值也取到负值, 从而在 $(0,\varepsilon)$ 中有无穷多个零点;
\item $\limsup_{t\downarrow0}\frac{B_t}{\sqrt{t}}=+\infty$, $\liminf_{t\downarrow0}\frac{B_t}{\sqrt{t}}=-\infty$.
\end{enumerate}
\end{corollary}

\begin{proof}
(1) 令 $\tau^+=\inf\{t>0: B_t>0\}$. 由轨道连续性,
\[
\{\tau^+=0\}=\bigcap_{n\ge1}\ \bigcup_{q\in\Q\cap(0,1/n)}\{B_q>0\},
\]
对每个 $m$, 右端等于 $\bigcap_{n\ge m}\bigcup_{q\in\Q\cap(0,1/n)}\{B_q>0\}\in\mathcal{F}^B_{1/m}$, 所以 $\{\tau^+=0\}\in\mathcal{F}^B_{0+}$. 另一方面, 由概率的连续性,
\[
\P(\tau^+=0)=\lim_{n\to\infty}\P\Big(\tau^+<\frac{1}{n}\Big)\ge\limsup_{n\to\infty}\P\Big(B_{\frac{1}{2n}}>0\Big)=\frac{1}{2}.
\]
由 Blumenthal 0-1 律, $\P(\tau^+=0)=1$. 对 $-B$ 应用同样的论证, $\inf\{t>0:B_t<0\}=0$ $\as$ 于是在任意 $(0,\varepsilon)$ 中有 $t_1$ 使 $B_{t_1}>0$, 在 $(0,t_1)$ 中又有 $t_2$ 使 $B_{t_2}<0$, 由介值定理 $(t_2,t_1)$ 中有零点; 如此继续, 得到无穷多个零点.

(2) 固定 $c>0$, 令 $E_c=\{\limsup_{t\downarrow0}B_t/\sqrt{t}\ge c\}$. 由于 $\limsup_{t\downarrow0}$ 只依赖于任意小的 $(0,\delta)$ 上的轨道 (并且由连续性, 可以只取有理时刻), $E_c\in\mathcal{F}^B_\delta$ 对一切 $\delta>0$ 成立, 即 $E_c\in\mathcal{F}^B_{0+}$. 令 $A_n=\{B_{1/n}\ge c/\sqrt{n}\}$, 则 $\limsup_nA_n\subseteq E_c$, 由集合的 Fatou 不等式 (命题~\ref{prop:fatou-sets}),
\[
\P(E_c)\ge\P\Big(\limsup_{n\to\infty}A_n\Big)\ge\limsup_{n\to\infty}\P(A_n)=\P(Z\ge c)>0,
\]
其中 $Z\sim N(0,1)$, 用到了 $\sqrt{n}B_{1/n}\sim N(0,1)$. 由 Blumenthal 0-1 律, $\P(E_c)=1$. 对 $c=1,2,\dots$ 取交即得 $\limsup=+\infty$; 对 $-B$ 应用得 $\liminf=-\infty$.
\end{proof}

(2) 说明, 尽管 $B_t$ 的典型大小是 $\sqrt{t}$, 轨道在 $0$ 附近会无穷多次远远超出 $\pm c\sqrt{t}$ 的范围. 精确的包络由 4.5 节的重对数律给出: $\sqrt{2t\log\log(1/t)}$. 再由时间反演 (定理~\ref{thm:bm-invariance} (4)), (2) 等价于 $\limsup_{t\to\infty}B_t/\sqrt{t}=+\infty$.

现在构造满足通常条件的信息流. 记 $\mathcal{N}=\{N\in\mathcal{F}:\P(N)=0\}$.

\begin{definition}[扩充信息流]\label{def:augmented}
称 $\mathcal{F}_t=\sigma\big(\mathcal{F}^B_t\cup\mathcal{N}\big)$ ($t\ge0$) 为布朗运动的\textbf{扩充信息流} (augmented filtration).
\end{definition}

\begin{lemma}\label{lem:null-augment}
设 $\mathcal{G}\subseteq\mathcal{F}$ 是子 $\sigma$-代数. 则 $A\in\sigma(\mathcal{G}\cup\mathcal{N})$ 当且仅当存在 $A'\in\mathcal{G}$ 使 $\P(A\triangle A')=0$.
\end{lemma}

\begin{proof}
令 $\mathcal{A}$ 为满足右端条件的 $A\in\mathcal{F}$ 全体. $\mathcal{A}\supseteq\mathcal{G}$ (取 $A'=A$), $\mathcal{A}\supseteq\mathcal{N}$ (取 $A'=\varnothing$). $\mathcal{A}$ 是 $\sigma$-代数: 由 $A^c\triangle A'^c=A\triangle A'$ 知它对补封闭; 由 $\big(\bigcup_kA_k\big)\triangle\big(\bigcup_kA'_k\big)\subseteq\bigcup_k(A_k\triangle A'_k)$ 知它对可列并封闭. 所以 $\sigma(\mathcal{G}\cup\mathcal{N})\subseteq\mathcal{A}$. 反之, 若 $A'\in\mathcal{G}$, $N=A\triangle A'\in\mathcal{N}$, 则 $A=A'\triangle N\in\sigma(\mathcal{G}\cup\mathcal{N})$.
\end{proof}

\begin{theorem}[扩充信息流满足通常条件]\label{thm:augmented}
设 $\{\mathcal{F}_t\}$ 是标准布朗运动 $B$ 的扩充信息流. 则
\begin{enumerate}
\item $\{\mathcal{F}_t\}$ 满足通常条件;
\item $B$ 是 $\{\mathcal{F}_t\}$-布朗运动.
\end{enumerate}
\end{theorem}

\begin{proof}
(2) 设 $s<t$, $A\in\mathcal{F}_s$, $C$ 是 Borel 集. 由引理~\ref{lem:null-augment}, 存在 $A'\in\mathcal{F}^B_s$ 使 $\P(A\triangle A')=0$. 由于 $B$ 是关于 $\{\mathcal{F}^B_t\}$ 的布朗运动 (4.2 节),
\begin{align*}
\P\big(A\cap\{B_t-B_s\in C\}\big)&=\P\big(A'\cap\{B_t-B_s\in C\}\big)\\
&=\P(A')\P(B_t-B_s\in C)=\P(A)\P(B_t-B_s\in C).
\end{align*}

(1) 完备性由定义保证. 下证右连续性: 设 $A\in\mathcal{F}_{s+}=\bigcap_{n}\mathcal{F}_{s+1/n}$. 由引理~\ref{lem:null-augment}, 对每个 $n$ 存在 $A_n\in\mathcal{F}^B_{s+1/n}$ 使 $\P(A\triangle A_n)=0$. 令 $A'=\limsup_nA_n$. 对每个 $m$, $A'=\limsup_{n\ge m}A_n$ 只涉及 $A_n\in\mathcal{F}^B_{s+1/n}\subseteq\mathcal{F}^B_{s+1/m}$ ($n\ge m$), 所以 $A'\in\mathcal{F}^B_{s+1/m}$; 对一切 $m$ 成立, 即 $A'\in\mathcal{F}^B_{s+}$. 又 $A\triangle A'\subseteq\bigcup_n(A\triangle A_n)$ (若 $\omega\in A\setminus A'$, 则有某个 $n$ 使 $\omega\notin A_n$; 若 $\omega\in A'\setminus A$, 则有某个 $n$ 使 $\omega\in A_n$), 所以 $\P(A\triangle A')=0$. 由引理~\ref{lem:germ}, $Y:=\E[\1_{A'}\,|\,\mathcal{F}^B_s]$ 满足 $Y=\E[\1_{A'}\,|\,\mathcal{F}^B_{s+}]=\1_{A'}$ $\as$ 令 $A''=\{Y=1\}\in\mathcal{F}^B_s$, 则 $\P(A'\triangle A'')=0$, 从而 $\P(A\triangle A'')=0$. 由引理~\ref{lem:null-augment}, $A\in\mathcal{F}_s$. 这就证明了 $\mathcal{F}_{s+}\subseteq\mathcal{F}_s$, 反向包含是显然的.
\end{proof}

\emph{今后, 若不另作说明, 布朗运动的信息流取为扩充信息流, 或更一般地, 取为满足通常条件且使 $B$ 成为 $\{\mathcal{F}_t\}$-布朗运动的信息流.} 在这样的信息流下, 第~3~章的全部连续时间结论都可以使用: 由命题~\ref{prop:hitting-continuous}, 首达闭集的时刻与首达开集的时刻都是停时; 由定理~\ref{thm:bm-martingales}, 三个基本鞅的停止过程仍是鞅 (推论~\ref{cor:ct-ost-bounded}).

% ---------------------------------------------------------------------
\subsection{强马尔可夫性}

马尔可夫性允许布朗运动在任意\emph{固定}时刻 ``重新开始''. 在应用中, 我们更希望在\emph{随机}时刻重新开始, 例如在首次到达水平 $a$ 的时刻 $\tau_a$. 对停时, 这是可以做到的.

\begin{theorem}[强马尔可夫性]\label{thm:strong-markov}
设 $B$ 是 $\{\mathcal{F}_t\}$-布朗运动, $\tau$ 是停时, $\P(\tau<\infty)>0$. 在 $\{\tau<\infty\}$ 上令
\[
W_t=B_{\tau+t}-B_\tau,\qquad t\ge0.
\]
则对每个 $A\in\mathcal{F}_\tau$ 与每个 $\Gamma\in\mathcal{C}$,
\begin{equation}\label{eq:strong-markov}
\P\big(A\cap\{\tau<\infty\}\cap\{W\in\Gamma\}\big)=\P\big(A\cap\{\tau<\infty\}\big)\,\mathbb{W}(\Gamma),
\end{equation}
其中 $\mathbb{W}$ 是 Wiener 测度. 即在条件概率 $\P(\cdot\,|\,\tau<\infty)$ 下, $W$ 是标准布朗运动, 且与 $\mathcal{F}_\tau$ 独立.
\end{theorem}

证明的思路是用取值于二进点的停时 $\tau_n\downarrow\tau$ 逼近 $\tau$ (引理~\ref{lem:dyadic-approx}). 对只取可数个值的停时, 结论是马尔可夫性对每个取值分别应用的结果; 再由轨道连续性取极限.

\begin{proof}
先说明 $W$ 有定义且可测: 设 $\tau_n=(\lfloor2^n\tau\rfloor+1)/2^n$ 如引理~\ref{lem:dyadic-approx}, 则 $\{\tau_n<\infty\}=\{\tau<\infty\}$, 且在其上 $B_{\tau_n+t}=\sum_{k\ge1}B_{k2^{-n}+t}\1_{\{\tau_n=k2^{-n}\}}$ 是随机变量; 由轨道连续性与 $\tau_n\downarrow\tau$, $B_{\tau+t}=\lim_nB_{\tau_n+t}$ 也是随机变量.

\emph{第一步: 离散化的停时.} 记 $G=A\cap\{\tau<\infty\}$, $W^{(n)}_t=B_{\tau_n+t}-B_{\tau_n}$. 对 $k\ge1$, 记 $s_k=k2^{-n}$. 由于 $A\in\mathcal{F}_\tau$,
\[
A\cap\{\tau_n=s_k\}=\big(A\cap\{\tau<s_k\}\big)\setminus\big(A\cap\{\tau<s_{k-1}\}\big),
\]
而 $A\cap\{\tau<t\}=\bigcup_m\big(A\cap\{\tau\le t-\frac{1}{m}\}\big)\in\mathcal{F}_t$, 所以 $A\cap\{\tau_n=s_k\}\in\mathcal{F}_{s_k}$. 在这个事件上 $W^{(n)}=B_{s_k+\cdot}-B_{s_k}$. 设 $F(w)=\prod_{i=1}^mf_i(w(u_i))$, $f_i$ 有界连续. 由马尔可夫性~\ref{thm:bm-markov} (1), $B_{s_k+\cdot}-B_{s_k}$ 是与 $\mathcal{F}_{s_k}$ 独立的标准布朗运动, 所以
\begin{align*}
\E\big[F(W^{(n)});G\big]&=\sum_{k\ge1}\E\big[F(B_{s_k+\cdot}-B_{s_k});A\cap\{\tau_n=s_k\}\big]\\
&=\sum_{k\ge1}\P\big(A\cap\{\tau_n=s_k\}\big)\,\E\big[F(B)\big]=\P(G)\,\E\big[F(B)\big].
\end{align*}

\emph{第二步: 取极限.} 在 $G$ 上, 由轨道连续性与 $\tau_n\downarrow\tau$, $W^{(n)}_{u_i}\to W_{u_i}$, 从而 $F(W^{(n)})\to F(W)$. 由有界收敛定理,
\[
\E\big[F(W);G\big]=\P(G)\,\E\big[F(B)\big].
\]

\emph{第三步: 推广.} 固定 $u_1<\cdots<u_m$. 满足 $\E[g(W_{u_1},\dots,W_{u_m});G]=\P(G)\E[g(B_{u_1},\dots,B_{u_m})]$ 的有界 Borel 函数 $g$ 构成包含常数的线性空间, 由单调收敛定理对单调有界极限封闭, 由第二步包含有界连续函数之积; 由引理~\ref{lem:cont-to-borel}, 它包含一切有界 Borel 函数. 取 $g$ 为示性函数, 即知 \eqref{eq:strong-markov} 对一切柱集 $\Gamma$ 成立. 满足 \eqref{eq:strong-markov} 的 $\Gamma$ 构成 $\lambda$-系, 柱集构成生成 $\mathcal{C}$ 的 $\pi$-系, 由 $\pi$-$\lambda$ 定理, \eqref{eq:strong-markov} 对一切 $\Gamma\in\mathcal{C}$ 成立. 取 $A=\Omega$ 得 $\P(W\in\Gamma\,|\,\tau<\infty)=\mathbb{W}(\Gamma)$, 再代回 \eqref{eq:strong-markov} 即得条件独立性.
\end{proof}

\begin{figure}[htbp]
\centering
\begin{tikzpicture}
\begin{axis}[SAaxis, axis x line=bottom, axis y line=left, width=12cm, height=6cm, xmin=0, xmax=1.08, ymin=-0.75, ymax=1.35,
  xtick={0,1}, ytick={0,0.6}, yticklabels={$0$,$a$}, xlabel={$t$}, clip=false]
  \addplot[SADark!20] coordinates {(0,0) (1,0)};
  \addplot[SAYellowText,dashed] coordinates {(0,0.6) (1.04,0.6)};
  \addplot[SABlue,thin] coordinates {(0.000,0.000) (0.002,0.028) (0.004,-0.005) (0.006,0.013) (0.008,-0.010) (0.010,0.048) (0.012,0.048) (0.014,-0.033) (0.016,-0.009) (0.018,-0.036) (0.020,-0.005) (0.022,0.020) (0.024,-0.001) (0.026,0.030) (0.028,0.031) (0.030,0.067) (0.032,0.078) (0.034,-0.004) (0.036,-0.003) (0.038,-0.023) (0.040,-0.036) (0.042,-0.010) (0.044,-0.045) (0.046,0.034) (0.048,-0.044) (0.050,-0.126) (0.052,-0.129) (0.054,-0.052) (0.056,-0.042) (0.058,-0.062) (0.060,-0.048) (0.062,-0.037) (0.064,-0.029) (0.066,-0.016) (0.068,0.022) (0.070,0.046) (0.072,0.045) (0.074,0.020) (0.076,0.055) (0.078,0.028) (0.080,0.065) (0.082,-0.016) (0.084,0.024) (0.086,0.049) (0.088,0.030) (0.090,0.043) (0.092,0.057) (0.094,-0.022) (0.096,0.055) (0.098,0.120) (0.100,0.173) (0.102,0.185) (0.104,0.169) (0.106,0.173) (0.108,0.211) (0.110,0.283) (0.112,0.280) (0.114,0.352) (0.116,0.323) (0.118,0.359) (0.120,0.349) (0.122,0.338) (0.124,0.331) (0.126,0.302) (0.128,0.374) (0.130,0.363) (0.132,0.320) (0.134,0.303) (0.136,0.251) (0.138,0.280) (0.140,0.257) (0.142,0.219) (0.144,0.252) (0.146,0.300) (0.148,0.278) (0.150,0.302) (0.152,0.253) (0.154,0.250) (0.156,0.270) (0.158,0.388) (0.160,0.386) (0.162,0.367) (0.164,0.459) (0.166,0.462) (0.168,0.550) (0.170,0.557) (0.172,0.566) (0.174,0.525) (0.176,0.504) (0.178,0.537) (0.180,0.535) (0.182,0.485) (0.184,0.398) (0.186,0.413) (0.188,0.414) (0.190,0.397) (0.192,0.434) (0.194,0.372) (0.196,0.387) (0.198,0.310) (0.200,0.333) (0.202,0.259) (0.204,0.260) (0.206,0.241) (0.208,0.186) (0.210,0.197) (0.212,0.219) (0.214,0.283) (0.216,0.214) (0.218,0.146) (0.220,0.154) (0.222,0.227) (0.224,0.250) (0.226,0.285) (0.228,0.227) (0.230,0.183) (0.232,0.124) (0.234,0.104) (0.236,0.137) (0.238,0.189) (0.240,0.236) (0.242,0.193) (0.244,0.191) (0.246,0.164) (0.248,0.166) (0.250,0.120) (0.252,0.147) (0.254,0.237) (0.256,0.273) (0.258,0.235) (0.260,0.219) (0.262,0.151) (0.264,0.202) (0.266,0.236) (0.268,0.258) (0.270,0.239) (0.272,0.345) (0.274,0.367) (0.276,0.413) (0.278,0.489) (0.280,0.477) (0.282,0.525) (0.284,0.509) (0.286,0.500) (0.288,0.532) (0.290,0.400) (0.292,0.384) (0.294,0.341) (0.296,0.294) (0.298,0.306) (0.300,0.321) (0.302,0.345) (0.304,0.373) (0.306,0.457) (0.308,0.492) (0.310,0.545) (0.312,0.517) (0.314,0.567) (0.316,0.506) (0.318,0.509) (0.320,0.589) (0.322,0.520) (0.324,0.600)};
  \addplot[SARed,thin] coordinates {(0.324,0.633) (0.326,0.610) (0.328,0.617) (0.330,0.559) (0.332,0.690) (0.334,0.598) (0.336,0.565) (0.338,0.551) (0.340,0.622) (0.342,0.671) (0.344,0.657) (0.346,0.660) (0.348,0.698) (0.350,0.714) (0.352,0.709) (0.354,0.727) (0.356,0.770) (0.358,0.795) (0.360,0.835) (0.362,0.854) (0.364,0.828) (0.366,0.896) (0.368,0.904) (0.370,0.876) (0.372,0.906) (0.374,0.954) (0.376,0.991) (0.378,1.010) (0.380,0.960) (0.382,1.004) (0.384,1.030) (0.386,1.056) (0.388,0.997) (0.390,1.017) (0.392,0.963) (0.394,0.969) (0.396,1.040) (0.398,1.000) (0.400,1.025) (0.402,1.029) (0.404,1.017) (0.406,1.040) (0.408,0.988) (0.410,1.036) (0.412,1.035) (0.414,1.026) (0.416,1.032) (0.418,1.024) (0.420,1.055) (0.422,1.048) (0.424,1.042) (0.426,1.042) (0.428,1.118) (0.430,1.110) (0.432,1.040) (0.434,1.032) (0.436,0.965) (0.438,0.994) (0.440,0.934) (0.442,0.860) (0.444,0.870) (0.446,0.958) (0.448,0.929) (0.450,0.911) (0.452,0.874) (0.454,0.935) (0.456,0.862) (0.458,0.897) (0.460,0.835) (0.462,0.850) (0.464,0.804) (0.466,0.867) (0.468,0.882) (0.470,0.945) (0.472,0.945) (0.474,0.928) (0.476,0.942) (0.478,0.873) (0.480,0.864) (0.482,0.907) (0.484,0.880) (0.486,0.865) (0.488,0.912) (0.490,0.843) (0.492,0.851) (0.494,0.765) (0.496,0.761) (0.498,0.717) (0.500,0.743) (0.502,0.783) (0.504,0.850) (0.506,0.819) (0.508,0.846) (0.510,0.859) (0.512,0.784) (0.514,0.751) (0.516,0.753) (0.518,0.721) (0.520,0.683) (0.522,0.705) (0.524,0.705) (0.526,0.749) (0.528,0.771) (0.530,0.781) (0.532,0.768) (0.534,0.770) (0.536,0.756) (0.538,0.716) (0.540,0.711) (0.542,0.657) (0.544,0.686) (0.546,0.715) (0.548,0.671) (0.550,0.633) (0.552,0.693) (0.554,0.752) (0.556,0.741) (0.558,0.720) (0.560,0.668) (0.562,0.743) (0.564,0.649) (0.566,0.672) (0.568,0.681) (0.570,0.720) (0.572,0.702) (0.574,0.696) (0.576,0.704) (0.578,0.805) (0.580,0.734) (0.582,0.707) (0.584,0.705) (0.586,0.808) (0.588,0.770) (0.590,0.711) (0.592,0.730) (0.594,0.719) (0.596,0.773) (0.598,0.790) (0.600,0.772) (0.602,0.741) (0.604,0.666) (0.606,0.631) (0.608,0.663) (0.610,0.653) (0.612,0.704) (0.614,0.721) (0.616,0.761) (0.618,0.698) (0.620,0.737) (0.622,0.688) (0.624,0.666) (0.626,0.671) (0.628,0.725) (0.630,0.765) (0.632,0.816) (0.634,0.813) (0.636,0.794) (0.638,0.709) (0.640,0.719) (0.642,0.715) (0.644,0.691) (0.646,0.715) (0.648,0.707) (0.650,0.692) (0.652,0.765) (0.654,0.814) (0.656,0.865) (0.658,0.945) (0.660,0.922) (0.662,0.978) (0.664,0.964) (0.666,0.943) (0.668,1.049) (0.670,0.988) (0.672,1.031) (0.674,1.043) (0.676,0.954) (0.678,0.948) (0.680,0.952) (0.682,0.898) (0.684,0.946) (0.686,0.909) (0.688,0.926) (0.690,0.996) (0.692,0.984) (0.694,1.008) (0.696,0.978) (0.698,0.996) (0.700,1.063) (0.702,0.966) (0.704,0.970) (0.706,0.950) (0.708,0.974) (0.710,0.985) (0.712,0.930) (0.714,0.957) (0.716,1.068) (0.718,1.027) (0.720,1.045) (0.722,0.957) (0.724,0.999) (0.726,0.991) (0.728,0.938) (0.730,0.929) (0.732,0.906) (0.734,0.964) (0.736,1.021) (0.738,1.002) (0.740,1.000) (0.742,1.026) (0.744,0.884) (0.746,0.891) (0.748,0.915) (0.750,0.963) (0.752,0.951) (0.754,0.891) (0.756,0.808) (0.758,0.794) (0.760,0.703) (0.762,0.704) (0.764,0.716) (0.766,0.811) (0.768,0.813) (0.770,0.770) (0.772,0.787) (0.774,0.751) (0.776,0.727) (0.778,0.691) (0.780,0.740) (0.782,0.711) (0.784,0.773) (0.786,0.748) (0.788,0.737) (0.790,0.758) (0.792,0.725) (0.794,0.693) (0.796,0.677) (0.798,0.645) (0.800,0.656) (0.802,0.658) (0.804,0.677) (0.806,0.717) (0.808,0.692) (0.810,0.678) (0.812,0.655) (0.814,0.753) (0.816,0.769) (0.818,0.766) (0.820,0.813) (0.822,0.780) (0.824,0.797) (0.826,0.785) (0.828,0.772) (0.830,0.815) (0.832,0.863) (0.834,0.827) (0.836,0.828) (0.838,0.763) (0.840,0.793) (0.842,0.786) (0.844,0.863) (0.846,0.887) (0.848,0.885) (0.850,0.957) (0.852,0.997) (0.854,0.932) (0.856,0.885) (0.858,0.829) (0.860,0.801) (0.862,0.780) (0.864,0.768) (0.866,0.780) (0.868,0.688) (0.870,0.646) (0.872,0.639) (0.874,0.695) (0.876,0.659) (0.878,0.646) (0.880,0.614) (0.882,0.549) (0.884,0.589) (0.886,0.555) (0.888,0.638) (0.890,0.664) (0.892,0.668) (0.894,0.617) (0.896,0.675) (0.898,0.675) (0.900,0.746) (0.902,0.730) (0.904,0.770) (0.906,0.816) (0.908,0.767) (0.910,0.723) (0.912,0.691) (0.914,0.678) (0.916,0.590) (0.918,0.528) (0.920,0.508) (0.922,0.507) (0.924,0.488) (0.926,0.435) (0.928,0.437) (0.930,0.388) (0.932,0.388) (0.934,0.404) (0.936,0.463) (0.938,0.421) (0.940,0.416) (0.942,0.393) (0.944,0.396) (0.946,0.405) (0.948,0.407) (0.950,0.390) (0.952,0.406) (0.954,0.367) (0.956,0.427) (0.958,0.283) (0.960,0.324) (0.962,0.312) (0.964,0.265) (0.966,0.216) (0.968,0.199) (0.970,0.205) (0.972,0.111) (0.974,0.067) (0.976,0.091) (0.978,0.083) (0.980,0.029) (0.982,0.024) (0.984,0.020) (0.986,-0.019) (0.988,-0.089) (0.990,-0.192) (0.992,-0.239) (0.994,-0.284) (0.996,-0.278) (0.998,-0.312) (1.000,-0.425)};
  \draw[SARed!70,->,thick] (axis cs:0.324,-0.6) -- (axis cs:0.324,1.3);
  \draw[SARed!70,->,thick] (axis cs:0.324,0.6) -- (axis cs:1.06,0.6) node[right,font=\small] {$s$};
  \node[SARed,font=\small,anchor=south west] at (axis cs:0.324,1.2) {$W_s$};
  \fill[SARed] (axis cs:0.324,0.6) circle (2pt);
  \draw[SADark!50,densely dotted] (axis cs:0.324,0) -- (axis cs:0.324,-0.6);
  \node[font=\small,below] at (axis cs:0.324,-0.62) {$\tau_a$};
\end{axis}
\end{tikzpicture}
\caption{强马尔可夫性. 布朗运动首次到达水平 $a$ 后 (时刻 $\tau_a$), 以 $(\tau_a,a)$ 为新原点 (红色坐标轴) 观察, 此后的轨道 $W_s=B_{\tau_a+s}-a$ (红色) 是标准布朗运动, 且与 $\tau_a$ 之前的历史 (蓝色) 独立}
\label{fig:strong-markov}
\end{figure}

\begin{remark}
(1) 与定理~\ref{thm:bm-markov} (2) 的证明相同, \eqref{eq:strong-markov} 可以改写为: 对有界 $\mathcal{C}$-可测函数 $F$, 在 $\{\tau<\infty\}$ 上
\[
\E\big[F(B_{\tau+\cdot})\,\big|\,\mathcal{F}_\tau\big]=\E_{B_\tau}[F]\qquad\as
\]
(2) 停时的条件是本质的. 设 $L=\sup\{t\le1: B_t=0\}$ 为 $[0,1]$ 中的最后一个零点 (它不是停时: 要判断 $L\le t$, 必须知道 $B$ 在 $(t,1]$ 上不再回到零). 由 $L$ 的定义, $B_{L+t}-B_L=B_{L+t}$ 在 $(0,1-L]$ 上没有零点; 而由推论~\ref{cor:bm-germ} (1), 标准布朗运动在任意 $(0,\varepsilon)$ 中都有零点. 所以 $B_{L+\cdot}-B_L$ 不是布朗运动.
\end{remark}

\begin{intuition}
强马尔可夫性是马尔可夫性在 ``停时'' 这类随机时刻上的推广. 停时是 ``不偷看未来'' 的随机时刻, 在停时处重新开始, 过程 ``不知道'' 自己是在一个特殊的时刻被观察的, 所以此后的行为与在固定时刻重新开始没有区别. 反之, 像最后零点 $L$ 这样需要预知未来才能确定的时刻, 本身就携带了关于未来的信息 (``此后不再回到零''), 重新开始的过程当然有偏.
\end{intuition}

% =====================================================================
\section{首达时、最大值与派生过程}
\label{sec:4-5}

本节把前面的工具付诸计算. 可选停时定理与三个基本鞅给出出口概率与期望出口时间 (4.5.1 小节); 强马尔可夫性给出反射原理, 从而给出最大值与首达时的分布 (4.5.2--4.5.3 小节); 马尔可夫性与反射原理结合, 给出最后零点的反正弦律 (4.5.4 小节). 4.5.5 小节介绍由布朗运动派生的几类重要过程, 最后在 4.5.6 小节证明重对数律. 本节中 $B$ 是标准布朗运动, 信息流满足通常条件, 对 $a\in\R$ 记
\[
\tau_a=\inf\{t\ge0: B_t=a\},\qquad M_t=\max_{s\le t}B_s.
\]
由命题~\ref{prop:hitting-continuous} (1), $\tau_a$ 是停时.

% ---------------------------------------------------------------------
\subsection{双边出口问题}

设 $a,b>0$, 布朗运动从 $0$ 出发, 首次离开区间 $(-a,b)$ 的时刻为
\[
\tau=\inf\{t\ge0: B_t\notin(-a,b)\}=\tau_{-a}\wedge\tau_b.
\]
这是第~3~章赌徒破产问题 (例~\ref{ex:gambler-ruin}) 的连续版本: 赌徒的财富是布朗运动, 输掉 $a$ 时破产, 赢到 $b$ 时离场.

\begin{theorem}[双边出口]\label{thm:exit}
\begin{enumerate}
\item 存在 $q<1$ 使 $\P(\tau>n)\le q^n$ ($n\ge1$). 特别地, $\tau<\infty$ $\as$, 且对一切 $p>0$, $\E\tau^p<\infty$.
\item $\P(B_\tau=b)=\frac{a}{a+b}$, $\P(B_\tau=-a)=\frac{b}{a+b}$.
\item $\E\tau=ab$.
\end{enumerate}
\end{theorem}

\begin{proof}
(1) 若 $\tau>n$, 则 $B$ 在 $[0,n]$ 上始终位于 $(-a,b)$ 中, 从而 $\abs{B_k-B_{k-1}}<a+b$ 对 $k=1,\dots,n$ 成立. 这 $n$ 个增量相互独立且都服从 $N(0,1)$, 所以
\[
\P(\tau>n)\le\P\big(\abs{Z}<a+b\big)^n=:q^n,\qquad q<1.
\]
于是 $\P(\tau=\infty)\le\lim_nq^n=0$. 对 $p>0$, 在 $\{n<\tau\le n+1\}$ 上 $\tau^p\le(n+1)^p$, 所以 $\E\tau^p\le\sum_{n\ge0}(n+1)^p\P(\tau>n)\le\sum_{n\ge0}(n+1)^pq^n<\infty$ (约定 $q^0=1$).

(2) 由定理~\ref{thm:bm-martingales} (1) 与推论~\ref{cor:ct-ost-bounded} (2), 停止过程 $B_{\tau\wedge t}$ 是鞅, 所以 $\E B_{\tau\wedge t}=0$. 由于 $\abs{B_{\tau\wedge t}}\le a\vee b$, 且由 (1), $B_{\tau\wedge t}\to B_\tau$ $\as$ ($t\to\infty$), 由有界收敛定理 $\E B_\tau=0$. 而 $B_\tau$ 只取 $b$ 与 $-a$ 两个值 (轨道连续), 记 $p=\P(B_\tau=b)$, 则
\[
bp-a(1-p)=0,\qquad p=\frac{a}{a+b}.
\]

(3) 由定理~\ref{thm:bm-martingales} (2), $B^2_{\tau\wedge t}-\tau\wedge t$ 是鞅, 所以 $\E\big[B^2_{\tau\wedge t}\big]=\E[\tau\wedge t]$. 令 $t\to\infty$: 左端由有界收敛定理趋于 $\E B_\tau^2$, 右端由单调收敛定理趋于 $\E\tau$. 由 (2),
\[
\E\tau=\E B_\tau^2=b^2\cdot\frac{a}{a+b}+a^2\cdot\frac{b}{a+b}=\frac{ab(a+b)}{a+b}=ab.\qedhere
\]
\end{proof}

这两个公式与简单随机游动的赌徒破产问题 (公平情形) 完全相同, 二者的证明都只用到 ``$S_n$ 与 $S_n^2-n$ 是鞅'', 而布朗运动是随机游动的缩放极限.

\begin{figure}[htbp]
\centering
\begin{tikzpicture}
\begin{axis}[SAaxis, axis x line=bottom, axis y line=left, width=12cm, height=5.6cm, xmin=0, xmax=1.25, ymin=-0.8, ymax=1.25,
  xtick={0,0.5,1}, ytick={-0.6,0,0.9}, yticklabels={$-a$,$0$,$b$}, xlabel={$t$}, clip=false]
  \fill[SAGreen!10] (axis cs:0,-0.6) rectangle (axis cs:1.2,0.9);
  \addplot[SAGreenText,thick] coordinates {(0,0.9) (1.2,0.9)};
  \addplot[SAGreenText,thick] coordinates {(0,-0.6) (1.2,-0.6)};
  \addplot[SADark!20] coordinates {(0,0) (1.2,0)};
  \addplot[SADark!30,thin] coordinates {(0.790,0.940) (0.792,0.961) (0.794,0.950) (0.796,1.051) (0.798,1.020) (0.800,1.046) (0.802,1.014) (0.804,1.049) (0.806,1.057) (0.808,1.050) (0.810,1.027) (0.812,1.036) (0.814,0.979) (0.816,0.971) (0.818,0.969) (0.820,0.991) (0.822,1.017) (0.824,0.983) (0.826,0.972) (0.828,0.978) (0.830,0.958) (0.832,0.898) (0.834,0.943) (0.836,0.952) (0.838,0.978) (0.840,0.990) (0.842,1.016) (0.844,1.037) (0.846,0.954) (0.848,1.010) (0.850,1.001) (0.852,1.011) (0.854,0.960) (0.856,0.885) (0.858,0.819) (0.860,0.835) (0.862,0.817) (0.864,0.888) (0.866,0.876) (0.868,0.860) (0.870,0.800) (0.872,0.819) (0.874,0.816) (0.876,0.805) (0.878,0.797) (0.880,0.868) (0.882,0.795) (0.884,0.771) (0.886,0.799) (0.888,0.794) (0.890,0.779) (0.892,0.687) (0.894,0.752) (0.896,0.726) (0.898,0.781) (0.900,0.844) (0.902,0.832) (0.904,0.833) (0.906,0.854) (0.908,0.833) (0.910,0.807) (0.912,0.833) (0.914,0.841) (0.916,0.861) (0.918,0.905) (0.920,0.912) (0.922,0.940) (0.924,0.910) (0.926,0.912) (0.928,0.992) (0.930,0.980) (0.932,0.937) (0.934,0.899) (0.936,0.920) (0.938,0.928) (0.940,0.976) (0.942,0.941) (0.944,0.899) (0.946,0.952) (0.948,0.952) (0.950,0.917) (0.952,0.881) (0.954,0.905) (0.956,0.889) (0.958,0.916) (0.960,0.888) (0.962,0.842) (0.964,0.834) (0.966,0.804) (0.968,0.783) (0.970,0.794) (0.972,0.800) (0.974,0.744) (0.976,0.706) (0.978,0.696) (0.980,0.726) (0.982,0.680) (0.984,0.635) (0.986,0.613) (0.988,0.597) (0.990,0.644) (0.992,0.675) (0.994,0.683) (0.996,0.669) (0.998,0.726) (1.000,0.715) (1.002,0.724) (1.004,0.734) (1.006,0.713) (1.008,0.777) (1.010,0.779) (1.012,0.812) (1.014,0.815) (1.016,0.775) (1.018,0.828) (1.020,0.870) (1.022,0.911) (1.024,0.941) (1.026,0.961) (1.028,0.993) (1.030,0.988) (1.032,1.010) (1.034,1.001) (1.036,1.019) (1.038,1.068) (1.040,1.121) (1.042,1.107) (1.044,1.067) (1.046,0.968) (1.048,0.942) (1.050,0.874) (1.052,0.930) (1.054,0.969) (1.056,1.028) (1.058,0.989) (1.060,0.889) (1.062,0.909) (1.064,0.854) (1.066,0.842) (1.068,0.880) (1.070,0.929) (1.072,0.923) (1.074,0.870) (1.076,0.829) (1.078,0.895) (1.080,0.979) (1.082,0.889) (1.084,0.842) (1.086,0.825) (1.088,0.856) (1.090,0.813) (1.092,0.773) (1.094,0.712) (1.096,0.653) (1.098,0.650) (1.100,0.674) (1.102,0.717) (1.104,0.735) (1.106,0.724) (1.108,0.758) (1.110,0.773) (1.112,0.834) (1.114,0.847) (1.116,0.797) (1.118,0.873) (1.120,0.914) (1.122,0.880) (1.124,0.835) (1.126,0.839) (1.128,0.822) (1.130,0.813) (1.132,0.825) (1.134,0.809) (1.136,0.791) (1.138,0.824) (1.140,0.724) (1.142,0.806) (1.144,0.744) (1.146,0.665) (1.148,0.654) (1.150,0.658) (1.152,0.682) (1.154,0.596) (1.156,0.554) (1.158,0.501) (1.160,0.504) (1.162,0.460) (1.164,0.521) (1.166,0.452) (1.168,0.417) (1.170,0.406) (1.172,0.492) (1.174,0.507) (1.176,0.449) (1.178,0.380) (1.180,0.392) (1.182,0.438) (1.184,0.410) (1.186,0.490) (1.188,0.538) (1.190,0.478) (1.192,0.546) (1.194,0.562) (1.196,0.502) (1.198,0.545) (1.200,0.524)};
  \addplot[SABlue,thin] coordinates {(0.000,0.000) (0.002,-0.028) (0.004,-0.002) (0.006,0.009) (0.008,0.006) (0.010,0.047) (0.012,0.066) (0.014,0.051) (0.016,0.085) (0.018,0.028) (0.020,0.021) (0.022,-0.034) (0.024,-0.082) (0.026,-0.093) (0.028,-0.120) (0.030,-0.219) (0.032,-0.149) (0.034,-0.132) (0.036,-0.100) (0.038,-0.107) (0.040,-0.107) (0.042,-0.051) (0.044,-0.012) (0.046,-0.108) (0.048,-0.028) (0.050,-0.042) (0.052,0.020) (0.054,0.033) (0.056,0.088) (0.058,-0.007) (0.060,-0.061) (0.062,-0.136) (0.064,-0.127) (0.066,-0.143) (0.068,-0.133) (0.070,-0.118) (0.072,-0.122) (0.074,-0.128) (0.076,-0.219) (0.078,-0.278) (0.080,-0.281) (0.082,-0.276) (0.084,-0.309) (0.086,-0.261) (0.088,-0.270) (0.090,-0.279) (0.092,-0.197) (0.094,-0.251) (0.096,-0.312) (0.098,-0.293) (0.100,-0.274) (0.102,-0.228) (0.104,-0.207) (0.106,-0.186) (0.108,-0.149) (0.110,-0.118) (0.112,-0.219) (0.114,-0.210) (0.116,-0.262) (0.118,-0.233) (0.120,-0.221) (0.122,-0.263) (0.124,-0.236) (0.126,-0.168) (0.128,-0.119) (0.130,-0.101) (0.132,-0.133) (0.134,-0.121) (0.136,-0.158) (0.138,-0.099) (0.140,-0.056) (0.142,-0.008) (0.144,0.056) (0.146,-0.021) (0.148,0.010) (0.150,0.004) (0.152,-0.055) (0.154,-0.071) (0.156,-0.115) (0.158,-0.024) (0.160,-0.062) (0.162,-0.046) (0.164,-0.025) (0.166,0.023) (0.168,0.039) (0.170,0.052) (0.172,-0.056) (0.174,0.023) (0.176,0.048) (0.178,0.029) (0.180,0.026) (0.182,0.076) (0.184,0.115) (0.186,0.102) (0.188,0.075) (0.190,0.086) (0.192,0.041) (0.194,0.102) (0.196,0.109) (0.198,0.089) (0.200,-0.019) (0.202,0.080) (0.204,0.099) (0.206,0.015) (0.208,0.019) (0.210,-0.018) (0.212,0.062) (0.214,0.021) (0.216,0.005) (0.218,-0.009) (0.220,-0.138) (0.222,-0.132) (0.224,-0.198) (0.226,-0.236) (0.228,-0.195) (0.230,-0.168) (0.232,-0.190) (0.234,-0.256) (0.236,-0.252) (0.238,-0.270) (0.240,-0.280) (0.242,-0.303) (0.244,-0.289) (0.246,-0.330) (0.248,-0.348) (0.250,-0.353) (0.252,-0.417) (0.254,-0.379) (0.256,-0.435) (0.258,-0.388) (0.260,-0.347) (0.262,-0.394) (0.264,-0.362) (0.266,-0.411) (0.268,-0.399) (0.270,-0.417) (0.272,-0.367) (0.274,-0.380) (0.276,-0.360) (0.278,-0.337) (0.280,-0.305) (0.282,-0.254) (0.284,-0.255) (0.286,-0.290) (0.288,-0.238) (0.290,-0.298) (0.292,-0.316) (0.294,-0.310) (0.296,-0.325) (0.298,-0.329) (0.300,-0.374) (0.302,-0.409) (0.304,-0.420) (0.306,-0.441) (0.308,-0.343) (0.310,-0.366) (0.312,-0.355) (0.314,-0.301) (0.316,-0.239) (0.318,-0.194) (0.320,-0.180) (0.322,-0.178) (0.324,-0.164) (0.326,-0.148) (0.328,-0.215) (0.330,-0.251) (0.332,-0.213) (0.334,-0.253) (0.336,-0.265) (0.338,-0.226) (0.340,-0.238) (0.342,-0.269) (0.344,-0.314) (0.346,-0.311) (0.348,-0.363) (0.350,-0.420) (0.352,-0.359) (0.354,-0.364) (0.356,-0.386) (0.358,-0.418) (0.360,-0.360) (0.362,-0.374) (0.364,-0.416) (0.366,-0.458) (0.368,-0.415) (0.370,-0.395) (0.372,-0.387) (0.374,-0.435) (0.376,-0.505) (0.378,-0.426) (0.380,-0.383) (0.382,-0.437) (0.384,-0.448) (0.386,-0.495) (0.388,-0.533) (0.390,-0.585) (0.392,-0.539) (0.394,-0.514) (0.396,-0.501) (0.398,-0.484) (0.400,-0.480) (0.402,-0.427) (0.404,-0.355) (0.406,-0.293) (0.408,-0.205) (0.410,-0.288) (0.412,-0.297) (0.414,-0.330) (0.416,-0.372) (0.418,-0.384) (0.420,-0.357) (0.422,-0.348) (0.424,-0.368) (0.426,-0.270) (0.428,-0.269) (0.430,-0.245) (0.432,-0.246) (0.434,-0.236) (0.436,-0.230) (0.438,-0.276) (0.440,-0.198) (0.442,-0.179) (0.444,-0.194) (0.446,-0.194) (0.448,-0.169) (0.450,-0.162) (0.452,-0.120) (0.454,-0.077) (0.456,-0.071) (0.458,-0.011) (0.460,-0.013) (0.462,-0.040) (0.464,-0.127) (0.466,-0.145) (0.468,-0.128) (0.470,-0.136) (0.472,-0.162) (0.474,-0.092) (0.476,-0.029) (0.478,-0.042) (0.480,-0.061) (0.482,-0.047) (0.484,-0.032) (0.486,-0.049) (0.488,-0.141) (0.490,-0.132) (0.492,-0.091) (0.494,0.015) (0.496,0.076) (0.498,0.092) (0.500,0.036) (0.502,0.033) (0.504,0.041) (0.506,0.004) (0.508,0.008) (0.510,0.064) (0.512,0.082) (0.514,0.107) (0.516,0.119) (0.518,0.145) (0.520,0.184) (0.522,0.165) (0.524,0.108) (0.526,0.081) (0.528,0.106) (0.530,0.091) (0.532,0.067) (0.534,0.075) (0.536,0.085) (0.538,0.178) (0.540,0.188) (0.542,0.222) (0.544,0.224) (0.546,0.290) (0.548,0.375) (0.550,0.415) (0.552,0.391) (0.554,0.397) (0.556,0.397) (0.558,0.410) (0.560,0.465) (0.562,0.425) (0.564,0.487) (0.566,0.453) (0.568,0.468) (0.570,0.511) (0.572,0.450) (0.574,0.427) (0.576,0.512) (0.578,0.590) (0.580,0.558) (0.582,0.590) (0.584,0.649) (0.586,0.650) (0.588,0.608) (0.590,0.573) (0.592,0.636) (0.594,0.569) (0.596,0.568) (0.598,0.638) (0.600,0.636) (0.602,0.619) (0.604,0.601) (0.606,0.577) (0.608,0.473) (0.610,0.511) (0.612,0.557) (0.614,0.521) (0.616,0.478) (0.618,0.449) (0.620,0.402) (0.622,0.377) (0.624,0.383) (0.626,0.438) (0.628,0.485) (0.630,0.617) (0.632,0.669) (0.634,0.608) (0.636,0.636) (0.638,0.658) (0.640,0.652) (0.642,0.607) (0.644,0.609) (0.646,0.563) (0.648,0.565) (0.650,0.553) (0.652,0.537) (0.654,0.572) (0.656,0.584) (0.658,0.576) (0.660,0.494) (0.662,0.561) (0.664,0.583) (0.666,0.535) (0.668,0.608) (0.670,0.629) (0.672,0.689) (0.674,0.727) (0.676,0.600) (0.678,0.632) (0.680,0.681) (0.682,0.696) (0.684,0.694) (0.686,0.795) (0.688,0.733) (0.690,0.719) (0.692,0.757) (0.694,0.699) (0.696,0.760) (0.698,0.715) (0.700,0.724) (0.702,0.773) (0.704,0.810) (0.706,0.824) (0.708,0.785) (0.710,0.772) (0.712,0.723) (0.714,0.691) (0.716,0.672) (0.718,0.716) (0.720,0.754) (0.722,0.777) (0.724,0.785) (0.726,0.826) (0.728,0.798) (0.730,0.786) (0.732,0.723) (0.734,0.655) (0.736,0.675) (0.738,0.653) (0.740,0.737) (0.742,0.724) (0.744,0.781) (0.746,0.753) (0.748,0.730) (0.750,0.854) (0.752,0.784) (0.754,0.746) (0.756,0.741) (0.758,0.734) (0.760,0.815) (0.762,0.834) (0.764,0.806) (0.766,0.695) (0.768,0.696) (0.770,0.656) (0.772,0.655) (0.774,0.758) (0.776,0.766) (0.778,0.790) (0.780,0.871) (0.782,0.895) (0.784,0.877) (0.786,0.864) (0.788,0.897) (0.790,0.900)};
  \fill[SARed] (axis cs:0.790,0.9) circle (2pt);
  \draw[SARed,densely dotted] (axis cs:0.790,0.9) -- (axis cs:0.790,-0.8);
  \node[SARed,font=\small,below] at (axis cs:0.790,-0.8) {$\tau$};
\end{axis}
\end{tikzpicture}
\caption{双边出口问题. 布朗运动在时刻 $\tau$ 首次离开带形区域 $(-a,b)$ (这里从上边界离开); 此后的轨道 (灰色) 与问题无关. 从上边界离开的概率为 $\frac{a}{a+b}$, 平均出口时间为 $ab$}
\label{fig:exit}
\end{figure}

\begin{example}[出口时间的高阶矩与 Laplace 变换]\label{ex:exit-moments}
取对称区间 $a=b$, 则 $\E\tau=a^2$. 用 Hermite 鞅 $H_4(B_t,t)=B_t^4-6tB_t^2+3t^2$ (命题~\ref{prop:hermite-martingales}) 计算二阶矩: 停止过程的期望为零,
\[
\E\big[B_{\tau\wedge t}^4-6(\tau\wedge t)B_{\tau\wedge t}^2+3(\tau\wedge t)^2\big]=0.
\]
括号内的量的绝对值不超过 $a^4+6a^2\tau+3\tau^2$, 由定理~\ref{thm:exit} (1) 它是可积的. 令 $t\to\infty$, 由控制收敛定理, 并注意 $B_\tau^2=a^2$,
\[
a^4-6a^2\E\tau+3\E\tau^2=0,\qquad \E\tau^2=\frac{6a^4-a^4}{3}=\frac{5a^4}{3},\qquad\Var(\tau)=\frac{2a^4}{3}.
\]
再用指数鞅求 Laplace 变换. 对 $\lambda>0$, 令 $\theta=\sqrt{2\lambda}$. 由定理~\ref{thm:bm-martingales} (3), $e^{\theta B_t-\lambda t}$ 与 $e^{-\theta B_t-\lambda t}$ 都是鞅, 所以它们的平均 $\cosh(\theta B_t)e^{-\lambda t}$ 也是鞅. 停止过程以 $\cosh(\theta a)$ 为界, 由有界收敛定理与 $\cosh(\theta B_\tau)=\cosh(\theta a)$,
\[
\cosh(\theta a)\,\E e^{-\lambda\tau}=1,\qquad \E e^{-\lambda\tau}=\frac{1}{\cosh\big(a\sqrt{2\lambda}\big)}.
\]
(可以验证, 对 $\lambda$ 求导后在 $\lambda=0$ 处取值, 得到的 $\E\tau$ 与 $\E\tau^2$ 与上面一致.)
\end{example}

\begin{corollary}[常返性]\label{cor:bm-recurrence}
\begin{enumerate}
\item 对每个 $a\in\R$, $\tau_a<\infty$ $\as$
\item 几乎必然地, $\sup_{t\ge0}B_t=+\infty$, $\inf_{t\ge0}B_t=-\infty$. 从而对每个 $x\in\R$, 集合 $\{t: B_t=x\}$ 无界: 布朗运动无穷多次回到每一个水平.
\end{enumerate}
\end{corollary}

\begin{proof}
(1) 设 $a>0$. 对每个 $n$, 若 $B$ 从上边界离开 $(-n,a)$, 则 $\tau_a<\infty$, 所以由定理~\ref{thm:exit} (2), $\P(\tau_a<\infty)\ge\frac{n}{n+a}\to1$. $a<0$ 的情形由对称性 (定理~\ref{thm:bm-invariance} (1)) 得到, $a=0$ 时 $\tau_0=0$.

(2) $\{\sup_tB_t=+\infty\}=\bigcap_{n\ge1}\{\tau_n<\infty\}$, 由 (1) 是可列个概率为 $1$ 的事件之交. 对 $\inf$ 同理. 在这两个事件之交上, 对任意 $T>0$, 由于 $\max_{s\le T}\abs{B_s}<\infty$, 轨道在 $[T,\infty)$ 上仍然无上界也无下界, 由介值定理, 它在 $[T,\infty)$ 上取到每一个值 $x$.
\end{proof}

% ---------------------------------------------------------------------
\subsection{反射原理}

设 $a>0$. 布朗运动首次到达 $a$ 之后, 由强马尔可夫性, 此后的过程是从 $a$ 出发的 ``新'' 布朗运动, 向上与向下的可能性相同. 把 $\tau_a$ 之后的轨道关于水平线 $y=a$ 作镜面反射, 得到的应当仍是布朗运动 (图~\ref{fig:reflection}).

\begin{theorem}[反射原理]\label{thm:reflection}
设 $a>0$. 令
\[
\widetilde{B}_t=\begin{cases}B_t,&t\le\tau_a,\\ 2a-B_t,&t>\tau_a.\end{cases}
\]
则 $\widetilde{B}$ 是标准布朗运动. 由此, 对 $t>0$ 与 $x\le a$,
\begin{equation}\label{eq:reflection-joint}
\P\big(M_t\ge a,\ B_t\le x\big)=\P\big(B_t\ge2a-x\big),
\end{equation}
特别地,
\begin{equation}\label{eq:reflection-max}
\P(M_t\ge a)=2\,\P(B_t\ge a)=\P\big(\abs{B_t}\ge a\big),
\end{equation}
即 $M_t$ 与 $\abs{B_t}$ 同分布.
\end{theorem}

\begin{proof}
由推论~\ref{cor:bm-recurrence}, $\tau_a<\infty$ $\as$; 在一个零概率集上修改轨道, 不妨设处处成立. 令 $Y_t=B_{t\wedge\tau_a}$ (停止在 $\tau_a$ 的轨道), $W_t=B_{\tau_a+t}-a$. 则
\[
B_t=Y_t+W_{(t-\tau_a)^+},\qquad \widetilde{B}_t=Y_t-W_{(t-\tau_a)^+},\qquad t\ge0.
\]
定义映射 $\Psi: C[0,\infty)\times[0,\infty)\times C[0,\infty)\to C[0,\infty)$, $\Psi(y,r,w)(t)=y(t)+w((t-r)^+)$. 每个坐标 $\Psi(\cdot)(t)$ 关于 $r$ 连续、关于 $(y,w)$ 可测, 因而 $\Psi$ 是可测映射. 于是
\[
B=\Psi(Y,\tau_a,W),\qquad\widetilde{B}=\Psi(Y,\tau_a,-W).
\]
$Y_t=B_{t\wedge\tau_a}$ 是 $\mathcal{F}_{t\wedge\tau_a}\subseteq\mathcal{F}_{\tau_a}$-可测的 (第~3~章定义~\ref{def:ct-stopping-time} 之后的说明与命题~\ref{prop:F-tau} (2)), $\tau_a$ 是 $\mathcal{F}_{\tau_a}$-可测的. 由强马尔可夫性~\ref{thm:strong-markov}, $W$ 是与 $\mathcal{F}_{\tau_a}$ 独立的标准布朗运动; 由对称性, $-W$ 也是标准布朗运动, 且同样与 $\mathcal{F}_{\tau_a}$ 独立. 所以 $(Y,\tau_a,W)$ 与 $(Y,\tau_a,-W)$ 的联合分布都等于 $(Y,\tau_a)$ 的分布与 Wiener 测度的乘积, 从而相同. 它们在 $\Psi$ 下的像 $B$ 与 $\widetilde{B}$ 也同分布. $\widetilde{B}$ 的轨道连续 (在 $\tau_a$ 处两段都等于 $a$), 有限维分布与 $B$ 相同, 所以是标准布朗运动.

现在证明 \eqref{eq:reflection-joint}. $\widetilde{B}$ 与 $B$ 在 $[0,\tau_a]$ 上相同, 所以 $\widetilde{B}$ 首次到达 $a$ 的时刻也是 $\tau_a$; 而在 $\{\tau_a\le t\}$ 上 $\widetilde{B}_t=2a-B_t$. 注意 $\{M_t\ge a\}=\{\tau_a\le t\}$, 于是
\[
\big\{M_t\ge a,\ B_t\le x\big\}=\big\{\tau_a\le t,\ \widetilde{B}_t\ge2a-x\big\}=\big\{\widetilde{B}_t\ge2a-x\big\},
\]
最后一个等号是因为 $2a-x\ge a$: 若 $\widetilde{B}_t\ge2a-x\ge a$, 由介值定理 $\widetilde{B}$ 在 $t$ 之前已到达 $a$, 即 $\tau_a\le t$. 由于 $\widetilde{B}_t\overset{d}{=}B_t$, 得 \eqref{eq:reflection-joint}.

对 \eqref{eq:reflection-max}, 把 $\{M_t\ge a\}$ 按 $B_t<a$ 与 $B_t\ge a$ 分成两部分 (后者蕴含 $M_t\ge a$):
\[
\P(M_t\ge a)=\P(M_t\ge a,\ B_t<a)+\P(B_t\ge a).
\]
在 \eqref{eq:reflection-joint} 中令 $x\uparrow a$, 由概率的连续性, $\P(M_t\ge a,B_t<a)=\P(B_t>a)=\P(B_t\ge a)$. 所以 $\P(M_t\ge a)=2\P(B_t\ge a)=\P(\abs{B_t}\ge a)$. 对 $a\le0$, $\P(M_t\ge a)=1=\P(\abs{B_t}\ge a)$ 显然. 所以 $M_t\overset{d}{=}\abs{B_t}$.
\end{proof}

\begin{figure}[htbp]
\centering
\begin{tikzpicture}
\begin{axis}[SAaxis, axis x line=bottom, axis y line=left, width=12cm, height=6.4cm, xmin=0, xmax=1.12, ymin=-0.55, ymax=1.85,
  xtick={0,1}, ytick={0,0.7,1.4}, yticklabels={$0$,$a$,$2a$}, xlabel={$t$}, clip=false]
  \addplot[SADark!20] coordinates {(0,0) (1,0)};
  \addplot[SAYellowText,dashed] coordinates {(0,0.7) (1.05,0.7)};
  \addplot[SABlue,thin] coordinates {(0.000,0.000) (0.002,0.069) (0.004,0.062) (0.006,0.058) (0.008,-0.004) (0.010,-0.036) (0.012,-0.013) (0.014,0.028) (0.016,-0.077) (0.018,-0.058) (0.020,-0.104) (0.022,-0.123) (0.024,-0.101) (0.026,-0.121) (0.028,-0.105) (0.030,-0.097) (0.032,-0.066) (0.034,0.028) (0.036,0.071) (0.038,0.057) (0.040,0.089) (0.042,0.133) (0.044,0.154) (0.046,0.196) (0.048,0.124) (0.050,0.100) (0.052,0.107) (0.054,0.128) (0.056,0.136) (0.058,0.167) (0.060,0.275) (0.062,0.286) (0.064,0.258) (0.066,0.227) (0.068,0.300) (0.070,0.228) (0.072,0.259) (0.074,0.298) (0.076,0.316) (0.078,0.273) (0.080,0.282) (0.082,0.278) (0.084,0.284) (0.086,0.333) (0.088,0.337) (0.090,0.282) (0.092,0.294) (0.094,0.293) (0.096,0.297) (0.098,0.383) (0.100,0.333) (0.102,0.313) (0.104,0.304) (0.106,0.347) (0.108,0.335) (0.110,0.390) (0.112,0.414) (0.114,0.399) (0.116,0.446) (0.118,0.493) (0.120,0.467) (0.122,0.463) (0.124,0.484) (0.126,0.476) (0.128,0.474) (0.130,0.445) (0.132,0.494) (0.134,0.517) (0.136,0.543) (0.138,0.573) (0.140,0.561) (0.142,0.550) (0.144,0.567) (0.146,0.532) (0.148,0.603) (0.150,0.536) (0.152,0.534) (0.154,0.528) (0.156,0.561) (0.158,0.531) (0.160,0.581) (0.162,0.550) (0.164,0.524) (0.166,0.486) (0.168,0.469) (0.170,0.484) (0.172,0.465) (0.174,0.385) (0.176,0.351) (0.178,0.321) (0.180,0.243) (0.182,0.202) (0.184,0.181) (0.186,0.228) (0.188,0.302) (0.190,0.329) (0.192,0.349) (0.194,0.371) (0.196,0.392) (0.198,0.385) (0.200,0.483) (0.202,0.460) (0.204,0.487) (0.206,0.468) (0.208,0.546) (0.210,0.606) (0.212,0.614) (0.214,0.509) (0.216,0.508) (0.218,0.538) (0.220,0.486) (0.222,0.480) (0.224,0.543) (0.226,0.589) (0.228,0.564) (0.230,0.587) (0.232,0.609) (0.234,0.564) (0.236,0.463) (0.238,0.447) (0.240,0.480) (0.242,0.418) (0.244,0.440) (0.246,0.423) (0.248,0.438) (0.250,0.443) (0.252,0.453) (0.254,0.438) (0.256,0.401) (0.258,0.355) (0.260,0.401) (0.262,0.402) (0.264,0.441) (0.266,0.468) (0.268,0.484) (0.270,0.460) (0.272,0.409) (0.274,0.316) (0.276,0.258) (0.278,0.220) (0.280,0.271) (0.282,0.278) (0.284,0.236) (0.286,0.248) (0.288,0.235) (0.290,0.154) (0.292,0.185) (0.294,0.282) (0.296,0.341) (0.298,0.349) (0.300,0.390) (0.302,0.456) (0.304,0.381) (0.306,0.420) (0.308,0.502) (0.310,0.451) (0.312,0.434) (0.314,0.469) (0.316,0.530) (0.318,0.526) (0.320,0.482) (0.322,0.566) (0.324,0.558) (0.326,0.621) (0.328,0.630) (0.330,0.649) (0.332,0.722)};
  \addplot[SABlue,thin] coordinates {(0.332,0.722) (0.334,0.712) (0.336,0.779) (0.338,0.830) (0.340,0.812) (0.342,0.874) (0.344,0.878) (0.346,0.906) (0.348,0.861) (0.350,0.920) (0.352,0.917) (0.354,0.957) (0.356,1.026) (0.358,1.001) (0.360,0.933) (0.362,0.884) (0.364,0.857) (0.366,0.806) (0.368,0.831) (0.370,0.808) (0.372,0.876) (0.374,0.896) (0.376,0.877) (0.378,0.866) (0.380,0.882) (0.382,0.879) (0.384,0.933) (0.386,0.966) (0.388,1.034) (0.390,0.998) (0.392,1.014) (0.394,1.029) (0.396,1.063) (0.398,1.107) (0.400,1.094) (0.402,1.025) (0.404,1.011) (0.406,1.048) (0.408,1.062) (0.410,1.016) (0.412,0.961) (0.414,1.010) (0.416,0.995) (0.418,1.026) (0.420,1.018) (0.422,0.965) (0.424,0.906) (0.426,0.854) (0.428,0.835) (0.430,0.915) (0.432,0.897) (0.434,0.992) (0.436,1.040) (0.438,1.054) (0.440,1.004) (0.442,1.041) (0.444,0.960) (0.446,0.989) (0.448,0.961) (0.450,0.913) (0.452,0.886) (0.454,0.876) (0.456,0.825) (0.458,0.840) (0.460,0.851) (0.462,0.821) (0.464,0.826) (0.466,0.915) (0.468,0.885) (0.470,0.975) (0.472,1.004) (0.474,1.045) (0.476,1.009) (0.478,1.017) (0.480,1.035) (0.482,1.008) (0.484,1.074) (0.486,1.126) (0.488,1.035) (0.490,1.014) (0.492,1.049) (0.494,1.086) (0.496,1.147) (0.498,1.165) (0.500,1.147) (0.502,1.147) (0.504,1.114) (0.506,1.057) (0.508,0.986) (0.510,1.038) (0.512,1.018) (0.514,1.060) (0.516,1.048) (0.518,1.023) (0.520,1.015) (0.522,1.027) (0.524,1.017) (0.526,1.018) (0.528,1.021) (0.530,1.046) (0.532,1.014) (0.534,0.978) (0.536,0.977) (0.538,0.997) (0.540,0.937) (0.542,0.869) (0.544,0.873) (0.546,0.777) (0.548,0.746) (0.550,0.668) (0.552,0.616) (0.554,0.634) (0.556,0.685) (0.558,0.674) (0.560,0.633) (0.562,0.592) (0.564,0.629) (0.566,0.625) (0.568,0.586) (0.570,0.612) (0.572,0.651) (0.574,0.663) (0.576,0.657) (0.578,0.604) (0.580,0.601) (0.582,0.590) (0.584,0.548) (0.586,0.448) (0.588,0.422) (0.590,0.479) (0.592,0.421) (0.594,0.455) (0.596,0.475) (0.598,0.421) (0.600,0.490) (0.602,0.490) (0.604,0.543) (0.606,0.537) (0.608,0.590) (0.610,0.613) (0.612,0.623) (0.614,0.617) (0.616,0.635) (0.618,0.640) (0.620,0.636) (0.622,0.625) (0.624,0.612) (0.626,0.653) (0.628,0.580) (0.630,0.554) (0.632,0.550) (0.634,0.590) (0.636,0.650) (0.638,0.700) (0.640,0.681) (0.642,0.677) (0.644,0.645) (0.646,0.588) (0.648,0.625) (0.650,0.667) (0.652,0.712) (0.654,0.738) (0.656,0.635) (0.658,0.694) (0.660,0.740) (0.662,0.777) (0.664,0.758) (0.666,0.720) (0.668,0.720) (0.670,0.768) (0.672,0.733) (0.674,0.739) (0.676,0.706) (0.678,0.735) (0.680,0.783) (0.682,0.751) (0.684,0.784) (0.686,0.811) (0.688,0.853) (0.690,0.815) (0.692,0.819) (0.694,0.773) (0.696,0.758) (0.698,0.739) (0.700,0.763) (0.702,0.793) (0.704,0.814) (0.706,0.791) (0.708,0.805) (0.710,0.828) (0.712,0.872) (0.714,0.898) (0.716,0.854) (0.718,0.831) (0.720,0.782) (0.722,0.780) (0.724,0.789) (0.726,0.817) (0.728,0.794) (0.730,0.807) (0.732,0.789) (0.734,0.765) (0.736,0.827) (0.738,0.839) (0.740,0.815) (0.742,0.775) (0.744,0.760) (0.746,0.794) (0.748,0.846) (0.750,0.808) (0.752,0.767) (0.754,0.723) (0.756,0.686) (0.758,0.636) (0.760,0.665) (0.762,0.634) (0.764,0.556) (0.766,0.493) (0.768,0.513) (0.770,0.488) (0.772,0.479) (0.774,0.419) (0.776,0.409) (0.778,0.339) (0.780,0.252) (0.782,0.212) (0.784,0.088) (0.786,0.090) (0.788,0.103) (0.790,0.099) (0.792,0.075) (0.794,0.017) (0.796,0.074) (0.798,0.004) (0.800,0.015) (0.802,0.041) (0.804,0.079) (0.806,0.088) (0.808,0.024) (0.810,0.043) (0.812,0.068) (0.814,0.086) (0.816,0.080) (0.818,0.083) (0.820,0.042) (0.822,0.090) (0.824,0.130) (0.826,0.125) (0.828,0.125) (0.830,0.075) (0.832,0.045) (0.834,0.022) (0.836,-0.000) (0.838,-0.062) (0.840,-0.044) (0.842,-0.004) (0.844,-0.004) (0.846,0.032) (0.848,0.083) (0.850,0.055) (0.852,0.081) (0.854,0.060) (0.856,0.107) (0.858,0.101) (0.860,0.173) (0.862,0.229) (0.864,0.248) (0.866,0.227) (0.868,0.290) (0.870,0.203) (0.872,0.204) (0.874,0.236) (0.876,0.299) (0.878,0.253) (0.880,0.250) (0.882,0.307) (0.884,0.306) (0.886,0.306) (0.888,0.315) (0.890,0.289) (0.892,0.275) (0.894,0.312) (0.896,0.236) (0.898,0.189) (0.900,0.176) (0.902,0.119) (0.904,0.111) (0.906,0.072) (0.908,0.108) (0.910,-0.017) (0.912,0.015) (0.914,0.030) (0.916,0.093) (0.918,0.041) (0.920,0.041) (0.922,0.034) (0.924,0.111) (0.926,0.109) (0.928,0.213) (0.930,0.235) (0.932,0.217) (0.934,0.312) (0.936,0.311) (0.938,0.373) (0.940,0.448) (0.942,0.416) (0.944,0.464) (0.946,0.398) (0.948,0.375) (0.950,0.383) (0.952,0.462) (0.954,0.508) (0.956,0.554) (0.958,0.565) (0.960,0.520) (0.962,0.466) (0.964,0.523) (0.966,0.454) (0.968,0.436) (0.970,0.390) (0.972,0.356) (0.974,0.440) (0.976,0.500) (0.978,0.484) (0.980,0.491) (0.982,0.465) (0.984,0.354) (0.986,0.467) (0.988,0.371) (0.990,0.328) (0.992,0.351) (0.994,0.340) (0.996,0.328) (0.998,0.336) (1.000,0.329)};
  \addplot[SARed,thin,densely dashed] coordinates {(0.332,0.678) (0.334,0.688) (0.336,0.621) (0.338,0.570) (0.340,0.588) (0.342,0.526) (0.344,0.522) (0.346,0.494) (0.348,0.539) (0.350,0.480) (0.352,0.483) (0.354,0.443) (0.356,0.374) (0.358,0.399) (0.360,0.467) (0.362,0.516) (0.364,0.543) (0.366,0.594) (0.368,0.569) (0.370,0.592) (0.372,0.524) (0.374,0.504) (0.376,0.523) (0.378,0.534) (0.380,0.518) (0.382,0.521) (0.384,0.467) (0.386,0.434) (0.388,0.366) (0.390,0.402) (0.392,0.386) (0.394,0.371) (0.396,0.337) (0.398,0.293) (0.400,0.306) (0.402,0.375) (0.404,0.389) (0.406,0.352) (0.408,0.338) (0.410,0.384) (0.412,0.439) (0.414,0.390) (0.416,0.405) (0.418,0.374) (0.420,0.382) (0.422,0.435) (0.424,0.494) (0.426,0.546) (0.428,0.565) (0.430,0.485) (0.432,0.503) (0.434,0.408) (0.436,0.360) (0.438,0.346) (0.440,0.396) (0.442,0.359) (0.444,0.440) (0.446,0.411) (0.448,0.439) (0.450,0.487) (0.452,0.514) (0.454,0.524) (0.456,0.575) (0.458,0.560) (0.460,0.549) (0.462,0.579) (0.464,0.574) (0.466,0.485) (0.468,0.515) (0.470,0.425) (0.472,0.396) (0.474,0.355) (0.476,0.391) (0.478,0.383) (0.480,0.365) (0.482,0.392) (0.484,0.326) (0.486,0.274) (0.488,0.365) (0.490,0.386) (0.492,0.351) (0.494,0.314) (0.496,0.253) (0.498,0.235) (0.500,0.253) (0.502,0.253) (0.504,0.286) (0.506,0.343) (0.508,0.414) (0.510,0.362) (0.512,0.382) (0.514,0.340) (0.516,0.352) (0.518,0.377) (0.520,0.385) (0.522,0.373) (0.524,0.383) (0.526,0.382) (0.528,0.379) (0.530,0.354) (0.532,0.386) (0.534,0.422) (0.536,0.423) (0.538,0.403) (0.540,0.463) (0.542,0.531) (0.544,0.527) (0.546,0.623) (0.548,0.654) (0.550,0.732) (0.552,0.784) (0.554,0.766) (0.556,0.715) (0.558,0.726) (0.560,0.767) (0.562,0.808) (0.564,0.771) (0.566,0.775) (0.568,0.814) (0.570,0.788) (0.572,0.749) (0.574,0.737) (0.576,0.743) (0.578,0.796) (0.580,0.799) (0.582,0.810) (0.584,0.852) (0.586,0.952) (0.588,0.978) (0.590,0.921) (0.592,0.979) (0.594,0.945) (0.596,0.925) (0.598,0.979) (0.600,0.910) (0.602,0.910) (0.604,0.857) (0.606,0.863) (0.608,0.810) (0.610,0.787) (0.612,0.777) (0.614,0.783) (0.616,0.765) (0.618,0.760) (0.620,0.764) (0.622,0.775) (0.624,0.788) (0.626,0.747) (0.628,0.820) (0.630,0.846) (0.632,0.850) (0.634,0.810) (0.636,0.750) (0.638,0.700) (0.640,0.719) (0.642,0.723) (0.644,0.755) (0.646,0.812) (0.648,0.775) (0.650,0.733) (0.652,0.688) (0.654,0.662) (0.656,0.765) (0.658,0.706) (0.660,0.660) (0.662,0.623) (0.664,0.642) (0.666,0.680) (0.668,0.680) (0.670,0.632) (0.672,0.667) (0.674,0.661) (0.676,0.694) (0.678,0.665) (0.680,0.617) (0.682,0.649) (0.684,0.616) (0.686,0.589) (0.688,0.547) (0.690,0.585) (0.692,0.581) (0.694,0.627) (0.696,0.642) (0.698,0.661) (0.700,0.637) (0.702,0.607) (0.704,0.586) (0.706,0.609) (0.708,0.595) (0.710,0.572) (0.712,0.528) (0.714,0.502) (0.716,0.546) (0.718,0.569) (0.720,0.618) (0.722,0.620) (0.724,0.611) (0.726,0.583) (0.728,0.606) (0.730,0.593) (0.732,0.611) (0.734,0.635) (0.736,0.573) (0.738,0.561) (0.740,0.585) (0.742,0.625) (0.744,0.640) (0.746,0.606) (0.748,0.554) (0.750,0.592) (0.752,0.633) (0.754,0.677) (0.756,0.714) (0.758,0.764) (0.760,0.735) (0.762,0.766) (0.764,0.844) (0.766,0.907) (0.768,0.887) (0.770,0.912) (0.772,0.921) (0.774,0.981) (0.776,0.991) (0.778,1.061) (0.780,1.148) (0.782,1.188) (0.784,1.312) (0.786,1.310) (0.788,1.297) (0.790,1.301) (0.792,1.325) (0.794,1.383) (0.796,1.326) (0.798,1.396) (0.800,1.385) (0.802,1.359) (0.804,1.321) (0.806,1.312) (0.808,1.376) (0.810,1.357) (0.812,1.332) (0.814,1.314) (0.816,1.320) (0.818,1.317) (0.820,1.358) (0.822,1.310) (0.824,1.270) (0.826,1.275) (0.828,1.275) (0.830,1.325) (0.832,1.355) (0.834,1.378) (0.836,1.400) (0.838,1.462) (0.840,1.444) (0.842,1.404) (0.844,1.404) (0.846,1.368) (0.848,1.317) (0.850,1.345) (0.852,1.319) (0.854,1.340) (0.856,1.293) (0.858,1.299) (0.860,1.227) (0.862,1.171) (0.864,1.152) (0.866,1.173) (0.868,1.110) (0.870,1.197) (0.872,1.196) (0.874,1.164) (0.876,1.101) (0.878,1.147) (0.880,1.150) (0.882,1.093) (0.884,1.094) (0.886,1.094) (0.888,1.085) (0.890,1.111) (0.892,1.125) (0.894,1.088) (0.896,1.164) (0.898,1.211) (0.900,1.224) (0.902,1.281) (0.904,1.289) (0.906,1.328) (0.908,1.292) (0.910,1.417) (0.912,1.385) (0.914,1.370) (0.916,1.307) (0.918,1.359) (0.920,1.359) (0.922,1.366) (0.924,1.289) (0.926,1.291) (0.928,1.187) (0.930,1.165) (0.932,1.183) (0.934,1.088) (0.936,1.089) (0.938,1.027) (0.940,0.952) (0.942,0.984) (0.944,0.936) (0.946,1.002) (0.948,1.025) (0.950,1.017) (0.952,0.938) (0.954,0.892) (0.956,0.846) (0.958,0.835) (0.960,0.880) (0.962,0.934) (0.964,0.877) (0.966,0.946) (0.968,0.964) (0.970,1.010) (0.972,1.044) (0.974,0.960) (0.976,0.900) (0.978,0.916) (0.980,0.909) (0.982,0.935) (0.984,1.046) (0.986,0.933) (0.988,1.029) (0.990,1.072) (0.992,1.049) (0.994,1.060) (0.996,1.072) (0.998,1.064) (1.000,1.071)};
  \fill[SADark] (axis cs:0.332,0.7) circle (1.8pt);
  \draw[SADark!50,densely dotted] (axis cs:0.332,0.7) -- (axis cs:0.332,-0.55) node[below,font=\small,SADark] {$\tau_a$};
  \fill[SABlue] (axis cs:1,0.329) circle (1.8pt) node[right,font=\small] {$B_t=x$};
  \fill[SARed] (axis cs:1,1.071) circle (1.8pt) node[right,font=\small] {$\widetilde{B}_t=2a-x$};
\end{axis}
\end{tikzpicture}
\caption{反射原理. 在首达时 $\tau_a$ 之后, 把轨道 (蓝) 关于水平线 $y=a$ 反射, 得到 $\widetilde{B}$ (红色虚线), 它仍是布朗运动. 这建立了 ``到达过 $a$ 且终点 $\le x$'' 的轨道与 ``终点 $\ge2a-x$'' 的轨道之间的一一对应}
\label{fig:reflection}
\end{figure}

\begin{corollary}[最大值与终点的联合分布]\label{cor:max-joint}
对 $t>0$, $(M_t,B_t)$ 有联合密度
\[
f_{M_t,B_t}(m,x)=\frac{2(2m-x)}{\sqrt{2\pi t^3}}\exp\Big(-\frac{(2m-x)^2}{2t}\Big),\qquad m\ge0,\ x\le m,
\]
在其余区域为零. $M_t$ 的密度为 $\frac{2}{\sqrt{2\pi t}}e^{-m^2/(2t)}$ ($m\ge0$).
\end{corollary}

\begin{proof}
记 $\varphi_t(y)=\frac{1}{\sqrt{2\pi t}}e^{-y^2/(2t)}$. 由 \eqref{eq:reflection-joint}, 对 $m>0$, $x\le m$,
\[
G(m,x):=\P(M_t\ge m,\ B_t\le x)=\int_{2m-x}^\infty\varphi_t(y)\dd y.
\]
若 $(M_t,B_t)$ 有密度 $f$, 则 $G(m,x)=\int_m^\infty\int_{-\infty}^xf(m',x')\dd x'\dd m'$, 所以 $f=-\partial_m\partial_xG$. 直接计算,
\[
\partial_xG=\varphi_t(2m-x),\qquad -\partial_m\partial_xG=-2\varphi_t'(2m-x)=\frac{2(2m-x)}{t}\varphi_t(2m-x),
\]
即为所述的 $f$. 反过来, 对上式定义的 $f$ 逐次积分可以验证 $\int_m^\infty\int_{-\infty}^xf=G(m,x)$, 而形如 $\{M_t\ge m,B_t\le x\}$ 的事件构成决定二维分布的 $\pi$-系 (由推论~\ref{cor:uniqueness}), 所以 $f$ 确实是联合密度 (注意 $M_t\ge B_t$, $M_t\ge0$, 且 $\P(M_t=0)=0$). $M_t$ 的密度由 \eqref{eq:reflection-max} 直接求导得到.
\end{proof}

\begin{remark}
由联合密度可以算出 $M_t-B_t$ 的分布: 它与 $M_t$ 同分布, 即与 $\abs{B_t}$ 同分布. 更深刻的 \textbf{L\'evy 定理} 断言: 作为过程, $\{M_t-B_t\}_{t\ge0}$ 与 $\{\abs{B_t}\}_{t\ge0}$ 同分布 (见~\cite{karatzas} 定理 3.6.17). 回到 4.1 节, 反射原理与 Donsker 不变原理~\ref{thm:donsker} 结合, 给出随机游动最大值的渐近分布 $\P(\max_{k\le n}S_k\ge a\sqrt{n})\to2(1-\Phi(a))$.
\end{remark}

% ---------------------------------------------------------------------
\subsection{首达时的分布}

\begin{theorem}[首达时]\label{thm:hitting-time}
设 $a>0$.
\begin{enumerate}
\item $\P(\tau_a\le t)=2\big(1-\Phi(a/\sqrt{t})\big)$, $\tau_a$ 的密度为
\[
f_{\tau_a}(t)=\frac{a}{\sqrt{2\pi t^3}}\exp\Big(-\frac{a^2}{2t}\Big),\qquad t>0.
\]
等价地, $\tau_a\overset{d}{=}a^2/Z^2$, 其中 $Z\sim N(0,1)$.
\item $\P(\tau_a<\infty)=1$, 但 $\E\tau_a=\infty$.
\item 对 $\lambda>0$, $\E e^{-\lambda\tau_a}=e^{-a\sqrt{2\lambda}}$.
\end{enumerate}
\end{theorem}

\begin{proof}
(1) $\{\tau_a\le t\}=\{M_t\ge a\}$, 由反射原理~\ref{thm:reflection},
\[
\P(\tau_a\le t)=\P(\abs{B_t}\ge a)=\P\big(\abs{Z}\ge a/\sqrt{t}\big)=2\big(1-\Phi(a/\sqrt{t})\big).
\]
另一方面 $\P(a^2/Z^2\le t)=\P(Z^2\ge a^2/t)$ 与之相同, 所以 $\tau_a\overset{d}{=}a^2/Z^2$. 对 $t$ 求导,
\[
f_{\tau_a}(t)=-2\varphi\Big(\frac{a}{\sqrt{t}}\Big)\cdot\frac{\dd}{\dd t}\frac{a}{\sqrt{t}}=2\varphi\Big(\frac{a}{\sqrt{t}}\Big)\cdot\frac{a}{2t^{3/2}}=\frac{a}{\sqrt{2\pi t^3}}e^{-a^2/(2t)}.
\]

(2) 第一个结论已在推论~\ref{cor:bm-recurrence} 中证明 (也可在 (1) 中令 $t\to\infty$). 对期望, 由尾积分公式 (例~\ref{ex:tail-formula}),
\[
\E\tau_a=\int_0^\infty\P(\tau_a>t)\dd t=\int_0^\infty\P\big(\abs{Z}<a/\sqrt{t}\big)\dd t.
\]
当 $t\ge a^2$ 时, $a/\sqrt{t}\le1$, 而在 $[-1,1]$ 上 $\varphi\ge\varphi(1)$, 所以 $\P(\abs{Z}<a/\sqrt{t})\ge2\varphi(1)\,a/\sqrt{t}$. 由于 $\int_{a^2}^\infty t^{-1/2}\dd t=\infty$, $\E\tau_a=\infty$.

(3) 令 $\theta=\sqrt{2\lambda}$. 由定理~\ref{thm:bm-martingales} (3), $N_t=\exp(\theta B_t-\lambda t)$ 是鞅, 其停止过程 $N_{t\wedge\tau_a}$ 也是鞅, 所以 $\E N_{t\wedge\tau_a}=1$. 在 $[0,\tau_a]$ 上 $B\le a$, 所以 $0<N_{t\wedge\tau_a}\le e^{\theta a}$. 由于 $\tau_a<\infty$ $\as$, 当 $t\to\infty$ 时 $N_{t\wedge\tau_a}\to e^{\theta a-\lambda\tau_a}$ $\as$ 由有界收敛定理,
\[
\E e^{\theta a-\lambda\tau_a}=1,\qquad\E e^{-\lambda\tau_a}=e^{-\theta a}=e^{-a\sqrt{2\lambda}}.\qedhere
\]
\end{proof}

\begin{figure}[htbp]
\centering
\begin{tikzpicture}
\begin{axis}[SAaxis, axis x line=bottom, axis y line=left, width=11cm, height=5.4cm, xmin=0, xmax=6.2, ymin=0, ymax=2,
  xtick={0,1,...,6}, ytick={0,0.5,1,1.5,2}, xlabel={$t$}, samples=300,
  legend style={font=\small,draw=none,fill=none,at={(0.99,0.98)},anchor=north east}, legend cell align=left]
  \addplot[SARed,thick,domain=0.004:6.2] {0.5/sqrt(2*pi*x^3)*exp(-0.25/(2*x))};
  \addplot[SABlue,thick,domain=0.02:6.2] {1/sqrt(2*pi*x^3)*exp(-1/(2*x))};
  \addplot[SAGreenText,thick,domain=0.05:6.2] {2/sqrt(2*pi*x^3)*exp(-4/(2*x))};
  \legend{$a=\frac12$, $a=1$, $a=2$}
\end{axis}
\end{tikzpicture}
\caption{首达时 $\tau_a$ 的密度. 峰值位于 $t=a^2/3$, 而右侧尾部只以 $t^{-3/2}$ 的速度衰减, 衰减太慢以致期望无穷. 由尺度不变性, $\tau_a\overset{d}{=}a^2\tau_1$}
\label{fig:hitting-density}
\end{figure}

\begin{intuition}
(2) 初看令人惊讶, 布朗运动必然会到达 $a$, 但平均需要无穷长的时间. 原因在于尾部 $\P(\tau_a>t)\sim a\sqrt{2/(\pi t)}$ 衰减太慢, 轨道有相当的概率先向下走得很远, 再花极长的时间返回. 若与双边出口比较, 在 $-n$ 处设一个吸收壁, 则 $\E(\tau_a\wedge\tau_{-n})=an$, 令 $n\to\infty$ 即趋于无穷. 用 Laplace 变换同样得到 $\E\tau_a=\lim_{\lambda\downarrow0}\frac{1-\E e^{-\lambda\tau_a}}{\lambda}=\lim_{\lambda\downarrow0}\frac{1-e^{-a\sqrt{2\lambda}}}{\lambda}=\infty$.
\end{intuition}

\begin{remark}[首达时过程]
对 $0<a<b$, 由于轨道连续, 布朗运动到达 $b$ 之前必先到达 $a$, 所以 $\tau_b=\tau_a+\inf\{t: B_{\tau_a+t}-a=b-a\}$. 由强马尔可夫性, $\tau_b-\tau_a$ 与 $\mathcal{F}_{\tau_a}$ 独立, 且与 $\tau_{b-a}$ 同分布. 所以过程 $\{\tau_a\}_{a\ge0}$ 单调不减, 且具有独立平稳增量, 其 Laplace 变换 $e^{-a\sqrt{2\lambda}}=\big(e^{-\sqrt{2\lambda}}\big)^a$ 关于 $a$ 呈指数形式. 这是指数为 $\frac{1}{2}$ 的\textbf{稳定从属过程}的标准例子, 它的轨道是纯跳的: 每一次跳跃对应于布朗运动的一段向下的 ``远足''.
\end{remark}

% ---------------------------------------------------------------------
\subsection{反正弦律}

设 $t>0$. 令
\[
L_t=\sup\{s\le t: B_s=0\}
\]
为 $[0,t]$ 中的最后一个零点. 由轨道连续性, 零点集是闭集且包含 $0$, 所以 $L_t\le s$ 当且仅当 $B$ 在 $(s,t]$ 上没有零点.

\begin{theorem}[L\'evy 反正弦律]\label{thm:arcsine}
对 $0\le s\le t$,
\[
\P(L_t\le s)=\frac{2}{\pi}\arcsin\sqrt{\frac{s}{t}}.
\]
特别地, $L_1$ 的密度为 $\frac{1}{\pi\sqrt{s(1-s)}}$, $0<s<1$.
\end{theorem}

\begin{proof}
$s=0$ 的情形由推论~\ref{cor:bm-germ} (1) 得到 (两端都是 $0$), $s=t$ 的情形是平凡的. 设 $0<s<t$. 令 $F(w)=\1\{w(v)\ne0,\ \forall v\in(0,t-s]\}$, 由轨道连续性它是 $\mathcal{C}$-可测的 (可以只用有理时刻表达). 则 $\{L_t\le s\}=\{F(B_{s+\cdot})=1\}$. 由马尔可夫性~\ref{thm:bm-markov} (2),
\[
\P(L_t\le s)=\E\big[\E[F(B_{s+\cdot})\,|\,\mathcal{F}_s]\big]=\E\big[h(B_s)\big],\qquad h(x)=\E_x[F].
\]
计算 $h$. 设 $x>0$, $W$ 是标准布朗运动. $x+W$ 在 $(0,t-s]$ 上没有零点, 当且仅当 $W$ 在 $[0,t-s]$ 上没有到达 $-x$, 即 $\min_{v\le t-s}W_v>-x$. 由对称性 $-W$ 也是布朗运动, 再由反射原理~\ref{thm:reflection},
\[
h(x)=\P\Big(\max_{v\le t-s}W_v<x\Big)=1-\P\big(\abs{W_{t-s}}\ge x\big)=\P\big(\abs{W_{t-s}}<x\big).
\]
由对称性, 对 $x<0$ 有 $h(x)=\P(\abs{W_{t-s}}<\abs{x})$; 由推论~\ref{cor:bm-germ} (1), $h(0)=0$, 上式对 $x=0$ 也成立. 取 $W$ 与 $B_s$ 独立, 由代入引理~\ref{lem:freezing},
\[
\P(L_t\le s)=\E h(B_s)=\P\big(\abs{W_{t-s}}<\abs{B_s}\big)=\P\big(\sqrt{t-s}\,\abs{Y}<\sqrt{s}\,\abs{X}\big),
\]
其中 $X=B_s/\sqrt{s}$ 与 $Y=W_{t-s}/\sqrt{t-s}$ 是独立的标准正态变量. 把 $(X,Y)$ 写成极坐标 $(R\cos\Theta,R\sin\Theta)$, 由例~\ref{ex:box-muller}, $\Theta$ 在 $(0,2\pi)$ 上均匀分布. 令 $\alpha\in(0,\frac{\pi}{2})$ 满足 $\tan\alpha=\sqrt{s/(t-s)}$, 即 $\sin\alpha=\sqrt{s/t}$. 于是
\[
\P\Big(\frac{\abs{Y}}{\abs{X}}<\tan\alpha\Big)=\P\big(\abs{\tan\Theta}<\tan\alpha\big)=\frac{4\alpha}{2\pi}=\frac{2}{\pi}\arcsin\sqrt{\frac{s}{t}},
\]
这里 $\abs{\tan\Theta}<\tan\alpha$ 当且仅当 $\Theta$ 落在 $0$ 或 $\pi$ 的 $\alpha$-邻域内 (模 $2\pi$), 总长度为 $4\alpha$. 对 $s$ 求导即得 $L_1$ 的密度.
\end{proof}

\begin{figure}[htbp]
\centering
\begin{tikzpicture}
\begin{axis}[SAaxis, axis x line=bottom, axis y line=left, width=7.8cm, height=5.4cm, xmin=0, xmax=1.04, ymin=-1.1, ymax=0.75,
  xtick={0,0.5,1}, ytick={-1,-0.5,0,0.5}, xlabel={$t$}, clip=false]
  \fill[SAYellow!30] (axis cs:0.259,-1.1) rectangle (axis cs:1,0.75);
  \addplot[SADark!30] coordinates {(0,0) (1,0)};
  \addplot[SABlue,thin] coordinates {(0.000,0.000) (0.003,0.018) (0.005,0.012) (0.007,0.009) (0.010,-0.046) (0.013,0.056) (0.015,0.018) (0.018,0.085) (0.020,0.071) (0.022,0.069) (0.025,0.107) (0.028,0.044) (0.030,-0.001) (0.033,-0.071) (0.035,-0.038) (0.037,-0.034) (0.040,0.060) (0.043,0.112) (0.045,0.142) (0.048,0.142) (0.050,0.113) (0.052,0.029) (0.055,0.016) (0.058,0.038) (0.060,0.120) (0.062,0.161) (0.065,0.243) (0.068,0.298) (0.070,0.333) (0.072,0.292) (0.075,0.242) (0.077,0.178) (0.080,0.218) (0.083,0.200) (0.085,0.081) (0.088,0.170) (0.090,0.159) (0.092,0.126) (0.095,0.131) (0.098,0.211) (0.100,0.246) (0.103,0.337) (0.105,0.352) (0.107,0.302) (0.110,0.308) (0.113,0.320) (0.115,0.412) (0.118,0.417) (0.120,0.416) (0.122,0.349) (0.125,0.301) (0.128,0.195) (0.130,0.148) (0.133,0.150) (0.135,0.122) (0.138,0.064) (0.140,0.109) (0.143,0.142) (0.145,0.084) (0.147,0.129) (0.150,0.092) (0.152,0.038) (0.155,0.093) (0.158,0.018) (0.160,-0.029) (0.163,0.017) (0.165,0.003) (0.168,0.032) (0.170,0.119) (0.173,0.088) (0.175,0.149) (0.177,0.167) (0.180,0.254) (0.182,0.215) (0.185,0.279) (0.188,0.232) (0.190,0.219) (0.193,0.239) (0.195,0.276) (0.198,0.340) (0.200,0.304) (0.203,0.256) (0.205,0.179) (0.208,0.112) (0.210,0.128) (0.212,0.063) (0.215,0.025) (0.217,-0.013) (0.220,-0.045) (0.223,-0.070) (0.225,-0.072) (0.228,-0.037) (0.230,-0.024) (0.233,-0.083) (0.235,0.031) (0.238,0.029) (0.240,0.010) (0.242,0.043) (0.245,0.016) (0.247,0.023) (0.250,-0.044) (0.253,-0.026) (0.255,0.002) (0.258,-0.006) (0.260,-0.025) (0.263,-0.061) (0.265,-0.120) (0.268,-0.132) (0.270,-0.161) (0.273,-0.155) (0.275,-0.199) (0.278,-0.291) (0.280,-0.249) (0.283,-0.227) (0.285,-0.161) (0.288,-0.257) (0.290,-0.184) (0.292,-0.163) (0.295,-0.219) (0.297,-0.233) (0.300,-0.152) (0.302,-0.184) (0.305,-0.160) (0.307,-0.114) (0.310,-0.174) (0.312,-0.275) (0.315,-0.370) (0.318,-0.294) (0.320,-0.258) (0.323,-0.239) (0.325,-0.307) (0.328,-0.310) (0.330,-0.322) (0.333,-0.380) (0.335,-0.433) (0.338,-0.387) (0.340,-0.333) (0.343,-0.282) (0.345,-0.233) (0.348,-0.147) (0.350,-0.224) (0.352,-0.255) (0.355,-0.252) (0.357,-0.240) (0.360,-0.169) (0.362,-0.183) (0.365,-0.192) (0.367,-0.201) (0.370,-0.194) (0.372,-0.203) (0.375,-0.220) (0.378,-0.241) (0.380,-0.212) (0.383,-0.264) (0.385,-0.303) (0.388,-0.299) (0.390,-0.347) (0.393,-0.307) (0.395,-0.318) (0.398,-0.352) (0.400,-0.374) (0.403,-0.504) (0.405,-0.454) (0.408,-0.471) (0.410,-0.424) (0.413,-0.406) (0.415,-0.413) (0.417,-0.410) (0.420,-0.425) (0.422,-0.349) (0.425,-0.338) (0.427,-0.223) (0.430,-0.239) (0.432,-0.285) (0.435,-0.297) (0.438,-0.368) (0.440,-0.443) (0.443,-0.496) (0.445,-0.626) (0.448,-0.601) (0.450,-0.645) (0.453,-0.686) (0.455,-0.713) (0.458,-0.737) (0.460,-0.752) (0.463,-0.688) (0.465,-0.730) (0.468,-0.801) (0.470,-0.757) (0.473,-0.756) (0.475,-0.759) (0.478,-0.731) (0.480,-0.756) (0.482,-0.758) (0.485,-0.686) (0.487,-0.786) (0.490,-0.689) (0.492,-0.639) (0.495,-0.681) (0.497,-0.746) (0.500,-0.763) (0.503,-0.733) (0.505,-0.681) (0.508,-0.701) (0.510,-0.709) (0.512,-0.679) (0.515,-0.677) (0.517,-0.686) (0.520,-0.735) (0.522,-0.733) (0.525,-0.708) (0.527,-0.634) (0.530,-0.730) (0.532,-0.855) (0.535,-0.813) (0.537,-0.851) (0.540,-0.875) (0.542,-0.887) (0.545,-0.804) (0.547,-0.717) (0.550,-0.751) (0.552,-0.701) (0.555,-0.673) (0.557,-0.647) (0.560,-0.702) (0.562,-0.695) (0.565,-0.681) (0.568,-0.645) (0.570,-0.716) (0.573,-0.729) (0.575,-0.740) (0.578,-0.757) (0.580,-0.826) (0.583,-0.871) (0.585,-0.846) (0.588,-0.857) (0.590,-0.865) (0.593,-0.845) (0.595,-0.967) (0.598,-0.957) (0.600,-1.002) (0.603,-0.911) (0.605,-0.825) (0.608,-0.779) (0.610,-0.782) (0.613,-0.740) (0.615,-0.641) (0.618,-0.661) (0.620,-0.759) (0.623,-0.740) (0.625,-0.774) (0.628,-0.747) (0.630,-0.687) (0.633,-0.647) (0.635,-0.601) (0.638,-0.667) (0.640,-0.653) (0.642,-0.634) (0.645,-0.633) (0.647,-0.580) (0.650,-0.592) (0.652,-0.592) (0.655,-0.683) (0.657,-0.589) (0.660,-0.693) (0.662,-0.664) (0.665,-0.679) (0.667,-0.670) (0.670,-0.699) (0.672,-0.666) (0.675,-0.678) (0.677,-0.629) (0.680,-0.699) (0.682,-0.639) (0.685,-0.610) (0.688,-0.578) (0.690,-0.671) (0.693,-0.640) (0.695,-0.708) (0.698,-0.688) (0.700,-0.635) (0.703,-0.645) (0.705,-0.557) (0.708,-0.485) (0.710,-0.526) (0.713,-0.483) (0.715,-0.418) (0.718,-0.368) (0.720,-0.332) (0.723,-0.303) (0.725,-0.259) (0.728,-0.202) (0.730,-0.300) (0.733,-0.361) (0.735,-0.388) (0.738,-0.413) (0.740,-0.382) (0.743,-0.307) (0.745,-0.246) (0.748,-0.200) (0.750,-0.176) (0.753,-0.125) (0.755,-0.073) (0.758,-0.037) (0.760,-0.068) (0.763,-0.072) (0.765,-0.101) (0.768,-0.128) (0.770,-0.173) (0.772,-0.118) (0.775,-0.288) (0.777,-0.321) (0.780,-0.346) (0.782,-0.343) (0.785,-0.338) (0.787,-0.354) (0.790,-0.342) (0.792,-0.281) (0.795,-0.348) (0.797,-0.319) (0.800,-0.337) (0.802,-0.266) (0.805,-0.266) (0.807,-0.330) (0.810,-0.371) (0.812,-0.426) (0.815,-0.404) (0.818,-0.417) (0.820,-0.375) (0.823,-0.268) (0.825,-0.317) (0.828,-0.362) (0.830,-0.357) (0.833,-0.409) (0.835,-0.396) (0.838,-0.368) (0.840,-0.321) (0.843,-0.296) (0.845,-0.357) (0.848,-0.386) (0.850,-0.418) (0.853,-0.471) (0.855,-0.515) (0.858,-0.638) (0.860,-0.706) (0.863,-0.699) (0.865,-0.670) (0.868,-0.606) (0.870,-0.577) (0.873,-0.569) (0.875,-0.590) (0.878,-0.580) (0.880,-0.593) (0.883,-0.621) (0.885,-0.589) (0.888,-0.579) (0.890,-0.507) (0.893,-0.496) (0.895,-0.409) (0.897,-0.482) (0.900,-0.507) (0.902,-0.495) (0.905,-0.499) (0.907,-0.482) (0.910,-0.439) (0.912,-0.487) (0.915,-0.586) (0.917,-0.614) (0.920,-0.615) (0.922,-0.596) (0.925,-0.629) (0.927,-0.688) (0.930,-0.655) (0.932,-0.692) (0.935,-0.650) (0.938,-0.618) (0.940,-0.629) (0.943,-0.647) (0.945,-0.681) (0.948,-0.671) (0.950,-0.587) (0.953,-0.610) (0.955,-0.574) (0.958,-0.599) (0.960,-0.670) (0.963,-0.612) (0.965,-0.601) (0.968,-0.568) (0.970,-0.552) (0.973,-0.502) (0.975,-0.578) (0.978,-0.664) (0.980,-0.677) (0.983,-0.645) (0.985,-0.609) (0.988,-0.557) (0.990,-0.513) (0.993,-0.574) (0.995,-0.597) (0.998,-0.509) (1.000,-0.616)};
  \fill[SARed] (axis cs:0.259,0) circle (2pt);
  \draw[SARed,densely dotted] (axis cs:0.259,0) -- (axis cs:0.259,-1.1) node[below,font=\small] {$L_1$};
\end{axis}
\end{tikzpicture}\hfill
\begin{tikzpicture}
\begin{axis}[SAaxis, axis x line=bottom, axis y line=left, width=7.8cm, height=5.4cm, xmin=0, xmax=1.04, ymin=0, ymax=4,
  xtick={0,0.1,0.5,1}, xticklabels={$0$,$0.1$,$0.5$,$1$}, ytick={0,1,2,3,4}, xlabel={$s$}, samples=200]
  \addplot[draw=none,fill=SARed!20,domain=0.0005:0.1] {1/(pi*sqrt(x*(1-x)))} \closedcycle;
  \addplot[SABlue,thick,domain=0.008:0.992] {1/(pi*sqrt(x*(1-x)))};
  \addplot[SADark!40,dashed] coordinates {(0,1) (1,1)};
  \node[font=\scriptsize,SARed,anchor=west] at (axis cs:0.12,2.4) {$\P(L_1\le0.1)\approx0.205$};
\end{axis}
\end{tikzpicture}
\caption{左: 一条布朗轨道及其在 $[0,1]$ 中的最后零点 $L_1$; 阴影区间内轨道不再回到零. 右: 反正弦密度 $\frac{1}{\pi\sqrt{s(1-s)}}$ (蓝) 与均匀密度 (灰色虚线) 的对比. 最后零点最可能出现在两端附近}
\label{fig:arcsine}
\end{figure}

\begin{intuition}
反正弦密度在两端趋于无穷, 在中点最小, 最后零点 $L_1$ 因而更可能出现在很早或很晚的时刻: $\P(L_1\le0.1)=\frac{2}{\pi}\arcsin\sqrt{0.1}\approx0.205$, 即有约五分之一的概率, 在后面 $90\%$ 的时间里布朗运动都停留在零的同一侧. 换成公平赌博的语言: 两个赌徒长时间对赌, 领先者易手的次数远比直觉中少, 一方长期领先是常态. 这与 ``大数定律使双方各领先一半时间'' 的直觉相反; 大数定律管的是平均值, 而 ``谁领先'' 由部分和的符号决定, 部分和的波动 (量级 $\sqrt{n}$) 并不会被平均掉.
\end{intuition}

\begin{remark}[另外两个反正弦律]
L\'evy 还证明了: 布朗运动在 $[0,1]$ 中处于正半轴的总时间 $\int_0^1\1_{\{B_s>0\}}\dd s$, 以及 $B$ 在 $[0,1]$ 上取到最大值的 (几乎必然唯一的) 时刻, 都服从同一个反正弦分布. 证明见~\cite{morters} 第 5 章或~\cite{karatzas} 4.4 节. 第一个结论的离散版本 (随机游动在正半轴的时间比例) 可由 Donsker 不变原理~\ref{thm:donsker} 得到.
\end{remark}

% ---------------------------------------------------------------------
\subsection{由布朗运动派生的过程}

第~5--6~章的随机微分方程的解都由布朗运动 ``生长'' 出来. 本小节介绍几个可以直接用布朗运动写出的例子, 它们分别是第~6~章中线性方程的解, 也是应用中最常见的模型.

\paragraph{带漂移的布朗运动.} 设 $\mu\in\R$, $\sigma>0$. 称 $X_t=\mu t+\sigma B_t$ 为漂移为 $\mu$、扩散系数为 $\sigma$ 的布朗运动. 它的增量独立, $X_t-X_s\sim N(\mu(t-s),\sigma^2(t-s))$. 由于 $\sigma B_t/t\to0$ (推论~\ref{cor:bm-slln}), 长期行为由漂移决定: $X_t/t\to\mu$ $\as$ 但首达时的分布对漂移的符号极为敏感. 由尺度变换, 不妨设 $\sigma=1$.

\begin{theorem}[带漂移布朗运动的首达时]\label{thm:drift-hitting}
设 $X_t=\mu t+B_t$, $a>0$, $\tau_a^\mu=\inf\{t: X_t=a\}$. 则对 $\lambda>0$,
\[
\E\Big[e^{-\lambda\tau^\mu_a}\1_{\{\tau^\mu_a<\infty\}}\Big]=\exp\Big(a\big(\mu-\sqrt{\mu^2+2\lambda}\big)\Big).
\]
因此 $\P(\tau^\mu_a<\infty)=1$ ($\mu\ge0$), $\P(\tau^\mu_a<\infty)=e^{-2\abs{\mu}a}$ ($\mu<0$); 当 $\mu>0$ 时 $\E\tau_a^\mu=a/\mu$.
\end{theorem}

\begin{proof}
对 $\theta\in\R$, 由定理~\ref{thm:bm-martingales} (3),
\[
\exp\Big(\theta B_t-\frac{1}{2}\theta^2t\Big)=\exp\Big(\theta X_t-\big(\theta\mu+\frac{1}{2}\theta^2\big)t\Big)
\]
是鞅. 取 $\theta=-\mu+\sqrt{\mu^2+2\lambda}>0$, 它是方程 $\theta\mu+\frac{1}{2}\theta^2=\lambda$ 的正根, 于是 $N_t=e^{\theta X_t-\lambda t}$ 是鞅. 停止过程满足 $\E N_{t\wedge\tau^\mu_a}=1$. 由于在 $[0,\tau^\mu_a]$ 上 $X\le a$, $0<N_{t\wedge\tau^\mu_a}\le e^{\theta a}$. 令 $t\to\infty$: 在 $\{\tau^\mu_a<\infty\}$ 上, $N_{t\wedge\tau^\mu_a}\to e^{\theta a-\lambda\tau^\mu_a}$; 在 $\{\tau^\mu_a=\infty\}$ 上, $N_t\le e^{\theta a-\lambda t}\to0$. 由有界收敛定理,
\[
\E\Big[e^{\theta a-\lambda\tau^\mu_a}\1_{\{\tau^\mu_a<\infty\}}\Big]=1,
\]
即所述公式. 令 $\lambda\downarrow0$, 由单调收敛定理, 左端趋于 $\P(\tau^\mu_a<\infty)$, 右端趋于 $e^{a(\mu-\abs{\mu})}$, 它在 $\mu\ge0$ 时为 $1$, 在 $\mu<0$ 时为 $e^{2\mu a}=e^{-2\abs{\mu}a}$.

设 $\mu>0$, 则 $\tau^\mu_a<\infty$ $\as$ 由于 $\frac{1-e^{-\lambda\tau}}{\lambda}$ 在 $\lambda\downarrow0$ 时单调上升到 $\tau$, 由单调收敛定理,
\[
\E\tau^\mu_a=\lim_{\lambda\downarrow0}\frac{1-e^{a(\mu-\sqrt{\mu^2+2\lambda})}}{\lambda}=a\cdot\frac{\dd}{\dd\lambda}\sqrt{\mu^2+2\lambda}\,\Big|_{\lambda=0}=\frac{a}{\mu},
\]
其中第二个等号是导数的定义 (函数 $\lambda\mapsto1-e^{a(\mu-\sqrt{\mu^2+2\lambda})}$ 在 $0$ 处取值为 $0$, 导数为 $a\cdot\frac{1}{\mu}$).
\end{proof}

$\mu=0$ 时公式退化为定理~\ref{thm:hitting-time} (3). 对 $\mu<0$, 结论可以改写为
\[
\P\Big(\sup_{t\ge0}X_t\ge a\Big)=e^{-2\abs{\mu}a},\qquad a>0,
\]
即带负漂移的布朗运动的全局最大值服从参数为 $2\abs{\mu}$ 的指数分布 (图~\ref{fig:drift}). 对一般的 $\sigma$, 参数为 $2\abs{\mu}/\sigma^2$. 在保险的风险理论中, 若把保险公司的盈余近似为带正漂移的布朗运动 $u+\mu t+\sigma B_t$, 则破产概率 (盈余跌到零以下的概率) 为 $e^{-2\mu u/\sigma^2}$: 它随初始资本 $u$ 指数下降.

\begin{figure}[htbp]
\centering
\begin{tikzpicture}
\begin{axis}[SAaxis, axis x line=bottom, axis y line=left, width=12cm, height=5.8cm, xmin=0, xmax=8.3, ymin=-6.6, ymax=1.7,
  xtick={0,2,4,6,8}, ytick={-6,-4,-2,0,1}, xlabel={$t$}, clip=false]
  \addplot[SADark!25] coordinates {(0,0) (8,0)};
  \addplot[SADark!60,dashed] coordinates {(0,0) (8,-4)};
  \node[SADark!70,font=\small,anchor=north] at (axis cs:7.4,-3.8) {$\mu t$};
  \addplot[SABlue,thin] coordinates {(0.000,0.000) (0.020,-0.197) (0.040,-0.291) (0.060,-0.323) (0.080,-0.008) (0.100,0.118) (0.120,0.048) (0.140,0.180) (0.160,0.165) (0.180,-0.144) (0.200,-0.026) (0.220,-0.170) (0.240,-0.153) (0.260,-0.428) (0.280,-0.594) (0.300,-0.352) (0.320,-0.530) (0.340,-0.601) (0.360,-0.356) (0.380,-0.259) (0.400,-0.176) (0.420,-0.270) (0.440,-0.293) (0.460,-0.401) (0.480,-0.474) (0.500,-0.647) (0.520,-0.604) (0.540,-0.583) (0.560,-0.522) (0.580,-0.447) (0.600,-0.514) (0.620,-0.781) (0.640,-1.015) (0.660,-1.104) (0.680,-1.224) (0.700,-0.994) (0.720,-0.720) (0.740,-0.698) (0.760,-0.637) (0.780,-0.629) (0.800,-0.847) (0.820,-0.695) (0.840,-0.607) (0.860,-0.693) (0.880,-0.741) (0.900,-0.720) (0.920,-0.885) (0.940,-0.673) (0.960,-0.837) (0.980,-1.042) (1.000,-0.930) (1.020,-0.781) (1.040,-0.885) (1.060,-0.893) (1.080,-0.774) (1.100,-0.731) (1.120,-0.709) (1.140,-0.778) (1.160,-0.909) (1.180,-1.143) (1.200,-1.036) (1.220,-1.079) (1.240,-1.270) (1.260,-1.347) (1.280,-1.252) (1.300,-1.276) (1.320,-1.367) (1.340,-1.408) (1.360,-1.518) (1.380,-1.665) (1.400,-1.531) (1.420,-1.447) (1.440,-1.354) (1.460,-1.106) (1.480,-1.196) (1.500,-1.268) (1.520,-1.116) (1.540,-1.253) (1.560,-1.460) (1.580,-1.647) (1.600,-1.564) (1.620,-1.449) (1.640,-1.581) (1.660,-1.845) (1.680,-1.911) (1.700,-1.875) (1.720,-1.712) (1.740,-1.680) (1.760,-1.677) (1.780,-1.300) (1.800,-1.118) (1.820,-1.253) (1.840,-1.347) (1.860,-1.381) (1.880,-1.210) (1.900,-1.019) (1.920,-0.839) (1.940,-0.847) (1.960,-0.962) (1.980,-1.078) (2.000,-0.852) (2.020,-0.817) (2.040,-0.504) (2.060,-0.702) (2.080,-0.600) (2.100,-0.312) (2.120,-0.418) (2.140,-0.625) (2.160,-0.533) (2.180,-0.598) (2.200,-0.635) (2.220,-0.457) (2.240,-0.543) (2.260,-0.427) (2.280,-0.300) (2.300,-0.175) (2.320,-0.156) (2.340,-0.010) (2.360,-0.099) (2.380,0.003) (2.400,0.074) (2.420,0.023) (2.440,0.070) (2.460,0.213) (2.480,0.406) (2.500,0.516) (2.520,0.712) (2.540,0.615) (2.560,0.499) (2.580,0.484) (2.600,0.473) (2.620,0.453) (2.640,0.629) (2.660,0.274) (2.680,0.103) (2.700,0.243) (2.720,0.221) (2.740,0.056) (2.760,0.311) (2.780,-0.083) (2.800,-0.279) (2.820,-0.186) (2.840,-0.264) (2.860,-0.180) (2.880,0.034) (2.900,0.091) (2.920,0.068) (2.940,0.126) (2.960,-0.122) (2.980,-0.025) (3.000,-0.188) (3.020,-0.310) (3.040,-0.286) (3.060,-0.187) (3.080,-0.261) (3.100,-0.327) (3.120,-0.314) (3.140,-0.166) (3.160,-0.319) (3.180,-0.518) (3.200,-0.690) (3.220,-0.606) (3.240,-0.407) (3.260,-0.550) (3.280,-0.602) (3.300,-0.558) (3.320,-0.670) (3.340,-0.616) (3.360,-0.596) (3.380,-0.782) (3.400,-0.836) (3.420,-0.874) (3.440,-1.037) (3.460,-0.546) (3.480,-0.687) (3.500,-0.506) (3.520,-0.639) (3.540,-0.585) (3.560,-0.401) (3.580,-0.202) (3.600,-0.533) (3.620,-0.565) (3.640,-0.640) (3.660,-0.634) (3.680,-0.532) (3.700,-0.330) (3.720,-0.568) (3.740,-0.589) (3.760,-0.495) (3.780,-0.556) (3.800,-0.879) (3.820,-0.825) (3.840,-0.962) (3.860,-0.975) (3.880,-0.720) (3.900,-0.717) (3.920,-0.627) (3.940,-0.844) (3.960,-0.828) (3.980,-0.854) (4.000,-0.684) (4.020,-0.487) (4.040,-0.718) (4.060,-0.544) (4.080,-0.808) (4.100,-0.837) (4.120,-0.778) (4.140,-0.969) (4.160,-1.052) (4.180,-1.132) (4.200,-1.053) (4.220,-0.969) (4.240,-1.042) (4.260,-1.182) (4.280,-1.306) (4.300,-1.378) (4.320,-1.286) (4.340,-0.950) (4.360,-0.951) (4.380,-0.955) (4.400,-0.695) (4.420,-0.849) (4.440,-0.919) (4.460,-1.221) (4.480,-1.335) (4.500,-1.169) (4.520,-1.308) (4.540,-1.321) (4.560,-1.271) (4.580,-1.247) (4.600,-1.244) (4.620,-1.187) (4.640,-1.126) (4.660,-0.937) (4.680,-0.913) (4.700,-1.015) (4.720,-0.898) (4.740,-0.910) (4.760,-0.840) (4.780,-0.848) (4.800,-0.968) (4.820,-0.832) (4.840,-0.899) (4.860,-0.803) (4.880,-0.763) (4.900,-0.625) (4.920,-0.922) (4.940,-0.960) (4.960,-0.898) (4.980,-1.068) (5.000,-1.164) (5.020,-1.022) (5.040,-0.934) (5.060,-1.002) (5.080,-0.998) (5.100,-0.827) (5.120,-0.897) (5.140,-1.197) (5.160,-1.093) (5.180,-1.034) (5.200,-1.082) (5.220,-1.102) (5.240,-1.314) (5.260,-1.409) (5.280,-1.488) (5.300,-1.629) (5.320,-1.497) (5.340,-1.469) (5.360,-1.587) (5.380,-1.782) (5.400,-1.849) (5.420,-1.881) (5.440,-2.115) (5.460,-2.016) (5.480,-2.028) (5.500,-1.917) (5.520,-1.948) (5.540,-2.029) (5.560,-2.278) (5.580,-2.190) (5.600,-2.241) (5.620,-2.328) (5.640,-2.216) (5.660,-2.102) (5.680,-1.963) (5.700,-2.087) (5.720,-2.185) (5.740,-2.100) (5.760,-2.078) (5.780,-2.015) (5.800,-2.184) (5.820,-2.205) (5.840,-2.346) (5.860,-2.198) (5.880,-2.081) (5.900,-1.922) (5.920,-1.846) (5.940,-1.982) (5.960,-2.039) (5.980,-1.929) (6.000,-1.899) (6.020,-2.160) (6.040,-2.234) (6.060,-2.445) (6.080,-2.352) (6.100,-2.393) (6.120,-2.323) (6.140,-2.420) (6.160,-2.533) (6.180,-2.364) (6.200,-2.336) (6.220,-2.478) (6.240,-2.508) (6.260,-2.657) (6.280,-2.781) (6.300,-2.995) (6.320,-3.060) (6.340,-3.045) (6.360,-2.952) (6.380,-2.939) (6.400,-3.012) (6.420,-3.066) (6.440,-2.823) (6.460,-2.753) (6.480,-2.429) (6.500,-2.208) (6.520,-2.533) (6.540,-2.639) (6.560,-2.569) (6.580,-2.560) (6.600,-2.752) (6.620,-2.769) (6.640,-2.737) (6.660,-2.715) (6.680,-2.745) (6.700,-2.726) (6.720,-2.534) (6.740,-2.403) (6.760,-2.534) (6.780,-2.445) (6.800,-2.532) (6.820,-2.633) (6.840,-2.575) (6.860,-2.414) (6.880,-2.476) (6.900,-2.519) (6.920,-2.403) (6.940,-2.293) (6.960,-2.393) (6.980,-2.366) (7.000,-2.353) (7.020,-2.530) (7.040,-2.543) (7.060,-2.783) (7.080,-2.839) (7.100,-3.042) (7.120,-3.146) (7.140,-3.121) (7.160,-2.995) (7.180,-2.980) (7.200,-3.167) (7.220,-3.060) (7.240,-3.075) (7.260,-3.112) (7.280,-3.097) (7.300,-3.143) (7.320,-3.090) (7.340,-3.113) (7.360,-3.383) (7.380,-3.386) (7.400,-3.549) (7.420,-3.542) (7.440,-3.748) (7.460,-3.784) (7.480,-3.608) (7.500,-3.815) (7.520,-3.733) (7.540,-3.725) (7.560,-3.895) (7.580,-3.951) (7.600,-3.942) (7.620,-3.705) (7.640,-3.979) (7.660,-3.897) (7.680,-3.892) (7.700,-3.889) (7.720,-3.645) (7.740,-3.756) (7.760,-3.779) (7.780,-4.045) (7.800,-3.959) (7.820,-4.094) (7.840,-4.109) (7.860,-3.888) (7.880,-3.945) (7.900,-4.045) (7.920,-3.913) (7.940,-3.641) (7.960,-3.524) (7.980,-3.652) (8.000,-3.759)};
  \addplot[SARed,thin] coordinates {(0.000,0.000) (0.020,-0.184) (0.040,-0.149) (0.060,-0.263) (0.080,-0.028) (0.100,0.057) (0.120,0.066) (0.140,0.115) (0.160,0.210) (0.180,-0.202) (0.200,-0.280) (0.220,-0.239) (0.240,-0.312) (0.260,-0.143) (0.280,-0.210) (0.300,-0.242) (0.320,-0.059) (0.340,-0.054) (0.360,-0.096) (0.380,-0.040) (0.400,0.012) (0.420,0.232) (0.440,0.064) (0.460,0.050) (0.480,-0.092) (0.500,-0.007) (0.520,-0.249) (0.540,-0.314) (0.560,-0.499) (0.580,-0.503) (0.600,-0.447) (0.620,-0.438) (0.640,-0.335) (0.660,-0.360) (0.680,-0.330) (0.700,-0.210) (0.720,-0.559) (0.740,-0.794) (0.760,-0.709) (0.780,-0.527) (0.800,-0.709) (0.820,-0.589) (0.840,-0.664) (0.860,-0.553) (0.880,-0.512) (0.900,-0.769) (0.920,-0.790) (0.940,-0.628) (0.960,-0.481) (0.980,-0.924) (1.000,-0.670) (1.020,-0.641) (1.040,-0.553) (1.060,-0.592) (1.080,-0.610) (1.100,-0.669) (1.120,-0.731) (1.140,-0.866) (1.160,-0.734) (1.180,-0.826) (1.200,-0.891) (1.220,-0.834) (1.240,-0.690) (1.260,-0.699) (1.280,-0.631) (1.300,-0.505) (1.320,-0.510) (1.340,-0.682) (1.360,-0.727) (1.380,-0.972) (1.400,-1.044) (1.420,-0.849) (1.440,-1.024) (1.460,-0.847) (1.480,-0.857) (1.500,-0.953) (1.520,-1.048) (1.540,-1.010) (1.560,-0.812) (1.580,-0.776) (1.600,-0.783) (1.620,-0.732) (1.640,-0.772) (1.660,-0.513) (1.680,-0.603) (1.700,-0.548) (1.720,-0.719) (1.740,-0.669) (1.760,-0.719) (1.780,-0.575) (1.800,-0.463) (1.820,-0.407) (1.840,-0.560) (1.860,-0.523) (1.880,-0.545) (1.900,-0.598) (1.920,-0.729) (1.940,-0.882) (1.960,-1.062) (1.980,-0.835) (2.000,-1.033) (2.020,-1.037) (2.040,-1.147) (2.060,-1.284) (2.080,-1.467) (2.100,-1.282) (2.120,-1.313) (2.140,-1.374) (2.160,-1.538) (2.180,-1.591) (2.200,-1.728) (2.220,-1.751) (2.240,-2.023) (2.260,-1.959) (2.280,-2.124) (2.300,-2.155) (2.320,-2.000) (2.340,-2.039) (2.360,-2.248) (2.380,-2.017) (2.400,-1.972) (2.420,-1.793) (2.440,-1.685) (2.460,-1.520) (2.480,-1.385) (2.500,-1.444) (2.520,-1.608) (2.540,-1.527) (2.560,-1.441) (2.580,-1.551) (2.600,-1.677) (2.620,-1.464) (2.640,-1.348) (2.660,-1.264) (2.680,-1.386) (2.700,-1.186) (2.720,-1.229) (2.740,-1.140) (2.760,-1.220) (2.780,-1.220) (2.800,-1.309) (2.820,-1.222) (2.840,-1.202) (2.860,-1.187) (2.880,-1.185) (2.900,-1.001) (2.920,-1.017) (2.940,-1.198) (2.960,-1.126) (2.980,-1.302) (3.000,-1.327) (3.020,-1.453) (3.040,-1.351) (3.060,-1.343) (3.080,-1.356) (3.100,-1.338) (3.120,-1.355) (3.140,-1.529) (3.160,-1.678) (3.180,-1.663) (3.200,-1.483) (3.220,-1.348) (3.240,-1.526) (3.260,-1.243) (3.280,-1.261) (3.300,-1.176) (3.320,-1.047) (3.340,-0.960) (3.360,-0.677) (3.380,-0.715) (3.400,-0.712) (3.420,-0.730) (3.440,-0.743) (3.460,-0.802) (3.480,-0.692) (3.500,-0.693) (3.520,-0.550) (3.540,-0.815) (3.560,-0.756) (3.580,-0.831) (3.600,-1.152) (3.620,-1.377) (3.640,-1.391) (3.660,-1.437) (3.680,-1.533) (3.700,-1.322) (3.720,-1.481) (3.740,-1.437) (3.760,-1.479) (3.780,-1.373) (3.800,-1.292) (3.820,-1.490) (3.840,-1.328) (3.860,-1.422) (3.880,-1.399) (3.900,-1.258) (3.920,-1.135) (3.940,-1.234) (3.960,-1.474) (3.980,-1.530) (4.000,-1.375) (4.020,-1.215) (4.040,-1.187) (4.060,-1.262) (4.080,-1.229) (4.100,-1.232) (4.120,-1.142) (4.140,-1.145) (4.160,-1.017) (4.180,-0.995) (4.200,-0.983) (4.220,-1.101) (4.240,-0.798) (4.260,-0.892) (4.280,-0.730) (4.300,-0.616) (4.320,-0.406) (4.340,-0.309) (4.360,-0.205) (4.380,-0.134) (4.400,-0.222) (4.420,-0.093) (4.440,-0.227) (4.460,-0.359) (4.480,-0.450) (4.500,-0.688) (4.520,-0.745) (4.540,-0.723) (4.560,-0.646) (4.580,-0.570) (4.600,-0.366) (4.620,-0.303) (4.640,-0.532) (4.660,-0.269) (4.680,-0.516) (4.700,-0.831) (4.720,-0.774) (4.740,-0.964) (4.760,-0.951) (4.780,-0.838) (4.800,-1.045) (4.820,-0.662) (4.840,-0.654) (4.860,-0.385) (4.880,-0.550) (4.900,-0.567) (4.920,-0.644) (4.940,-0.734) (4.960,-0.773) (4.980,-0.739) (5.000,-0.711) (5.020,-0.677) (5.040,-0.627) (5.060,-0.661) (5.080,-0.658) (5.100,-0.380) (5.120,-0.330) (5.140,-0.418) (5.160,-0.244) (5.180,-0.393) (5.200,-0.701) (5.220,-0.843) (5.240,-0.959) (5.260,-1.091) (5.280,-1.221) (5.300,-1.412) (5.320,-1.362) (5.340,-1.544) (5.360,-1.430) (5.380,-1.391) (5.400,-1.161) (5.420,-1.336) (5.440,-1.375) (5.460,-1.439) (5.480,-1.220) (5.500,-1.369) (5.520,-1.430) (5.540,-1.384) (5.560,-1.487) (5.580,-1.500) (5.600,-1.560) (5.620,-1.410) (5.640,-1.459) (5.660,-1.412) (5.680,-1.584) (5.700,-1.503) (5.720,-1.626) (5.740,-1.518) (5.760,-1.714) (5.780,-1.637) (5.800,-1.967) (5.820,-2.351) (5.840,-2.124) (5.860,-2.080) (5.880,-1.846) (5.900,-1.838) (5.920,-1.942) (5.940,-2.013) (5.960,-1.992) (5.980,-2.146) (6.000,-2.133) (6.020,-2.045) (6.040,-1.999) (6.060,-1.945) (6.080,-2.004) (6.100,-2.221) (6.120,-2.069) (6.140,-2.121) (6.160,-2.029) (6.180,-2.010) (6.200,-2.344) (6.220,-2.237) (6.240,-2.271) (6.260,-2.434) (6.280,-2.415) (6.300,-2.709) (6.320,-2.862) (6.340,-2.832) (6.360,-2.846) (6.380,-2.738) (6.400,-2.902) (6.420,-2.773) (6.440,-3.030) (6.460,-2.961) (6.480,-2.906) (6.500,-2.847) (6.520,-2.795) (6.540,-2.637) (6.560,-2.800) (6.580,-2.673) (6.600,-2.624) (6.620,-2.582) (6.640,-2.358) (6.660,-2.317) (6.680,-2.426) (6.700,-2.712) (6.720,-2.871) (6.740,-2.751) (6.760,-2.749) (6.780,-2.650) (6.800,-2.510) (6.820,-2.311) (6.840,-2.314) (6.860,-2.302) (6.880,-2.371) (6.900,-2.749) (6.920,-2.590) (6.940,-2.603) (6.960,-2.511) (6.980,-2.507) (7.000,-2.455) (7.020,-2.506) (7.040,-2.586) (7.060,-2.591) (7.080,-2.632) (7.100,-2.628) (7.120,-2.731) (7.140,-2.513) (7.160,-2.559) (7.180,-2.733) (7.200,-2.654) (7.220,-2.573) (7.240,-2.444) (7.260,-2.408) (7.280,-2.238) (7.300,-2.177) (7.320,-2.277) (7.340,-2.200) (7.360,-2.361) (7.380,-2.292) (7.400,-2.178) (7.420,-2.121) (7.440,-2.029) (7.460,-2.052) (7.480,-2.198) (7.500,-2.338) (7.520,-2.605) (7.540,-2.585) (7.560,-2.669) (7.580,-2.784) (7.600,-3.085) (7.620,-3.328) (7.640,-3.442) (7.660,-3.430) (7.680,-3.325) (7.700,-3.352) (7.720,-3.318) (7.740,-3.280) (7.760,-3.287) (7.780,-3.300) (7.800,-3.267) (7.820,-3.273) (7.840,-3.546) (7.860,-3.799) (7.880,-3.902) (7.900,-3.964) (7.920,-4.029) (7.940,-4.074) (7.960,-3.761) (7.980,-3.789) (8.000,-3.766)};
  \addplot[SAGreenText,thin] coordinates {(0.000,0.000) (0.020,-0.117) (0.040,-0.110) (0.060,-0.038) (0.080,-0.199) (0.100,-0.440) (0.120,-0.137) (0.140,-0.216) (0.160,-0.150) (0.180,-0.158) (0.200,-0.239) (0.220,-0.218) (0.240,-0.482) (0.260,-0.415) (0.280,-0.324) (0.300,-0.333) (0.320,0.015) (0.340,0.241) (0.360,0.169) (0.380,0.098) (0.400,0.363) (0.420,0.503) (0.440,0.602) (0.460,0.722) (0.480,0.642) (0.500,0.613) (0.520,0.486) (0.540,0.714) (0.560,0.772) (0.580,0.693) (0.600,0.721) (0.620,0.763) (0.640,0.828) (0.660,0.487) (0.680,0.474) (0.700,0.628) (0.720,0.459) (0.740,0.142) (0.760,0.339) (0.780,0.227) (0.800,0.132) (0.820,0.164) (0.840,0.184) (0.860,0.073) (0.880,0.033) (0.900,0.054) (0.920,0.154) (0.940,0.313) (0.960,0.189) (0.980,0.275) (1.000,0.394) (1.020,0.339) (1.040,0.252) (1.060,-0.042) (1.080,0.010) (1.100,-0.050) (1.120,-0.223) (1.140,-0.441) (1.160,-0.391) (1.180,-0.574) (1.200,-0.598) (1.220,-0.455) (1.240,-0.806) (1.260,-0.814) (1.280,-0.756) (1.300,-0.660) (1.320,-0.584) (1.340,-0.467) (1.360,-0.322) (1.380,-0.125) (1.400,-0.054) (1.420,-0.005) (1.440,-0.058) (1.460,-0.189) (1.480,-0.092) (1.500,-0.106) (1.520,-0.085) (1.540,-0.222) (1.560,-0.196) (1.580,-0.313) (1.600,-0.356) (1.620,-0.128) (1.640,-0.095) (1.660,0.005) (1.680,0.216) (1.700,0.214) (1.720,0.469) (1.740,0.299) (1.760,0.042) (1.780,0.230) (1.800,0.220) (1.820,0.365) (1.840,0.449) (1.860,0.571) (1.880,0.583) (1.900,0.471) (1.920,0.646) (1.940,0.395) (1.960,0.233) (1.980,0.326) (2.000,0.408) (2.020,0.457) (2.040,0.838) (2.060,0.786) (2.080,1.032) (2.100,1.066) (2.120,1.124) (2.140,1.008) (2.160,1.035) (2.180,1.153) (2.200,1.256) (2.220,1.046) (2.240,0.889) (2.260,1.020) (2.280,0.947) (2.300,0.904) (2.320,0.634) (2.340,0.220) (2.360,0.426) (2.380,0.052) (2.400,0.067) (2.420,0.062) (2.440,-0.099) (2.460,-0.185) (2.480,-0.307) (2.500,-0.545) (2.520,-0.442) (2.540,-0.437) (2.560,-0.290) (2.580,-0.358) (2.600,-0.572) (2.620,-0.312) (2.640,-0.278) (2.660,-0.420) (2.680,-0.477) (2.700,-0.604) (2.720,-0.695) (2.740,-0.482) (2.760,-0.445) (2.780,-0.626) (2.800,-0.473) (2.820,-0.364) (2.840,-0.530) (2.860,-0.541) (2.880,-0.693) (2.900,-0.730) (2.920,-0.723) (2.940,-0.617) (2.960,-0.762) (2.980,-0.819) (3.000,-0.718) (3.020,-0.351) (3.040,-0.167) (3.060,-0.246) (3.080,-0.124) (3.100,-0.274) (3.120,-0.346) (3.140,-0.448) (3.160,-0.560) (3.180,-0.240) (3.200,-0.339) (3.220,-0.256) (3.240,-0.194) (3.260,-0.323) (3.280,-0.566) (3.300,-0.345) (3.320,-0.523) (3.340,-0.575) (3.360,-0.732) (3.380,-0.609) (3.400,-0.792) (3.420,-0.636) (3.440,-0.708) (3.460,-0.796) (3.480,-0.896) (3.500,-0.789) (3.520,-0.681) (3.540,-0.657) (3.560,-0.427) (3.580,-0.329) (3.600,-0.548) (3.620,-0.841) (3.640,-0.598) (3.660,-0.803) (3.680,-0.912) (3.700,-0.851) (3.720,-0.844) (3.740,-0.941) (3.760,-0.776) (3.780,-0.939) (3.800,-1.130) (3.820,-1.156) (3.840,-0.879) (3.860,-0.998) (3.880,-0.845) (3.900,-0.980) (3.920,-0.865) (3.940,-0.858) (3.960,-0.739) (3.980,-0.575) (4.000,-0.325) (4.020,-0.285) (4.040,-0.444) (4.060,-0.629) (4.080,-0.649) (4.100,-0.863) (4.120,-0.726) (4.140,-0.532) (4.160,-0.418) (4.180,-0.561) (4.200,-0.757) (4.220,-0.929) (4.240,-0.884) (4.260,-0.959) (4.280,-0.809) (4.300,-0.844) (4.320,-0.922) (4.340,-0.868) (4.360,-0.686) (4.380,-0.635) (4.400,-0.804) (4.420,-0.828) (4.440,-1.160) (4.460,-1.000) (4.480,-1.025) (4.500,-0.899) (4.520,-1.042) (4.540,-1.225) (4.560,-1.213) (4.580,-1.330) (4.600,-1.417) (4.620,-1.390) (4.640,-1.364) (4.660,-1.476) (4.680,-1.723) (4.700,-1.640) (4.720,-1.479) (4.740,-1.401) (4.760,-1.672) (4.780,-1.657) (4.800,-1.601) (4.820,-1.557) (4.840,-1.475) (4.860,-1.195) (4.880,-1.118) (4.900,-1.274) (4.920,-1.275) (4.940,-1.297) (4.960,-1.369) (4.980,-1.206) (5.000,-1.108) (5.020,-0.993) (5.040,-0.754) (5.060,-1.015) (5.080,-1.162) (5.100,-0.952) (5.120,-0.811) (5.140,-0.970) (5.160,-1.069) (5.180,-1.288) (5.200,-1.378) (5.220,-1.401) (5.240,-1.306) (5.260,-1.128) (5.280,-0.794) (5.300,-0.748) (5.320,-0.802) (5.340,-0.458) (5.360,-0.670) (5.380,-0.909) (5.400,-0.707) (5.420,-0.869) (5.440,-1.168) (5.460,-1.420) (5.480,-1.465) (5.500,-1.506) (5.520,-1.303) (5.540,-1.456) (5.560,-1.461) (5.580,-1.378) (5.600,-1.697) (5.620,-1.572) (5.640,-1.512) (5.660,-1.585) (5.680,-1.281) (5.700,-1.151) (5.720,-1.063) (5.740,-1.086) (5.760,-1.179) (5.780,-1.149) (5.800,-1.295) (5.820,-1.314) (5.840,-1.412) (5.860,-1.595) (5.880,-1.442) (5.900,-1.567) (5.920,-1.548) (5.940,-1.736) (5.960,-1.935) (5.980,-1.928) (6.000,-1.955) (6.020,-1.765) (6.040,-1.548) (6.060,-1.696) (6.080,-1.749) (6.100,-1.586) (6.120,-1.558) (6.140,-1.428) (6.160,-1.578) (6.180,-1.245) (6.200,-1.001) (6.220,-1.089) (6.240,-1.205) (6.260,-1.086) (6.280,-0.975) (6.300,-0.821) (6.320,-0.715) (6.340,-0.778) (6.360,-0.836) (6.380,-0.903) (6.400,-0.781) (6.420,-0.810) (6.440,-0.876) (6.460,-0.720) (6.480,-0.460) (6.500,-0.064) (6.520,-0.242) (6.540,-0.350) (6.560,-0.322) (6.580,-0.268) (6.600,-0.239) (6.620,-0.431) (6.640,-0.543) (6.660,-0.396) (6.680,-0.399) (6.700,-0.404) (6.720,-0.472) (6.740,-0.366) (6.760,-0.358) (6.780,-0.495) (6.800,-0.648) (6.820,-0.690) (6.840,-0.709) (6.860,-0.539) (6.880,-0.494) (6.900,-0.543) (6.920,-0.668) (6.940,-0.919) (6.960,-1.035) (6.980,-1.277) (7.000,-1.584) (7.020,-1.633) (7.040,-1.672) (7.060,-1.377) (7.080,-1.424) (7.100,-1.472) (7.120,-1.615) (7.140,-1.541) (7.160,-1.343) (7.180,-1.346) (7.200,-1.540) (7.220,-1.505) (7.240,-1.451) (7.260,-1.359) (7.280,-1.319) (7.300,-1.401) (7.320,-1.625) (7.340,-1.622) (7.360,-1.851) (7.380,-1.818) (7.400,-1.788) (7.420,-1.725) (7.440,-1.710) (7.460,-1.646) (7.480,-1.725) (7.500,-1.534) (7.520,-1.384) (7.540,-1.162) (7.560,-1.238) (7.580,-1.118) (7.600,-1.208) (7.620,-1.202) (7.640,-1.310) (7.660,-1.287) (7.680,-1.573) (7.700,-1.307) (7.720,-1.378) (7.740,-1.359) (7.760,-1.098) (7.780,-1.183) (7.800,-1.367) (7.820,-1.284) (7.840,-1.119) (7.860,-1.305) (7.880,-1.411) (7.900,-1.323) (7.920,-1.343) (7.940,-1.374) (7.960,-1.264) (7.980,-1.529) (8.000,-1.268)};
  \addplot[SAYellowText,thin] coordinates {(0.000,0.000) (0.020,0.041) (0.040,0.220) (0.060,-0.199) (0.080,-0.044) (0.100,-0.024) (0.120,0.139) (0.140,0.419) (0.160,0.492) (0.180,0.291) (0.200,0.170) (0.220,0.380) (0.240,0.387) (0.260,0.568) (0.280,0.622) (0.300,0.648) (0.320,0.565) (0.340,0.796) (0.360,0.686) (0.380,0.538) (0.400,0.721) (0.420,0.751) (0.440,0.751) (0.460,0.693) (0.480,0.479) (0.500,0.491) (0.520,0.324) (0.540,0.318) (0.560,0.461) (0.580,0.547) (0.600,0.484) (0.620,0.453) (0.640,0.515) (0.660,0.664) (0.680,0.636) (0.700,0.671) (0.720,0.641) (0.740,0.853) (0.760,0.767) (0.780,0.612) (0.800,0.858) (0.820,0.983) (0.840,0.830) (0.860,0.770) (0.880,0.532) (0.900,0.452) (0.920,0.148) (0.940,0.259) (0.960,0.260) (0.980,0.292) (1.000,0.349) (1.020,0.433) (1.040,0.443) (1.060,0.494) (1.080,0.394) (1.100,0.557) (1.120,0.304) (1.140,0.495) (1.160,0.608) (1.180,0.420) (1.200,0.278) (1.220,0.354) (1.240,0.466) (1.260,0.499) (1.280,0.353) (1.300,0.579) (1.320,0.440) (1.340,0.169) (1.360,0.072) (1.380,0.034) (1.400,-0.191) (1.420,-0.107) (1.440,-0.166) (1.460,-0.159) (1.480,-0.109) (1.500,-0.090) (1.520,0.065) (1.540,0.126) (1.560,-0.039) (1.580,0.040) (1.600,-0.150) (1.620,-0.199) (1.640,-0.069) (1.660,-0.098) (1.680,-0.071) (1.700,-0.183) (1.720,-0.447) (1.740,-0.491) (1.760,-0.388) (1.780,-0.204) (1.800,-0.111) (1.820,-0.114) (1.840,-0.208) (1.860,-0.256) (1.880,-0.565) (1.900,-0.636) (1.920,-0.548) (1.940,-0.229) (1.960,-0.469) (1.980,-0.412) (2.000,-0.601) (2.020,-0.612) (2.040,-0.742) (2.060,-0.656) (2.080,-0.421) (2.100,-0.278) (2.120,-0.537) (2.140,-0.373) (2.160,-0.586) (2.180,-0.593) (2.200,-0.725) (2.220,-0.536) (2.240,-0.459) (2.260,-0.484) (2.280,-0.507) (2.300,-0.530) (2.320,-0.536) (2.340,-0.287) (2.360,-0.344) (2.380,-0.586) (2.400,-0.745) (2.420,-0.624) (2.440,-0.822) (2.460,-0.804) (2.480,-1.081) (2.500,-1.027) (2.520,-1.052) (2.540,-1.280) (2.560,-1.229) (2.580,-1.533) (2.600,-1.429) (2.620,-1.605) (2.640,-1.809) (2.660,-1.855) (2.680,-1.631) (2.700,-1.720) (2.720,-1.708) (2.740,-1.687) (2.760,-2.032) (2.780,-2.104) (2.800,-2.115) (2.820,-2.118) (2.840,-2.215) (2.860,-2.093) (2.880,-2.359) (2.900,-2.264) (2.920,-2.286) (2.940,-2.254) (2.960,-2.402) (2.980,-2.666) (3.000,-2.776) (3.020,-2.529) (3.040,-2.357) (3.060,-2.616) (3.080,-2.613) (3.100,-2.453) (3.120,-2.422) (3.140,-2.230) (3.160,-2.120) (3.180,-2.164) (3.200,-2.363) (3.220,-2.540) (3.240,-2.524) (3.260,-2.593) (3.280,-2.655) (3.300,-2.568) (3.320,-2.405) (3.340,-2.386) (3.360,-2.449) (3.380,-2.574) (3.400,-2.341) (3.420,-2.324) (3.440,-2.419) (3.460,-2.275) (3.480,-2.411) (3.500,-2.458) (3.520,-2.272) (3.540,-2.389) (3.560,-2.674) (3.580,-2.659) (3.600,-2.492) (3.620,-2.458) (3.640,-2.477) (3.660,-2.765) (3.680,-2.700) (3.700,-2.894) (3.720,-2.950) (3.740,-2.936) (3.760,-2.877) (3.780,-3.159) (3.800,-3.325) (3.820,-3.281) (3.840,-3.148) (3.860,-3.152) (3.880,-2.910) (3.900,-3.027) (3.920,-2.993) (3.940,-2.743) (3.960,-2.858) (3.980,-2.856) (4.000,-2.931) (4.020,-2.974) (4.040,-2.816) (4.060,-2.727) (4.080,-2.725) (4.100,-2.611) (4.120,-2.635) (4.140,-2.589) (4.160,-2.398) (4.180,-2.314) (4.200,-2.017) (4.220,-1.996) (4.240,-2.122) (4.260,-1.951) (4.280,-1.958) (4.300,-1.834) (4.320,-1.986) (4.340,-1.850) (4.360,-1.707) (4.380,-1.844) (4.400,-1.566) (4.420,-1.525) (4.440,-1.571) (4.460,-1.604) (4.480,-1.722) (4.500,-1.765) (4.520,-1.832) (4.540,-1.632) (4.560,-1.710) (4.580,-1.600) (4.600,-1.499) (4.620,-1.480) (4.640,-1.400) (4.660,-1.417) (4.680,-1.030) (4.700,-1.075) (4.720,-0.838) (4.740,-0.566) (4.760,-0.653) (4.780,-0.794) (4.800,-0.605) (4.820,-0.825) (4.840,-0.980) (4.860,-1.040) (4.880,-1.113) (4.900,-0.923) (4.920,-1.140) (4.940,-1.295) (4.960,-1.413) (4.980,-1.470) (5.000,-1.450) (5.020,-1.124) (5.040,-1.168) (5.060,-1.100) (5.080,-0.980) (5.100,-0.772) (5.120,-0.741) (5.140,-0.769) (5.160,-0.851) (5.180,-1.047) (5.200,-1.209) (5.220,-1.332) (5.240,-1.550) (5.260,-1.728) (5.280,-1.724) (5.300,-1.995) (5.320,-2.158) (5.340,-2.126) (5.360,-2.408) (5.380,-2.341) (5.400,-2.255) (5.420,-2.222) (5.440,-2.324) (5.460,-2.526) (5.480,-2.562) (5.500,-2.760) (5.520,-2.813) (5.540,-2.894) (5.560,-2.887) (5.580,-2.824) (5.600,-2.737) (5.620,-3.041) (5.640,-3.197) (5.660,-3.122) (5.680,-3.157) (5.700,-3.152) (5.720,-3.305) (5.740,-3.711) (5.760,-3.729) (5.780,-3.579) (5.800,-3.478) (5.820,-3.361) (5.840,-3.356) (5.860,-3.411) (5.880,-3.468) (5.900,-3.235) (5.920,-3.106) (5.940,-3.275) (5.960,-3.402) (5.980,-3.306) (6.000,-3.365) (6.020,-3.274) (6.040,-3.344) (6.060,-3.346) (6.080,-3.674) (6.100,-3.614) (6.120,-3.525) (6.140,-3.560) (6.160,-3.533) (6.180,-3.623) (6.200,-3.770) (6.220,-4.014) (6.240,-4.091) (6.260,-4.064) (6.280,-4.198) (6.300,-4.484) (6.320,-4.695) (6.340,-4.823) (6.360,-5.092) (6.380,-5.188) (6.400,-5.068) (6.420,-5.067) (6.440,-5.085) (6.460,-4.948) (6.480,-4.907) (6.500,-4.989) (6.520,-5.138) (6.540,-4.881) (6.560,-5.029) (6.580,-5.263) (6.600,-5.140) (6.620,-5.169) (6.640,-5.179) (6.660,-4.975) (6.680,-5.045) (6.700,-5.065) (6.720,-5.172) (6.740,-5.346) (6.760,-5.514) (6.780,-5.741) (6.800,-5.995) (6.820,-6.215) (6.840,-6.161) (6.860,-6.279) (6.880,-6.408) (6.900,-6.164) (6.920,-5.973) (6.940,-6.005) (6.960,-6.056) (6.980,-5.973) (7.000,-5.987) (7.020,-5.942) (7.040,-5.754) (7.060,-5.940) (7.080,-5.889) (7.100,-5.848) (7.120,-5.913) (7.140,-5.688) (7.160,-5.626) (7.180,-5.425) (7.200,-5.403) (7.220,-5.325) (7.240,-5.361) (7.260,-5.524) (7.280,-5.581) (7.300,-5.424) (7.320,-5.495) (7.340,-5.368) (7.360,-5.328) (7.380,-5.312) (7.400,-5.189) (7.420,-5.356) (7.440,-5.206) (7.460,-5.253) (7.480,-5.278) (7.500,-5.422) (7.520,-5.541) (7.540,-5.455) (7.560,-5.723) (7.580,-5.645) (7.600,-5.734) (7.620,-5.870) (7.640,-5.549) (7.660,-5.315) (7.680,-5.183) (7.700,-5.209) (7.720,-5.281) (7.740,-5.291) (7.760,-5.191) (7.780,-5.280) (7.800,-5.400) (7.820,-5.498) (7.840,-5.343) (7.860,-5.113) (7.880,-5.160) (7.900,-5.266) (7.920,-4.989) (7.940,-4.993) (7.960,-5.065) (7.980,-5.172) (8.000,-5.148)};
  \fill[SABlue] (axis cs:2.520,0.712) circle (2pt);
  \fill[SARed] (axis cs:0.150,0.251) circle (2pt);
  \fill[SAGreenText] (axis cs:2.200,1.256) circle (2pt);
  \fill[SAYellowText] (axis cs:0.810,0.990) circle (2pt);
\end{axis}
\end{tikzpicture}
\caption{带负漂移的布朗运动 $X_t=-\frac12t+B_t$ 的四条轨道. 圆点标出每条轨道的全局最大值, 它服从参数为 $2\abs{\mu}=1$ 的指数分布; 此后轨道被漂移 (虚线) 带向 $-\infty$, 不再回来}
\label{fig:drift}
\end{figure}

\paragraph{几何布朗运动.} 设 $s_0>0$, $\mu\in\R$, $\sigma>0$. 称
\[
S_t=s_0\exp\Big(\sigma B_t+\big(\mu-\frac{1}{2}\sigma^2\big)t\Big)
\]
为\textbf{几何布朗运动}. 它是金融中 Black-Scholes 模型的股价过程: 对数收益 $\log(S_t/S_s)$ 服从正态分布, 且与过去独立. 第~5~章将证明它是随机微分方程 $\dd S_t=\mu S_t\dd t+\sigma S_t\dd B_t$ 的解, 修正项 $-\frac{1}{2}\sigma^2$ 正是 It\^o 公式的产物.

\begin{proposition}[几何布朗运动的性质]\label{prop:gbm}
\begin{enumerate}
\item $\log S_t\sim N\big(\log s_0+(\mu-\frac{1}{2}\sigma^2)t,\ \sigma^2t\big)$, $\E S_t=s_0e^{\mu t}$, $\Var(S_t)=s_0^2e^{2\mu t}\big(e^{\sigma^2t}-1\big)$.
\item $e^{-\mu t}S_t$ 是鞅; 特别地, $\mu=0$ 时 $S$ 是鞅.
\item 对 $s<t$, $S_t/S_s$ 与 $\mathcal{F}_s$ 独立, 且与 $S_{t-s}/s_0$ 同分布.
\item $\frac{1}{t}\log S_t\to\mu-\frac{1}{2}\sigma^2$ $\as$ 特别地, 若 $\mu<\frac{1}{2}\sigma^2$, 则 $S_t\to0$ $\as$
\end{enumerate}
\end{proposition}

\begin{proof}
(1) 第一式由 $B_t\sim N(0,t)$ 得到. 由例~\ref{ex:normal-mgf}, $\E e^{\sigma B_t}=e^{\sigma^2t/2}$, $\E e^{2\sigma B_t}=e^{2\sigma^2t}$, 所以
\[
\E S_t=s_0e^{(\mu-\frac{1}{2}\sigma^2)t}e^{\frac{1}{2}\sigma^2t}=s_0e^{\mu t},\qquad\E S_t^2=s_0^2e^{(2\mu-\sigma^2)t}e^{2\sigma^2t}=s_0^2e^{2\mu t}e^{\sigma^2t},
\]
相减即得方差. (2) $e^{-\mu t}S_t=s_0\exp(\sigma B_t-\frac{1}{2}\sigma^2t)$ 是指数鞅 (定理~\ref{thm:bm-martingales} (3), $\theta=\sigma$). (3) $S_t/S_s=\exp\big(\sigma(B_t-B_s)+(\mu-\frac{1}{2}\sigma^2)(t-s)\big)$ 是增量 $B_t-B_s$ 的函数. (4) $\frac{1}{t}\log S_t=\frac{\log s_0}{t}+\sigma\frac{B_t}{t}+\mu-\frac{1}{2}\sigma^2$, 由布朗运动的强大数定律 (推论~\ref{cor:bm-slln}) 即得.
\end{proof}

\begin{figure}[htbp]
\centering
\begin{tikzpicture}
\begin{axis}[SAaxis, axis x line=bottom, axis y line=left, width=12cm, height=6cm, xmin=0, xmax=10.3, ymin=0.1, ymax=20, ymode=log,
  xtick={0,2,4,6,8,10}, ytick={0.1,0.5,1,2,5,10}, yticklabels={$0.1$,$0.5$,$1$,$2$,$5$,$10$}, xlabel={$t$}, ylabel={$S_t$},
  legend style={font=\small,draw=none,fill=none,at={(0.5,-0.2)},anchor=north}, legend columns=2, legend cell align=left]
  \addplot[SARed,thick,domain=0:10,samples=40] {exp(0.1*x)};
  \addplot[SADark,thick,dashed,domain=0:10,samples=40] {exp(-0.08*x)};
  \addplot[SABlue!80,thin,forget plot] coordinates {(0.000,1.000) (0.040,0.968) (0.080,1.093) (0.120,1.320) (0.160,1.221) (0.200,1.157) (0.240,1.034) (0.280,0.849) (0.320,0.855) (0.360,0.867) (0.400,0.890) (0.440,0.782) (0.480,0.662) (0.520,0.701) (0.560,0.628) (0.600,0.547) (0.640,0.525) (0.680,0.586) (0.720,0.741) (0.760,0.680) (0.800,0.692) (0.840,0.672) (0.880,0.699) (0.920,0.582) (0.960,0.631) (1.000,0.608) (1.040,0.619) (1.080,0.596) (1.120,0.631) (1.160,0.604) (1.200,0.643) (1.240,0.594) (1.280,0.673) (1.320,0.755) (1.360,0.711) (1.400,0.588) (1.440,0.601) (1.480,0.707) (1.520,0.812) (1.560,0.897) (1.600,0.904) (1.640,1.033) (1.680,1.222) (1.720,1.139) (1.760,0.964) (1.800,0.868) (1.840,0.867) (1.880,0.818) (1.920,0.707) (1.960,0.729) (2.000,0.766) (2.040,0.720) (2.080,0.721) (2.120,0.772) (2.160,0.834) (2.200,0.864) (2.240,0.980) (2.280,0.954) (2.320,1.126) (2.360,1.064) (2.400,1.044) (2.440,1.162) (2.480,1.090) (2.520,1.017) (2.560,1.155) (2.600,1.144) (2.640,1.003) (2.680,1.125) (2.720,1.315) (2.760,1.466) (2.800,1.658) (2.840,1.938) (2.880,2.085) (2.920,2.042) (2.960,1.999) (3.000,1.998) (3.040,1.791) (3.080,1.743) (3.120,1.790) (3.160,1.646) (3.200,1.995) (3.240,2.169) (3.280,2.569) (3.320,2.008) (3.360,2.097) (3.400,2.248) (3.440,2.365) (3.480,2.443) (3.520,1.809) (3.560,1.920) (3.600,2.699) (3.640,2.999) (3.680,2.960) (3.720,3.176) (3.760,3.886) (3.800,3.438) (3.840,3.360) (3.880,3.003) (3.920,3.622) (3.960,3.796) (4.000,3.452) (4.040,3.024) (4.080,2.910) (4.120,2.939) (4.160,2.780) (4.200,2.822) (4.240,1.917) (4.280,2.113) (4.320,2.399) (4.360,2.696) (4.400,2.720) (4.440,2.827) (4.480,2.432) (4.520,3.250) (4.560,3.373) (4.600,3.181) (4.640,3.030) (4.680,3.323) (4.720,3.271) (4.760,3.051) (4.800,3.252) (4.840,3.526) (4.880,4.943) (4.920,5.559) (4.960,6.335) (5.000,6.733) (5.040,7.219) (5.080,8.254) (5.120,6.713) (5.160,8.475) (5.200,8.629) (5.240,9.291) (5.280,9.110) (5.320,8.301) (5.360,8.092) (5.400,8.103) (5.440,7.537) (5.480,7.936) (5.520,7.970) (5.560,8.456) (5.600,9.540) (5.640,8.108) (5.680,10.904) (5.720,8.622) (5.760,9.582) (5.800,9.194) (5.840,9.278) (5.880,8.042) (5.920,8.533) (5.960,8.158) (6.000,10.544) (6.040,10.438) (6.080,10.557) (6.120,10.220) (6.160,8.249) (6.200,7.292) (6.240,6.322) (6.280,8.484) (6.320,7.160) (6.360,6.581) (6.400,6.508) (6.440,5.996) (6.480,5.194) (6.520,6.279) (6.560,5.919) (6.600,6.305) (6.640,6.780) (6.680,7.737) (6.720,9.078) (6.760,8.364) (6.800,11.022) (6.840,12.865) (6.880,11.888) (6.920,9.903) (6.960,10.435) (7.000,13.610) (7.040,12.718) (7.080,12.226) (7.120,10.970) (7.160,12.855) (7.200,13.033) (7.240,13.995) (7.280,10.618) (7.320,8.129) (7.360,9.691) (7.400,8.317) (7.440,7.339) (7.480,5.361) (7.520,4.892) (7.560,5.064) (7.600,5.251) (7.640,5.497) (7.680,4.851) (7.720,5.890) (7.760,5.747) (7.800,5.472) (7.840,5.647) (7.880,4.644) (7.920,5.109) (7.960,5.883) (8.000,5.534) (8.040,5.406) (8.080,5.559) (8.120,5.308) (8.160,5.133) (8.200,4.602) (8.240,4.303) (8.280,5.031) (8.320,5.013) (8.360,5.238) (8.400,5.842) (8.440,6.775) (8.480,5.567) (8.520,4.933) (8.560,4.026) (8.600,4.844) (8.640,5.678) (8.680,5.459) (8.720,6.446) (8.760,5.875) (8.800,7.174) (8.840,6.943) (8.880,5.485) (8.920,4.828) (8.960,5.453) (9.000,5.201) (9.040,4.825) (9.080,5.697) (9.120,5.797) (9.160,5.339) (9.200,4.114) (9.240,3.213) (9.280,2.804) (9.320,2.740) (9.360,2.616) (9.400,2.314) (9.440,2.191) (9.480,2.192) (9.520,2.145) (9.560,1.984) (9.600,2.133) (9.640,2.654) (9.680,2.591) (9.720,2.413) (9.760,2.211) (9.800,2.525) (9.840,2.032) (9.880,2.023) (9.920,1.657) (9.960,1.375) (10.000,1.209)};
  \addplot[SABlue!80,thin,forget plot] coordinates {(0.000,1.000) (0.040,1.178) (0.080,0.898) (0.120,0.854) (0.160,0.925) (0.200,0.938) (0.240,0.902) (0.280,1.040) (0.320,1.138) (0.360,1.232) (0.400,1.356) (0.440,1.240) (0.480,1.126) (0.520,1.257) (0.560,1.198) (0.600,1.414) (0.640,1.155) (0.680,0.952) (0.720,0.974) (0.760,1.019) (0.800,0.883) (0.840,1.097) (0.880,1.010) (0.920,1.216) (0.960,1.297) (1.000,1.209) (1.040,1.152) (1.080,1.282) (1.120,1.624) (1.160,2.016) (1.200,1.948) (1.240,2.325) (1.280,1.944) (1.320,1.973) (1.360,1.885) (1.400,1.586) (1.440,1.715) (1.480,1.473) (1.520,1.354) (1.560,1.428) (1.600,1.372) (1.640,1.261) (1.680,1.467) (1.720,1.272) (1.760,1.172) (1.800,1.086) (1.840,1.119) (1.880,1.189) (1.920,1.180) (1.960,1.074) (2.000,1.066) (2.040,1.490) (2.080,1.353) (2.120,1.670) (2.160,1.767) (2.200,2.073) (2.240,1.854) (2.280,1.784) (2.320,1.451) (2.360,1.427) (2.400,1.506) (2.440,1.383) (2.480,1.547) (2.520,1.619) (2.560,1.336) (2.600,1.190) (2.640,1.252) (2.680,1.515) (2.720,1.655) (2.760,1.787) (2.800,1.677) (2.840,1.791) (2.880,1.402) (2.920,1.482) (2.960,1.390) (3.000,1.461) (3.040,1.896) (3.080,2.039) (3.120,1.960) (3.160,1.777) (3.200,1.763) (3.240,1.847) (3.280,2.646) (3.320,2.648) (3.360,2.020) (3.400,2.025) (3.440,1.632) (3.480,1.995) (3.520,2.001) (3.560,1.555) (3.600,1.924) (3.640,1.897) (3.680,1.764) (3.720,1.947) (3.760,2.212) (3.800,1.907) (3.840,2.169) (3.880,1.916) (3.920,1.887) (3.960,1.895) (4.000,2.180) (4.040,2.044) (4.080,2.416) (4.120,2.103) (4.160,1.837) (4.200,2.082) (4.240,1.749) (4.280,1.914) (4.320,2.033) (4.360,1.661) (4.400,1.796) (4.440,1.613) (4.480,1.402) (4.520,1.231) (4.560,1.165) (4.600,1.064) (4.640,1.082) (4.680,0.821) (4.720,0.785) (4.760,0.927) (4.800,0.699) (4.840,0.699) (4.880,0.622) (4.920,0.664) (4.960,0.670) (5.000,0.807) (5.040,0.676) (5.080,0.621) (5.120,0.595) (5.160,0.516) (5.200,0.551) (5.240,0.596) (5.280,0.580) (5.320,0.493) (5.360,0.514) (5.400,0.455) (5.440,0.352) (5.480,0.312) (5.520,0.302) (5.560,0.254) (5.600,0.259) (5.640,0.302) (5.680,0.337) (5.720,0.346) (5.760,0.385) (5.800,0.488) (5.840,0.602) (5.880,0.753) (5.920,0.598) (5.960,0.568) (6.000,0.470) (6.040,0.491) (6.080,0.497) (6.120,0.520) (6.160,0.474) (6.200,0.505) (6.240,0.432) (6.280,0.350) (6.320,0.315) (6.360,0.338) (6.400,0.314) (6.440,0.321) (6.480,0.258) (6.520,0.264) (6.560,0.254) (6.600,0.322) (6.640,0.314) (6.680,0.314) (6.720,0.332) (6.760,0.326) (6.800,0.354) (6.840,0.318) (6.880,0.286) (6.920,0.244) (6.960,0.210) (7.000,0.177) (7.040,0.194) (7.080,0.234) (7.120,0.217) (7.160,0.224) (7.200,0.195) (7.240,0.175) (7.280,0.184) (7.320,0.221) (7.360,0.208) (7.400,0.222) (7.440,0.197) (7.480,0.193) (7.520,0.232) (7.560,0.281) (7.600,0.294) (7.640,0.255) (7.680,0.265) (7.720,0.327) (7.760,0.308) (7.800,0.340) (7.840,0.332) (7.880,0.352) (7.920,0.283) (7.960,0.337) (8.000,0.418) (8.040,0.414) (8.080,0.447) (8.120,0.403) (8.160,0.435) (8.200,0.459) (8.240,0.456) (8.280,0.374) (8.320,0.378) (8.360,0.366) (8.400,0.326) (8.440,0.318) (8.480,0.351) (8.520,0.330) (8.560,0.309) (8.600,0.271) (8.640,0.267) (8.680,0.300) (8.720,0.263) (8.760,0.314) (8.800,0.272) (8.840,0.304) (8.880,0.319) (8.920,0.436) (8.960,0.477) (9.000,0.548) (9.040,0.563) (9.080,0.612) (9.120,0.525) (9.160,0.484) (9.200,0.460) (9.240,0.322) (9.280,0.288) (9.320,0.336) (9.360,0.341) (9.400,0.301) (9.440,0.243) (9.480,0.206) (9.520,0.218) (9.560,0.225) (9.600,0.225) (9.640,0.230) (9.680,0.243) (9.720,0.214) (9.760,0.195) (9.800,0.224) (9.840,0.172) (9.880,0.196) (9.920,0.184) (9.960,0.185) (10.000,0.216)};
  \addplot[SABlue!80,thin,forget plot] coordinates {(0.000,1.000) (0.040,0.920) (0.080,0.731) (0.120,0.728) (0.160,0.752) (0.200,0.680) (0.240,0.685) (0.280,0.711) (0.320,0.674) (0.360,0.692) (0.400,0.865) (0.440,0.965) (0.480,0.921) (0.520,0.904) (0.560,0.953) (0.600,0.810) (0.640,0.702) (0.680,0.716) (0.720,0.659) (0.760,0.693) (0.800,0.648) (0.840,0.582) (0.880,0.700) (0.920,0.692) (0.960,0.684) (1.000,0.639) (1.040,0.656) (1.080,0.550) (1.120,0.517) (1.160,0.515) (1.200,0.519) (1.240,0.568) (1.280,0.626) (1.320,0.571) (1.360,0.540) (1.400,0.532) (1.440,0.492) (1.480,0.497) (1.520,0.572) (1.560,0.563) (1.600,0.572) (1.640,0.571) (1.680,0.474) (1.720,0.496) (1.760,0.408) (1.800,0.414) (1.840,0.468) (1.880,0.458) (1.920,0.570) (1.960,0.490) (2.000,0.517) (2.040,0.514) (2.080,0.438) (2.120,0.371) (2.160,0.421) (2.200,0.387) (2.240,0.448) (2.280,0.626) (2.320,0.463) (2.360,0.473) (2.400,0.599) (2.440,0.507) (2.480,0.477) (2.520,0.468) (2.560,0.462) (2.600,0.489) (2.640,0.508) (2.680,0.614) (2.720,0.567) (2.760,0.499) (2.800,0.431) (2.840,0.505) (2.880,0.429) (2.920,0.455) (2.960,0.418) (3.000,0.368) (3.040,0.314) (3.080,0.301) (3.120,0.357) (3.160,0.404) (3.200,0.361) (3.240,0.264) (3.280,0.249) (3.320,0.255) (3.360,0.301) (3.400,0.250) (3.440,0.244) (3.480,0.231) (3.520,0.218) (3.560,0.233) (3.600,0.231) (3.640,0.267) (3.680,0.275) (3.720,0.198) (3.760,0.176) (3.800,0.176) (3.840,0.162) (3.880,0.169) (3.920,0.196) (3.960,0.199) (4.000,0.215) (4.040,0.213) (4.080,0.262) (4.120,0.284) (4.160,0.270) (4.200,0.243) (4.240,0.219) (4.280,0.216) (4.320,0.202) (4.360,0.209) (4.400,0.204) (4.440,0.181) (4.480,0.178) (4.520,0.185) (4.560,0.196) (4.600,0.191) (4.640,0.179) (4.680,0.174) (4.720,0.198) (4.760,0.218) (4.800,0.214) (4.840,0.266) (4.880,0.256) (4.920,0.215) (4.960,0.232) (5.000,0.211) (5.040,0.220) (5.080,0.204) (5.120,0.223) (5.160,0.196) (5.200,0.170) (5.240,0.141) (5.280,0.168) (5.320,0.183) (5.360,0.181) (5.400,0.256) (5.440,0.308) (5.480,0.279) (5.520,0.264) (5.560,0.333) (5.600,0.337) (5.640,0.396) (5.680,0.411) (5.720,0.427) (5.760,0.402) (5.800,0.440) (5.840,0.428) (5.880,0.526) (5.920,0.562) (5.960,0.510) (6.000,0.438) (6.040,0.407) (6.080,0.402) (6.120,0.420) (6.160,0.340) (6.200,0.346) (6.240,0.296) (6.280,0.281) (6.320,0.314) (6.360,0.369) (6.400,0.384) (6.440,0.410) (6.480,0.427) (6.520,0.480) (6.560,0.615) (6.600,0.549) (6.640,0.512) (6.680,0.606) (6.720,0.511) (6.760,0.535) (6.800,0.442) (6.840,0.529) (6.880,0.500) (6.920,0.475) (6.960,0.484) (7.000,0.435) (7.040,0.479) (7.080,0.489) (7.120,0.420) (7.160,0.479) (7.200,0.549) (7.240,0.506) (7.280,0.377) (7.320,0.404) (7.360,0.317) (7.400,0.306) (7.440,0.261) (7.480,0.270) (7.520,0.263) (7.560,0.345) (7.600,0.340) (7.640,0.332) (7.680,0.319) (7.720,0.317) (7.760,0.331) (7.800,0.290) (7.840,0.275) (7.880,0.249) (7.920,0.241) (7.960,0.229) (8.000,0.272) (8.040,0.283) (8.080,0.279) (8.120,0.298) (8.160,0.410) (8.200,0.459) (8.240,0.528) (8.280,0.533) (8.320,0.494) (8.360,0.494) (8.400,0.417) (8.440,0.358) (8.480,0.382) (8.520,0.346) (8.560,0.373) (8.600,0.412) (8.640,0.343) (8.680,0.337) (8.720,0.267) (8.760,0.283) (8.800,0.254) (8.840,0.203) (8.880,0.239) (8.920,0.230) (8.960,0.253) (9.000,0.264) (9.040,0.249) (9.080,0.265) (9.120,0.306) (9.160,0.379) (9.200,0.312) (9.240,0.360) (9.280,0.367) (9.320,0.308) (9.360,0.302) (9.400,0.294) (9.440,0.322) (9.480,0.288) (9.520,0.279) (9.560,0.350) (9.600,0.306) (9.640,0.308) (9.680,0.355) (9.720,0.334) (9.760,0.381) (9.800,0.361) (9.840,0.344) (9.880,0.353) (9.920,0.373) (9.960,0.437) (10.000,0.466)};
  \addplot[SABlue!80,thin,forget plot] coordinates {(0.000,1.000) (0.040,1.075) (0.080,0.903) (0.120,0.977) (0.160,0.801) (0.200,1.081) (0.240,1.131) (0.280,0.946) (0.320,1.024) (0.360,1.107) (0.400,0.832) (0.440,0.870) (0.480,0.668) (0.520,0.548) (0.560,0.575) (0.600,0.619) (0.640,0.607) (0.680,0.641) (0.720,0.762) (0.760,0.891) (0.800,0.897) (0.840,0.915) (0.880,0.993) (0.920,0.807) (0.960,0.770) (1.000,0.915) (1.040,0.907) (1.080,0.803) (1.120,0.891) (1.160,0.810) (1.200,0.845) (1.240,0.741) (1.280,0.692) (1.320,0.670) (1.360,0.803) (1.400,0.763) (1.440,0.762) (1.480,0.685) (1.520,0.835) (1.560,0.731) (1.600,0.664) (1.640,0.603) (1.680,0.683) (1.720,0.656) (1.760,0.616) (1.800,0.590) (1.840,0.629) (1.880,0.632) (1.920,0.676) (1.960,0.655) (2.000,0.504) (2.040,0.385) (2.080,0.334) (2.120,0.266) (2.160,0.271) (2.200,0.198) (2.240,0.203) (2.280,0.228) (2.320,0.227) (2.360,0.212) (2.400,0.175) (2.440,0.217) (2.480,0.194) (2.520,0.187) (2.560,0.150) (2.600,0.154) (2.640,0.149) (2.680,0.152) (2.720,0.159) (2.760,0.208) (2.800,0.227) (2.840,0.187) (2.880,0.210) (2.920,0.192) (2.960,0.170) (3.000,0.154) (3.040,0.187) (3.080,0.209) (3.120,0.219) (3.160,0.272) (3.200,0.283) (3.240,0.273) (3.280,0.316) (3.320,0.318) (3.360,0.378) (3.400,0.451) (3.440,0.469) (3.480,0.383) (3.520,0.307) (3.560,0.293) (3.600,0.270) (3.640,0.263) (3.680,0.250) (3.720,0.309) (3.760,0.348) (3.800,0.381) (3.840,0.454) (3.880,0.496) (3.920,0.433) (3.960,0.444) (4.000,0.457) (4.040,0.475) (4.080,0.472) (4.120,0.459) (4.160,0.534) (4.200,0.527) (4.240,0.470) (4.280,0.472) (4.320,0.472) (4.360,0.450) (4.400,0.465) (4.440,0.459) (4.480,0.486) (4.520,0.434) (4.560,0.500) (4.600,0.472) (4.640,0.572) (4.680,0.590) (4.720,0.576) (4.760,0.528) (4.800,0.521) (4.840,0.497) (4.880,0.418) (4.920,0.430) (4.960,0.497) (5.000,0.525) (5.040,0.610) (5.080,0.730) (5.120,0.685) (5.160,0.744) (5.200,0.851) (5.240,0.908) (5.280,0.787) (5.320,0.762) (5.360,0.827) (5.400,0.807) (5.440,0.879) (5.480,0.856) (5.520,0.911) (5.560,0.989) (5.600,0.979) (5.640,1.017) (5.680,1.140) (5.720,1.196) (5.760,1.220) (5.800,1.246) (5.840,1.079) (5.880,0.806) (5.920,0.850) (5.960,1.046) (6.000,0.755) (6.040,1.023) (6.080,1.031) (6.120,1.204) (6.160,1.125) (6.200,1.238) (6.240,1.690) (6.280,1.621) (6.320,1.656) (6.360,1.536) (6.400,1.311) (6.440,1.261) (6.480,1.132) (6.520,1.132) (6.560,1.309) (6.600,1.467) (6.640,1.571) (6.680,1.622) (6.720,1.367) (6.760,1.459) (6.800,1.485) (6.840,1.494) (6.880,1.219) (6.920,1.226) (6.960,1.331) (7.000,1.349) (7.040,1.401) (7.080,1.385) (7.120,1.652) (7.160,1.459) (7.200,1.618) (7.240,1.507) (7.280,1.322) (7.320,1.398) (7.360,1.686) (7.400,2.028) (7.440,2.265) (7.480,1.923) (7.520,2.296) (7.560,2.568) (7.600,2.402) (7.640,2.411) (7.680,2.388) (7.720,2.809) (7.760,2.456) (7.800,2.624) (7.840,2.797) (7.880,2.913) (7.920,2.935) (7.960,3.521) (8.000,3.514) (8.040,2.959) (8.080,2.848) (8.120,3.570) (8.160,3.225) (8.200,3.118) (8.240,3.148) (8.280,2.656) (8.320,2.531) (8.360,2.809) (8.400,2.778) (8.440,3.381) (8.480,3.422) (8.520,3.583) (8.560,4.065) (8.600,4.681) (8.640,4.966) (8.680,5.599) (8.720,5.385) (8.760,5.504) (8.800,6.129) (8.840,5.473) (8.880,4.938) (8.920,4.847) (8.960,4.845) (9.000,4.906) (9.040,4.165) (9.080,4.084) (9.120,4.266) (9.160,4.060) (9.200,3.970) (9.240,3.416) (9.280,3.270) (9.320,2.819) (9.360,3.043) (9.400,3.393) (9.440,3.118) (9.480,2.958) (9.520,3.056) (9.560,2.715) (9.600,2.964) (9.640,2.721) (9.680,3.112) (9.720,2.718) (9.760,2.911) (9.800,2.617) (9.840,2.666) (9.880,2.388) (9.920,2.740) (9.960,2.441) (10.000,3.304)};
  \addplot[SABlue!80,thin,forget plot] coordinates {(0.000,1.000) (0.040,0.867) (0.080,0.892) (0.120,0.866) (0.160,0.846) (0.200,0.870) (0.240,0.821) (0.280,0.867) (0.320,0.832) (0.360,0.762) (0.400,1.010) (0.440,0.810) (0.480,1.006) (0.520,0.854) (0.560,0.837) (0.600,0.659) (0.640,0.605) (0.680,0.561) (0.720,0.629) (0.760,0.662) (0.800,0.602) (0.840,0.530) (0.880,0.569) (0.920,0.499) (0.960,0.565) (1.000,0.525) (1.040,0.605) (1.080,0.595) (1.120,0.520) (1.160,0.452) (1.200,0.410) (1.240,0.425) (1.280,0.456) (1.320,0.477) (1.360,0.465) (1.400,0.467) (1.440,0.512) (1.480,0.433) (1.520,0.443) (1.560,0.358) (1.600,0.410) (1.640,0.399) (1.680,0.399) (1.720,0.422) (1.760,0.392) (1.800,0.429) (1.840,0.478) (1.880,0.413) (1.920,0.452) (1.960,0.408) (2.000,0.370) (2.040,0.384) (2.080,0.333) (2.120,0.317) (2.160,0.320) (2.200,0.307) (2.240,0.331) (2.280,0.367) (2.320,0.398) (2.360,0.379) (2.400,0.451) (2.440,0.513) (2.480,0.622) (2.520,0.719) (2.560,0.748) (2.600,0.560) (2.640,0.601) (2.680,0.566) (2.720,0.612) (2.760,0.557) (2.800,0.514) (2.840,0.431) (2.880,0.441) (2.920,0.526) (2.960,0.486) (3.000,0.445) (3.040,0.490) (3.080,0.514) (3.120,0.480) (3.160,0.452) (3.200,0.480) (3.240,0.544) (3.280,0.502) (3.320,0.462) (3.360,0.434) (3.400,0.380) (3.440,0.395) (3.480,0.461) (3.520,0.413) (3.560,0.387) (3.600,0.336) (3.640,0.315) (3.680,0.269) (3.720,0.256) (3.760,0.283) (3.800,0.343) (3.840,0.338) (3.880,0.306) (3.920,0.256) (3.960,0.229) (4.000,0.196) (4.040,0.220) (4.080,0.254) (4.120,0.207) (4.160,0.206) (4.200,0.181) (4.240,0.161) (4.280,0.162) (4.320,0.155) (4.360,0.201) (4.400,0.224) (4.440,0.225) (4.480,0.258) (4.520,0.278) (4.560,0.252) (4.600,0.241) (4.640,0.253) (4.680,0.285) (4.720,0.328) (4.760,0.305) (4.800,0.254) (4.840,0.200) (4.880,0.228) (4.920,0.193) (4.960,0.241) (5.000,0.286) (5.040,0.296) (5.080,0.302) (5.120,0.275) (5.160,0.268) (5.200,0.277) (5.240,0.225) (5.280,0.228) (5.320,0.218) (5.360,0.193) (5.400,0.264) (5.440,0.307) (5.480,0.297) (5.520,0.289) (5.560,0.267) (5.600,0.260) (5.640,0.286) (5.680,0.310) (5.720,0.298) (5.760,0.310) (5.800,0.339) (5.840,0.304) (5.880,0.311) (5.920,0.284) (5.960,0.336) (6.000,0.379) (6.040,0.361) (6.080,0.417) (6.120,0.428) (6.160,0.405) (6.200,0.391) (6.240,0.511) (6.280,0.565) (6.320,0.627) (6.360,0.543) (6.400,0.585) (6.440,0.688) (6.480,0.834) (6.520,0.925) (6.560,0.932) (6.600,0.857) (6.640,0.806) (6.680,0.889) (6.720,0.956) (6.760,0.854) (6.800,0.907) (6.840,1.061) (6.880,0.955) (6.920,1.043) (6.960,1.025) (7.000,1.056) (7.040,0.891) (7.080,1.045) (7.120,0.900) (7.160,0.940) (7.200,0.985) (7.240,0.922) (7.280,0.817) (7.320,0.755) (7.360,0.727) (7.400,0.675) (7.440,0.960) (7.480,0.865) (7.520,0.821) (7.560,0.767) (7.600,0.715) (7.640,0.657) (7.680,0.739) (7.720,0.763) (7.760,0.772) (7.800,0.691) (7.840,0.656) (7.880,0.735) (7.920,0.673) (7.960,0.823) (8.000,0.742) (8.040,0.881) (8.080,0.959) (8.120,1.000) (8.160,1.092) (8.200,1.049) (8.240,1.191) (8.280,1.048) (8.320,0.894) (8.360,0.882) (8.400,0.914) (8.440,1.009) (8.480,1.096) (8.520,1.163) (8.560,1.200) (8.600,1.384) (8.640,1.482) (8.680,1.316) (8.720,1.411) (8.760,1.362) (8.800,1.396) (8.840,1.492) (8.880,1.433) (8.920,1.307) (8.960,1.223) (9.000,1.050) (9.040,0.901) (9.080,0.897) (9.120,0.849) (9.160,0.888) (9.200,1.057) (9.240,1.304) (9.280,1.096) (9.320,1.350) (9.360,1.640) (9.400,1.709) (9.440,1.339) (9.480,1.453) (9.520,1.581) (9.560,1.698) (9.600,1.417) (9.640,1.354) (9.680,1.513) (9.720,1.370) (9.760,1.342) (9.800,1.268) (9.840,1.439) (9.880,1.260) (9.920,1.252) (9.960,1.407) (10.000,1.485)};
  \legend{均值 $\E S_t=e^{\mu t}$\quad, 中位数 $e^{(\mu-\sigma^2/2)t}$}
\end{axis}
\end{tikzpicture}
\caption{几何布朗运动 ($s_0=1$, $\mu=0.1$, $\sigma=0.6$, 纵轴为对数刻度). 由于 $\mu<\frac12\sigma^2$, 均值指数增长, 典型轨道 (中位数) 却指数下降并最终趋于零; 均值被少数涨幅极大的轨道撑起}
\label{fig:gbm}
\end{figure}

\begin{intuition}
(4) 与 (1) 的对比称为 ``波动率拖累'': 当 $0<\mu<\frac{1}{2}\sigma^2$ 时, $\E S_t\to\infty$, 而 $S_t\to0$ $\as$ 均值由极少数涨幅巨大的轨道撑起, 典型的轨道却在下降 (图~\ref{fig:gbm}). 原因是收益以乘法方式复合: 先涨 $50\%$ 再跌 $50\%$, 结果是 $1.5\times0.5=0.75$, 并未回到原点. 长期增长率由对数收益的均值 $\mu-\frac{1}{2}\sigma^2$ 决定, 而非收益的均值 $\mu$.
\end{intuition}

\paragraph{布朗桥.} 布朗桥是 ``终点被固定'' 的布朗运动. 它出现在统计检验 (经验分布函数的极限, Kolmogorov-Smirnov 检验)、模拟算法 (先生成终点再填充中间, 即 L\'evy-Ciesielski 构造的逐层加细) 以及路径积分中.

\begin{definition}[布朗桥]\label{def:bridge}
$[0,1]$ 上轨道连续、均值为零的高斯过程 $X$ 称为 (从 $0$ 到 $0$ 的) \textbf{布朗桥}, 如果
\[
\Cov(X_s,X_t)=s\wedge t-st,\qquad s,t\in[0,1].
\]
\end{definition}

由于高斯过程的有限维分布由均值与协方差决定, 布朗桥的分布 (作为 $C[0,1]$ 上的测度) 是唯一的. 特别地 $\Var(X_t)=t(1-t)$, $X_0=X_1=0$ $\as$

\begin{theorem}[布朗桥的三种构造]\label{thm:bridge}
设 $B$, $W$ 是标准布朗运动.
\begin{enumerate}
\item $X_t=B_t-tB_1$ ($t\in[0,1]$) 是布朗桥, 并且过程 $X$ 与 $B_1$ 独立.
\item 对 $0<t_1<\cdots<t_n<1$ 与 $y\in\R$, $(B_{t_1},\dots,B_{t_n})$ 在给定 $B_1=y$ 下的条件分布, 是 $(X_{t_1}+t_1y,\dots,X_{t_n}+t_ny)$ 的分布. 特别地, 给定 $B_1=0$ 时, 布朗运动在 $[0,1]$ 上的条件分布就是布朗桥.
\item $Z_t=(1-t)W_{t/(1-t)}$ ($0\le t<1$), $Z_1=0$ 是布朗桥.
\end{enumerate}
\end{theorem}

\begin{proof}
(1) $X$ 的轨道连续, 是 $B$ 的线性变换, 因而是均值为零的高斯过程 (命题~\ref{prop:gaussian-operations}). 对 $s,t\in[0,1]$,
\begin{align*}
\Cov(X_s,X_t)&=\Cov(B_s,B_t)-t\Cov(B_s,B_1)-s\Cov(B_1,B_t)+st\Var(B_1)\\
&=s\wedge t-st-st+st=s\wedge t-st.
\end{align*}
又 $\Cov(X_t,B_1)=t-t\cdot1=0$. 对任意 $t_1,\dots,t_n$, $(X_{t_1},\dots,X_{t_n},B_1)$ 是正态向量, 而 $B_1$ 与每个 $X_{t_i}$ 不相关, 由定理~\ref{thm:gaussian-uncorrelated}, $(X_{t_1},\dots,X_{t_n})$ 与 $B_1$ 独立. 这样的事件构成生成 $\sigma(X)$ 的 $\pi$-系, 由定理~\ref{thm:indep-pi}, $\sigma(X)$ 与 $B_1$ 独立.

(2) $B_{t_i}=X_{t_i}+t_iB_1$, 其中 $(X_{t_1},\dots,X_{t_n})$ 与 $B_1$ 独立. 由代入引理~\ref{lem:freezing}, 对有界 Borel 函数 $g$,
\[
\E\big[g(B_{t_1},\dots,B_{t_n})\,\big|\,B_1\big]=\psi(B_1),\qquad\psi(y)=\E\big[g(X_{t_1}+t_1y,\dots,X_{t_n}+t_ny)\big].
\]
因此 $(X_{t_i}+t_iy)_i$ 的分布是给定 $B_1=y$ 时的正则条件分布. (这与定理~\ref{thm:gaussian-conditional} 的结论一致: 条件均值为 $(t_iy)_i$, 条件协方差为 $t_i\wedge t_j-t_it_j$, 参见例~\ref{ex:bridge-preview}.)

(3) $Z$ 在 $[0,1)$ 上轨道连续, 是均值为零的高斯过程. 对 $0\le s\le t<1$, 由于 $u\mapsto\frac{u}{1-u}$ 单调递增,
\[
\Cov(Z_s,Z_t)=(1-s)(1-t)\Big(\frac{s}{1-s}\wedge\frac{t}{1-t}\Big)=(1-s)(1-t)\frac{s}{1-s}=s(1-t)=s\wedge t-st.
\]
在 $t=1$ 处的连续性: 令 $u=\frac{t}{1-t}$, 则 $1-t=\frac{1}{1+u}$, $Z_t=\frac{W_u}{1+u}$. 当 $t\uparrow1$ 时 $u\to\infty$, 由布朗运动的强大数定律 (推论~\ref{cor:bm-slln}), $\frac{W_u}{1+u}=\frac{W_u}{u}\cdot\frac{u}{1+u}\to0$ $\as$ 在这个零概率集之外, $Z$ 在 $[0,1]$ 上连续. 涉及 $Z_1=0$ 的协方差为零, 也等于 $s\wedge1-s\cdot1$.
\end{proof}

\begin{figure}[htbp]
\centering
\begin{tikzpicture}
\begin{axis}[SAaxis, axis x line=bottom, axis y line=left, width=12cm, height=5.6cm, xmin=0, xmax=1.03, ymin=-1.3, ymax=1.3,
  xtick={0,0.25,0.5,0.75,1}, ytick={-1,-0.5,0,0.5,1}, xlabel={$t$}]
  \addplot[draw=none,fill=SAYellow!25,domain=0:1,samples=100] {sqrt(abs(x*(1-x)))} \closedcycle;
  \addplot[draw=none,fill=SAYellow!25,domain=0:1,samples=100] {-sqrt(abs(x*(1-x)))} \closedcycle;
  \addplot[SAYellowText,dashed,domain=0:1,samples=100] {sqrt(abs(x*(1-x)))};
  \addplot[SAYellowText,dashed,domain=0:1,samples=100] {-sqrt(abs(x*(1-x)))};
  \addplot[SABlue,thin] coordinates {(0.000,0.000) (0.004,0.119) (0.008,0.126) (0.012,0.067) (0.016,0.127) (0.020,0.242) (0.024,0.237) (0.028,0.253) (0.032,0.291) (0.036,0.210) (0.040,0.117) (0.044,0.091) (0.048,0.074) (0.052,0.032) (0.056,0.041) (0.060,0.033) (0.064,0.054) (0.068,0.138) (0.072,0.161) (0.076,0.100) (0.080,0.134) (0.084,0.054) (0.088,-0.012) (0.092,-0.029) (0.096,0.048) (0.100,0.014) (0.104,0.036) (0.108,0.096) (0.112,0.184) (0.116,0.143) (0.120,0.160) (0.124,0.098) (0.128,0.039) (0.132,0.070) (0.136,0.073) (0.140,0.098) (0.144,0.040) (0.148,0.132) (0.152,0.185) (0.156,0.221) (0.160,0.191) (0.164,0.214) (0.168,0.233) (0.172,0.294) (0.176,0.248) (0.180,0.282) (0.184,0.265) (0.188,0.285) (0.192,0.194) (0.196,0.314) (0.200,0.331) (0.204,0.310) (0.208,0.196) (0.212,0.157) (0.216,0.150) (0.220,0.152) (0.224,0.101) (0.228,0.012) (0.232,0.015) (0.236,0.013) (0.240,0.029) (0.244,0.048) (0.248,-0.069) (0.252,-0.095) (0.256,-0.126) (0.260,-0.088) (0.264,-0.085) (0.268,-0.208) (0.272,-0.251) (0.276,-0.373) (0.280,-0.316) (0.284,-0.214) (0.288,-0.219) (0.292,-0.236) (0.296,-0.258) (0.300,-0.198) (0.304,-0.225) (0.308,-0.207) (0.312,-0.188) (0.316,-0.059) (0.320,-0.047) (0.324,0.124) (0.328,0.175) (0.332,0.207) (0.336,0.178) (0.340,0.135) (0.344,0.095) (0.348,0.187) (0.352,0.100) (0.356,0.121) (0.360,-0.059) (0.364,-0.090) (0.368,-0.098) (0.372,-0.103) (0.376,-0.046) (0.380,-0.007) (0.384,-0.104) (0.388,-0.097) (0.392,-0.203) (0.396,-0.316) (0.400,-0.334) (0.404,-0.453) (0.408,-0.468) (0.412,-0.474) (0.416,-0.444) (0.420,-0.494) (0.424,-0.577) (0.428,-0.533) (0.432,-0.499) (0.436,-0.409) (0.440,-0.376) (0.444,-0.390) (0.448,-0.313) (0.452,-0.315) (0.456,-0.246) (0.460,-0.215) (0.464,-0.188) (0.468,-0.156) (0.472,-0.167) (0.476,-0.284) (0.480,-0.248) (0.484,-0.313) (0.488,-0.184) (0.492,-0.111) (0.496,-0.099) (0.500,-0.161) (0.504,-0.152) (0.508,-0.174) (0.512,-0.107) (0.516,0.008) (0.520,-0.039) (0.524,-0.045) (0.528,0.028) (0.532,0.056) (0.536,0.062) (0.540,0.029) (0.544,0.023) (0.548,-0.009) (0.552,0.073) (0.556,0.110) (0.560,0.118) (0.564,0.093) (0.568,0.218) (0.572,0.155) (0.576,0.130) (0.580,0.172) (0.584,0.144) (0.588,0.123) (0.592,0.171) (0.596,0.215) (0.600,0.119) (0.604,0.224) (0.608,0.230) (0.612,0.234) (0.616,0.196) (0.620,0.192) (0.624,0.172) (0.628,0.163) (0.632,0.144) (0.636,0.076) (0.640,0.021) (0.644,0.056) (0.648,0.051) (0.652,0.046) (0.656,0.011) (0.660,0.095) (0.664,0.008) (0.668,-0.035) (0.672,-0.197) (0.676,-0.246) (0.680,-0.420) (0.684,-0.455) (0.688,-0.499) (0.692,-0.445) (0.696,-0.424) (0.700,-0.410) (0.704,-0.491) (0.708,-0.500) (0.712,-0.544) (0.716,-0.559) (0.720,-0.599) (0.724,-0.627) (0.728,-0.695) (0.732,-0.705) (0.736,-0.640) (0.740,-0.740) (0.744,-0.783) (0.748,-0.793) (0.752,-0.752) (0.756,-0.882) (0.760,-0.932) (0.764,-0.920) (0.768,-0.740) (0.772,-0.662) (0.776,-0.604) (0.780,-0.492) (0.784,-0.481) (0.788,-0.489) (0.792,-0.541) (0.796,-0.492) (0.800,-0.520) (0.804,-0.473) (0.808,-0.458) (0.812,-0.517) (0.816,-0.435) (0.820,-0.378) (0.824,-0.336) (0.828,-0.190) (0.832,-0.194) (0.836,-0.235) (0.840,-0.182) (0.844,-0.219) (0.848,-0.224) (0.852,-0.170) (0.856,-0.038) (0.860,-0.055) (0.864,-0.160) (0.868,-0.225) (0.872,-0.199) (0.876,-0.266) (0.880,-0.218) (0.884,-0.192) (0.888,-0.132) (0.892,-0.154) (0.896,-0.169) (0.900,-0.190) (0.904,-0.164) (0.908,-0.239) (0.912,-0.190) (0.916,-0.078) (0.920,-0.079) (0.924,-0.144) (0.928,-0.129) (0.932,-0.031) (0.936,-0.104) (0.940,-0.119) (0.944,-0.158) (0.948,-0.124) (0.952,-0.082) (0.956,-0.097) (0.960,-0.135) (0.964,-0.159) (0.968,-0.188) (0.972,-0.179) (0.976,-0.052) (0.980,-0.194) (0.984,-0.120) (0.988,-0.155) (0.992,-0.115) (0.996,-0.106) (1.000,0.000)};
  \addplot[SARed,thin] coordinates {(0.000,0.000) (0.004,-0.037) (0.008,-0.007) (0.012,0.001) (0.016,-0.034) (0.020,0.006) (0.024,-0.048) (0.028,-0.138) (0.032,-0.060) (0.036,-0.077) (0.040,-0.126) (0.044,-0.079) (0.048,-0.005) (0.052,0.090) (0.056,-0.015) (0.060,-0.081) (0.064,-0.113) (0.068,-0.209) (0.072,-0.198) (0.076,-0.208) (0.080,-0.272) (0.084,-0.395) (0.088,-0.263) (0.092,-0.293) (0.096,-0.282) (0.100,-0.220) (0.104,-0.231) (0.108,-0.205) (0.112,-0.268) (0.116,-0.205) (0.120,-0.209) (0.124,-0.211) (0.128,-0.086) (0.132,-0.071) (0.136,-0.001) (0.140,-0.069) (0.144,-0.054) (0.148,-0.166) (0.152,-0.193) (0.156,-0.194) (0.160,-0.178) (0.164,-0.151) (0.168,-0.099) (0.172,-0.201) (0.176,-0.165) (0.180,-0.130) (0.184,-0.082) (0.188,-0.075) (0.192,-0.151) (0.196,-0.177) (0.200,-0.237) (0.204,-0.150) (0.208,-0.179) (0.212,-0.234) (0.216,-0.193) (0.220,-0.131) (0.224,-0.063) (0.228,-0.109) (0.232,0.014) (0.236,-0.070) (0.240,0.094) (0.244,0.041) (0.248,-0.007) (0.252,-0.013) (0.256,0.012) (0.260,-0.037) (0.264,-0.085) (0.268,-0.037) (0.272,0.008) (0.276,-0.059) (0.280,0.035) (0.284,0.025) (0.288,0.028) (0.292,0.029) (0.296,-0.005) (0.300,-0.016) (0.304,-0.136) (0.308,-0.059) (0.312,-0.009) (0.316,0.109) (0.320,0.083) (0.324,0.195) (0.328,0.180) (0.332,0.120) (0.336,0.080) (0.340,-0.020) (0.344,-0.081) (0.348,-0.166) (0.352,-0.143) (0.356,-0.070) (0.360,-0.105) (0.364,-0.096) (0.368,-0.214) (0.372,-0.305) (0.376,-0.388) (0.380,-0.407) (0.384,-0.471) (0.388,-0.406) (0.392,-0.379) (0.396,-0.371) (0.400,-0.414) (0.404,-0.423) (0.408,-0.447) (0.412,-0.471) (0.416,-0.337) (0.420,-0.380) (0.424,-0.384) (0.428,-0.348) (0.432,-0.281) (0.436,-0.355) (0.440,-0.356) (0.444,-0.405) (0.448,-0.423) (0.452,-0.471) (0.456,-0.553) (0.460,-0.548) (0.464,-0.513) (0.468,-0.602) (0.472,-0.527) (0.476,-0.585) (0.480,-0.612) (0.484,-0.609) (0.488,-0.565) (0.492,-0.578) (0.496,-0.568) (0.500,-0.631) (0.504,-0.614) (0.508,-0.673) (0.512,-0.732) (0.516,-0.724) (0.520,-0.794) (0.524,-0.633) (0.528,-0.642) (0.532,-0.604) (0.536,-0.621) (0.540,-0.616) (0.544,-0.604) (0.548,-0.710) (0.552,-0.737) (0.556,-0.716) (0.560,-0.688) (0.564,-0.565) (0.568,-0.531) (0.572,-0.514) (0.576,-0.583) (0.580,-0.644) (0.584,-0.661) (0.588,-0.634) (0.592,-0.622) (0.596,-0.774) (0.600,-0.770) (0.604,-0.721) (0.608,-0.648) (0.612,-0.545) (0.616,-0.621) (0.620,-0.490) (0.624,-0.432) (0.628,-0.503) (0.632,-0.454) (0.636,-0.506) (0.640,-0.514) (0.644,-0.406) (0.648,-0.428) (0.652,-0.328) (0.656,-0.408) (0.660,-0.440) (0.664,-0.369) (0.668,-0.366) (0.672,-0.417) (0.676,-0.424) (0.680,-0.446) (0.684,-0.434) (0.688,-0.509) (0.692,-0.574) (0.696,-0.648) (0.700,-0.495) (0.704,-0.509) (0.708,-0.560) (0.712,-0.567) (0.716,-0.572) (0.720,-0.553) (0.724,-0.540) (0.728,-0.511) (0.732,-0.408) (0.736,-0.428) (0.740,-0.361) (0.744,-0.378) (0.748,-0.210) (0.752,-0.183) (0.756,-0.309) (0.760,-0.334) (0.764,-0.329) (0.768,-0.389) (0.772,-0.440) (0.776,-0.318) (0.780,-0.340) (0.784,-0.375) (0.788,-0.490) (0.792,-0.373) (0.796,-0.347) (0.800,-0.307) (0.804,-0.211) (0.808,-0.169) (0.812,-0.301) (0.816,-0.338) (0.820,-0.345) (0.824,-0.281) (0.828,-0.393) (0.832,-0.406) (0.836,-0.314) (0.840,-0.316) (0.844,-0.388) (0.848,-0.317) (0.852,-0.297) (0.856,-0.243) (0.860,-0.228) (0.864,-0.230) (0.868,-0.218) (0.872,-0.236) (0.876,-0.122) (0.880,-0.163) (0.884,-0.068) (0.888,-0.086) (0.892,-0.162) (0.896,-0.117) (0.900,-0.093) (0.904,-0.184) (0.908,-0.354) (0.912,-0.354) (0.916,-0.435) (0.920,-0.327) (0.924,-0.212) (0.928,-0.185) (0.932,-0.100) (0.936,-0.082) (0.940,-0.033) (0.944,0.019) (0.948,0.032) (0.952,0.034) (0.956,-0.044) (0.960,-0.108) (0.964,-0.037) (0.968,0.084) (0.972,0.076) (0.976,0.035) (0.980,0.139) (0.984,0.099) (0.988,0.127) (0.992,0.031) (0.996,0.029) (1.000,0.000)};
  \addplot[SAGreenText,thin] coordinates {(0.000,0.000) (0.004,0.006) (0.008,-0.060) (0.012,-0.106) (0.016,-0.096) (0.020,-0.186) (0.024,-0.226) (0.028,-0.324) (0.032,-0.341) (0.036,-0.366) (0.040,-0.392) (0.044,-0.371) (0.048,-0.493) (0.052,-0.490) (0.056,-0.446) (0.060,-0.474) (0.064,-0.471) (0.068,-0.473) (0.072,-0.437) (0.076,-0.452) (0.080,-0.495) (0.084,-0.506) (0.088,-0.418) (0.092,-0.407) (0.096,-0.398) (0.100,-0.322) (0.104,-0.339) (0.108,-0.360) (0.112,-0.385) (0.116,-0.392) (0.120,-0.415) (0.124,-0.390) (0.128,-0.530) (0.132,-0.505) (0.136,-0.563) (0.140,-0.566) (0.144,-0.639) (0.148,-0.608) (0.152,-0.609) (0.156,-0.490) (0.160,-0.462) (0.164,-0.516) (0.168,-0.655) (0.172,-0.634) (0.176,-0.705) (0.180,-0.645) (0.184,-0.631) (0.188,-0.684) (0.192,-0.616) (0.196,-0.565) (0.200,-0.610) (0.204,-0.538) (0.208,-0.507) (0.212,-0.525) (0.216,-0.548) (0.220,-0.649) (0.224,-0.622) (0.228,-0.675) (0.232,-0.716) (0.236,-0.673) (0.240,-0.736) (0.244,-0.678) (0.248,-0.749) (0.252,-0.712) (0.256,-0.783) (0.260,-0.867) (0.264,-0.771) (0.268,-0.649) (0.272,-0.574) (0.276,-0.640) (0.280,-0.586) (0.284,-0.662) (0.288,-0.580) (0.292,-0.585) (0.296,-0.628) (0.300,-0.658) (0.304,-0.685) (0.308,-0.650) (0.312,-0.700) (0.316,-0.724) (0.320,-0.753) (0.324,-0.677) (0.328,-0.639) (0.332,-0.552) (0.336,-0.570) (0.340,-0.646) (0.344,-0.628) (0.348,-0.684) (0.352,-0.648) (0.356,-0.696) (0.360,-0.774) (0.364,-0.690) (0.368,-0.641) (0.372,-0.589) (0.376,-0.485) (0.380,-0.488) (0.384,-0.483) (0.388,-0.535) (0.392,-0.516) (0.396,-0.543) (0.400,-0.453) (0.404,-0.363) (0.408,-0.338) (0.412,-0.429) (0.416,-0.468) (0.420,-0.494) (0.424,-0.477) (0.428,-0.490) (0.432,-0.476) (0.436,-0.396) (0.440,-0.297) (0.444,-0.178) (0.448,-0.182) (0.452,-0.238) (0.456,-0.140) (0.460,-0.130) (0.464,-0.067) (0.468,-0.115) (0.472,-0.087) (0.476,-0.071) (0.480,-0.017) (0.484,-0.098) (0.488,-0.045) (0.492,0.030) (0.496,0.049) (0.500,-0.007) (0.504,0.068) (0.508,0.008) (0.512,0.124) (0.516,0.132) (0.520,0.178) (0.524,0.059) (0.528,0.084) (0.532,-0.003) (0.536,0.024) (0.540,-0.007) (0.544,-0.020) (0.548,0.155) (0.552,0.167) (0.556,0.087) (0.560,0.223) (0.564,0.324) (0.568,0.170) (0.572,0.179) (0.576,0.136) (0.580,0.208) (0.584,0.165) (0.588,0.125) (0.592,0.041) (0.596,0.167) (0.600,0.263) (0.604,0.263) (0.608,0.324) (0.612,0.314) (0.616,0.298) (0.620,0.316) (0.624,0.285) (0.628,0.322) (0.632,0.349) (0.636,0.166) (0.640,0.286) (0.644,0.369) (0.648,0.242) (0.652,0.289) (0.656,0.257) (0.660,0.272) (0.664,0.227) (0.668,0.240) (0.672,0.380) (0.676,0.355) (0.680,0.369) (0.684,0.356) (0.688,0.323) (0.692,0.329) (0.696,0.321) (0.700,0.254) (0.704,0.215) (0.708,0.252) (0.712,0.281) (0.716,0.357) (0.720,0.317) (0.724,0.313) (0.728,0.339) (0.732,0.328) (0.736,0.348) (0.740,0.351) (0.744,0.384) (0.748,0.486) (0.752,0.459) (0.756,0.363) (0.760,0.302) (0.764,0.266) (0.768,0.156) (0.772,0.200) (0.776,0.215) (0.780,0.428) (0.784,0.450) (0.788,0.468) (0.792,0.418) (0.796,0.455) (0.800,0.532) (0.804,0.542) (0.808,0.610) (0.812,0.606) (0.816,0.617) (0.820,0.736) (0.824,0.750) (0.828,0.729) (0.832,0.785) (0.836,0.806) (0.840,0.778) (0.844,0.659) (0.848,0.629) (0.852,0.657) (0.856,0.487) (0.860,0.468) (0.864,0.451) (0.868,0.341) (0.872,0.386) (0.876,0.276) (0.880,0.232) (0.884,0.352) (0.888,0.195) (0.892,0.144) (0.896,0.186) (0.900,0.082) (0.904,0.048) (0.908,0.133) (0.912,0.133) (0.916,0.151) (0.920,0.155) (0.924,0.070) (0.928,-0.016) (0.932,-0.043) (0.936,-0.078) (0.940,-0.120) (0.944,-0.004) (0.948,-0.063) (0.952,-0.128) (0.956,-0.099) (0.960,-0.153) (0.964,-0.218) (0.968,-0.163) (0.972,-0.134) (0.976,-0.105) (0.980,-0.064) (0.984,0.007) (0.988,-0.038) (0.992,0.027) (0.996,-0.046) (1.000,0.000)};
  \addplot[SADark,thin] coordinates {(0.000,0.000) (0.004,-0.028) (0.008,0.037) (0.012,0.087) (0.016,-0.031) (0.020,-0.069) (0.024,-0.113) (0.028,-0.011) (0.032,0.055) (0.036,0.067) (0.040,0.101) (0.044,0.025) (0.048,0.019) (0.052,0.077) (0.056,0.086) (0.060,0.035) (0.064,0.046) (0.068,0.058) (0.072,0.149) (0.076,-0.021) (0.080,0.073) (0.084,0.092) (0.088,0.179) (0.092,0.119) (0.096,0.160) (0.100,0.085) (0.104,0.131) (0.108,0.084) (0.112,0.115) (0.116,0.105) (0.120,0.041) (0.124,0.073) (0.128,-0.052) (0.132,-0.083) (0.136,-0.012) (0.140,0.043) (0.144,-0.072) (0.148,-0.043) (0.152,-0.069) (0.156,-0.051) (0.160,-0.114) (0.164,-0.068) (0.168,-0.101) (0.172,0.012) (0.176,-0.103) (0.180,-0.056) (0.184,-0.062) (0.188,-0.124) (0.192,-0.083) (0.196,-0.062) (0.200,-0.051) (0.204,-0.056) (0.208,-0.028) (0.212,0.020) (0.216,-0.002) (0.220,0.155) (0.224,0.171) (0.228,0.226) (0.232,0.337) (0.236,0.378) (0.240,0.408) (0.244,0.353) (0.248,0.348) (0.252,0.397) (0.256,0.321) (0.260,0.340) (0.264,0.365) (0.268,0.247) (0.272,0.182) (0.276,0.111) (0.280,0.155) (0.284,0.196) (0.288,0.249) (0.292,0.260) (0.296,0.358) (0.300,0.432) (0.304,0.383) (0.308,0.385) (0.312,0.401) (0.316,0.406) (0.320,0.497) (0.324,0.671) (0.328,0.702) (0.332,0.668) (0.336,0.585) (0.340,0.611) (0.344,0.604) (0.348,0.602) (0.352,0.559) (0.356,0.708) (0.360,0.676) (0.364,0.713) (0.368,0.732) (0.372,0.758) (0.376,0.783) (0.380,0.800) (0.384,0.822) (0.388,0.652) (0.392,0.598) (0.396,0.585) (0.400,0.530) (0.404,0.561) (0.408,0.530) (0.412,0.622) (0.416,0.634) (0.420,0.538) (0.424,0.663) (0.428,0.567) (0.432,0.500) (0.436,0.518) (0.440,0.508) (0.444,0.463) (0.448,0.434) (0.452,0.369) (0.456,0.316) (0.460,0.174) (0.464,0.219) (0.468,0.162) (0.472,0.194) (0.476,0.214) (0.480,0.196) (0.484,0.172) (0.488,0.210) (0.492,0.188) (0.496,0.223) (0.500,0.247) (0.504,0.174) (0.508,0.133) (0.512,0.158) (0.516,0.172) (0.520,0.228) (0.524,0.219) (0.528,0.301) (0.532,0.251) (0.536,0.221) (0.540,0.183) (0.544,0.302) (0.548,0.271) (0.552,0.243) (0.556,0.154) (0.560,0.117) (0.564,0.046) (0.568,0.123) (0.572,0.096) (0.576,0.087) (0.580,0.200) (0.584,0.169) (0.588,0.295) (0.592,0.300) (0.596,0.241) (0.600,0.188) (0.604,0.207) (0.608,0.251) (0.612,0.278) (0.616,0.334) (0.620,0.222) (0.624,0.183) (0.628,0.086) (0.632,0.104) (0.636,0.128) (0.640,0.132) (0.644,0.218) (0.648,0.234) (0.652,0.228) (0.656,0.125) (0.660,0.174) (0.664,0.163) (0.668,0.173) (0.672,0.150) (0.676,0.176) (0.680,0.108) (0.684,0.102) (0.688,-0.020) (0.692,0.013) (0.696,0.075) (0.700,0.100) (0.704,0.084) (0.708,0.101) (0.712,0.088) (0.716,0.031) (0.720,-0.013) (0.724,0.094) (0.728,0.165) (0.732,0.179) (0.736,0.312) (0.740,0.268) (0.744,0.306) (0.748,0.286) (0.752,0.147) (0.756,0.073) (0.760,0.119) (0.764,0.196) (0.768,0.144) (0.772,0.160) (0.776,0.032) (0.780,0.031) (0.784,0.063) (0.788,0.065) (0.792,0.009) (0.796,0.011) (0.800,0.074) (0.804,0.036) (0.808,0.085) (0.812,-0.007) (0.816,0.024) (0.820,0.170) (0.824,0.134) (0.828,0.140) (0.832,0.177) (0.836,0.225) (0.840,0.092) (0.844,0.182) (0.848,0.241) (0.852,0.243) (0.856,0.256) (0.860,0.206) (0.864,0.233) (0.868,0.306) (0.872,0.262) (0.876,0.292) (0.880,0.374) (0.884,0.261) (0.888,0.208) (0.892,0.199) (0.896,0.305) (0.900,0.243) (0.904,0.300) (0.908,0.274) (0.912,0.234) (0.916,0.252) (0.920,0.292) (0.924,0.304) (0.928,0.351) (0.932,0.280) (0.936,0.325) (0.940,0.298) (0.944,0.264) (0.948,0.365) (0.952,0.276) (0.956,0.270) (0.960,0.288) (0.964,0.298) (0.968,0.222) (0.972,0.189) (0.976,0.240) (0.980,0.262) (0.984,0.197) (0.988,0.195) (0.992,0.053) (0.996,0.097) (1.000,0.000)};
\end{axis}
\end{tikzpicture}
\caption{布朗桥 $X_t=B_t-tB_1$ 的四条轨道. 阴影为一倍标准差范围 $\abs{x}\le\sqrt{t(1-t)}$: 不确定性在中点最大, 在两端被 ``钉住'' 为零}
\label{fig:bridge}
\end{figure}

\begin{remark}
更一般地, 在 $[0,T]$ 上从 $x$ 到 $y$ 的布朗桥可以写为 $x+\frac{t}{T}(y-x)+\big(W_t-\frac{t}{T}W_T\big)$. 构造 (3) 说明布朗桥可以由布朗运动经过时间变换与缩放得到; 第~6~章将看到, 它是随机微分方程 $\dd X_t=-\frac{X_t}{1-t}\dd t+\dd W_t$ 的解, 漂移项 $-\frac{X_t}{1-t}$ 随着 $t\to1$ 越来越强地把过程拉回零点.
\end{remark}

\begin{example}[Ornstein-Uhlenbeck 过程]\label{ex:ou}
令 $U_t=e^{-t}W_{e^{2t}}$ ($t\in\R$). 它是均值为零的连续高斯过程, 且对 $s\le t$,
\[
\Cov(U_s,U_t)=e^{-s-t}\big(e^{2s}\wedge e^{2t}\big)=e^{-s-t}e^{2s}=e^{-(t-s)}.
\]
协方差只依赖于时间差, 所以 $U$ 是平稳过程, 每个 $U_t\sim N(0,1)$. 它是第~2~章正态自回归过程 (例~\ref{ex:ar1}) 的连续时间版本: 在等间隔 $h$ 的时刻上观察, $U_{(k+1)h}=e^{-h}U_{kh}+e^{-(k+1)h}\big(W_{e^{2(k+1)h}}-W_{e^{2kh}}\big)$, 其中后一项与 $\{U_{jh}\}_{j\le k}$ 独立, 服从 $N(0,1-e^{-2h})$. 所以 $\{U_{kh}\}_k$ 正是系数为 $a=e^{-h}$ 的正态自回归过程. 第~6~章将看到它是方程 $\dd U_t=-U_t\dd t+\sqrt{2}\dd\widetilde{W}_t$ 的平稳解, 描述受到回复力作用的粒子速度, 在金融中则是利率的 Vasicek 模型.
\end{example}

% ---------------------------------------------------------------------
\subsection{重对数律}

由推论~\ref{cor:bm-germ} 与时间反演, $B_t/\sqrt{t}$ 当 $t\to\infty$ 时上下无界; 另一方面, 由强大数定律, $B_t/t\to0$. 布朗运动的精确振幅介于 $\sqrt{t}$ 与 $t$ 之间. Khinchin 的重对数律给出了答案: 只比 $\sqrt{t}$ 多出一个增长极其缓慢的因子 $\sqrt{2\log\log t}$.

\begin{theorem}[重对数律]\label{thm:lil}
记 $h(t)=\sqrt{2t\log\log t}$ ($t>e$). 则
\[
\limsup_{t\to\infty}\frac{B_t}{h(t)}=1,\qquad\liminf_{t\to\infty}\frac{B_t}{h(t)}=-1\qquad\as
\]
\end{theorem}

证明分上下两半. 上界: 沿几何序列 $t_n=q^n$ ($q>1$ 接近 $1$) 控制每一段上的最大值, 反射原理把最大值的尾概率化为正态尾概率, 由 Borel-Cantelli 第一引理, 只有有限段会越界. 下界: 沿 $t_n=q^n$ ($q$ 很大) 考虑独立增量 $B_{t_n}-B_{t_{n-1}}$, 它几乎占据了 $B_{t_n}$ 的全部方差, 由 Borel-Cantelli 第二引理, 无穷多个增量会接近 $h(t_n)$. 两半都依赖于 Mills 比估计 (定理~\ref{thm:mills}) 的精确性.

\begin{proof}
函数 $t\log\log t$ 在 $(e,\infty)$ 上严格递增 (其导数为 $\log\log t+\frac{1}{\log t}>0$), 所以 $h$ 在 $(e,\infty)$ 上严格递增. 由对称性 (对 $-B$ 应用关于 $\limsup$ 的结论), 只需证明 $\limsup$ 的结论.

\emph{上界: $\limsup_{t\to\infty}B_t/h(t)\le1$ $\as$} 固定 $\varepsilon>0$, 令 $q=1+\varepsilon$, $t_n=q^n$. 令
\[
A_n=\Big\{\max_{s\le t_{n+1}}B_s\ge(1+\varepsilon)h(t_n)\Big\}.
\]
由反射原理~\ref{thm:reflection}, $\P(A_n)=2\P\big(B_{t_{n+1}}\ge(1+\varepsilon)h(t_n)\big)=2\big(1-\Phi(x_n)\big)$, 其中
\[
x_n=\frac{(1+\varepsilon)h(t_n)}{\sqrt{t_{n+1}}}=(1+\varepsilon)\sqrt{\frac{2\log\log t_n}{q}}.
\]
$x_n\to\infty$, 所以当 $n$ 充分大时 $x_n\ge1$, 由 Mills 比估计 (定理~\ref{thm:mills}), $1-\Phi(x_n)\le\varphi(x_n)/x_n\le\varphi(x_n)$. 于是
\[
\P(A_n)\le\frac{2}{\sqrt{2\pi}}\exp\Big(-\frac{(1+\varepsilon)^2}{q}\log\log t_n\Big)=\frac{2}{\sqrt{2\pi}}\big(n\log q\big)^{-(1+\varepsilon)},
\]
这里用到 $(1+\varepsilon)^2/q=1+\varepsilon$ 与 $\log t_n=n\log q$. 右端关于 $n$ 可和, 由 Borel-Cantelli 第一引理~\ref{thm:bc1}, 几乎必然地只有有限个 $A_n$ 发生. 于是存在 $N(\omega)$, 当 $n\ge N(\omega)$ 时 (取 $N$ 足够大使 $t_N>e$), 对一切 $t\in[t_n,t_{n+1}]$,
\[
B_t\le\max_{s\le t_{n+1}}B_s<(1+\varepsilon)h(t_n)\le(1+\varepsilon)h(t),
\]
最后一步用到 $h$ 的单调性. 所以 $\limsup_{t\to\infty}B_t/h(t)\le1+\varepsilon$ $\as$ 令 $\varepsilon$ 取遍 $\{1/k\}$ 并取交, 得 $\limsup\le1$ $\as$

\emph{下界: $\limsup_{t\to\infty}B_t/h(t)\ge1$ $\as$} 固定 $\varepsilon\in(0,1)$, 取 $q>1$ 充分大, 使得
\[
\beta:=\frac{(1-\varepsilon)^2}{1-1/q}<1.
\]
令 $t_n=q^n$, $\Delta_n=B_{t_n}-B_{t_{n-1}}\sim N\big(0,t_n(1-\frac{1}{q})\big)$. 由定义~\ref{def:bm} (2), $\Delta_1,\Delta_2,\dots$ 相互独立. 令
\[
D_n=\big\{\Delta_n\ge(1-\varepsilon)h(t_n)\big\},\qquad y_n=\frac{(1-\varepsilon)h(t_n)}{\sqrt{t_n(1-1/q)}}=\sqrt{2\beta\log\log t_n}.
\]
则 $\P(D_n)=1-\Phi(y_n)$. 当 $n$ 充分大时 $y_n\ge1$, 由 Mills 比估计的下界, $1-\Phi(y_n)\ge\frac{y_n}{1+y_n^2}\varphi(y_n)\ge\frac{\varphi(y_n)}{2y_n}$, 所以
\[
\P(D_n)\ge\frac{1}{2\sqrt{2\pi}}\cdot\frac{e^{-\beta\log\log t_n}}{\sqrt{2\beta\log\log t_n}}=\frac{1}{2\sqrt{2\pi}}\cdot\frac{(n\log q)^{-\beta}}{\sqrt{2\beta\log(n\log q)}}.
\]
由于 $\beta<1$, 右端的量级为 $n^{-\beta}/\sqrt{\log n}$, 当 $n$ 充分大时不小于 $\frac{1}{n}$, 所以 $\sum_n\P(D_n)=\infty$. 由 Borel-Cantelli 第二引理~\ref{thm:bc2}, 几乎必然地有无穷多个 $D_n$ 发生.

另一方面, 对 $-B$ 应用已证的上界, 几乎必然地, 当 $n$ 充分大时 $-B_{t_{n-1}}\le2h(t_{n-1})$. 而
\[
h(t_{n-1})^2=\frac{2t_n}{q}\log\log t_{n-1}\le\frac{1}{q}\cdot2t_n\log\log t_n=\frac{h(t_n)^2}{q}.
\]
所以对无穷多个 $n$,
\[
B_{t_n}=B_{t_{n-1}}+\Delta_n\ge-\frac{2}{\sqrt{q}}h(t_n)+(1-\varepsilon)h(t_n)=\Big(1-\varepsilon-\frac{2}{\sqrt{q}}\Big)h(t_n).
\]
于是 $\limsup_{t\to\infty}B_t/h(t)\ge1-\varepsilon-2/\sqrt{q}$ $\as$ 让 $\varepsilon$ 取遍 $\{1/k\}$, $q$ 取遍充分大的整数, 对可列个概率为 $1$ 的事件取交, 得 $\limsup\ge1$ $\as$
\end{proof}

\begin{figure}[htbp]
\centering
\begin{tikzpicture}
\begin{axis}[SAaxis, axis x line=bottom, axis y line=left, width=12cm, height=5.6cm, xmin=1.2, xmax=8.1, ymin=-1.45, ymax=1.45,
  xtick={2,4,6,8}, xticklabels={$10^2$,$10^4$,$10^6$,$10^8$}, ytick={-1,0,1}, xlabel={$t$ (对数刻度)}, ylabel={$B_t/h(t)$}]
  \fill[SAGreen!10] (axis cs:1.2,-1) rectangle (axis cs:8.1,1);
  \addplot[SAGreenText,thick] coordinates {(1.2,1) (8.1,1)};
  \addplot[SAGreenText,thick] coordinates {(1.2,-1) (8.1,-1)};
  \addplot[SADark!30] coordinates {(1.2,0) (8.1,0)};
  \addplot[SABlue,thin] coordinates {(1.204,-0.997) (1.209,-0.970) (1.213,-0.980) (1.218,-1.079) (1.222,-1.021) (1.227,-0.994) (1.231,-0.977) (1.236,-0.948) (1.240,-0.906) (1.245,-0.933) (1.249,-0.792) (1.254,-0.756) (1.259,-0.839) (1.263,-0.877) (1.268,-0.790) (1.272,-0.618) (1.277,-0.620) (1.281,-0.745) (1.286,-0.676) (1.290,-0.764) (1.295,-0.731) (1.299,-0.670) (1.304,-0.673) (1.308,-0.716) (1.313,-0.643) (1.317,-0.614) (1.322,-0.571) (1.326,-0.512) (1.331,-0.532) (1.336,-0.592) (1.340,-0.652) (1.345,-0.645) (1.349,-0.715) (1.354,-0.718) (1.358,-0.737) (1.363,-0.842) (1.367,-0.888) (1.372,-0.907) (1.376,-0.833) (1.381,-0.881) (1.385,-0.867) (1.390,-0.980) (1.394,-0.962) (1.399,-1.006) (1.404,-0.903) (1.408,-0.856) (1.413,-0.786) (1.417,-0.719) (1.422,-0.734) (1.426,-0.809) (1.431,-0.798) (1.435,-0.893) (1.440,-0.851) (1.444,-0.937) (1.449,-0.885) (1.453,-0.814) (1.458,-0.725) (1.462,-0.665) (1.467,-0.557) (1.472,-0.512) (1.476,-0.576) (1.481,-0.520) (1.485,-0.500) (1.490,-0.484) (1.494,-0.507) (1.499,-0.614) (1.503,-0.588) (1.508,-0.603) (1.512,-0.657) (1.517,-0.630) (1.521,-0.623) (1.526,-0.724) (1.530,-0.693) (1.535,-0.645) (1.539,-0.707) (1.544,-0.753) (1.549,-0.767) (1.553,-0.707) (1.558,-0.699) (1.562,-0.716) (1.567,-0.813) (1.571,-0.851) (1.576,-0.830) (1.580,-0.852) (1.585,-0.829) (1.589,-0.829) (1.594,-0.727) (1.598,-0.648) (1.603,-0.633) (1.607,-0.559) (1.612,-0.514) (1.617,-0.451) (1.621,-0.522) (1.626,-0.537) (1.630,-0.445) (1.635,-0.497) (1.639,-0.445) (1.644,-0.459) (1.648,-0.372) (1.653,-0.403) (1.657,-0.526) (1.662,-0.522) (1.666,-0.503) (1.671,-0.446) (1.675,-0.391) (1.680,-0.400) (1.685,-0.384) (1.689,-0.453) (1.694,-0.506) (1.698,-0.486) (1.703,-0.465) (1.707,-0.568) (1.712,-0.570) (1.716,-0.560) (1.721,-0.581) (1.725,-0.608) (1.730,-0.542) (1.734,-0.467) (1.739,-0.535) (1.743,-0.546) (1.748,-0.430) (1.753,-0.349) (1.757,-0.206) (1.762,-0.218) (1.766,-0.248) (1.771,-0.155) (1.775,-0.172) (1.780,-0.211) (1.784,-0.285) (1.789,-0.166) (1.793,-0.208) (1.798,-0.320) (1.802,-0.330) (1.807,-0.345) (1.811,-0.351) (1.816,-0.275) (1.820,-0.262) (1.825,-0.148) (1.830,-0.184) (1.834,-0.107) (1.839,-0.139) (1.843,-0.237) (1.848,-0.156) (1.852,-0.091) (1.857,-0.155) (1.861,-0.194) (1.866,-0.254) (1.870,-0.273) (1.875,-0.348) (1.879,-0.356) (1.884,-0.404) (1.888,-0.392) (1.893,-0.425) (1.898,-0.441) (1.902,-0.471) (1.907,-0.542) (1.911,-0.671) (1.916,-0.711) (1.920,-0.694) (1.925,-0.793) (1.929,-0.848) (1.934,-0.814) (1.938,-0.912) (1.943,-0.821) (1.947,-0.795) (1.952,-0.756) (1.956,-0.778) (1.961,-0.742) (1.966,-0.807) (1.970,-0.728) (1.975,-0.664) (1.979,-0.674) (1.984,-0.588) (1.988,-0.489) (1.993,-0.568) (1.997,-0.608) (2.002,-0.547) (2.006,-0.464) (2.011,-0.417) (2.015,-0.400) (2.020,-0.422) (2.024,-0.396) (2.029,-0.372) (2.033,-0.362) (2.038,-0.444) (2.043,-0.507) (2.047,-0.383) (2.052,-0.445) (2.056,-0.422) (2.061,-0.527) (2.065,-0.608) (2.070,-0.574) (2.074,-0.654) (2.079,-0.745) (2.083,-0.731) (2.088,-0.716) (2.092,-0.692) (2.097,-0.644) (2.101,-0.543) (2.106,-0.538) (2.111,-0.436) (2.115,-0.413) (2.120,-0.483) (2.124,-0.491) (2.129,-0.546) (2.133,-0.588) (2.138,-0.520) (2.142,-0.567) (2.147,-0.586) (2.151,-0.693) (2.156,-0.667) (2.160,-0.565) (2.165,-0.609) (2.169,-0.621) (2.174,-0.677) (2.179,-0.653) (2.183,-0.601) (2.188,-0.666) (2.192,-0.637) (2.197,-0.599) (2.201,-0.632) (2.206,-0.603) (2.210,-0.605) (2.215,-0.599) (2.219,-0.602) (2.224,-0.505) (2.228,-0.477) (2.233,-0.452) (2.237,-0.389) (2.242,-0.415) (2.247,-0.355) (2.251,-0.397) (2.256,-0.389) (2.260,-0.466) (2.265,-0.464) (2.269,-0.306) (2.274,-0.259) (2.278,-0.379) (2.283,-0.428) (2.287,-0.492) (2.292,-0.482) (2.296,-0.467) (2.301,-0.406) (2.305,-0.393) (2.310,-0.354) (2.314,-0.340) (2.319,-0.374) (2.324,-0.399) (2.328,-0.330) (2.333,-0.275) (2.337,-0.180) (2.342,-0.099) (2.346,-0.152) (2.351,-0.096) (2.355,-0.040) (2.360,-0.043) (2.364,-0.052) (2.369,-0.035) (2.373,-0.122) (2.378,-0.177) (2.382,-0.224) (2.387,-0.260) (2.392,-0.291) (2.396,-0.238) (2.401,-0.215) (2.405,-0.274) (2.410,-0.254) (2.414,-0.407) (2.419,-0.424) (2.423,-0.446) (2.428,-0.466) (2.432,-0.512) (2.437,-0.444) (2.441,-0.561) (2.446,-0.619) (2.450,-0.570) (2.455,-0.632) (2.460,-0.615) (2.464,-0.563) (2.469,-0.424) (2.473,-0.483) (2.478,-0.334) (2.482,-0.379) (2.487,-0.382) (2.491,-0.432) (2.496,-0.491) (2.500,-0.470) (2.505,-0.509) (2.509,-0.519) (2.514,-0.554) (2.518,-0.444) (2.523,-0.417) (2.527,-0.436) (2.532,-0.532) (2.537,-0.515) (2.541,-0.482) (2.546,-0.459) (2.550,-0.411) (2.555,-0.438) (2.559,-0.419) (2.564,-0.341) (2.568,-0.342) (2.573,-0.330) (2.577,-0.342) (2.582,-0.209) (2.586,-0.255) (2.591,-0.194) (2.595,-0.250) (2.600,-0.197) (2.605,-0.151) (2.609,-0.187) (2.614,-0.271) (2.618,-0.325) (2.623,-0.342) (2.627,-0.246) (2.632,-0.309) (2.636,-0.260) (2.641,-0.232) (2.645,-0.226) (2.650,-0.178) (2.654,-0.197) (2.659,-0.204) (2.663,-0.274) (2.668,-0.438) (2.673,-0.438) (2.677,-0.373) (2.682,-0.363) (2.686,-0.327) (2.691,-0.368) (2.695,-0.425) (2.700,-0.515) (2.704,-0.445) (2.709,-0.433) (2.713,-0.398) (2.718,-0.365) (2.722,-0.371) (2.727,-0.329) (2.731,-0.220) (2.736,-0.161) (2.741,-0.078) (2.745,-0.119) (2.750,-0.180) (2.754,-0.191) (2.759,-0.252) (2.763,-0.234) (2.768,-0.192) (2.772,-0.258) (2.777,-0.243) (2.781,-0.286) (2.786,-0.313) (2.790,-0.351) (2.795,-0.295) (2.799,-0.278) (2.804,-0.170) (2.808,-0.268) (2.813,-0.163) (2.818,-0.202) (2.822,-0.161) (2.827,-0.114) (2.831,-0.164) (2.836,-0.131) (2.840,-0.107) (2.845,0.000) (2.849,-0.052) (2.854,-0.060) (2.858,-0.104) (2.863,-0.138) (2.867,-0.095) (2.872,-0.042) (2.876,-0.048) (2.881,-0.045) (2.886,-0.085) (2.890,-0.048) (2.895,-0.011) (2.899,0.004) (2.904,-0.002) (2.908,0.053) (2.913,0.025) (2.917,0.040) (2.922,0.015) (2.926,0.054) (2.931,0.139) (2.935,0.150) (2.940,0.080) (2.944,0.040) (2.949,0.090) (2.954,0.122) (2.958,0.121) (2.963,0.001) (2.967,-0.002) (2.972,-0.006) (2.976,0.052) (2.981,0.070) (2.985,-0.014) (2.990,0.030) (2.994,-0.044) (2.999,0.026) (3.003,0.015) (3.008,0.010) (3.012,0.131) (3.017,0.144) (3.021,0.216) (3.026,0.205) (3.031,0.202) (3.035,0.207) (3.040,0.230) (3.044,0.209) (3.049,0.220) (3.053,0.267) (3.058,0.156) (3.062,0.174) (3.067,0.221) (3.071,0.258) (3.076,0.256) (3.080,0.249) (3.085,0.238) (3.089,0.229) (3.094,0.234) (3.099,0.327) (3.103,0.363) (3.108,0.343) (3.112,0.316) (3.117,0.470) (3.121,0.456) (3.126,0.425) (3.130,0.434) (3.135,0.446) (3.139,0.374) (3.144,0.406) (3.148,0.403) (3.153,0.279) (3.157,0.364) (3.162,0.459) (3.167,0.392) (3.171,0.422) (3.176,0.472) (3.180,0.359) (3.185,0.284) (3.189,0.356) (3.194,0.326) (3.198,0.362) (3.203,0.404) (3.207,0.396) (3.212,0.449) (3.216,0.490) (3.221,0.500) (3.225,0.515) (3.230,0.534) (3.234,0.502) (3.239,0.469) (3.244,0.449) (3.248,0.461) (3.253,0.484) (3.257,0.441) (3.262,0.402) (3.266,0.441) (3.271,0.388) (3.275,0.401) (3.280,0.441) (3.284,0.387) (3.289,0.380) (3.293,0.390) (3.298,0.368) (3.302,0.335) (3.307,0.326) (3.312,0.285) (3.316,0.314) (3.321,0.376) (3.325,0.401) (3.330,0.384) (3.334,0.324) (3.339,0.311) (3.343,0.437) (3.348,0.442) (3.352,0.419) (3.357,0.419) (3.361,0.418) (3.366,0.419) (3.370,0.508) (3.375,0.450) (3.380,0.375) (3.384,0.407) (3.389,0.399) (3.393,0.424) (3.398,0.333) (3.402,0.205) (3.407,0.281) (3.411,0.289) (3.416,0.333) (3.420,0.382) (3.425,0.304) (3.429,0.249) (3.434,0.180) (3.438,0.154) (3.443,0.042) (3.448,0.066) (3.452,0.133) (3.457,0.236) (3.461,0.253) (3.466,0.248) (3.470,0.197) (3.475,0.263) (3.479,0.101) (3.484,0.046) (3.488,0.096) (3.493,0.149) (3.497,0.187) (3.502,0.218) (3.506,0.266) (3.511,0.335) (3.515,0.321) (3.520,0.292) (3.525,0.324) (3.529,0.262) (3.534,0.244) (3.538,0.225) (3.543,0.244) (3.547,0.255) (3.552,0.276) (3.556,0.283) (3.561,0.259) (3.565,0.231) (3.570,0.171) (3.574,0.193) (3.579,0.181) (3.583,0.224) (3.588,0.234) (3.593,0.121) (3.597,0.229) (3.602,0.221) (3.606,0.237) (3.611,0.334) (3.615,0.332) (3.620,0.359) (3.624,0.337) (3.629,0.351) (3.633,0.412) (3.638,0.521) (3.642,0.539) (3.647,0.493) (3.651,0.546) (3.656,0.548) (3.661,0.565) (3.665,0.669) (3.670,0.644) (3.674,0.639) (3.679,0.616) (3.683,0.591) (3.688,0.631) (3.692,0.621) (3.697,0.556) (3.701,0.495) (3.706,0.454) (3.710,0.532) (3.715,0.621) (3.719,0.701) (3.724,0.641) (3.728,0.562) (3.733,0.598) (3.738,0.650) (3.742,0.698) (3.747,0.675) (3.751,0.728) (3.756,0.806) (3.760,0.806) (3.765,0.846) (3.769,0.900) (3.774,0.944) (3.778,0.968) (3.783,0.949) (3.787,0.898) (3.792,0.916) (3.796,0.963) (3.801,1.007) (3.806,1.017) (3.810,1.005) (3.815,1.006) (3.819,1.016) (3.824,1.141) (3.828,1.169) (3.833,1.141) (3.837,1.127) (3.842,1.119) (3.846,1.061) (3.851,1.118) (3.855,1.100) (3.860,1.056) (3.864,1.066) (3.869,1.000) (3.874,1.036) (3.878,1.066) (3.883,1.047) (3.887,1.131) (3.892,1.113) (3.896,1.117) (3.901,1.148) (3.905,1.182) (3.910,1.134) (3.914,1.112) (3.919,1.103) (3.923,1.069) (3.928,1.130) (3.932,1.088) (3.937,1.109) (3.942,1.067) (3.946,1.085) (3.951,1.095) (3.955,1.168) (3.960,1.085) (3.964,0.984) (3.969,1.008) (3.973,1.066) (3.978,0.996) (3.982,0.944) (3.987,0.935) (3.991,0.867) (3.996,0.849) (4.000,0.903) (4.005,0.879) (4.009,0.909) (4.014,0.954) (4.019,0.907) (4.023,0.939) (4.028,0.941) (4.032,1.000) (4.037,1.024) (4.041,1.034) (4.046,1.054) (4.050,1.158) (4.055,1.108) (4.059,1.097) (4.064,1.176) (4.068,1.262) (4.073,1.228) (4.077,1.238) (4.082,1.187) (4.087,1.188) (4.091,1.169) (4.096,1.193) (4.100,1.241) (4.105,1.164) (4.109,1.175) (4.114,1.235) (4.118,1.155) (4.123,1.181) (4.127,1.141) (4.132,1.193) (4.136,1.174) (4.141,1.043) (4.145,1.055) (4.150,1.100) (4.155,1.050) (4.159,1.061) (4.164,1.053) (4.168,1.087) (4.173,1.052) (4.177,1.072) (4.182,1.024) (4.186,1.036) (4.191,1.068) (4.195,1.122) (4.200,1.163) (4.204,1.231) (4.209,1.219) (4.213,1.159) (4.218,1.131) (4.222,1.192) (4.227,1.138) (4.232,1.045) (4.236,1.006) (4.241,1.078) (4.245,1.070) (4.250,1.013) (4.254,1.001) (4.259,1.004) (4.263,1.094) (4.268,1.093) (4.272,1.146) (4.277,1.150) (4.281,1.075) (4.286,1.031) (4.290,1.052) (4.295,0.991) (4.300,1.001) (4.304,0.977) (4.309,0.932) (4.313,0.983) (4.318,0.947) (4.322,0.922) (4.327,0.987) (4.331,0.931) (4.336,0.848) (4.340,0.905) (4.345,0.965) (4.349,1.018) (4.354,1.059) (4.358,1.045) (4.363,1.078) (4.368,1.116) (4.372,1.090) (4.377,1.079) (4.381,1.081) (4.386,1.112) (4.390,1.033) (4.395,1.041) (4.399,1.031) (4.404,1.048) (4.408,1.023) (4.413,0.983) (4.417,0.945) (4.422,0.843) (4.426,0.878) (4.431,0.830) (4.436,0.735) (4.440,0.785) (4.445,0.784) (4.449,0.685) (4.454,0.755) (4.458,0.782) (4.463,0.769) (4.467,0.864) (4.472,0.913) (4.476,0.918) (4.481,0.909) (4.485,0.872) (4.490,0.845) (4.494,0.931) (4.499,0.938) (4.503,0.931) (4.508,0.951) (4.513,0.954) (4.517,0.908) (4.522,0.938) (4.526,0.881) (4.531,0.907) (4.535,0.893) (4.540,0.922) (4.544,0.870) (4.549,0.927) (4.553,0.968) (4.558,0.942) (4.562,0.880) (4.567,0.943) (4.571,0.903) (4.576,0.940) (4.581,0.945) (4.585,0.946) (4.590,0.943) (4.594,0.960) (4.599,0.959) (4.603,0.910) (4.608,0.925) (4.612,0.858) (4.617,0.894) (4.621,0.880) (4.626,0.845) (4.630,0.787) (4.635,0.892) (4.639,0.935) (4.644,0.825) (4.649,0.772) (4.653,0.814) (4.658,0.837) (4.662,0.855) (4.667,0.805) (4.671,0.742) (4.676,0.722) (4.680,0.650) (4.685,0.691) (4.689,0.713) (4.694,0.688) (4.698,0.644) (4.703,0.665) (4.707,0.718) (4.712,0.728) (4.716,0.805) (4.721,0.746) (4.726,0.736) (4.730,0.713) (4.735,0.701) (4.739,0.762) (4.744,0.779) (4.748,0.780) (4.753,0.810) (4.757,0.904) (4.762,0.903) (4.766,0.922) (4.771,0.861) (4.775,0.881) (4.780,0.910) (4.784,0.847) (4.789,0.880) (4.794,0.894) (4.798,0.917) (4.803,0.880) (4.807,0.867) (4.812,0.903) (4.816,0.822) (4.821,0.798) (4.825,0.723) (4.830,0.694) (4.834,0.731) (4.839,0.708) (4.843,0.677) (4.848,0.665) (4.852,0.635) (4.857,0.669) (4.862,0.626) (4.866,0.696) (4.871,0.736) (4.875,0.781) (4.880,0.758) (4.884,0.740) (4.889,0.692) (4.893,0.725) (4.898,0.577) (4.902,0.579) (4.907,0.604) (4.911,0.672) (4.916,0.656) (4.920,0.689) (4.925,0.617) (4.930,0.617) (4.934,0.577) (4.939,0.539) (4.943,0.564) (4.948,0.540) (4.952,0.567) (4.957,0.545) (4.961,0.564) (4.966,0.575) (4.970,0.660) (4.975,0.573) (4.979,0.495) (4.984,0.517) (4.988,0.542) (4.993,0.501) (4.997,0.491) (5.002,0.510) (5.007,0.488) (5.011,0.540) (5.016,0.476) (5.020,0.433) (5.025,0.463) (5.029,0.493) (5.034,0.562) (5.038,0.539) (5.043,0.536) (5.047,0.527) (5.052,0.474) (5.056,0.505) (5.061,0.442) (5.065,0.475) (5.070,0.437) (5.075,0.396) (5.079,0.405) (5.084,0.394) (5.088,0.435) (5.093,0.493) (5.097,0.558) (5.102,0.587) (5.106,0.563) (5.111,0.595) (5.115,0.594) (5.120,0.591) (5.124,0.525) (5.129,0.478) (5.133,0.513) (5.138,0.522) (5.143,0.544) (5.147,0.580) (5.152,0.545) (5.156,0.436) (5.161,0.381) (5.165,0.439) (5.170,0.438) (5.174,0.444) (5.179,0.432) (5.183,0.441) (5.188,0.450) (5.192,0.431) (5.197,0.384) (5.201,0.348) (5.206,0.322) (5.210,0.348) (5.215,0.331) (5.220,0.323) (5.224,0.342) (5.229,0.403) (5.233,0.357) (5.238,0.367) (5.242,0.360) (5.247,0.393) (5.251,0.365) (5.256,0.372) (5.260,0.395) (5.265,0.437) (5.269,0.371) (5.274,0.397) (5.278,0.346) (5.283,0.336) (5.288,0.303) (5.292,0.273) (5.297,0.249) (5.301,0.218) (5.306,0.187) (5.310,0.143) (5.315,0.108) (5.319,0.106) (5.324,0.194) (5.328,0.242) (5.333,0.211) (5.337,0.105) (5.342,-0.006) (5.346,-0.088) (5.351,-0.101) (5.356,-0.044) (5.360,0.016) (5.365,0.008) (5.369,0.014) (5.374,0.020) (5.378,-0.067) (5.383,0.000) (5.387,0.155) (5.392,0.089) (5.396,-0.012) (5.401,0.099) (5.405,0.056) (5.410,-0.013) (5.414,-0.031) (5.419,-0.020) (5.424,0.046) (5.428,0.036) (5.433,0.050) (5.437,0.015) (5.442,0.021) (5.446,0.014) (5.451,-0.010) (5.455,0.002) (5.460,-0.007) (5.464,0.033) (5.469,0.021) (5.473,0.067) (5.478,0.094) (5.482,0.045) (5.487,0.044) (5.491,-0.067) (5.496,-0.046) (5.501,0.023) (5.505,0.051) (5.510,0.121) (5.514,0.170) (5.519,0.090) (5.523,0.152) (5.528,0.095) (5.532,0.138) (5.537,0.064) (5.541,0.061) (5.546,0.113) (5.550,0.097) (5.555,0.101) (5.559,0.056) (5.564,0.020) (5.569,-0.002) (5.573,-0.076) (5.578,-0.105) (5.582,-0.055) (5.587,-0.045) (5.591,-0.134) (5.596,-0.129) (5.600,-0.152) (5.605,-0.250) (5.609,-0.183) (5.614,-0.312) (5.618,-0.298) (5.623,-0.208) (5.627,-0.169) (5.632,-0.171) (5.637,-0.177) (5.641,-0.201) (5.646,-0.204) (5.650,-0.281) (5.655,-0.287) (5.659,-0.353) (5.664,-0.347) (5.668,-0.373) (5.673,-0.355) (5.677,-0.325) (5.682,-0.279) (5.686,-0.169) (5.691,-0.232) (5.695,-0.249) (5.700,-0.227) (5.704,-0.276) (5.709,-0.262) (5.714,-0.185) (5.718,-0.221) (5.723,-0.223) (5.727,-0.224) (5.732,-0.195) (5.736,-0.158) (5.741,-0.197) (5.745,-0.245) (5.750,-0.275) (5.754,-0.231) (5.759,-0.210) (5.763,-0.195) (5.768,-0.174) (5.772,-0.192) (5.777,-0.129) (5.782,-0.161) (5.786,-0.236) (5.791,-0.170) (5.795,-0.188) (5.800,-0.098) (5.804,-0.068) (5.809,-0.108) (5.813,-0.163) (5.818,-0.244) (5.822,-0.182) (5.827,-0.165) (5.831,-0.228) (5.836,-0.298) (5.840,-0.287) (5.845,-0.288) (5.850,-0.286) (5.854,-0.274) (5.859,-0.232) (5.863,-0.269) (5.868,-0.255) (5.872,-0.263) (5.877,-0.165) (5.881,-0.100) (5.886,-0.153) (5.890,-0.147) (5.895,-0.171) (5.899,-0.253) (5.904,-0.272) (5.908,-0.328) (5.913,-0.370) (5.918,-0.352) (5.922,-0.337) (5.927,-0.349) (5.931,-0.309) (5.936,-0.289) (5.940,-0.230) (5.945,-0.289) (5.949,-0.291) (5.954,-0.199) (5.958,-0.208) (5.963,-0.284) (5.967,-0.298) (5.972,-0.301) (5.976,-0.314) (5.981,-0.247) (5.985,-0.200) (5.990,-0.211) (5.995,-0.287) (5.999,-0.304) (6.004,-0.322) (6.008,-0.302) (6.013,-0.293) (6.017,-0.247) (6.022,-0.211) (6.026,-0.264) (6.031,-0.291) (6.035,-0.267) (6.040,-0.224) (6.044,-0.229) (6.049,-0.246) (6.053,-0.259) (6.058,-0.230) (6.063,-0.186) (6.067,-0.208) (6.072,-0.163) (6.076,-0.110) (6.081,-0.092) (6.085,-0.047) (6.090,-0.041) (6.094,-0.020) (6.099,-0.010) (6.103,-0.053) (6.108,-0.073) (6.112,-0.104) (6.117,-0.069) (6.121,0.017) (6.126,0.007) (6.131,-0.038) (6.135,-0.032) (6.140,-0.043) (6.144,-0.021) (6.149,-0.017) (6.153,-0.007) (6.158,-0.037) (6.162,-0.042) (6.167,-0.066) (6.171,-0.024) (6.176,-0.039) (6.180,-0.054) (6.185,-0.075) (6.189,-0.113) (6.194,-0.008) (6.198,-0.054) (6.203,-0.058) (6.208,-0.086) (6.212,-0.041) (6.217,-0.125) (6.221,-0.103) (6.226,-0.079) (6.230,-0.094) (6.235,-0.111) (6.239,-0.151) (6.244,-0.119) (6.248,-0.157) (6.253,-0.221) (6.257,-0.186) (6.262,-0.180) (6.266,-0.239) (6.271,-0.240) (6.276,-0.243) (6.280,-0.313) (6.285,-0.279) (6.289,-0.319) (6.294,-0.334) (6.298,-0.280) (6.303,-0.197) (6.307,-0.148) (6.312,-0.174) (6.316,-0.196) (6.321,-0.135) (6.325,-0.188) (6.330,-0.193) (6.334,-0.181) (6.339,-0.167) (6.344,-0.268) (6.348,-0.259) (6.353,-0.258) (6.357,-0.215) (6.362,-0.196) (6.366,-0.128) (6.371,-0.116) (6.375,-0.027) (6.380,-0.018) (6.384,0.059) (6.389,0.033) (6.393,0.035) (6.398,0.025) (6.402,0.054) (6.407,0.140) (6.411,0.037) (6.416,0.038) (6.421,0.010) (6.425,0.033) (6.430,0.027) (6.434,-0.011) (6.439,0.102) (6.443,0.101) (6.448,0.238) (6.452,0.260) (6.457,0.260) (6.461,0.113) (6.466,0.124) (6.470,0.142) (6.475,0.172) (6.479,0.183) (6.484,0.210) (6.489,0.170) (6.493,0.191) (6.498,0.270) (6.502,0.284) (6.507,0.361) (6.511,0.397) (6.516,0.426) (6.520,0.441) (6.525,0.434) (6.529,0.395) (6.534,0.485) (6.538,0.511) (6.543,0.466) (6.547,0.476) (6.552,0.401) (6.557,0.375) (6.561,0.409) (6.566,0.423) (6.570,0.455) (6.575,0.445) (6.579,0.483) (6.584,0.445) (6.588,0.450) (6.593,0.399) (6.597,0.403) (6.602,0.448) (6.606,0.390) (6.611,0.399) (6.615,0.389) (6.620,0.437) (6.625,0.437) (6.629,0.440) (6.634,0.433) (6.638,0.409) (6.643,0.379) (6.647,0.385) (6.652,0.421) (6.656,0.481) (6.661,0.431) (6.665,0.451) (6.670,0.433) (6.674,0.435) (6.679,0.448) (6.683,0.384) (6.688,0.352) (6.692,0.306) (6.697,0.257) (6.702,0.223) (6.706,0.141) (6.711,0.110) (6.715,0.201) (6.720,0.246) (6.724,0.334) (6.729,0.398) (6.733,0.414) (6.738,0.402) (6.742,0.374) (6.747,0.464) (6.751,0.511) (6.756,0.521) (6.760,0.552) (6.765,0.560) (6.770,0.507) (6.774,0.517) (6.779,0.590) (6.783,0.599) (6.788,0.634) (6.792,0.627) (6.797,0.533) (6.801,0.460) (6.806,0.435) (6.810,0.498) (6.815,0.453) (6.819,0.405) (6.824,0.442) (6.828,0.427) (6.833,0.491) (6.838,0.504) (6.842,0.414) (6.847,0.476) (6.851,0.390) (6.856,0.356) (6.860,0.268) (6.865,0.355) (6.869,0.314) (6.874,0.303) (6.878,0.353) (6.883,0.305) (6.887,0.278) (6.892,0.246) (6.896,0.248) (6.901,0.265) (6.905,0.205) (6.910,0.200) (6.915,0.243) (6.919,0.218) (6.924,0.179) (6.928,0.166) (6.933,0.172) (6.937,0.181) (6.942,0.211) (6.946,0.230) (6.951,0.303) (6.955,0.219) (6.960,0.148) (6.964,0.203) (6.969,0.168) (6.973,0.204) (6.978,0.196) (6.983,0.234) (6.987,0.123) (6.992,0.082) (6.996,0.021) (7.001,0.026) (7.005,-0.038) (7.010,-0.096) (7.014,-0.055) (7.019,-0.004) (7.023,0.034) (7.028,0.047) (7.032,0.001) (7.037,-0.068) (7.041,-0.046) (7.046,-0.041) (7.051,-0.080) (7.055,-0.156) (7.060,-0.121) (7.064,-0.109) (7.069,-0.115) (7.073,-0.074) (7.078,-0.124) (7.082,-0.164) (7.087,-0.243) (7.091,-0.241) (7.096,-0.270) (7.100,-0.216) (7.105,-0.172) (7.109,-0.211) (7.114,-0.213) (7.119,-0.255) (7.123,-0.336) (7.128,-0.362) (7.132,-0.270) (7.137,-0.223) (7.141,-0.296) (7.146,-0.287) (7.150,-0.222) (7.155,-0.148) (7.159,-0.106) (7.164,-0.088) (7.168,-0.071) (7.173,-0.077) (7.177,-0.033) (7.182,-0.002) (7.186,-0.017) (7.191,0.035) (7.196,0.033) (7.200,0.050) (7.205,0.078) (7.209,0.066) (7.214,0.020) (7.218,0.053) (7.223,0.077) (7.227,0.117) (7.232,0.085) (7.236,0.103) (7.241,0.174) (7.245,0.161) (7.250,0.131) (7.254,0.158) (7.259,0.126) (7.264,0.154) (7.268,0.145) (7.273,0.089) (7.277,0.131) (7.282,0.049) (7.286,0.032) (7.291,0.047) (7.295,0.084) (7.300,0.111) (7.304,0.116) (7.309,0.087) (7.313,0.168) (7.318,0.114) (7.322,0.117) (7.327,0.169) (7.332,0.134) (7.336,0.100) (7.341,0.076) (7.345,0.045) (7.350,0.016) (7.354,0.019) (7.359,0.040) (7.363,-0.010) (7.368,-0.039) (7.372,0.029) (7.377,0.030) (7.381,0.051) (7.386,0.073) (7.390,0.064) (7.395,0.113) (7.399,0.110) (7.404,0.146) (7.409,0.221) (7.413,0.266) (7.418,0.294) (7.422,0.273) (7.427,0.266) (7.431,0.231) (7.436,0.142) (7.440,0.126) (7.445,0.130) (7.449,0.139) (7.454,0.026) (7.458,0.084) (7.463,0.098) (7.467,0.082) (7.472,0.035) (7.477,0.019) (7.481,0.054) (7.486,0.031) (7.490,0.058) (7.495,-0.003) (7.499,-0.055) (7.504,-0.007) (7.508,0.017) (7.513,0.013) (7.517,-0.017) (7.522,0.026) (7.526,-0.010) (7.531,0.009) (7.535,-0.025) (7.540,-0.016) (7.545,-0.026) (7.549,-0.057) (7.554,-0.022) (7.558,-0.087) (7.563,-0.132) (7.567,-0.172) (7.572,-0.185) (7.576,-0.238) (7.581,-0.268) (7.585,-0.236) (7.590,-0.232) (7.594,-0.253) (7.599,-0.278) (7.603,-0.332) (7.608,-0.365) (7.613,-0.416) (7.617,-0.460) (7.622,-0.340) (7.626,-0.341) (7.631,-0.340) (7.635,-0.342) (7.640,-0.352) (7.644,-0.311) (7.649,-0.296) (7.653,-0.293) (7.658,-0.310) (7.662,-0.288) (7.667,-0.286) (7.671,-0.269) (7.676,-0.289) (7.680,-0.313) (7.685,-0.289) (7.690,-0.280) (7.694,-0.305) (7.699,-0.275) (7.703,-0.250) (7.708,-0.210) (7.712,-0.184) (7.717,-0.225) (7.721,-0.210) (7.726,-0.252) (7.730,-0.242) (7.735,-0.248) (7.739,-0.202) (7.744,-0.194) (7.748,-0.213) (7.753,-0.197) (7.758,-0.245) (7.762,-0.223) (7.767,-0.226) (7.771,-0.261) (7.776,-0.343) (7.780,-0.288) (7.785,-0.315) (7.789,-0.267) (7.794,-0.254) (7.798,-0.253) (7.803,-0.204) (7.807,-0.186) (7.812,-0.129) (7.816,-0.162) (7.821,-0.203) (7.826,-0.178) (7.830,-0.166) (7.835,-0.182) (7.839,-0.180) (7.844,-0.208) (7.848,-0.201) (7.853,-0.276) (7.857,-0.252) (7.862,-0.301) (7.866,-0.351) (7.871,-0.332) (7.875,-0.319) (7.880,-0.293) (7.884,-0.302) (7.889,-0.282) (7.893,-0.211) (7.898,-0.259) (7.903,-0.226) (7.907,-0.301) (7.912,-0.305) (7.916,-0.328) (7.921,-0.314) (7.925,-0.403) (7.930,-0.349) (7.934,-0.247) (7.939,-0.208) (7.943,-0.172) (7.948,-0.083) (7.952,-0.073) (7.957,-0.092) (7.961,-0.127) (7.966,-0.147) (7.971,-0.047) (7.975,-0.020) (7.980,0.055) (7.984,0.096) (7.989,0.119) (7.993,0.131) (7.998,0.084)};
\end{axis}
\end{tikzpicture}
\caption{一条布朗轨道的 $B_t/\sqrt{2t\log\log t}$, 时间跨度 $t\in[16,10^8]$ (横轴为对数刻度). 比值在 $[-1,1]$ 附近缓慢振荡; 由于 $\log\log t$ 增长极慢, 收敛到重对数律所描述的行为也极慢: 在有限时间内比值仍可能超出 $\pm1$}
\label{fig:lil}
\end{figure}

\begin{corollary}\label{cor:lil-cluster}
几乎必然地, $\{B_t/h(t): t\to\infty\}$ 的极限点集恰为 $[-1,1]$; 并且 (局部重对数律)
\[
\limsup_{t\downarrow0}\frac{B_t}{\sqrt{2t\log\log(1/t)}}=1,\qquad\liminf_{t\downarrow0}\frac{B_t}{\sqrt{2t\log\log(1/t)}}=-1\qquad\as
\]
\end{corollary}

\begin{proof}
第一个结论: $t\mapsto B_t/h(t)$ 在 $(e,\infty)$ 上连续, 由重对数律, 它无穷多次接近 $1$ 又无穷多次接近 $-1$, 由介值定理, 它无穷多次经过 $[-1,1]$ 中的每一个值, 所以每个 $c\in[-1,1]$ 都是极限点; 而 $[-1,1]$ 之外没有极限点.

局部形式: 由时间反演 (定理~\ref{thm:bm-invariance} (4)), $X_u=uB_{1/u}$ 是标准布朗运动. 令 $u=1/t$, 则 $B_t=tX_{1/t}=X_u/u$, 且
\[
\frac{B_t}{\sqrt{2t\log\log(1/t)}}=\frac{X_u/u}{\sqrt{2u^{-1}\log\log u}}=\frac{X_u}{\sqrt{2u\log\log u}}.
\]
$t\downarrow0$ 等价于 $u\to\infty$, 对 $X$ 应用重对数律即得.
\end{proof}

由马尔可夫性, 局部重对数律对任意\emph{固定}时刻 $s$ 处的增量 $B_{s+h}-B_s$ 也成立. 但它不能对所有 $s$ 同时成立: 由 L\'evy 连续模定理 (4.3 节末的注), 存在 (随机的) 时刻, 使增量的振幅达到更大的 $\sqrt{2h\log(1/h)}$. 这与处处不可微的论证中 ``对 $n$ 个可能位置求并'' 的代价是同一回事.

\begin{intuition}
为什么是 $\sqrt{2t\log\log t}$? 在几何序列 $t_n=q^n$ 上, 各段 $[t_{n-1},t_n]$ 上的增量 (近似) 独立, 每一段的标准化增量都是标准正态的. 正态变量超过 $x$ 的概率约为 $e^{-x^2/2}$, 要使无穷多段中有一段超过 $x_n$, 需要 $\sum_ne^{-x_n^2/2}=\infty$. 取 $x_n=\sqrt{2c\log n}$, 则 $e^{-x_n^2/2}=n^{-c}$, 临界值是 $c=1$. 而 $\log n\approx\log\log t_n$. 所以重对数律中的 ``$2\log\log$'' 恰好是几何时间尺度上 ``独立正态序列的最大值'' 的增长速度 (参见例~\ref{ex:normal-max}): 时间呈几何增长, 有效的独立试验次数只有对数多个, 最大值的增长因此只有 ``对数的对数'' 那么慢.
\end{intuition}

\begin{remark}
由 Donsker 不变原理的强形式 (强逼近定理), 重对数律对满足 $\E\xi_1=0$, $\Var(\xi_1)=1$ 的随机游动同样成立: $\limsup_nS_n/\sqrt{2n\log\log n}=1$ $\as$ (Hartman-Wintner 定理). 这是独立随机变量部分和的精确振幅, 介于中心极限定理 ($\sqrt{n}$) 与强大数定律 ($n$) 之间. Strassen 进一步证明了泛函形式: 轨道 $\{B_{nt}/h(n)\}_{t\in[0,1]}$ 的极限点集恰为满足 $\int_0^1f'(t)^2\dd t\le1$ 的函数 $f$ 全体 (见~\cite{morters} 第 5 章).
\end{remark}

% ---------------------------------------------------------------------
\subsection*{本章小结}
\addcontentsline{toc}{subsection}{本章小结}

\begin{center}
\begin{minipage}{\textwidth}
\centering
\captionof{table}{布朗运动的主要结论}
\label{tab:ch4-summary}
\renewcommand{\arraystretch}{1.5}
\begin{tabular}{>{\raggedright\arraybackslash}p{0.24\textwidth} >{\raggedright\arraybackslash}p{0.44\textwidth} >{\raggedright\arraybackslash}p{0.22\textwidth}}
\toprule
结论 & 内容 & 后续应用 \\
\midrule
Donsker 不变原理 & 缩放随机游动 $\to$ 布朗运动 & 随机游动问题的极限 \\
存在性 (K-C 定理) & 连续修正, $\gamma<\frac12$ 阶 H\"older 连续 & 构造 SDE 的连续解 \\
不变性 & 对称、尺度、平移、时间反演 & 化归计算 \\
二次变差 & $\ip{B}_t=t$, 即 $(\dd B_t)^2=\dd t$ & It\^o 公式 (第~5~章) \\
处处不可微 & 无穷变差, $\dd B_t/\dd t$ 不存在 & 必须定义随机积分 \\
基本鞅 & $B_t$, $B_t^2-t$, $e^{\theta B_t-\theta^2t/2}$ & L\'evy 刻画, Girsanov 定理 \\
(强) 马尔可夫性 & 在停时处重新开始 & 扩散过程, 热方程 (第~6~章) \\
反射原理 & $M_t\overset{d}{=}\abs{B_t}$, $\tau_a\overset{d}{=}a^2/Z^2$ & 障碍期权定价 \\
反正弦律 & 最后零点的密度 $\frac{1}{\pi\sqrt{s(1-s)}}$ & 领先变换的规律 \\
重对数律 & $\limsup B_t/\sqrt{2t\log\log t}=1$ & 轨道的精确振幅 \\
\bottomrule
\end{tabular}
\end{minipage}
\end{center}

% =====================================================================
%  第 5 章  随机微积分
% =====================================================================
\chapter{随机微积分}
\label{ch:5}

普通微积分的核心是两件事: 积分 $\int f\dd g$ 与链式法则 $\dd f(g)=f'(g)\dd g$. 本章为布朗运动建立同样的两件事. 如第~4~章所示, 布朗轨道处处不可微 (定理~\ref{thm:nowhere-diff}), 在任何区间上都是无穷变差的 (推论~\ref{cor:bm-infinite-variation}), 所以 $\int H\dd B$ 无法按 Riemann-Stieltjes 积分逐条轨道地定义; 而二次变差 $\ip{B}_t=t$ (定理~\ref{thm:bm-qv}) 又意味着 Taylor 展开中的二阶项不能略去, 普通的链式法则必须修正.

It\^o 的解决办法是把积分定义为 ``不偷看未来'' 的 Riemann 和的 $L^2$ 极限. 这个定义有两个支柱: 一是第~3~章的鞅变换 (定义~\ref{def:martingale-transform}), 它保证积分是鞅; 二是 It\^o 等距, 它把随机积分的 $L^2$ 范数化为普通积分, 从而可以用 Hilbert 空间的完备性完成极限过程. 有了积分, Taylor 展开加上 $(\dd B_t)^2=\dd t$ 就给出 It\^o 公式, 它是随机分析中最重要的计算工具.

本章的结构如下. 5.1 节用 $\int B\dd B$ 的例子说明为什么必须仔细选择 Riemann 和中的取值点; 5.2 节构造 It\^o 积分; 5.3 节证明 It\^o 公式并给出一系列应用; 5.4 节以 It\^o 公式为基础, 介绍 L\'evy 刻画、Girsanov 定理与鞅表示定理, 并引入局部时. 本章中 $(\Omega,\mathcal{F},\P)$ 上的信息流 $\{\mathcal{F}_t\}$ 满足通常条件, $B$ 是 $\{\mathcal{F}_t\}$-布朗运动 (定义~\ref{def:Ft-bm}), $T>0$ 是固定的时间范围.

% =====================================================================
\section{随机积分的动机}
\label{sec:5-1}

% ---------------------------------------------------------------------
\subsection{从常微分方程到随机微分方程}

考虑一个受随机扰动的动力系统. 若没有扰动, 状态 $X_t$ 满足常微分方程 $\frac{\dd X_t}{\dd t}=b(X_t)$. 加上扰动后, 一个自然的模型是
\begin{equation}\label{eq:sde-formal}
\frac{\dd X_t}{\dd t}=b(X_t)+\sigma(X_t)\,\xi_t,
\end{equation}
其中 $\xi_t$ 是 ``白噪声'': 不同时刻相互独立、均值为零的随机扰动. 第~4~章的直观 (4.3.3 小节末) 已经指出, 白噪声应当是 $\frac{\dd B_t}{\dd t}$, 但这个导数并不存在. 于是把 \eqref{eq:sde-formal} 改写为积分形式
\[
X_t=X_0+\int_0^tb(X_s)\dd s+\int_0^t\sigma(X_s)\dd B_s.
\]
第一个积分是普通的 Lebesgue 积分. 本章的首要任务是赋予第二个积分以意义. 最简单的情形是 $\sigma\equiv1$, 此时 $\int_0^t\dd B_s=B_t$; 困难出现在被积函数随机且依赖于 $B$ 本身的时候, 例如 $\int_0^tB_s\dd B_s$.

% ---------------------------------------------------------------------
\subsection{Riemann-Stieltjes 积分为何失效}

对确定的函数, 若 $g$ 在 $[0,t]$ 上连续, $f$ 是有界变差函数, 则 Riemann-Stieltjes 和
\[
\sum_ig(t_i^*)\big(f(t_i)-f(t_{i-1})\big),\qquad t_i^*\in[t_{i-1},t_i]
\]
当分割的步长趋于零时收敛, 且极限与取值点 $t_i^*$ 的选择无关. 换一种取值点, 和式的变化不超过
\[
\max_i\sup_{\abs{u-v}\le\abs{\Pi}}\abs{g(u)-g(v)}\cdot\sum_i\abs{f(t_i)-f(t_{i-1})},
\]
第一个因子由 $g$ 的一致连续性趋于零, 第二个因子不超过 $f$ 的全变差. 对布朗运动, 第二个因子趋于无穷 (推论~\ref{cor:bm-infinite-variation}), 这个论证失效. 这种失效源于布朗运动本身的性质. 由一致有界原理可以证明, 若对 \emph{一切} 连续函数 $g$, 和式 $\sum_ig(t_{i-1})(f(t_i)-f(t_{i-1}))$ 都收敛, 则 $f$ 必是有界变差的 (见~\cite{revuz} 第 IV 章开头). 下面的例子说明, 对布朗运动, 取值点的选择确实会改变极限.

% ---------------------------------------------------------------------
\subsection{\texorpdfstring{例子: $\int_0^tB\dd B$}{例子: ∫B dB}}

设 $\Pi_n: 0=t_0<t_1<\cdots<t_m=t$ 是 $[0,t]$ 的分割序列 (分点个数 $m$ 依赖于 $n$), $\abs{\Pi_n}\to0$. 记 $\Delta B_i=B_{t_i}-B_{t_{i-1}}$, 并考虑三种 Riemann 和:
\[
L_n=\sum_iB_{t_{i-1}}\Delta B_i,\qquad R_n=\sum_iB_{t_i}\Delta B_i,\qquad S_n=\sum_iB_{\frac{t_{i-1}+t_i}{2}}\Delta B_i,
\]
分别在左端点、右端点与中点取值.

\begin{proposition}[三种 Riemann 和的极限]\label{prop:three-sums}
当 $n\to\infty$ 时, 在 $L^2$ 中
\[
L_n\to\frac{1}{2}\big(B_t^2-t\big),\qquad R_n\to\frac{1}{2}\big(B_t^2+t\big),\qquad S_n\to\frac{1}{2}B_t^2.
\]
\end{proposition}

\begin{proof}
记 $Q_n=\sum_i(\Delta B_i)^2$. 两个代数恒等式是关键:
\begin{align*}
R_n-L_n&=\sum_i(B_{t_i}-B_{t_{i-1}})\Delta B_i=Q_n,\\
R_n+L_n&=\sum_i(B_{t_i}+B_{t_{i-1}})(B_{t_i}-B_{t_{i-1}})=\sum_i\big(B_{t_i}^2-B_{t_{i-1}}^2\big)=B_t^2.
\end{align*}
于是 $L_n=\frac{1}{2}(B_t^2-Q_n)$, $R_n=\frac{1}{2}(B_t^2+Q_n)$. 由布朗运动的二次变差定理~\ref{thm:bm-qv}, $Q_n\to t$ 于 $L^2$, 前两个结论得证.

对 $S_n$, 记 $m_i=\frac{t_{i-1}+t_i}{2}$, 把 $\Delta B_i$ 拆为 $(B_{m_i}-B_{t_{i-1}})+(B_{t_i}-B_{m_i})$,
\[
S_n=L_n+\sum_i(B_{m_i}-B_{t_{i-1}})\Delta B_i=L_n+\sum_i(B_{m_i}-B_{t_{i-1}})^2+\sum_i(B_{m_i}-B_{t_{i-1}})(B_{t_i}-B_{m_i}).
\]
第二项是一组独立随机变量之和, 均值为 $\sum_i\frac{t_i-t_{i-1}}{2}=\frac{t}{2}$, 方差为 $\sum_i2\big(\frac{t_i-t_{i-1}}{2}\big)^2\le\frac{t}{2}\abs{\Pi_n}\to0$ (与定理~\ref{thm:bm-qv} 的证明相同), 所以它在 $L^2$ 中趋于 $\frac{t}{2}$. 第三项中各被加项是两个独立的均值为零的正态变量之积, 不同的 $i$ 之间也相互独立, 所以它的均值为零, 方差为 $\sum_i\frac{(t_i-t_{i-1})^2}{4}\le\frac{t}{4}\abs{\Pi_n}\to0$. 因此 $S_n\to\frac{1}{2}(B_t^2-t)+\frac{t}{2}=\frac{1}{2}B_t^2$.
\end{proof}

\begin{figure}[htbp]
\centering
\begin{tikzpicture}
\begin{axis}[SAaxis, axis x line=bottom, axis y line=left, width=12cm, height=6cm, xmin=0.5, xmax=14.5, ymin=-0.4, ymax=1.25,
  xtick={1,2,...,14}, ytick={0,0.5,1}, xlabel={$n$ (分割 $[0,1]$ 为 $2^n$ 段)},
  legend style={font=\small,draw=none,fill=none,at={(0.5,-0.22)},anchor=north}, legend columns=3, legend cell align=left]
  \addplot[SARed,thick,mark=*,mark size=1.4pt] coordinates {(1,0.7979) (2,1.2868) (3,1.3292) (4,1.0124) (5,0.9557) (6,0.9597) (7,1.0437) (8,1.0514) (9,1.0031) (10,1.0051) (11,1.0113) (12,1.0205) (13,1.0244) (14,1.0258)};
  \addplot[SAGreenText,thick,mark=square*,mark size=1.4pt] coordinates {(1,0.4537) (2,1.2350) (3,0.5783) (4,0.6212) (5,0.5728) (6,0.4672) (7,0.5082) (8,0.5365) (9,0.5327) (10,0.5214) (11,0.5146) (12,0.5126) (13,0.5181) (14,0.5229)};
  \addplot[SABlue,thick,mark=triangle*,mark size=1.6pt] coordinates {(1,0.2518) (2,-0.2370) (3,-0.2795) (4,0.0373) (5,0.0940) (6,0.0901) (7,0.0060) (8,-0.0017) (9,0.0466) (10,0.0446) (11,0.0384) (12,0.0293) (13,0.0253) (14,0.0239)};
  \addplot[SARed!60,dashed,forget plot] coordinates {(0.5,1.0249) (14.5,1.0249)};
  \addplot[SAGreenText!60,dashed,forget plot] coordinates {(0.5,0.5249) (14.5,0.5249)};
  \addplot[SABlue!60,dashed,forget plot] coordinates {(0.5,0.0249) (14.5,0.0249)};
  \legend{右端点 $R_n$, 中点 $S_n$, 左端点 $L_n$}
\end{axis}
\end{tikzpicture}
\caption{同一条布朗轨道上 $\int_0^1B\dd B$ 的三种 Riemann 和 (二进分割). 它们分别收敛到三个不同的值 (虚线): $\frac12(B_1^2+1)$, $\frac12B_1^2$, $\frac12(B_1^2-1)$. 相邻两条虚线的间距恰为二次变差的一半, 即 $\frac12$}
\label{fig:three-sums}
\end{figure}

与普通微积分对照: 若 $f$ 连续可微且 $f(0)=0$, 则 $\int_0^tf\dd f=\frac{1}{2}f(t)^2$, 与取值点无关. 对布朗运动, 三种取值点给出三个不同的答案, 差别恰好来自二次变差 $Q_n\to t$ (图~\ref{fig:three-sums}). 因此 ``$\int_0^tB\dd B$'' 本身没有确定的含义, 必须事先约定取值点.

% ---------------------------------------------------------------------
\subsection{为什么选左端点}

左端点的选择 (It\^o 积分) 与中点的选择 (\textbf{Stratonovich 积分}, 记作 $\int B\circ\dd B$) 各有优点. 中点给出的 $\frac{1}{2}B_t^2$ 符合普通微积分的链式法则; 左端点给出的 $\frac{1}{2}(B_t^2-t)$ 多了一个修正项, 但它是一个\emph{鞅} (定理~\ref{thm:bm-martingales} (2)), 期望为零. 计算左、右端点和的期望可见,
\[
\E L_n=\sum_i\E\big[B_{t_{i-1}}\big]\E\big[\Delta B_i\big]=0,\qquad\E R_n=\E L_n+\E Q_n=t.
\]
$L_n$ 正是第~3~章的鞅变换: 若把 $\Delta B_i$ 看作第 $i$ 局的输赢, $B_{t_{i-1}}$ 看作第 $i$ 局开始\emph{之前}决定的赌注, 则 $L_n$ 是累计收益, 由鞅变换定理~\ref{thm:martingale-transform}, 它是鞅. 右端点和则相当于 ``看到第 $i$ 局的结果之后再决定这一局的赌注'', 这是作弊, 平均收益为正.

\begin{intuition}
在建模中, 被积函数 $H_t$ 代表时刻 $t$ 做出的决策 (持有多少股票、施加多大的控制) 或系统在时刻 $t$ 对噪声的响应强度. 这些量只能依赖于到时刻 $t$ 为止的信息, 而噪声增量 $\dd B_t=B_{t+\dd t}-B_t$ 则是此后才揭晓的. 左端点取值精确地表达了这种因果关系: 先有 $H_t$, 后有 $\dd B_t$, 二者独立. 正因为如此, It\^o 积分是鞅, 且有 It\^o 等距 $\E\big(\int H\dd B\big)^2=\E\int H^2\dd t$. 代价是链式法则多出二阶项, 即 5.3 节的 It\^o 公式. Stratonovich 积分保持了普通的链式法则, 但失去了鞅性质; 二者可以互相转换 (5.3 节末的注), 本书以 It\^o 积分为主.
\end{intuition}

把上面的讨论推广到一般的被积过程 $H$, 就是下一节的构造: 先对 ``阶梯形'' 的 $H$ (在每个小区间上取左端点的值) 定义积分, 证明 It\^o 等距, 再用 $L^2$ 极限推广到一般的 $H$.

% =====================================================================
\section{It\^o 积分的构造与性质}
\label{sec:5-2}

构造分三步: (1) 确定被积过程的范围, 它们必须 ``不偷看未来''; (2) 对简单过程 (阶梯形的被积过程) 直接定义积分, 并证明 It\^o 等距; (3) 证明简单过程在适当的 $L^2$ 空间中稠密, 由等距把积分延拓到整个空间. 之后再用停时把积分推广到只满足 $\int_0^TH_s^2\dd s<\infty$ $\as$ 的被积过程.

% ---------------------------------------------------------------------
\subsection{被积过程}

被积过程 $H$ 应当是适应的. 但积分 $\int H_s^2\dd s$ 还要求 $H$ 作为 $(s,\omega)$ 的二元函数是可测的, 而且这种可测性要与信息流相容.

\begin{definition}[循序可测]\label{def:progressive}
过程 $H=\{H_t\}_{t\in[0,T]}$ 称为\textbf{循序可测的} (progressively measurable), 如果对每个 $t\in[0,T]$, 映射 $(s,\omega)\mapsto H_s(\omega)$ 限制在 $[0,t]\times\Omega$ 上是 $\mathcal{B}([0,t])\otimes\mathcal{F}_t$-可测的.
\end{definition}

循序可测过程是适应的: 对固定的 $s$, 可测函数的截口是可测的, 所以 $\omega\mapsto H_s(\omega)$ 是 $\mathcal{F}_s$-可测的 (在上述定义中取 $t=s$). 由 Fubini 定理~\ref{thm:fubini}, 若 $H$ 循序可测且 $\int_0^t\abs{H_s}\dd s<\infty$, 则 $\int_0^tH_s\dd s$ 是 $\mathcal{F}_t$-可测的. 循序可测过程的和、积、连续函数仍是循序可测的. 下面的命题说明, 实际中遇到的过程都是循序可测的.

\begin{proposition}\label{prop:progressive}
设 $H$ 是适应过程, 且每条轨道右连续 (或每条轨道左连续). 则 $H$ 循序可测. 特别地, 对停时 $\tau$, $\1_{[0,\tau]}(t)=\1_{\{t\le\tau\}}$ 是循序可测的.
\end{proposition}

\begin{proof}
设轨道右连续. 固定 $t$, 对 $n\ge1$ 与 $s\in[0,t]$, 令 $H^{(n)}_s=H_{(k+1)t2^{-n}}$, 当 $s\in[kt2^{-n},(k+1)t2^{-n})$ ($0\le k<2^n$), 并令 $H^{(n)}_t=H_t$. 则
\[
H^{(n)}_s(\omega)=\sum_{k=0}^{2^n-1}\1_{[kt2^{-n},(k+1)t2^{-n})}(s)\,H_{(k+1)t2^{-n}}(\omega)+\1_{\{t\}}(s)H_t(\omega),
\]
每一项都是 $\mathcal{B}([0,t])$-可测函数与 $\mathcal{F}_t$-可测函数之积, 所以 $H^{(n)}$ 是 $\mathcal{B}([0,t])\otimes\mathcal{F}_t$-可测的. 对 $s<t$, $(k+1)t2^{-n}\downarrow s$, 由右连续性 $H^{(n)}_s\to H_s$; 对 $s=t$ 显然. 可测函数的逐点极限可测. 左连续的情形相同, 只需改用 $H_{kt2^{-n}}$ (左端点) 并让 $s\in(kt2^{-n},(k+1)t2^{-n}]$. 最后, $\1_{\{t\le\tau\}}$ 关于 $t$ 左连续, 且 $\{t\le\tau\}=\{\tau<t\}^c\in\mathcal{F}_t$ (引理~\ref{lem:optional-time} 的证明), 所以它是适应的.
\end{proof}

\begin{definition}[被积过程空间]\label{def:L2-integrands}
记 $\mathscr{L}^2_T$ 为满足
\[
\norm{H}_T^2:=\E\int_0^TH_s^2\dd s<\infty
\]
的循序可测过程 $H$ 全体, 并把 $\norm{H-H'}_T=0$ 的两个过程视为同一元素.
\end{definition}

令 $\mathcal{P}_T$ 为 $[0,T]\times\Omega$ 上使示性函数循序可测的集合全体, 可以验证它是 $\sigma$-代数 (称为\textbf{循序 $\sigma$-代数}), 且 $H$ 循序可测当且仅当 $H$ 是 $\mathcal{P}_T$-可测的. 于是 $\mathscr{L}^2_T$ 就是测度空间 $([0,T]\times\Omega,\mathcal{P}_T,\lambda\otimes\P)$ 上的 $L^2$ 空间 ($\lambda$ 为 Lebesgue 测度), 由 Riesz-Fischer 定理~\ref{thm:riesz-fischer} (其证明对一般的有限测度空间逐字成立) 它是 Hilbert 空间, 内积为 $\ip{H,K}_T=\E\int_0^TH_sK_s\dd s$.

% ---------------------------------------------------------------------
\subsection{简单过程的积分}

\begin{definition}[简单过程]\label{def:simple-process}
设 $0=t_0<t_1<\cdots<t_m=T$, $\xi_{i-1}$ 是有界的 $\mathcal{F}_{t_{i-1}}$-可测随机变量 ($i=1,\dots,m$). 称
\[
H_t=\xi_0\1_{\{0\}}(t)+\sum_{i=1}^m\xi_{i-1}\1_{(t_{i-1},t_i]}(t),\qquad t\in[0,T]
\]
为\textbf{简单过程}, 全体记为 $\mathscr{S}_T$. 定义 $H$ 关于 $B$ 的积分为
\begin{equation}\label{eq:simple-integral}
(H\cdot B)_t=\int_0^tH_s\dd B_s=\sum_{i=1}^m\xi_{i-1}\big(B_{t_i\wedge t}-B_{t_{i-1}\wedge t}\big),\qquad t\in[0,T].
\end{equation}
\end{definition}

简单过程在每个区间 $(t_{i-1},t_i]$ 上取常值 $\xi_{i-1}$, 这个值在区间的\emph{左端点}就已确定. 它左连续且适应, 由命题~\ref{prop:progressive} 循序可测; 又有界, 所以 $\mathscr{S}_T\subseteq\mathscr{L}^2_T$. 定义 \eqref{eq:simple-integral} 与表示方式无关: 在分割中插入一个新的分点 $u\in(t_{i-1},t_i)$, 相当于把 $\xi_{i-1}\1_{(t_{i-1},t_i]}$ 写成 $\xi_{i-1}\1_{(t_{i-1},u]}+\xi_{i-1}\1_{(u,t_i]}$ (由于 $\xi_{i-1}$ 也是 $\mathcal{F}_u$-可测的, 这仍是简单过程的表示), 而 \eqref{eq:simple-integral} 中相应的项 $\xi_{i-1}(B_{t_i\wedge t}-B_{t_{i-1}\wedge t})$ 恰好拆成两项之和; 两种表示总有公共的加细. 同理, 两个简单过程可以写在同一个分割上, 所以 $\mathscr{S}_T$ 是线性空间, 积分是线性的. \eqref{eq:simple-integral} 正是第~3~章的鞅变换 (定义~\ref{def:martingale-transform}) 在时刻 $t_0<t_1<\cdots$ 上的形式.

\begin{proposition}[简单过程的 It\^o 积分]\label{prop:simple-integral}
设 $H\in\mathscr{S}_T$, $0\le s\le t\le T$. 则
\begin{enumerate}
\item $H\cdot B$ 是连续鞅, $(H\cdot B)_0=0$;
\item (\textbf{条件 It\^o 等距})
\[
\E\Big[\big((H\cdot B)_t-(H\cdot B)_s\big)^2\,\Big|\,\mathcal{F}_s\Big]=\E\Big[\int_s^tH_u^2\dd u\,\Big|\,\mathcal{F}_s\Big];
\]
\item $(H\cdot B)_t^2-\int_0^tH_u^2\dd u$ 是鞅. 特别地 (\textbf{It\^o 等距}), $\E\big[(H\cdot B)_t^2\big]=\E\int_0^tH_u^2\dd u$.
\end{enumerate}
\end{proposition}

\begin{proof}
必要时加细分割, 不妨设 $s=t_j$, $t=t_k$ 都是分点 ($j\le k$). 记 $\Delta B_i=B_{t_i}-B_{t_{i-1}}$. 则 $(H\cdot B)_t-(H\cdot B)_s=\sum_{i=j+1}^k\xi_{i-1}\Delta B_i$. 由于 $\xi_{i-1}$ 有界, 各项平方可积.

(1) 轨道连续性与 $(H\cdot B)_0=0$ 是显然的, 适应性也是. 对 $j<i\le k$, $\Delta B_i$ 与 $\mathcal{F}_{t_{i-1}}$ 独立, 由塔性质与提出已知因子,
\[
\E\big[\xi_{i-1}\Delta B_i\,\big|\,\mathcal{F}_s\big]=\E\big[\xi_{i-1}\E[\Delta B_i\,|\,\mathcal{F}_{t_{i-1}}]\,\big|\,\mathcal{F}_s\big]=0.
\]
对 $i$ 求和即得鞅性质.

(2) 展开平方,
\[
\Big(\sum_{i=j+1}^k\xi_{i-1}\Delta B_i\Big)^2=\sum_{i=j+1}^k\xi_{i-1}^2(\Delta B_i)^2+2\sum_{j<i<l\le k}\xi_{i-1}\xi_{l-1}\Delta B_i\Delta B_l.
\]
对平方项, $\Delta B_i$ 与 $\mathcal{F}_{t_{i-1}}$ 独立且 $\E(\Delta B_i)^2=t_i-t_{i-1}$, 所以 $\E[\xi_{i-1}^2(\Delta B_i)^2\,|\,\mathcal{F}_{t_{i-1}}]=\xi_{i-1}^2(t_i-t_{i-1})$. 对交叉项 ($i<l$), $\xi_{i-1}\xi_{l-1}\Delta B_i$ 是 $\mathcal{F}_{t_{l-1}}$-可测的, 而 $\E[\Delta B_l\,|\,\mathcal{F}_{t_{l-1}}]=0$, 所以 $\E[\xi_{i-1}\xi_{l-1}\Delta B_i\Delta B_l\,|\,\mathcal{F}_{t_{l-1}}]=0$. 再对 $\mathcal{F}_s$ 取条件期望 (塔性质),
\[
\E\Big[\Big(\sum_{i=j+1}^k\xi_{i-1}\Delta B_i\Big)^2\,\Big|\,\mathcal{F}_s\Big]=\E\Big[\sum_{i=j+1}^k\xi_{i-1}^2(t_i-t_{i-1})\,\Big|\,\mathcal{F}_s\Big]=\E\Big[\int_s^tH_u^2\dd u\,\Big|\,\mathcal{F}_s\Big].
\]

(3) 对任意平方可积鞅 $X$, $\E[X_t^2-X_s^2\,|\,\mathcal{F}_s]=\E[(X_t-X_s)^2\,|\,\mathcal{F}_s]$, 因为 $\E[X_tX_s\,|\,\mathcal{F}_s]=X_s^2$. 由 (2),
\[
\E\Big[(H\cdot B)_t^2-\int_0^tH_u^2\dd u\,\Big|\,\mathcal{F}_s\Big]=(H\cdot B)_s^2-\int_0^sH_u^2\dd u.
\]
取 $s=0$ 并取期望即得 It\^o 等距.
\end{proof}

It\^o 等距说: 映射 $H\mapsto(H\cdot B)_T$ 是从 $(\mathscr{S}_T,\norm{\cdot}_T)$ 到 $L^2(\Omega,\mathcal{F}_T,\P)$ 的\emph{等距}线性映射. 它是第~3~章离散 It\^o 等距 (推论~\ref{cor:discrete-ito-isometry}) 的连续版本: 布朗运动的 ``可料二次变差'' 是 $\dd t$. 等距映射把 Cauchy 列映为 Cauchy 列, 所以只要 $\mathscr{S}_T$ 在 $\mathscr{L}^2_T$ 中稠密, 积分就可以唯一地延拓到 $\mathscr{L}^2_T$ 上.

% ---------------------------------------------------------------------
\subsection{简单过程的稠密性}

\begin{lemma}[简单过程的稠密性]\label{lem:simple-dense}
$\mathscr{S}_T$ 在 $\mathscr{L}^2_T$ 中稠密: 对每个 $H\in\mathscr{L}^2_T$, 存在 $H^n\in\mathscr{S}_T$ 使 $\norm{H^n-H}_T\to0$.
\end{lemma}

\begin{proof}
记 $\overline{\mathscr{S}}$ 为 $\mathscr{S}_T$ 在 $\mathscr{L}^2_T$ 中的闭包. 它是闭的线性子空间. 分三步证明 $\overline{\mathscr{S}}=\mathscr{L}^2_T$.

\emph{第一步: 有界、轨道连续的适应过程属于 $\overline{\mathscr{S}}$.} 设 $\abs{H}\le C$. 令 $s_k=kT2^{-n}$,
\[
H^n_t=H_0\1_{\{0\}}(t)+\sum_{k=1}^{2^n}H_{s_{k-1}}\1_{(s_{k-1},s_k]}(t).
\]
$H_{s_{k-1}}$ 有界且 $\mathcal{F}_{s_{k-1}}$-可测, 所以 $H^n\in\mathscr{S}_T$. 对每个 $(t,\omega)$, 由轨道连续性 $H^n_t(\omega)\to H_t(\omega)$; 又 $\abs{H^n_t-H_t}^2\le4C^2$. 由有界收敛定理 (在有限测度空间 $([0,T]\times\Omega,\lambda\otimes\P)$ 上), $\norm{H^n-H}_T\to0$ (图~\ref{fig:simple-approx}).

\emph{第二步: 有界循序可测过程属于 $\overline{\mathscr{S}}$.} 设 $\abs{H}\le C$. 对 $m\ge1$ 令
\[
G^m_t=m\int_{(t-1/m)\vee0}^tH_s\dd s,\qquad t\in[0,T].
\]
(a) $G^m$ 适应: 由循序可测性, $(s,\omega)\mapsto H_s(\omega)\1_{[(t-1/m)\vee0,t]}(s)$ 在 $[0,t]\times\Omega$ 上是 $\mathcal{B}([0,t])\otimes\mathcal{F}_t$-可测的, 由 Fubini 定理, 它对 $s$ 的积分是 $\mathcal{F}_t$-可测的. (b) $G^m$ 有界且轨道连续: $\abs{G^m_t}\le C$, 且 $\abs{G^m_t-G^m_u}\le2mC\abs{t-u}$. 由第一步, $G^m\in\overline{\mathscr{S}}$. (c) $\norm{G^m-H}_T\to0$: 固定 $\omega$, $s\mapsto H_s(\omega)$ 是 $[0,T]$ 上的有界 Borel 函数. 由 Lebesgue 微分定理 (见~\cite{folland} 定理 3.20), 对 Lebesgue 几乎每个 $t\in(0,T]$,
\[
\abs{G^m_t(\omega)-H_t(\omega)}\le m\int_{t-1/m}^t\abs{H_s(\omega)-H_t(\omega)}\dd s\to0\qquad(m\to\infty,\ m>1/t).
\]
由于 $\abs{G^m-H}^2\le4C^2$, 由有界收敛定理先对 $t$ 积分, 得 $\int_0^T\abs{G^m_t(\omega)-H_t(\omega)}^2\dd t\to0$; 再对 $\omega$ 用有界收敛定理, 得 $\norm{G^m-H}_T\to0$. 由于 $\overline{\mathscr{S}}$ 是闭的, $H\in\overline{\mathscr{S}}$.

\emph{第三步: 一般情形.} 设 $H\in\mathscr{L}^2_T$, 令 $H^{(n)}=(H\wedge n)\vee(-n)$. 它有界、循序可测, 由第二步属于 $\overline{\mathscr{S}}$. 又 $\abs{H^{(n)}-H}^2\le H^2$, 且 $H^{(n)}\to H$ 逐点成立, 由控制收敛定理 $\norm{H^{(n)}-H}_T\to0$. 所以 $H\in\overline{\mathscr{S}}$.
\end{proof}

\begin{figure}[htbp]
\centering
\begin{tikzpicture}
\begin{axis}[SAaxis, axis x line=bottom, axis y line=left, width=12cm, height=5.4cm, xmin=0, xmax=1.04, ymin=-0.15, ymax=1.35,
  xtick={0,0.125,0.25,0.375,0.5,0.625,0.75,0.875,1}, xticklabels={$0$,,$\frac14$,,$\frac12$,,$\frac34$,,$T$}, ytick={0,0.5,1}, xlabel={$t$}]
  \addplot[SADark!45,thin] coordinates {(0.000,0.300) (0.003,0.302) (0.005,0.370) (0.007,0.431) (0.010,0.405) (0.013,0.391) (0.015,0.364) (0.018,0.393) (0.020,0.390) (0.022,0.427) (0.025,0.335) (0.028,0.413) (0.030,0.408) (0.033,0.442) (0.035,0.436) (0.037,0.417) (0.040,0.440) (0.043,0.481) (0.045,0.471) (0.048,0.463) (0.050,0.497) (0.052,0.454) (0.055,0.378) (0.058,0.398) (0.060,0.364) (0.062,0.268) (0.065,0.228) (0.068,0.204) (0.070,0.145) (0.072,0.070) (0.075,0.072) (0.077,0.117) (0.080,0.105) (0.083,0.068) (0.085,0.087) (0.088,0.123) (0.090,0.108) (0.092,0.135) (0.095,0.187) (0.098,0.177) (0.100,0.136) (0.103,0.154) (0.105,0.166) (0.107,0.221) (0.110,0.157) (0.113,0.124) (0.115,0.082) (0.118,-0.005) (0.120,0.001) (0.122,0.028) (0.125,-0.009) (0.128,0.060) (0.130,0.101) (0.133,0.133) (0.135,0.153) (0.138,0.201) (0.140,0.134) (0.143,0.165) (0.145,0.195) (0.147,0.106) (0.150,0.124) (0.152,0.111) (0.155,0.150) (0.158,0.128) (0.160,0.127) (0.163,0.145) (0.165,0.101) (0.168,0.131) (0.170,0.125) (0.173,0.150) (0.175,0.124) (0.177,0.178) (0.180,0.209) (0.182,0.200) (0.185,0.231) (0.188,0.294) (0.190,0.384) (0.193,0.305) (0.195,0.349) (0.198,0.373) (0.200,0.368) (0.203,0.318) (0.205,0.380) (0.208,0.317) (0.210,0.346) (0.212,0.411) (0.215,0.331) (0.217,0.316) (0.220,0.250) (0.223,0.262) (0.225,0.338) (0.228,0.439) (0.230,0.350) (0.233,0.322) (0.235,0.357) (0.238,0.436) (0.240,0.457) (0.242,0.419) (0.245,0.434) (0.247,0.434) (0.250,0.423) (0.253,0.387) (0.255,0.406) (0.258,0.421) (0.260,0.417) (0.263,0.406) (0.265,0.341) (0.268,0.317) (0.270,0.377) (0.273,0.368) (0.275,0.296) (0.278,0.363) (0.280,0.389) (0.283,0.495) (0.285,0.498) (0.288,0.475) (0.290,0.402) (0.292,0.468) (0.295,0.597) (0.297,0.556) (0.300,0.523) (0.302,0.553) (0.305,0.512) (0.307,0.498) (0.310,0.481) (0.312,0.490) (0.315,0.545) (0.318,0.546) (0.320,0.592) (0.323,0.571) (0.325,0.588) (0.328,0.481) (0.330,0.408) (0.333,0.448) (0.335,0.418) (0.338,0.447) (0.340,0.475) (0.343,0.541) (0.345,0.581) (0.348,0.632) (0.350,0.626) (0.352,0.592) (0.355,0.555) (0.357,0.531) (0.360,0.474) (0.362,0.447) (0.365,0.442) (0.367,0.455) (0.370,0.438) (0.372,0.342) (0.375,0.338) (0.378,0.349) (0.380,0.403) (0.383,0.432) (0.385,0.400) (0.388,0.364) (0.390,0.464) (0.393,0.502) (0.395,0.594) (0.398,0.700) (0.400,0.659) (0.403,0.679) (0.405,0.702) (0.408,0.730) (0.410,0.757) (0.413,0.767) (0.415,0.776) (0.417,0.700) (0.420,0.692) (0.422,0.655) (0.425,0.661) (0.427,0.638) (0.430,0.669) (0.432,0.710) (0.435,0.725) (0.438,0.741) (0.440,0.745) (0.443,0.723) (0.445,0.715) (0.448,0.690) (0.450,0.709) (0.453,0.710) (0.455,0.739) (0.458,0.673) (0.460,0.717) (0.463,0.679) (0.465,0.642) (0.468,0.632) (0.470,0.604) (0.473,0.619) (0.475,0.590) (0.478,0.537) (0.480,0.494) (0.482,0.560) (0.485,0.562) (0.487,0.504) (0.490,0.504) (0.492,0.426) (0.495,0.516) (0.497,0.440) (0.500,0.464) (0.503,0.491) (0.505,0.418) (0.508,0.433) (0.510,0.483) (0.512,0.507) (0.515,0.520) (0.517,0.567) (0.520,0.575) (0.522,0.592) (0.525,0.586) (0.527,0.617) (0.530,0.570) (0.532,0.610) (0.535,0.584) (0.537,0.530) (0.540,0.548) (0.542,0.633) (0.545,0.681) (0.547,0.655) (0.550,0.690) (0.552,0.667) (0.555,0.661) (0.557,0.666) (0.560,0.550) (0.562,0.559) (0.565,0.508) (0.568,0.473) (0.570,0.399) (0.573,0.323) (0.575,0.276) (0.578,0.318) (0.580,0.401) (0.583,0.400) (0.585,0.454) (0.588,0.441) (0.590,0.345) (0.593,0.353) (0.595,0.375) (0.598,0.354) (0.600,0.369) (0.603,0.397) (0.605,0.354) (0.608,0.280) (0.610,0.269) (0.613,0.259) (0.615,0.241) (0.618,0.291) (0.620,0.377) (0.623,0.356) (0.625,0.391) (0.628,0.438) (0.630,0.474) (0.633,0.526) (0.635,0.507) (0.638,0.544) (0.640,0.645) (0.642,0.685) (0.645,0.654) (0.647,0.684) (0.650,0.747) (0.652,0.766) (0.655,0.738) (0.657,0.815) (0.660,0.752) (0.662,0.778) (0.665,0.777) (0.667,0.717) (0.670,0.778) (0.672,0.795) (0.675,0.737) (0.677,0.767) (0.680,0.746) (0.682,0.650) (0.685,0.616) (0.688,0.629) (0.690,0.661) (0.693,0.669) (0.695,0.672) (0.698,0.695) (0.700,0.684) (0.703,0.743) (0.705,0.751) (0.708,0.765) (0.710,0.792) (0.713,0.740) (0.715,0.704) (0.718,0.781) (0.720,0.797) (0.723,0.783) (0.725,0.800) (0.728,0.776) (0.730,0.790) (0.733,0.757) (0.735,0.771) (0.738,0.692) (0.740,0.758) (0.743,0.733) (0.745,0.654) (0.748,0.642) (0.750,0.624) (0.753,0.631) (0.755,0.574) (0.758,0.594) (0.760,0.775) (0.763,0.711) (0.765,0.728) (0.768,0.713) (0.770,0.722) (0.772,0.632) (0.775,0.574) (0.777,0.597) (0.780,0.595) (0.782,0.677) (0.785,0.642) (0.787,0.650) (0.790,0.796) (0.792,0.758) (0.795,0.713) (0.797,0.711) (0.800,0.709) (0.802,0.751) (0.805,0.757) (0.807,0.720) (0.810,0.730) (0.812,0.862) (0.815,0.926) (0.818,0.787) (0.820,0.779) (0.823,0.736) (0.825,0.767) (0.828,0.758) (0.830,0.855) (0.833,0.901) (0.835,0.946) (0.838,0.955) (0.840,0.953) (0.843,0.970) (0.845,1.035) (0.848,1.062) (0.850,1.044) (0.853,1.108) (0.855,1.117) (0.858,1.114) (0.860,1.074) (0.863,1.046) (0.865,1.014) (0.868,0.874) (0.870,0.925) (0.873,0.953) (0.875,0.943) (0.878,0.995) (0.880,1.017) (0.883,1.035) (0.885,0.912) (0.888,0.900) (0.890,0.927) (0.893,1.007) (0.895,1.099) (0.897,1.171) (0.900,1.112) (0.902,1.003) (0.905,1.031) (0.907,1.032) (0.910,1.081) (0.912,1.074) (0.915,1.085) (0.917,1.010) (0.920,1.033) (0.922,1.010) (0.925,1.028) (0.927,1.051) (0.930,1.045) (0.932,1.061) (0.935,1.060) (0.938,1.098) (0.940,1.123) (0.943,1.096) (0.945,1.066) (0.948,1.154) (0.950,1.149) (0.953,1.198) (0.955,1.218) (0.958,1.190) (0.960,1.152) (0.963,1.158) (0.965,1.171) (0.968,1.202) (0.970,1.138) (0.973,1.109) (0.975,1.104) (0.978,1.184) (0.980,1.066) (0.983,1.088) (0.985,1.061) (0.988,1.082) (0.990,1.058) (0.993,1.076) (0.995,1.154) (0.998,1.172) (1.000,1.194)};
  \addplot[SABlue,thick] coordinates {(0.000,0.300) (0.125,0.300)};
  \addplot[SABlue,thick] coordinates {(0.125,-0.009) (0.250,-0.009)};
  \addplot[SABlue,thick] coordinates {(0.250,0.423) (0.375,0.423)};
  \addplot[SABlue,thick] coordinates {(0.375,0.338) (0.500,0.338)};
  \addplot[SABlue,thick] coordinates {(0.500,0.464) (0.625,0.464)};
  \addplot[SABlue,thick] coordinates {(0.625,0.391) (0.750,0.391)};
  \addplot[SABlue,thick] coordinates {(0.750,0.624) (0.875,0.624)};
  \addplot[SABlue,thick] coordinates {(0.875,0.943) (1.000,0.943)};
  \fill[SABlue] (axis cs:0.000,0.300) circle (1.5pt);
  \fill[SABlue] (axis cs:0.125,-0.009) circle (1.5pt);
  \fill[SABlue] (axis cs:0.250,0.423) circle (1.5pt);
  \fill[SABlue] (axis cs:0.375,0.338) circle (1.5pt);
  \fill[SABlue] (axis cs:0.500,0.464) circle (1.5pt);
  \fill[SABlue] (axis cs:0.625,0.391) circle (1.5pt);
  \fill[SABlue] (axis cs:0.750,0.624) circle (1.5pt);
  \fill[SABlue] (axis cs:0.875,0.943) circle (1.5pt);
\end{axis}
\end{tikzpicture}
\caption{用简单过程逼近连续的适应过程 (灰色). 在每个区间 $(s_{k-1},s_k]$ 上, 简单过程 (蓝色) 取\emph{左端点}的值 $H_{s_{k-1}}$ (实心点), 它在区间开始时就已知. 分割加细时, 由轨道连续性, 阶梯函数逐点收敛到原过程}
\label{fig:simple-approx}
\end{figure}

\begin{remark}
文献中常把被积过程的条件写成 ``适应且 $(t,\omega)$ 联合可测''. 这样的过程总有一个循序可测的修正 (见~\cite{karatzas} 命题 1.1.12), 所以两种写法定义出的积分相同. 第二步中循序可测性的作用是保证 $G^m$ 适应.
\end{remark}

% ---------------------------------------------------------------------
\subsection{It\^o 积分及其性质}

\begin{theorem}[It\^o 积分]\label{thm:ito-integral}
设 $H\in\mathscr{L}^2_T$. 则存在 (在不可区分的意义下) 唯一的连续过程 $\{(H\cdot B)_t\}_{t\in[0,T]}$, 使得对任意满足 $\norm{H^n-H}_T\to0$ 的简单过程序列 $H^n\in\mathscr{S}_T$,
\begin{equation}\label{eq:ito-integral-limit}
\E\Big[\sup_{t\le T}\big\lvert(H^n\cdot B)_t-(H\cdot B)_t\big\rvert^2\Big]\to0.
\end{equation}
$H\cdot B$ 是连续鞅, 称为 $H$ 关于 $B$ 的 \textbf{It\^o 积分}, 记作 $(H\cdot B)_t=\int_0^tH_s\dd B_s$. 对 $s\le t$, 记 $\int_s^tH_u\dd B_u=(H\cdot B)_t-(H\cdot B)_s$.
\end{theorem}

\begin{proof}
由引理~\ref{lem:simple-dense}, 存在简单过程 $H^n\to H$. 记 $X^n=H^n\cdot B$. 由命题~\ref{prop:simple-integral}, $X^n-X^m=(H^n-H^m)\cdot B$ 是连续鞅, 由连续时间 Doob 不等式~\ref{thm:ct-doob} 与 It\^o 等距,
\begin{equation}\label{eq:ito-doob}
\E\Big[\sup_{t\le T}\abs{X^n_t-X^m_t}^2\Big]\le4\E\big[(X^n_T-X^m_T)^2\big]=4\norm{H^n-H^m}_T^2.
\end{equation}
由于 $\{H^n\}$ 是 Cauchy 列, 可取子列 $n_k$ 使 $\norm{H^{n_{k+1}}-H^{n_k}}_T\le2^{-k}$. 于是
\[
\E\Big[\sum_k\sup_{t\le T}\big\lvert X^{n_{k+1}}_t-X^{n_k}_t\big\rvert\Big]\le\sum_k\Big(\E\sup_{t\le T}\big\lvert X^{n_{k+1}}_t-X^{n_k}_t\big\rvert^2\Big)^{1/2}\le\sum_k2\cdot2^{-k}<\infty.
\]
所以几乎必然地 $\sum_k\sup_t\abs{X^{n_{k+1}}_t-X^{n_k}_t}<\infty$, 在这个事件上 $X^{n_k}$ 在 $[0,T]$ 上一致收敛, 其极限 $X$ 是连续的; 在零概率的补集上令 $X\equiv0$. 由 Fatou 引理与 \eqref{eq:ito-doob},
\[
\E\Big[\sup_t\abs{X^n_t-X_t}^2\Big]\le\liminf_{k\to\infty}\E\Big[\sup_t\abs{X^n_t-X^{n_k}_t}^2\Big]\le4\limsup_{k\to\infty}\norm{H^n-H^{n_k}}_T^2,
\]
当 $n\to\infty$ 时右端趋于零. 这就证明了 \eqref{eq:ito-integral-limit} 对这个序列成立.

与序列无关: 若 $K^n$ 是另一个逼近 $H$ 的简单过程序列, 则交错序列 $H^1$, $K^1$, $H^2$, $K^2$, $\dots$ 也逼近 $H$, 由上面的论证它对应的积分一致收敛到某个 $X'$, 而其子列分别收敛到 $X$ 与 (由 $K^n$ 构造的) 极限, 所以二者不可区分. 唯一性: 若两个连续过程都满足 \eqref{eq:ito-integral-limit}, 则它们之差的上确界在均方意义下为零.

鞅性质: $X^n_t$ 是 $\mathcal{F}_t$-可测的, $X_t$ 是它的几乎必然极限, 由于 $\mathcal{F}_t$ 完备, $X_t$ 是 $\mathcal{F}_t$-可测的. 由 \eqref{eq:ito-integral-limit}, $X^n_t\to X_t$ 于 $L^2$. 条件期望是 $L^2$ 上的压缩 (推论~\ref{cor:ce-Lp}), 所以对 $s\le t$,
\[
\E[X_t\,|\,\mathcal{F}_s]=\lim_{n}\E[X^n_t\,|\,\mathcal{F}_s]=\lim_nX^n_s=X_s\qquad(L^2\text{ 极限}).\qedhere
\]
\end{proof}

\begin{theorem}[It\^o 积分的性质]\label{thm:ito-properties}
设 $H,K\in\mathscr{L}^2_T$, $a,b\in\R$, $0\le s\le t\le T$. 则
\begin{enumerate}
\item (线性) $(aH+bK)\cdot B=a(H\cdot B)+b(K\cdot B)$;
\item (鞅性) $H\cdot B$ 是连续鞅, $\E\int_0^tH\dd B=0$;
\item (It\^o 等距) $\E\big[(\int_s^tH\dd B)^2\,|\,\mathcal{F}_s\big]=\E\big[\int_s^tH_u^2\dd u\,|\,\mathcal{F}_s\big]$, 特别地 $\E(\int_0^tH\dd B)^2=\E\int_0^tH_u^2\dd u$;
\item (二次变差) $(H\cdot B)^2-\int_0^\cdot H_u^2\dd u$ 与 $(H\cdot B)(K\cdot B)-\int_0^\cdot H_uK_u\dd u$ 是鞅. 因此 $\ip{H\cdot B}_t=\int_0^tH_u^2\dd u$, 且 $\E\big[(H\cdot B)_t(K\cdot B)_t\big]=\E\int_0^tH_uK_u\dd u$;
\item (提出已知因子) 若 $\xi$ 是有界的 $\mathcal{F}_s$-可测随机变量, 则 $\int_0^T\xi\1_{(s,t]}(u)H_u\dd B_u=\xi\int_s^tH_u\dd B_u$;
\item (停止) 若 $\tau$ 是停时, 则几乎必然地, 对一切 $t\le T$, $(H\cdot B)_{t\wedge\tau}=\big((H\1_{[0,\tau]})\cdot B\big)_t$.
\end{enumerate}
\end{theorem}

\begin{proof}
设 $H^n,K^n\in\mathscr{S}_T$ 分别逼近 $H,K$.

(1) 对简单过程成立, $aH^n+bK^n$ 逼近 $aH+bK$, 取极限即可. (2) 已在定理~\ref{thm:ito-integral} 中证明.

(3) 记 $Y_n=\int_s^tH^n\dd B$, $Y=\int_s^tH\dd B$, $A_n=\int_s^t(H^n_u)^2\dd u$, $A=\int_s^tH_u^2\dd u$. 由 \eqref{eq:ito-integral-limit}, $Y_n\to Y$ 于 $L^2$, 所以由 $\abs{Y_n^2-Y^2}\le\abs{Y_n-Y}\abs{Y_n+Y}$ 与 Cauchy-Schwarz 不等式, $Y_n^2\to Y^2$ 于 $L^1$. 同理 $\E\abs{A_n-A}\le\norm{H^n-H}_T\norm{H^n+H}_T\to0$. 条件期望是 $L^1$ 上的压缩, 在命题~\ref{prop:simple-integral} (2) 中取极限即得.

(4) 与命题~\ref{prop:simple-integral} (3) 的证明相同, 由 (2)(3) 得到第一个鞅. 对第二个, 用极化恒等式
\[
(H\cdot B)(K\cdot B)=\tfrac{1}{4}\big[((H+K)\cdot B)^2-((H-K)\cdot B)^2\big],\qquad HK=\tfrac{1}{4}\big[(H+K)^2-(H-K)^2\big],
\]
它是两个鞅之差. 由二次变差的唯一性 (定理~\ref{thm:quadratic-variation}), $\ip{H\cdot B}_t=\int_0^tH_u^2\dd u$.

(5) 对简单过程 $H^n$, 把 $s,t$ 加入分割后, $\xi\1_{(s,t]}H^n$ 仍是简单过程 (在 $(t_{i-1},t_i]\subseteq(s,t]$ 上取值 $\xi\xi_{i-1}$, 它是 $\mathcal{F}_{t_{i-1}}$-可测的), 由 \eqref{eq:simple-integral} 直接验证等式. 由于 $\xi$ 有界, $\norm{\xi\1_{(s,t]}H^n-\xi\1_{(s,t]}H}_T\le\sup\abs{\xi}\cdot\norm{H^n-H}_T\to0$, 两端在 $L^2$ 中取极限即可.

(6) 由命题~\ref{prop:progressive}, $H\1_{[0,\tau]}\in\mathscr{L}^2_T$. 记 $M=H\cdot B$, $N=(H\1_{[0,\tau]})\cdot B$. 固定 $t$. 由 (4), $M^2-\int_0^\cdot H^2\dd u$, $N^2-\int_0^\cdot H^2\1_{[0,\tau]}\dd u$ 与 $MN-\int_0^\cdot H^2\1_{[0,\tau]}\dd u$ 都是连续鞅. 对有界停时 $t\wedge\tau$ 应用可选停时定理 (推论~\ref{cor:ct-ost-bounded}),
\[
\E M_{t\wedge\tau}^2=\E\int_0^{t\wedge\tau}H_u^2\dd u,\qquad\E\big[M_{t\wedge\tau}N_{t\wedge\tau}\big]=\E\int_0^{t\wedge\tau}H_u^2\dd u,
\]
而 $\E N_t^2=\E\int_0^tH_u^2\1_{\{u\le\tau\}}\dd u=\E\int_0^{t\wedge\tau}H_u^2\dd u$. 又由可选停时定理 $\E[N_t\,|\,\mathcal{F}_{t\wedge\tau}]=N_{t\wedge\tau}$, 而 $M_{t\wedge\tau}$ 是 $\mathcal{F}_{t\wedge\tau}$-可测的 (第~3~章定义~\ref{def:ct-stopping-time} 之后的说明), 所以 $\E[M_{t\wedge\tau}N_t]=\E[M_{t\wedge\tau}N_{t\wedge\tau}]$. 于是
\[
\E\big[(M_{t\wedge\tau}-N_t)^2\big]=\E M_{t\wedge\tau}^2-2\E\big[M_{t\wedge\tau}N_t\big]+\E N_t^2=0.
\]
所以对每个 $t$, $M_{t\wedge\tau}=N_t$ $\as$; 两端关于 $t$ 连续, 在有理数 $t$ 上同时成立即对一切 $t$ 成立.
\end{proof}

\begin{intuition}
It\^o 积分的全部性质都可以追溯到一个简单的事实: 在 $(t,t+\dd t]$ 上, 赌注 $H_t$ 已知, 而噪声 $\dd B_t$ 与之独立, 均值为零、方差为 $\dd t$. 于是每一小段的 ``收益'' $H_t\dd B_t$ 的条件均值为零 (鞅性), 条件方差为 $H_t^2\dd t$ (等距与二次变差). 不同小段的收益互不相关, 所以总方差就是各段方差之和 $\int H_t^2\dd t$.
\end{intuition}

% ---------------------------------------------------------------------
\subsection{局部化}

在应用中, 被积过程往往只满足 $\int_0^TH_s^2\dd s<\infty$ $\as$, 而不满足 $\E\int_0^TH_s^2\dd s<\infty$. 例如 $H_s=e^{B_s^2}$: 由于 $\E e^{2B_s^2}=\E e^{2sZ^2}=\infty$ 当 $s\ge\frac14$ ($Z\sim N(0,1)$), 若 $T\ge\frac14$ 则 $H\notin\mathscr{L}^2_T$; 但 $H$ 的轨道连续, 所以 $\int_0^TH_s^2\dd s<\infty$ 对每条轨道成立. 用停时截断, 可以把 It\^o 积分推广到这类过程.

\begin{definition}\label{def:L2-loc}
记 $\mathscr{L}^2_{\mathrm{loc}}$ 为满足 $\int_0^TH_s^2\dd s<\infty$ $\as$ 的循序可测过程 $H$ 全体.
\end{definition}

设 $H\in\mathscr{L}^2_{\mathrm{loc}}$. 令 $A_t=\int_0^tH_s^2\dd s$, 它是连续适应过程 (在零概率集上令 $A\equiv0$). 对 $n\ge1$ 令
\[
\tau_n=\inf\{t\in[0,T]: A_t\ge n\}\qquad(\inf\varnothing=T).
\]
由命题~\ref{prop:hitting-continuous}, $\tau_n$ 是停时; $\tau_n$ 关于 $n$ 单调不减, 且当 $n>A_T$ 时 $\tau_n=T$. 令 $H^{(n)}=H\1_{[0,\tau_n]}$, 则 $\int_0^T(H^{(n)}_s)^2\dd s=A_{\tau_n}\le n$, 所以 $H^{(n)}\in\mathscr{L}^2_T$. 对 $m\ge n$, 由于 $H^{(m)}\1_{[0,\tau_n]}=H^{(n)}$, 由定理~\ref{thm:ito-properties} (6),
\begin{equation}\label{eq:localization-consistent}
\big(H^{(m)}\cdot B\big)_{t\wedge\tau_n}=\big(H^{(n)}\cdot B\big)_t,\qquad t\le T,\quad\as
\end{equation}
即在 $\{t\le\tau_n\}$ 上, $(H^{(m)}\cdot B)_t$ 与 $(H^{(n)}\cdot B)_t$ 相等. 这就使下面的定义合理.

\begin{definition}[局部化的 It\^o 积分]\label{def:ito-local}
设 $H\in\mathscr{L}^2_{\mathrm{loc}}$, $\tau_n$, $H^{(n)}$ 如上. 定义
\[
(H\cdot B)_t=\int_0^tH_s\dd B_s=\big(H^{(n)}\cdot B\big)_t,\qquad\text{当 }t\le\tau_n.
\]
\end{definition}

由于 $\bigcup_n\{\tau_n=T\}$ 的概率为 $1$, 这在一个概率为 $1$ 的事件上定义了 $(H\cdot B)_t$ ($t\le T$), 且与 $n$ 的选取无关. 若 $H\in\mathscr{L}^2_T$, 则由定理~\ref{thm:ito-properties} (6), 新定义与旧定义一致.

\begin{theorem}[局部化 It\^o 积分的性质]\label{thm:ito-local}
设 $H,K\in\mathscr{L}^2_{\mathrm{loc}}$. 则
\begin{enumerate}
\item $H\cdot B$ 是连续局部鞅, 局部化序列为 $\{\tau_n\}$: $(H\cdot B)_{\cdot\wedge\tau_n}=H^{(n)}\cdot B$ 是鞅;
\item $(H\cdot B)^2-\int_0^\cdot H_s^2\dd s$ 是连续局部鞅, 所以 $\ip{H\cdot B}_t=\int_0^tH_s^2\dd s$;
\item 线性与停止性质 (定理~\ref{thm:ito-properties} (6)) 仍成立.
\end{enumerate}
\end{theorem}

\begin{proof}
(1) 由 \eqref{eq:localization-consistent} (令 $m\to\infty$, 即用 $H\cdot B$ 代替 $H^{(m)}\cdot B$), $(H\cdot B)_{t\wedge\tau_n}=(H^{(n)}\cdot B)_t$, 右端是连续鞅. (这里局部化序列满足 $\P(\tau_n=T)\to1$, 这是定义~\ref{def:local-martingale} 在有限时间区间 $[0,T]$ 上的对应形式.) (2) 停止在 $\tau_n$ 后, $(H\cdot B)^2_{\cdot\wedge\tau_n}-\int_0^{\cdot\wedge\tau_n}H^2\dd s=(H^{(n)}\cdot B)^2-\int_0^\cdot(H^{(n)})^2\dd s$, 由定理~\ref{thm:ito-properties} (4) 是鞅. (3) 线性: 对 $H$, $K$ 与 $H+K$ 取公共的局部化序列 (各自序列的最小者), 在其上用 $\mathscr{L}^2_T$ 中的线性. 停止性质: 对 $H^{(n)}$ 应用定理~\ref{thm:ito-properties} (6) 后令 $n\to\infty$.
\end{proof}

局部化的代价是$H\cdot B$ 一般只是局部鞅, 不一定是鞅, $\E\int_0^tH\dd B=0$ 与 It\^o 等距都可能不成立. 若 $H\in\mathscr{L}^2_T$, 它们当然仍成立. 下面两个引理在局部化的框架下代替 It\^o 等距.

\begin{lemma}[局部性]\label{lem:locality}
设 $H,K\in\mathscr{L}^2_{\mathrm{loc}}$, $A\in\mathcal{F}$, 且对 $\omega\in A$ 与几乎每个 $t\in[0,T]$ 有 $H_t(\omega)=K_t(\omega)$. 则几乎必然地在 $A$ 上, 对一切 $t\le T$, $(H\cdot B)_t=(K\cdot B)_t$.
\end{lemma}

\begin{proof}
令 $D=H-K$, $\sigma=\inf\{t: \int_0^tD_s^2\dd s>0\}\wedge T$. 它是停时 (由命题~\ref{prop:hitting-continuous} (2) 与通常条件), 且由连续性 $\int_0^\sigma D_s^2\dd s=0$. 所以 $D\1_{[0,\sigma]}\in\mathscr{L}^2_T$ 且 $\norm{D\1_{[0,\sigma]}}_T=0$, 其积分恒为零. 由停止性质, $(D\cdot B)_{t\wedge\sigma}=((D\1_{[0,\sigma]})\cdot B)_t=0$. 在 $A$ 上 $\int_0^TD^2\dd s=0$, 从而 $\sigma=T$, 于是 $(D\cdot B)_t=0$ 对一切 $t\le T$ 成立. 由线性, $D\cdot B=H\cdot B-K\cdot B$.
\end{proof}

例如, 设 $\xi$ 是 $\mathcal{F}_s$-可测的 (不必有界), $H\in\mathscr{L}^2_{\mathrm{loc}}$. 令 $\xi_K=(\xi\wedge K)\vee(-K)$. 由定理~\ref{thm:ito-properties} (5) (对 $H^{(n)}$ 应用后令 $n\to\infty$), $\int_0^T\xi_K\1_{(s,t]}H\dd B=\xi_K\int_s^tH\dd B$; 在 $\{\abs{\xi}\le K\}$ 上 $\xi_K=\xi$, 由局部性引理, 那里也有 $\int_0^T\xi\1_{(s,t]}H\dd B=\xi\int_s^tH\dd B$. 对 $K$ 取并, 提出已知因子的性质对不必有界的 $\xi$ 也成立.

\begin{lemma}[依概率收敛]\label{lem:ito-prob}
设 $H,K\in\mathscr{L}^2_{\mathrm{loc}}$. 则对任意 $\varepsilon,\delta>0$,
\[
\P\Big(\sup_{t\le T}\big\lvert(H\cdot B)_t-(K\cdot B)_t\big\rvert>\varepsilon\Big)\le\P\Big(\int_0^T(H_s-K_s)^2\dd s\ge\delta\Big)+\frac{\delta}{\varepsilon^2}.
\]
特别地, 若 $H^n,H\in\mathscr{L}^2_{\mathrm{loc}}$ 且 $\int_0^T(H^n_s-H_s)^2\dd s\to0$ 依概率成立, 则 $\sup_{t\le T}\abs{(H^n\cdot B)_t-(H\cdot B)_t}\to0$ 依概率成立.
\end{lemma}

\begin{proof}
令 $D=H-K$, $M=D\cdot B$, $\sigma=\inf\{t: \int_0^tD_s^2\dd s\ge\delta\}\wedge T$. 在事件 $\{\int_0^TD^2\dd s<\delta\}$ 上 $\sigma=T$, 所以
\[
\Big\{\sup_{t\le T}\abs{M_t}>\varepsilon\Big\}\subseteq\Big\{\int_0^TD_s^2\dd s\ge\delta\Big\}\cup\Big\{\sup_{t\le T}\abs{M_{t\wedge\sigma}}>\varepsilon\Big\}.
\]
由停止性质, $M_{\cdot\wedge\sigma}=(D\1_{[0,\sigma]})\cdot B$, 而 $\norm{D\1_{[0,\sigma]}}_T^2\le\delta$. 对非负下鞅 $M_{\cdot\wedge\sigma}^2$ 应用 Doob 不等式~\ref{thm:ct-doob} 与 It\^o 等距,
\[
\P\Big(\sup_{t\le T}\abs{M_{t\wedge\sigma}}>\varepsilon\Big)\le\frac{\E M_{T\wedge\sigma}^2}{\varepsilon^2}\le\frac{\delta}{\varepsilon^2}.
\]
第二个结论: 对给定的 $\varepsilon$ 与 $\eta>0$, 取 $\delta=\eta\varepsilon^2$, 则 $\limsup_n\P(\sup_t\abs{\cdot}>\varepsilon)\le\eta$.
\end{proof}

% ---------------------------------------------------------------------
\subsection{Riemann 和与 Wiener 积分}

It\^o 积分是用 $L^2$ 极限抽象地定义的. 对连续的被积过程, 它就是左端点 Riemann 和的极限, 这与 5.1 节的动机相吻合.

\begin{proposition}[Riemann 和逼近]\label{prop:riemann-approx}
设 $H$ 是轨道连续的适应过程, $\Pi_n: 0=t_0<\cdots<t_m=T$ 是 $[0,T]$ 的分割序列, $\abs{\Pi_n}\to0$. 则
\[
\sup_{t\le T}\Big\lvert\sum_{i}H_{t_{i-1}}\big(B_{t_i\wedge t}-B_{t_{i-1}\wedge t}\big)-\int_0^tH_s\dd B_s\Big\rvert\to0\qquad\text{依概率}.
\]
若 $H$ 还是有界的, 则收敛在 $\E\sup_{t\le T}\abs{\cdot}^2$ 的意义下成立.
\end{proposition}

\begin{proof}
令 $H^{\Pi}=H_0\1_{\{0\}}+\sum_iH_{t_{i-1}}\1_{(t_{i-1},t_i]}$. 它左连续且适应, 轨道有界, 所以 $H^\Pi\in\mathscr{L}^2_{\mathrm{loc}}$. 先说明 $(H^\Pi\cdot B)_t$ 等于 Riemann 和: 对 $K>0$, 令 $\xi^K_{i-1}=(H_{t_{i-1}}\wedge K)\vee(-K)$, 则 $H^{\Pi,K}=\xi_0^K\1_{\{0\}}+\sum_i\xi^K_{i-1}\1_{(t_{i-1},t_i]}\in\mathscr{S}_T$, 其积分是以 $\xi^K_{i-1}$ 代替 $H_{t_{i-1}}$ 的 Riemann 和. 在事件 $A_K=\{\max_i\abs{H_{t_{i-1}}}\le K\}$ 上 $H^\Pi=H^{\Pi,K}$, 且两个 Riemann 和相等, 由局部性引理~\ref{lem:locality}, 在 $A_K$ 上 $(H^\Pi\cdot B)_t$ 等于原来的 Riemann 和. 而 $\bigcup_KA_K=\Omega$.

由轨道在 $[0,T]$ 上一致连续, 对每个 $\omega$, $\sup_s\abs{H^{\Pi_n}_s-H_s}\to0$, 从而 $\int_0^T(H^{\Pi_n}_s-H_s)^2\dd s\to0$. 由引理~\ref{lem:ito-prob} 即得依概率收敛. 若 $\abs{H}\le C$, 则 $H^{\Pi_n}\in\mathscr{S}_T$ (无需截断), 由有界收敛定理 $\norm{H^{\Pi_n}-H}_T\to0$, 由定理~\ref{thm:ito-integral} 得 $\E\sup\abs{\cdot}^2\to0$.
\end{proof}

结合命题~\ref{prop:three-sums}, 得到第一个具体的 It\^o 积分:
\begin{equation}\label{eq:int-BdB}
\int_0^tB_s\dd B_s=\frac{1}{2}\big(B_t^2-t\big).
\end{equation}
(严格地说: 由命题~\ref{prop:riemann-approx}, 左端点和依概率收敛到左端, 由命题~\ref{prop:three-sums}, 它们在 $L^2$ 中收敛到右端, 依概率收敛的极限几乎必然唯一.) 修正项 $-\frac{t}{2}$ 使右端成为鞅, 这与定理~\ref{thm:ito-properties} (2) 一致.

被积函数是确定函数时, It\^o 积分是正态的, 称为 \textbf{Wiener 积分}.

\begin{proposition}[Wiener 积分]\label{prop:wiener-integral}
设 $f\in L^2[0,T]$ 是确定的函数, $X_t=\int_0^tf(s)\dd B_s$. 则 $X$ 是均值为零的连续高斯过程, 增量独立 ($X_t-X_s$ 与 $\mathcal{F}_s$ 独立), 且
\[
\Cov(X_s,X_t)=\int_0^{s\wedge t}f(u)^2\dd u,\qquad X_t\sim N\Big(0,\int_0^tf(u)^2\dd u\Big).
\]
\end{proposition}

\begin{proof}
确定的函数是循序可测的, $\norm{f}_T^2=\int_0^Tf^2\dd u<\infty$, 所以 $f\in\mathscr{L}^2_T$. 先设 $f=\sum_ic_i\1_{(t_{i-1},t_i]}$ 是阶梯函数, 则它是简单过程 (常数是 $\mathcal{F}_0$-可测的), $X_t=\sum_ic_i(B_{t_i\wedge t}-B_{t_{i-1}\wedge t})$ 是布朗运动在有限个时刻上取值的线性组合. 所以任意有限个 $X_{s_1},\dots,X_{s_k}$ 构成正态向量. 阶梯函数在 $L^2[0,T]$ 中稠密 (见~\cite{folland} 第 2, 6 章), 取阶梯函数 $f_n\to f$, 则由 It\^o 等距 $X^n_s\to X_s$ 于 $L^2$. 由推论~\ref{cor:gaussian-space} ($L^2$ 极限保持联合正态性), $X$ 是高斯过程, 均值为零. 协方差: 对 $s\le t$, 由鞅性与 It\^o 等距,
\[
\E[X_sX_t]=\E\big[X_s\E[X_t\,|\,\mathcal{F}_s]\big]=\E X_s^2=\int_0^sf^2\dd u.
\]
增量的独立性: 对 $A\in\mathcal{F}_s$ 与 $u\in\R$, 阶梯函数情形下 $X^n_t-X^n_s$ 是 $s$ 之后的布朗增量的函数, 与 $\mathcal{F}_s$ 独立, 所以 $\E[e^{\mathrm{i}u(X^n_t-X^n_s)};A]=\P(A)\E e^{\mathrm{i}u(X^n_t-X^n_s)}$. 令 $n\to\infty$, 由有界收敛定理 (沿几乎必然收敛的子列), $\E[e^{\mathrm{i}u(X_t-X_s)};A]=\P(A)\E e^{\mathrm{i}u(X_t-X_s)}$. $\P(A)>0$ 时, 上式表明 $X_t-X_s$ 在 $\P(\cdot\,|\,A)$ 下的特征函数与在 $\P$ 下相同, 由唯一性定理~\ref{thm:cf-uniqueness}, $\P(X_t-X_s\in C\,|\,A)=\P(X_t-X_s\in C)$, 即独立性.
\end{proof}

\begin{example}[分部积分与 $\int_0^tB_s\dd s$ 的分布]\label{ex:int-sdB}
取 $H_s=s$. 由 Abel 求和 (分部求和) 公式,
\[
\sum_it_{i-1}\big(B_{t_i}-B_{t_{i-1}}\big)=t\,B_t-\sum_iB_{t_i}(t_i-t_{i-1}).
\]
令 $\abs{\Pi_n}\to0$: 左端依概率收敛到 $\int_0^ts\dd B_s$ (命题~\ref{prop:riemann-approx}), 右端的和是连续函数 $s\mapsto B_s$ 的 Riemann 和, 逐条轨道收敛到 $\int_0^tB_s\dd s$. 于是
\[
\int_0^ts\dd B_s=tB_t-\int_0^tB_s\dd s,\qquad\text{即}\qquad\int_0^tB_s\dd s=\int_0^t(t-s)\dd B_s.
\]
由命题~\ref{prop:wiener-integral}, $\int_0^tB_s\dd s\sim N\big(0,\int_0^t(t-s)^2\dd s\big)=N\big(0,\frac{t^3}{3}\big)$. 这个结论也可以直接由 $\Var\int_0^tB_s\dd s=\int_0^t\int_0^t(s\wedge u)\dd s\dd u=\frac{t^3}{3}$ 验证, 但正态性由 Wiener 积分的表示一目了然.
\end{example}

\begin{remark}
本节的构造对任意 $T$ 成立, 而且不同的 $T$ 给出的积分在公共区间上一致 (由定理~\ref{thm:ito-properties} (6), 取 $\tau=T'<T$). 所以若 $H$ 在 $[0,\infty)$ 上循序可测, 且对每个 $T$ 有 $\int_0^TH_s^2\dd s<\infty$ $\as$, 则 $(H\cdot B)_t$ 对一切 $t\ge0$ 有定义, 是 $[0,\infty)$ 上的连续局部鞅. 今后不再区分有限与无限的时间范围.
\end{remark}

% =====================================================================
\section{It\^o 公式}
\label{sec:5-3}

普通微积分的链式法则 $\dd f(x_t)=f'(x_t)\dd x_t$ 来自 Taylor 展开中略去二阶项. 对布朗运动, $(\dd B_t)^2=\dd t$ 与 $\dd t$ 同阶, 二阶项必须保留:
\[
\dd f(B_t)=f'(B_t)\dd B_t+\frac{1}{2}f''(B_t)(\dd B_t)^2=f'(B_t)\dd B_t+\frac{1}{2}f''(B_t)\dd t.
\]
这就是 It\^o 公式. 本节先引入一类对光滑变换封闭的过程, 即 It\^o 过程, 然后严格证明 It\^o 公式, 并给出它的一系列应用.

% ---------------------------------------------------------------------
\subsection{It\^o 过程}

\begin{definition}[It\^o 过程]\label{def:ito-process}
称过程 $X$ 为 \textbf{It\^o 过程}, 如果
\begin{equation}\label{eq:ito-process}
X_t=X_0+\int_0^tb_s\dd s+\int_0^t\sigma_s\dd B_s,\qquad t\in[0,T],
\end{equation}
其中 $X_0$ 是 $\mathcal{F}_0$-可测的, $b$ 循序可测且 $\int_0^T\abs{b_s}\dd s<\infty$ $\as$, $\sigma\in\mathscr{L}^2_{\mathrm{loc}}$. 常把 \eqref{eq:ito-process} 简记为微分形式
\[
\dd X_t=b_t\dd t+\sigma_t\dd B_t,
\]
称 $b$ 为\textbf{漂移系数}, $\sigma$ 为\textbf{扩散系数}. 定义 $X$ 的二次变差为 $\ip{X}_t=\int_0^t\sigma_s^2\dd s$; 对两个 It\^o 过程 $X$, $Y$ (扩散系数分别为 $\sigma$, $\eta$), 定义 $\ip{X,Y}_t=\int_0^t\sigma_s\eta_s\dd s$.
\end{definition}

It\^o 过程是一个连续的有限变差过程 $A_t=\int_0^tb_s\dd s$ 与一个连续局部鞅 $M_t=\int_0^t\sigma_s\dd B_s$ 之和 (加上初值). 前者描述 ``趋势'', 后者描述 ``波动''. 这种分解是唯一的.

\begin{proposition}[It\^o 过程的性质]\label{prop:ito-process}
设 $X$, $Y$ 是 It\^o 过程.
\begin{enumerate}
\item (分解唯一) 若 $X$ 有两种形如 \eqref{eq:ito-process} 的表示, 系数分别为 $(b,\sigma)$ 与 $(b',\sigma')$, 则 $b=b'$, $\sigma=\sigma'$ 对 $\lambda\otimes\P$-几乎每个 $(t,\omega)$ 成立.
\item (二次变差) 对 $[0,t]$ 的分割序列 $\Pi_n$, $\abs{\Pi_n}\to0$, 记 $\Delta X_i=X_{t_i}-X_{t_{i-1}}$, 则
\[
\sum_i(\Delta X_i)^2\to\ip{X}_t,\qquad\sum_i\Delta X_i\Delta Y_i\to\ip{X,Y}_t\qquad\text{依概率}.
\]
\end{enumerate}
\end{proposition}

\begin{proof}
(1) 两种表示相减, $\int_0^t(b_s-b'_s)\dd s=\int_0^t(\sigma'_s-\sigma_s)\dd B_s$. 右端是从 $0$ 出发的连续局部鞅, 左端是连续的有限变差过程. 由停时截断, 使二者及其全变差有界后, 由定理~\ref{thm:no-finite-variation}, 它们恒为零. 于是 $\int_0^t(\sigma'_s-\sigma_s)\dd B_s\equiv0$, 其二次变差 $\int_0^t(\sigma'_s-\sigma_s)^2\dd s$ 也恒为零 (定理~\ref{thm:ito-local} (2) 与二次变差的唯一性), 所以 $\sigma=\sigma'$ 几乎处处. 对几乎每个 $\omega$, 函数 $t\mapsto\int_0^t(b_s-b'_s)\dd s$ 恒为零, 从而被积函数对几乎每个 $t$ 为零.

(2) 记 $\Delta A_i=\int_{t_{i-1}}^{t_i}b_s\dd s$, $\Delta M_i=\int_{t_{i-1}}^{t_i}\sigma_s\dd B_s$. 则
\[
\sum_i(\Delta X_i)^2=\sum_i(\Delta M_i)^2+\sum_i\Delta A_i\big(\Delta A_i+2\Delta M_i\big).
\]
由连续局部鞅的二次变差定理~\ref{thm:quadratic-variation} 与 $\ip{M}_t=\int_0^t\sigma_s^2\dd s$ (定理~\ref{thm:ito-local} (2)), 第一项依概率收敛到 $\ip{X}_t$. 第二项的绝对值不超过
\[
\Big(\max_i\abs{\Delta A_i}+2\max_i\abs{\Delta M_i}\Big)\int_0^t\abs{b_s}\dd s,
\]
由 $A$, $M$ 的轨道一致连续, 它逐条轨道地趋于零. 对 $\ip{X,Y}$, 用极化恒等式 $\Delta X\Delta Y=\frac{1}{4}[(\Delta X+\Delta Y)^2-(\Delta X-\Delta Y)^2]$, 并注意 $X\pm Y$ 是扩散系数为 $\sigma\pm\eta$ 的 It\^o 过程.
\end{proof}

对 It\^o 过程 $X$ 与循序可测过程 $H$, 若 $\int_0^T\abs{H_sb_s}\dd s<\infty$ 且 $\int_0^TH_s^2\sigma_s^2\dd s<\infty$ $\as$ (例如 $H$ 轨道连续), 定义
\[
\int_0^tH_s\dd X_s=\int_0^tH_sb_s\dd s+\int_0^tH_s\sigma_s\dd B_s,
\]
即 $H\dd X=Hb\dd t+H\sigma\dd B$. 它仍是 It\^o 过程. 二次变差的计算可以归结为下面的形式乘法表, 其中 $\dd\ip{X}_t=(\dd X_t)^2$:
\begin{equation}\label{eq:mult-table}
\dd t\cdot\dd t=0,\qquad\dd t\cdot\dd B_t=0,\qquad\dd B_t\cdot\dd B_t=\dd t;\qquad\text{从而}\quad(\dd X_t)^2=\sigma_t^2\dd t.
\end{equation}

下面两个引理是 It\^o 公式证明的主要步骤: 它们分别把 Taylor 展开中的一阶项与二阶项识别为积分.

\begin{lemma}[关于 It\^o 过程的 Riemann 和]\label{lem:riemann-dX}
设 $X$ 是 It\^o 过程, $g$ 是轨道连续的适应过程, $\Pi_n$ 是 $[0,t]$ 的分割序列, $\abs{\Pi_n}\to0$. 则
\[
\sum_ig_{t_{i-1}}\Delta X_i\to\int_0^tg_s\dd X_s\qquad\text{依概率}.
\]
\end{lemma}

\begin{proof}
沿用命题~\ref{prop:ito-process} 证明中的记号. 对漂移部分,
\[
\Big\lvert\sum_ig_{t_{i-1}}\Delta A_i-\int_0^tg_sb_s\dd s\Big\rvert\le\sup_{\abs{u-v}\le\abs{\Pi_n}}\abs{g_u-g_v}\cdot\int_0^t\abs{b_s}\dd s\to0.
\]
对鞅部分, 令 $g^n=g_0\1_{\{0\}}+\sum_ig_{t_{i-1}}\1_{(t_{i-1},t_i]}$. 由提出已知因子的性质 (定理~\ref{thm:ito-properties} (5) 及其在局部化后对无界因子的推广, 见引理~\ref{lem:locality} 之后的说明), $\int_0^tg^n_s\sigma_s\dd B_s=\sum_ig_{t_{i-1}}\Delta M_i$. 由于 $\int_0^t(g^n_s-g_s)^2\sigma_s^2\dd s\le\sup_s\abs{g^n_s-g_s}^2\int_0^t\sigma_s^2\dd s\to0$ 逐条轨道成立, 由引理~\ref{lem:ito-prob}, $\sum_ig_{t_{i-1}}\Delta M_i\to\int_0^tg_s\sigma_s\dd B_s$ 依概率成立.
\end{proof}

\begin{lemma}[带权的二次变差]\label{lem:weighted-qv}
设 $X$ 是 It\^o 过程, $g$ 是轨道连续的适应过程, $\Pi_n$ 是 $[0,t]$ 的\emph{嵌套}分割序列 ($\Pi_n\subseteq\Pi_{n+1}$, 例如二进分割), $\abs{\Pi_n}\to0$. 则
\[
\sum_ig_{t_{i-1}}(\Delta X_i)^2\to\int_0^tg_s\dd\ip{X}_s=\int_0^tg_s\sigma_s^2\dd s\qquad\text{依概率}.
\]
\end{lemma}

\begin{proof}
记 $S_n=\sum_ig_{t_{i-1}}(\Delta X_i)^2$, $I=\int_0^tg_s\sigma_s^2\dd s$, $V=\ip{X}_t$, $Q_n=\sum_i(\Delta X_i)^2$, 并记 $w(\delta)=\sup_{u,v\le t,\,\abs{u-v}\le\delta}\abs{g_u-g_v}$, 由 $g$ 的一致连续性, 对每个 $\omega$, $w(\delta)\to0$ ($\delta\to0$).

固定 $m$, 记 $\Pi_m=\{s_0<s_1<\cdots<s_k\}$. 对 $n\ge m$, $\Pi_n$ 加细了 $\Pi_m$, 于是 $\Pi_n$ 的每个小区间 $[t_{i-1},t_i]$ 落在某个 $[s_{j-1},s_j]$ 中. 记 $Q_n^{(j)}$ 为 $(\Delta X_i)^2$ 在落入 $[s_{j-1},s_j]$ 的那些 $i$ 上的和. 则
\[
\Big\lvert S_n-\sum_{j=1}^kg_{s_{j-1}}Q_n^{(j)}\Big\rvert\le w(\abs{\Pi_m})\,Q_n,\qquad\Big\lvert\sum_{j=1}^kg_{s_{j-1}}\big(\ip{X}_{s_j}-\ip{X}_{s_{j-1}}\big)-I\Big\rvert\le w(\abs{\Pi_m})\,V.
\]
由命题~\ref{prop:ito-process} (2) (分别应用于 $[0,s_j]$ 与 $[0,s_{j-1}]$ 上由 $\Pi_n$ 诱导的分割), 当 $n\to\infty$ 时 $Q_n^{(j)}\to\ip{X}_{s_j}-\ip{X}_{s_{j-1}}$ 依概率成立, 而 $\max_j\abs{g_{s_{j-1}}}<\infty$, 所以有限和
\[
D_n^{(m)}:=\sum_{j=1}^kg_{s_{j-1}}\Big(Q_n^{(j)}-\big(\ip{X}_{s_j}-\ip{X}_{s_{j-1}}\big)\Big)\to0\qquad\text{依概率}.
\]
综合三式, $\abs{S_n-I}\le w(\abs{\Pi_m})(Q_n+V)+\abs{D_n^{(m)}}$. 对 $\varepsilon>0$,
\[
\P(\abs{S_n-I}>2\varepsilon)\le\P\big(\abs{Q_n-V}>1\big)+\P\big(w(\abs{\Pi_m})(2V+1)>\varepsilon\big)+\P\big(\abs{D_n^{(m)}}>\varepsilon\big).
\]
令 $n\to\infty$, 第一项与第三项趋于零, 所以 $\limsup_n\P(\abs{S_n-I}>2\varepsilon)\le\P\big(w(\abs{\Pi_m})(2V+1)>\varepsilon\big)$. 再令 $m\to\infty$, 由于 $w(\abs{\Pi_m})(2V+1)\to0$ 逐条轨道成立, 右端趋于零.
\end{proof}

% ---------------------------------------------------------------------
\subsection{It\^o 公式}

记 $C^{1,2}$ 为 $[0,T]\times\R$ 上关于 $t$ 一阶、关于 $x$ 二阶连续可微的函数全体, 偏导数记为 $\partial_tf$, $\partial_xf$, $\partial_{xx}f$.

\begin{theorem}[It\^o 公式]\label{thm:ito-formula}
设 $X$ 是 It\^o 过程, $\dd X_t=b_t\dd t+\sigma_t\dd B_t$, $f\in C^{1,2}$. 则几乎必然地, 对一切 $t\in[0,T]$,
\begin{equation}\label{eq:ito-formula}
f(t,X_t)=f(0,X_0)+\int_0^t\partial_tf(s,X_s)\dd s+\int_0^t\partial_xf(s,X_s)\dd X_s+\frac{1}{2}\int_0^t\partial_{xx}f(s,X_s)\dd\ip{X}_s.
\end{equation}
写成微分形式, 并代入 $\dd X$ 与 $\dd\ip{X}$,
\[
\dd f(t,X_t)=\Big(\partial_tf+b_t\,\partial_xf+\frac{1}{2}\sigma_t^2\,\partial_{xx}f\Big)(t,X_t)\dd t+\sigma_t\,\partial_xf(t,X_t)\dd B_t.
\]
特别地, 对 $X=B$ 与 $f=f(x)$, $f(B_t)=f(B_0)+\int_0^tf'(B_s)\dd B_s+\frac{1}{2}\int_0^tf''(B_s)\dd s$.
\end{theorem}

\eqref{eq:ito-formula} 右端的每一项都有意义: $s\mapsto\partial_tf(s,X_s)$ 等都是轨道连续的适应过程, 在 $[0,t]$ 上有界. 证明的思路与普通链式法则相同, 只是 Taylor 展开要展开到二阶, 并用引理~\ref{lem:weighted-qv} 处理二阶项.

\begin{proof}
固定 $t\in(0,T]$, 取 $[0,t]$ 的二进分割 $\Pi_n$: $t_i=it2^{-n}$. 固定 $\omega$ (在下面用到的几乎必然事件中), 轨道 $s\mapsto X_s(\omega)$ 在 $[0,t]$ 上连续, 取值于某个紧区间 $K=[-R,R]$. $\partial_tf$ 与 $\partial_{xx}f$ 在紧集 $[0,t]\times K$ 上一致连续, 记它们的连续模为
\[
\rho_1(\delta)=\sup\big\{\abs{\partial_tf(s,x)-\partial_tf(s,y)}: s\le t,\ x,y\in K,\ \abs{x-y}\le\delta\big\},
\]
$\rho_2(\delta)$ 类似 (把 $\partial_tf$ 换成 $\partial_{xx}f$), 则 $\rho_1(\delta),\rho_2(\delta)\to0$ ($\delta\to0$). 记 $\delta_n=\max_{\abs{u-v}\le\abs{\Pi_n}}\abs{X_u-X_v}$, 由一致连续性 $\delta_n\to0$.

把总增量拆为
\[
f(t,X_t)-f(0,X_0)=\sum_i\Big[f(t_i,X_{t_i})-f(t_{i-1},X_{t_i})\Big]+\sum_i\Big[f(t_{i-1},X_{t_i})-f(t_{i-1},X_{t_{i-1}})\Big].
\]
由微积分基本定理, 第一个和中的每一项等于 $\int_{t_{i-1}}^{t_i}\partial_tf(s,X_{t_i})\dd s$. 对第二个和中的每一项, 用带积分余项的 Taylor 公式 $\phi(y)-\phi(x)=\phi'(x)(y-x)+(y-x)^2\int_0^1(1-\theta)\phi''(x+\theta(y-x))\dd\theta$, 其中 $\phi=f(t_{i-1},\cdot)$, $x=X_{t_{i-1}}$, $y=X_{t_i}$. 注意 $\int_0^1(1-\theta)\dd\theta=\frac{1}{2}$, 得
\[
f(t_{i-1},X_{t_i})-f(t_{i-1},X_{t_{i-1}})=\partial_xf(t_{i-1},X_{t_{i-1}})\Delta X_i+\frac{1}{2}\partial_{xx}f(t_{i-1},X_{t_{i-1}})(\Delta X_i)^2+r_i,
\]
其中
\[
r_i=(\Delta X_i)^2\int_0^1(1-\theta)\Big[\partial_{xx}f\big(t_{i-1},X_{t_{i-1}}+\theta\Delta X_i\big)-\partial_{xx}f\big(t_{i-1},X_{t_{i-1}}\big)\Big]\dd\theta,
\]
满足 $\abs{r_i}\le\frac{1}{2}\rho_2(\delta_n)(\Delta X_i)^2$.
于是 $f(t,X_t)-f(0,X_0)=I_1^n+I_2^n+\frac{1}{2}I_3^n+R^n$, 其中
\begin{align*}
I_1^n&=\sum_i\int_{t_{i-1}}^{t_i}\partial_tf(s,X_{t_i})\dd s, & I_2^n&=\sum_i\partial_xf(t_{i-1},X_{t_{i-1}})\Delta X_i,\\
I_3^n&=\sum_i\partial_{xx}f(t_{i-1},X_{t_{i-1}})(\Delta X_i)^2, & R^n&=\sum_ir_i.
\end{align*}
令 $n\to\infty$, 对各项分别求极限:
\begin{itemize}
\item $\abs{I_1^n-\int_0^t\partial_tf(s,X_s)\dd s}\le\sum_i\int_{t_{i-1}}^{t_i}\abs{\partial_tf(s,X_{t_i})-\partial_tf(s,X_s)}\dd s\le t\,\rho_1(\delta_n)\to0$, 逐条轨道成立.
\item 由引理~\ref{lem:riemann-dX} (取 $g_s=\partial_xf(s,X_s)$), $I_2^n\to\int_0^t\partial_xf(s,X_s)\dd X_s$ 依概率.
\item 由引理~\ref{lem:weighted-qv} (取 $g_s=\partial_{xx}f(s,X_s)$, 二进分割是嵌套的), $I_3^n\to\int_0^t\partial_{xx}f(s,X_s)\dd\ip{X}_s$ 依概率.
\item $\abs{R^n}\le\frac{1}{2}\rho_2(\delta_n)\sum_i(\Delta X_i)^2$. 由命题~\ref{prop:ito-process} (2), $\sum_i(\Delta X_i)^2\to\ip{X}_t$ 依概率, 而 $\rho_2(\delta_n)\to0$ 逐条轨道, 所以 $R^n\to0$ 依概率.
\end{itemize}
左端与 $n$ 无关, 依概率收敛的极限几乎必然唯一, 所以 \eqref{eq:ito-formula} 对固定的 $t$ 几乎必然成立. 对有理数 $t$ 取可列交, 而两端都关于 $t$ 连续, 所以几乎必然地 \eqref{eq:ito-formula} 对一切 $t\in[0,T]$ 成立.
\end{proof}

\begin{figure}[htbp]
\centering
\begin{tikzpicture}
\begin{axis}[SAaxis, axis x line=bottom, axis y line=left, width=12cm, height=6cm, xmin=0, xmax=2.06, ymin=-1.6, ymax=5.3,
  xtick={0,0.5,1,1.5,2}, ytick={-1,0,1,2,3,4,5}, xlabel={$t$},
  legend style={font=\small,draw=none,fill=none,at={(0.5,-0.2)},anchor=north}, legend columns=3, legend cell align=left]
  \addplot[SADark,thin] coordinates {(0.000,0.000) (0.005,0.003) (0.010,0.000) (0.015,0.001) (0.020,0.014) (0.025,0.013) (0.030,0.035) (0.035,0.035) (0.040,0.026) (0.045,0.033) (0.050,0.015) (0.055,0.041) (0.060,0.008) (0.065,0.002) (0.070,0.004) (0.075,0.001) (0.080,0.005) (0.085,0.020) (0.090,0.041) (0.095,0.029) (0.100,0.024) (0.105,0.008) (0.110,0.002) (0.115,0.002) (0.120,0.001) (0.125,0.001) (0.130,0.000) (0.135,0.017) (0.140,0.046) (0.145,0.030) (0.150,0.071) (0.155,0.057) (0.160,0.058) (0.165,0.062) (0.170,0.066) (0.175,0.053) (0.180,0.068) (0.185,0.050) (0.190,0.090) (0.195,0.088) (0.200,0.070) (0.205,0.135) (0.210,0.163) (0.215,0.327) (0.220,0.443) (0.225,0.432) (0.230,0.640) (0.235,0.515) (0.240,0.475) (0.245,0.599) (0.250,0.670) (0.255,0.636) (0.260,0.692) (0.265,0.640) (0.270,0.629) (0.275,0.592) (0.280,0.623) (0.285,0.758) (0.290,0.983) (0.295,0.994) (0.300,0.935) (0.305,0.797) (0.310,0.682) (0.315,0.700) (0.320,0.637) (0.325,0.720) (0.330,0.496) (0.335,0.465) (0.340,0.379) (0.345,0.339) (0.350,0.349) (0.355,0.332) (0.360,0.147) (0.365,0.111) (0.370,0.059) (0.375,0.116) (0.380,0.070) (0.385,0.032) (0.390,0.049) (0.395,0.040) (0.400,0.104) (0.405,0.070) (0.410,0.050) (0.415,0.007) (0.420,0.014) (0.425,0.032) (0.430,0.004) (0.435,0.002) (0.440,0.024) (0.445,0.018) (0.450,0.009) (0.455,0.004) (0.460,0.000) (0.465,0.001) (0.470,0.003) (0.475,0.006) (0.480,0.006) (0.485,0.002) (0.490,0.000) (0.495,0.002) (0.500,0.003) (0.505,0.014) (0.510,0.006) (0.515,0.003) (0.520,0.005) (0.525,0.004) (0.530,0.004) (0.535,0.036) (0.540,0.059) (0.545,0.042) (0.550,0.033) (0.555,0.041) (0.560,0.040) (0.565,0.051) (0.570,0.084) (0.575,0.051) (0.580,0.057) (0.585,0.098) (0.590,0.109) (0.595,0.134) (0.600,0.114) (0.605,0.082) (0.610,0.093) (0.615,0.097) (0.620,0.070) (0.625,0.070) (0.630,0.068) (0.635,0.119) (0.640,0.120) (0.645,0.129) (0.650,0.090) (0.655,0.057) (0.660,0.044) (0.665,0.027) (0.670,0.076) (0.675,0.044) (0.680,0.090) (0.685,0.055) (0.690,0.063) (0.695,0.099) (0.700,0.129) (0.705,0.156) (0.710,0.183) (0.715,0.235) (0.720,0.281) (0.725,0.327) (0.730,0.443) (0.735,0.475) (0.740,0.457) (0.745,0.347) (0.750,0.550) (0.755,0.563) (0.760,0.643) (0.765,0.738) (0.770,0.632) (0.775,0.841) (0.780,0.852) (0.785,1.128) (0.790,1.153) (0.795,1.141) (0.800,1.118) (0.805,1.109) (0.810,1.160) (0.815,1.100) (0.820,0.979) (0.825,0.890) (0.830,1.058) (0.835,1.202) (0.840,1.434) (0.845,1.511) (0.850,1.527) (0.855,1.474) (0.860,1.681) (0.865,1.600) (0.870,1.545) (0.875,1.401) (0.880,1.412) (0.885,1.733) (0.890,2.103) (0.895,2.023) (0.900,2.431) (0.905,2.083) (0.910,2.163) (0.915,2.036) (0.920,2.093) (0.925,2.099) (0.930,2.439) (0.935,2.376) (0.940,2.288) (0.945,2.044) (0.950,2.092) (0.955,2.080) (0.960,2.291) (0.965,2.184) (0.970,2.458) (0.975,2.568) (0.980,2.093) (0.985,1.810) (0.990,1.981) (0.995,2.034) (1.000,2.241) (1.005,2.225) (1.010,2.027) (1.015,1.874) (1.020,1.769) (1.025,1.676) (1.030,1.988) (1.035,1.801) (1.040,1.654) (1.045,1.878) (1.050,2.004) (1.055,2.077) (1.060,1.704) (1.065,1.914) (1.070,1.738) (1.075,1.774) (1.080,1.980) (1.085,1.836) (1.090,1.666) (1.095,1.405) (1.100,1.265) (1.105,1.040) (1.110,1.134) (1.115,1.323) (1.120,1.114) (1.125,1.281) (1.130,1.204) (1.135,1.292) (1.140,1.243) (1.145,1.156) (1.150,1.145) (1.155,1.144) (1.160,1.307) (1.165,1.015) (1.170,1.154) (1.175,1.064) (1.180,0.873) (1.185,0.937) (1.190,0.674) (1.195,0.735) (1.200,0.796) (1.205,0.846) (1.210,0.745) (1.215,0.960) (1.220,1.077) (1.225,1.068) (1.230,0.936) (1.235,1.200) (1.240,1.049) (1.245,1.018) (1.250,0.684) (1.255,0.615) (1.260,0.470) (1.265,0.387) (1.270,0.419) (1.275,0.469) (1.280,0.412) (1.285,0.331) (1.290,0.347) (1.295,0.394) (1.300,0.385) (1.305,0.335) (1.310,0.345) (1.315,0.353) (1.320,0.267) (1.325,0.314) (1.330,0.207) (1.335,0.234) (1.340,0.144) (1.345,0.257) (1.350,0.219) (1.355,0.248) (1.360,0.167) (1.365,0.213) (1.370,0.292) (1.375,0.278) (1.380,0.316) (1.385,0.246) (1.390,0.304) (1.395,0.355) (1.400,0.389) (1.405,0.345) (1.410,0.257) (1.415,0.266) (1.420,0.260) (1.425,0.339) (1.430,0.336) (1.435,0.308) (1.440,0.391) (1.445,0.305) (1.450,0.317) (1.455,0.282) (1.460,0.218) (1.465,0.193) (1.470,0.234) (1.475,0.218) (1.480,0.157) (1.485,0.169) (1.490,0.136) (1.495,0.133) (1.500,0.136) (1.505,0.150) (1.510,0.156) (1.515,0.137) (1.520,0.046) (1.525,0.036) (1.530,0.011) (1.535,0.028) (1.540,0.020) (1.545,0.053) (1.550,0.041) (1.555,0.024) (1.560,0.066) (1.565,0.154) (1.570,0.220) (1.575,0.313) (1.580,0.267) (1.585,0.220) (1.590,0.285) (1.595,0.269) (1.600,0.305) (1.605,0.411) (1.610,0.616) (1.615,0.582) (1.620,0.536) (1.625,0.547) (1.630,0.554) (1.635,0.553) (1.640,0.616) (1.645,0.573) (1.650,0.754) (1.655,0.797) (1.660,0.827) (1.665,0.782) (1.670,0.784) (1.675,0.933) (1.680,0.987) (1.685,1.125) (1.690,1.194) (1.695,1.055) (1.700,1.147) (1.705,1.088) (1.710,1.167) (1.715,1.676) (1.720,1.473) (1.725,1.510) (1.730,1.648) (1.735,1.612) (1.740,1.662) (1.745,1.576) (1.750,1.941) (1.755,1.925) (1.760,1.785) (1.765,1.970) (1.770,2.216) (1.775,2.416) (1.780,2.323) (1.785,2.401) (1.790,2.453) (1.795,2.463) (1.800,2.527) (1.805,2.836) (1.810,3.215) (1.815,3.433) (1.820,3.302) (1.825,3.433) (1.830,3.815) (1.835,3.782) (1.840,4.175) (1.845,4.499) (1.850,4.355) (1.855,3.754) (1.860,3.513) (1.865,3.650) (1.870,3.351) (1.875,3.214) (1.880,2.924) (1.885,2.784) (1.890,3.174) (1.895,3.382) (1.900,3.551) (1.905,3.367) (1.910,3.042) (1.915,3.250) (1.920,3.495) (1.925,3.569) (1.930,3.558) (1.935,3.275) (1.940,3.221) (1.945,3.064) (1.950,3.286) (1.955,3.431) (1.960,4.054) (1.965,4.167) (1.970,4.874) (1.975,4.681) (1.980,4.641) (1.985,4.752) (1.990,5.105) (1.995,4.973) (2.000,4.715)};
  \addplot[SABlue,thin] coordinates {(0.000,0.000) (0.005,-0.005) (0.010,-0.012) (0.015,-0.019) (0.020,-0.016) (0.025,-0.018) (0.030,0.000) (0.035,-0.004) (0.040,-0.014) (0.045,-0.010) (0.050,-0.030) (0.055,-0.011) (0.060,-0.049) (0.065,-0.056) (0.070,-0.056) (0.075,-0.060) (0.080,-0.059) (0.085,-0.047) (0.090,-0.029) (0.095,-0.046) (0.100,-0.054) (0.105,-0.077) (0.110,-0.098) (0.115,-0.098) (0.120,-0.104) (0.125,-0.107) (0.130,-0.116) (0.135,-0.104) (0.140,-0.082) (0.145,-0.100) (0.150,-0.063) (0.155,-0.078) (0.160,-0.083) (0.165,-0.081) (0.170,-0.078) (0.175,-0.092) (0.180,-0.081) (0.185,-0.100) (0.190,-0.065) (0.195,-0.072) (0.200,-0.091) (0.205,-0.034) (0.210,-0.010) (0.215,0.142) (0.220,0.255) (0.225,0.244) (0.230,0.446) (0.235,0.315) (0.240,0.270) (0.245,0.390) (0.250,0.460) (0.255,0.423) (0.260,0.476) (0.265,0.420) (0.270,0.406) (0.275,0.367) (0.280,0.397) (0.285,0.525) (0.290,0.741) (0.295,0.746) (0.300,0.686) (0.305,0.545) (0.310,0.429) (0.315,0.445) (0.320,0.376) (0.325,0.456) (0.330,0.227) (0.335,0.193) (0.340,0.106) (0.345,0.063) (0.350,0.066) (0.355,0.046) (0.360,-0.151) (0.365,-0.188) (0.370,-0.246) (0.375,-0.196) (0.380,-0.246) (0.385,-0.287) (0.390,-0.273) (0.395,-0.285) (0.400,-0.227) (0.405,-0.264) (0.410,-0.295) (0.415,-0.344) (0.420,-0.346) (0.425,-0.334) (0.430,-0.369) (0.435,-0.379) (0.440,-0.360) (0.445,-0.375) (0.450,-0.386) (0.455,-0.396) (0.460,-0.400) (0.465,-0.402) (0.470,-0.405) (0.475,-0.415) (0.480,-0.417) (0.485,-0.422) (0.490,-0.427) (0.495,-0.432) (0.500,-0.433) (0.505,-0.426) (0.510,-0.438) (0.515,-0.442) (0.520,-0.451) (0.525,-0.462) (0.530,-0.464) (0.535,-0.438) (0.540,-0.420) (0.545,-0.440) (0.550,-0.454) (0.555,-0.447) (0.560,-0.460) (0.565,-0.455) (0.570,-0.428) (0.575,-0.469) (0.580,-0.469) (0.585,-0.435) (0.590,-0.424) (0.595,-0.403) (0.600,-0.429) (0.605,-0.469) (0.610,-0.467) (0.615,-0.467) (0.620,-0.495) (0.625,-0.497) (0.630,-0.500) (0.635,-0.460) (0.640,-0.468) (0.645,-0.460) (0.650,-0.501) (0.655,-0.538) (0.660,-0.553) (0.665,-0.575) (0.670,-0.538) (0.675,-0.572) (0.680,-0.529) (0.685,-0.568) (0.690,-0.567) (0.695,-0.535) (0.700,-0.514) (0.705,-0.499) (0.710,-0.474) (0.715,-0.425) (0.720,-0.381) (0.725,-0.344) (0.730,-0.232) (0.735,-0.201) (0.740,-0.221) (0.745,-0.334) (0.750,-0.140) (0.755,-0.132) (0.760,-0.054) (0.765,0.038) (0.770,-0.075) (0.775,0.123) (0.780,0.130) (0.785,0.400) (0.790,0.418) (0.795,0.401) (0.800,0.375) (0.805,0.360) (0.810,0.406) (0.815,0.341) (0.820,0.216) (0.825,0.124) (0.830,0.281) (0.835,0.420) (0.840,0.642) (0.845,0.718) (0.850,0.730) (0.855,0.676) (0.860,0.881) (0.865,0.798) (0.870,0.739) (0.875,0.591) (0.880,0.598) (0.885,0.910) (0.890,1.275) (0.895,1.191) (0.900,1.593) (0.905,1.241) (0.910,1.317) (0.915,1.184) (0.920,1.239) (0.925,1.237) (0.930,1.556) (0.935,1.492) (0.940,1.400) (0.945,1.149) (0.950,1.194) (0.955,1.180) (0.960,1.387) (0.965,1.279) (0.970,1.549) (0.975,1.656) (0.980,1.169) (0.985,0.882) (0.990,1.051) (0.995,1.098) (1.000,1.299) (1.005,1.281) (1.010,1.075) (1.015,0.920) (1.020,0.814) (1.025,0.719) (1.030,1.020) (1.035,0.827) (1.040,0.676) (1.045,0.897) (1.050,1.019) (1.055,1.091) (1.060,0.712) (1.065,0.918) (1.070,0.740) (1.075,0.774) (1.080,0.975) (1.085,0.828) (1.090,0.652) (1.095,0.385) (1.100,0.242) (1.105,0.013) (1.110,0.105) (1.115,0.289) (1.120,0.077) (1.125,0.241) (1.130,0.157) (1.135,0.238) (1.140,0.184) (1.145,0.090) (1.150,0.076) (1.155,0.071) (1.160,0.229) (1.165,-0.068) (1.170,0.069) (1.175,-0.024) (1.180,-0.219) (1.185,-0.156) (1.190,-0.429) (1.195,-0.371) (1.200,-0.314) (1.205,-0.273) (1.210,-0.376) (1.215,-0.164) (1.220,-0.049) (1.225,-0.061) (1.230,-0.198) (1.235,0.056) (1.240,-0.096) (1.245,-0.129) (1.250,-0.476) (1.255,-0.551) (1.260,-0.702) (1.265,-0.790) (1.270,-0.760) (1.275,-0.712) (1.280,-0.771) (1.285,-0.856) (1.290,-0.841) (1.295,-0.796) (1.300,-0.809) (1.305,-0.859) (1.310,-0.851) (1.315,-0.846) (1.320,-0.938) (1.325,-0.892) (1.330,-1.005) (1.335,-0.981) (1.340,-1.081) (1.345,-0.974) (1.350,-1.018) (1.355,-0.991) (1.360,-1.082) (1.365,-1.041) (1.370,-0.968) (1.375,-0.985) (1.380,-0.952) (1.385,-1.031) (1.390,-0.976) (1.395,-0.925) (1.400,-0.893) (1.405,-0.939) (1.410,-1.033) (1.415,-1.025) (1.420,-1.033) (1.425,-0.958) (1.430,-0.965) (1.435,-0.998) (1.440,-0.918) (1.445,-1.006) (1.450,-0.998) (1.455,-1.037) (1.460,-1.109) (1.465,-1.136) (1.470,-1.101) (1.475,-1.120) (1.480,-1.186) (1.485,-1.179) (1.490,-1.218) (1.495,-1.223) (1.500,-1.221) (1.505,-1.210) (1.510,-1.209) (1.515,-1.232) (1.520,-1.334) (1.525,-1.346) (1.530,-1.374) (1.535,-1.359) (1.540,-1.368) (1.545,-1.339) (1.550,-1.354) (1.555,-1.373) (1.560,-1.335) (1.565,-1.254) (1.570,-1.189) (1.575,-1.101) (1.580,-1.152) (1.585,-1.203) (1.590,-1.148) (1.595,-1.164) (1.600,-1.130) (1.605,-1.030) (1.610,-0.831) (1.615,-0.871) (1.620,-0.921) (1.625,-0.913) (1.630,-0.908) (1.635,-0.914) (1.640,-0.855) (1.645,-0.903) (1.650,-0.725) (1.655,-0.691) (1.660,-0.669) (1.665,-0.718) (1.670,-0.723) (1.675,-0.579) (1.680,-0.532) (1.685,-0.396) (1.690,-0.330) (1.695,-0.473) (1.700,-0.388) (1.705,-0.454) (1.710,-0.384) (1.715,0.113) (1.720,-0.093) (1.725,-0.062) (1.730,0.073) (1.735,0.032) (1.740,0.081) (1.745,-0.007) (1.750,0.353) (1.755,0.335) (1.760,0.193) (1.765,0.370) (1.770,0.606) (1.775,0.802) (1.780,0.702) (1.785,0.774) (1.790,0.824) (1.795,0.829) (1.800,0.889) (1.805,1.192) (1.810,1.568) (1.815,1.779) (1.820,1.646) (1.825,1.767) (1.830,2.144) (1.835,2.103) (1.840,2.489) (1.845,2.806) (1.850,2.659) (1.855,2.051) (1.860,1.804) (1.865,1.935) (1.870,1.632) (1.875,1.494) (1.880,1.198) (1.885,1.055) (1.890,1.440) (1.895,1.644) (1.900,1.807) (1.905,1.619) (1.910,1.288) (1.915,1.492) (1.920,1.733) (1.925,1.795) (1.930,1.779) (1.935,1.493) (1.940,1.436) (1.945,1.275) (1.950,1.495) (1.955,1.636) (1.960,2.250) (1.965,2.357) (1.970,3.053) (1.975,2.858) (1.980,2.817) (1.985,2.926) (1.990,3.274) (1.995,3.138) (2.000,2.876)};
  \addplot[SARed,thick,dashed] coordinates {(0,0) (2,2)};
  \legend{$B_t^2$, $2\int_0^tB_s\dd B_s$, $t$}
\end{axis}
\end{tikzpicture}
\caption{It\^o 公式 $B_t^2=2\int_0^tB_s\dd B_s+t$ 的一个实现. 鞅部分 $2\int_0^tB\dd B$ (蓝色, 用左端点 Riemann 和计算) 围绕零上下波动, 均值恒为零; 两者之差恰为确定性的修正项 $t$ (红色虚线): 它来自布朗运动的二次变差}
\label{fig:ito-square}
\end{figure}

\begin{intuition}
It\^o 公式表明, $f(t,X_t)$ 仍是 It\^o 过程, 它的扩散系数是 ``链式法则'' 给出的 $\sigma\,\partial_xf$, 而漂移系数除了链式法则给出的 $\partial_tf+b\,\partial_xf$ 之外, 多出一项 $\frac{1}{2}\sigma^2\partial_{xx}f$. 这一项来自 Jensen 式的凸性效应. 若 $f$ 是凸的 ($f''>0$), 则在对称的随机扰动下, $f$ 的平均值会上升. 例如 $f(x)=x^2$: $B_t$ 的均值不变, 但 $B_t^2$ 的均值以速率 $1$ 增长 (图~\ref{fig:ito-square}).
\end{intuition}

% ---------------------------------------------------------------------
\subsection{多维 It\^o 公式}

设 $\rvec{B}=(B^1,\dots,B^d)^\T$ 是 $d$ 维 $\{\mathcal{F}_t\}$-布朗运动: 它适应、轨道连续、$\rvec{B}_0=\rvec{0}$, 且对 $s<t$, $\rvec{B}_t-\rvec{B}_s$ 与 $\mathcal{F}_s$ 独立, 服从 $N(\rvec{0},(t-s)\mat{I}_d)$. 每个分量 $B^k$ 都是 $\{\mathcal{F}_t\}$-布朗运动, 所以对 $H\in\mathscr{L}^2_{\mathrm{loc}}$, 积分 $\int H\dd B^k$ 已有定义.

\begin{proposition}[独立分量的交互变差]\label{prop:cross-bm}
设 $H,K\in\mathscr{L}^2_T$, $k\ne l$. 则 $(H\cdot B^k)(K\cdot B^l)$ 是鞅. 若 $H,K\in\mathscr{L}^2_{\mathrm{loc}}$, 则它是局部鞅.
\end{proposition}

\begin{proof}
先设 $H,K\in\mathscr{S}_T$, 不妨写在同一分割上: $H=\sum_i\xi_{i-1}\1_{(t_{i-1},t_i]}$, $K=\sum_i\eta_{i-1}\1_{(t_{i-1},t_i]}$ (时刻 $0$ 处的值不影响积分). 与命题~\ref{prop:simple-integral} (2) 的证明相同, 只需验证: 对 $i<j$, $\E[\xi_{i-1}\Delta B^k_i\,\eta_{j-1}\Delta B^l_j\,|\,\mathcal{F}_{t_{j-1}}]=0$, 这由 $\E[\Delta B^l_j\,|\,\mathcal{F}_{t_{j-1}}]=0$ 得到; 对 $i=j$,
\[
\E\big[\xi_{i-1}\eta_{i-1}\Delta B^k_i\Delta B^l_i\,\big|\,\mathcal{F}_{t_{i-1}}\big]=\xi_{i-1}\eta_{i-1}\E\big[\Delta B^k_i\Delta B^l_i\big]=0,
\]
因为 $(\Delta B^k_i,\Delta B^l_i)$ 与 $\mathcal{F}_{t_{i-1}}$ 独立, 且两个分量独立、均值为零. 于是 $\E[(H\cdot B^k)_t(K\cdot B^l)_t-(H\cdot B^k)_s(K\cdot B^l)_s\,|\,\mathcal{F}_s]=0$. 一般情形由 $L^2$ 逼近 (与定理~\ref{thm:ito-properties} (3) 的证明相同) 与局部化得到.
\end{proof}

设 $\rvec{X}=(X^1,\dots,X^n)^\T$, 其中
\begin{equation}\label{eq:ito-process-multi}
X^i_t=X^i_0+\int_0^tb^i_s\dd s+\sum_{k=1}^d\int_0^t\sigma^{ik}_s\dd B^k_s,\qquad i=1,\dots,n,
\end{equation}
$b^i$ 与 $\sigma^{ik}$ 满足定义~\ref{def:ito-process} 中的条件. 写成向量形式 $\dd\rvec{X}_t=\rvec{b}_t\dd t+\mat{\Sigma}_t\dd\rvec{B}_t$, 其中 $\mat{\Sigma}_t=(\sigma^{ik}_t)$ 是 $n\times d$ 矩阵值过程. 由命题~\ref{prop:cross-bm}, $X^i$ 的鞅部分 $M^i=\sum_k\int\sigma^{ik}\dd B^k$ 满足 $M^iM^j-\int_0^\cdot\sum_k\sigma^{ik}\sigma^{jk}\dd s$ 是局部鞅. 与命题~\ref{prop:ito-process} (2) 相同, 对 $[0,t]$ 的分割, $\sum\Delta X^i\Delta X^j$ 依概率收敛到
\[
\ip{X^i,X^j}_t=\int_0^t\sum_{k=1}^d\sigma^{ik}_s\sigma^{jk}_s\dd s=\int_0^t\big(\mat{\Sigma}_s\mat{\Sigma}_s^\T\big)_{ij}\dd s.
\]
乘法表 \eqref{eq:mult-table} 相应地推广为 $\dd B^k_t\cdot\dd B^l_t=\delta_{kl}\dd t$: 独立的布朗运动之间没有交互变差.

\begin{theorem}[多维 It\^o 公式]\label{thm:ito-multi}
设 $\rvec{X}$ 如 \eqref{eq:ito-process-multi}, $f:[0,T]\times\R^n\to\R$ 关于 $t$ 一阶、关于 $\rvec{x}$ 二阶连续可微. 则几乎必然地, 对一切 $t\le T$,
\[
f(t,\rvec{X}_t)=f(0,\rvec{X}_0)+\int_0^t\partial_tf\dd s+\sum_{i=1}^n\int_0^t\partial_if\dd X^i_s+\frac{1}{2}\sum_{i,j=1}^n\int_0^t\partial_{ij}f\dd\ip{X^i,X^j}_s,
\]
其中偏导数都在 $(s,\rvec{X}_s)$ 处取值, $\partial_i=\partial/\partial x_i$.
\end{theorem}

\begin{proof}
与定理~\ref{thm:ito-formula} 的证明完全相同, 只有两处改动: Taylor 公式换成多元形式, 二阶项为 $\frac{1}{2}\sum_{i,j}\partial_{ij}f\,\Delta X^i\Delta X^j$; 引理~\ref{lem:weighted-qv} 换成交互变差的版本 $\sum g_{t_{l-1}}\Delta X^i_l\Delta X^j_l\to\int g\dd\ip{X^i,X^j}$, 它由引理~\ref{lem:weighted-qv} 应用于 It\^o 过程 $X^i\pm X^j$ 并作极化得到 (引理~\ref{lem:weighted-qv} 的证明对由多维布朗运动驱动的 It\^o 过程逐字成立).
\end{proof}

\begin{corollary}[乘积法则]\label{cor:product-rule}
设 $X$, $Y$ 是 It\^o 过程. 则
\[
X_tY_t=X_0Y_0+\int_0^tX_s\dd Y_s+\int_0^tY_s\dd X_s+\ip{X,Y}_t,
\]
即 $\dd(XY)=X\dd Y+Y\dd X+\dd\ip{X,Y}$.
特别地, 若 $Y$ 是有限变差的 ($Y$ 的扩散系数为零), 则 $\dd(XY)=X\dd Y+Y\dd X$, 与普通的分部积分相同.
\end{corollary}

\begin{proof}
对 $f(x,y)=xy$ 应用定理~\ref{thm:ito-multi}: $\partial_xf=y$, $\partial_yf=x$, $\partial_{xy}f=\partial_{yx}f=1$, $\partial_{xx}f=\partial_{yy}f=0$.
\end{proof}

\begin{remark}[Stratonovich 积分]
对 It\^o 过程 $X$, $Y$, 定义 Stratonovich 积分 $\int_0^tY_s\circ\dd X_s=\int_0^tY_s\dd X_s+\frac{1}{2}\ip{Y,X}_t$; 可以证明它是中点 Riemann 和 $\sum\frac{1}{2}(Y_{t_{i-1}}+Y_{t_i})\Delta X_i$ 的极限 (与命题~\ref{prop:three-sums} 比较). 若 $f\in C^3$, 对 $Y=f'(X)$ 应用 It\^o 公式得 $\ip{f'(X),X}_t=\int_0^tf''(X_s)\dd\ip{X}_s$, 于是 It\^o 公式改写为
\[
f(X_t)=f(X_0)+\int_0^tf'(X_s)\circ\dd X_s,
\]
这正是普通的链式法则. 所以 Stratonovich 积分适合在流形上做几何 (坐标变换的规则与普通微积分相同), 也是光滑噪声逼近白噪声时的自然极限 (Wong-Zakai 定理); It\^o 积分则保持鞅性质, 适合计算期望. 二者可以按上式互相转换.
\end{remark}

% ---------------------------------------------------------------------
\subsection{应用举例}

\begin{example}[多项式与矩]\label{ex:ito-moments}
对 $f(x)=x^n$ ($n\ge2$), It\^o 公式给出
\[
B_t^n=n\int_0^tB_s^{n-1}\dd B_s+\frac{n(n-1)}{2}\int_0^tB_s^{n-2}\dd s.
\]
$n=2$ 时就是 \eqref{eq:int-BdB}. 由于 $\E\int_0^tB_s^{2n-2}\dd s<\infty$, 被积过程 $B^{n-1}\in\mathscr{L}^2_T$, 随机积分的期望为零 (定理~\ref{thm:ito-properties} (2)). 由 Fubini 定理,
\[
\E B_t^n=\frac{n(n-1)}{2}\int_0^t\E B_s^{n-2}\dd s.
\]
由 $\E B_t^0=1$, $\E B_t=0$ 归纳得 $\E B_t^{2k+1}=0$, $\E B_t^{2k}=\frac{(2k)(2k-1)}{2}\int_0^t(2k-3)!!\,s^{k-1}\dd s=(2k-1)!!\,t^k$, 与例~\ref{ex:normal-mgf} 一致. 这里没有用到正态密度, 只用了 It\^o 公式.
\end{example}

\begin{proposition}[随机指数]\label{prop:stochastic-exp}
设 $\theta\in\mathscr{L}^2_{\mathrm{loc}}$, 令
\[
Z_t=\exp\Big(\int_0^t\theta_s\dd B_s-\frac{1}{2}\int_0^t\theta_s^2\dd s\Big).
\]
则 $Z$ 满足 $\dd Z_t=\theta_tZ_t\dd B_t$, 即 $Z_t=1+\int_0^tZ_s\theta_s\dd B_s$. 因此 $Z$ 是非负连续局部鞅, 从而是上鞅, $\E Z_t\le1$; $Z$ 在 $[0,T]$ 上是鞅当且仅当 $\E Z_T=1$.
\end{proposition}

\begin{proof}
$X_t=\int_0^t\theta\dd B-\frac{1}{2}\int_0^t\theta^2\dd s$ 是 It\^o 过程, $\dd\ip{X}_t=\theta_t^2\dd t$. 对 $f(x)=e^x$ 应用 It\^o 公式,
\[
\dd Z_t=Z_t\dd X_t+\frac{1}{2}Z_t\dd\ip{X}_t=Z_t\theta_t\dd B_t-\frac{1}{2}Z_t\theta_t^2\dd t+\frac{1}{2}Z_t\theta_t^2\dd t=Z_t\theta_t\dd B_t.
\]
$Z\theta$ 轨道上平方可积, 所以 $Z$ 是连续局部鞅 (定理~\ref{thm:ito-local}). 非负连续局部鞅是上鞅 (命题~\ref{prop:local-martingale} (4)), 所以 $\E Z_t\le\E Z_0=1$. 若 $\E Z_T=1$, 则对 $t\le T$, $1\ge\E Z_t\ge\E Z_T=1$, 而 $Z_t-\E[Z_T\,|\,\mathcal{F}_t]\ge0$ 的期望为零, 所以 $Z_t=\E[Z_T\,|\,\mathcal{F}_t]$ $\as$, 即 $Z$ 是鞅; 反之显然.
\end{proof}

$Z$ 记作 $\mathscr{E}(\theta\cdot B)$, 称为\textbf{随机指数} (Dol\'eans-Dade 指数): 它是线性方程 $\dd Z=Z\,\theta\dd B$, $Z_0=1$ 的解, 是普通指数函数 ($\dd z=z\,\dd x$ 的解) 的随机类比, 修正项 $-\frac{1}{2}\int\theta^2$ 来自 It\^o 公式. 当 $\theta$ 是常数时, 它就是第~4~章的指数鞅. 一般情况下 $Z$ 可能只是严格的上鞅 ($\E Z_T<1$), 5.4 节将给出保证 $\E Z_T=1$ 的条件.

\begin{example}[几何布朗运动与 Ornstein-Uhlenbeck 过程]\label{ex:gbm-ou}
(1) 设 $S_t=s_0\exp\big(\sigma B_t+(\mu-\frac{1}{2}\sigma^2)t\big)$ (4.5.5 小节). 对 $f(t,x)=s_0\exp\big(\sigma x+(\mu-\frac{1}{2}\sigma^2)t\big)$, $\partial_tf=(\mu-\frac{1}{2}\sigma^2)f$, $\partial_xf=\sigma f$, $\partial_{xx}f=\sigma^2f$. 由 It\^o 公式,
\[
\dd S_t=\Big(\mu-\frac{1}{2}\sigma^2\Big)S_t\dd t+\sigma S_t\dd B_t+\frac{1}{2}\sigma^2S_t\dd t=\mu S_t\dd t+\sigma S_t\dd B_t.
\]
这就是 Black-Scholes 模型的随机微分方程: 股价的瞬时收益率 $\frac{\dd S_t}{S_t}$ 等于 $\mu\dd t$ 加上噪声 $\sigma\dd B_t$. 第~4~章所说的 ``波动率拖累'' 在这里看得很清楚: $\dd\log S_t=(\mu-\frac{1}{2}\sigma^2)\dd t+\sigma\dd B_t$, 对数的漂移比收益率的漂移少 $\frac{1}{2}\sigma^2$.

(2) 设 $\theta,\sigma>0$, $U_0$ 是 $\mathcal{F}_0$-可测的. 令 $Y_t=U_0+\sigma\int_0^te^{\theta s}\dd B_s$, $U_t=e^{-\theta t}Y_t$. 由乘积法则 ($e^{-\theta t}$ 有限变差, $\dd Y_t=\sigma e^{\theta t}\dd B_t$),
\[
\dd U_t=Y_t\dd\big(e^{-\theta t}\big)+e^{-\theta t}\dd Y_t=-\theta U_t\dd t+\sigma\dd B_t.
\]
所以 $U$ 满足 Ornstein-Uhlenbeck 方程 $\dd U_t=-\theta U_t\dd t+\sigma\dd B_t$: 漂移 $-\theta U_t$ 把过程拉回零点. 若 $U_0=u$ 是常数, 由 Wiener 积分 (命题~\ref{prop:wiener-integral}), $U_t\sim N\big(ue^{-\theta t},\frac{\sigma^2}{2\theta}(1-e^{-2\theta t})\big)$; 当 $t\to\infty$ 时趋于平稳分布 $N(0,\frac{\sigma^2}{2\theta})$. 取 $\theta=1$, $\sigma=\sqrt{2}$, 就得到例~\ref{ex:ou} 中的过程.
\end{example}

\begin{proposition}[时空调和函数与鞅]\label{prop:space-time-harmonic}
设 $h\in C^{1,2}$ 满足向后热方程
\[
\partial_th+\frac{1}{2}\partial_{xx}h=0.
\]
则 $h(t,B_t)=h(0,0)+\int_0^t\partial_xh(s,B_s)\dd B_s$ 是连续局部鞅; 若 $\E\int_0^T\big(\partial_xh(s,B_s)\big)^2\dd s<\infty$, 它是鞅.
\end{proposition}

\begin{proof}
由 It\^o 公式, $\dd h(t,B_t)=(\partial_th+\frac{1}{2}\partial_{xx}h)\dd t+\partial_xh\dd B_t=\partial_xh\dd B_t$. 后一结论由定理~\ref{thm:ito-properties} (2) 得到.
\end{proof}

这就严格化了第~4~章的直观: $B_t$, $B_t^2-t$, $e^{\theta B_t-\theta^2t/2}$ 以及全部 Hermite 鞅 $H_n(B_t,t)$ 都由满足向后热方程的 $h$ 给出. 反过来, 第~6~章将从这个命题出发, 用鞅表示偏微分方程的解 (Feynman-Kac 公式).

\begin{example}[多维布朗运动的常返与暂留]\label{ex:bm-annulus}
设 $d\ge2$, $\rvec{x}\in\R^d$, $0<r<\abs{\rvec{x}}<R$, $\rvec{W}=\rvec{x}+\rvec{B}$ 是从 $\rvec{x}$ 出发的 $d$ 维布朗运动. 令
\[
\tau_r=\inf\{t:\abs{\rvec{W}_t}=r\},\qquad\tau_R=\inf\{t:\abs{\rvec{W}_t}=R\},\qquad\tau=\tau_r\wedge\tau_R.
\]
我们计算 $p=\P(\tau_r<\tau_R)$, 即在走出大球之前先进入小球的概率 (图~\ref{fig:annulus}). 关键是找一个在环形区域上满足 $\Delta f=\sum_i\partial_{ii}f=0$ 的径向函数 (调和函数):
\[
f(\rvec{y})=\begin{cases}\log\abs{\rvec{y}},&d=2,\\ \abs{\rvec{y}}^{2-d},&d\ge3.\end{cases}
\]
直接计算可验证 $\Delta f=0$ 在 $\R^d\setminus\{\rvec{0}\}$ 上成立. 把 $f$ 在 $\{\frac{r}{2}<\abs{\rvec{y}}<2R\}$ 之外修改为 $\R^d$ 上的 $C^2$ 函数 $\tilde{f}$ (不影响下面的论证, 因为停止过程不离开闭环形区域). 由多维 It\^o 公式 ($\dd\ip{W^i,W^j}_t=\delta_{ij}\dd t$),
\begin{align*}
f(\rvec{W}_{t\wedge\tau})&=f(\rvec{x})+\sum_{i=1}^d\int_0^{t\wedge\tau}\partial_if(\rvec{W}_s)\dd B^i_s+\frac{1}{2}\int_0^{t\wedge\tau}\Delta f(\rvec{W}_s)\dd s\\
&=f(\rvec{x})+\sum_{i=1}^d\int_0^t\1_{\{s\le\tau\}}\partial_if(\rvec{W}_s)\dd B^i_s.
\end{align*}
被积过程有界, 属于 $\mathscr{L}^2_T$, 所以 $\E f(\rvec{W}_{t\wedge\tau})=f(\rvec{x})$. 与定理~\ref{thm:exit} (1) 相同可证 $\tau<\infty$ $\as$ (考察第一个分量 $W^1$ 离开 $(-R,R)$ 的时刻即可). 令 $t\to\infty$, 由有界收敛定理, $f(\rvec{x})=p\,f(r)+(1-p)\,f(R)$ (这里 $f(r)$ 指 $f$ 在半径为 $r$ 的球面上的值). 于是
\[
p=\frac{f(R)-f(\abs{\rvec{x}})}{f(R)-f(r)}=\begin{cases}\dfrac{\log R-\log\abs{\rvec{x}}}{\log R-\log r},&d=2,\\[2ex]\dfrac{R^{2-d}-\abs{\rvec{x}}^{2-d}}{R^{2-d}-r^{2-d}},&d\ge3.\end{cases}
\]
\begin{itemize}
\item 令 $R\to\infty$: $d=2$ 时 $p\to1$, 平面布朗运动必然进入任何小圆盘 (\textbf{邻域常返}); $d\ge3$ 时 $p\to(r/\abs{\rvec{x}})^{d-2}<1$, 布朗运动有正概率永不进入小球 (\textbf{暂留}). 可以进一步证明, $d\ge3$ 时 $\abs{\rvec{W}_t}\to\infty$ $\as$ (见~\cite{morters} 第 3 章).
\item 令 $r\to0$ ($R$ 固定): 对一切 $d\ge2$, $p\to0$. 所以在走出大球之前击中原点的概率为零; 再令 $R\to\infty$, 得到 $d\ge2$ 时布朗运动 (从 $\rvec{x}\ne\rvec{0}$ 出发) 几乎必然永不击中原点: 单点是 ``极集'' (polar).
\end{itemize}
这就是 P\'olya 关于随机游动的著名结论的连续版本: ``醉汉总能回家, 醉鸟却可能永远迷路.''
\end{example}

\begin{figure}[htbp]
\centering
\begin{tikzpicture}[scale=2.6]
  \fill[SABlue!6] (0,0) circle (1);
  \fill[white] (0,0) circle (0.25);
  \draw[SADark!70,thick] (0,0) circle (1);
  \draw[SARed,thick] (0,0) circle (0.25);
  \fill[SADark] (0,0) circle (0.6pt) node[below right,font=\scriptsize] {$\rvec{0}$};
  \draw[SABlue,very thin] plot coordinates {(0.550,0.000) (0.545,0.057) (0.575,0.056) (0.531,0.028) (0.571,0.019) (0.592,0.013) (0.600,-0.044) (0.635,-0.042) (0.613,0.006) (0.596,0.031) (0.617,0.055) (0.577,0.070) (0.523,0.074) (0.556,0.113) (0.571,0.083) (0.547,0.101) (0.572,0.125) (0.538,0.158) (0.564,0.122) (0.521,0.096) (0.505,0.080) (0.515,0.061) (0.492,0.127) (0.517,0.110) (0.542,0.082) (0.572,0.075) (0.604,0.143) (0.601,0.144) (0.595,0.107) (0.566,0.095) (0.573,0.122) (0.543,0.139) (0.556,0.095) (0.599,0.116) (0.579,0.112) (0.574,0.133) (0.568,0.142) (0.546,0.108) (0.565,0.122) (0.543,0.106) (0.560,0.080) (0.574,0.114) (0.584,0.145) (0.574,0.129) (0.551,0.106) (0.537,0.125) (0.528,0.158) (0.462,0.180) (0.452,0.161) (0.480,0.144) (0.456,0.161) (0.459,0.128) (0.464,0.129) (0.475,0.131) (0.477,0.135) (0.483,0.118) (0.464,0.110) (0.471,0.087) (0.467,0.051) (0.515,0.096) (0.547,0.104) (0.520,0.126) (0.508,0.154) (0.496,0.136) (0.515,0.123) (0.532,0.128) (0.542,0.091) (0.530,0.103) (0.540,0.067) (0.538,0.075) (0.529,0.058) (0.521,0.077) (0.508,0.058) (0.493,0.062) (0.462,0.056) (0.502,0.023) (0.501,0.030) (0.549,0.048) (0.566,0.054) (0.618,0.077) (0.592,0.067) (0.622,0.083) (0.607,0.071) (0.628,0.048) (0.624,0.060) (0.602,0.055) (0.627,0.029) (0.611,0.023) (0.629,-0.000) (0.612,0.035) (0.690,-0.008) (0.706,0.003) (0.711,0.039) (0.710,0.032) (0.711,0.047) (0.705,0.046) (0.739,0.052) (0.756,0.061) (0.748,0.045) (0.732,0.069) (0.729,0.068) (0.727,0.123) (0.669,0.113) (0.696,0.114) (0.681,0.153) (0.663,0.172) (0.708,0.212) (0.739,0.255) (0.733,0.232) (0.689,0.200) (0.696,0.208) (0.683,0.234) (0.669,0.260) (0.645,0.238) (0.651,0.233) (0.634,0.254) (0.644,0.272) (0.675,0.281) (0.671,0.228) (0.676,0.180) (0.694,0.153) (0.664,0.132) (0.702,0.126) (0.669,0.130) (0.677,0.143) (0.695,0.094) (0.652,0.106) (0.649,0.096) (0.634,0.099) (0.609,0.104) (0.561,0.123) (0.561,0.129) (0.544,0.114) (0.558,0.139) (0.559,0.161) (0.557,0.181) (0.539,0.190) (0.560,0.192) (0.538,0.174) (0.534,0.195) (0.539,0.200) (0.482,0.219) (0.439,0.215) (0.446,0.187) (0.472,0.193) (0.480,0.211) (0.485,0.207) (0.441,0.216) (0.452,0.231) (0.428,0.241) (0.406,0.244) (0.414,0.288) (0.392,0.319) (0.376,0.328) (0.378,0.341) (0.359,0.343) (0.388,0.317) (0.394,0.365) (0.403,0.344) (0.386,0.341) (0.434,0.353) (0.476,0.377) (0.527,0.372) (0.566,0.380) (0.559,0.379) (0.557,0.406) (0.575,0.425) (0.570,0.448) (0.586,0.441) (0.614,0.416) (0.590,0.434) (0.586,0.453) (0.579,0.405) (0.568,0.402) (0.522,0.445) (0.502,0.452) (0.484,0.434) (0.508,0.462) (0.511,0.448) (0.528,0.472) (0.542,0.474) (0.543,0.507) (0.532,0.530) (0.549,0.549) (0.533,0.538) (0.545,0.547) (0.529,0.549) (0.547,0.553) (0.556,0.570) (0.546,0.583) (0.548,0.578) (0.535,0.597) (0.515,0.598) (0.490,0.625) (0.470,0.623) (0.529,0.627) (0.537,0.642) (0.514,0.622) (0.521,0.625) (0.534,0.643) (0.577,0.632) (0.569,0.634) (0.551,0.648) (0.541,0.639) (0.555,0.603) (0.575,0.602) (0.595,0.582) (0.581,0.602) (0.614,0.632) (0.594,0.598) (0.569,0.575) (0.505,0.608) (0.576,0.616) (0.587,0.627) (0.586,0.602) (0.612,0.576) (0.575,0.586) (0.610,0.648) (0.633,0.622) (0.632,0.679) (0.617,0.662) (0.603,0.619) (0.555,0.620) (0.571,0.647) (0.546,0.625) (0.542,0.644) (0.556,0.653) (0.509,0.640) (0.498,0.663) (0.509,0.677) (0.540,0.661) (0.553,0.669) (0.515,0.630) (0.495,0.652) (0.533,0.627) (0.512,0.578) (0.505,0.589) (0.492,0.587) (0.487,0.539) (0.485,0.566) (0.498,0.538) (0.498,0.529) (0.506,0.498) (0.495,0.470) (0.485,0.475) (0.465,0.503) (0.485,0.496) (0.478,0.475) (0.427,0.457) (0.460,0.442) (0.441,0.440) (0.414,0.428) (0.442,0.440) (0.404,0.503) (0.423,0.484) (0.432,0.466) (0.466,0.476) (0.490,0.509) (0.502,0.480) (0.554,0.527) (0.588,0.561) (0.595,0.565) (0.592,0.572) (0.597,0.532) (0.600,0.536) (0.597,0.525) (0.590,0.533) (0.613,0.549) (0.584,0.549) (0.566,0.559) (0.577,0.547) (0.545,0.550) (0.551,0.523) (0.552,0.519) (0.536,0.518) (0.517,0.532) (0.503,0.550) (0.520,0.530) (0.490,0.482) (0.475,0.484) (0.464,0.474) (0.443,0.469) (0.457,0.494) (0.464,0.487) (0.493,0.485) (0.501,0.527) (0.489,0.525) (0.526,0.533) (0.501,0.515) (0.490,0.562) (0.475,0.592) (0.518,0.530) (0.485,0.521) (0.520,0.519) (0.547,0.551) (0.565,0.556) (0.545,0.573) (0.511,0.588) (0.547,0.600) (0.562,0.612) (0.596,0.575) (0.598,0.576) (0.606,0.554) (0.615,0.561) (0.638,0.532) (0.661,0.514) (0.666,0.494) (0.652,0.486) (0.655,0.448) (0.657,0.450) (0.660,0.428) (0.656,0.445) (0.675,0.448) (0.672,0.409) (0.665,0.365) (0.674,0.347) (0.665,0.362) (0.648,0.344) (0.655,0.336) (0.696,0.352) (0.668,0.318) (0.660,0.292) (0.656,0.308) (0.653,0.304) (0.653,0.309) (0.642,0.280) (0.643,0.283) (0.692,0.318) (0.707,0.325) (0.741,0.313) (0.756,0.310) (0.742,0.324) (0.744,0.328) (0.738,0.342) (0.723,0.337) (0.717,0.310) (0.678,0.330) (0.676,0.319) (0.653,0.314) (0.638,0.303) (0.637,0.312) (0.633,0.285) (0.605,0.350) (0.641,0.364) (0.610,0.363) (0.613,0.378) (0.650,0.352) (0.645,0.315) (0.649,0.304) (0.639,0.310) (0.610,0.357) (0.617,0.313) (0.623,0.269) (0.595,0.253) (0.621,0.239) (0.654,0.253) (0.637,0.299) (0.608,0.284) (0.570,0.250) (0.573,0.233) (0.561,0.263) (0.540,0.266) (0.511,0.256) (0.516,0.280) (0.550,0.272) (0.613,0.312) (0.634,0.329) (0.630,0.350) (0.608,0.385) (0.617,0.397) (0.579,0.386) (0.599,0.382) (0.618,0.414) (0.561,0.441) (0.581,0.402) (0.565,0.415) (0.538,0.430) (0.559,0.392) (0.580,0.409) (0.623,0.434) (0.644,0.372) (0.650,0.353) (0.616,0.316) (0.630,0.319) (0.591,0.337) (0.555,0.317) (0.599,0.283) (0.627,0.286) (0.628,0.315) (0.631,0.281) (0.604,0.260) (0.613,0.255) (0.661,0.236) (0.664,0.210) (0.667,0.179) (0.676,0.173) (0.688,0.196) (0.659,0.239) (0.679,0.209) (0.671,0.222) (0.713,0.230) (0.718,0.203) (0.714,0.156) (0.677,0.125) (0.673,0.120) (0.688,0.103) (0.699,0.115) (0.674,0.134) (0.636,0.115) (0.652,0.143) (0.673,0.127) (0.660,0.109) (0.657,0.111) (0.610,0.105) (0.620,0.100) (0.583,0.070) (0.594,0.078) (0.554,0.041) (0.521,0.053) (0.552,0.046) (0.597,0.050) (0.597,0.023) (0.565,0.044) (0.554,0.075) (0.575,0.108) (0.558,0.113) (0.566,0.085) (0.543,0.122) (0.550,0.104) (0.526,0.124) (0.543,0.147) (0.483,0.150) (0.446,0.177) (0.450,0.188) (0.403,0.205) (0.411,0.194) (0.412,0.208) (0.388,0.256) (0.399,0.276) (0.402,0.334) (0.400,0.354) (0.431,0.323) (0.456,0.362) (0.401,0.323) (0.408,0.342) (0.410,0.341) (0.391,0.307) (0.415,0.289) (0.420,0.276) (0.434,0.248) (0.412,0.278) (0.437,0.269) (0.463,0.297) (0.477,0.283) (0.451,0.285) (0.464,0.265) (0.466,0.272) (0.450,0.280) (0.433,0.277) (0.405,0.280) (0.357,0.269) (0.355,0.293) (0.360,0.278) (0.371,0.303) (0.373,0.341) (0.369,0.335) (0.361,0.316) (0.344,0.320) (0.336,0.320) (0.353,0.249) (0.375,0.244) (0.362,0.250) (0.362,0.232) (0.413,0.240) (0.388,0.251) (0.408,0.272) (0.428,0.277) (0.435,0.292) (0.442,0.322) (0.395,0.284) (0.407,0.311) (0.399,0.289) (0.373,0.299) (0.394,0.331) (0.400,0.288) (0.405,0.269) (0.415,0.264) (0.442,0.226) (0.433,0.257) (0.413,0.262) (0.391,0.240) (0.391,0.226) (0.375,0.202) (0.402,0.203) (0.393,0.212) (0.423,0.205) (0.385,0.223) (0.378,0.233) (0.412,0.244) (0.366,0.203) (0.301,0.200) (0.294,0.178) (0.297,0.173) (0.311,0.189) (0.297,0.166) (0.284,0.161) (0.235,0.138) (0.222,0.139) (0.207,0.140)};
  \fill[SAGreenText] (0.55,0) circle (0.9pt) node[above right,font=\small] {$\rvec{x}$};
  \fill[SARed] (0.207,0.140) circle (0.9pt);
  \draw[SADark!60,->] (0,0) -- (-0.6,0.8) node[midway,left,font=\small] {$R$};
  \draw[SARed,->] (0,0) -- (-0.177,-0.177) node[midway,below left=-1pt,font=\scriptsize] {$r$};
\end{tikzpicture}
\caption{平面布朗运动从 $\rvec{x}$ 出发, 在走出半径为 $R$ 的大圆之前先进入半径为 $r$ 的小圆 (红点为进入点). 这一事件的概率为 $\frac{\log R-\log\abs{\rvec{x}}}{\log R-\log r}$, 当 $R\to\infty$ 时趋于 $1$, 当 $r\to0$ 时趋于 $0$}
\label{fig:annulus}
\end{figure}

% =====================================================================
\section{随机微积分进阶专题}
\label{sec:5-4}

本节介绍 It\^o 公式的几个深刻推论. L\'evy 刻画表明布朗运动由 ``连续鞅 + 二次变差为 $t$'' 唯一决定; 改变概率测度可以给布朗运动加上漂移, 这是 Girsanov 定理的内容; 鞅表示定理则断言, 布朗信息流下的一切平方可积鞅都能写成关于 $B$ 的随机积分. 它们是第~6~章与金融数学的基本工具. 最后简要介绍局部时, 它处理 It\^o 公式不能直接应用的非光滑函数 $\abs{x}$.

% ---------------------------------------------------------------------
\subsection{L\'evy 刻画}

\begin{theorem}[L\'evy 刻画]\label{thm:levy-char}
设 $\rvec{B}$ 是 $d$ 维 $\{\mathcal{F}_t\}$-布朗运动, $\sigma^{ik}\in\mathscr{L}^2_{\mathrm{loc}}$ ($i=1,\dots,n$, $k=1,\dots,d$), 令
\[
M^i_t=\sum_{k=1}^d\int_0^t\sigma^{ik}_s\dd B^k_s,\qquad i=1,\dots,n.
\]
若 $\ip{M^i,M^j}_t=\delta_{ij}t$, 即矩阵 $\mat{\Sigma}_s=(\sigma^{ik}_s)$ 对几乎每个 $(s,\omega)$ 满足 $\mat{\Sigma}_s\mat{\Sigma}_s^\T=\mat{I}_n$, 则 $\rvec{M}=(M^1,\dots,M^n)^\T$ 是 $n$ 维 $\{\mathcal{F}_t\}$-布朗运动.
\end{theorem}

\begin{proof}
固定 $\rvec{u}\in\R^n$. 令 $f(t,\rvec{x})=\exp\big(\mathrm{i}\rvec{u}^\T\rvec{x}+\frac{1}{2}\abs{\rvec{u}}^2t\big)$, 对其实部与虚部分别应用多维 It\^o 公式 (定理~\ref{thm:ito-multi}). 由 $\partial_tf=\frac{1}{2}\abs{\rvec{u}}^2f$, $\partial_if=\mathrm{i}u_if$, $\partial_{ij}f=-u_iu_jf$ 与 $\dd\ip{M^i,M^j}_t=\delta_{ij}\dd t$,
\[
\dd f(t,\rvec{M}_t)=\frac{1}{2}\abs{\rvec{u}}^2f\dd t+\sum_i\mathrm{i}u_if\dd M^i_t-\frac{1}{2}\sum_iu_i^2f\dd t=\mathrm{i}\sum_{i,k}u_if(t,\rvec{M}_t)\sigma^{ik}_t\dd B^k_t.
\]
由于 $\abs{f(t,\rvec{M}_t)}=e^{\abs{\rvec{u}}^2t/2}\le e^{\abs{\rvec{u}}^2T/2}$, 且 $\sum_k(\sigma^{ik})^2=1$, 被积过程 (实部与虚部) 属于 $\mathscr{L}^2_T$, 所以 $f(t,\rvec{M}_t)$ 的实部与虚部都是鞅. 对 $s<t$, 用 $\mathcal{F}_s$-可测的有界因子 $e^{-\mathrm{i}\rvec{u}^\T\rvec{M}_s-\frac{1}{2}\abs{\rvec{u}}^2t}$ 乘鞅等式 $\E[f(t,\rvec{M}_t)\,|\,\mathcal{F}_s]=f(s,\rvec{M}_s)$, 得
\[
\E\Big[e^{\mathrm{i}\rvec{u}^\T(\rvec{M}_t-\rvec{M}_s)}\,\Big|\,\mathcal{F}_s\Big]=e^{-\frac{1}{2}\abs{\rvec{u}}^2(t-s)}.
\]
对 $A\in\mathcal{F}_s$, $\P(A)>0$, 上式表明 $\rvec{M}_t-\rvec{M}_s$ 在 $\P(\cdot\,|\,A)$ 下的特征函数是 $N(\rvec{0},(t-s)\mat{I}_n)$ 的特征函数. 由特征函数的唯一性 (第~1~章定理~\ref{thm:indep-cf} 之前的说明), 对一切 Borel 集 $C$, $\P(\rvec{M}_t-\rvec{M}_s\in C\,|\,A)=N(\rvec{0},(t-s)\mat{I}_n)(C)$. 于是 $\rvec{M}_t-\rvec{M}_s$ 与 $\mathcal{F}_s$ 独立, 服从 $N(\rvec{0},(t-s)\mat{I}_n)$. 又 $\rvec{M}$ 连续、适应、$\rvec{M}_0=\rvec{0}$, 所以是 $n$ 维 $\{\mathcal{F}_t\}$-布朗运动.
\end{proof}

更一般地, L\'evy 刻画对任意连续局部鞅成立: 若 $\rvec{M}$ 是 $n$ 维连续局部鞅, $\rvec{M}_0=\rvec{0}$, 且 $M^iM^j-\delta_{ij}t$ 都是局部鞅, 则 $\rvec{M}$ 是布朗运动 (见~\cite{karatzas} 定理 3.3.16). 证明与上面相同, 只是需要关于一般连续局部鞅的随机积分与 It\^o 公式. 这个结论是第~4~章定理~\ref{thm:bm-martingales} (1)(2) 的逆命题: 在连续鞅中, 布朗运动由二次变差 $\ip{M}_t=t$ 唯一刻画, 正态分布可由此推出.

\begin{example}\label{ex:levy-applications}
(1) 令 $\sgn(x)=1$ ($x>0$), $-1$ ($x\le0$), $W_t=\int_0^t\sgn(B_s)\dd B_s$. 由于 $\sgn(B_s)^2=1$, $W$ 是布朗运动. 注意 $W$ 与 $B$ 并不独立, 事实上 (5.4.4 小节) $\abs{B_t}=W_t+L_t$, 其中 $L$ 是局部时.

(2) 设 $\mat{O}_s$ 是取值于 $d\times d$ 正交矩阵的循序可测过程, 则 $\int_0^t\mat{O}_s\dd\rvec{B}_s$ 是 $d$ 维布朗运动. 这是第~4~章旋转不变性 (命题~\ref{prop:bm-rotation}) 的随机版本: 即使旋转随时间随机变化 (但不偷看未来), 布朗运动仍然是布朗运动.

(3) (Bessel 过程) 设 $d\ge2$, $\rvec{W}=\rvec{x}+\rvec{B}$, $\rvec{x}\ne\rvec{0}$. 由例~\ref{ex:bm-annulus}, $\rvec{W}$ 几乎必然永不击中原点, 而 $\rvec{y}\mapsto\abs{\rvec{y}}$ 在 $\R^d\setminus\{\rvec{0}\}$ 上光滑, $\partial_i\abs{\rvec{y}}=\frac{y_i}{\abs{\rvec{y}}}$, $\Delta\abs{\rvec{y}}=\frac{d-1}{\abs{\rvec{y}}}$. 对 $R_t=\abs{\rvec{W}_t}$ 应用 It\^o 公式 (用停时截断到 $\abs{\rvec{W}}$ 远离零与无穷的时段, 再取极限),
\[
R_t=\abs{\rvec{x}}+\beta_t+\frac{d-1}{2}\int_0^t\frac{\dd s}{R_s},\qquad\beta_t=\sum_{i=1}^d\int_0^t\frac{W^i_s}{R_s}\dd B^i_s.
\]
由于 $\sum_i(W^i_s/R_s)^2=1$, $\beta$ 是一维布朗运动. 所以 $d$ 维布朗运动到原点的距离是一个一维布朗运动加上一个向外的漂移 $\frac{d-1}{2R_t}$, 称为 $d$ 维 \textbf{Bessel 过程}. 维数越高, 向外的漂移越强, 这从另一个角度解释了例~\ref{ex:bm-annulus} 中的暂留性.
\end{example}

\begin{intuition}
二次变差 $\ip{M}_t$ 是鞅 $M$ 的 ``内在时钟'' (第~3~章). L\'evy 刻画说: 内在时钟与真实时间同步的连续鞅就是布朗运动. 更一般地, 任何连续局部鞅都是 ``按照自己的时钟运行的布朗运动'', 这就是下面的定理.
\end{intuition}

\begin{theorem}[Dambis-Dubins-Schwarz]\label{thm:dds}
设 $M$ 是连续局部鞅, $M_0=0$, $\ip{M}_\infty=\infty$ $\as$ 对 $s\ge0$ 令 $T_s=\inf\{t:\ip{M}_t>s\}$. 则 $W_s=M_{T_s}$ 是布朗运动 (关于信息流 $\{\mathcal{F}_{T_s}\}$), 且 $M_t=W_{\ip{M}_t}$ 对一切 $t$ 成立. (证明见~\cite{karatzas} 定理 3.4.6.)
\end{theorem}

例如, 对确定的 $\sigma$, Wiener 积分 $M_t=\int_0^t\sigma(s)\dd B_s$ 与 $W_{\int_0^t\sigma^2\dd s}$ 同分布 (二者都是均值为零、协方差为 $\int_0^{s\wedge t}\sigma^2$ 的高斯过程), 这正是 DDS 定理在高斯情形的特例. DDS 定理把关于连续鞅的许多问题 (例如重对数律、零点集的结构) 化归为布朗运动的相应问题.

与此相关, 连续鞅的极大值与二次变差的各阶矩是可比的.

\begin{theorem}[Burkholder-Davis-Gundy 不等式]\label{thm:bdg}
对每个 $p>0$, 存在只依赖于 $p$ 的常数 $0<c_p\le C_p<\infty$, 使得对任意连续局部鞅 $M$ ($M_0=0$) 与停时 $\tau$,
\[
c_p\,\E\big[\ip{M}_\tau^{p/2}\big]\le\E\Big[\sup_{t\le\tau}\abs{M_t}^p\Big]\le C_p\,\E\big[\ip{M}_\tau^{p/2}\big].
\]
(证明见~\cite{karatzas} 定理 3.3.28.)
\end{theorem}

$p=2$ 的情形可由 Doob 不等式与 It\^o 等距直接得到 ($c_2=1$, $C_2=4$). 对 It\^o 积分 $M=H\cdot B$, BDG 不等式说 $\E\sup_{t\le T}\abs{\int_0^tH\dd B}^p\le C_p\E(\int_0^TH^2\dd s)^{p/2}$; 估计随机微分方程的解的高阶矩时要用到它 (第~6~章 Feynman-Kac 公式之后的注).

% ---------------------------------------------------------------------
\subsection{Girsanov 定理}

先看一个有限维的例子. 设 $X\sim N(0,1)$ (在 $\P$ 下). 定义新的概率测度 $\widetilde{\P}(A)=\E[e^{\mu X-\mu^2/2};A]$. 则在 $\widetilde{\P}$ 下, $X$ 的密度为 $\varphi(x)e^{\mu x-\mu^2/2}=\varphi(x-\mu)$, 即 $X\sim N(\mu,1)$: 用似然比 $e^{\mu X-\mu^2/2}$ 重新加权, 就把正态变量的均值从 $0$ 平移到 $\mu$ (图~\ref{fig:girsanov}). Girsanov 定理是这个事实的过程版本: 用随机指数重新加权, 可以给布朗运动加上 (随机的) 漂移.

\begin{figure}[htbp]
\centering
\begin{tikzpicture}
\begin{axis}[SAaxis, axis x line=bottom, axis y line=left, width=11cm, height=5.4cm, xmin=-3.6, xmax=4.6, ymin=0, ymax=0.48,
  xtick={-3,-2,-1,0,1,2,3,4}, ytick={0,0.2,0.4}, xlabel={$x$}, samples=161, domain=-3.6:4.6,
  legend style={font=\small,draw=none,fill=none,at={(0.5,-0.22)},anchor=north}, legend columns=3, legend cell align=left]
  \addplot[SABlue,thick] {exp(-x^2/2)/sqrt(2*pi)};
  \addplot[SARed,thick] {exp(-(x-1)^2/2)/sqrt(2*pi)};
  \addplot[SADark!60,dashed,domain=-3.6:1.9] {0.1*exp(x-0.5)};
  \legend{{$\P$: $N(0,1)$}, {$\widetilde{\P}$: $N(1,1)$}, {似然比 $\frac{\dd\widetilde{\P}}{\dd\P}=e^{x-\frac12}$ (缩小 $10$ 倍)}}
\end{axis}
\end{tikzpicture}
\caption{测度变换 ($\mu=1$). 似然比 $e^{\mu x-\mu^2/2}$ 给右侧的取值更大的权重, 把 $N(0,1)$ 变为 $N(\mu,1)$. Girsanov 定理把这个 ``指数倾斜'' 应用于整条布朗轨道}
\label{fig:girsanov}
\end{figure}

\begin{theorem}[Girsanov 定理]\label{thm:girsanov}
设 $\theta\in\mathscr{L}^2_{\mathrm{loc}}$, $Z_t=\exp\big(\int_0^t\theta_s\dd B_s-\frac{1}{2}\int_0^t\theta_s^2\dd s\big)$, 并设 $\E Z_T=1$ (等价地, $Z$ 在 $[0,T]$ 上是鞅, 命题~\ref{prop:stochastic-exp}). 在 $\mathcal{F}_T$ 上定义概率测度
\[
\widetilde{\P}(A)=\E\big[Z_T\1_A\big],\qquad A\in\mathcal{F}_T.
\]
则
\[
\widetilde{B}_t=B_t-\int_0^t\theta_s\dd s,\qquad t\in[0,T]
\]
在 $\widetilde{\P}$ 下是 $\{\mathcal{F}_t\}$-布朗运动.
\end{theorem}

\begin{proof}
由于 $Z_T>0$ 且 $\E Z_T=1$, $\widetilde{\P}$ 是与 $\P$ 等价的概率测度 (二者的零概率集相同). 对 $t\le T$ 与有界 $\mathcal{F}_t$-可测的 $\xi$, 由 $Z$ 的鞅性质,
\begin{equation}\label{eq:girsanov-bayes}
\widetilde{\E}[\xi]=\E[Z_T\xi]=\E\big[\E[Z_T\,|\,\mathcal{F}_t]\,\xi\big]=\E[Z_t\xi].
\end{equation}

固定 $u\in\R$, 令 $Y_t=Z_t\exp\big(\mathrm{i}u\widetilde{B}_t+\frac{1}{2}u^2t\big)$. 合并指数, 并利用 $\frac{1}{2}(\theta+\mathrm{i}u)^2=\frac{1}{2}\theta^2+\mathrm{i}u\theta-\frac{1}{2}u^2$,
\[
Y_t=\exp\Big(\int_0^t(\theta_s+\mathrm{i}u)\dd B_s-\frac{1}{2}\int_0^t(\theta_s+\mathrm{i}u)^2\dd s\Big).
\]
与命题~\ref{prop:stochastic-exp} 的计算相同 (对实部与虚部应用二维 It\^o 公式, 或直接把 $e^x$ 的 It\^o 公式用于复值的 It\^o 过程), $\dd Y_t=(\theta_t+\mathrm{i}u)Y_t\dd B_t$, 所以 $Y$ (的实部与虚部) 是连续局部鞅. 下证 $Y$ 是鞅. 取局部化序列 $\rho_n\uparrow T$ ($\P(\rho_n=T)\to1$), 使 $Y_{\cdot\wedge\rho_n}$ 是鞅. 对 $t\le T$, 由可选停时定理 (推论~\ref{cor:ct-ost-bounded}) 应用于鞅 $Z$,
\[
\abs{Y_{t\wedge\rho_n}}=Z_{t\wedge\rho_n}e^{\frac{1}{2}u^2(t\wedge\rho_n)}\le e^{\frac{1}{2}u^2T}\,\E\big[Z_T\,\big|\,\mathcal{F}_{t\wedge\rho_n}\big].
\]
由引理~\ref{lem:ce-family-ui}, $\{Y_{t\wedge\rho_n}\}_{n\ge1}$ 一致可积, 而 $Y_{t\wedge\rho_n}\to Y_t$ $\as$ 由 Vitali 收敛定理~\ref{thm:vitali}, 收敛也在 $L^1$ 中成立, 所以对 $s\le t$, $\E[Y_t\,|\,\mathcal{F}_s]=\lim_n\E[Y_{t\wedge\rho_n}\,|\,\mathcal{F}_s]=\lim_nY_{s\wedge\rho_n}=Y_s$.

现在对 $s<t\le T$ 与 $A\in\mathcal{F}_s$, 由 \eqref{eq:girsanov-bayes} (取 $\xi=e^{\mathrm{i}u(\widetilde{B}_t-\widetilde{B}_s)}\1_A$, 实部虚部分别处理) 与 $Y$ 的鞅性质,
\begin{align*}
\widetilde{\E}\Big[e^{\mathrm{i}u(\widetilde{B}_t-\widetilde{B}_s)};A\Big]&=\E\Big[Z_te^{\mathrm{i}u(\widetilde{B}_t-\widetilde{B}_s)};A\Big]=\E\Big[Y_t\,e^{-\mathrm{i}u\widetilde{B}_s-\frac{1}{2}u^2t};A\Big]\\
&=\E\Big[Y_s\,e^{-\mathrm{i}u\widetilde{B}_s-\frac{1}{2}u^2t};A\Big]=e^{-\frac{1}{2}u^2(t-s)}\,\E[Z_s;A]=e^{-\frac{1}{2}u^2(t-s)}\,\widetilde{\P}(A).
\end{align*}
与定理~\ref{thm:levy-char} 的证明相同, 由此可知在 $\widetilde{\P}$ 下, $\widetilde{B}_t-\widetilde{B}_s$ 与 $\mathcal{F}_s$ 独立, 服从 $N(0,t-s)$. 又 $\widetilde{B}$ 连续、适应、从 $0$ 出发, 所以它是 $\widetilde{\P}$ 下的 $\{\mathcal{F}_t\}$-布朗运动.
\end{proof}

换一种说法: 在 $\widetilde{\P}$ 下, $B_t=\widetilde{B}_t+\int_0^t\theta_s\dd s$ 是 ``布朗运动加漂移 $\theta$''. 测度变换不改变轨道, 只改变各条轨道的权重; 它改变漂移, 但不改变扩散系数 (二次变差 $\ip{B}_t=t$ 是逐条轨道的性质, 在等价测度下不变).

定理的条件 $\E Z_T=1$ 需要验证. 最常用的充分条件如下.

\begin{proposition}\label{prop:novikov}
\begin{enumerate}
\item 若 $\theta$ 有界, 则 $\E Z_T=1$.
\item (\textbf{Novikov 条件}) 若 $\E\exp\big(\frac{1}{2}\int_0^T\theta_s^2\dd s\big)<\infty$, 则 $\E Z_T=1$. (证明见~\cite{karatzas} 推论 3.5.13.)
\end{enumerate}
\end{proposition}

\begin{proof}[(1) 的证明]
设 $\abs{\theta}\le K$. 由命题~\ref{prop:stochastic-exp}, $Z=1+\int Z\theta\dd B$; 只需证明 $Z\theta\in\mathscr{L}^2_T$, 则 $Z$ 是鞅. 记 $Z^{(2\theta)}_s=\exp(\int_0^s2\theta\dd B-\frac{1}{2}\int_0^s4\theta^2\dd u)$, 它也是随机指数, 所以 $\E Z^{(2\theta)}_s\le1$. 而
\[
Z_s^2=\exp\Big(2\int_0^s\theta\dd B-\int_0^s\theta^2\dd u\Big)=Z^{(2\theta)}_s\exp\Big(\int_0^s\theta^2\dd u\Big)\le e^{K^2T}Z^{(2\theta)}_s,
\]
所以 $\E\int_0^TZ_s^2\theta_s^2\dd s\le K^2Te^{K^2T}<\infty$.
\end{proof}

\begin{example}[带漂移布朗运动的首达时密度]\label{ex:girsanov-drift}
取常数 $\theta\equiv\mu$, 则 $Z_t=e^{\mu B_t-\mu^2t/2}$, 由命题~\ref{prop:novikov} (1), Girsanov 定理适用: 在 $\widetilde{\P}$ 下, $B_t=\widetilde{B}_t+\mu t$ 是漂移为 $\mu$ 的布朗运动. 设 $a>0$, $\tau_a=\inf\{t: B_t=a\}$. 事件 $\{\tau_a\le t\}$ 属于 $\mathcal{F}_{\tau_a\wedge t}$ (对 $r\ge t$, $\{\tau_a\le t\}\cap\{\tau_a\wedge t\le r\}=\{\tau_a\le t\}\in\mathcal{F}_t\subseteq\mathcal{F}_r$; 对 $r<t$, 它等于 $\{\tau_a\le r\}\in\mathcal{F}_r$). 由 \eqref{eq:girsanov-bayes}、可选停时定理 ($\E[Z_t\,|\,\mathcal{F}_{\tau_a\wedge t}]=Z_{\tau_a\wedge t}$) 与 $B_{\tau_a}=a$,
\[
\widetilde{\P}(\tau_a\le t)=\E\big[Z_t;\tau_a\le t\big]=\E\big[Z_{\tau_a};\tau_a\le t\big]=\E\Big[e^{\mu a-\frac{1}{2}\mu^2\tau_a};\tau_a\le t\Big].
\]
代入 $\P$ 下 $\tau_a$ 的密度 (定理~\ref{thm:hitting-time}), 得到 $\widetilde{\P}$ 下 $\tau_a$ 的密度
\[
e^{\mu a-\frac{1}{2}\mu^2s}\cdot\frac{a}{\sqrt{2\pi s^3}}e^{-\frac{a^2}{2s}}=\frac{a}{\sqrt{2\pi s^3}}\exp\Big(-\frac{(a-\mu s)^2}{2s}\Big),\qquad s>0.
\]
这就是 \textbf{逆高斯分布}. 对它积分, $\widetilde{\P}(\tau_a<\infty)=e^{\mu a}\E e^{-\frac{1}{2}\mu^2\tau_a}=e^{\mu a-a\abs{\mu}}$ (定理~\ref{thm:hitting-time} (3)), 与定理~\ref{thm:drift-hitting} 一致. 第~4~章用指数鞅与可选停时定理只得到了 Laplace 变换; 这里 Girsanov 定理直接给出了整个分布.
\end{example}

\begin{remark}
(1) Girsanov 定理对 $d$ 维布朗运动同样成立: 若 $\rvec{\theta}$ 取值于 $\R^d$, $Z=\exp(\sum_k\int\theta^k\dd B^k-\frac{1}{2}\int\abs{\rvec{\theta}}^2\dd s)$, $\E Z_T=1$, 则 $\rvec{B}-\int\rvec{\theta}\dd s$ 是 $\widetilde{\P}$ 下的 $d$ 维布朗运动. 证明相同.

(2) 在 Black-Scholes 模型 $\dd S_t=\mu S_t\dd t+\sigma S_t\dd B_t$ 中, 取 $\theta=-\frac{\mu-r}{\sigma}$ ($r$ 为无风险利率, $\frac{\mu-r}{\sigma}$ 称为风险的市场价格), 则在 $\widetilde{\P}$ 下 $\dd S_t=rS_t\dd t+\sigma S_t\dd\widetilde{B}_t$, 贴现价格 $e^{-rt}S_t$ 是鞅. $\widetilde{\P}$ 称为\textbf{风险中性测度}, 衍生品的无套利价格是贴现收益在 $\widetilde{\P}$ 下的期望. 第~6~章将由此推出 Black-Scholes 公式.

(3) $\widetilde{\P}$ 只定义在 $\mathcal{F}_T$ 上. 在无限时间区间上, 即使 $\theta$ 是常数, $\widetilde{\P}$ 与 $\P$ 一般也互相奇异: 在 $\P$ 下 $B_t/t\to0$, 在 $\widetilde{\P}$ 下 $B_t/t\to\mu$.
\end{remark}

\begin{intuition}
Girsanov 定理的离散版本是第~3~章的似然比鞅 (例~\ref{ex:likelihood-ratio}): 把每一步的增量 $\Delta B$ 的分布从 $N(0,\Delta t)$ 换成 $N(\theta\Delta t,\Delta t)$, 似然比为 $\exp(\theta\Delta B-\frac{1}{2}\theta^2\Delta t)$, 把各步乘起来, 取极限, 就是随机指数 $Z_T$. 由于漂移 $\theta\Delta t$ 远小于噪声的标准差 $\sqrt{\Delta t}$, 单步的似然比接近 $1$, 两个测度在有限时间内仍然等价; 而扩散系数的改变会使每一步的似然比偏离 $1$ 的程度为 $O(1)$, 累积起来使两个测度奇异, 这就是 Girsanov 定理只能改变漂移的原因.
\end{intuition}

% ---------------------------------------------------------------------
\subsection{鞅表示定理}

本小节中 $\{\mathcal{F}_t\}$ 是布朗运动 $B$ 的扩充信息流 (定义~\ref{def:augmented}): 除了 $B$ 本身, 没有其他随机性来源. 此时 $\mathcal{F}_0$ 只含零概率集及其补集, $\mathcal{F}_0$-可测的随机变量几乎必然是常数.

\begin{theorem}[It\^o 鞅表示定理]\label{thm:mart-rep}
设 $F\in L^2(\Omega,\mathcal{F}_T,\P)$. 则存在唯一的 $H\in\mathscr{L}^2_T$, 使得
\[
F=\E F+\int_0^TH_s\dd B_s\qquad\as
\]
\end{theorem}

\begin{proof}
令 $\mathscr{R}$ 为可以如此表示的 $F\in L^2(\mathcal{F}_T)$ 全体.

\emph{唯一性.} 若 $\E F+\int H\dd B=\E F+\int H'\dd B$, 则由 It\^o 等距 $\norm{H-H'}_T^2=\E(\int(H-H')\dd B)^2=0$.

\emph{第一步: $\mathscr{R}$ 是闭线性子空间.} 线性是显然的. 设 $F_n=\E F_n+\int H^n\dd B\in\mathscr{R}$, $F_n\to F$ 于 $L^2$. 则 $\E F_n\to\E F$, 从而 $\int H^n\dd B=F_n-\E F_n$ 是 $L^2$ 中的 Cauchy 列. 由 It\^o 等距, $\{H^n\}$ 是 $\mathscr{L}^2_T$ 中的 Cauchy 列, 由完备性收敛到某个 $H$, 再由等距 $\int H^n\dd B\to\int H\dd B$. 所以 $F=\E F+\int H\dd B\in\mathscr{R}$.

\emph{第二步: 指数泛函属于 $\mathscr{R}$.} 设 $h=\sum_{j=1}^n\lambda_j\1_{(t_{j-1},t_j]}$ 是 $[0,T]$ 上的阶梯函数 ($\lambda_j\in\R$), 令 $\mathscr{E}^h_t=\exp\big(\int_0^th\dd B-\frac{1}{2}\int_0^th^2\dd s\big)$. 由命题~\ref{prop:stochastic-exp}, $\mathscr{E}^h_T=1+\int_0^T\mathscr{E}^h_sh(s)\dd B_s$. 由于 $\int_0^sh\dd B\sim N(0,\int_0^sh^2)$ (命题~\ref{prop:wiener-integral}) 与例~\ref{ex:normal-mgf},
\[
\E\big(\mathscr{E}^h_s\big)^2=e^{-\int_0^sh^2}\,\E e^{2\int_0^sh\dd B}=e^{-\int_0^sh^2}e^{2\int_0^sh^2}=e^{\int_0^sh^2}\le e^{\int_0^Th^2},
\]
所以 $\mathscr{E}^hh\in\mathscr{L}^2_T$, 且 $\E\mathscr{E}^h_T=1$. 于是 $\mathscr{E}^h_T\in\mathscr{R}$.

\emph{第三步: 指数泛函的线性张成在 $L^2(\mathcal{F}_T)$ 中稠密.} 设 $G\in L^2(\mathcal{F}_T)$ 与一切 $\mathscr{E}^h_T$ 正交. 由于 $\mathscr{E}^h_T$ 等于正的常数乘以 $\exp(\sum_j\lambda_j\Delta B_j)$ ($\Delta B_j=B_{t_j}-B_{t_{j-1}}$), 对任意分割 $0=t_0<\cdots<t_n=T$ 与 $\rvec{\lambda}\in\R^n$,
\[
\phi(\rvec{z}):=\E\Big[G\exp\Big(\sum_{j=1}^nz_j\Delta B_j\Big)\Big]
\]
在 $\rvec{z}=\rvec{\lambda}\in\R^n$ 处为零. 对 $\rvec{z}\in\mathbb{C}^n$, 由 Cauchy-Schwarz 不等式, $\E\big[\abs{G}e^{\sum_j\abs{z_j}\abs{\Delta B_j}}\big]<\infty$, 所以 $\phi$ 在 $\mathbb{C}^n$ 上有定义; 由期望号下求导 (定理~\ref{thm:diff-under-E} 对复变量同样成立, 控制函数取 $\abs{G}\abs{\Delta B_j}e^{\sum_j(\abs{z_j}+1)\abs{\Delta B_j}}$), $\phi$ 关于每个变量都是整函数. 固定 $z_2,\dots,z_n\in\R$, $z_1\mapsto\phi$ 是在实轴上为零的整函数, 由解析函数的唯一性定理, 它在 $\mathbb{C}$ 上恒为零; 再固定 $z_1\in\mathbb{C}$ 与 $z_3,\dots,z_n\in\R$, 对 $z_2$ 重复这一论证; 依此类推, $\phi\equiv0$ 于 $\mathbb{C}^n$. 特别地, 对一切 $\rvec{u}\in\R^n$,
\[
\E\big[G^+e^{\mathrm{i}\rvec{u}^\T\Delta\rvec{B}}\big]=\E\big[G^-e^{\mathrm{i}\rvec{u}^\T\Delta\rvec{B}}\big],\qquad\Delta\rvec{B}=(\Delta B_1,\dots,\Delta B_n)^\T.
\]
令 $\mu^\pm(C)=\E[G^\pm;\Delta\rvec{B}\in C]$, 它们是 $\R^n$ 上的有限测度, 由 $\rvec{u}=\rvec{0}$ 知总质量相等, 记为 $c$. 若 $c=0$, 则 $G=0$ $\as$; 若 $c>0$, 则 $\mu^\pm/c$ 是特征函数相同的概率测度, 由唯一性定理它们相等. 总之 $\E[G;\Delta\rvec{B}\in C]=0$ 对一切 Borel 集 $C$ 成立, 即 $\E[G;A]=0$ 对一切 $A\in\sigma(B_{t_1},\dots,B_{t_n})$ 成立. 所有这样的 $A$ (对一切有限分割) 构成生成 $\mathcal{F}^B_T$ 的 $\pi$-系 (两个分割取公共加细即可), 而使 $\E[G;A]=0$ 的 $A$ 构成 $\lambda$-系 (由控制收敛定理), 由 $\pi$-$\lambda$ 定理, $\E[G;A]=0$ 对一切 $A\in\mathcal{F}^B_T$ 成立. 由引理~\ref{lem:null-augment}, $\mathcal{F}_T$ 中的事件与 $\mathcal{F}^B_T$ 中的某个事件只相差零概率集, 所以 $\E[G;A]=0$ 对一切 $A\in\mathcal{F}_T$ 成立. 取 $A=\{G>0\}$ 与 $\{G<0\}$, 得 $G=0$ $\as$

\emph{结论.} 由第二步, 指数泛函的线性张成包含于 $\mathscr{R}$; 由第一步, 其闭包也包含于 $\mathscr{R}$. 由第三步, 这个闭包的正交补为 $\{0\}$, 所以闭包就是 $L^2(\mathcal{F}_T)$ (投影定理~\ref{thm:projection}). 因此 $\mathscr{R}=L^2(\mathcal{F}_T)$.
\end{proof}

\begin{corollary}[布朗信息流下的鞅]\label{cor:mart-rep}
设 $M$ 是关于 $B$ 的扩充信息流的平方可积鞅. 则存在唯一的 $H\in\mathscr{L}^2_T$, 使得对一切 $t\le T$, $M_t=\E M_0+\int_0^tH_s\dd B_s$ $\as$ 特别地, $M$ 有连续修正.
\end{corollary}

\begin{proof}
对 $F=M_T$ 应用定理~\ref{thm:mart-rep}, 并注意 $\E M_T=\E M_0$. 由鞅性与 It\^o 积分的鞅性, $M_t=\E[M_T\,|\,\mathcal{F}_t]=\E M_0+\E[\int_0^TH\dd B\,|\,\mathcal{F}_t]=\E M_0+\int_0^tH\dd B$. 右端是连续过程.
\end{proof}

\begin{example}\label{ex:mart-rep}
表示中的被积过程常常可以用 It\^o 公式求出.
(1) $B_T^2=T+\int_0^T2B_s\dd B_s$ (由 \eqref{eq:int-BdB}).
(2) 由 It\^o 公式与例~\ref{ex:int-sdB}, $B_T^3=3\int_0^TB_s^2\dd B_s+3\int_0^TB_s\dd s=\int_0^T3\big(B_s^2+T-s\big)\dd B_s$.
(3) $N_t=\E[e^{B_T}\,|\,\mathcal{F}_t]=e^{B_t+\frac{1}{2}(T-t)}$ 满足 $\dd N_t=N_t\dd B_t$ (It\^o 公式), 所以 $e^{B_T}=e^{T/2}+\int_0^Te^{B_s+\frac{1}{2}(T-s)}\dd B_s$.
一般地, 若 $F=f(B_T)$, 则 $u(t,x)=\E f(x+B_{T-t})$ 满足向后热方程 (定理~\ref{thm:heat-equation}), 由命题~\ref{prop:space-time-harmonic}, $H_t=\partial_xu(t,B_t)$.
\end{example}

\begin{remark}
(1) 推论~\ref{cor:mart-rep} 对局部鞅也成立: 布朗信息流下的每个局部鞅都形如 $M_0+\int H\dd B$, $H\in\mathscr{L}^2_{\mathrm{loc}}$. 这是一个结构性的结论: 在布朗信息流下, 随机性只有 ``一个来源'', 每个鞅的增量都与 $\dd B$ 成比例. 若信息流中还有与 $B$ 独立的随机性 (例如一个独立的 Poisson 过程), 结论就不成立了.
(2) 被积过程 $H$ 的显式公式由 Malliavin 分析中的 \textbf{Clark-Ocone 公式} 给出: $H_t=\E[D_tF\,|\,\mathcal{F}_t]$, 其中 $D_tF$ 表示 $F$ 对轨道扰动的 ``导数'', 扰动发生在时刻 $t$. 例如 $F=B_T^3$ 时 $D_tF=3B_T^2$, $\E[3B_T^2\,|\,\mathcal{F}_t]=3(B_t^2+T-t)$, 与例~\ref{ex:mart-rep} (2) 一致.
(3) 在金融中, 鞅表示定理保证 Black-Scholes 市场是\textbf{完备}的: 任何 (平方可积的) 衍生品收益 $F$ 都可以由股票与债券的自融资组合完全复制, 复制组合的股票头寸由 $H$ 决定.
\end{remark}

% ---------------------------------------------------------------------
\subsection{局部时与 Tanaka 公式}

It\^o 公式要求 $f\in C^2$. 对 $f(x)=\abs{x}$, 形式上 $f'(x)=\sgn(x)$, $f''(x)=2\delta_0(x)$ ($\delta_0$ 为 Dirac 函数), 若不顾一切地套用 It\^o 公式, 得到
\[
\abs{B_t}=\int_0^t\sgn(B_s)\dd B_s+\int_0^t\delta_0(B_s)\dd s.
\]
最后一项应理解为布朗运动在零点 ``停留的时间密度''. 布朗运动在零点停留的总时间 $\int_0^t\1_{\{B_s=0\}}\dd s$ 为零 (由 Fubini 定理, 其期望为 $\int_0^t\P(B_s=0)\dd s=0$), 但把 $\delta_0$ 换成 $\frac{1}{2\varepsilon}\1_{(-\varepsilon,\varepsilon)}$ 再令 $\varepsilon\to0$, 极限并不为零.

\begin{definition}[局部时]\label{def:local-time}
布朗运动在 $0$ 处的\textbf{局部时}定义为
\[
L_t=\lim_{\varepsilon\downarrow0}\frac{1}{2\varepsilon}\int_0^t\1_{\{\abs{B_s}<\varepsilon\}}\dd s,
\]
可以证明, 极限几乎必然存在, 且关于 $t\in[0,T]$ 一致成立 (见~\cite{karatzas} 3.6 节).
\end{definition}

由定义, $L$ 连续、单调不减、$L_0=0$, 并且只在 $B$ 的零点集 $\{t: B_t=0\}$ 上增长 (若 $B$ 在 $[s,t]$ 上不为零, 则当 $\varepsilon$ 足够小时被积函数在 $[s,t]$ 上为零). 由于零点集的 Lebesgue 测度为零, $L$ 是一个 ``奇异'' 的增函数, 类似于 Cantor 函数.

\begin{theorem}[Tanaka 公式]\label{thm:tanaka}
几乎必然地, 对一切 $t\le T$,
\[
\abs{B_t}=\int_0^t\sgn(B_s)\dd B_s+L_t.
\]
\end{theorem}

\begin{proof}[证明梗概]
对 $\varepsilon>0$ 令 $f_\varepsilon(x)=\abs{x}$ ($\abs{x}\ge\varepsilon$), $f_\varepsilon(x)=\frac{x^2+\varepsilon^2}{2\varepsilon}$ ($\abs{x}<\varepsilon$). 则 $f_\varepsilon\in C^1$, $f_\varepsilon'(x)=\big(\frac{x}{\varepsilon}\wedge1\big)\vee(-1)$, $f''_\varepsilon=\frac{1}{\varepsilon}\1_{(-\varepsilon,\varepsilon)}$ 除两点外存在. It\^o 公式可以推广到这类函数 (用光滑函数逼近 $f_\varepsilon$), 给出
\[
f_\varepsilon(B_t)=f_\varepsilon(0)+\int_0^tf'_\varepsilon(B_s)\dd B_s+\frac{1}{2\varepsilon}\int_0^t\1_{\{\abs{B_s}<\varepsilon\}}\dd s.
\]
令 $\varepsilon\to0$: $f_\varepsilon(x)\to\abs{x}$ 一致成立; 由 It\^o 等距, $\E\big(\int_0^t(f'_\varepsilon(B_s)-\sgn(B_s))\dd B_s\big)^2\le\E\int_0^t\1_{\{\abs{B_s}<\varepsilon\}}\dd s\to0$. 于是最后一项在 $L^2$ 中收敛, 其极限就是 $L_t$, 公式成立.
\end{proof}

\begin{figure}[htbp]
\centering
\begin{tikzpicture}
\begin{axis}[SAaxis, axis x line=bottom, axis y line=left, width=12cm, height=6cm, xmin=0, xmax=1.03, ymin=-1.35, ymax=1.35,
  xtick={0,0.25,0.5,0.75,1}, ytick={-1,-0.5,0,0.5,1}, xlabel={$t$},
  legend style={font=\small,draw=none,fill=none,at={(0.5,-0.2)},anchor=north}, legend columns=3, legend cell align=left]
  \addplot[SADark!20,forget plot] coordinates {(0,0) (1,0)};
  \addplot[SABlue,thin] coordinates {(0.000,0.000) (0.002,0.013) (0.004,0.003) (0.006,0.089) (0.008,0.115) (0.010,0.115) (0.012,0.145) (0.014,0.145) (0.016,0.056) (0.018,0.051) (0.020,0.017) (0.022,0.010) (0.024,0.043) (0.026,0.003) (0.028,0.029) (0.030,0.058) (0.032,0.087) (0.034,0.094) (0.036,0.172) (0.038,0.188) (0.040,0.203) (0.042,0.216) (0.044,0.202) (0.046,0.179) (0.048,0.198) (0.050,0.212) (0.052,0.271) (0.054,0.266) (0.056,0.199) (0.058,0.227) (0.060,0.242) (0.062,0.198) (0.064,0.116) (0.066,0.041) (0.068,0.034) (0.070,0.066) (0.072,0.016) (0.074,0.042) (0.076,0.018) (0.078,0.036) (0.080,0.074) (0.082,0.115) (0.084,0.012) (0.086,0.016) (0.088,0.011) (0.090,0.085) (0.092,0.227) (0.094,0.291) (0.096,0.233) (0.098,0.221) (0.100,0.151) (0.102,0.161) (0.104,0.131) (0.106,0.163) (0.108,0.116) (0.110,0.122) (0.112,0.176) (0.114,0.263) (0.116,0.352) (0.118,0.348) (0.120,0.303) (0.122,0.243) (0.124,0.213) (0.126,0.228) (0.128,0.247) (0.130,0.261) (0.132,0.226) (0.134,0.181) (0.136,0.158) (0.138,0.171) (0.140,0.162) (0.142,0.150) (0.144,0.117) (0.146,0.128) (0.148,0.100) (0.150,0.190) (0.152,0.268) (0.154,0.280) (0.156,0.266) (0.158,0.357) (0.160,0.378) (0.162,0.343) (0.164,0.346) (0.166,0.318) (0.168,0.327) (0.170,0.264) (0.172,0.281) (0.174,0.293) (0.176,0.359) (0.178,0.366) (0.180,0.396) (0.182,0.410) (0.184,0.342) (0.186,0.335) (0.188,0.257) (0.190,0.238) (0.192,0.223) (0.194,0.153) (0.196,0.182) (0.198,0.154) (0.200,0.121) (0.202,0.095) (0.204,0.118) (0.206,0.130) (0.208,0.092) (0.210,0.080) (0.212,0.155) (0.214,0.135) (0.216,0.134) (0.218,0.113) (0.220,0.176) (0.222,0.207) (0.224,0.230) (0.226,0.177) (0.228,0.143) (0.230,0.122) (0.232,0.110) (0.234,0.132) (0.236,0.047) (0.238,0.080) (0.240,0.089) (0.242,0.061) (0.244,0.060) (0.246,0.058) (0.248,0.112) (0.250,0.092) (0.252,0.102) (0.254,0.088) (0.256,0.079) (0.258,0.084) (0.260,0.036) (0.262,0.061) (0.264,0.091) (0.266,0.157) (0.268,0.226) (0.270,0.182) (0.272,0.202) (0.274,0.313) (0.276,0.369) (0.278,0.378) (0.280,0.422) (0.282,0.411) (0.284,0.436) (0.286,0.470) (0.288,0.549) (0.290,0.592) (0.292,0.532) (0.294,0.553) (0.296,0.591) (0.298,0.612) (0.300,0.635) (0.302,0.599) (0.304,0.641) (0.306,0.643) (0.308,0.629) (0.310,0.651) (0.312,0.654) (0.314,0.640) (0.316,0.639) (0.318,0.587) (0.320,0.537) (0.322,0.519) (0.324,0.455) (0.326,0.501) (0.328,0.438) (0.330,0.442) (0.332,0.450) (0.334,0.400) (0.336,0.388) (0.338,0.363) (0.340,0.333) (0.342,0.311) (0.344,0.353) (0.346,0.317) (0.348,0.259) (0.350,0.228) (0.352,0.215) (0.354,0.275) (0.356,0.349) (0.358,0.403) (0.360,0.402) (0.362,0.402) (0.364,0.439) (0.366,0.412) (0.368,0.414) (0.370,0.469) (0.372,0.411) (0.374,0.432) (0.376,0.411) (0.378,0.424) (0.380,0.418) (0.382,0.415) (0.384,0.367) (0.386,0.408) (0.388,0.428) (0.390,0.341) (0.392,0.366) (0.394,0.406) (0.396,0.409) (0.398,0.373) (0.400,0.351) (0.402,0.383) (0.404,0.361) (0.406,0.331) (0.408,0.324) (0.410,0.291) (0.412,0.320) (0.414,0.267) (0.416,0.275) (0.418,0.266) (0.420,0.262) (0.422,0.329) (0.424,0.316) (0.426,0.257) (0.428,0.261) (0.430,0.209) (0.432,0.263) (0.434,0.305) (0.436,0.288) (0.438,0.299) (0.440,0.256) (0.442,0.233) (0.444,0.407) (0.446,0.341) (0.448,0.297) (0.450,0.196) (0.452,0.168) (0.454,0.135) (0.456,0.172) (0.458,0.212) (0.460,0.247) (0.462,0.282) (0.464,0.317) (0.466,0.243) (0.468,0.268) (0.470,0.266) (0.472,0.284) (0.474,0.358) (0.476,0.350) (0.478,0.271) (0.480,0.218) (0.482,0.188) (0.484,0.140) (0.486,0.080) (0.488,0.192) (0.490,0.187) (0.492,0.191) (0.494,0.209) (0.496,0.125) (0.498,0.089) (0.500,0.028) (0.502,0.032) (0.504,0.092) (0.506,0.050) (0.508,0.133) (0.510,0.179) (0.512,0.245) (0.514,0.284) (0.516,0.355) (0.518,0.312) (0.520,0.302) (0.522,0.273) (0.524,0.227) (0.526,0.217) (0.528,0.177) (0.530,0.193) (0.532,0.146) (0.534,0.113) (0.536,0.111) (0.538,0.187) (0.540,0.199) (0.542,0.175) (0.544,0.191) (0.546,0.123) (0.548,0.144) (0.550,0.172) (0.552,0.213) (0.554,0.257) (0.556,0.289) (0.558,0.339) (0.560,0.267) (0.562,0.211) (0.564,0.175) (0.566,0.117) (0.568,0.149) (0.570,0.106) (0.572,0.162) (0.574,0.085) (0.576,0.059) (0.578,0.098) (0.580,0.065) (0.582,0.033) (0.584,0.030) (0.586,0.021) (0.588,0.009) (0.590,0.018) (0.592,0.009) (0.594,0.130) (0.596,0.168) (0.598,0.106) (0.600,0.123) (0.602,0.078) (0.604,0.024) (0.606,0.009) (0.608,0.026) (0.610,0.054) (0.612,0.038) (0.614,0.027) (0.616,0.009) (0.618,0.108) (0.620,0.072) (0.622,0.020) (0.624,0.010) (0.626,0.029) (0.628,0.049) (0.630,0.042) (0.632,0.075) (0.634,0.081) (0.636,0.093) (0.638,0.062) (0.640,0.103) (0.642,0.090) (0.644,0.088) (0.646,0.140) (0.648,0.148) (0.650,0.107) (0.652,0.128) (0.654,0.101) (0.656,0.035) (0.658,0.014) (0.660,0.071) (0.662,0.037) (0.664,0.017) (0.666,0.032) (0.668,0.026) (0.670,0.003) (0.672,0.028) (0.674,0.080) (0.676,0.040) (0.678,0.072) (0.680,0.084) (0.682,0.049) (0.684,0.016) (0.686,0.055) (0.688,0.034) (0.690,0.003) (0.692,0.044) (0.694,0.006) (0.696,0.027) (0.698,0.031) (0.700,0.046) (0.702,0.022) (0.704,0.033) (0.706,0.003) (0.708,0.033) (0.710,0.046) (0.712,0.063) (0.714,0.016) (0.716,0.089) (0.718,0.095) (0.720,0.145) (0.722,0.144) (0.724,0.116) (0.726,0.118) (0.728,0.146) (0.730,0.152) (0.732,0.080) (0.734,0.104) (0.736,0.103) (0.738,0.193) (0.740,0.227) (0.742,0.281) (0.744,0.239) (0.746,0.218) (0.748,0.233) (0.750,0.251) (0.752,0.233) (0.754,0.262) (0.756,0.332) (0.758,0.374) (0.760,0.442) (0.762,0.405) (0.764,0.446) (0.766,0.481) (0.768,0.454) (0.770,0.496) (0.772,0.497) (0.774,0.469) (0.776,0.469) (0.778,0.481) (0.780,0.418) (0.782,0.414) (0.784,0.417) (0.786,0.431) (0.788,0.434) (0.790,0.411) (0.792,0.471) (0.794,0.468) (0.796,0.497) (0.798,0.556) (0.800,0.616) (0.802,0.685) (0.804,0.638) (0.806,0.742) (0.808,0.717) (0.810,0.711) (0.812,0.734) (0.814,0.844) (0.816,0.829) (0.818,0.835) (0.820,0.803) (0.822,0.706) (0.824,0.738) (0.826,0.779) (0.828,0.796) (0.830,0.777) (0.832,0.814) (0.834,0.781) (0.836,0.751) (0.838,0.734) (0.840,0.670) (0.842,0.678) (0.844,0.672) (0.846,0.724) (0.848,0.748) (0.850,0.811) (0.852,0.807) (0.854,0.771) (0.856,0.674) (0.858,0.706) (0.860,0.635) (0.862,0.563) (0.864,0.491) (0.866,0.436) (0.868,0.462) (0.870,0.483) (0.872,0.446) (0.874,0.424) (0.876,0.379) (0.878,0.419) (0.880,0.430) (0.882,0.420) (0.884,0.394) (0.886,0.391) (0.888,0.419) (0.890,0.416) (0.892,0.468) (0.894,0.477) (0.896,0.461) (0.898,0.529) (0.900,0.538) (0.902,0.576) (0.904,0.506) (0.906,0.504) (0.908,0.628) (0.910,0.588) (0.912,0.546) (0.914,0.537) (0.916,0.433) (0.918,0.371) (0.920,0.357) (0.922,0.327) (0.924,0.290) (0.926,0.287) (0.928,0.324) (0.930,0.357) (0.932,0.350) (0.934,0.329) (0.936,0.304) (0.938,0.282) (0.940,0.312) (0.942,0.284) (0.944,0.339) (0.946,0.316) (0.948,0.242) (0.950,0.303) (0.952,0.331) (0.954,0.345) (0.956,0.326) (0.958,0.407) (0.960,0.429) (0.962,0.453) (0.964,0.432) (0.966,0.383) (0.968,0.369) (0.970,0.372) (0.972,0.412) (0.974,0.407) (0.976,0.335) (0.978,0.310) (0.980,0.333) (0.982,0.342) (0.984,0.278) (0.986,0.284) (0.988,0.242) (0.990,0.202) (0.992,0.211) (0.994,0.226) (0.996,0.265) (0.998,0.263) (1.000,0.214)};
  \addplot[SARed,thick] coordinates {(0.000,0.000) (0.002,0.063) (0.004,0.112) (0.006,0.118) (0.008,0.118) (0.010,0.118) (0.012,0.118) (0.014,0.118) (0.016,0.118) (0.018,0.120) (0.020,0.146) (0.022,0.189) (0.024,0.214) (0.026,0.216) (0.028,0.231) (0.030,0.231) (0.032,0.231) (0.034,0.231) (0.036,0.231) (0.038,0.231) (0.040,0.231) (0.042,0.231) (0.044,0.231) (0.046,0.231) (0.048,0.231) (0.050,0.231) (0.052,0.231) (0.054,0.231) (0.056,0.231) (0.058,0.231) (0.060,0.231) (0.062,0.231) (0.064,0.231) (0.066,0.231) (0.068,0.252) (0.070,0.252) (0.072,0.257) (0.074,0.257) (0.076,0.270) (0.078,0.287) (0.080,0.287) (0.082,0.287) (0.084,0.306) (0.086,0.371) (0.088,0.398) (0.090,0.430) (0.092,0.430) (0.094,0.430) (0.096,0.430) (0.098,0.430) (0.100,0.430) (0.102,0.430) (0.104,0.430) (0.106,0.430) (0.108,0.430) (0.110,0.430) (0.112,0.430) (0.114,0.430) (0.116,0.430) (0.118,0.430) (0.120,0.430) (0.122,0.430) (0.124,0.430) (0.126,0.430) (0.128,0.430) (0.130,0.430) (0.132,0.430) (0.134,0.430) (0.136,0.430) (0.138,0.430) (0.140,0.430) (0.142,0.430) (0.144,0.430) (0.146,0.430) (0.148,0.430) (0.150,0.430) (0.152,0.430) (0.154,0.430) (0.156,0.430) (0.158,0.430) (0.160,0.430) (0.162,0.430) (0.164,0.430) (0.166,0.430) (0.168,0.430) (0.170,0.430) (0.172,0.430) (0.174,0.430) (0.176,0.430) (0.178,0.430) (0.180,0.430) (0.182,0.430) (0.184,0.430) (0.186,0.430) (0.188,0.430) (0.190,0.430) (0.192,0.430) (0.194,0.430) (0.196,0.430) (0.198,0.430) (0.200,0.430) (0.202,0.430) (0.204,0.430) (0.206,0.430) (0.208,0.430) (0.210,0.430) (0.212,0.430) (0.214,0.430) (0.216,0.430) (0.218,0.430) (0.220,0.430) (0.222,0.430) (0.224,0.430) (0.226,0.430) (0.228,0.430) (0.230,0.430) (0.232,0.430) (0.234,0.430) (0.236,0.430) (0.238,0.430) (0.240,0.430) (0.242,0.430) (0.244,0.430) (0.246,0.436) (0.248,0.436) (0.250,0.436) (0.252,0.436) (0.254,0.436) (0.256,0.436) (0.258,0.436) (0.260,0.443) (0.262,0.448) (0.264,0.448) (0.266,0.448) (0.268,0.448) (0.270,0.448) (0.272,0.448) (0.274,0.448) (0.276,0.448) (0.278,0.448) (0.280,0.448) (0.282,0.448) (0.284,0.448) (0.286,0.448) (0.288,0.448) (0.290,0.448) (0.292,0.448) (0.294,0.448) (0.296,0.448) (0.298,0.448) (0.300,0.448) (0.302,0.448) (0.304,0.448) (0.306,0.448) (0.308,0.448) (0.310,0.448) (0.312,0.448) (0.314,0.448) (0.316,0.448) (0.318,0.448) (0.320,0.448) (0.322,0.448) (0.324,0.448) (0.326,0.448) (0.328,0.448) (0.330,0.448) (0.332,0.448) (0.334,0.448) (0.336,0.448) (0.338,0.448) (0.340,0.448) (0.342,0.448) (0.344,0.448) (0.346,0.448) (0.348,0.448) (0.350,0.448) (0.352,0.448) (0.354,0.448) (0.356,0.448) (0.358,0.448) (0.360,0.448) (0.362,0.448) (0.364,0.448) (0.366,0.448) (0.368,0.448) (0.370,0.448) (0.372,0.448) (0.374,0.448) (0.376,0.448) (0.378,0.448) (0.380,0.448) (0.382,0.448) (0.384,0.448) (0.386,0.448) (0.388,0.448) (0.390,0.448) (0.392,0.448) (0.394,0.448) (0.396,0.448) (0.398,0.448) (0.400,0.448) (0.402,0.448) (0.404,0.448) (0.406,0.448) (0.408,0.448) (0.410,0.448) (0.412,0.448) (0.414,0.448) (0.416,0.448) (0.418,0.448) (0.420,0.448) (0.422,0.448) (0.424,0.448) (0.426,0.448) (0.428,0.448) (0.430,0.448) (0.432,0.448) (0.434,0.448) (0.436,0.448) (0.438,0.448) (0.440,0.448) (0.442,0.448) (0.444,0.448) (0.446,0.448) (0.448,0.448) (0.450,0.448) (0.452,0.448) (0.454,0.448) (0.456,0.448) (0.458,0.448) (0.460,0.448) (0.462,0.448) (0.464,0.448) (0.466,0.448) (0.468,0.448) (0.470,0.448) (0.472,0.448) (0.474,0.448) (0.476,0.448) (0.478,0.448) (0.480,0.448) (0.482,0.448) (0.484,0.448) (0.486,0.448) (0.488,0.448) (0.490,0.448) (0.492,0.448) (0.494,0.448) (0.496,0.448) (0.498,0.448) (0.500,0.454) (0.502,0.454) (0.504,0.454) (0.506,0.454) (0.508,0.454) (0.510,0.454) (0.512,0.454) (0.514,0.454) (0.516,0.454) (0.518,0.454) (0.520,0.454) (0.522,0.454) (0.524,0.454) (0.526,0.454) (0.528,0.454) (0.530,0.454) (0.532,0.454) (0.534,0.454) (0.536,0.454) (0.538,0.454) (0.540,0.454) (0.542,0.454) (0.544,0.454) (0.546,0.454) (0.548,0.454) (0.550,0.454) (0.552,0.454) (0.554,0.454) (0.556,0.454) (0.558,0.454) (0.560,0.454) (0.562,0.454) (0.564,0.454) (0.566,0.454) (0.568,0.454) (0.570,0.454) (0.572,0.454) (0.574,0.454) (0.576,0.454) (0.578,0.454) (0.580,0.454) (0.582,0.454) (0.584,0.454) (0.586,0.519) (0.588,0.562) (0.590,0.589) (0.592,0.598) (0.594,0.602) (0.596,0.602) (0.598,0.602) (0.600,0.602) (0.602,0.602) (0.604,0.603) (0.606,0.616) (0.608,0.636) (0.610,0.636) (0.612,0.666) (0.614,0.701) (0.616,0.716) (0.618,0.734) (0.620,0.734) (0.622,0.741) (0.624,0.765) (0.626,0.808) (0.628,0.846) (0.630,0.886) (0.632,0.886) (0.634,0.886) (0.636,0.886) (0.638,0.886) (0.640,0.886) (0.642,0.886) (0.644,0.886) (0.646,0.886) (0.648,0.886) (0.650,0.886) (0.652,0.886) (0.654,0.886) (0.656,0.886) (0.658,0.923) (0.660,0.923) (0.662,0.928) (0.664,0.964) (0.666,0.971) (0.668,1.026) (0.670,1.068) (0.672,1.084) (0.674,1.091) (0.676,1.091) (0.678,1.091) (0.680,1.091) (0.682,1.091) (0.684,1.092) (0.686,1.092) (0.688,1.092) (0.690,1.099) (0.692,1.106) (0.694,1.114) (0.696,1.127) (0.698,1.133) (0.700,1.133) (0.702,1.149) (0.704,1.161) (0.706,1.203) (0.708,1.232) (0.710,1.232) (0.712,1.232) (0.714,1.234) (0.716,1.244) (0.718,1.244) (0.720,1.244) (0.722,1.244) (0.724,1.244) (0.726,1.244) (0.728,1.244) (0.730,1.244) (0.732,1.244) (0.734,1.244) (0.736,1.244) (0.738,1.244) (0.740,1.244) (0.742,1.244) (0.744,1.244) (0.746,1.244) (0.748,1.244) (0.750,1.244) (0.752,1.244) (0.754,1.244) (0.756,1.244) (0.758,1.244) (0.760,1.244) (0.762,1.244) (0.764,1.244) (0.766,1.244) (0.768,1.244) (0.770,1.244) (0.772,1.244) (0.774,1.244) (0.776,1.244) (0.778,1.244) (0.780,1.244) (0.782,1.244) (0.784,1.244) (0.786,1.244) (0.788,1.244) (0.790,1.244) (0.792,1.244) (0.794,1.244) (0.796,1.244) (0.798,1.244) (0.800,1.244) (0.802,1.244) (0.804,1.244) (0.806,1.244) (0.808,1.244) (0.810,1.244) (0.812,1.244) (0.814,1.244) (0.816,1.244) (0.818,1.244) (0.820,1.244) (0.822,1.244) (0.824,1.244) (0.826,1.244) (0.828,1.244) (0.830,1.244) (0.832,1.244) (0.834,1.244) (0.836,1.244) (0.838,1.244) (0.840,1.244) (0.842,1.244) (0.844,1.244) (0.846,1.244) (0.848,1.244) (0.850,1.244) (0.852,1.244) (0.854,1.244) (0.856,1.244) (0.858,1.244) (0.860,1.244) (0.862,1.244) (0.864,1.244) (0.866,1.244) (0.868,1.244) (0.870,1.244) (0.872,1.244) (0.874,1.244) (0.876,1.244) (0.878,1.244) (0.880,1.244) (0.882,1.244) (0.884,1.244) (0.886,1.244) (0.888,1.244) (0.890,1.244) (0.892,1.244) (0.894,1.244) (0.896,1.244) (0.898,1.244) (0.900,1.244) (0.902,1.244) (0.904,1.244) (0.906,1.244) (0.908,1.244) (0.910,1.244) (0.912,1.244) (0.914,1.244) (0.916,1.244) (0.918,1.244) (0.920,1.244) (0.922,1.244) (0.924,1.244) (0.926,1.244) (0.928,1.244) (0.930,1.244) (0.932,1.244) (0.934,1.244) (0.936,1.244) (0.938,1.244) (0.940,1.244) (0.942,1.244) (0.944,1.244) (0.946,1.244) (0.948,1.244) (0.950,1.244) (0.952,1.244) (0.954,1.244) (0.956,1.244) (0.958,1.244) (0.960,1.244) (0.962,1.244) (0.964,1.244) (0.966,1.244) (0.968,1.244) (0.970,1.244) (0.972,1.244) (0.974,1.244) (0.976,1.244) (0.978,1.244) (0.980,1.244) (0.982,1.244) (0.984,1.244) (0.986,1.244) (0.988,1.244) (0.990,1.244) (0.992,1.244) (0.994,1.244) (0.996,1.244) (0.998,1.244) (1.000,1.244)};
  \addplot[SADark!45,thin] coordinates {(0.000,0.000) (0.002,-0.050) (0.004,-0.109) (0.006,-0.029) (0.008,-0.003) (0.010,-0.003) (0.012,0.027) (0.014,0.027) (0.016,-0.062) (0.018,-0.069) (0.020,-0.129) (0.022,-0.179) (0.024,-0.171) (0.026,-0.213) (0.028,-0.201) (0.030,-0.172) (0.032,-0.144) (0.034,-0.137) (0.036,-0.059) (0.038,-0.042) (0.040,-0.028) (0.042,-0.015) (0.044,-0.029) (0.046,-0.052) (0.048,-0.033) (0.050,-0.019) (0.052,0.040) (0.054,0.035) (0.056,-0.032) (0.058,-0.004) (0.060,0.011) (0.062,-0.033) (0.064,-0.115) (0.066,-0.190) (0.068,-0.218) (0.070,-0.186) (0.072,-0.241) (0.074,-0.215) (0.076,-0.252) (0.078,-0.251) (0.080,-0.213) (0.082,-0.172) (0.084,-0.294) (0.086,-0.355) (0.088,-0.387) (0.090,-0.345) (0.092,-0.203) (0.094,-0.139) (0.096,-0.197) (0.098,-0.209) (0.100,-0.279) (0.102,-0.269) (0.104,-0.299) (0.106,-0.267) (0.108,-0.314) (0.110,-0.308) (0.112,-0.254) (0.114,-0.167) (0.116,-0.078) (0.118,-0.082) (0.120,-0.127) (0.122,-0.187) (0.124,-0.217) (0.126,-0.202) (0.128,-0.183) (0.130,-0.169) (0.132,-0.204) (0.134,-0.249) (0.136,-0.272) (0.138,-0.259) (0.140,-0.268) (0.142,-0.280) (0.144,-0.313) (0.146,-0.302) (0.148,-0.330) (0.150,-0.240) (0.152,-0.162) (0.154,-0.150) (0.156,-0.164) (0.158,-0.073) (0.160,-0.052) (0.162,-0.087) (0.164,-0.084) (0.166,-0.112) (0.168,-0.103) (0.170,-0.166) (0.172,-0.149) (0.174,-0.137) (0.176,-0.071) (0.178,-0.064) (0.180,-0.034) (0.182,-0.020) (0.184,-0.088) (0.186,-0.095) (0.188,-0.173) (0.190,-0.192) (0.192,-0.207) (0.194,-0.277) (0.196,-0.248) (0.198,-0.276) (0.200,-0.309) (0.202,-0.335) (0.204,-0.312) (0.206,-0.300) (0.208,-0.338) (0.210,-0.350) (0.212,-0.275) (0.214,-0.295) (0.216,-0.296) (0.218,-0.317) (0.220,-0.254) (0.222,-0.223) (0.224,-0.200) (0.226,-0.253) (0.228,-0.287) (0.230,-0.308) (0.232,-0.320) (0.234,-0.298) (0.236,-0.383) (0.238,-0.350) (0.240,-0.341) (0.242,-0.369) (0.244,-0.370) (0.246,-0.377) (0.248,-0.324) (0.250,-0.343) (0.252,-0.334) (0.254,-0.348) (0.256,-0.356) (0.258,-0.352) (0.260,-0.407) (0.262,-0.387) (0.264,-0.357) (0.266,-0.292) (0.268,-0.223) (0.270,-0.266) (0.272,-0.246) (0.274,-0.135) (0.276,-0.080) (0.278,-0.070) (0.280,-0.026) (0.282,-0.038) (0.284,-0.013) (0.286,0.021) (0.288,0.100) (0.290,0.144) (0.292,0.083) (0.294,0.105) (0.296,0.143) (0.298,0.164) (0.300,0.187) (0.302,0.151) (0.304,0.193) (0.306,0.195) (0.308,0.181) (0.310,0.203) (0.312,0.206) (0.314,0.192) (0.316,0.191) (0.318,0.138) (0.320,0.089) (0.322,0.071) (0.324,0.007) (0.326,0.053) (0.328,-0.011) (0.330,-0.006) (0.332,0.002) (0.334,-0.048) (0.336,-0.060) (0.338,-0.085) (0.340,-0.116) (0.342,-0.137) (0.344,-0.095) (0.346,-0.131) (0.348,-0.189) (0.350,-0.220) (0.352,-0.233) (0.354,-0.173) (0.356,-0.099) (0.358,-0.045) (0.360,-0.046) (0.362,-0.046) (0.364,-0.010) (0.366,-0.036) (0.368,-0.034) (0.370,0.021) (0.372,-0.038) (0.374,-0.017) (0.376,-0.037) (0.378,-0.024) (0.380,-0.030) (0.382,-0.033) (0.384,-0.081) (0.386,-0.040) (0.388,-0.020) (0.390,-0.107) (0.392,-0.083) (0.394,-0.043) (0.396,-0.039) (0.398,-0.075) (0.400,-0.098) (0.402,-0.066) (0.404,-0.088) (0.406,-0.117) (0.408,-0.124) (0.410,-0.157) (0.412,-0.128) (0.414,-0.182) (0.416,-0.174) (0.418,-0.182) (0.420,-0.187) (0.422,-0.119) (0.424,-0.132) (0.426,-0.191) (0.428,-0.187) (0.430,-0.240) (0.432,-0.185) (0.434,-0.143) (0.436,-0.160) (0.438,-0.150) (0.440,-0.192) (0.442,-0.215) (0.444,-0.041) (0.446,-0.107) (0.448,-0.151) (0.450,-0.253) (0.452,-0.280) (0.454,-0.314) (0.456,-0.277) (0.458,-0.236) (0.460,-0.201) (0.462,-0.166) (0.464,-0.131) (0.466,-0.205) (0.468,-0.180) (0.470,-0.182) (0.472,-0.164) (0.474,-0.090) (0.476,-0.098) (0.478,-0.177) (0.480,-0.230) (0.482,-0.260) (0.484,-0.309) (0.486,-0.368) (0.488,-0.257) (0.490,-0.261) (0.492,-0.257) (0.494,-0.240) (0.496,-0.324) (0.498,-0.359) (0.500,-0.427) (0.502,-0.422) (0.504,-0.363) (0.506,-0.404) (0.508,-0.321) (0.510,-0.275) (0.512,-0.209) (0.514,-0.171) (0.516,-0.099) (0.518,-0.143) (0.520,-0.152) (0.522,-0.182) (0.524,-0.227) (0.526,-0.237) (0.528,-0.277) (0.530,-0.262) (0.532,-0.309) (0.534,-0.341) (0.536,-0.343) (0.538,-0.268) (0.540,-0.256) (0.542,-0.280) (0.544,-0.263) (0.546,-0.332) (0.548,-0.310) (0.550,-0.283) (0.552,-0.241) (0.554,-0.197) (0.556,-0.165) (0.558,-0.115) (0.560,-0.187) (0.562,-0.244) (0.564,-0.280) (0.566,-0.337) (0.568,-0.305) (0.570,-0.348) (0.572,-0.293) (0.574,-0.370) (0.576,-0.395) (0.578,-0.357) (0.580,-0.389) (0.582,-0.422) (0.584,-0.425) (0.586,-0.498) (0.588,-0.553) (0.590,-0.572) (0.592,-0.589) (0.594,-0.472) (0.596,-0.434) (0.598,-0.496) (0.600,-0.479) (0.602,-0.524) (0.604,-0.579) (0.606,-0.608) (0.608,-0.610) (0.610,-0.581) (0.612,-0.628) (0.614,-0.674) (0.616,-0.707) (0.618,-0.626) (0.620,-0.662) (0.622,-0.721) (0.624,-0.754) (0.626,-0.779) (0.628,-0.797) (0.630,-0.844) (0.632,-0.811) (0.634,-0.805) (0.636,-0.793) (0.638,-0.824) (0.640,-0.783) (0.642,-0.796) (0.644,-0.798) (0.646,-0.746) (0.648,-0.738) (0.650,-0.779) (0.652,-0.758) (0.654,-0.785) (0.656,-0.851) (0.658,-0.909) (0.660,-0.852) (0.662,-0.892) (0.664,-0.947) (0.666,-0.939) (0.668,-1.000) (0.670,-1.065) (0.672,-1.056) (0.674,-1.011) (0.676,-1.051) (0.678,-1.019) (0.680,-1.007) (0.682,-1.042) (0.684,-1.076) (0.686,-1.037) (0.688,-1.058) (0.690,-1.096) (0.692,-1.063) (0.694,-1.107) (0.696,-1.101) (0.698,-1.103) (0.700,-1.088) (0.702,-1.126) (0.704,-1.129) (0.706,-1.200) (0.708,-1.199) (0.710,-1.186) (0.712,-1.169) (0.714,-1.218) (0.716,-1.155) (0.718,-1.150) (0.720,-1.099) (0.722,-1.101) (0.724,-1.128) (0.726,-1.126) (0.728,-1.098) (0.730,-1.093) (0.732,-1.164) (0.734,-1.140) (0.736,-1.141) (0.738,-1.051) (0.740,-1.018) (0.742,-0.963) (0.744,-1.006) (0.746,-1.027) (0.748,-1.012) (0.750,-0.994) (0.752,-1.011) (0.754,-0.983) (0.756,-0.912) (0.758,-0.871) (0.760,-0.803) (0.762,-0.840) (0.764,-0.798) (0.766,-0.764) (0.768,-0.790) (0.770,-0.749) (0.772,-0.747) (0.774,-0.775) (0.776,-0.776) (0.778,-0.763) (0.780,-0.826) (0.782,-0.830) (0.784,-0.828) (0.786,-0.814) (0.788,-0.811) (0.790,-0.833) (0.792,-0.773) (0.794,-0.777) (0.796,-0.747) (0.798,-0.689) (0.800,-0.629) (0.802,-0.560) (0.804,-0.606) (0.806,-0.502) (0.808,-0.527) (0.810,-0.533) (0.812,-0.510) (0.814,-0.400) (0.816,-0.415) (0.818,-0.410) (0.820,-0.442) (0.822,-0.538) (0.824,-0.507) (0.826,-0.466) (0.828,-0.448) (0.830,-0.467) (0.832,-0.430) (0.834,-0.463) (0.836,-0.494) (0.838,-0.510) (0.840,-0.574) (0.842,-0.567) (0.844,-0.573) (0.846,-0.520) (0.848,-0.496) (0.850,-0.433) (0.852,-0.437) (0.854,-0.473) (0.856,-0.571) (0.858,-0.538) (0.860,-0.610) (0.862,-0.682) (0.864,-0.753) (0.866,-0.808) (0.868,-0.782) (0.870,-0.761) (0.872,-0.799) (0.874,-0.820) (0.876,-0.866) (0.878,-0.826) (0.880,-0.815) (0.882,-0.824) (0.884,-0.851) (0.886,-0.854) (0.888,-0.826) (0.890,-0.828) (0.892,-0.776) (0.894,-0.767) (0.896,-0.783) (0.898,-0.716) (0.900,-0.707) (0.902,-0.669) (0.904,-0.738) (0.906,-0.741) (0.908,-0.617) (0.910,-0.656) (0.912,-0.698) (0.914,-0.707) (0.916,-0.811) (0.918,-0.873) (0.920,-0.888) (0.922,-0.917) (0.924,-0.955) (0.926,-0.958) (0.928,-0.921) (0.930,-0.888) (0.932,-0.895) (0.934,-0.915) (0.936,-0.941) (0.938,-0.962) (0.940,-0.932) (0.942,-0.961) (0.944,-0.905) (0.946,-0.929) (0.948,-1.003) (0.950,-0.941) (0.952,-0.913) (0.954,-0.899) (0.956,-0.919) (0.958,-0.837) (0.960,-0.815) (0.962,-0.792) (0.964,-0.812) (0.966,-0.862) (0.968,-0.876) (0.970,-0.873) (0.972,-0.832) (0.974,-0.837) (0.976,-0.909) (0.978,-0.934) (0.980,-0.911) (0.982,-0.902) (0.984,-0.966) (0.986,-0.960) (0.988,-1.002) (0.990,-1.042) (0.992,-1.033) (0.994,-1.019) (0.996,-0.979) (0.998,-0.982) (1.000,-1.031)};
  \legend{$\abs{B_t}$, 局部时 $L_t$, $W_t=\int_0^t\sgn(B_s)\dd B_s$}
\end{axis}
\end{tikzpicture}
\caption{Tanaka 公式 $\abs{B_t}=W_t+L_t$. 局部时 $L$ (红色) 只在 $\abs{B}$ 触及零时增长, 在 $\abs{B}$ 远离零的时段保持不变; 它恰好把布朗运动 $W$ (灰色) ``推回'' 非负区域. 图中 $L$ 用 $\frac{1}{2\varepsilon}\int_0^t\1_{\{\abs{B_s}<\varepsilon\}}\dd s$ ($\varepsilon=0.01$) 近似}
\label{fig:local-time}
\end{figure}

\begin{example}[局部时的简单性质]\label{ex:local-time}
(1) \emph{期望.} 在 Tanaka 公式两端取期望, 由于 $\sgn(B)\in\mathscr{L}^2_T$, 随机积分的期望为零, 所以
\[
\E L_t=\E\abs{B_t}=\sqrt{\frac{2t}{\pi}}.
\]
也可以从定义直接验证: $\frac{1}{2\varepsilon}\int_0^t\P(\abs{B_s}<\varepsilon)\dd s\to\int_0^t\frac{1}{\sqrt{2\pi s}}\dd s=\sqrt{\frac{2t}{\pi}}$, 因为 $\frac{1}{2\varepsilon}\P(\abs{B_s}<\varepsilon)$ 趋于 $B_s$ 在 $0$ 处的密度 $\frac{1}{\sqrt{2\pi s}}$.

(2) \emph{反射布朗运动.} 由例~\ref{ex:levy-applications} (1), $W_t=\int_0^t\sgn(B_s)\dd B_s$ 是布朗运动. 所以 Tanaka 公式表明, 反射布朗运动 $\abs{B}$ 等于布朗运动 $W$ 与增过程 $L$ 之和, 后者只在零点增长, 恰好足以使 $W+L\ge0$ (图~\ref{fig:local-time}). 可以证明 (Skorokhod 引理), 这样的 $L$ 是唯一的, 且 $L_t=\max_{s\le t}(-W_s)$. 再由 $-W$ 也是布朗运动, 得到 L\'evy 定理: $(\abs{B_t},L_t)_{t\ge0}$ 与 $(M_t-B_t,M_t)_{t\ge0}$ 同分布, 其中 $M_t=\max_{s\le t}B_s$ (参见 4.5.2 小节末的注). 特别地, $L_t\overset{d}{=}M_t\overset{d}{=}\abs{B_t}$, 这与 (1) 中的期望一致.

(3) \emph{一般水平的局部时.} 对 $a\in\R$, 同样定义 $L^a_t=\lim_{\varepsilon\downarrow0}\frac{1}{2\varepsilon}\int_0^t\1_{\{\abs{B_s-a}<\varepsilon\}}\dd s$. 则 Tanaka 公式推广为 $\abs{B_t-a}=\abs{a}+\int_0^t\sgn(B_s-a)\dd B_s+L^a_t$, 并且有\textbf{占位时公式}: 对有界 Borel 函数 $g$,
\[
\int_0^tg(B_s)\dd s=\int_{\R}g(a)L^a_t\dd a,
\]
即 $a\mapsto L^a_t$ 是布朗运动在 $[0,t]$ 上的 ``占位时间'' 的密度 (见~\cite{karatzas} 3.6 节). 若 $f$ 是两个凸函数之差, 由此得到 It\^o-Tanaka 公式 $f(B_t)=f(B_0)+\int_0^tf'_-(B_s)\dd B_s+\frac{1}{2}\int_{\R}L^a_t\,f''(\dd a)$, 其中 $f''$ 理解为广义二阶导数 (一个测度).
\end{example}

% ---------------------------------------------------------------------
\subsection*{本章小结}
\addcontentsline{toc}{subsection}{本章小结}

\begin{center}
\begin{minipage}{\textwidth}
\centering
\captionof{table}{随机微积分的主要结论}
\label{tab:ch5-summary}
\renewcommand{\arraystretch}{1.5}
\begin{tabular}{>{\raggedright\arraybackslash}p{0.22\textwidth} >{\raggedright\arraybackslash}p{0.46\textwidth} >{\raggedright\arraybackslash}p{0.22\textwidth}}
\toprule
结论 & 内容 & 后续应用 \\
\midrule
It\^o 等距 & $\E(\int_0^tH\dd B)^2=\E\int_0^tH^2\dd s$ & 积分的构造与估计 \\
It\^o 积分的性质 & $H\in\mathscr{L}^2_T$ 时是连续鞅, $\ip{H\cdot B}=\int H^2\dd s$ & SDE 解的鞅部分 \\
局部化 & $\int_0^TH^2\dd s<\infty$ $\as$ 时是连续局部鞅 & 一般的被积过程 \\
It\^o 公式 & $\dd f=\partial_tf\dd t+\partial_xf\dd X+\frac12\partial_{xx}f\dd\ip{X}$ & 一切计算 \\
随机指数 & $\dd Z=\theta Z\dd B$, $Z=\mathscr{E}(\theta\cdot B)$ & Girsanov, 线性 SDE \\
L\'evy 刻画 & $\ip{M}_t=t$ 的连续鞅是布朗运动 & 识别布朗运动 \\
Girsanov 定理 & 测度变换 $\dd\widetilde{\P}=Z_T\dd\P$ 给 $B$ 加漂移 $\theta$ & 风险中性定价, 弱解 \\
鞅表示定理 & 布朗泛函 $F=\E F+\int H\dd B$ & 市场完备性, 对冲 \\
Tanaka 公式 & $\abs{B_t}=\int\sgn(B)\dd B+L_t$ & 反射, 非光滑函数 \\
\bottomrule
\end{tabular}
\end{minipage}
\end{center}

% =====================================================================
%  第 6 章  随机微分方程与扩散过程
% =====================================================================
\chapter{随机微分方程与扩散过程}
\label{ch:6}

第~5~章开头提出的问题, 现在可以回答了. 一个受白噪声扰动的动力系统
\[
\frac{\dd X_t}{\dd t}=b(t,X_t)+\sigma(t,X_t)\,\text{``}\frac{\dd B_t}{\dd t}\text{''}
\]
应当理解为积分方程 $X_t=X_0+\int_0^tb(s,X_s)\dd s+\int_0^t\sigma(s,X_s)\dd B_s$, 其中第二个积分是 It\^o 积分. 这样的方程称为\textbf{随机微分方程} (SDE). 它的解是一类特别重要的马尔可夫过程, 称为\textbf{扩散过程}: 局部地看, 它在 $(t,t+\dd t]$ 内的增量近似服从均值为 $b\dd t$、方差为 $\sigma^2\dd t$ 的正态分布.

本章是全书的综合. 存在唯一性定理的证明结合了 Doob 不等式 (第~3~章) 与 It\^o 等距 (第~5~章), 马尔可夫性的推导依赖代入引理与单调类定理 (第~2,~4~章). 生成元与 Kolmogorov 方程推广了布朗运动与热方程的联系 (第~4~章), 而 It\^o 公式与可选停时定理又为 Feynman-Kac 公式及 Dirichlet 问题提供了工具. 期权定价除了使用 Girsanov 定理与鞅表示定理, 还要用到反射原理和首达时的分布.

本章的结构如下. 6.1 节给出定义并求解线性方程; 6.2 节证明 Lipschitz 条件下强解的存在唯一性, 并讨论弱解; 6.3 节研究解的马尔可夫性与生成元, 推导 Kolmogorov 方程并求解一维扩散的出口问题; 6.4 节建立随机微分方程与偏微分方程的联系, 并以 Black-Scholes 理论、障碍期权与美式期权为应用; 6.5 节介绍数值方法. 为了使记号简洁, 理论部分只讨论一维方程 ($X$ 与 $B$ 都是一维的); 多维的情形只需把绝对值换成范数, 证明逐字相同. 本章中信息流 $\{\mathcal{F}_t\}$ 满足通常条件, $B$ 是 $\{\mathcal{F}_t\}$-布朗运动.

% =====================================================================
\section{随机微分方程}
\label{sec:6-1}

% ---------------------------------------------------------------------
\subsection{强解与弱解}

设 $b,\sigma:[0,T]\times\R\to\R$ 是 Borel 函数, 分别称为\textbf{漂移系数}与\textbf{扩散系数}. 考虑方程
\begin{equation}\label{eq:sde}
\dd X_t=b(t,X_t)\dd t+\sigma(t,X_t)\dd B_t,\qquad X_0=\xi.
\end{equation}

\begin{definition}[强解]\label{def:strong-solution}
设 $\xi$ 是 $\mathcal{F}_0$-可测的随机变量. 称轨道连续的适应过程 $X=\{X_t\}_{t\in[0,T]}$ 为方程 \eqref{eq:sde} 的\textbf{强解}, 如果 $\int_0^T\big(\abs{b(s,X_s)}+\sigma(s,X_s)^2\big)\dd s<\infty$ $\as$, 且几乎必然地, 对一切 $t\in[0,T]$,
\begin{equation}\label{eq:sde-integral}
X_t=\xi+\int_0^tb(s,X_s)\dd s+\int_0^t\sigma(s,X_s)\dd B_s.
\end{equation}
若方程的任意两个强解 (对同一个 $B$ 与 $\xi$) 都不可区分, 称方程有\textbf{轨道唯一性}.
\end{definition}

条件 $\int_0^T(\abs{b}+\sigma^2)\dd s<\infty$ 保证 \eqref{eq:sde-integral} 右端有意义: $s\mapsto\sigma(s,X_s)$ 循序可测 (Borel 函数与循序可测过程的复合) 且属于 $\mathscr{L}^2_{\mathrm{loc}}$. 所以强解是一个 It\^o 过程 (定义~\ref{def:ito-process}), 系数为 $b(t,X_t)$ 与 $\sigma(t,X_t)$. ``强'' 是指解在\emph{给定的}概率空间上、相对于\emph{给定的}布朗运动构造的, 它是 ``输入'' $(\xi,B)$ 的函数. 有时我们只关心解的分布, 此时可以放宽要求.

\begin{definition}[弱解]\label{def:weak-solution}
称方程 \eqref{eq:sde} (初始分布为 $\mu$) 有\textbf{弱解}, 如果存在一个概率空间 $(\Omega,\mathcal{F},\P)$、其上满足通常条件的信息流 $\{\mathcal{F}_t\}$、一个 $\{\mathcal{F}_t\}$-布朗运动 $W$ 与一个轨道连续的适应过程 $X$, 使得 $X_0\sim\mu$, 且 \eqref{eq:sde-integral} 在 $B$ 换成 $W$ 后成立. 若任意两个弱解 (可能在不同的概率空间上) 的 $X$ 有相同的有限维分布, 称方程有\textbf{分布唯一性}.
\end{definition}

强解当然是弱解. 二者的区别在于, 弱解中的布朗运动 $W$ 是和解 $X$ 一起被构造出来的, $X$ 不必是 $W$ 的函数. 6.2.2 小节将给出一个有弱解而没有强解的例子.

% ---------------------------------------------------------------------
\subsection{线性方程}

当 $b$ 与 $\sigma$ 关于 $x$ 是线性函数时, 方程可以显式求解. 方法与常微分方程的 ``常数变易法'' 相同, 只是要用 It\^o 公式代替普通的链式法则.

\begin{proposition}[一维线性方程]\label{prop:linear-sde}
设 $a_1,a_2,\sigma_1,\sigma_2$ 是常数, $x\in\R$. 方程
\begin{equation}\label{eq:linear-sde}
\dd X_t=(a_1X_t+a_2)\dd t+(\sigma_1X_t+\sigma_2)\dd B_t,\qquad X_0=x
\end{equation}
的解为
\begin{align*}
X_t&=\Phi_t\Big(x+(a_2-\sigma_1\sigma_2)\int_0^t\Phi_s^{-1}\dd s+\sigma_2\int_0^t\Phi_s^{-1}\dd B_s\Big),\\
\Phi_t&=\exp\Big(\big(a_1-\tfrac{1}{2}\sigma_1^2\big)t+\sigma_1B_t\Big).
\end{align*}
\end{proposition}

\begin{proof}
由 It\^o 公式 (例~\ref{ex:gbm-ou} (1)), $\Phi$ 满足齐次方程 $\dd\Phi_t=a_1\Phi_t\dd t+\sigma_1\Phi_t\dd B_t$, $\Phi_0=1$. 对 $f(x)=x^{-1}$ 应用 It\^o 公式 ($f'=-x^{-2}$, $f''=2x^{-3}$, $\dd\ip{\Phi}_t=\sigma_1^2\Phi_t^2\dd t$),
\[
\dd\Phi_t^{-1}=-\Phi_t^{-2}\dd\Phi_t+\Phi_t^{-3}\dd\ip{\Phi}_t=\Phi_t^{-1}\big((\sigma_1^2-a_1)\dd t-\sigma_1\dd B_t\big).
\]
记 $Y_t$ 为括号中的过程, 则 $\dd Y_t=(a_2-\sigma_1\sigma_2)\Phi_t^{-1}\dd t+\sigma_2\Phi_t^{-1}\dd B_t$, $\ip{\Phi,Y}$ 的密度为 $\sigma_1\Phi_t\cdot\sigma_2\Phi_t^{-1}=\sigma_1\sigma_2$. 由乘积法则 (推论~\ref{cor:product-rule}),
\begin{align*}
\dd(\Phi_tY_t)&=Y_t\dd\Phi_t+\Phi_t\dd Y_t+\dd\ip{\Phi,Y}_t\\
&=a_1\Phi_tY_t\dd t+\sigma_1\Phi_tY_t\dd B_t+(a_2-\sigma_1\sigma_2)\dd t+\sigma_2\dd B_t+\sigma_1\sigma_2\dd t\\
&=(a_1X_t+a_2)\dd t+(\sigma_1X_t+\sigma_2)\dd B_t,
\end{align*}
其中 $X_t=\Phi_tY_t$, $X_0=x$. (解的唯一性由下一节的定理~\ref{thm:sde-existence} 保证.)
\end{proof}

由于 $\Phi$ 是可以显式写出的正过程, 命题~\ref{prop:linear-sde} 给出的解是 $B$ 的显式泛函. 几个重要的特例如下.

\begin{example}[可解方程]\label{ex:linear-examples}
(1) \emph{几何布朗运动.} $a_2=\sigma_2=0$, $a_1=\mu$, $\sigma_1=\sigma$: $X_t=x\Phi_t=x\exp((\mu-\frac{1}{2}\sigma^2)t+\sigma B_t)$, 这就是 4.5.5 小节的几何布朗运动.

(2) \emph{Ornstein-Uhlenbeck 过程 (Vasicek 模型).} 设 $\kappa,\sigma>0$, $\theta\in\R$, 方程
\[
\dd X_t=\kappa(\theta-X_t)\dd t+\sigma\dd B_t
\]
描述一个被拉向长期均值 $\theta$ 的过程, 回复速度为 $\kappa$. 这里 $\sigma_1=0$, $\Phi_t=e^{-\kappa t}$, 所以
\[
X_t=\theta+(x-\theta)e^{-\kappa t}+\sigma\int_0^te^{-\kappa(t-s)}\dd B_s.
\]
由 Wiener 积分 (命题~\ref{prop:wiener-integral}), $X_t\sim N\big(\theta+(x-\theta)e^{-\kappa t},\frac{\sigma^2}{2\kappa}(1-e^{-2\kappa t})\big)$. 当 $t\to\infty$ 时, 分布收敛到 $N(\theta,\frac{\sigma^2}{2\kappa})$, 与初值无关 (图~\ref{fig:ou-paths}). 在金融中, 这是短期利率的 Vasicek 模型.

(3) \emph{布朗桥.} 方程 $\dd X_t=-\frac{X_t}{1-t}\dd t+\dd B_t$ ($0\le t<1$), $X_0=0$ 的系数依赖于时间. 同样的常数变易法 (以 $\Phi_t=1-t$ 代替) 给出
\[
X_t=(1-t)\int_0^t\frac{\dd B_s}{1-s},\qquad0\le t<1.
\]
由命题~\ref{prop:wiener-integral}, 对 $s\le t<1$, $\Cov(X_s,X_t)=(1-s)(1-t)\int_0^s\frac{\dd u}{(1-u)^2}=(1-s)(1-t)\big(\frac{1}{1-s}-1\big)=s(1-t)$, 这正是布朗桥的协方差 (定义~\ref{def:bridge}). 漂移 $-\frac{X_t}{1-t}$ 在 $t\to1$ 时越来越强, 把过程 ``钉'' 回零点.
\end{example}

\begin{figure}[htbp]
\centering
\begin{tikzpicture}
\begin{axis}[SAaxis, axis x line=bottom, axis y line=left, width=12cm, height=5.6cm, xmin=0, xmax=6.1, ymin=-2.2, ymax=3.4,
  xtick={0,1,...,6}, ytick={-2,-1,0,1,2,3}, xlabel={$t$}]
  \addplot[draw=none,fill=SAYellow!25] coordinates {(0.000,3.000) (0.050,3.290) (0.100,3.317) (0.150,3.302) (0.200,3.268) (0.250,3.223) (0.300,3.172) (0.350,3.117) (0.400,3.060) (0.450,3.002) (0.500,2.944) (0.550,2.886) (0.600,2.829) (0.650,2.772) (0.700,2.717) (0.750,2.664) (0.800,2.611) (0.850,2.561) (0.900,2.512) (0.950,2.464) (1.000,2.419) (1.050,2.375) (1.100,2.332) (1.150,2.291) (1.200,2.252) (1.250,2.214) (1.300,2.178) (1.350,2.144) (1.400,2.110) (1.450,2.078) (1.500,2.048) (1.550,2.019) (1.600,1.991) (1.650,1.964) (1.700,1.938) (1.750,1.914) (1.800,1.891) (1.850,1.868) (1.900,1.847) (1.950,1.827) (2.000,1.807) (2.050,1.789) (2.100,1.771) (2.150,1.754) (2.200,1.738) (2.250,1.723) (2.300,1.708) (2.350,1.694) (2.400,1.681) (2.450,1.668) (2.500,1.656) (2.550,1.644) (2.600,1.633) (2.650,1.623) (2.700,1.613) (2.750,1.603) (2.800,1.594) (2.850,1.585) (2.900,1.577) (2.950,1.569) (3.000,1.562) (3.050,1.555) (3.100,1.548) (3.150,1.541) (3.200,1.535) (3.250,1.529) (3.300,1.524) (3.350,1.519) (3.400,1.514) (3.450,1.509) (3.500,1.504) (3.550,1.500) (3.600,1.496) (3.650,1.492) (3.700,1.488) (3.750,1.484) (3.800,1.481) (3.850,1.478) (3.900,1.475) (3.950,1.472) (4.000,1.469) (4.050,1.466) (4.100,1.464) (4.150,1.461) (4.200,1.459) (4.250,1.457) (4.300,1.455) (4.350,1.453) (4.400,1.451) (4.450,1.449) (4.500,1.447) (4.550,1.446) (4.600,1.444) (4.650,1.443) (4.700,1.441) (4.750,1.440) (4.800,1.439) (4.850,1.438) (4.900,1.437) (4.950,1.435) (5.000,1.434) (5.050,1.433) (5.100,1.432) (5.150,1.432) (5.200,1.431) (5.250,1.430) (5.300,1.429) (5.350,1.428) (5.400,1.428) (5.450,1.427) (5.500,1.426) (5.550,1.426) (5.600,1.425) (5.650,1.425) (5.700,1.424) (5.750,1.424) (5.800,1.423) (5.850,1.423) (5.900,1.422) (5.950,1.422) (6.000,1.422) (6.000,-1.407) (5.950,-1.406) (5.900,-1.406) (5.850,-1.406) (5.800,-1.405) (5.750,-1.405) (5.700,-1.404) (5.650,-1.404) (5.600,-1.403) (5.550,-1.403) (5.500,-1.402) (5.450,-1.401) (5.400,-1.401) (5.350,-1.400) (5.300,-1.399) (5.250,-1.398) (5.200,-1.398) (5.150,-1.397) (5.100,-1.396) (5.050,-1.395) (5.000,-1.394) (4.950,-1.393) (4.900,-1.392) (4.850,-1.391) (4.800,-1.389) (4.750,-1.388) (4.700,-1.387) (4.650,-1.385) (4.600,-1.384) (4.550,-1.382) (4.500,-1.381) (4.450,-1.379) (4.400,-1.377) (4.350,-1.375) (4.300,-1.373) (4.250,-1.371) (4.200,-1.369) (4.150,-1.367) (4.100,-1.364) (4.050,-1.362) (4.000,-1.359) (3.950,-1.356) (3.900,-1.353) (3.850,-1.350) (3.800,-1.347) (3.750,-1.343) (3.700,-1.340) (3.650,-1.336) (3.600,-1.332) (3.550,-1.327) (3.500,-1.323) (3.450,-1.318) (3.400,-1.313) (3.350,-1.308) (3.300,-1.303) (3.250,-1.297) (3.200,-1.291) (3.150,-1.284) (3.100,-1.278) (3.050,-1.271) (3.000,-1.263) (2.950,-1.255) (2.900,-1.247) (2.850,-1.238) (2.800,-1.229) (2.750,-1.220) (2.700,-1.209) (2.650,-1.199) (2.600,-1.187) (2.550,-1.176) (2.500,-1.163) (2.450,-1.150) (2.400,-1.136) (2.350,-1.122) (2.300,-1.106) (2.250,-1.090) (2.200,-1.073) (2.150,-1.055) (2.100,-1.036) (2.050,-1.016) (2.000,-0.995) (1.950,-0.973) (1.900,-0.950) (1.850,-0.925) (1.800,-0.899) (1.750,-0.871) (1.700,-0.842) (1.650,-0.812) (1.600,-0.779) (1.550,-0.745) (1.500,-0.709) (1.450,-0.671) (1.400,-0.631) (1.350,-0.588) (1.300,-0.543) (1.250,-0.495) (1.200,-0.445) (1.150,-0.392) (1.100,-0.335) (1.050,-0.275) (1.000,-0.211) (0.950,-0.144) (0.900,-0.072) (0.850,0.004) (0.800,0.085) (0.750,0.171) (0.700,0.262) (0.650,0.360) (0.600,0.464) (0.550,0.576) (0.500,0.695) (0.450,0.823) (0.400,0.962) (0.350,1.111) (0.300,1.273) (0.250,1.449) (0.200,1.644) (0.150,1.862) (0.100,2.112) (0.050,2.417) (0.000,3.000)} -- cycle;
  \addplot[SADark!20] coordinates {(0,0) (6,0)};
  \addplot[SABlue,thin] coordinates {(0.000,3.000) (0.010,3.173) (0.020,2.887) (0.030,2.900) (0.040,2.815) (0.050,2.742) (0.060,2.693) (0.070,2.465) (0.080,2.418) (0.090,2.308) (0.100,2.615) (0.110,2.612) (0.120,2.551) (0.130,2.497) (0.140,2.406) (0.150,2.277) (0.160,2.215) (0.170,2.241) (0.180,2.195) (0.190,2.269) (0.200,2.226) (0.210,2.207) (0.220,2.338) (0.230,2.369) (0.240,2.296) (0.250,2.255) (0.260,2.286) (0.270,2.456) (0.280,2.404) (0.290,2.356) (0.300,2.433) (0.310,2.320) (0.320,2.268) (0.330,2.333) (0.340,2.368) (0.350,2.353) (0.360,2.397) (0.370,2.091) (0.380,2.172) (0.390,2.055) (0.400,1.869) (0.410,1.877) (0.420,1.929) (0.430,1.865) (0.440,1.739) (0.450,1.725) (0.460,1.702) (0.470,1.825) (0.480,1.881) (0.490,1.882) (0.500,1.974) (0.510,1.934) (0.520,1.822) (0.530,1.862) (0.540,1.902) (0.550,1.862) (0.560,1.765) (0.570,1.770) (0.580,1.505) (0.590,1.558) (0.600,1.592) (0.610,1.413) (0.620,1.405) (0.630,1.295) (0.640,1.357) (0.650,1.141) (0.660,1.039) (0.670,1.099) (0.680,1.203) (0.690,0.977) (0.700,0.918) (0.710,0.941) (0.720,0.871) (0.730,1.021) (0.740,0.892) (0.750,0.918) (0.760,0.805) (0.770,0.937) (0.780,0.925) (0.790,0.879) (0.800,0.699) (0.810,0.860) (0.820,0.926) (0.830,0.992) (0.840,1.095) (0.850,1.119) (0.860,1.044) (0.870,0.954) (0.880,0.865) (0.890,0.993) (0.900,0.838) (0.910,0.770) (0.920,0.730) (0.930,0.746) (0.940,0.795) (0.950,0.663) (0.960,0.484) (0.970,0.479) (0.980,0.595) (0.990,0.665) (1.000,0.679) (1.010,0.641) (1.020,0.664) (1.030,0.633) (1.040,0.708) (1.050,0.622) (1.060,0.629) (1.070,0.612) (1.080,0.660) (1.090,0.676) (1.100,0.923) (1.110,1.063) (1.120,1.201) (1.130,0.986) (1.140,0.942) (1.150,0.872) (1.160,0.917) (1.170,0.681) (1.180,0.791) (1.190,0.889) (1.200,0.751) (1.210,0.646) (1.220,0.560) (1.230,0.559) (1.240,0.617) (1.250,0.814) (1.260,0.787) (1.270,0.855) (1.280,0.862) (1.290,1.029) (1.300,1.092) (1.310,1.218) (1.320,1.098) (1.330,1.068) (1.340,0.977) (1.350,1.117) (1.360,1.171) (1.370,1.129) (1.380,1.073) (1.390,1.110) (1.400,1.029) (1.410,0.927) (1.420,0.965) (1.430,1.201) (1.440,1.164) (1.450,1.097) (1.460,0.970) (1.470,0.827) (1.480,0.871) (1.490,0.947) (1.500,0.939) (1.510,0.963) (1.520,0.965) (1.530,0.969) (1.540,0.807) (1.550,0.754) (1.560,0.757) (1.570,0.672) (1.580,0.618) (1.590,0.530) (1.600,0.492) (1.610,0.572) (1.620,0.526) (1.630,0.505) (1.640,0.581) (1.650,0.639) (1.660,0.802) (1.670,0.586) (1.680,0.665) (1.690,0.610) (1.700,0.618) (1.710,0.695) (1.720,0.796) (1.730,0.891) (1.740,0.898) (1.750,1.049) (1.760,1.011) (1.770,0.986) (1.780,1.056) (1.790,0.991) (1.800,1.196) (1.810,1.285) (1.820,1.489) (1.830,1.472) (1.840,1.419) (1.850,1.421) (1.860,1.480) (1.870,1.407) (1.880,1.431) (1.890,1.415) (1.900,1.562) (1.910,1.480) (1.920,1.570) (1.930,1.490) (1.940,1.380) (1.950,1.295) (1.960,1.164) (1.970,1.167) (1.980,1.258) (1.990,1.262) (2.000,1.311) (2.010,1.461) (2.020,1.473) (2.030,1.478) (2.040,1.436) (2.050,1.261) (2.060,1.324) (2.070,1.136) (2.080,1.190) (2.090,1.177) (2.100,1.277) (2.110,1.258) (2.120,1.163) (2.130,1.187) (2.140,1.120) (2.150,1.091) (2.160,1.084) (2.170,1.060) (2.180,1.030) (2.190,0.937) (2.200,0.909) (2.210,0.687) (2.220,0.665) (2.230,0.539) (2.240,0.645) (2.250,0.765) (2.260,0.563) (2.270,0.572) (2.280,0.554) (2.290,0.444) (2.300,0.493) (2.310,0.442) (2.320,0.262) (2.330,0.232) (2.340,0.215) (2.350,0.224) (2.360,0.099) (2.370,0.159) (2.380,0.231) (2.390,0.115) (2.400,0.048) (2.410,0.040) (2.420,-0.017) (2.430,0.156) (2.440,0.176) (2.450,0.073) (2.460,-0.006) (2.470,-0.011) (2.480,0.217) (2.490,0.197) (2.500,0.208) (2.510,0.257) (2.520,0.250) (2.530,0.475) (2.540,0.417) (2.550,0.487) (2.560,0.498) (2.570,0.555) (2.580,0.433) (2.590,0.604) (2.600,0.586) (2.610,0.584) (2.620,0.485) (2.630,0.407) (2.640,0.459) (2.650,0.553) (2.660,0.622) (2.670,0.737) (2.680,0.800) (2.690,0.795) (2.700,0.871) (2.710,0.921) (2.720,0.902) (2.730,0.965) (2.740,1.083) (2.750,1.096) (2.760,1.050) (2.770,1.111) (2.780,1.289) (2.790,1.255) (2.800,1.233) (2.810,1.207) (2.820,1.317) (2.830,1.121) (2.840,1.146) (2.850,1.254) (2.860,1.160) (2.870,1.297) (2.880,1.337) (2.890,1.269) (2.900,1.277) (2.910,1.112) (2.920,1.044) (2.930,1.216) (2.940,1.115) (2.950,1.287) (2.960,1.266) (2.970,1.353) (2.980,1.343) (2.990,1.107) (3.000,1.039) (3.010,1.048) (3.020,0.923) (3.030,0.779) (3.040,0.808) (3.050,0.735) (3.060,0.559) (3.070,0.486) (3.080,0.567) (3.090,0.514) (3.100,0.602) (3.110,0.748) (3.120,0.743) (3.130,0.624) (3.140,0.488) (3.150,0.480) (3.160,0.563) (3.170,0.585) (3.180,0.633) (3.190,0.583) (3.200,0.607) (3.210,0.554) (3.220,0.516) (3.230,0.382) (3.240,0.222) (3.250,0.172) (3.260,0.060) (3.270,-0.036) (3.280,0.031) (3.290,0.020) (3.300,0.310) (3.310,0.399) (3.320,0.332) (3.330,0.406) (3.340,0.438) (3.350,0.365) (3.360,0.471) (3.370,0.389) (3.380,0.081) (3.390,0.159) (3.400,0.096) (3.410,0.243) (3.420,0.176) (3.430,0.057) (3.440,0.210) (3.450,0.101) (3.460,0.120) (3.470,0.263) (3.480,0.272) (3.490,0.262) (3.500,0.269) (3.510,0.350) (3.520,0.266) (3.530,0.324) (3.540,0.378) (3.550,0.607) (3.560,0.567) (3.570,0.471) (3.580,0.539) (3.590,0.494) (3.600,0.518) (3.610,0.526) (3.620,0.637) (3.630,0.565) (3.640,0.551) (3.650,0.291) (3.660,0.200) (3.670,0.351) (3.680,0.279) (3.690,0.207) (3.700,0.084) (3.710,0.140) (3.720,0.157) (3.730,0.037) (3.740,0.103) (3.750,0.128) (3.760,0.202) (3.770,0.136) (3.780,0.307) (3.790,0.345) (3.800,0.330) (3.810,0.275) (3.820,0.248) (3.830,0.331) (3.840,0.301) (3.850,0.405) (3.860,0.459) (3.870,0.588) (3.880,0.664) (3.890,0.651) (3.900,0.454) (3.910,0.479) (3.920,0.553) (3.930,0.564) (3.940,0.614) (3.950,0.561) (3.960,0.517) (3.970,0.418) (3.980,0.550) (3.990,0.560) (4.000,0.505) (4.010,0.368) (4.020,0.349) (4.030,0.489) (4.040,0.530) (4.050,0.453) (4.060,0.590) (4.070,0.727) (4.080,0.817) (4.090,0.887) (4.100,0.994) (4.110,1.038) (4.120,1.105) (4.130,1.016) (4.140,1.155) (4.150,1.167) (4.160,1.356) (4.170,1.140) (4.180,1.198) (4.190,1.284) (4.200,1.237) (4.210,1.270) (4.220,1.332) (4.230,1.376) (4.240,1.309) (4.250,1.282) (4.260,1.370) (4.270,1.484) (4.280,1.479) (4.290,1.472) (4.300,1.477) (4.310,1.566) (4.320,1.445) (4.330,1.298) (4.340,1.298) (4.350,1.175) (4.360,1.105) (4.370,1.102) (4.380,1.139) (4.390,1.021) (4.400,1.091) (4.410,0.904) (4.420,0.949) (4.430,0.782) (4.440,0.648) (4.450,0.705) (4.460,0.658) (4.470,0.572) (4.480,0.522) (4.490,0.563) (4.500,0.552) (4.510,0.544) (4.520,0.643) (4.530,0.566) (4.540,0.513) (4.550,0.456) (4.560,0.568) (4.570,0.661) (4.580,0.782) (4.590,0.771) (4.600,0.748) (4.610,0.825) (4.620,0.778) (4.630,0.853) (4.640,0.916) (4.650,0.838) (4.660,0.821) (4.670,0.801) (4.680,0.966) (4.690,1.213) (4.700,1.191) (4.710,1.181) (4.720,0.942) (4.730,0.954) (4.740,0.944) (4.750,0.828) (4.760,0.898) (4.770,0.733) (4.780,0.704) (4.790,0.711) (4.800,0.580) (4.810,0.659) (4.820,0.758) (4.830,0.612) (4.840,0.466) (4.850,0.521) (4.860,0.513) (4.870,0.387) (4.880,0.354) (4.890,0.340) (4.900,0.322) (4.910,0.401) (4.920,0.436) (4.930,0.607) (4.940,0.738) (4.950,0.798) (4.960,0.897) (4.970,1.020) (4.980,1.080) (4.990,0.948) (5.000,1.024) (5.010,1.021) (5.020,1.172) (5.030,1.169) (5.040,1.210) (5.050,1.231) (5.060,1.386) (5.070,1.390) (5.080,1.349) (5.090,1.366) (5.100,1.309) (5.110,1.312) (5.120,1.298) (5.130,1.312) (5.140,1.491) (5.150,1.435) (5.160,1.425) (5.170,1.391) (5.180,1.444) (5.190,1.312) (5.200,1.230) (5.210,1.214) (5.220,1.256) (5.230,1.239) (5.240,1.177) (5.250,1.079) (5.260,1.119) (5.270,1.251) (5.280,1.271) (5.290,1.277) (5.300,1.309) (5.310,1.202) (5.320,1.159) (5.330,1.121) (5.340,1.139) (5.350,1.124) (5.360,0.992) (5.370,1.094) (5.380,1.252) (5.390,1.213) (5.400,1.241) (5.410,1.255) (5.420,1.321) (5.430,1.160) (5.440,1.309) (5.450,1.250) (5.460,1.363) (5.470,1.357) (5.480,1.305) (5.490,1.278) (5.500,1.355) (5.510,1.255) (5.520,1.206) (5.530,1.146) (5.540,0.963) (5.550,1.083) (5.560,1.130) (5.570,1.117) (5.580,1.221) (5.590,1.114) (5.600,1.096) (5.610,1.256) (5.620,1.341) (5.630,1.107) (5.640,0.764) (5.650,0.653) (5.660,0.752) (5.670,0.698) (5.680,0.669) (5.690,0.723) (5.700,0.628) (5.710,0.643) (5.720,0.715) (5.730,0.755) (5.740,0.786) (5.750,0.719) (5.760,0.763) (5.770,0.713) (5.780,0.835) (5.790,1.083) (5.800,0.998) (5.810,0.772) (5.820,0.722) (5.830,0.752) (5.840,0.849) (5.850,0.948) (5.860,0.861) (5.870,0.773) (5.880,0.892) (5.890,0.774) (5.900,0.875) (5.910,0.844) (5.920,0.912) (5.930,0.973) (5.940,1.040) (5.950,1.112) (5.960,1.002) (5.970,1.022) (5.980,1.013) (5.990,1.034) (6.000,1.146)};
  \addplot[SARed,thin] coordinates {(0.000,3.000) (0.010,2.943) (0.020,2.843) (0.030,2.742) (0.040,2.684) (0.050,2.732) (0.060,2.624) (0.070,2.713) (0.080,2.715) (0.090,2.666) (0.100,2.507) (0.110,2.506) (0.120,2.550) (0.130,2.349) (0.140,2.444) (0.150,2.584) (0.160,2.592) (0.170,2.615) (0.180,2.829) (0.190,2.923) (0.200,2.805) (0.210,2.775) (0.220,2.732) (0.230,2.690) (0.240,2.428) (0.250,2.479) (0.260,2.391) (0.270,2.384) (0.280,2.346) (0.290,2.410) (0.300,2.356) (0.310,2.254) (0.320,2.284) (0.330,2.401) (0.340,2.346) (0.350,2.459) (0.360,2.460) (0.370,2.383) (0.380,2.511) (0.390,2.404) (0.400,2.481) (0.410,2.407) (0.420,2.428) (0.430,2.281) (0.440,2.175) (0.450,2.259) (0.460,2.255) (0.470,2.237) (0.480,2.171) (0.490,2.194) (0.500,2.033) (0.510,1.901) (0.520,1.866) (0.530,1.765) (0.540,1.678) (0.550,1.654) (0.560,1.624) (0.570,1.585) (0.580,1.735) (0.590,1.819) (0.600,1.923) (0.610,1.901) (0.620,1.832) (0.630,1.669) (0.640,1.618) (0.650,1.394) (0.660,1.322) (0.670,1.385) (0.680,1.533) (0.690,1.613) (0.700,1.512) (0.710,1.546) (0.720,1.452) (0.730,1.519) (0.740,1.493) (0.750,1.563) (0.760,1.804) (0.770,1.878) (0.780,2.091) (0.790,2.088) (0.800,2.034) (0.810,2.147) (0.820,2.297) (0.830,2.219) (0.840,2.205) (0.850,2.040) (0.860,2.000) (0.870,2.068) (0.880,2.039) (0.890,1.962) (0.900,1.803) (0.910,1.821) (0.920,1.882) (0.930,1.923) (0.940,1.929) (0.950,1.783) (0.960,1.941) (0.970,1.916) (0.980,1.883) (0.990,1.937) (1.000,1.702) (1.010,1.679) (1.020,1.649) (1.030,1.669) (1.040,1.787) (1.050,1.816) (1.060,1.615) (1.070,1.566) (1.080,1.571) (1.090,1.570) (1.100,1.426) (1.110,1.247) (1.120,1.289) (1.130,1.310) (1.140,1.295) (1.150,1.111) (1.160,1.276) (1.170,0.975) (1.180,0.929) (1.190,1.128) (1.200,1.101) (1.210,1.113) (1.220,0.985) (1.230,1.110) (1.240,0.970) (1.250,0.986) (1.260,0.990) (1.270,0.875) (1.280,0.955) (1.290,0.804) (1.300,0.728) (1.310,0.718) (1.320,0.670) (1.330,0.710) (1.340,0.795) (1.350,0.758) (1.360,0.696) (1.370,0.638) (1.380,0.595) (1.390,0.670) (1.400,0.641) (1.410,0.676) (1.420,0.660) (1.430,0.569) (1.440,0.649) (1.450,0.540) (1.460,0.535) (1.470,0.678) (1.480,0.562) (1.490,0.551) (1.500,0.388) (1.510,0.333) (1.520,0.181) (1.530,0.143) (1.540,0.089) (1.550,0.101) (1.560,0.257) (1.570,0.283) (1.580,0.346) (1.590,0.337) (1.600,0.309) (1.610,0.223) (1.620,0.042) (1.630,0.048) (1.640,0.210) (1.650,0.255) (1.660,0.200) (1.670,0.362) (1.680,0.315) (1.690,0.250) (1.700,0.127) (1.710,0.030) (1.720,0.005) (1.730,0.204) (1.740,0.182) (1.750,0.077) (1.760,0.016) (1.770,-0.006) (1.780,-0.176) (1.790,-0.162) (1.800,-0.091) (1.810,-0.058) (1.820,0.003) (1.830,-0.126) (1.840,-0.294) (1.850,-0.163) (1.860,-0.052) (1.870,-0.205) (1.880,-0.182) (1.890,-0.178) (1.900,-0.122) (1.910,-0.100) (1.920,-0.149) (1.930,-0.128) (1.940,-0.111) (1.950,-0.115) (1.960,-0.298) (1.970,-0.593) (1.980,-0.666) (1.990,-0.822) (2.000,-0.789) (2.010,-0.710) (2.020,-0.676) (2.030,-0.706) (2.040,-0.692) (2.050,-0.569) (2.060,-0.636) (2.070,-0.538) (2.080,-0.462) (2.090,-0.432) (2.100,-0.283) (2.110,-0.351) (2.120,-0.254) (2.130,-0.344) (2.140,-0.417) (2.150,-0.518) (2.160,-0.469) (2.170,-0.558) (2.180,-0.547) (2.190,-0.654) (2.200,-0.542) (2.210,-0.354) (2.220,-0.328) (2.230,-0.161) (2.240,-0.104) (2.250,0.113) (2.260,0.122) (2.270,-0.035) (2.280,0.113) (2.290,0.100) (2.300,-0.004) (2.310,-0.070) (2.320,-0.231) (2.330,-0.245) (2.340,-0.349) (2.350,-0.198) (2.360,-0.130) (2.370,-0.294) (2.380,-0.441) (2.390,-0.493) (2.400,-0.680) (2.410,-0.684) (2.420,-0.682) (2.430,-0.752) (2.440,-0.799) (2.450,-0.741) (2.460,-0.592) (2.470,-0.438) (2.480,-0.430) (2.490,-0.340) (2.500,-0.410) (2.510,-0.365) (2.520,-0.457) (2.530,-0.380) (2.540,-0.500) (2.550,-0.591) (2.560,-0.586) (2.570,-0.524) (2.580,-0.662) (2.590,-0.709) (2.600,-0.688) (2.610,-0.659) (2.620,-0.736) (2.630,-0.719) (2.640,-0.738) (2.650,-0.591) (2.660,-0.530) (2.670,-0.656) (2.680,-0.869) (2.690,-0.850) (2.700,-0.724) (2.710,-0.700) (2.720,-0.605) (2.730,-0.647) (2.740,-0.776) (2.750,-0.694) (2.760,-0.678) (2.770,-0.603) (2.780,-0.734) (2.790,-0.709) (2.800,-0.764) (2.810,-0.676) (2.820,-0.727) (2.830,-0.780) (2.840,-0.831) (2.850,-0.697) (2.860,-0.687) (2.870,-0.756) (2.880,-0.850) (2.890,-0.946) (2.900,-0.791) (2.910,-0.963) (2.920,-1.038) (2.930,-0.965) (2.940,-0.838) (2.950,-0.844) (2.960,-0.758) (2.970,-0.674) (2.980,-0.570) (2.990,-0.559) (3.000,-0.389) (3.010,-0.364) (3.020,-0.206) (3.030,-0.312) (3.040,-0.213) (3.050,-0.334) (3.060,-0.477) (3.070,-0.420) (3.080,-0.496) (3.090,-0.443) (3.100,-0.542) (3.110,-0.374) (3.120,-0.365) (3.130,-0.471) (3.140,-0.445) (3.150,-0.348) (3.160,-0.417) (3.170,-0.357) (3.180,-0.145) (3.190,-0.053) (3.200,-0.048) (3.210,-0.099) (3.220,-0.036) (3.230,-0.181) (3.240,-0.476) (3.250,-0.533) (3.260,-0.630) (3.270,-0.444) (3.280,-0.421) (3.290,-0.509) (3.300,-0.614) (3.310,-0.552) (3.320,-0.623) (3.330,-0.669) (3.340,-0.665) (3.350,-0.644) (3.360,-0.621) (3.370,-0.597) (3.380,-0.455) (3.390,-0.275) (3.400,-0.430) (3.410,-0.162) (3.420,-0.141) (3.430,-0.009) (3.440,-0.075) (3.450,-0.035) (3.460,-0.006) (3.470,0.013) (3.480,0.105) (3.490,0.174) (3.500,0.080) (3.510,0.074) (3.520,0.154) (3.530,0.228) (3.540,0.200) (3.550,0.177) (3.560,0.168) (3.570,0.064) (3.580,-0.058) (3.590,0.107) (3.600,0.074) (3.610,0.090) (3.620,0.022) (3.630,0.159) (3.640,0.339) (3.650,0.321) (3.660,0.377) (3.670,0.305) (3.680,0.485) (3.690,0.580) (3.700,0.498) (3.710,0.625) (3.720,0.662) (3.730,0.704) (3.740,0.740) (3.750,0.674) (3.760,0.621) (3.770,0.684) (3.780,0.679) (3.790,0.585) (3.800,0.789) (3.810,0.767) (3.820,0.711) (3.830,0.642) (3.840,0.711) (3.850,0.808) (3.860,0.831) (3.870,0.916) (3.880,0.765) (3.890,0.727) (3.900,0.680) (3.910,0.680) (3.920,0.729) (3.930,0.525) (3.940,0.837) (3.950,0.749) (3.960,0.782) (3.970,0.863) (3.980,0.945) (3.990,0.966) (4.000,0.673) (4.010,0.564) (4.020,0.639) (4.030,0.546) (4.040,0.441) (4.050,0.207) (4.060,0.331) (4.070,0.383) (4.080,0.295) (4.090,0.283) (4.100,0.274) (4.110,0.164) (4.120,0.139) (4.130,0.166) (4.140,0.222) (4.150,0.349) (4.160,0.351) (4.170,0.304) (4.180,0.344) (4.190,0.280) (4.200,0.155) (4.210,0.270) (4.220,0.231) (4.230,0.243) (4.240,0.307) (4.250,0.196) (4.260,0.260) (4.270,0.304) (4.280,0.293) (4.290,0.332) (4.300,0.255) (4.310,0.035) (4.320,0.008) (4.330,-0.081) (4.340,0.008) (4.350,0.036) (4.360,-0.015) (4.370,-0.008) (4.380,0.128) (4.390,0.242) (4.400,0.278) (4.410,0.320) (4.420,0.291) (4.430,0.276) (4.440,0.334) (4.450,0.321) (4.460,0.349) (4.470,0.236) (4.480,0.069) (4.490,-0.099) (4.500,-0.140) (4.510,-0.005) (4.520,0.115) (4.530,0.016) (4.540,-0.138) (4.550,-0.056) (4.560,-0.007) (4.570,-0.142) (4.580,-0.034) (4.590,-0.015) (4.600,0.119) (4.610,0.153) (4.620,0.242) (4.630,0.273) (4.640,0.177) (4.650,0.154) (4.660,0.086) (4.670,0.169) (4.680,0.070) (4.690,-0.082) (4.700,-0.012) (4.710,0.104) (4.720,0.248) (4.730,0.356) (4.740,0.204) (4.750,0.133) (4.760,0.079) (4.770,0.077) (4.780,-0.137) (4.790,-0.135) (4.800,-0.105) (4.810,-0.198) (4.820,-0.183) (4.830,-0.061) (4.840,-0.055) (4.850,0.052) (4.860,0.154) (4.870,0.194) (4.880,0.231) (4.890,0.193) (4.900,0.316) (4.910,0.267) (4.920,0.196) (4.930,0.265) (4.940,0.192) (4.950,0.233) (4.960,0.337) (4.970,0.343) (4.980,0.351) (4.990,0.441) (5.000,0.512) (5.010,0.355) (5.020,0.223) (5.030,0.354) (5.040,0.391) (5.050,0.181) (5.060,0.047) (5.070,0.050) (5.080,-0.046) (5.090,0.003) (5.100,0.012) (5.110,0.062) (5.120,-0.053) (5.130,0.012) (5.140,0.208) (5.150,0.129) (5.160,0.260) (5.170,0.358) (5.180,0.420) (5.190,0.412) (5.200,0.294) (5.210,0.206) (5.220,0.282) (5.230,0.379) (5.240,0.499) (5.250,0.592) (5.260,0.678) (5.270,0.631) (5.280,0.506) (5.290,0.524) (5.300,0.571) (5.310,0.576) (5.320,0.422) (5.330,0.368) (5.340,0.220) (5.350,0.235) (5.360,0.130) (5.370,0.115) (5.380,0.148) (5.390,0.163) (5.400,0.300) (5.410,0.189) (5.420,-0.027) (5.430,0.070) (5.440,0.211) (5.450,0.354) (5.460,0.382) (5.470,0.434) (5.480,0.445) (5.490,0.314) (5.500,0.257) (5.510,0.307) (5.520,0.250) (5.530,0.440) (5.540,0.285) (5.550,0.242) (5.560,0.353) (5.570,0.292) (5.580,0.180) (5.590,0.046) (5.600,-0.024) (5.610,0.129) (5.620,0.172) (5.630,0.225) (5.640,0.174) (5.650,0.259) (5.660,0.271) (5.670,0.180) (5.680,0.136) (5.690,0.158) (5.700,0.276) (5.710,0.271) (5.720,0.321) (5.730,0.522) (5.740,0.441) (5.750,0.316) (5.760,0.415) (5.770,0.457) (5.780,0.446) (5.790,0.472) (5.800,0.438) (5.810,0.488) (5.820,0.453) (5.830,0.532) (5.840,0.413) (5.850,0.332) (5.860,0.343) (5.870,0.297) (5.880,0.366) (5.890,0.362) (5.900,0.308) (5.910,0.330) (5.920,0.277) (5.930,0.123) (5.940,0.013) (5.950,-0.063) (5.960,-0.106) (5.970,-0.091) (5.980,0.004) (5.990,0.063) (6.000,0.244)};
  \addplot[SAGreenText,thin] coordinates {(0.000,3.000) (0.010,3.101) (0.020,3.056) (0.030,3.194) (0.040,3.284) (0.050,3.210) (0.060,3.214) (0.070,3.295) (0.080,3.291) (0.090,3.364) (0.100,3.374) (0.110,3.415) (0.120,3.354) (0.130,3.363) (0.140,3.298) (0.150,3.337) (0.160,3.140) (0.170,3.147) (0.180,3.245) (0.190,3.275) (0.200,3.245) (0.210,3.271) (0.220,3.159) (0.230,3.049) (0.240,2.988) (0.250,2.948) (0.260,3.020) (0.270,2.905) (0.280,2.981) (0.290,2.799) (0.300,2.698) (0.310,2.672) (0.320,2.765) (0.330,2.681) (0.340,2.446) (0.350,2.468) (0.360,2.445) (0.370,2.318) (0.380,2.195) (0.390,2.109) (0.400,2.102) (0.410,2.152) (0.420,1.988) (0.430,2.201) (0.440,2.180) (0.450,2.125) (0.460,2.176) (0.470,2.170) (0.480,2.385) (0.490,2.475) (0.500,2.418) (0.510,2.360) (0.520,2.291) (0.530,2.239) (0.540,2.288) (0.550,2.205) (0.560,2.348) (0.570,2.321) (0.580,2.288) (0.590,2.364) (0.600,2.148) (0.610,2.038) (0.620,2.081) (0.630,2.066) (0.640,2.048) (0.650,2.072) (0.660,1.996) (0.670,1.963) (0.680,1.812) (0.690,1.612) (0.700,1.726) (0.710,1.620) (0.720,1.537) (0.730,1.664) (0.740,1.757) (0.750,1.692) (0.760,1.626) (0.770,1.606) (0.780,1.499) (0.790,1.542) (0.800,1.545) (0.810,1.547) (0.820,1.610) (0.830,1.670) (0.840,1.840) (0.850,1.977) (0.860,1.791) (0.870,1.866) (0.880,1.740) (0.890,1.584) (0.900,1.643) (0.910,1.561) (0.920,1.518) (0.930,1.613) (0.940,1.535) (0.950,1.434) (0.960,1.456) (0.970,1.567) (0.980,1.608) (0.990,1.441) (1.000,1.366) (1.010,1.341) (1.020,1.524) (1.030,1.601) (1.040,1.674) (1.050,1.781) (1.060,1.736) (1.070,1.770) (1.080,1.805) (1.090,1.694) (1.100,1.679) (1.110,1.723) (1.120,1.602) (1.130,1.670) (1.140,1.733) (1.150,1.822) (1.160,1.731) (1.170,1.728) (1.180,1.859) (1.190,1.924) (1.200,1.830) (1.210,1.838) (1.220,1.965) (1.230,1.941) (1.240,1.964) (1.250,2.203) (1.260,2.181) (1.270,2.153) (1.280,2.035) (1.290,1.953) (1.300,2.085) (1.310,2.053) (1.320,1.918) (1.330,1.810) (1.340,1.796) (1.350,1.857) (1.360,1.768) (1.370,1.825) (1.380,1.761) (1.390,1.782) (1.400,1.818) (1.410,1.863) (1.420,1.804) (1.430,1.772) (1.440,1.807) (1.450,1.662) (1.460,1.665) (1.470,1.701) (1.480,1.789) (1.490,1.769) (1.500,1.734) (1.510,1.732) (1.520,1.643) (1.530,1.607) (1.540,1.632) (1.550,1.645) (1.560,1.525) (1.570,1.478) (1.580,1.253) (1.590,1.320) (1.600,1.254) (1.610,1.265) (1.620,1.143) (1.630,1.123) (1.640,1.155) (1.650,1.241) (1.660,1.092) (1.670,1.164) (1.680,1.135) (1.690,1.045) (1.700,1.081) (1.710,1.039) (1.720,0.998) (1.730,0.917) (1.740,0.928) (1.750,0.921) (1.760,0.893) (1.770,0.787) (1.780,0.726) (1.790,0.757) (1.800,0.875) (1.810,0.862) (1.820,0.992) (1.830,0.979) (1.840,0.859) (1.850,0.911) (1.860,0.804) (1.870,0.834) (1.880,0.930) (1.890,0.888) (1.900,0.869) (1.910,1.104) (1.920,1.114) (1.930,1.128) (1.940,1.192) (1.950,1.169) (1.960,1.146) (1.970,1.171) (1.980,0.920) (1.990,1.049) (2.000,0.931) (2.010,0.813) (2.020,0.856) (2.030,0.849) (2.040,0.714) (2.050,0.726) (2.060,0.604) (2.070,0.527) (2.080,0.371) (2.090,0.283) (2.100,0.318) (2.110,0.325) (2.120,0.351) (2.130,0.437) (2.140,0.427) (2.150,0.355) (2.160,0.322) (2.170,0.481) (2.180,0.478) (2.190,0.503) (2.200,0.487) (2.210,0.504) (2.220,0.540) (2.230,0.532) (2.240,0.510) (2.250,0.425) (2.260,0.335) (2.270,0.371) (2.280,0.439) (2.290,0.525) (2.300,0.523) (2.310,0.434) (2.320,0.491) (2.330,0.545) (2.340,0.389) (2.350,0.463) (2.360,0.514) (2.370,0.673) (2.380,0.552) (2.390,0.418) (2.400,0.490) (2.410,0.305) (2.420,0.228) (2.430,0.116) (2.440,0.012) (2.450,-0.179) (2.460,-0.202) (2.470,-0.105) (2.480,-0.061) (2.490,-0.018) (2.500,-0.055) (2.510,0.088) (2.520,0.138) (2.530,0.240) (2.540,0.082) (2.550,0.115) (2.560,-0.001) (2.570,0.033) (2.580,0.051) (2.590,-0.062) (2.600,-0.045) (2.610,0.070) (2.620,0.118) (2.630,0.175) (2.640,0.088) (2.650,0.019) (2.660,-0.092) (2.670,-0.078) (2.680,0.040) (2.690,0.083) (2.700,0.082) (2.710,-0.030) (2.720,0.014) (2.730,-0.010) (2.740,-0.198) (2.750,-0.368) (2.760,-0.398) (2.770,-0.482) (2.780,-0.314) (2.790,-0.307) (2.800,-0.407) (2.810,-0.451) (2.820,-0.288) (2.830,-0.379) (2.840,-0.299) (2.850,-0.252) (2.860,-0.104) (2.870,-0.050) (2.880,-0.089) (2.890,-0.036) (2.900,-0.028) (2.910,0.069) (2.920,0.063) (2.930,0.259) (2.940,0.247) (2.950,0.288) (2.960,0.221) (2.970,0.232) (2.980,0.234) (2.990,0.249) (3.000,0.337) (3.010,0.281) (3.020,0.196) (3.030,0.305) (3.040,0.275) (3.050,0.207) (3.060,0.320) (3.070,0.396) (3.080,0.131) (3.090,0.094) (3.100,0.172) (3.110,0.255) (3.120,0.203) (3.130,0.357) (3.140,0.277) (3.150,0.253) (3.160,0.265) (3.170,0.264) (3.180,0.343) (3.190,0.443) (3.200,0.436) (3.210,0.572) (3.220,0.501) (3.230,0.477) (3.240,0.645) (3.250,0.651) (3.260,0.753) (3.270,0.705) (3.280,0.674) (3.290,0.660) (3.300,0.673) (3.310,0.636) (3.320,0.634) (3.330,0.697) (3.340,0.812) (3.350,0.818) (3.360,0.877) (3.370,0.962) (3.380,0.849) (3.390,0.688) (3.400,0.844) (3.410,0.970) (3.420,0.969) (3.430,0.998) (3.440,0.898) (3.450,0.755) (3.460,0.782) (3.470,0.631) (3.480,0.718) (3.490,0.647) (3.500,0.673) (3.510,0.861) (3.520,1.042) (3.530,1.008) (3.540,1.107) (3.550,0.997) (3.560,0.932) (3.570,0.963) (3.580,1.093) (3.590,1.029) (3.600,1.075) (3.610,1.077) (3.620,0.944) (3.630,0.924) (3.640,1.067) (3.650,1.050) (3.660,1.092) (3.670,1.166) (3.680,1.084) (3.690,1.198) (3.700,1.338) (3.710,1.485) (3.720,1.603) (3.730,1.544) (3.740,1.473) (3.750,1.383) (3.760,1.282) (3.770,1.193) (3.780,1.281) (3.790,1.416) (3.800,1.341) (3.810,1.410) (3.820,1.273) (3.830,1.230) (3.840,1.121) (3.850,0.950) (3.860,0.948) (3.870,0.948) (3.880,0.900) (3.890,0.841) (3.900,0.736) (3.910,0.717) (3.920,0.652) (3.930,0.671) (3.940,0.546) (3.950,0.564) (3.960,0.611) (3.970,0.729) (3.980,0.828) (3.990,0.839) (4.000,0.764) (4.010,0.576) (4.020,0.570) (4.030,0.706) (4.040,0.865) (4.050,0.604) (4.060,0.524) (4.070,0.485) (4.080,0.377) (4.090,0.393) (4.100,0.394) (4.110,0.339) (4.120,0.333) (4.130,0.172) (4.140,0.074) (4.150,-0.015) (4.160,0.075) (4.170,-0.087) (4.180,-0.106) (4.190,-0.027) (4.200,-0.032) (4.210,0.050) (4.220,-0.015) (4.230,0.029) (4.240,0.119) (4.250,0.341) (4.260,0.410) (4.270,0.549) (4.280,0.545) (4.290,0.551) (4.300,0.702) (4.310,0.754) (4.320,0.889) (4.330,0.952) (4.340,1.154) (4.350,1.381) (4.360,1.286) (4.370,1.459) (4.380,1.490) (4.390,1.524) (4.400,1.499) (4.410,1.557) (4.420,1.433) (4.430,1.349) (4.440,1.261) (4.450,1.280) (4.460,1.361) (4.470,1.423) (4.480,1.448) (4.490,1.377) (4.500,1.191) (4.510,1.222) (4.520,1.227) (4.530,1.254) (4.540,1.386) (4.550,1.357) (4.560,1.326) (4.570,1.086) (4.580,1.133) (4.590,1.120) (4.600,1.038) (4.610,1.069) (4.620,0.958) (4.630,0.897) (4.640,1.087) (4.650,1.153) (4.660,1.237) (4.670,1.302) (4.680,1.356) (4.690,1.364) (4.700,1.581) (4.710,1.561) (4.720,1.621) (4.730,1.586) (4.740,1.537) (4.750,1.282) (4.760,1.080) (4.770,1.226) (4.780,1.364) (4.790,1.309) (4.800,1.295) (4.810,1.345) (4.820,1.235) (4.830,1.401) (4.840,1.285) (4.850,1.203) (4.860,0.917) (4.870,1.061) (4.880,1.058) (4.890,1.210) (4.900,1.100) (4.910,1.124) (4.920,1.122) (4.930,1.132) (4.940,1.329) (4.950,1.543) (4.960,1.666) (4.970,1.611) (4.980,1.685) (4.990,1.804) (5.000,1.842) (5.010,1.882) (5.020,1.913) (5.030,1.805) (5.040,1.813) (5.050,1.955) (5.060,1.838) (5.070,1.829) (5.080,1.895) (5.090,1.741) (5.100,1.711) (5.110,1.527) (5.120,1.452) (5.130,1.298) (5.140,1.197) (5.150,1.137) (5.160,1.099) (5.170,1.068) (5.180,0.819) (5.190,0.969) (5.200,0.933) (5.210,0.815) (5.220,0.744) (5.230,0.768) (5.240,0.660) (5.250,0.613) (5.260,0.560) (5.270,0.608) (5.280,0.527) (5.290,0.567) (5.300,0.596) (5.310,0.576) (5.320,0.693) (5.330,0.689) (5.340,0.565) (5.350,0.589) (5.360,0.595) (5.370,0.585) (5.380,0.602) (5.390,0.830) (5.400,0.845) (5.410,0.882) (5.420,0.849) (5.430,0.971) (5.440,0.956) (5.450,0.866) (5.460,0.736) (5.470,0.781) (5.480,0.847) (5.490,0.681) (5.500,0.602) (5.510,0.660) (5.520,0.756) (5.530,0.617) (5.540,0.570) (5.550,0.647) (5.560,0.587) (5.570,0.536) (5.580,0.477) (5.590,0.342) (5.600,0.343) (5.610,0.370) (5.620,0.227) (5.630,0.234) (5.640,0.342) (5.650,0.359) (5.660,0.450) (5.670,0.390) (5.680,0.522) (5.690,0.439) (5.700,0.622) (5.710,0.571) (5.720,0.493) (5.730,0.399) (5.740,0.354) (5.750,0.301) (5.760,0.229) (5.770,0.194) (5.780,0.258) (5.790,0.343) (5.800,0.306) (5.810,0.425) (5.820,0.322) (5.830,0.216) (5.840,0.101) (5.850,0.159) (5.860,0.003) (5.870,-0.136) (5.880,-0.008) (5.890,0.162) (5.900,0.294) (5.910,0.142) (5.920,0.141) (5.930,0.222) (5.940,0.229) (5.950,0.065) (5.960,0.030) (5.970,-0.140) (5.980,-0.262) (5.990,-0.046) (6.000,0.021)};
  \addplot[SADark,thick,dashed,domain=0:6,samples=80] {3*exp(-x)};
\end{axis}
\end{tikzpicture}
\caption{Ornstein-Uhlenbeck 过程 $\dd X_t=-X_t\dd t+\dd B_t$ ($\kappa=\sigma=1$, $\theta=0$) 从 $X_0=3$ 出发的三条轨道. 虚线为均值 $3e^{-t}$, 阴影为均值 $\pm2$ 倍标准差的范围. 均值指数衰减, 方差增长到平稳值 $\frac12$}
\label{fig:ou-paths}
\end{figure}

% ---------------------------------------------------------------------
\subsection{可约化为布朗运动的方程}

若扩散系数 $\sigma(x)>0$ 不依赖于时间且足够光滑, 适当的变量替换可以把扩散系数化为 $1$. 设 $F$ 满足 $F'(x)=\frac{1}{\sigma(x)}$. 由 It\^o 公式, $Y_t=F(X_t)$ 满足
\begin{equation}\label{eq:lamperti}
\dd Y_t=F'(X_t)\dd X_t+\frac{1}{2}F''(X_t)\sigma(X_t)^2\dd t=\Big(\frac{b(X_t)}{\sigma(X_t)}-\frac{1}{2}\sigma'(X_t)\Big)\dd t+\dd B_t,
\end{equation}
因为 $F''=-\frac{\sigma'}{\sigma^2}$. 这称为 \textbf{Lamperti 变换}. 若 \eqref{eq:lamperti} 中的漂移为零, 则 $Y$ 是 (平移的) 布朗运动, 原方程的解为 $X_t=F^{-1}(F(x)+B_t)$.

\begin{example}\label{ex:lamperti}
方程 $\dd X_t=\frac{1}{2}X_t\dd t+\sqrt{1+X_t^2}\dd B_t$: 取 $F(x)=\operatorname{arsinh}x$, 则 $F'=\frac{1}{\sqrt{1+x^2}}$, 漂移 $\frac{b}{\sigma}-\frac{1}{2}\sigma'=\frac{x/2}{\sqrt{1+x^2}}-\frac{1}{2}\cdot\frac{x}{\sqrt{1+x^2}}=0$. 所以 $X_t=\sinh\big(\operatorname{arsinh}x+B_t\big)$. 可以用 It\^o 公式直接验证.
\end{example}

% =====================================================================
\section{存在性与唯一性}
\label{sec:6-2}

常微分方程 $\dot{x}=b(x)$ 的经典理论告诉我们: 若 $b$ 满足 Lipschitz 条件, 则解存在唯一, 证明方法是 Picard 迭代. 随机微分方程的情形完全类似, 只是距离要用 $\E\sup_t\abs{\cdot}^2$ 来度量, 随机积分项要用 Doob 不等式与 It\^o 等距来估计.

% ---------------------------------------------------------------------
\subsection{Lipschitz 条件下的强解}

本节假设存在常数 $K$, 使得对一切 $t\in[0,T]$ 与 $x,y\in\R$,
\begin{align}
&\abs{b(t,x)-b(t,y)}+\abs{\sigma(t,x)-\sigma(t,y)}\le K\abs{x-y},\tag{H1}\label{eq:H1}\\
&\abs{b(t,x)}+\abs{\sigma(t,x)}\le K(1+\abs{x}).\tag{H2}\label{eq:H2}
\end{align}
\eqref{eq:H1} 是关于 $x$ 的一致 Lipschitz 条件, \eqref{eq:H2} 是线性增长条件 (若 $b(t,0)$ 与 $\sigma(t,0)$ 有界, 它由 \eqref{eq:H1} 推出). 记 $\mathscr{C}^2_T$ 为满足 $\E\sup_{t\le T}\abs{X_t}^2<\infty$ 的连续适应过程全体.

\begin{lemma}[Gronwall 不等式]\label{lem:gronwall}
设 $u$ 是 $[0,T]$ 上的非负有界 Borel 函数, $a,c\ge0$, 且 $u(t)\le a+c\int_0^tu(s)\dd s$ 对一切 $t\in[0,T]$ 成立. 则 $u(t)\le ae^{ct}$.
\end{lemma}

\begin{proof}
对 $n$ 归纳证明
\[
u(t)\le a\sum_{k=0}^{n-1}\frac{(ct)^k}{k!}+c^n\int_0^t\frac{(t-s)^{n-1}}{(n-1)!}u(s)\dd s.
\]
$n=1$ 即假设. 设它对 $n$ 成立, 把 $u(s)\le a+c\int_0^su(r)\dd r$ 代入最后一项, 并由 Fubini 定理交换积分次序,
\begin{align*}
c^n\int_0^t\frac{(t-s)^{n-1}}{(n-1)!}u(s)\dd s&\le a\frac{(ct)^n}{n!}+c^{n+1}\int_0^tu(r)\int_r^t\frac{(t-s)^{n-1}}{(n-1)!}\dd s\dd r\\
&=a\frac{(ct)^n}{n!}+c^{n+1}\int_0^t\frac{(t-r)^n}{n!}u(r)\dd r.
\end{align*}
归纳完成. 最后一项不超过 $\sup u\cdot\frac{(ct)^n}{n!}\to0$, 令 $n\to\infty$ 即得.
\end{proof}

对连续适应过程 $X$ 与 $\mathcal{F}_0$-可测的 $\xi$, 记
\[
\Psi(\xi,X)_t=\xi+\int_0^tb(s,X_s)\dd s+\int_0^t\sigma(s,X_s)\dd B_s,\qquad t\in[0,T].
\]
方程 \eqref{eq:sde} 的强解就是映射 $X\mapsto\Psi(\xi,X)$ 的不动点. 下面的估计是全部证明的核心.

\begin{lemma}[基本估计]\label{lem:sde-estimate}
设 \eqref{eq:H1} 成立, $X,Y$ 是循序可测过程, $\E\int_0^T(X_s^2+Y_s^2)\dd s<\infty$, $\xi,\eta\in L^2(\mathcal{F}_0)$, 且 $b(\cdot,X),\sigma(\cdot,X)$ 等满足 $\E\int_0^T(b^2+\sigma^2)\dd s<\infty$ (在 \eqref{eq:H2} 下自动成立). 则对 $t\le T$,
\[
\E\Big[\sup_{s\le t}\abs{\Psi(\xi,X)_s-\Psi(\eta,Y)_s}^2\Big]\le3\E\abs{\xi-\eta}^2+L\int_0^t\E\abs{X_s-Y_s}^2\dd s,
\]
其中 $L=3K^2(T+4)$.
\end{lemma}

\begin{proof}
记 $\Delta b_s=b(s,X_s)-b(s,Y_s)$, $\Delta\sigma_s=\sigma(s,X_s)-\sigma(s,Y_s)$. 由 $(a_1+a_2+a_3)^2\le3(a_1^2+a_2^2+a_3^2)$,
\[
\sup_{s\le t}\abs{\Psi(\xi,X)_s-\Psi(\eta,Y)_s}^2\le3\abs{\xi-\eta}^2+3\sup_{s\le t}\Big\lvert\int_0^s\Delta b_u\dd u\Big\rvert^2+3\sup_{s\le t}\Big\lvert\int_0^s\Delta\sigma_u\dd B_u\Big\rvert^2.
\]
对第二项, 由 Cauchy-Schwarz 不等式与 \eqref{eq:H1}, $\big(\int_0^t\abs{\Delta b_u}\dd u\big)^2\le t\int_0^t\abs{\Delta b_u}^2\dd u\le TK^2\int_0^t\abs{X_u-Y_u}^2\dd u$. 对第三项, $\Delta\sigma\in\mathscr{L}^2_T$, 由 Doob 不等式~\ref{thm:ct-doob} 与 It\^o 等距,
\[
\E\sup_{s\le t}\Big\lvert\int_0^s\Delta\sigma_u\dd B_u\Big\rvert^2\le4\E\int_0^t\abs{\Delta\sigma_u}^2\dd u\le4K^2\int_0^t\E\abs{X_u-Y_u}^2\dd u.
\]
取期望并相加即得.
\end{proof}

\begin{theorem}[强解的存在唯一性]\label{thm:sde-existence}
设 \eqref{eq:H1}\eqref{eq:H2} 成立, $\xi\in L^2(\mathcal{F}_0)$. 则方程 \eqref{eq:sde} 有强解 $X\in\mathscr{C}^2_T$, 并且有轨道唯一性; 同时存在只依赖于 $K,T$ 的常数 $C$, 使
\begin{equation}\label{eq:sde-moment}
\E\Big[\sup_{t\le T}\abs{X_t}^2\Big]\le C\big(1+\E\xi^2\big).
\end{equation}
\end{theorem}

\begin{proof}
\emph{唯一性.} 设 $X$, $Y$ 都是强解 (不必属于 $\mathscr{C}^2_T$). 对 $N\ge1$, 令 $\tau_N=\inf\{t:\abs{X_t}\vee\abs{Y_t}\ge N\}\wedge T$. 由停止性质 (定理~\ref{thm:ito-local}),
\[
X_{t\wedge\tau_N}=\xi+\int_0^t\1_{\{s\le\tau_N\}}b(s,X_s)\dd s+\int_0^t\1_{\{s\le\tau_N\}}\sigma(s,X_s)\dd B_s,
\]
对 $Y$ 同理. 在 $\{s\le\tau_N\}$ 上 $X_s=X_{s\wedge\tau_N}$, 所以 $\1_{\{s\le\tau_N\}}\abs{b(s,X_s)-b(s,Y_s)}\le K\abs{X_{s\wedge\tau_N}-Y_{s\wedge\tau_N}}$, 对 $\sigma$ 同理; 并且在 $(0,\tau_N]$ 上 $\abs{X},\abs{Y}\le N$, 所以被积过程有界. 与引理~\ref{lem:sde-estimate} 的证明相同 (此处初值相同), 函数 $e(t)=\E\sup_{s\le t}\abs{X_{s\wedge\tau_N}-Y_{s\wedge\tau_N}}^2$ 满足 $e(t)\le L\int_0^te(s)\dd s$. $e\le4N^2$ 有界 (若 $\abs{\xi}\ge N$, 则 $\tau_N=0$, 差为零; 否则停止后的两个过程都被 $N$ 控制), 由 Gronwall 不等式 $e\equiv0$. 所以 $X=Y$ 在 $[0,\tau_N]$ 上成立. 由轨道连续性, $\tau_N\uparrow T$ 且最终等于 $T$, 所以 $X$ 与 $Y$ 不可区分.

\emph{存在性.} 令 $X^0_t\equiv\xi$, $X^{k+1}=\Psi(\xi,X^k)$. 若 $X^k\in\mathscr{C}^2_T$, 则由 \eqref{eq:H2}, $b(\cdot,X^k)^2+\sigma(\cdot,X^k)^2\le2K^2(1+\abs{X^k}^2)$ 可积, 与引理~\ref{lem:sde-estimate} 的证明相同可得 $\E\sup_t\abs{X^{k+1}_t}^2<\infty$, 所以每个 $X^k\in\mathscr{C}^2_T$. 记 $D_k(t)=\E\sup_{s\le t}\abs{X^{k+1}_s-X^k_s}^2$. 由引理~\ref{lem:sde-estimate} (初值相同),
\[
D_k(t)\le L\int_0^t\E\abs{X^k_s-X^{k-1}_s}^2\dd s\le L\int_0^tD_{k-1}(s)\dd s.
\]
归纳得 $D_k(t)\le D_0(T)\frac{(Lt)^k}{k!}$. 记 $\norm{Z}=\big(\E\sup_{t\le T}\abs{Z_t}^2\big)^{1/2}$, 则 $\sum_k\norm{X^{k+1}-X^k}\le\sum_k\big(D_0(T)\frac{(LT)^k}{k!}\big)^{1/2}<\infty$. 于是 $\E\sum_k\sup_t\abs{X^{k+1}_t-X^k_t}<\infty$, 几乎必然地 $X^k$ 在 $[0,T]$ 上一致收敛, 极限 $X$ 连续、适应 (在零概率集上令 $X\equiv0$). 由 Fatou 引理与 Minkowski 不等式, $\norm{X-X^k}\le\sum_{j\ge k}\norm{X^{j+1}-X^j}\to0$, 特别地 $X\in\mathscr{C}^2_T$. 再由引理~\ref{lem:sde-estimate},
\[
\norm{\Psi(\xi,X)-X^{k+1}}^2=\norm{\Psi(\xi,X)-\Psi(\xi,X^k)}^2\le LT\norm{X-X^k}^2\to0,
\]
而 $X^{k+1}\to X$, 所以 $\norm{\Psi(\xi,X)-X}=0$: $X$ 与 $\Psi(\xi,X)$ 都连续, 二者不可区分, 即 $X$ 是强解.

\emph{矩估计.} 记 $m(t)=\E\sup_{s\le t}\abs{X_s}^2<\infty$. 与引理~\ref{lem:sde-estimate} 的证明相同, 用 \eqref{eq:H2} 代替 \eqref{eq:H1} ($\abs{b}^2,\abs{\sigma}^2\le2K^2(1+\abs{x}^2)$),
\begin{align*}
m(t)&\le3\E\xi^2+6K^2(T+4)\int_0^t\big(1+m(s)\big)\dd s\\
&\le3\E\xi^2+6K^2(T+4)T+6K^2(T+4)\int_0^tm(s)\dd s.
\end{align*}
由 Gronwall 不等式, $m(T)\le\big(3\E\xi^2+6K^2(T+4)T\big)e^{6K^2(T+4)T}$, 即得 \eqref{eq:sde-moment}.
\end{proof}

同样的论证给出解对初值的连续依赖.

\begin{proposition}[对初值的连续依赖]\label{prop:sde-continuity}
设 \eqref{eq:H1}\eqref{eq:H2} 成立, $X^\xi$, $X^\eta$ 分别是初值为 $\xi,\eta\in L^2(\mathcal{F}_0)$ 的解. 则
\[
\E\Big[\sup_{t\le T}\abs{X^\xi_t-X^\eta_t}^2\Big]\le3e^{LT}\,\E\abs{\xi-\eta}^2.
\]
\end{proposition}

\begin{proof}
由引理~\ref{lem:sde-estimate}, $u(t)=\E\sup_{s\le t}\abs{X^\xi_s-X^\eta_s}^2$ 满足 $u(t)\le3\E\abs{\xi-\eta}^2+L\int_0^tu(s)\dd s$, 再用 Gronwall 不等式.
\end{proof}

\begin{remark}
(1) \emph{初值.} 条件 $\xi\in L^2$ 可以去掉: 对一般的 $\mathcal{F}_0$-可测的 $\xi$, 在 $\{\abs{\xi}\le n\}\in\mathcal{F}_0$ 上用初值 $\xi\1_{\{\abs{\xi}\le n\}}$ 的解, 由局部性 (引理~\ref{lem:locality}) 与唯一性, 这些解可以拼接起来.

(2) \emph{Lipschitz 条件的作用.} 若只有局部 Lipschitz 条件, 解在 ``爆炸时刻'' 之前存在唯一. 例如 $\dot{x}=x^2$, $x_0=1$ 的解 $\frac{1}{1-t}$ 在 $t=1$ 爆炸, 线性增长条件 \eqref{eq:H2} 正是用来排除这种情况的. 若连局部 Lipschitz 条件也不满足, 唯一性可能失效: $\dot{x}=3x^{2/3}$, $x_0=0$ 有解 $x\equiv0$ 与 $x=t^3$.

(3) \emph{Yamada-Watanabe 条件.} 对一维方程, Lipschitz 条件对 $\sigma$ 可以放宽: 若 $b$ 满足 Lipschitz 条件, 而 $\abs{\sigma(t,x)-\sigma(t,y)}\le K\abs{x-y}^{1/2}$, 则轨道唯一性仍成立 (见~\cite{karatzas} 命题 5.2.13). 这适用于金融中的 Cox-Ingersoll-Ross 利率模型 $\dd X_t=\kappa(\theta-X_t)\dd t+\sigma\sqrt{X_t}\dd B_t$, 其扩散系数 $\sigma\sqrt{x}$ 在 $0$ 附近不是 Lipschitz 的.

(4) \emph{解是噪声的泛函.} 由构造, 每个 Picard 迭代 $X^k$ 是 $(\xi,B)$ 的 (同一个) 可测泛函, 所以解 $X$ 也是. 特别地, $X$ 的分布由 $\xi$ 的分布与 $(b,\sigma)$ 唯一决定, 与概率空间的选取无关.
\end{remark}

% ---------------------------------------------------------------------
\subsection{弱解}

弱解的概念允许我们处理系数不满足 Lipschitz 条件的方程. Girsanov 定理提供了构造弱解的一般方法: 先取一个布朗运动作为 ``解'', 再改变测度, 使它获得需要的漂移.

\begin{proposition}[用 Girsanov 定理构造弱解]\label{prop:weak-girsanov}
设 $b:[0,T]\times\R\to\R$ 是有界 Borel 函数, $x\in\R$. 则方程
\[
\dd X_t=b(t,X_t)\dd t+\dd W_t,\qquad X_0=x
\]
有弱解.
\end{proposition}

\begin{proof}
在某个概率空间上取标准布朗运动 $B$ 及其扩充信息流, 令 $X_t=x+B_t$. 过程 $\theta_t=b(t,X_t)$ 有界且循序可测, 由命题~\ref{prop:novikov} (1), $Z_t=\exp(\int_0^t\theta\dd B-\frac{1}{2}\int_0^t\theta^2\dd s)$ 在 $[0,T]$ 上是鞅. 令 $\widetilde{\P}=Z_T\cdot\P$. 由 Girsanov 定理~\ref{thm:girsanov}, $W_t=B_t-\int_0^tb(s,X_s)\dd s$ 是 $\widetilde{\P}$ 下的布朗运动, 而
\[
X_t=x+B_t=x+\int_0^tb(s,X_s)\dd s+W_t.
\]
所以 $(\Omega,\mathcal{F},\widetilde{\P},\{\mathcal{F}_t\},W,X)$ 是一个弱解 ($\widetilde{\P}$ 与 $\P$ 等价, 信息流在 $\widetilde{\P}$ 下仍满足通常条件).
\end{proof}

这里 $b$ 可以是不连续的, 例如 $b(x)=-\sgn(x)$ (总是把过程拉向原点的 ``常数力''). 可以证明, 这样构造的弱解在分布意义下唯一 (见~\cite{karatzas} 命题 5.3.10); 事实上, Zvonkin 与 Veretennikov 证明了此时还有强解. 下面的例子说明, 强解确实可能不存在.

\begin{example}[Tanaka 方程]\label{ex:tanaka-sde}
考虑方程
\[
\dd X_t=\sgn(X_t)\dd W_t,\qquad X_0=0,
\]
其中 $\sgn(x)=1$ ($x>0$), $-1$ ($x\le0$), 与例~\ref{ex:levy-applications} 相同.

(1) \emph{弱解存在.} 设 $X$ 是标准布朗运动, 令 $W_t=\int_0^t\sgn(X_s)\dd X_s$. 由 L\'evy 刻画 (例~\ref{ex:levy-applications} (1)), $W$ 是布朗运动, 且 $\int_0^t\sgn(X_s)\dd W_s=\int_0^t\sgn(X_s)^2\dd X_s=X_t$. 所以 $(X,W)$ 是弱解.

(2) \emph{分布唯一.} 若 $(X,W)$ 是任一弱解, 则 $X=\int\sgn(X)\dd W$ 满足 $\ip{X}_t=\int_0^t\sgn(X_s)^2\dd s=t$, 由 L\'evy 刻画~\ref{thm:levy-char}, $X$ 是布朗运动. 所以一切弱解有相同的分布.

(3) \emph{轨道唯一性不成立.} 设 $(X,W)$ 是弱解. 由于 $\sgn(-x)=-\sgn(x)$ 对 $x\ne0$ 成立, 而 $X$ 是布朗运动, $\int_0^T\1_{\{X_s=0\}}\dd s=0$ $\as$, 所以 $\sgn(-X_s)=-\sgn(X_s)$ 对几乎每个 $s$ 成立, 从而 $\int_0^t\sgn(-X_s)\dd W_s=-X_t$: $-X$ 也是 (同一个 $W$ 驱动的) 解.

(4) \emph{强解不存在.} 设 $(X,W)$ 是弱解. 由 Tanaka 公式~\ref{thm:tanaka} (应用于布朗运动 $X$),
\[
W_t=\int_0^t\sgn(X_s)\dd X_s=\abs{X_t}-L_t,\qquad L_t=\lim_{\varepsilon\downarrow0}\frac{1}{2\varepsilon}\int_0^t\1_{\{\abs{X_s}<\varepsilon\}}\dd s.
\]
右端只依赖于 $\abs{X}$, 所以 $\mathcal{F}^W_t\subseteq\mathcal{F}^{\abs{X}}_t$ (扩充后). 若 $X$ 是强解, 即 $X_t$ 是 $\mathcal{F}^W_t$-可测的, 则 $X_t$ 是 $\mathcal{F}^{\abs{X}}_t$-可测的, 即存在可测泛函 $G$ 使 $X_t=G(\abs{X_s}:s\le t)$ $\as$ 但 $-X$ 也是布朗运动, $(X,\abs{X})$ 与 $(-X,\abs{X})$ 同分布, 所以 $-X_t=G(\abs{X_s}:s\le t)$ 也几乎必然成立, 从而 $X_t=0$ $\as$, 对 $t>0$ 矛盾.
\end{example}

\begin{intuition}
强解是 ``可以由噪声计算出来'' 的解: 知道了驱动的布朗运动 $W$ 的整条轨道, 就知道了 $X$. 对 Tanaka 方程, $W$ 只记录了 $\abs{X}$ 的信息, 每当 $X$ 回到零点, 它下一段 ``向上走还是向下走'' 是 $W$ 无法决定的额外随机性. 弱解则只要求 $(X,W)$ 的联合分布正确, 允许 $X$ 携带 $W$ 之外的随机性. Yamada-Watanabe 定理把两种概念联系起来: 弱解存在 + 轨道唯一性 $\Rightarrow$ 强解存在且分布唯一 (见~\cite{karatzas} 推论 5.3.23).
\end{intuition}

% =====================================================================
\section{扩散过程: 马尔可夫性与生成元}
\label{sec:6-3}

本节假设系数不依赖于时间: $b=b(x)$, $\sigma=\sigma(x)$, 且满足 \eqref{eq:H1}\eqref{eq:H2} (于是 $b,\sigma$ 连续). 对 $s\ge0$ 与 $\mathcal{F}_s$-可测的 $\eta\in L^2$, 记 $X^{s,\eta}=\{X^{s,\eta}_t\}_{t\ge s}$ 为从时刻 $s$、位置 $\eta$ 出发的解:
\[
X^{s,\eta}_t=\eta+\int_s^tb(X^{s,\eta}_u)\dd u+\int_s^t\sigma(X^{s,\eta}_u)\dd B_u,\qquad t\ge s.
\]
定理~\ref{thm:sde-existence} 与命题~\ref{prop:sde-continuity} 对初始时刻 $s$ 同样成立 (证明只需平移时间). 记 $X^x=X^{0,x}$, 并用 $\P_x$, $\E_x$ 表示与 $X^x$ 有关的概率与期望. 时齐方程的解称为 (It\^o) \textbf{扩散过程}.

% ---------------------------------------------------------------------
\subsection{马尔可夫性}

\begin{lemma}\label{lem:sde-shift}
设 $s\ge0$, $y,z\in\R$, 记 $W_u=B_{s+u}-B_s$.
\begin{enumerate}
\item $\{X^{s,y}_{s+u}\}_{u\ge0}$ 关于 $\sigma(W_u:u\ge0)$ (扩充后) 可测, 从而与 $\mathcal{F}_s$ 独立;
\item $\{X^{s,y}_{s+u}\}_{u\ge0}$ 与 $\{X^y_u\}_{u\ge0}$ 有相同的有限维分布;
\item 对有界连续函数 $f$ 与 $h\ge0$, $y\mapsto P_hf(y):=\E f(X^y_h)$ 是连续函数.
\end{enumerate}
\end{lemma}

\begin{proof}
(1)(2) 由定理~\ref{thm:bm-markov}, $W$ 是与 $\mathcal{F}_s$ 独立的布朗运动. $Y_u=X^{s,y}_{s+u}$ 满足以 $W$ 为驱动的方程 $Y_u=y+\int_0^ub(Y_r)\dd r+\int_0^u\sigma(Y_r)\dd W_r$ (随机积分的时间平移可由 Riemann 和逼近验证, 命题~\ref{prop:riemann-approx}). 对这个方程作 Picard 迭代: 每个迭代都是连续过程, 由命题~\ref{prop:riemann-approx}, 关于 $W$ 的随机积分是 $W$ 与被积过程的 Riemann 和的依概率极限; 对 $k$ 归纳可知, 每个迭代都关于 $\sigma(W)$ 可测, 而且它作为 $W$ 的泛函, 与以 $B$ 为驱动、从 $y$ 出发的迭代是 ``同一个'' 泛函 (Riemann 和的表达式相同). 由于 $W$ 与 $B$ 同分布, 二者的迭代 (连同驱动的布朗运动) 有相同的有限维分布; 取极限 (几乎必然收敛保持同分布) 即得 (1)(2).

(3) 由命题~\ref{prop:sde-continuity}, $z\to y$ 时 $X^z_h\to X^y_h$ 于 $L^2$, 从而依概率收敛. 由于 $f$ 有界连续, $f(X^z_h)\to f(X^y_h)$ 依概率成立, 由有界收敛定理 $P_hf(z)\to P_hf(y)$.
\end{proof}

\begin{theorem}[扩散过程的马尔可夫性]\label{thm:sde-markov}
设 $X=X^{0,\xi}$, $\xi\in L^2(\mathcal{F}_0)$. 则对 $s,h\ge0$ 与有界 Borel 函数 $f$,
\[
\E\big[f(X_{s+h})\,\big|\,\mathcal{F}_s\big]=P_hf(X_s),\qquad P_hf(y)=\E f(X^y_h).
\]
所以 $X$ 是时齐马尔可夫过程, 转移概率为 $P_h(y,A)=\P(X^y_h\in A)$.
\end{theorem}

\begin{proof}
\emph{第一步.} 由唯一性, $X_{s+\cdot}$ 就是从时刻 $s$、位置 $X_s$ 出发的解: $X_t=X^{s,X_s}_t$ ($t\ge s$), 因为二者满足同一个方程.

\emph{第二步: 初值取可数多个值.} 设 $\eta=\sum_jy_j\1_{A_j}$, 其中 $\{A_j\}\subseteq\mathcal{F}_s$ 是 $\Omega$ 的可数分割. 令 $Y=\sum_j\1_{A_j}X^{s,y_j}$. 在 $A_j$ 上, $Y=X^{s,y_j}$, 由局部性引理~\ref{lem:locality}, 在 $A_j$ 上 $\int_s^t\sigma(Y_u)\dd B_u=\int_s^t\sigma(X^{s,y_j}_u)\dd B_u$; 所以 $Y$ 满足初值为 $\eta$ 的方程, 由唯一性 $Y=X^{s,\eta}$. 于是, 由提出已知因子、引理~\ref{lem:sde-shift} (1)(2),
\begin{align*}
\E\big[f(X^{s,\eta}_{s+h})\,\big|\,\mathcal{F}_s\big]&=\sum_j\1_{A_j}\E\big[f(X^{s,y_j}_{s+h})\,\big|\,\mathcal{F}_s\big]=\sum_j\1_{A_j}\E f(X^{s,y_j}_{s+h})\\
&=\sum_j\1_{A_j}P_hf(y_j)=P_hf(\eta).
\end{align*}

\emph{第三步: 一般初值, $f$ 有界连续.} 令 $\eta_n=2^{-n}\lfloor2^nX_s\rfloor$, 则 $\abs{\eta_n-X_s}\le2^{-n}$. 由命题~\ref{prop:sde-continuity} (初始时刻为 $s$), $X^{s,\eta_n}_{s+h}\to X^{s,X_s}_{s+h}=X_{s+h}$ 于 $L^2$. 由有界收敛定理, 第二步等式的左端在 $L^1$ 中收敛到 $\E[f(X_{s+h})\,|\,\mathcal{F}_s]$; 由引理~\ref{lem:sde-shift} (3), 右端 $P_hf(\eta_n)\to P_hf(X_s)$.

\emph{第四步: $f$ 有界 Borel.} 满足结论 (且 $P_hf$ 为 Borel 函数) 的有界 Borel 函数 $f$ 全体是包含常数的线性空间, 由单调收敛定理与条件单调收敛定理对有界单调极限封闭, 由第三步包含有界连续函数. 由引理~\ref{lem:cont-to-borel} ($n=1$), 它包含全体有界 Borel 函数.
\end{proof}

与布朗运动相同 (定理~\ref{thm:bm-markov} 与定理~\ref{thm:strong-markov}), 扩散过程对整条未来轨道的泛函也有马尔可夫性, 并且有强马尔可夫性: 对停时 $\tau$, 在 $\{\tau<\infty\}$ 上 $\E[F(X_{\tau+\cdot})\,|\,\mathcal{F}_\tau]=\E_{X_\tau}[F]$ (见~\cite{karatzas} 定理 5.4.20). 引理~\ref{lem:sde-shift} (3) 说明 $P_h$ 把有界连续函数映为有界连续函数, 这称为 \textbf{Feller 性}.

% ---------------------------------------------------------------------
\subsection{生成元与 Dynkin 公式}

马尔可夫过程的转移概率满足 Chapman-Kolmogorov 方程 $P_{t+s}=P_tP_s$ (定理~\ref{thm:chapman-kolmogorov}), 所以整个半群 $\{P_t\}$ 由它在 $t=0$ 附近的行为决定, 就像 $e^{tA}$ 由 $A$ 决定一样.

\begin{definition}[生成元]\label{def:generator}
扩散过程的\textbf{无穷小生成元}定义为
\[
\mathcal{A}f(x)=\lim_{t\downarrow0}\frac{\E_xf(X_t)-f(x)}{t},
\]
定义域为使极限对一切 $x$ 存在的函数 $f$. 对二阶可微的 $f$, 记
\begin{equation}\label{eq:L-operator}
Lf(x)=b(x)f'(x)+\frac{1}{2}\sigma(x)^2f''(x).
\end{equation}
\end{definition}

记 $C^2_c$ 为二阶连续可微且支撑紧的函数全体.

\begin{proposition}[Dynkin 公式]\label{prop:dynkin}
设 $f\in C^2_c$, $\tau$ 是停时, $\E_x\tau<\infty$. 则
\[
\E_xf(X_\tau)=f(x)+\E_x\int_0^\tau Lf(X_s)\dd s.
\]
\end{proposition}

\begin{proof}
由 It\^o 公式, $f(X_t)=f(x)+\int_0^tLf(X_s)\dd s+\int_0^tf'(X_s)\sigma(X_s)\dd B_s$. 由于 $f'$ 支撑紧而 $\sigma$ 连续, $f'\sigma$ 有界, 所以随机积分是鞅 (定理~\ref{thm:ito-properties}). 对有界停时 $\tau\wedge n$ 应用可选停时定理 (推论~\ref{cor:ct-ost-bounded}),
\[
\E_xf(X_{\tau\wedge n})=f(x)+\E_x\int_0^{\tau\wedge n}Lf(X_s)\dd s.
\]
令 $n\to\infty$: 左端由有界收敛定理收敛 ($\tau<\infty$ $\as$); 右端被积函数 $Lf$ 有界, $\abs{\int_0^{\tau\wedge n}Lf\dd s}\le\norm{Lf}_\infty\tau\in L^1$, 由控制收敛定理收敛.
\end{proof}

\begin{corollary}[生成元的形式]\label{cor:generator}
对 $f\in C^2_c$, $\mathcal{A}f=Lf$.
\end{corollary}

\begin{proof}
在 Dynkin 公式中取 $\tau=t$, 由 Fubini 定理,
\[
\frac{\E_xf(X_t)-f(x)}{t}=\frac{1}{t}\int_0^t\E_xLf(X_s)\dd s.
\]
$Lf$ 有界连续, $X$ 的轨道连续, 由有界收敛定理, $s\mapsto\E_xLf(X_s)$ 在 $s=0$ 处连续, 其值为 $Lf(x)$. 所以右端趋于 $Lf(x)$.
\end{proof}

\begin{intuition}
生成元把扩散过程归结为两个 ``局部特征'': 在状态 $x$, 过程在短时间 $\Delta t$ 内的增量近似满足
\[
\E[\Delta X\,|\,X_t=x]\approx b(x)\Delta t,\qquad\Var(\Delta X\,|\,X_t=x)\approx\sigma(x)^2\Delta t.
\]
$Lf$ 中的一阶项来自平均位移, 二阶项来自方差 (这正是 It\^o 公式的修正项). 对布朗运动, $L=\frac{1}{2}\frac{\dd^2}{\dd x^2}$ (第~4~章的直观); 对 OU 过程 $\dd X=-\kappa X\dd t+\sigma\dd B$, $L=-\kappa x\frac{\dd}{\dd x}+\frac{\sigma^2}{2}\frac{\dd^2}{\dd x^2}$; 对几何布朗运动, $L=\mu x\frac{\dd}{\dd x}+\frac{\sigma^2x^2}{2}\frac{\dd^2}{\dd x^2}$. 多维情形 $\dd\rvec{X}=\rvec{b}\dd t+\mat{\Sigma}\dd\rvec{B}$, 生成元为 $L=\sum_ib_i\partial_i+\frac{1}{2}\sum_{i,j}a_{ij}\partial_{ij}$, 其中 $\mat{A}=(a_{ij})=\mat{\Sigma}\mat{\Sigma}^\T$ 是 ``扩散矩阵''.
\end{intuition}

% ---------------------------------------------------------------------
\subsection{Kolmogorov 方程与平稳分布}

第~4~章证明了布朗运动的转移密度满足热方程 (定理~\ref{thm:heat-equation}). 对一般的扩散过程, 热方程中的 $\frac{1}{2}\partial_{xx}$ 换成生成元 $L$ 或其伴随算子.

\begin{theorem}[Kolmogorov 向后方程]\label{thm:backward}
设 $f$ 有界连续, $u\in C^{1,2}((0,\infty)\times\R)\cap C([0,\infty)\times\R)$ 有界, 且对每个 $T$, $\partial_xu$ 在 $(0,T]\times\R$ 上有界. 若
\[
\partial_tu=Lu\quad\text{于 }(0,\infty)\times\R,\qquad u(0,x)=f(x),
\]
则 $u(t,x)=\E_xf(X_t)$.
\end{theorem}

\begin{proof}
固定 $t>0$ 与 $x$, 令 $v(s,y)=u(t-s,y)$ ($0\le s<t$), 则 $\partial_sv+Lv=0$. 对 $0<\varepsilon<t$, 在 $[0,t-\varepsilon]$ 上对 $v(s,X_s)$ ($X=X^x$) 应用 It\^o 公式,
\[
v(t-\varepsilon,X_{t-\varepsilon})=v(0,x)+\int_0^{t-\varepsilon}(\partial_sv+Lv)(s,X_s)\dd s+\int_0^{t-\varepsilon}\partial_yv(s,X_s)\sigma(X_s)\dd B_s.
\]
中间的积分为零. 由于 $\partial_yv$ 有界而 $\E\int_0^t\sigma(X_s)^2\dd s<\infty$ (由 \eqref{eq:H2} 与 \eqref{eq:sde-moment}), 随机积分的期望为零, 所以 $\E u(\varepsilon,X_{t-\varepsilon})=u(t,x)$. 令 $\varepsilon\downarrow0$, 由 $u$ 在 $[0,\infty)\times\R$ 上的连续性与轨道连续性, $u(\varepsilon,X_{t-\varepsilon})\to u(0,X_t)=f(X_t)$; 由有界收敛定理, $u(t,x)=\E_xf(X_t)$.
\end{proof}

定理说: 向后方程的 (足够正则的) 解必定由扩散过程表示, 特别地解是唯一的. 反过来, 若 $b,\sigma$ 与 $f$ 足够光滑 (例如具有有界的各阶导数), 则 $u(t,x)=\E_xf(X_t)$ 确实是 $C^{1,2}$ 的, 从而满足向后方程 (见~\cite{karatzas} 5.7 节). 它描述期望 $\E_xf(X_t)$ 如何随\emph{起点} $x$ 变化. 与之对偶, 向前方程描述 $X_t$ 的\emph{分布}如何随时间演化.

\begin{theorem}[Kolmogorov 向前方程 (Fokker-Planck 方程)]\label{thm:forward}
设 $b\in C^1$, $\sigma\in C^2$, $X_t$ ($X_0=x$) 对 $t>0$ 有密度 $p(t,y)$, $p\in C^{1,2}((0,\infty)\times\R)$, 并且可以在积分号下对 $t$ 求导 (即 $\frac{\dd}{\dd t}\int f(y)p(t,y)\dd y=\int f(y)\partial_tp(t,y)\dd y$ 对 $f\in C^2_c$ 成立). 则
\[
\partial_tp=L^*p:=-\partial_y\big(b\,p\big)+\frac{1}{2}\partial_{yy}\big(\sigma^2p\big).
\]
\end{theorem}

\begin{proof}
设 $f\in C^2_c$. 由 Dynkin 公式 (取 $\tau=t$), $\E_xf(X_t)=f(x)+\int_0^t\E_xLf(X_s)\dd s$, 被积函数关于 $s$ 连续, 所以
\[
\int f(y)\partial_tp(t,y)\dd y=\frac{\dd}{\dd t}\E_xf(X_t)=\E_xLf(X_t)=\int\big(b(y)f'(y)+\tfrac{1}{2}\sigma(y)^2f''(y)\big)p(t,y)\dd y.
\]
由于 $f$ 支撑紧, 分部积分没有边界项: $\int bf'p\dd y=-\int f\,\partial_y(bp)\dd y$, $\int\sigma^2f''p\dd y=\int f\,\partial_{yy}(\sigma^2p)\dd y$. 于是 $\int f\,(\partial_tp-L^*p)\dd y=0$ 对一切 $f\in C^2_c$ 成立. 连续函数若与一切 $C^2_c$ 函数的积分都为零, 则恒为零, 所以 $\partial_tp=L^*p$.
\end{proof}

\begin{figure}[htbp]
\centering
\begin{tikzpicture}
\begin{axis}[SAaxis, axis x line=bottom, axis y line=left, width=12cm, height=5.6cm, xmin=-2.6, xmax=4.4, ymin=0, ymax=1.42,
  xtick={-2,-1,0,1,2,3,4}, ytick={0,0.5,1}, xlabel={$y$}, samples=200, domain=-2.6:4.4,
  legend style={font=\small,draw=none,fill=none,at={(0.5,-0.22)},anchor=north}, legend columns=5, legend cell align=left]
  \addplot[SARed,thick] {exp(-(x-3*exp(-0.1))^2/(1-exp(-0.2)))/sqrt(pi*(1-exp(-0.2)))};
  \addplot[SAYellowText,thick] {exp(-(x-3*exp(-0.5))^2/(1-exp(-1)))/sqrt(pi*(1-exp(-1)))};
  \addplot[SAGreenText,thick] {exp(-(x-3*exp(-1))^2/(1-exp(-2)))/sqrt(pi*(1-exp(-2)))};
  \addplot[SABlue,thick] {exp(-(x-3*exp(-3))^2/(1-exp(-6)))/sqrt(pi*(1-exp(-6)))};
  \addplot[SADark,thick,dashed] {exp(-x^2)/sqrt(pi)};
  \legend{$t=0.1$, $t=0.5$, $t=1$, $t=3$, 平稳分布}
\end{axis}
\end{tikzpicture}
\caption{Fokker-Planck 方程的解: OU 过程 $\dd X_t=-X_t\dd t+\dd B_t$ 从 $x=3$ 出发时 $X_t$ 的密度 $N\big(3e^{-t},\frac12(1-e^{-2t})\big)$. 密度从点质量出发, 一边扩散一边向原点漂移, 最终收敛到平稳分布 $N(0,\frac12)$ (虚线)}
\label{fig:fokker-planck}
\end{figure}

若概率分布 $\pi$ 满足: 当 $X_0\sim\pi$ (与 $B$ 独立) 时, 对一切 $t$ 都有 $X_t\sim\pi$, 则称 $\pi$ 为\textbf{平稳分布}. 由向前方程, 平稳分布的密度应满足 $L^*\pi=0$. 在一维情形, 若进一步要求 ``概率流'' $b\pi-\frac{1}{2}(\sigma^2\pi)'$ 为零 (这在 $\pi$ 及其导数在 $\pm\infty$ 处趋于零时是自然的), 则得到一阶方程 $(\frac{1}{2}\sigma^2\pi)'=b\pi$, 其解为
\begin{equation}\label{eq:stationary-density}
\pi(y)=\frac{C}{\sigma(y)^2}\exp\Big(\int_0^y\frac{2b(z)}{\sigma(z)^2}\dd z\Big),
\end{equation}
其中 $C$ 是归一化常数 (要求右端可积). 直接验证: 令 $m=\frac{1}{2}\sigma^2\pi=\frac{C}{2}\exp(\int_0^y\frac{2b}{\sigma^2})$, 则 $m'=\frac{2b}{\sigma^2}m=b\pi$, 从而 $L^*\pi=-(b\pi)'+m''=0$. 可以证明, 在适当的条件下 (例如 $\sigma$ 有正下界且 \eqref{eq:stationary-density} 可积), 这就是唯一的平稳分布, 且对任意初值, $X_t$ 的分布收敛到 $\pi$.

\begin{example}[OU 过程与 Langevin 方程]\label{ex:stationary}
(1) 对 OU 过程 $\dd X=\kappa(\theta-X)\dd t+\sigma\dd B$, $\frac{2b}{\sigma^2}=\frac{2\kappa(\theta-y)}{\sigma^2}$, 由 \eqref{eq:stationary-density}, $\pi\propto e^{-\kappa(y-\theta)^2/\sigma^2}$, 即 $\pi=N(\theta,\frac{\sigma^2}{2\kappa})$. 这可以直接验证: 若 $X_0\sim N(\theta,\frac{\sigma^2}{2\kappa})$ 与 $B$ 独立, 由例~\ref{ex:linear-examples} (2), $X_t$ 是正态的, 均值为 $\theta$, 方差为 $e^{-2\kappa t}\frac{\sigma^2}{2\kappa}+\frac{\sigma^2}{2\kappa}(1-e^{-2\kappa t})=\frac{\sigma^2}{2\kappa}$.

(2) 设 $V$ 是光滑的 ``势能'' 函数, 考虑 \textbf{Langevin 方程} $\dd X_t=-V'(X_t)\dd t+\sqrt{2}\dd B_t$: 粒子沿势能下降的方向运动, 同时受到热噪声的扰动. 由 \eqref{eq:stationary-density}, 平稳分布为 \textbf{Gibbs 分布}
\[
\pi(y)=\frac{1}{Z}e^{-V(y)},\qquad Z=\int_{\R}e^{-V(y)}\dd y.
\]
当 $V$ 有多个极小点时 (例如双势阱 $V(y)=(y^2-1)^2$), 过程大部分时间停留在某个势阱附近, 偶尔借助噪声越过势垒跳到另一个势阱 (图~\ref{fig:langevin}). 反过来, 若想从一个已知 (至多相差一个常数因子) 的密度 $\pi$ 中抽样, 可以令 $V=-\log\pi$, 模拟 Langevin 方程足够长时间: 这是 Langevin Monte Carlo 方法的基础, 也是现代扩散生成模型的出发点之一.
\end{example}

\begin{figure}[htbp]
\centering
\begin{tikzpicture}
\begin{axis}[SAaxis, axis x line=bottom, axis y line=left, width=6.2cm, height=5.2cm, xmin=-2.1, xmax=2.1, ymin=0, ymax=1.6,
  xtick={-2,-1,0,1,2}, ytick={0,0.5,1,1.5}, xlabel={$y$}, samples=161, domain=-2:2,
  legend style={font=\scriptsize,draw=none,fill=none,at={(0.5,1.02)},anchor=south}, legend columns=2, legend cell align=left]
  \addplot[SADark,thick] {(x^2-1)^2};
  \addplot[SARed,thick,fill=SARed!12] {exp(-(x^2-1)^2)/1.97373} \closedcycle;
  \legend{$V(y)$, $\pi(y)$}
\end{axis}
\end{tikzpicture}\hfill
\begin{tikzpicture}
\begin{axis}[SAaxis, axis x line=bottom, axis y line=left, width=9cm, height=5.2cm, xmin=0, xmax=203, ymin=-2.1, ymax=2.1,
  xtick={0,50,100,150,200}, ytick={-1,0,1}, xlabel={$t$}]
  \addplot[SADark!25] coordinates {(0,0) (200,0)};
  \addplot[SABlue,very thin] coordinates {(0.000,-1.000) (0.250,-0.826) (0.500,-0.306) (0.750,-0.724) (1.000,-0.592) (1.250,0.757) (1.500,0.760) (1.750,1.278) (2.000,1.055) (2.250,1.144) (2.500,0.969) (2.750,1.230) (3.000,1.443) (3.250,0.577) (3.500,1.312) (3.750,1.028) (4.000,1.450) (4.250,1.280) (4.500,1.055) (4.750,0.932) (5.000,1.130) (5.250,0.572) (5.500,0.476) (5.750,0.522) (6.000,0.880) (6.250,0.493) (6.500,0.833) (6.750,0.934) (7.000,0.082) (7.250,0.841) (7.500,1.437) (7.750,1.173) (8.000,1.676) (8.250,1.006) (8.500,1.286) (8.750,0.432) (9.000,1.208) (9.250,1.375) (9.500,1.048) (9.750,1.058) (10.000,1.309) (10.250,1.232) (10.500,0.820) (10.750,0.475) (11.000,0.436) (11.250,-0.645) (11.500,-0.221) (11.750,-0.540) (12.000,-0.922) (12.250,-0.592) (12.500,-0.266) (12.750,-1.062) (13.000,-1.002) (13.250,-0.123) (13.500,1.088) (13.750,1.199) (14.000,1.007) (14.250,1.008) (14.500,0.669) (14.750,0.023) (15.000,-1.223) (15.250,-0.546) (15.500,-0.997) (15.750,-1.043) (16.000,-0.863) (16.250,0.192) (16.500,1.342) (16.750,0.603) (17.000,1.413) (17.250,1.052) (17.500,1.398) (17.750,1.075) (18.000,0.685) (18.250,0.378) (18.500,0.967) (18.750,1.613) (19.000,-0.181) (19.250,0.541) (19.500,0.809) (19.750,0.219) (20.000,1.127) (20.250,1.072) (20.500,1.143) (20.750,0.456) (21.000,1.080) (21.250,0.704) (21.500,1.485) (21.750,1.230) (22.000,0.260) (22.250,0.638) (22.500,0.042) (22.750,-0.679) (23.000,-0.801) (23.250,-1.043) (23.500,-1.240) (23.750,-0.575) (24.000,-0.246) (24.250,0.597) (24.500,1.048) (24.750,-0.201) (25.000,-0.872) (25.250,-1.169) (25.500,-1.195) (25.750,-0.787) (26.000,-0.386) (26.250,0.009) (26.500,-0.008) (26.750,-0.427) (27.000,-1.057) (27.250,-0.341) (27.500,-0.404) (27.750,-0.231) (28.000,-0.223) (28.250,-0.462) (28.500,-0.492) (28.750,-0.631) (29.000,-1.015) (29.250,-0.106) (29.500,1.487) (29.750,0.890) (30.000,-0.475) (30.250,-1.051) (30.500,-1.070) (30.750,-0.869) (31.000,-0.807) (31.250,-1.230) (31.500,-0.941) (31.750,-0.477) (32.000,-1.088) (32.250,-0.438) (32.500,-0.131) (32.750,-0.808) (33.000,-1.263) (33.250,-1.121) (33.500,-1.393) (33.750,-0.964) (34.000,-1.259) (34.250,-0.170) (34.500,-0.372) (34.750,0.928) (35.000,1.486) (35.250,1.121) (35.500,0.754) (35.750,1.147) (36.000,0.501) (36.250,0.747) (36.500,1.293) (36.750,0.701) (37.000,1.452) (37.250,1.053) (37.500,0.529) (37.750,-0.187) (38.000,-0.882) (38.250,-1.117) (38.500,0.258) (38.750,0.462) (39.000,0.447) (39.250,-0.426) (39.500,0.566) (39.750,1.057) (40.000,0.445) (40.250,0.290) (40.500,0.021) (40.750,0.641) (41.000,0.581) (41.250,0.015) (41.500,-0.335) (41.750,-0.593) (42.000,-1.044) (42.250,-1.113) (42.500,-0.867) (42.750,-0.472) (43.000,0.438) (43.250,0.529) (43.500,1.298) (43.750,0.903) (44.000,0.283) (44.250,-0.894) (44.500,-0.833) (44.750,-0.153) (45.000,0.463) (45.250,0.249) (45.500,0.563) (45.750,1.131) (46.000,0.940) (46.250,0.124) (46.500,1.223) (46.750,1.343) (47.000,0.930) (47.250,0.117) (47.500,0.001) (47.750,0.772) (48.000,0.616) (48.250,0.049) (48.500,0.146) (48.750,0.552) (49.000,1.326) (49.250,1.222) (49.500,0.855) (49.750,1.052) (50.000,1.249) (50.250,1.428) (50.500,1.196) (50.750,0.945) (51.000,1.231) (51.250,0.987) (51.500,1.018) (51.750,1.164) (52.000,0.557) (52.250,0.743) (52.500,1.042) (52.750,0.857) (53.000,0.260) (53.250,0.825) (53.500,0.897) (53.750,0.736) (54.000,0.952) (54.250,1.023) (54.500,0.528) (54.750,-0.409) (55.000,0.133) (55.250,-0.847) (55.500,-1.132) (55.750,-1.050) (56.000,-1.269) (56.250,-1.066) (56.500,-0.643) (56.750,-1.187) (57.000,-1.142) (57.250,-0.972) (57.500,-1.108) (57.750,-0.882) (58.000,-0.588) (58.250,-0.825) (58.500,-0.765) (58.750,-1.162) (59.000,-0.529) (59.250,-0.338) (59.500,-0.252) (59.750,-0.681) (60.000,-1.126) (60.250,-1.270) (60.500,-1.242) (60.750,-1.266) (61.000,-0.725) (61.250,-0.870) (61.500,-1.439) (61.750,-0.757) (62.000,-1.204) (62.250,-0.571) (62.500,-1.320) (62.750,-0.738) (63.000,-0.636) (63.250,-1.291) (63.500,-0.114) (63.750,0.070) (64.000,-0.401) (64.250,-0.579) (64.500,-0.574) (64.750,-0.358) (65.000,-1.066) (65.250,-0.943) (65.500,-1.187) (65.750,-0.633) (66.000,0.242) (66.250,0.625) (66.500,0.826) (66.750,0.804) (67.000,0.414) (67.250,0.725) (67.500,0.694) (67.750,0.872) (68.000,0.778) (68.250,0.857) (68.500,0.198) (68.750,0.230) (69.000,-0.803) (69.250,-0.920) (69.500,-1.430) (69.750,-0.845) (70.000,-0.369) (70.250,-0.805) (70.500,-1.132) (70.750,-0.809) (71.000,-1.188) (71.250,-0.631) (71.500,-0.937) (71.750,-1.040) (72.000,-0.660) (72.250,-1.013) (72.500,-0.742) (72.750,-1.494) (73.000,-1.100) (73.250,-1.104) (73.500,-0.480) (73.750,-0.492) (74.000,0.990) (74.250,0.559) (74.500,0.988) (74.750,1.117) (75.000,1.198) (75.250,0.969) (75.500,1.012) (75.750,0.702) (76.000,0.847) (76.250,-0.043) (76.500,0.489) (76.750,1.074) (77.000,-0.204) (77.250,-0.738) (77.500,-0.677) (77.750,-1.286) (78.000,-0.937) (78.250,-0.998) (78.500,-0.856) (78.750,-1.052) (79.000,-1.179) (79.250,-0.986) (79.500,-0.986) (79.750,-0.538) (80.000,-1.035) (80.250,-0.744) (80.500,-0.553) (80.750,-0.750) (81.000,-0.280) (81.250,-1.239) (81.500,-0.733) (81.750,-1.470) (82.000,-0.556) (82.250,1.106) (82.500,1.097) (82.750,1.467) (83.000,1.129) (83.250,-0.061) (83.500,-0.539) (83.750,-0.345) (84.000,-1.061) (84.250,-0.665) (84.500,-1.530) (84.750,-1.165) (85.000,-1.112) (85.250,-1.248) (85.500,-1.297) (85.750,-0.294) (86.000,0.431) (86.250,0.983) (86.500,1.304) (86.750,0.508) (87.000,0.588) (87.250,0.447) (87.500,1.218) (87.750,1.338) (88.000,1.269) (88.250,0.655) (88.500,1.013) (88.750,0.223) (89.000,-0.139) (89.250,0.044) (89.500,-0.356) (89.750,-1.023) (90.000,-0.926) (90.250,-0.640) (90.500,-1.205) (90.750,0.079) (91.000,-0.886) (91.250,-0.968) (91.500,-0.503) (91.750,0.391) (92.000,0.834) (92.250,-0.166) (92.500,-0.776) (92.750,-0.544) (93.000,-0.690) (93.250,-0.350) (93.500,-0.442) (93.750,0.985) (94.000,1.024) (94.250,0.242) (94.500,0.920) (94.750,0.825) (95.000,-0.064) (95.250,-0.656) (95.500,-0.340) (95.750,-0.441) (96.000,-1.212) (96.250,-0.531) (96.500,-0.765) (96.750,-0.899) (97.000,-1.271) (97.250,-1.157) (97.500,-1.097) (97.750,-0.856) (98.000,-0.096) (98.250,0.035) (98.500,1.029) (98.750,0.270) (99.000,0.541) (99.250,-0.367) (99.500,0.931) (99.750,1.191) (100.000,1.087) (100.250,1.095) (100.500,1.186) (100.750,1.262) (101.000,1.083) (101.250,0.964) (101.500,1.198) (101.750,0.967) (102.000,1.367) (102.250,1.193) (102.500,1.061) (102.750,1.102) (103.000,1.371) (103.250,0.265) (103.500,1.175) (103.750,0.904) (104.000,1.001) (104.250,0.027) (104.500,-0.937) (104.750,-0.663) (105.000,-1.074) (105.250,-0.941) (105.500,-0.870) (105.750,-0.784) (106.000,-0.438) (106.250,0.504) (106.500,0.507) (106.750,-0.136) (107.000,0.213) (107.250,1.320) (107.500,0.898) (107.750,1.546) (108.000,1.708) (108.250,0.151) (108.500,1.062) (108.750,1.176) (109.000,0.311) (109.250,0.386) (109.500,0.412) (109.750,0.468) (110.000,0.753) (110.250,0.972) (110.500,1.179) (110.750,1.235) (111.000,1.205) (111.250,0.896) (111.500,1.384) (111.750,0.954) (112.000,0.964) (112.250,1.438) (112.500,1.384) (112.750,1.321) (113.000,1.780) (113.250,-0.488) (113.500,-1.144) (113.750,-1.004) (114.000,-1.284) (114.250,-0.934) (114.500,-1.375) (114.750,-0.528) (115.000,-0.439) (115.250,-0.775) (115.500,-0.603) (115.750,-0.730) (116.000,-0.799) (116.250,-0.888) (116.500,-0.767) (116.750,-1.222) (117.000,-1.189) (117.250,-0.945) (117.500,-0.672) (117.750,-0.791) (118.000,-0.133) (118.250,-0.429) (118.500,-1.160) (118.750,-1.513) (119.000,-1.035) (119.250,-0.982) (119.500,-0.791) (119.750,-0.948) (120.000,-0.990) (120.250,-0.807) (120.500,-0.963) (120.750,-0.271) (121.000,-0.000) (121.250,-0.123) (121.500,-0.630) (121.750,-1.317) (122.000,-0.806) (122.250,-0.086) (122.500,-0.199) (122.750,0.659) (123.000,1.572) (123.250,1.210) (123.500,0.631) (123.750,0.243) (124.000,0.970) (124.250,1.092) (124.500,0.729) (124.750,1.071) (125.000,0.996) (125.250,0.765) (125.500,1.294) (125.750,1.250) (126.000,1.249) (126.250,0.888) (126.500,1.355) (126.750,1.153) (127.000,1.106) (127.250,0.771) (127.500,0.556) (127.750,0.363) (128.000,-0.304) (128.250,0.014) (128.500,-0.870) (128.750,-1.243) (129.000,-0.182) (129.250,-0.285) (129.500,-0.111) (129.750,-0.710) (130.000,-0.832) (130.250,-0.671) (130.500,-1.137) (130.750,-1.433) (131.000,-0.965) (131.250,-1.279) (131.500,-0.594) (131.750,-1.151) (132.000,-0.989) (132.250,-0.717) (132.500,-1.300) (132.750,-1.213) (133.000,-0.544) (133.250,-1.095) (133.500,-1.122) (133.750,-1.370) (134.000,-1.400) (134.250,-0.480) (134.500,-1.130) (134.750,-1.164) (135.000,-0.377) (135.250,-0.957) (135.500,-0.411) (135.750,-1.202) (136.000,-0.616) (136.250,-0.911) (136.500,-1.032) (136.750,-1.196) (137.000,-1.034) (137.250,-0.860) (137.500,-1.046) (137.750,-1.205) (138.000,-1.201) (138.250,-1.133) (138.500,-0.824) (138.750,-0.979) (139.000,-1.154) (139.250,-0.795) (139.500,-1.365) (139.750,-0.772) (140.000,-0.549) (140.250,-0.939) (140.500,-1.183) (140.750,-1.212) (141.000,0.018) (141.250,-0.318) (141.500,0.916) (141.750,0.737) (142.000,0.308) (142.250,-0.519) (142.500,-1.082) (142.750,-1.521) (143.000,-1.305) (143.250,-0.045) (143.500,-0.883) (143.750,-0.553) (144.000,-1.021) (144.250,-0.432) (144.500,-1.126) (144.750,-0.693) (145.000,-0.691) (145.250,0.131) (145.500,-0.432) (145.750,-0.933) (146.000,-0.717) (146.250,-0.738) (146.500,-0.852) (146.750,-0.652) (147.000,-0.930) (147.250,-1.362) (147.500,-1.461) (147.750,-0.800) (148.000,-1.299) (148.250,-0.954) (148.500,-1.239) (148.750,-1.019) (149.000,-1.111) (149.250,-1.124) (149.500,-0.987) (149.750,-0.873) (150.000,-0.993) (150.250,-1.179) (150.500,-0.895) (150.750,-0.620) (151.000,0.109) (151.250,0.241) (151.500,0.775) (151.750,-0.666) (152.000,-0.372) (152.250,-0.803) (152.500,0.226) (152.750,-0.464) (153.000,-0.716) (153.250,-0.860) (153.500,-0.938) (153.750,-0.474) (154.000,-1.468) (154.250,-0.974) (154.500,-1.169) (154.750,-1.107) (155.000,-1.034) (155.250,-1.186) (155.500,-0.668) (155.750,-1.109) (156.000,-0.509) (156.250,-0.884) (156.500,-0.893) (156.750,-0.717) (157.000,-1.054) (157.250,-0.948) (157.500,1.305) (157.750,1.708) (158.000,1.290) (158.250,1.639) (158.500,0.914) (158.750,1.102) (159.000,1.417) (159.250,1.135) (159.500,0.082) (159.750,0.168) (160.000,1.037) (160.250,0.934) (160.500,-0.190) (160.750,-0.645) (161.000,-1.358) (161.250,0.480) (161.500,-1.176) (161.750,-0.924) (162.000,-1.240) (162.250,-0.002) (162.500,-0.438) (162.750,-0.877) (163.000,-1.308) (163.250,-1.393) (163.500,-0.930) (163.750,-1.099) (164.000,-0.562) (164.250,-0.191) (164.500,0.886) (164.750,1.057) (165.000,0.896) (165.250,1.376) (165.500,1.009) (165.750,0.804) (166.000,0.807) (166.250,0.464) (166.500,1.135) (166.750,0.517) (167.000,1.370) (167.250,0.670) (167.500,1.440) (167.750,1.110) (168.000,1.290) (168.250,1.108) (168.500,0.711) (168.750,0.513) (169.000,0.538) (169.250,1.295) (169.500,0.967) (169.750,0.779) (170.000,0.477) (170.250,0.667) (170.500,0.870) (170.750,1.300) (171.000,1.082) (171.250,1.064) (171.500,1.275) (171.750,0.640) (172.000,1.252) (172.250,1.262) (172.500,1.160) (172.750,0.768) (173.000,-0.051) (173.250,0.092) (173.500,1.408) (173.750,0.969) (174.000,0.767) (174.250,0.866) (174.500,0.294) (174.750,0.980) (175.000,0.608) (175.250,-0.007) (175.500,-0.385) (175.750,-0.227) (176.000,-0.307) (176.250,-0.708) (176.500,-0.433) (176.750,0.061) (177.000,-0.000) (177.250,-0.278) (177.500,-0.630) (177.750,-1.054) (178.000,-0.960) (178.250,-0.951) (178.500,-0.373) (178.750,0.018) (179.000,-1.348) (179.250,0.122) (179.500,0.932) (179.750,1.389) (180.000,1.273) (180.250,1.264) (180.500,1.100) (180.750,1.159) (181.000,1.231) (181.250,0.641) (181.500,1.081) (181.750,0.650) (182.000,0.887) (182.250,0.709) (182.500,0.714) (182.750,0.939) (183.000,0.574) (183.250,0.371) (183.500,0.944) (183.750,1.439) (184.000,0.860) (184.250,-1.050) (184.500,-1.079) (184.750,-1.108) (185.000,-0.598) (185.250,-0.630) (185.500,-0.341) (185.750,-1.349) (186.000,-0.870) (186.250,-0.937) (186.500,-0.979) (186.750,-0.701) (187.000,-0.177) (187.250,-0.834) (187.500,-1.209) (187.750,-0.919) (188.000,0.432) (188.250,0.844) (188.500,0.972) (188.750,1.237) (189.000,1.112) (189.250,1.216) (189.500,0.395) (189.750,0.091) (190.000,-0.417) (190.250,-0.755) (190.500,-1.255) (190.750,-1.779) (191.000,-0.950) (191.250,-1.144) (191.500,-0.849) (191.750,-1.019) (192.000,-1.359) (192.250,-0.609) (192.500,-0.042) (192.750,-0.079) (193.000,-1.617) (193.250,-0.808) (193.500,-0.838) (193.750,-0.872) (194.000,-1.064) (194.250,-0.172) (194.500,-0.913) (194.750,-1.338) (195.000,-1.186) (195.250,-0.986) (195.500,0.866) (195.750,0.975) (196.000,1.606) (196.250,1.077) (196.500,0.458) (196.750,-0.172) (197.000,-0.625) (197.250,-0.149) (197.500,0.710) (197.750,0.071) (198.000,0.944) (198.250,0.698) (198.500,0.554) (198.750,0.516) (199.000,0.574) (199.250,0.600) (199.500,-1.106) (199.750,1.208) (200.000,0.978)};
\end{axis}
\end{tikzpicture}
\caption{双势阱 $V(y)=(y^2-1)^2$ 的 Langevin 方程. 左: 势能与 Gibbs 平稳密度 $\pi\propto e^{-V}$, 密度集中在两个势阱 $y=\pm1$ 附近. 右: 一条轨道, 它在两个势阱之间随机地来回跳跃}
\label{fig:langevin}
\end{figure}

% ---------------------------------------------------------------------
\subsection{一维扩散的出口问题}

第~4~章用鞅方法计算了布朗运动的出口概率与出口时间 (定理~\ref{thm:exit}). 对一般的一维扩散, 关键是找到使 $\phi(X_t)$ 成为鞅的函数 $\phi$, 即 $L\phi=0$ 的解.

设 $\ell<r$, 且 $\sigma>0$ 在 $[\ell,r]$ 上成立. 定义\textbf{尺度函数}
\begin{equation}\label{eq:scale}
s(x)=\int_{c}^x\exp\Big(-\int_{c}^y\frac{2b(z)}{\sigma(z)^2}\dd z\Big)\dd y,
\end{equation}
其中 $c\in[\ell,r]$ 是任取的参考点. 则 $s'>0$, 且 $s''=-\frac{2b}{\sigma^2}s'$, 即
\[
Ls=bs'+\frac{1}{2}\sigma^2s''=0.
\]
所以 $s(X_t)$ (在离开 $[\ell,r]$ 之前) 没有漂移: 尺度函数把扩散过程变成了 ``时间变换后的布朗运动'', 它的作用相当于布朗运动的恒等函数. 令 $\tau=\inf\{t: X_t\notin(\ell,r)\}$.

\begin{lemma}\label{lem:exit-finite}
对 $x\in(\ell,r)$, $\E_x\tau<\infty$.
\end{lemma}

\begin{proof}
记 $\sigma_0^2=\min_{[\ell,r]}\sigma^2>0$, $\beta=\max_{[\ell,r]}\abs{b}$, 取 $\lambda=\frac{2(\beta+1)}{\sigma_0^2}$. 令 $f(y)=e^{\lambda y}$ (在 $[\ell-1,r+1]$ 之外修改为 $C^2_c$ 函数, 不影响下面的论证). 在 $[\ell,r]$ 上,
\[
Lf=\Big(\lambda b+\frac{1}{2}\lambda^2\sigma^2\Big)e^{\lambda y}\ge\lambda\Big(\frac{1}{2}\lambda\sigma_0^2-\beta\Big)e^{\lambda\ell}=\lambda e^{\lambda\ell}=:c_0>0.
\]
对有界停时 $\tau\wedge n$ 应用 Dynkin 公式, $\E_xf(X_{\tau\wedge n})-f(x)=\E_x\int_0^{\tau\wedge n}Lf(X_s)\dd s\ge c_0\E_x(\tau\wedge n)$, 而左端不超过 $2e^{\lambda r}$. 令 $n\to\infty$, 由单调收敛定理 $\E_x\tau\le2e^{\lambda r}/c_0$.
\end{proof}

\begin{theorem}[一维扩散的出口问题]\label{thm:exit-diffusion}
设 $x\in(\ell,r)$. 则
\begin{enumerate}
\item $\P_x(X_\tau=r)=\dfrac{s(x)-s(\ell)}{s(r)-s(\ell)}$;
\item $\E_x\tau=v(x)$, 其中 $v$ 是边值问题 $Lv=-1$ ($\ell<x<r$), $v(\ell)=v(r)=0$ 的解.
\end{enumerate}
\end{theorem}

\begin{proof}
(1) 把 $s$ 在 $[\ell,r]$ 的一个邻域之外修改为 $C^2_c$ 函数. 由引理~\ref{lem:exit-finite} 与 Dynkin 公式 ($Ls=0$ 于 $[\ell,r]$), $\E_xs(X_\tau)=s(x)$. 而 $X_\tau\in\{\ell,r\}$, 记 $p=\P_x(X_\tau=r)$, 则 $p\,s(r)+(1-p)s(\ell)=s(x)$.

(2) 边值问题是二阶线性常微分方程的两点边值问题 (首项系数 $\frac{1}{2}\sigma^2>0$), 有唯一的 $C^2$ 解 (可用 $s$ 显式写出). 对 $v$ 应用 Dynkin 公式, $\E_xv(X_\tau)-v(x)=\E_x\int_0^\tau(-1)\dd s$, 而 $v(X_\tau)=0$, 所以 $\E_x\tau=v(x)$.
\end{proof}

\begin{example}\label{ex:exit-examples}
(1) 布朗运动: $s(x)=x$, $Lv=\frac{1}{2}v''=-1$ 给出 $v(x)=(x-\ell)(r-x)$, 与定理~\ref{thm:exit} 一致.

(2) 带漂移 $\mu\neq0$ 的布朗运动: $s(x)=\frac{1-e^{-2\mu x}}{2\mu}$ (取 $c=0$). 于是
\[
\P_x(X_\tau=r)=\frac{e^{-2\mu\ell}-e^{-2\mu x}}{e^{-2\mu\ell}-e^{-2\mu r}}.
\]
设 $\mu<0$, $x=0$, $r=a$, 令 $\ell\to-\infty$, 则 $e^{-2\mu\ell}\to0$, 右端趋于 $e^{2\mu a}$: 这就是定理~\ref{thm:drift-hitting} 中 $\P(\tau^\mu_a<\infty)=e^{-2\abs{\mu}a}$.

(3) 几何布朗运动 $\dd S=\mu S\dd t+\sigma S\dd B$ (在 $(0,\infty)$ 上): $\frac{2b}{\sigma^2}=\frac{2\mu}{\sigma^2y}$, $s'(y)=y^{-2\mu/\sigma^2}$ (取 $c=1$). 若 $\mu<\frac{1}{2}\sigma^2$, 则 $1-\frac{2\mu}{\sigma^2}>0$, $s(y)=\frac{y^{1-2\mu/\sigma^2}-1}{1-2\mu/\sigma^2}$, $s(0+)$ 有限. 令 $\ell\downarrow0$, 从 $x$ 出发在到达 $r$ 之前 ``跌到 $0$ 附近'' 的概率趋于 $1-\big(\frac{x}{r}\big)^{1-2\mu/\sigma^2}$. 当 $r\to\infty$ 时它趋于 $1$, 这与命题~\ref{prop:gbm} (4) 中 $\mu<\frac{1}{2}\sigma^2$ 时 $S_t\to0$ 一致.
\end{example}

\begin{remark}[边界分类与常返性]
尺度函数也决定了一维扩散的长期行为. 若 $\sigma>0$ 处处成立, 则扩散 $X$ (在 $\R$ 上) 常返 (即以概率 $1$ 无穷多次到达每一点) 当且仅当 $s(-\infty)=-\infty$ 且 $s(+\infty)=+\infty$; 若 $s(+\infty)<\infty$, 则以正概率 $X_t\to+\infty$. 更精细的 Feller 边界分类还需要 \textbf{速度测度} $m(\dd x)=\frac{2\dd x}{\sigma(x)^2s'(x)}$, 它满足 $\E_x\tau=\int_\ell^rG(x,y)\,m(\dd y)$, 其中 $G$ 是 Green 函数 (见~\cite{karatzas} 5.5 节). 例如 OU 过程的尺度函数 $s'(y)=e^{\kappa y^2/\sigma^2}$ (取 $\theta=0$) 在 $\pm\infty$ 处的积分发散, 所以 OU 过程常返.
\end{remark}

% =====================================================================
\section{随机微分方程与偏微分方程}
\label{sec:6-4}

向后方程 (定理~\ref{thm:backward}) 说明: 抛物型偏微分方程 $\partial_tu=Lu$ 的解可以表示为扩散过程的期望. 本节把这种 ``概率表示'' 推广到椭圆型方程的边值问题 (Dirichlet 问题) 与带位势和源项的抛物型方程 (Feynman-Kac 公式), 并以金融数学中的期权定价作为综合应用.

% ---------------------------------------------------------------------
\subsection{Dirichlet 问题}

设 $D\subseteq\R^d$ 是有界开集, $g$ 是 $\partial D$ 上的连续函数. \textbf{Dirichlet 问题}是求 $u\in C^2(D)\cap C(\overline{D})$, 使得
\[
\Delta u=0\quad\text{于 }D,\qquad u=g\quad\text{于 }\partial D.
\]
设 $\rvec{W}$ 是从 $\rvec{x}\in D$ 出发的 $d$ 维布朗运动 (概率与期望记为 $\P_{\rvec{x}}$, $\E_{\rvec{x}}$), $\tau_D=\inf\{t:\rvec{W}_t\notin D\}$.

\begin{theorem}[Dirichlet 问题的概率表示]\label{thm:dirichlet}
设 $k$ 是 $\overline{D}$ 上的有界连续函数. 若 $u\in C^2(D)\cap C(\overline{D})$ 满足
\[
\frac{1}{2}\Delta u=-k\quad\text{于 }D,\qquad u=g\quad\text{于 }\partial D,
\]
则 $u(\rvec{x})=\E_{\rvec{x}}\big[g(\rvec{W}_{\tau_D})+\int_0^{\tau_D}k(\rvec{W}_s)\dd s\big]$. 特别地 ($k=0$), 调和函数 $u$ 满足 $u(\rvec{x})=\E_{\rvec{x}}g(\rvec{W}_{\tau_D})$.
\end{theorem}

\begin{proof}
设 $D$ 包含于半径为 $R$ 的球中. 与定理~\ref{thm:exit} (1) 相同 (考察第一个分量离开 $(-R,R)$ 的时刻), $\E_{\rvec{x}}\tau_D<\infty$. 令 $D_n=\{\rvec{y}\in D:\operatorname{dist}(\rvec{y},\partial D)>\frac{1}{n}\}$, $\tau_n=\inf\{t:\rvec{W}_t\notin D_n\}$. $u$ 在 $\overline{D_n}$ 的一个邻域上是 $C^2$ 的 (把它在此邻域之外修改为 $\R^d$ 上的 $C^2$ 函数, 不影响下面的论证). 由多维 It\^o 公式~\ref{thm:ito-multi},
\[
u(\rvec{W}_{t\wedge\tau_n})=u(\rvec{x})+\int_0^{t\wedge\tau_n}\frac{1}{2}\Delta u(\rvec{W}_s)\dd s+\sum_{i=1}^d\int_0^t\1_{\{s\le\tau_n\}}\partial_iu(\rvec{W}_s)\dd W^i_s.
\]
被积过程 $\1_{\{s\le\tau_n\}}\partial_iu(\rvec{W}_s)$ 有界, 随机积分的期望为零, 所以
\[
u(\rvec{x})=\E_{\rvec{x}}\Big[u(\rvec{W}_{t\wedge\tau_n})+\int_0^{t\wedge\tau_n}k(\rvec{W}_s)\dd s\Big].
\]
先令 $t\to\infty$, 再令 $n\to\infty$. 由轨道连续性, $\tau_n\uparrow\tau_D$, $\rvec{W}_{\tau_n}\to\rvec{W}_{\tau_D}\in\partial D$. 由于 $u$ 在 $\overline{D}$ 上连续有界, $k$ 有界且 $\E_{\rvec{x}}\tau_D<\infty$, 由控制收敛定理即得结论.
\end{proof}

\begin{corollary}\label{cor:dirichlet}
\begin{enumerate}
\item (唯一性与最大值原理) Dirichlet 问题至多有一个解, 且解满足 $\min_{\partial D}g\le u\le\max_{\partial D}g$.
\item (平均值性质) 若 $u$ 在包含闭球 $\overline{B(\rvec{x},\rho)}$ 的开集上调和, 则 $u(\rvec{x})$ 等于 $u$ 在球面 $\partial B(\rvec{x},\rho)$ 上 (关于面积测度) 的平均值.
\item (球内的期望出口时间) 对 $D=B(\rvec{0},R)$, $\E_{\rvec{x}}\tau_D=\frac{R^2-\abs{\rvec{x}}^2}{d}$.
\end{enumerate}
\end{corollary}

\begin{proof}
(1) 由表示公式立即得到. (2) 由布朗运动的旋转不变性 (命题~\ref{prop:bm-rotation}), $\rvec{W}_{\tau}$ ($\tau$ 为离开球的时刻) 的分布在绕 $\rvec{x}$ 的旋转下不变, 所以是球面上的均匀分布. (3) 取 $u(\rvec{y})=\frac{R^2-\abs{\rvec{y}}^2}{d}$, 则 $\frac{1}{2}\Delta u=\frac{1}{2}\cdot\frac{-2d}{d}=-1$, 在球面上 $u=0$; 对 $k\equiv1$, $g\equiv0$ 应用定理.
\end{proof}

定理~\ref{thm:dirichlet} 假设了解的存在性. 反过来, $u(\rvec{x}):=\E_{\rvec{x}}g(\rvec{W}_{\tau_D})$ 总是 $D$ 内的调和函数 (由强马尔可夫性与平均值性质), 但它在边界上不一定连续. 若边界满足适当的正则性 (例如每个边界点处都有位于 $D$ 外部的小圆锥), 则 $u$ 在边界上连续, 从而是 Dirichlet 问题的解 (见~\cite{karatzas} 4.2 节). 正则性不能去掉: 设 $D=\{\rvec{y}\in\R^2:0<\abs{\rvec{y}}<1\}$ 是去心圆盘, 在单位圆上 $g=0$, 在原点 $g(\rvec{0})=1$. 由例~\ref{ex:bm-annulus}, 平面布朗运动不击中原点, 所以 $\E_{\rvec{x}}g(\rvec{W}_{\tau_D})=0$; 而 $u\equiv0$ 在原点处不等于 $g(\rvec{0})=1$, 由唯一性, 这个 Dirichlet 问题无解. 原点是一个 ``非正则'' 边界点: 布朗运动从它附近出发, 几乎必然不会从它离开.

% ---------------------------------------------------------------------
\subsection{Feynman-Kac 公式}

\begin{theorem}[Feynman-Kac 公式]\label{thm:feynman-kac}
设 $b,\sigma$ 满足 \eqref{eq:H1}\eqref{eq:H2}, $L$ 如 \eqref{eq:L-operator}, $c\ge0$ 与 $g$ 是 $\R$ 上的连续函数, $f$ 连续. 设 $v\in C^{1,2}([0,T)\times\R)\cap C([0,T]\times\R)$ 满足
\begin{equation}\label{eq:feynman-kac-pde}
\partial_tv+Lv-cv+g=0\quad\text{于 }[0,T)\times\R,\qquad v(T,x)=f(x),
\end{equation}
并设存在常数 $C$, 使 $\abs{v(t,x)}+\abs{g(x)}\le C(1+x^2)$. 则对 $(t,x)\in[0,T]\times\R$,
\[
v(t,x)=\E\Big[f(X^{t,x}_T)e^{-\int_t^Tc(X^{t,x}_u)\dd u}+\int_t^Tg(X^{t,x}_s)e^{-\int_t^sc(X^{t,x}_u)\dd u}\dd s\Big].
\]
\end{theorem}

\begin{proof}
记 $X=X^{t,x}$, $\Gamma_s=e^{-\int_t^sc(X_u)\dd u}\in(0,1]$, 并令
\[
Y_s=v(s,X_s)\Gamma_s+\int_t^sg(X_u)\Gamma_u\dd u,\qquad s\in[t,T].
\]
由 It\^o 公式与乘积法则 ($\dd\Gamma_s=-c(X_s)\Gamma_s\dd s$ 是有限变差的),
\[
\dd Y_s=\Gamma_s\big(\partial_tv+Lv-cv+g\big)(s,X_s)\dd s+\Gamma_s\,\partial_xv(s,X_s)\sigma(X_s)\dd B_s=\Gamma_s\,\partial_xv(s,X_s)\sigma(X_s)\dd B_s.
\]
令 $\tau_n=\inf\{s\ge t:\abs{X_s}\ge n\}\wedge(T-\frac{1}{n})$. 在 $[t,\tau_n]$ 上被积过程有界, 所以停止后的随机积分是鞅, $v(t,x)=Y_t=\E Y_{\tau_n}$. 令 $n\to\infty$: 由 $v$ 在 $[0,T]\times\R$ 上的连续性, $Y_{\tau_n}\to Y_T$; 又由增长条件,
\[
\abs{Y_{\tau_n}}\le C\Big(1+\sup_{s\le T}X_s^2\Big)+TC\Big(1+\sup_{s\le T}X_s^2\Big),
\]
由矩估计 \eqref{eq:sde-moment} 右端可积. 由控制收敛定理, $v(t,x)=\E Y_T$, 这就是所述公式 ($v(T,X_T)=f(X_T)$).
\end{proof}

\begin{intuition}
把 $c(x)$ 理解为粒子在位置 $x$ 处被 ``杀死'' (吸收) 的速率, $g(x)$ 理解为单位时间的 ``收益'', $f$ 理解为终端 ``奖励''. 则 $v(t,x)$ 是从 $(t,x)$ 出发、在被杀死之前累积的期望总收益. 在金融中, $c=r$ 是贴现率, $f$ 是期权的收益, $v$ 就是期权的价格. Feynman 最初用路径积分 (形式上的 ``Wiener 测度上的积分'') 表示 Schr\"odinger 方程的解; Kac 意识到对热方程这可以严格化, 并用它计算了布朗运动的占位时分布 (例如由 Feynman-Kac 公式可以推出 4.5.4 小节提到的第二个反正弦律).
\end{intuition}

\begin{remark}
(1) 定理同样假设了解的存在性. 若 $b,\sigma,c,g,f$ 足够光滑, 则右端的期望确实是 $C^{1,2}$ 的, 从而是 \eqref{eq:feynman-kac-pde} 的解 (见~\cite{karatzas} 定理 5.7.6). (2) 增长条件可以放宽为多项式增长, 此时需要解的高阶矩估计 $\E\sup_t\abs{X_t}^{2m}<\infty$, 它可以由 BDG 不等式~\ref{thm:bdg} 与 Gronwall 不等式得到.
\end{remark}

% ---------------------------------------------------------------------
\subsection{Black-Scholes 理论}

考虑一个由两种资产构成的市场: 无风险债券 $\beta_t=e^{rt}$ ($r\ge0$ 为利率) 与股票, 后者的价格满足
\[
\dd S_t=\mu S_t\dd t+\sigma S_t\dd B_t,\qquad S_0>0,
\]
$\{\mathcal{F}_t\}$ 是 $B$ 的扩充信息流. 一个\textbf{交易策略}是一对循序可测过程 $(\varphi_t,\psi_t)$, 分别表示持有的股票与债券的份数, 其价值为 $V_t=\varphi_tS_t+\psi_t\beta_t$. 若
\begin{equation}\label{eq:self-financing}
\dd V_t=\varphi_t\dd S_t+\psi_t\dd\beta_t,
\end{equation}
即价值的变化完全来自资产价格的变化 (没有外部资金的注入或抽出), 称策略是\textbf{自融资}的. 记贴现价格 $\widetilde{S}_t=e^{-rt}S_t$, $\widetilde{V}_t=e^{-rt}V_t$. 由乘积法则, 自融资条件等价于
\begin{equation}\label{eq:self-financing-discounted}
\dd\widetilde{V}_t=e^{-rt}\big(\dd V_t-rV_t\dd t\big)=e^{-rt}\varphi_t\big(\dd S_t-rS_t\dd t\big)=\varphi_t\dd\widetilde{S}_t.
\end{equation}

\emph{风险中性测度.} 令 $\theta=-\frac{\mu-r}{\sigma}$ (常数), $\widetilde{\P}=e^{\theta B_T-\frac{1}{2}\theta^2T}\cdot\P$. 由 Girsanov 定理~\ref{thm:girsanov}, $\widetilde{B}_t=B_t+\frac{\mu-r}{\sigma}t$ 是 $\widetilde{\P}$ 下的布朗运动, 且
\[
\dd S_t=rS_t\dd t+\sigma S_t\dd\widetilde{B}_t,\qquad\dd\widetilde{S}_t=\sigma\widetilde{S}_t\dd\widetilde{B}_t.
\]
所以在 $\widetilde{\P}$ 下, 股票的平均收益率等于无风险利率, 贴现价格 $\widetilde{S}$ 是鞅 (它是指数鞅). 由 \eqref{eq:self-financing-discounted}, 若 $\widetilde{\E}\int_0^T\varphi_t^2\widetilde{S}_t^2\dd t<\infty$ (称这样的策略为\textbf{容许}的), 则贴现价值 $\widetilde{V}$ 是 $\widetilde{\P}$-鞅. 这推出\textbf{无套利}: 不存在容许的自融资策略使 $V_0=0$, $V_T\ge0$ 且 $\P(V_T>0)>0$, 因为 $\widetilde{\E}\widetilde{V}_T=\widetilde{V}_0=0$ 迫使 $V_T=0$ $\widetilde{\P}$-$\as$, 从而 $\P$-$\as$ ($\widetilde{\P}$ 与 $\P$ 等价).

\begin{theorem}[Black-Scholes]\label{thm:black-scholes}
设欧式衍生品在时刻 $T$ 的收益为 $F=h(S_T)$, $h\ge0$ 连续且 $h(s)\le C(1+s)$.
\begin{enumerate}
\item (复制) 存在容许的自融资策略, 其价值满足 $V_T=F$. 任何这样的策略在时刻 $t$ 的价值都是
\begin{align*}
V_t&=e^{-r(T-t)}\widetilde{\E}\big[F\,\big|\,\mathcal{F}_t\big]=v(t,S_t),\\
v(t,s)&=e^{-r(T-t)}\E\,h\Big(s\,e^{(r-\frac{1}{2}\sigma^2)(T-t)+\sigma\sqrt{T-t}\,Z}\Big),
\end{align*}
其中 $Z\sim N(0,1)$. $V_0=v(0,S_0)$ 称为衍生品的 (无套利) 价格. 它与 $\mu$ 无关.
\item (看涨期权) 对 $h(s)=(s-K)^+$,
\[
v(t,s)=s\,\Phi(d_+)-Ke^{-r(T-t)}\Phi(d_-),\quad d_\pm=\frac{\log(s/K)+(r\pm\frac{1}{2}\sigma^2)(T-t)}{\sigma\sqrt{T-t}}.
\]
\item (偏微分方程与对冲) 若 $v\in C^{1,2}([0,T)\times(0,\infty))$ (对看涨期权由 (2) 显然), 则 $v$ 满足 \textbf{Black-Scholes 方程}
\[
\partial_tv+rs\,\partial_sv+\frac{1}{2}\sigma^2s^2\,\partial_{ss}v-rv=0,\qquad v(T,s)=h(s),
\]
且复制策略的股票头寸为 $\varphi_t=\partial_sv(t,S_t)$ (\textbf{Delta 对冲}).
\end{enumerate}
\end{theorem}

\begin{proof}
(1) 由于 $S_T$ 是对数正态的, $F\in L^2(\widetilde{\P})$. $B$ 与 $\widetilde{B}$ 只相差一个确定的函数, 所以它们生成相同的扩充信息流. 对 $\widetilde{\P}$ 与 $\widetilde{\P}$-布朗运动 $\widetilde{B}$ 应用鞅表示定理 (推论~\ref{cor:mart-rep}), 鞅 $M_t=\widetilde{\E}[e^{-rT}F\,|\,\mathcal{F}_t]$ 可以写成 $M_t=M_0+\int_0^tH_s\dd\widetilde{B}_s$, $H\in\mathscr{L}^2_T$ (关于 $\widetilde{\P}$). 令
\[
\varphi_t=\frac{H_t}{\sigma\widetilde{S}_t},\qquad\psi_t=M_t-\varphi_t\widetilde{S}_t.
\]
则 $\widetilde{V}_t=e^{-rt}(\varphi_tS_t+\psi_te^{rt})=\varphi_t\widetilde{S}_t+\psi_t=M_t$, 且 $\dd\widetilde{V}_t=H_t\dd\widetilde{B}_t=\varphi_t\sigma\widetilde{S}_t\dd\widetilde{B}_t=\varphi_t\dd\widetilde{S}_t$, 即策略自融资、容许, 且 $V_T=e^{rT}M_T=F$. 若另一个容许的自融资策略也复制 $F$, 其贴现价值是 $\widetilde{\P}$-鞅, 终值为 $e^{-rT}F$, 所以它等于 $M_t$. 最后, $S_T=S_t\exp\big((r-\frac{1}{2}\sigma^2)(T-t)+\sigma(\widetilde{B}_T-\widetilde{B}_t)\big)$, 其中 $\widetilde{B}_T-\widetilde{B}_t\sim N(0,T-t)$ 与 $\mathcal{F}_t$ 独立, 而 $S_t$ 是 $\mathcal{F}_t$-可测的, 由代入引理~\ref{lem:freezing}, $V_t=e^{rt}M_t=v(t,S_t)$.

(2) 记 $\tau=T-t$, $X=s\,e^{(r-\frac{1}{2}\sigma^2)\tau+\sigma\sqrt{\tau}Z}$. $X>K$ 当且仅当 $Z>-d_-$. 由配方 $\E[e^{aZ};Z>-d]=e^{a^2/2}\P(Z>-d-a)=e^{a^2/2}\Phi(d+a)$ (取 $a=\sigma\sqrt{\tau}$, 并注意 $d_-+\sigma\sqrt{\tau}=d_+$),
\[
\E\big[(X-K)^+\big]=s\,e^{(r-\frac{1}{2}\sigma^2)\tau}e^{\frac{1}{2}\sigma^2\tau}\Phi(d_+)-K\,\Phi(d_-)=se^{r\tau}\Phi(d_+)-K\Phi(d_-).
\]
乘以 $e^{-r\tau}$ 即得.

(3) $e^{-rt}v(t,S_t)=M_t$ 是 $\widetilde{\P}$-鞅. 对 $t<T$ 应用 It\^o 公式 (在 $\widetilde{\P}$ 下, $\dd S=rS\dd t+\sigma S\dd\widetilde{B}$),
\[
\dd\big(e^{-rt}v(t,S_t)\big)=e^{-rt}G(t,S_t)\dd t+e^{-rt}\partial_sv(t,S_t)\sigma S_t\dd\widetilde{B}_t,
\]
其中 $G=\partial_tv+rs\,\partial_sv+\frac{1}{2}\sigma^2s^2\partial_{ss}v-rv$.
于是 $\int_0^te^{-ru}G(u,S_u)\dd u$ 是连续的有限变差过程, 又等于两个连续局部鞅之差, 由 It\^o 过程分解的唯一性 (命题~\ref{prop:ito-process} (1)), 它恒为零, 且 $H_t=e^{-rt}\partial_sv(t,S_t)\sigma S_t$, 即 $\varphi_t=\partial_sv(t,S_t)$. 若在某点 $(t_0,s_0)$ 处 $G\ne0$, 由连续性, $G$ 在某个邻域 $U$ 内不为零 (不妨设为正); 而 $S_u$ 对每个 $u>0$ 有处处为正的密度, 所以 $\widetilde{\E}\int_0^{T}\1_U(u,S_u)\dd u>0$; 而 $\int_0^te^{-ru}G(u,S_u)\dd u\equiv0$ 推出 $G(u,S_u)=0$ 对几乎每个 $u$ 几乎必然成立, 矛盾. 所以 $G\equiv0$.
\end{proof}

\begin{figure}[htbp]
\centering
\begin{tikzpicture}
\begin{axis}[SAaxis, axis x line=bottom, axis y line=left, width=12cm, height=5.8cm, xmin=50, xmax=162, ymin=0, ymax=65,
  xtick={60,80,100,120,140,160}, ytick={0,20,40,60}, xlabel={$s$},
  legend style={font=\small,draw=none,fill=none,at={(0.02,0.98)},anchor=north west}, legend cell align=left]
  \addplot[SADark,thick,dashed] coordinates {(50,0) (100,0) (160,60)};
  \addplot[SABlue,thick] coordinates {(50.00,0.00) (51.00,0.00) (52.00,0.00) (53.00,0.00) (54.00,0.00) (55.00,0.00) (56.00,0.00) (57.00,0.00) (58.00,0.00) (59.00,0.00) (60.00,0.00) (61.00,0.00) (62.00,0.00) (63.00,0.00) (64.00,0.01) (65.00,0.01) (66.00,0.01) (67.00,0.02) (68.00,0.03) (69.00,0.03) (70.00,0.05) (71.00,0.06) (72.00,0.08) (73.00,0.10) (74.00,0.13) (75.00,0.17) (76.00,0.21) (77.00,0.26) (78.00,0.33) (79.00,0.40) (80.00,0.48) (81.00,0.58) (82.00,0.70) (83.00,0.83) (84.00,0.98) (85.00,1.14) (86.00,1.33) (87.00,1.54) (88.00,1.77) (89.00,2.03) (90.00,2.31) (91.00,2.61) (92.00,2.95) (93.00,3.31) (94.00,3.69) (95.00,4.11) (96.00,4.55) (97.00,5.01) (98.00,5.51) (99.00,6.03) (100.00,6.58) (101.00,7.16) (102.00,7.76) (103.00,8.39) (104.00,9.04) (105.00,9.71) (106.00,10.41) (107.00,11.13) (108.00,11.87) (109.00,12.63) (110.00,13.40) (111.00,14.20) (112.00,15.01) (113.00,15.84) (114.00,16.68) (115.00,17.53) (116.00,18.40) (117.00,19.28) (118.00,20.17) (119.00,21.08) (120.00,21.99) (121.00,22.91) (122.00,23.83) (123.00,24.77) (124.00,25.71) (125.00,26.66) (126.00,27.61) (127.00,28.57) (128.00,29.53) (129.00,30.50) (130.00,31.47) (131.00,32.44) (132.00,33.41) (133.00,34.39) (134.00,35.38) (135.00,36.36) (136.00,37.34) (137.00,38.33) (138.00,39.32) (139.00,40.31) (140.00,41.30) (141.00,42.29) (142.00,43.29) (143.00,44.28) (144.00,45.28) (145.00,46.27) (146.00,47.27) (147.00,48.26) (148.00,49.26) (149.00,50.26) (150.00,51.26) (151.00,52.25) (152.00,53.25) (153.00,54.25) (154.00,55.25) (155.00,56.25) (156.00,57.25) (157.00,58.25) (158.00,59.25) (159.00,60.25) (160.00,61.25)};
  \addplot[SARed,thick] coordinates {(50.00,0.12) (51.00,0.14) (52.00,0.17) (53.00,0.20) (54.00,0.24) (55.00,0.29) (56.00,0.34) (57.00,0.39) (58.00,0.46) (59.00,0.53) (60.00,0.61) (61.00,0.69) (62.00,0.79) (63.00,0.90) (64.00,1.01) (65.00,1.14) (66.00,1.27) (67.00,1.42) (68.00,1.58) (69.00,1.76) (70.00,1.94) (71.00,2.14) (72.00,2.35) (73.00,2.58) (74.00,2.81) (75.00,3.07) (76.00,3.33) (77.00,3.62) (78.00,3.91) (79.00,4.23) (80.00,4.55) (81.00,4.90) (82.00,5.25) (83.00,5.63) (84.00,6.01) (85.00,6.42) (86.00,6.84) (87.00,7.27) (88.00,7.72) (89.00,8.18) (90.00,8.66) (91.00,9.15) (92.00,9.66) (93.00,10.19) (94.00,10.72) (95.00,11.27) (96.00,11.84) (97.00,12.42) (98.00,13.01) (99.00,13.61) (100.00,14.23) (101.00,14.86) (102.00,15.50) (103.00,16.16) (104.00,16.83) (105.00,17.51) (106.00,18.19) (107.00,18.90) (108.00,19.61) (109.00,20.33) (110.00,21.06) (111.00,21.80) (112.00,22.55) (113.00,23.32) (114.00,24.09) (115.00,24.86) (116.00,25.65) (117.00,26.45) (118.00,27.25) (119.00,28.06) (120.00,28.88) (121.00,29.71) (122.00,30.54) (123.00,31.38) (124.00,32.23) (125.00,33.08) (126.00,33.94) (127.00,34.80) (128.00,35.67) (129.00,36.55) (130.00,37.43) (131.00,38.31) (132.00,39.20) (133.00,40.10) (134.00,41.00) (135.00,41.90) (136.00,42.81) (137.00,43.72) (138.00,44.64) (139.00,45.56) (140.00,46.48) (141.00,47.41) (142.00,48.34) (143.00,49.27) (144.00,50.21) (145.00,51.14) (146.00,52.09) (147.00,53.03) (148.00,53.98) (149.00,54.92) (150.00,55.88) (151.00,56.83) (152.00,57.79) (153.00,58.74) (154.00,59.70) (155.00,60.66) (156.00,61.63) (157.00,62.59) (158.00,63.56) (159.00,64.53) (160.00,65.50)};
  \legend{收益 $(s-K)^+$, $T-t=0.25$, $T-t=1$}
\end{axis}
\end{tikzpicture}
\caption{看涨期权的 Black-Scholes 价格 $v(t,s)$ ($K=100$, $r=0.05$, $\sigma=0.3$). 离到期越远, 价格越高于内在价值 $(s-K)^+$; 当 $t\to T$ 时价格曲线收敛到收益函数. 曲线的斜率 $\partial_sv=\Phi(d_+)\in(0,1)$ 就是对冲所需持有的股票份数}
\label{fig:black-scholes}
\end{figure}

\begin{intuition}
Black-Scholes 理论最令人意外的结论, 是期权价格与股票的期望收益率 $\mu$ 无关. 原因是期权可以用股票和债券完全复制, 复制成本只取决于股价的\emph{波动} (二次变差 $\sigma^2S^2\dd t$ 是逐条轨道的性质, 与概率测度无关), 而与股价的\emph{趋势}无关. Girsanov 定理恰好只改变趋势而不改变波动, 所以可以在一个 ``假想'' 的测度 $\widetilde{\P}$ 下计算, 在那里所有资产的平均收益率都是 $r$, 期权价格就是贴现收益的期望. 鞅表示定理则保证复制总是可能的 (市场完备).
\end{intuition}

由 Feynman-Kac 公式 (取 $c=r$, $g=0$, $L$ 为 $\widetilde{\P}$ 下 $S$ 的生成元 $rs\partial_s+\frac{1}{2}\sigma^2s^2\partial_{ss}$), Black-Scholes 方程与期望表示是同一件事的两个侧面: 前者是 ``对冲'' 的观点 (局部无风险组合的收益率等于 $r$), 后者是 ``风险中性定价'' 的观点.

% ---------------------------------------------------------------------
\subsection{障碍期权}

\textbf{障碍期权}的收益依赖于股价是否触及某个水平. 例如\textbf{向上敲出看涨期权} (up-and-out call) 的收益为
\[
F=(S_T-K)^+\1_{\{\max_{t\le T}S_t<H\}},\qquad S_0<H,\ K<H:
\]
若股价在到期前曾经涨到障碍 $H$, 期权作废. 它比普通看涨期权便宜, 适合预期股价温和上涨的投资者. 由于 $F$ 依赖于整条轨道, 定价需要 $(\max S,S_T)$ 的联合分布. 在 $\widetilde{\P}$ 下,
\[
S_t=S_0e^{\sigma\widehat{W}_t},\qquad\widehat{W}_t=\widetilde{B}_t+\nu t,\qquad\nu=\frac{r-\frac{1}{2}\sigma^2}{\sigma},
\]
所以问题归结为带漂移布朗运动的最大值与终点的联合分布. 第~4~章的反射原理处理了无漂移的情形, Girsanov 定理可以把漂移加上.

\begin{proposition}[带漂移布朗运动的镜像原理]\label{prop:image}
设 $W$ 是标准布朗运动, $\widehat{W}_t=W_t+\nu t$, $\widehat{M}_T=\max_{t\le T}\widehat{W}_t$, $h>0$. 则对任意满足 $f(x)=0$ ($x\ge h$) 的有界 Borel 函数 $f$,
\[
\E\big[f(\widehat{W}_T);\widehat{M}_T<h\big]=\E f(\widehat{W}_T)-e^{2\nu h}\,\E f(\widehat{W}_T+2h).
\]
\end{proposition}

\begin{proof}
由 Girsanov 定理 (取 $\theta\equiv\nu$), 在测度 $Z_T\cdot\P$ ($Z_T=e^{\nu W_T-\frac{1}{2}\nu^2T}$) 下 $W$ 是漂移为 $\nu$ 的布朗运动, 所以对轨道的有界可测泛函 $\Psi$, $\E\Psi(\widehat{W})=\E[\Psi(W)Z_T]$. 记 $M_T=\max_{t\le T}W_t$. 由反射原理 \eqref{eq:reflection-joint}, 对 $x\le h$, $\P(M_T\ge h,W_T\le x)=\P(W_T\ge2h-x)$, 即 $\{M_T\ge h\}$ 上 $W_T$ 在 $(-\infty,h]$ 中的 (次) 密度为 $\varphi_T(2h-x)$, $\varphi_T(y)=\frac{1}{\sqrt{2\pi T}}e^{-y^2/(2T)}$. 由于 $f$ 在 $[h,\infty)$ 上为零,
\begin{align*}
\E\big[f(\widehat{W}_T);\widehat{M}_T\ge h\big]&=\E\big[f(W_T)e^{\nu W_T-\frac{1}{2}\nu^2T};M_T\ge h\big]\\
&=\int_{-\infty}^hf(x)e^{\nu x-\frac{1}{2}\nu^2T}\varphi_T(2h-x)\dd x.
\end{align*}
展开平方可以验证恒等式
\[
e^{\nu x-\frac{1}{2}\nu^2T}\varphi_T(2h-x)=e^{2\nu h}\varphi_T(x-2h-\nu T),
\]
(两端的指数都等于 $-\frac{(2h-x)^2}{2T}+\nu x-\frac{1}{2}\nu^2T$), 而 $\varphi_T(y-\nu T)$ 正是 $\widehat{W}_T\sim N(\nu T,T)$ 的密度. 所以上式等于 $e^{2\nu h}\int f(x)\varphi_T(x-2h-\nu T)\dd x=e^{2\nu h}\E f(\widehat{W}_T+2h)$ (积分可以扩展到全体 $x$, 因为 $f$ 在 $[h,\infty)$ 上为零). 用 $\E f(\widehat{W}_T)$ 减去它即得结论.
\end{proof}

\begin{theorem}[向上敲出看涨期权]\label{thm:barrier}
设看涨期权的初始股价为 $s$, 执行价为 $k$, 到期时间为 $T$, 记其 Black-Scholes 价格为 $C(s,k)$ (定理~\ref{thm:black-scholes} (2), $t=0$), 并令
\[
U(s)=C(s,K)-C(s,H)-(H-K)e^{-rT}\Phi\big(d_-(s,H)\big),
\]
其中 $d_-(s,H)=\frac{\log(s/H)+(r-\frac{1}{2}\sigma^2)T}{\sigma\sqrt{T}}$. 则向上敲出看涨期权在时刻 $0$ 的价格为
\[
V_0=U(S_0)-\Big(\frac{H}{S_0}\Big)^{\frac{2r}{\sigma^2}-1}U\Big(\frac{H^2}{S_0}\Big).
\]
\end{theorem}

\begin{proof}
与定理~\ref{thm:black-scholes} (1) 相同 ($F$ 有界, 鞅表示定理仍适用), 价格为 $V_0=e^{-rT}\widetilde{\E}F$. 令 $g(s)=(s-K)^+\1_{\{s<H\}}$, $h=\frac{1}{\sigma}\log\frac{H}{S_0}>0$, $f(x)=g(S_0e^{\sigma x})$, 则 $f(x)=0$ 当 $x\ge h$, 且 $\{\max S<H\}=\{\widehat{M}_T<h\}$. 由命题~\ref{prop:image},
\[
\widetilde{\E}F=\widetilde{\E}g(S_T)-e^{2\nu h}\,\widetilde{\E}\,g\big(S_0e^{\sigma\widehat{W}_T+2\sigma h}\big)=\widetilde{\E}g(S_T)-\Big(\frac{H}{S_0}\Big)^{2\nu/\sigma}\widetilde{\E}\,g\Big(\frac{H^2}{S_0}e^{\sigma\widehat{W}_T}\Big),
\]
因为 $e^{\sigma h}=\frac{H}{S_0}$, $S_0e^{2\sigma h}=\frac{H^2}{S_0}$. 这里 $\frac{2\nu}{\sigma}=\frac{2r}{\sigma^2}-1$, 而 $\frac{H^2}{S_0}e^{\sigma\widehat{W}_T}$ 正是从 $\frac{H^2}{S_0}$ 出发的 (风险中性) 股价在 $T$ 时刻的值. 最后, 由 $g(s)=(s-K)^+-(s-H)^+-(H-K)\1_{\{s\ge H\}}$ 与 $\widetilde{\P}(S_T\ge H\,|\,S_0=s)=\Phi(d_-(s,H))$ (定理~\ref{thm:black-scholes} 证明中的计算), 从 $s$ 出发时 $e^{-rT}\widetilde{\E}g(S_T)=U(s)$.
\end{proof}

\begin{figure}[htbp]
\centering
\begin{tikzpicture}
\begin{axis}[SAaxis, axis x line=bottom, axis y line=left, width=12cm, height=5.8cm, xmin=50, xmax=136, ymin=0, ymax=6,
  xtick={60,80,100,120}, ytick={0,1,2,3,4,5,6}, xlabel={$S_0$}, clip=false,
  legend style={font=\small,draw=none,fill=none,at={(0.02,0.98)},anchor=north west}, legend cell align=left]
  \addplot[SADark!60,thick,dashed,restrict y to domain=0:6] coordinates {(50.00,0.12) (50.50,0.13) (51.00,0.14) (51.50,0.16) (52.00,0.17) (52.50,0.19) (53.00,0.20) (53.50,0.22) (54.00,0.24) (54.50,0.26) (55.00,0.29) (55.50,0.31) (56.00,0.34) (56.50,0.36) (57.00,0.39) (57.50,0.42) (58.00,0.46) (58.50,0.49) (59.00,0.53) (59.50,0.57) (60.00,0.61) (60.50,0.65) (61.00,0.69) (61.50,0.74) (62.00,0.79) (62.50,0.84) (63.00,0.90) (63.50,0.95) (64.00,1.01) (64.50,1.07) (65.00,1.14) (65.50,1.20) (66.00,1.27) (66.50,1.35) (67.00,1.42) (67.50,1.50) (68.00,1.58) (68.50,1.67) (69.00,1.76) (69.50,1.85) (70.00,1.94) (70.50,2.04) (71.00,2.14) (71.50,2.24) (72.00,2.35) (72.50,2.46) (73.00,2.58) (73.50,2.69) (74.00,2.81) (74.50,2.94) (75.00,3.07) (75.50,3.20) (76.00,3.33) (76.50,3.47) (77.00,3.62) (77.50,3.76) (78.00,3.91) (78.50,4.07) (79.00,4.23) (79.50,4.39) (80.00,4.55) (80.50,4.72) (81.00,4.90) (81.50,5.07) (82.00,5.25) (82.50,5.44) (83.00,5.63) (83.50,5.82) (84.00,6.01) (84.50,6.21) (85.00,6.42) (85.50,6.62) (86.00,6.84) (86.50,7.05) (87.00,7.27) (87.50,7.49) (88.00,7.72) (88.50,7.95) (89.00,8.18) (89.50,8.42) (90.00,8.66) (90.50,8.91) (91.00,9.15) (91.50,9.41) (92.00,9.66) (92.50,9.92) (93.00,10.19) (93.50,10.45) (94.00,10.72) (94.50,11.00) (95.00,11.27) (95.50,11.55) (96.00,11.84) (96.50,12.13) (97.00,12.42) (97.50,12.71) (98.00,13.01) (98.50,13.31) (99.00,13.61) (99.50,13.92) (100.00,14.23) (100.50,14.54) (101.00,14.86) (101.50,15.18) (102.00,15.50) (102.50,15.83) (103.00,16.16) (103.50,16.49) (104.00,16.83) (104.50,17.16) (105.00,17.51) (105.50,17.85) (106.00,18.19) (106.50,18.54) (107.00,18.90) (107.50,19.25) (108.00,19.61) (108.50,19.97) (109.00,20.33) (109.50,20.69) (110.00,21.06) (110.50,21.43) (111.00,21.80) (111.50,22.18) (112.00,22.55) (112.50,22.93) (113.00,23.32) (113.50,23.70) (114.00,24.09) (114.50,24.47) (115.00,24.86) (115.50,25.26) (116.00,25.65) (116.50,26.05) (117.00,26.45) (117.50,26.85) (118.00,27.25) (118.50,27.66) (119.00,28.06) (119.50,28.47) (120.00,28.88) (120.50,29.29) (121.00,29.71) (121.50,30.12) (122.00,30.54) (122.50,30.96) (123.00,31.38) (123.50,31.80) (124.00,32.23) (124.50,32.65) (125.00,33.08) (125.50,33.51) (126.00,33.94) (126.50,34.37) (127.00,34.80) (127.50,35.23) (128.00,35.67) (128.50,36.11) (129.00,36.55) (129.50,36.99) (130.00,37.43)};
  \addplot[SARed,thick] coordinates {(50.00,0.07) (50.50,0.08) (51.00,0.09) (51.50,0.09) (52.00,0.10) (52.50,0.11) (53.00,0.12) (53.50,0.13) (54.00,0.14) (54.50,0.15) (55.00,0.16) (55.50,0.17) (56.00,0.18) (56.50,0.19) (57.00,0.21) (57.50,0.22) (58.00,0.23) (58.50,0.25) (59.00,0.26) (59.50,0.28) (60.00,0.30) (60.50,0.31) (61.00,0.33) (61.50,0.35) (62.00,0.37) (62.50,0.38) (63.00,0.40) (63.50,0.42) (64.00,0.44) (64.50,0.46) (65.00,0.48) (65.50,0.51) (66.00,0.53) (66.50,0.55) (67.00,0.57) (67.50,0.59) (68.00,0.62) (68.50,0.64) (69.00,0.66) (69.50,0.69) (70.00,0.71) (70.50,0.74) (71.00,0.76) (71.50,0.79) (72.00,0.81) (72.50,0.83) (73.00,0.86) (73.50,0.88) (74.00,0.91) (74.50,0.93) (75.00,0.96) (75.50,0.98) (76.00,1.01) (76.50,1.03) (77.00,1.05) (77.50,1.08) (78.00,1.10) (78.50,1.12) (79.00,1.14) (79.50,1.17) (80.00,1.19) (80.50,1.21) (81.00,1.23) (81.50,1.25) (82.00,1.27) (82.50,1.29) (83.00,1.31) (83.50,1.33) (84.00,1.35) (84.50,1.36) (85.00,1.38) (85.50,1.39) (86.00,1.41) (86.50,1.42) (87.00,1.44) (87.50,1.45) (88.00,1.46) (88.50,1.47) (89.00,1.49) (89.50,1.49) (90.00,1.50) (90.50,1.51) (91.00,1.52) (91.50,1.53) (92.00,1.53) (92.50,1.54) (93.00,1.54) (93.50,1.54) (94.00,1.55) (94.50,1.55) (95.00,1.55) (95.50,1.55) (96.00,1.55) (96.50,1.54) (97.00,1.54) (97.50,1.54) (98.00,1.53) (98.50,1.53) (99.00,1.52) (99.50,1.51) (100.00,1.50) (100.50,1.49) (101.00,1.48) (101.50,1.47) (102.00,1.46) (102.50,1.45) (103.00,1.44) (103.50,1.42) (104.00,1.41) (104.50,1.39) (105.00,1.38) (105.50,1.36) (106.00,1.34) (106.50,1.32) (107.00,1.30) (107.50,1.28) (108.00,1.26) (108.50,1.24) (109.00,1.22) (109.50,1.20) (110.00,1.18) (110.50,1.15) (111.00,1.13) (111.50,1.10) (112.00,1.08) (112.50,1.05) (113.00,1.03) (113.50,1.00) (114.00,0.98) (114.50,0.95) (115.00,0.92) (115.50,0.89) (116.00,0.86) (116.50,0.84) (117.00,0.81) (117.50,0.78) (118.00,0.75) (118.50,0.72) (119.00,0.69) (119.50,0.66) (120.00,0.63) (120.50,0.60) (121.00,0.57) (121.50,0.53) (122.00,0.50) (122.50,0.47) (123.00,0.44) (123.50,0.41) (124.00,0.38) (124.50,0.35) (125.00,0.31) (125.50,0.28) (126.00,0.25) (126.50,0.22) (127.00,0.19) (127.50,0.16) (128.00,0.12) (128.50,0.09) (129.00,0.06) (129.50,0.03) (130.00,0.00)};
  \addplot[SADark!30,forget plot] coordinates {(130,0) (130,6)};
  \node[font=\small,SADark!70,anchor=south] at (axis cs:130,6) {$H$};
  \legend{普通看涨期权, 向上敲出看涨期权}
\end{axis}
\end{tikzpicture}
\caption{向上敲出看涨期权 ($K=100$, $H=130$, $T=1$, $r=0.05$, $\sigma=0.3$) 与普通看涨期权 (虚线, 超出纵轴范围的部分未画出) 的价格. 敲出条款使 $S_0=100$ 处的期权价格从普通期权的 $14.23$ 降至 $1.503$, 而且价格关于 $S_0$ 不再单调, 股价越接近障碍, 期权越可能作废}
\label{fig:barrier}
\end{figure}

\begin{remark}
(1) 这里假设障碍是\emph{连续}监测的. 实际合约常在离散时刻 (例如每个交易日收盘) 检查障碍, 此时价格略高; 6.5 节将看到, 用 Monte Carlo 方法逐步模拟时, 离散监测带来的偏差只以 $O(\sqrt{h})$ 的速度消失. (2) 镜像原理的名字来自偏微分方程: $V$ 满足 Black-Scholes 方程, 边界条件为 $V(t,H)=0$; 公式中的第二项恰是在 ``镜像点'' $H^2/S_0$ 处放置的一个反号的源, 使边界上的值相互抵消. 这与静电学中用镜像电荷求解边值问题的方法完全相同.
\end{remark}

% ---------------------------------------------------------------------
\subsection{美式期权与最优停时}

\textbf{美式期权}可以在到期日 $T$ 之前的任何时刻行权. 持有者选择行权时刻 $\tau$, 它必须是停时 (行权决定只能依赖于已有的信息), 行权时获得 $g(S_\tau)$, 例如看跌期权 $g(s)=(K-s)^+$. 可以证明 (见~\cite{shreve} 第 8 章), 美式期权的无套利价格是一个\textbf{最优停时问题}的值:
\begin{equation}\label{eq:american}
V_0=\sup_{\tau\le T}\widetilde{\E}\big[e^{-r\tau}g(S_\tau)\big],
\end{equation}
上确界取遍取值于 $[0,T]$ 的停时. 在离散时间中, 这个问题可以用逆向归纳求解: 在每个时刻比较 ``立即行权的收益'' 与 ``继续持有的价值'', 二者的较大者就是该时刻的价值 (第~3~章可选停时定理的一个应用, 称为 Snell 包络). 连续时间的情形更为微妙. 下面两个结果分别展示了鞅方法与偏微分方程方法.

\begin{proposition}[不付红利的美式看涨期权]\label{prop:american-call}
对 $g(s)=(s-K)^+$, $\sup_{\tau\le T}\widetilde{\E}[e^{-r\tau}(S_\tau-K)^+]=\widetilde{\E}[e^{-rT}(S_T-K)^+]$. 即提前行权永远不是最优的, 美式看涨期权与欧式看涨期权价格相同.
\end{proposition}

\begin{proof}
我们证明 $Y_t=e^{-rt}(S_t-K)^+=(\widetilde{S}_t-Ke^{-rt})^+$ 是 $\widetilde{\P}$-下鞅. 设 $t\le u\le T$. 由于 $\widetilde{S}$ 是鞅, 函数 $x\mapsto(x-Ke^{-rt})^+$ 是凸的, 由条件 Jensen 不等式~\ref{thm:cjensen},
\begin{align*}
Y_t&=\big(\widetilde{\E}[\widetilde{S}_u\,|\,\mathcal{F}_t]-Ke^{-rt}\big)^+\le\widetilde{\E}\big[(\widetilde{S}_u-Ke^{-rt})^+\,\big|\,\mathcal{F}_t\big]\\
&\le\widetilde{\E}\big[(\widetilde{S}_u-Ke^{-ru})^+\,\big|\,\mathcal{F}_t\big]=\widetilde{\E}[Y_u\,|\,\mathcal{F}_t],
\end{align*}
其中第二个不等号用到 $Ke^{-ru}\le Ke^{-rt}$ ($r\ge0$). 由可选停时定理 (对下鞅, 有界停时 $\tau\le T$; 与推论~\ref{cor:ct-ost-bounded} 的证明相同, 离散化后由定理~\ref{thm:ost-bounded} 得到), $\widetilde{\E}Y_\tau\le\widetilde{\E}Y_T$.
\end{proof}

直观上, 提前行权意味着提前支付执行价 $K$, 损失了 $K$ 的利息, 同时放弃了股价下跌时不行权的权利. 看跌期权的情形则不同, 股价很低时立即行权得到 $K-S_t$ 并获取利息可能更有利. 有限期限的美式看跌期权没有显式解, 但永久 ($T=\infty$) 的美式看跌期权可以完全求解.

\begin{theorem}[永久美式看跌期权]\label{thm:perpetual-put}
设 $r>0$, $g(s)=(K-s)^+$. 令 $\gamma=\frac{2r}{\sigma^2}$, $s^*=\frac{\gamma K}{1+\gamma}$,
\[
V^*(s)=\begin{cases}K-s,&0<s\le s^*,\\(K-s^*)\big(s/s^*\big)^{-\gamma},&s>s^*.\end{cases}
\]
则对一切 $s>0$,
\[
V^*(s)=\sup_\tau\widetilde{\E}_s\big[e^{-r\tau}(K-S_\tau)^+\1_{\{\tau<\infty\}}\big],
\]
上确界取遍一切停时, 且上确界在 $\tau^*=\inf\{t:S_t\le s^*\}$ 处达到. ($\widetilde{\E}_s$ 表示 $S_0=s$.)
\end{theorem}

\begin{proof}
\emph{第一步: 阈值策略.} 对 $0<b\le s$, 令 $\tau_b=\inf\{t:S_t\le b\}$. 由于 $\log\frac{S_t}{s}=\sigma\widehat{W}_t$, $\tau_b$ 是漂移为 $-\nu$ 的布朗运动 $-\widehat{W}$ 首次到达 $a=\frac{1}{\sigma}\log\frac{s}{b}$ 的时刻. 由定理~\ref{thm:drift-hitting} (取 $\mu=-\nu$, $\lambda=r$) 与 $\nu^2+2r=\frac{(r+\frac{1}{2}\sigma^2)^2}{\sigma^2}$,
\[
\widetilde{\E}_s\big[e^{-r\tau_b}\1_{\{\tau_b<\infty\}}\big]=\exp\Big(-a\Big(\nu+\frac{r+\frac{1}{2}\sigma^2}{\sigma}\Big)\Big)=\exp\Big(-\frac{2r}{\sigma}a\Big)=\Big(\frac{s}{b}\Big)^{-\gamma}.
\]
所以阈值策略 $\tau_b$ ($b<K$) 的价值为 $(K-b)(s/b)^{-\gamma}$. 对 $b$ 求最大值: $\frac{\dd}{\dd b}\big[(K-b)b^\gamma\big]=b^{\gamma-1}\big(\gamma(K-b)-b\big)$, 在 $b=s^*$ 处取得最大值. 于是对 $s>s^*$, 策略 $\tau^*$ 的价值为 $V^*(s)$; 对 $s\le s^*$, 立即行权 ($\tau^*=0$) 的价值为 $K-s=V^*(s)$. 所以上确界 $\ge V^*(s)$.

\emph{第二步: 验证 $V^*$ 是上界.} (a) $V^*$ 在 $(0,\infty)$ 上连续可微: 在 $s^*$ 处两段的值都是 $K-s^*$, 右导数为 $-\gamma\frac{K-s^*}{s^*}=-1$, 等于左导数 (\textbf{光滑粘贴}条件, 这正是 $s^*$ 的最优性条件). $V^*$ 在 $s^*$ 之外二阶连续可微, 且 $V^{*\prime\prime}$ 在 $s^*$ 附近有界. (b) $V^*\ge g$: 在 $(0,s^*]$ 上相等; 在 $(s^*,\infty)$ 上 $V^*>0$, 且 $V^*$ 是凸函数, 其在 $s^*$ 处的切线就是 $K-s$, 所以 $V^*\ge K-s$. (c) 记 $\widetilde{L}=rs\partial_s+\frac{1}{2}\sigma^2s^2\partial_{ss}$ 为 $\widetilde{\P}$ 下 $S$ 的生成元. 对 $s>s^*$, 由于 $\gamma$ 满足 $\frac{1}{2}\sigma^2\gamma(\gamma+1)-r\gamma-r=0$, $\widetilde{L}V^*-rV^*=0$; 对 $s<s^*$, $\widetilde{L}V^*-rV^*=-rs-r(K-s)=-rK<0$. 所以 $\widetilde{L}V^*-rV^*\le0$ (除 $s^*$ 一点外).

It\^o 公式对 $C^1$ 且除有限个点外 $C^2$、二阶导数局部有界的函数仍然成立 (例外点处二阶导数任意取值), 因为 $S$ 在任何一点停留的总时间为零; 证明与 Tanaka 公式相同, 用光滑函数逼近 (见~\cite{karatzas} 3.6 节). 于是
\[
e^{-rt}V^*(S_t)=V^*(s)+\int_0^te^{-ru}\big(\widetilde{L}V^*-rV^*\big)(S_u)\dd u+\int_0^te^{-ru}V^{*\prime}(S_u)\sigma S_u\dd\widetilde{B}_u.
\]
由于 $\abs{V^{*\prime}}\le1$ 且 $\widetilde{\E}\int_0^tS_u^2\dd u<\infty$, 随机积分是鞅. 对任意停时 $\tau$ 与 $t>0$, 由停止性质与 (c),
\begin{align*}
\widetilde{\E}_s\big[e^{-r(\tau\wedge t)}g(S_{\tau\wedge t})\big]&\le\widetilde{\E}_s\big[e^{-r(\tau\wedge t)}V^*(S_{\tau\wedge t})\big]\\
&=V^*(s)+\widetilde{\E}_s\int_0^{\tau\wedge t}e^{-ru}\big(\widetilde{L}V^*-rV^*\big)(S_u)\dd u\le V^*(s).
\end{align*}
令 $t\to\infty$, 由 Fatou 引理 ($g\ge0$), $\widetilde{\E}_s[e^{-r\tau}g(S_\tau)\1_{\{\tau<\infty\}}]\le V^*(s)$.
\end{proof}

\begin{figure}[htbp]
\centering
\begin{tikzpicture}
\begin{axis}[SAaxis, axis x line=bottom, axis y line=left, width=11cm, height=5.6cm, xmin=0, xmax=205, ymin=0, ymax=105,
  xtick={0,52.63,100,150,200}, xticklabels={$0$,$s^*$,$K$,$150$,$200$}, ytick={0,25,50,75,100}, xlabel={$s$},
  legend style={font=\small,draw=none,fill=none,at={(0.98,0.98)},anchor=north east}, legend cell align=left]
  \addplot[SADark,thick,dashed] coordinates {(0,100) (100,0) (200,0)};
  \addplot[SARed,thick,domain=52.63:200,samples=100,forget plot] {47.368*(x/52.63)^(-1.1111)};
  \addplot[SARed,thick] coordinates {(0,100) (52.63,47.37)};
  \fill[SARed] (axis cs:52.63,47.37) circle (2pt);
  \addplot[SADark!40,densely dotted,forget plot] coordinates {(52.63,0) (52.63,47.37)};
  \legend{行权收益 $(K-s)^+$, 期权价值 $V^*(s)$}
\end{axis}
\end{tikzpicture}
\caption{永久美式看跌期权 ($K=100$, $r=0.05$, $\sigma=0.3$, $s^*\approx52.6$). 在行权区域 $s\le s^*$ 价值等于收益; 在继续持有区域 $s>s^*$ 价值严格大于收益. 两段在 $s^*$ 处光滑粘贴 (函数值与导数都连续)}
\label{fig:perpetual-put}
\end{figure}

\begin{remark}[有限期限的美式看跌期权]
对 $T<\infty$, 价值函数 $v(t,s)$ 与一条随时间变化的\textbf{最优行权边界} $s^*(t)$ (当 $t\uparrow T$ 时上升到 $K$) 由下面的\textbf{自由边界问题}决定: 在继续持有区域 $\{s>s^*(t)\}$ 中 $v$ 满足 Black-Scholes 方程, 在行权区域中 $v=K-s$, 在边界上 $v$ 与 $\partial_sv$ 连续 (光滑粘贴). 等价地, $v$ 满足变分不等式
\[
\max\Big(\partial_tv+rs\partial_sv+\frac{1}{2}\sigma^2s^2\partial_{ss}v-rv,\ (K-s)^+-v\Big)=0.
\]
这个问题没有显式解, 实际中用二叉树、有限差分或回归型 Monte Carlo 方法 (Longstaff-Schwartz) 数值求解 (见~\cite{shreve} 第 8 章, \cite{glasserman} 第 8 章). 定理~\ref{thm:perpetual-put} 的证明结构, 即 ``先猜测阈值策略并优化阈值, 再验证上鞅性质'', 是求解最优停时问题的一般方法.
\end{remark}

% =====================================================================
\section{数值方法}
\label{sec:6-5}

大多数随机微分方程没有显式解 (有限期限的美式期权就是一个例子), 需要数值求解. 数值方法有两类目标: 逼近解的\emph{轨道} (强逼近, 例如模拟一个具体的系统、做滤波), 或逼近解的\emph{分布}的泛函 $\E f(X_T)$ (弱逼近, 例如 Monte Carlo 定价). 本节介绍最基本的 Euler-Maruyama 格式, 证明它的强收敛阶, 并讨论 Monte Carlo 方法.

% ---------------------------------------------------------------------
\subsection{Euler-Maruyama 格式}

设 $b,\sigma$ 不依赖于时间, 满足 \eqref{eq:H1}\eqref{eq:H2}. 取步长 $h=T/N$, 格点 $t_k=kh$. 模仿常微分方程的 Euler 方法, 在每一步把系数 ``冻结'' 在左端点:
\begin{equation}\label{eq:euler-maruyama}
\widehat{X}_{t_{k+1}}=\widehat{X}_{t_k}+b(\widehat{X}_{t_k})\,h+\sigma(\widehat{X}_{t_k})\,\Delta B_k,\qquad\Delta B_k=B_{t_{k+1}}-B_{t_k},\qquad\widehat{X}_0=\xi.
\end{equation}
模拟时, $\Delta B_k$ 是独立的 $N(0,h)$ 随机数. 注意系数必须取在左端点: 这正是 It\^o 积分的约定 (5.1 节), 取在中点的格式会收敛到 Stratonovich 方程的解.

\begin{theorem}[Euler-Maruyama 格式的强收敛]\label{thm:em}
设 $\xi\in L^2(\mathcal{F}_0)$, $X$ 是方程的解, $\widehat{X}$ 由 \eqref{eq:euler-maruyama} 定义. 则存在只依赖于 $K,T,\E\xi^2$ 的常数 $C$, 使得对一切 $h\le1$,
\[
\E\Big[\max_{0\le k\le N}\big\lvert X_{t_k}-\widehat{X}_{t_k}\big\rvert^2\Big]\le C\,h.
\]
\end{theorem}

\begin{proof}
记 $\eta(s)=t_k$ ($s\in[t_k,t_{k+1})$), 并定义连续时间的插值
\[
\widehat{X}_t=\xi+\int_0^tb(\widehat{X}_{\eta(s)})\dd s+\int_0^t\sigma(\widehat{X}_{\eta(s)})\dd B_s,\qquad t\in[0,T].
\]
在 $[t_k,t_{k+1}]$ 上被积过程是常数, 所以在格点上它与 \eqref{eq:euler-maruyama} 一致, 且在格点之间 $\widehat{X}_t-\widehat{X}_{\eta(t)}=b(\widehat{X}_{\eta(t)})(t-\eta(t))+\sigma(\widehat{X}_{\eta(t)})(B_t-B_{\eta(t)})$.

\emph{第一步: 矩估计.} 由 \eqref{eq:euler-maruyama} 与 \eqref{eq:H2} 对 $k$ 归纳可知 $\E\sup_t\abs{\widehat{X}_t}^2<\infty$. 记 $\widehat{m}(t)=\E\sup_{u\le t}\abs{\widehat{X}_u}^2$. 与定理~\ref{thm:sde-existence} 中矩估计的推导相同, $\widehat{m}(t)\le3\E\xi^2+6K^2(T+4)\int_0^t(1+\widehat{m}(s))\dd s$, 由 Gronwall 不等式, $\widehat{m}(T)\le M$, 其中 $M$ 与定理~\ref{thm:sde-existence} 中的上界相同, 不依赖于 $h$.

\emph{第二步: 局部误差.} 由于 $B_s-B_{\eta(s)}$ 与 $\mathcal{F}_{\eta(s)}$ 独立且方差不超过 $h$, 由 \eqref{eq:H2},
\[
\E\abs{\widehat{X}_s-\widehat{X}_{\eta(s)}}^2\le2h^2\E b(\widehat{X}_{\eta(s)})^2+2h\,\E\sigma(\widehat{X}_{\eta(s)})^2\le2(h^2+h)\cdot2K^2(1+M)\le C_1h.
\]

\emph{第三步: 全局误差.} $X=\Psi(\xi,X)$, 而 $\widehat{X}=\Psi(\xi,\widehat{X}_{\eta(\cdot)})$, 其中 $\widehat{X}_{\eta(\cdot)}$ 表示过程 $s\mapsto\widehat{X}_{\eta(s)}$ (它右连续且适应). 由基本估计引理~\ref{lem:sde-estimate}, $e(t)=\E\sup_{u\le t}\abs{X_u-\widehat{X}_u}^2$ 满足
\begin{align*}
e(t)&\le L\int_0^t\E\abs{X_s-\widehat{X}_{\eta(s)}}^2\dd s\le2L\int_0^t\Big(\E\abs{X_s-\widehat{X}_s}^2+\E\abs{\widehat{X}_s-\widehat{X}_{\eta(s)}}^2\Big)\dd s\\
&\le2L\int_0^te(s)\dd s+2LTC_1h.
\end{align*}
由 Gronwall 不等式, $e(T)\le2LTC_1e^{2LT}h$.
\end{proof}

所以 Euler-Maruyama 格式的\textbf{强收敛阶}为 $\frac{1}{2}$: 误差 $(\E\abs{X_T-\widehat{X}_T}^2)^{1/2}=O(h^{1/2})$. 这比常微分方程的 Euler 方法 (一阶) 差, 原因是布朗增量的量级为 $h^{1/2}$ 而非 $h$. 在 It\^o-Taylor 展开中, 被略去的最低阶项是 $\frac{1}{2}\sigma\sigma'\big((\Delta B)^2-h\big)$, 其量级为 $O(h)$, 高于 $O(h^2)$. 把这一项加入格式, 就得到 \textbf{Milstein 格式}
\[
\widehat{X}_{t_{k+1}}=\widehat{X}_{t_k}+b\,h+\sigma\,\Delta B_k+\frac{1}{2}\sigma\sigma'\big((\Delta B_k)^2-h\big)\qquad(\text{系数都在 }\widehat{X}_{t_k}\text{ 处取值}),
\]
它的强收敛阶为 $1$ (见~\cite{kloeden} 第 10 章). 若 $\sigma$ 是常数 (加性噪声), 二者相同, Euler-Maruyama 格式本身就有一阶强收敛. 图~\ref{fig:em-error} 用几何布朗运动 (有显式解, 可以精确计算误差) 验证了这两个收敛阶.

\begin{figure}[htbp]
\centering
\begin{tikzpicture}
\begin{axis}[SAaxis, axis x line=bottom, axis y line=left, width=6.6cm, height=5.6cm, xmode=log, ymode=log, xmin=0.003, xmax=0.65, ymin=0.0015, ymax=0.6,
  xlabel={步长 $h$}, ylabel={$\E\abs{X_1-\widehat{X}_1}$}, log basis x=2,
  legend style={font=\scriptsize,draw=none,fill=none,at={(0.5,-0.24)},anchor=north}, legend columns=2, legend cell align=left]
  \addplot[SABlue,thick,mark=*,mark size=1.4pt] coordinates {(0.5,0.35689) (0.25,0.27216) (0.125,0.20235) (0.0625,0.14587) (0.03125,0.10395) (0.015625,0.07426) (0.0078125,0.05224) (0.00390625,0.03759)};
  \addplot[SARed,thick,mark=square*,mark size=1.4pt] coordinates {(0.5,0.21145) (0.25,0.12737) (0.125,0.07035) (0.0625,0.03726) (0.03125,0.01925) (0.015625,0.00973) (0.0078125,0.00486) (0.00390625,0.00244)};
  \addplot[SABlue!50,dashed,domain=0.0039:0.5] {0.25*x^0.5};
  \addplot[SARed!50,dashed,domain=0.0039:0.5] {0.2*x};
  \legend{Euler-Maruyama, Milstein, 斜率 $\frac12$, 斜率 $1$}
\end{axis}
\end{tikzpicture}\hfill
\begin{tikzpicture}
\begin{axis}[SAaxis, axis x line=bottom, axis y line=left, width=8cm, height=5.6cm, xmin=0, xmax=1.03, ymin=0.7, ymax=2.6,
  xtick={0,0.25,0.5,0.75,1}, ytick={1,1.5,2,2.5}, xlabel={$t$},
  legend style={font=\scriptsize,draw=none,fill=none,at={(0.5,-0.24)},anchor=north}, legend columns=2, legend cell align=left]
  \addplot[SADark,thin] coordinates {(0.000,1.000) (0.004,1.039) (0.008,0.986) (0.012,0.998) (0.016,1.039) (0.020,0.988) (0.023,0.963) (0.027,0.934) (0.031,0.901) (0.035,0.917) (0.039,0.895) (0.043,0.902) (0.047,0.864) (0.051,0.902) (0.055,0.891) (0.059,0.889) (0.062,0.952) (0.066,0.965) (0.070,0.963) (0.074,0.999) (0.078,0.889) (0.082,0.939) (0.086,0.922) (0.090,0.815) (0.094,0.830) (0.098,0.869) (0.102,0.906) (0.105,0.933) (0.109,0.938) (0.113,0.975) (0.117,0.927) (0.121,1.020) (0.125,1.127) (0.129,1.213) (0.133,1.229) (0.137,1.158) (0.141,1.190) (0.145,1.159) (0.148,1.178) (0.152,1.186) (0.156,1.174) (0.160,1.175) (0.164,1.139) (0.168,1.150) (0.172,1.254) (0.176,1.312) (0.180,1.341) (0.184,1.327) (0.188,1.291) (0.191,1.223) (0.195,1.316) (0.199,1.312) (0.203,1.282) (0.207,1.226) (0.211,1.284) (0.215,1.328) (0.219,1.484) (0.223,1.412) (0.227,1.382) (0.230,1.397) (0.234,1.348) (0.238,1.440) (0.242,1.489) (0.246,1.554) (0.250,1.657) (0.254,1.797) (0.258,1.813) (0.262,1.900) (0.266,1.844) (0.270,1.863) (0.273,1.900) (0.277,1.997) (0.281,2.010) (0.285,2.012) (0.289,2.297) (0.293,2.265) (0.297,2.203) (0.301,2.283) (0.305,2.162) (0.309,2.101) (0.312,2.001) (0.316,2.039) (0.320,2.039) (0.324,2.108) (0.328,2.022) (0.332,2.019) (0.336,2.097) (0.340,2.229) (0.344,2.386) (0.348,2.417) (0.352,2.329) (0.355,2.283) (0.359,2.295) (0.363,2.150) (0.367,2.019) (0.371,1.879) (0.375,1.845) (0.379,1.715) (0.383,1.545) (0.387,1.539) (0.391,1.509) (0.395,1.457) (0.398,1.426) (0.402,1.506) (0.406,1.473) (0.410,1.503) (0.414,1.434) (0.418,1.392) (0.422,1.384) (0.426,1.495) (0.430,1.440) (0.434,1.324) (0.438,1.241) (0.441,1.154) (0.445,1.027) (0.449,0.994) (0.453,1.019) (0.457,1.064) (0.461,1.059) (0.465,1.151) (0.469,1.243) (0.473,1.158) (0.477,1.213) (0.480,1.189) (0.484,1.127) (0.488,1.144) (0.492,1.092) (0.496,1.188) (0.500,1.215) (0.504,1.191) (0.508,1.357) (0.512,1.373) (0.516,1.374) (0.520,1.321) (0.523,1.334) (0.527,1.274) (0.531,1.211) (0.535,1.179) (0.539,1.190) (0.543,1.188) (0.547,1.198) (0.551,1.204) (0.555,1.191) (0.559,1.201) (0.562,1.276) (0.566,1.274) (0.570,1.219) (0.574,1.339) (0.578,1.351) (0.582,1.300) (0.586,1.316) (0.590,1.405) (0.594,1.434) (0.598,1.415) (0.602,1.402) (0.605,1.348) (0.609,1.358) (0.613,1.430) (0.617,1.443) (0.621,1.391) (0.625,1.315) (0.629,1.270) (0.633,1.189) (0.637,1.229) (0.641,1.171) (0.645,1.259) (0.648,1.285) (0.652,1.331) (0.656,1.203) (0.660,1.175) (0.664,1.102) (0.668,1.192) (0.672,1.219) (0.676,1.197) (0.680,1.241) (0.684,1.217) (0.688,1.090) (0.691,1.159) (0.695,1.266) (0.699,1.304) (0.703,1.290) (0.707,1.221) (0.711,1.204) (0.715,1.249) (0.719,1.274) (0.723,1.230) (0.727,1.235) (0.730,1.237) (0.734,1.146) (0.738,1.173) (0.742,1.174) (0.746,1.212) (0.750,1.157) (0.754,1.139) (0.758,1.194) (0.762,1.226) (0.766,1.318) (0.770,1.257) (0.773,1.242) (0.777,1.179) (0.781,1.150) (0.785,1.122) (0.789,1.115) (0.793,1.082) (0.797,1.129) (0.801,1.079) (0.805,1.159) (0.809,1.153) (0.812,1.179) (0.816,1.168) (0.820,1.150) (0.824,1.164) (0.828,1.255) (0.832,1.266) (0.836,1.313) (0.840,1.271) (0.844,1.233) (0.848,1.239) (0.852,1.235) (0.855,1.109) (0.859,1.123) (0.863,1.074) (0.867,1.082) (0.871,1.000) (0.875,0.925) (0.879,0.876) (0.883,0.965) (0.887,0.889) (0.891,0.871) (0.895,0.948) (0.898,0.905) (0.902,0.954) (0.906,0.926) (0.910,0.930) (0.914,0.902) (0.918,0.948) (0.922,1.011) (0.926,0.932) (0.930,0.908) (0.934,0.958) (0.938,0.938) (0.941,0.999) (0.945,0.957) (0.949,0.973) (0.953,0.938) (0.957,0.928) (0.961,0.993) (0.965,1.062) (0.969,0.984) (0.973,0.989) (0.977,0.933) (0.980,0.897) (0.984,0.904) (0.988,0.993) (0.992,0.963) (0.996,0.986) (1.000,0.963)};
  \addplot[SARed,thick,mark=*,mark size=1.3pt] coordinates {(0.000,1.000) (0.125,1.160) (0.250,1.653) (0.375,1.897) (0.500,1.181) (0.625,1.321) (0.750,1.204) (0.875,0.983) (1.000,1.062)};
  \legend{精确解, Euler-Maruyama ($h=\frac18$)}
\end{axis}
\end{tikzpicture}
\caption{左: 几何布朗运动 $\dd X=\frac12X\dd t+0.8X\dd B$ 在 $T=1$ 处的强误差 (双对数坐标, $2\times10^4$ 条轨道), Euler-Maruyama 格式与 Milstein 格式的误差分别以 $h^{1/2}$ 与 $h$ 的速度下降. 右: 同一条布朗轨道驱动下的精确解与粗步长的 Euler-Maruyama 近似}
\label{fig:em-error}
\end{figure}

% ---------------------------------------------------------------------
\subsection{Monte Carlo 方法}

为计算 $\E f(X_T)$ (例如 Feynman-Kac 表示中的期望, 或期权价格), 用 Euler-Maruyama 格式独立地模拟 $n$ 条轨道, 取样本平均 $\frac{1}{n}\sum_{j=1}^nf(\widehat{X}^{(j)}_T)$. 其误差由两部分组成:
\[
\frac{1}{n}\sum_{j=1}^nf\big(\widehat{X}^{(j)}_T\big)-\E f(X_T)=\underbrace{\Big(\frac{1}{n}\sum_{j=1}^nf\big(\widehat{X}^{(j)}_T\big)-\E f\big(\widehat{X}_T\big)\Big)}_{\text{统计误差}}+\underbrace{\Big(\E f\big(\widehat{X}_T\big)-\E f(X_T)\Big)}_{\text{离散化偏差 (弱误差)}}.
\]
由中心极限定理, 统计误差的量级为 $\frac{\sigma_f}{\sqrt{n}}$ (第~1~章), 可以用样本标准差估计并给出置信区间. 对光滑的 $f$, 弱误差是 $O(h)$, 比强误差好 (\textbf{弱收敛阶}为 $1$, 见~\cite{kloeden} 第 14 章): 它只要求分布接近, 无需轨道接近. 要使总误差为 $\varepsilon$, 需要 $h\sim\varepsilon$, $n\sim\varepsilon^{-2}$, 计算量为 $n/h\sim\varepsilon^{-3}$. Giles 的\textbf{多层 Monte Carlo} 方法把不同步长的模拟组合起来, 可以把计算量降到 $\varepsilon^{-2}(\log\varepsilon)^2$ (见~\cite{glasserman}).

对依赖于整条轨道的泛函, 离散化偏差可能更大. 以 6.4.4 小节的向上敲出看涨期权为例: 若只在格点上检查股价是否超过障碍, 就会漏掉格点之间的越界, 从而高估期权价格. 图~\ref{fig:barrier-mc} 显示, 这种偏差只以 $O(\sqrt{h})$ 的速度消失: 在步长为 $h$ 的两个格点之间, 布朗轨道的振幅是 $\sqrt{h}$ 的量级.

\begin{figure}[htbp]
\centering
\begin{tikzpicture}
\begin{axis}[SAaxis, axis x line=bottom, axis y line=left, width=11cm, height=5.4cm, xmode=log, log basis x=2, xmin=3, xmax=700, ymin=1.3, ymax=3.2,
  xtick={4,16,64,256}, xticklabels={$4$,$16$,$64$,$256$}, ytick={1.5,2,2.5,3}, xlabel={$N$},
  legend style={font=\small,draw=none,fill=none,at={(0.98,0.98)},anchor=north east}, legend cell align=left]
  \addplot[SABlue,thick,mark=*,mark size=1.5pt,error bars/.cd,y dir=both,y explicit] coordinates {(4,3.0184) +- (0,0.0203) (8,2.6183) +- (0,0.0187) (16,2.2899) +- (0,0.0174) (32,2.0749) +- (0,0.0164) (64,1.9025) +- (0,0.0156) (128,1.7839) +- (0,0.0150) (256,1.6978) +- (0,0.0146) (512,1.6394) +- (0,0.0143)};
  \addplot[SARed,thick,dashed] coordinates {(3,1.503) (700,1.503)};
  \legend{离散监测的 Monte Carlo 估计, 连续监测的精确价格}
\end{axis}
\end{tikzpicture}
\caption{向上敲出看涨期权 ($S_0=K=100$, $H=130$, 其余参数同图~\ref{fig:barrier}) 的 Monte Carlo 价格 (横轴为每年的监测次数 $N=1/h$, $4\times10^5$ 条轨道, 误差棒为 $95\%$ 置信区间). 只在 $N$ 个等距时刻检查障碍时, 价格随 $N$ 增大而缓慢地下降到连续监测的价格 (定理~\ref{thm:barrier}), 偏差约与 $N^{-1/2}$ 成正比}
\label{fig:barrier-mc}
\end{figure}

\begin{remark}[布朗桥修正]
反射原理给出了消除这一偏差的办法. 由推论~\ref{cor:max-joint}, 对从 $x$ 出发、在时刻 $\Delta$ 到达 $y$ 的布朗运动 ($x,y<h$), 中间曾经到达 $h$ 的条件概率为
\[
\P\Big(\max_{u\le\Delta}W_u\ge h\,\Big|\,W_0=x,W_\Delta=y\Big)=\frac{\varphi_\Delta(2h-x-y)}{\varphi_\Delta(y-x)}=\exp\Big(-\frac{2(h-x)(h-y)}{\Delta}\Big)
\]
(先对 $x=0$ 计算: 联合密度与 $W_\Delta$ 的密度之比为 $\exp\big(-\frac{(2h-y)^2-y^2}{2\Delta}\big)=\exp\big(-\frac{2h(h-y)}{\Delta}\big)$, 一般情形平移即可). 漂移不改变这个条件概率 (给定两端点后, 漂移项被吸收), 扩散系数为 $\sigma$ 时把 $\Delta$ 换成 $\sigma^2\Delta$. 模拟时, 在每一步以这个概率判定 ``格点之间越界'', 即可得到连续监测的无偏估计. 用这种方法, 即使每年只监测 $50$ 次, $2\times10^6$ 条轨道给出的估计 $1.5038\pm0.0060$ 也与精确值 $1.503$ 吻合.
\end{remark}

% ---------------------------------------------------------------------
\subsection*{本章小结}
\addcontentsline{toc}{subsection}{本章小结}

\begin{center}
\begin{minipage}{\textwidth}
\centering
\captionof{table}{随机微分方程与扩散过程的主要结论}
\label{tab:ch6-summary}
\renewcommand{\arraystretch}{1.5}
\begin{tabular}{>{\raggedright\arraybackslash}p{0.24\textwidth} >{\raggedright\arraybackslash}p{0.46\textwidth} >{\raggedright\arraybackslash}p{0.20\textwidth}}
\toprule
结论 & 内容 & 用到的工具 \\
\midrule
线性方程 & 常数变易法给出显式解 (GBM, OU, 布朗桥) & It\^o 公式, 乘积法则 \\
存在唯一性 & Lipschitz + 线性增长 $\Rightarrow$ 唯一强解 & Picard 迭代, Doob, Gronwall \\
弱解 & Girsanov 构造; Tanaka 方程无强解 & Girsanov, L\'evy 刻画, Tanaka \\
马尔可夫性 & $\E[f(X_{s+h})\,|\,\mathcal{F}_s]=P_hf(X_s)$ & 唯一性, 代入引理, 单调类 \\
生成元 & $L=b\partial_x+\frac12\sigma^2\partial_{xx}$, Dynkin 公式 & It\^o 公式, 可选停时 \\
Kolmogorov 方程 & $\partial_tu=Lu$, $\partial_tp=L^*p$, 平稳分布 & It\^o 公式, 分部积分 \\
出口问题 & 尺度函数 $Ls=0$, 期望出口时间 $Lv=-1$ & Dynkin 公式 \\
Dirichlet 问题 & $u(\rvec{x})=\E_{\rvec{x}}g(\rvec{W}_{\tau_D})$ & 多维 It\^o 公式 \\
Feynman-Kac & 带贴现与源项的抛物方程的概率表示 & It\^o 公式, 矩估计 \\
期权定价 & Black-Scholes 公式, 障碍期权, 美式期权 & Girsanov, 鞅表示, 反射原理 \\
Euler-Maruyama & 强收敛阶 $\frac12$, 弱收敛阶 $1$ & 基本估计, Gronwall \\
\bottomrule
\end{tabular}
\end{minipage}
\end{center}

\subsection*{结语}
\addcontentsline{toc}{subsection}{结语}

回顾全书, 第~1~章的测度论提供了严格的语言, 第~2~章用条件期望刻画信息; 第~3~章的鞅论描述了 ``公平'' 的随机演化并提供了最有力的不等式与收敛定理, 第~4~章的布朗运动则是连续随机性的普适模型. 在此基础上, It\^o 积分与 It\^o 公式建立了随机微积分 (第~5~章), 随机微分方程把它们用于描述和计算一般的连续随机系统 (第~6~章). 进一步的学习可以沿几个方向展开: 一般半鞅与带跳过程的随机积分 (\cite{revuz}), 随机控制与最优停时, 滤波理论, Malliavin 分析, 随机偏微分方程, 以及更深入的金融数学 (\cite{shreve}). 扩散过程与偏微分方程之间的联系 (\cite{karatzas}, \cite{oksendal}) 则贯穿了所有这些方向.


\backmatter
% =====================================================================
%  参考文献
% =====================================================================
\begin{thebibliography}{99}
\addcontentsline{toc}{chapter}{参考文献}
\bibitem{durrett} R. Durrett. \emph{Probability: Theory and Examples}, 5th ed. Cambridge University Press, 2019.
\bibitem{billingsley} P. Billingsley. \emph{Probability and Measure}, anniversary ed. Wiley, 2012.
\bibitem{williams} D. Williams. \emph{Probability with Martingales}. Cambridge University Press, 1991.
\bibitem{kallenberg} O. Kallenberg. \emph{Foundations of Modern Probability}, 3rd ed. Springer, 2021.
\bibitem{folland} G. B. Folland. \emph{Real Analysis: Modern Techniques and Their Applications}, 2nd ed. Wiley, 1999.
\bibitem{karatzas} I. Karatzas, S. E. Shreve. \emph{Brownian Motion and Stochastic Calculus}, 2nd ed. Springer, 1991.
\bibitem{morters} P. M\"orters, Y. Peres. \emph{Brownian Motion}. Cambridge University Press, 2010.
\bibitem{revuz} D. Revuz, M. Yor. \emph{Continuous Martingales and Brownian Motion}, 3rd ed. Springer, 1999.
\bibitem{oksendal} B. \O ksendal. \emph{Stochastic Differential Equations: An Introduction with Applications}, 6th ed. Springer, 2003.
\bibitem{shreve} S. E. Shreve. \emph{Stochastic Calculus for Finance II: Continuous-Time Models}. Springer, 2004.
\bibitem{kloeden} P. E. Kloeden, E. Platen. \emph{Numerical Solution of Stochastic Differential Equations}. Springer, 1992.
\bibitem{glasserman} P. Glasserman. \emph{Monte Carlo Methods in Financial Engineering}. Springer, 2003.
\end{thebibliography}


\end{document}
